\documentclass[11pt,a4paper]{article}
\usepackage[margin=27mm,top=24mm,bottom=26mm]{geometry}
\usepackage{amsmath,amssymb,amsthm,mathtools,booktabs,array,xcolor}
\usepackage[unicode,colorlinks=true,linkcolor=black,citecolor=blue!45!black,urlcolor=blue!48!black]{hyperref}
\usepackage{microtype}
\definecolor{sourceblue}{RGB}{32,75,105}
\hypersetup{pdftitle={From theta lifts to the critical line},pdfauthor={Codex, an independent exposition of work by Levent Alpöge with Claude},pdfsubject={A source-based reconstruction of the Lean proof of RH for Dirichlet L-functions}}
\urlstyle{same}
\newcommand{\src}[3]{\href{https://github.com/leanprover-community/riemann-hypothesis/blob/40e56bab026552e089d7afcfeff08d2defb9c005/\detokenize{#1}\#L#2}{\textcolor{sourceblue}{#3}}}
\newcommand{\code}[1]{\texttt{#1}}
\newcommand{\Sha}{\operatorname{Sha}}
\newtheorem{theorem}{Theorem}[section]
\newtheorem{lemma}[theorem]{Lemma}
\newtheorem{proposition}[theorem]{Proposition}
\newtheorem{corollary}[theorem]{Corollary}
\theoremstyle{definition}
\newtheorem{definition}[theorem]{Definition}
\newtheorem{formalinput}[theorem]{Formal input}
\theoremstyle{remark}
\newtheorem{remark}[theorem]{Remark}
\numberwithin{equation}{section}
\setcounter{tocdepth}{1}
\setlength{\emergencystretch}{2em}
\title{\textbf{From theta lifts to the critical line}\\[4pt]
\large Arithmetic reconstruction of the Alpöge--Claude argument for Dirichlet \(L\)-functions}
\author{Codex\\[-1pt]\small Independent exposition of the November 2026 formalization}
\date{Revised 22 November 2026}
\begin{document}
\maketitle
\begin{abstract}
Alpöge with Claude deduce the critical-line theorem for Dirichlet
\(L\)-functions from one-sided bounds for weighted prime-power sums
on arbitrarily late intervals \([Y,Y^{3/2}]\). We give a paper proof
of this deduction and compute the signed average over translates
as four explicit twists of each retained character. A direct sum indexed by translates gives a single slope comparison with explicit costs for lattice denominators and exceptional target slopes. In the
arithmetic construction we attach Galois representations to the
standard definite-unitary systems and construct integral cycles
detecting the twisted norms of divisors. We give classical cuspidal
attachments on specified ternary and quartic branches. For each
fixed target we bound the entire dyadic field contribution on a
jointly selected interval and the opposite-parity contribution at
all odd primes. The remaining Selmer budget reduces to a sum of
matching-parity character lengths. We bound three ranges in terms
of order and twist and all of Case IV. For the zeta target we compute
finite quadratic layers and exhibit a positive Case II length. We prove both recursions for
the theta periods at the allowed signatures, and use the doubled
theta kernel to span the bounded forms on the marked period charts.
For the final weight we compute the two integral Lie determinants
at a product CM point and compare their finite orders on the
indefinite marked components. This gives bounded spanning sections
on each good cover satisfying the stage bounds. Over the field
constructed for the deep level we realize the principal marked components
and their level maps, and identify the unmodified lattices at good primes in central weight zero.
For the signed mean we give a test using only primes up to the window
endpoint and verify one complete window for a quartic character.
We compute the finite corrections to the relative determinant and
compare the recorded and mixed component volumes. A rank test on one
target quotient suffices. We bound the slope cost of a missed line.
The local rank input bounds both prime losses. For any fixed
\(r\ge1\), one signed relative determinant at one translate for
all but a power-saving fraction of the integer \(r\)-th powers in
arbitrarily late fixed-ratio intervals suffices for both estimates
at the original exponent. For every primitive target of conductor
at most \(400\,000\), both estimates hold on a subinterval of each
sufficiently late such interval inside a supplied window. The
full cofinal signed estimate remains open.
\end{abstract}
\tableofcontents

\section{The theorem and its arithmetic content}\label{intro:main}
In the proof announced by Alpöge with Claude \cite{announce}, the length of a chain of theta lifts is chosen at approximately square-root scale. This becomes a statement about primes only after comparing integral lattices, passing to finite Fourier coefficients and estimating the terms removed by averaging. The final analytic step must work with a signed sum on disjoint windows. It has no assumption on the gap between consecutive windows.

\begin{theorem}[Alpöge with Claude, in the published formalization]\label{main:thm}
Let \(N\ge1\) and let \(\chi\) be a Dirichlet character modulo \(N\). If \(\Re s>1/2\) and \((\chi\ne1\lor s\ne1)\), then \(L(s,\chi)\ne0\). Every zero with \(0<\Re s<1\) has real part \(1/2\). Every zero of \(\zeta\) other than \(-2,-4,\ldots\), with \(s\ne1\), has real part \(1/2\).
\end{theorem}
The exclusion at \(s=1\) accommodates the total-function convention of the formal statement at a principal pole. The five declarations in the submission give Mathlib's RH, a second spelling of the same zeta statement, the half-plane for every character, the strip for every character, and the strip for primitive characters. The final proof in Section~\ref{main:end} treats their classical content. \src{src/Main.lean}{4}{Five exported declarations}, \src{src/Statements.lean}{3}{submission statements} and \cite{mathlib}.

Write \(\Lambda(n)\) for the von Mangoldt function, \(t_+=\max(t,0)\), and let \(1\) in the following family always denote the character of modulus one. For \(y\ge1\) put \(P_\chi(y)=\sum_{1\le n\le y}\Lambda(n)\chi(n)(y/n-1)_+\). Our centred functions and their signed real combination are
\begin{equation}\label{eq:centred}
\begin{aligned}
 G_1(y)&=P_1(y)-y\log y+(1+\gamma)y,\\
 G_\chi(y)&=P_\chi(y)+y\frac{L'(1,\chi)}{L(1,\chi)}
                     &&(\chi\text{ nonprincipal}),\\
 G_{\mathcal F}(y)&=\sum_{i\in J}a_i\,\Re G_{\chi_i}(y),
                     && |J|<\infty,\quad a_i>0 .
\end{aligned}
\end{equation}
Here \(\gamma\) is Euler's constant. A family is admissible if every principal character in it has modulus one. For a primitive \(\psi\), let \(o\) be the order of the cyclic group generated by \(\psi\) among characters of its modulus. Take the primitive inducer of each power. The resulting finite set contains \(\psi\) and the trivial character. Its trivial weight is one and each other weight is \((3o+1)^{-1}\). Denote this family by \(\mathcal F_\psi\). Raising characters at a fixed modulus then taking primitive inducers produces no duplicates among distinct characters, since each inducer has a unique lift to that modulus. This identifies the precise family in the subsequent analytic transform. \src{src/Development/AutomorphicArithmetic/DedekindLFunctions/CharFamily.lean}{15}{centred functions}, \src{src/Formalism/SharpRestriction/RamificationFrobenius/TotientGaloisClasses/ExplicitFormulaErrorOfHostConditions.lean}{36}{weights}, \src{src/Formalism/SharpRestriction/RamificationFrobenius/TotientGaloisClasses/ExplicitFormulaErrorOfHostConditions.lean}{105}{cyclic level} and \src{src/Development/AutomorphicArithmetic/PrimitiveCosets/PrimInj.lean}{6}{injectivity at fixed host}.

\begin{theorem}[Arithmetic window extracted from the formalization]\label{main:window}
For each primitive \(\psi\) there is \(\varepsilon_0>0\) such that, for each \(0<\varepsilon<\varepsilon_0\) and each \(Y_0\), there exists \(Y\ge\max(Y_0,1)\) with
\begin{equation}\label{eq:window}
  G_{\mathcal F_\psi}(x)\le x^{1/2+9\varepsilon}
  \qquad\text{for all }Y\le x\le Y^{3/2}.
\end{equation}
\end{theorem}
The source for this edition is the 20 November release at revision
\texttt{40e56bab0265} of the community mirror~\cite{source}.
The arithmetic result cited for Theorem~\ref{main:window} is stated
as Formal input~\ref{red:design}. At fixed \(\psi,\varepsilon\) it
chooses one family with arbitrarily late complete windows. The
cohomological budget and the character bounds hold on the same hosts.
At each good \(x\) the geometric comparison supplies a single set of
scales, ranks and weights for all its translates. Substitution of
its two geometric budgets gives \eqref{scope:missing}. Section~\ref{main:assembly}
gives the selection order and the three inequalities used in that
substitution. When \(\psi=1\) the left side of \eqref{eq:window} is a
single real function, although the construction still uses three
quadratic twists. Section~\ref{an:section} shows why a zero to the right
would force a positive excursion in every sufficiently late window.

\begin{theorem}[Conditional conclusion of this paper]\label{main:conditional}
Suppose that, for every primitive \(\psi\) and every sufficiently small
\(\varepsilon>0\), there are arbitrarily large \(Y\) with the following
property. At every \(x\in[Y,Y^{3/2}]\) there are data satisfying the
finite-character, rank, scale and excluded-character hypotheses stated
with \eqref{scope:missing}, with the indicated constants uniform in
\(x\) and \(Y\), and \eqref{scope:missing} holds for every
translate of those data. Then all three assertions of
Theorem~\ref{main:thm} hold.
\end{theorem}
\begin{proof}
The finite Fourier identity and its four error estimates in
Section~\ref{ar:section} give \eqref{eq:window} at each such \(x\).
Theorem~\ref{an:zero} applies to these intervals at every exponent
strictly between \(1/2\) and \(1\), by choosing \(\varepsilon\) small
enough. Section~\ref{main:end} applies this to each primitive character,
passes to arbitrary conductors and reflects the resulting zero-free
half-plane across the critical line.
\end{proof}

Theorem~\ref{main:conditional} also follows from the weaker signed
average \eqref{scope:average}. Sections~\ref{geo:translate} and~\ref{red:translate} give a single slope comparison for this average and state its arithmetic costs. Proposition~\ref{ar:mean} computes
that average and Corollary~\ref{ar:meanhost} bounds its three
auxiliary twists. Section~\ref{an:section} then excludes a zero
by its contribution to the signed Mellin transform. On the
arithmetic side Proposition~\ref{famgal:rank} extracts the seed
rank. Proposition~\ref{addg:unitary} supplies the classical attachment
on the standard definite unitary group. Theorem~\ref{addc:cycle} constructs
the homology classes for the twisted-norm count.
Corollary~\ref{hr:imprimitive} treats the imprimitive
quartic branch, and Lemma~\ref{hr:towerfin} identifies the finite
Hecke quotient. Section~\ref{famsel:complement} sums the Selmer
estimates over the required range. Corollary~\ref{addo:remainder}
gives a single residual character sum after the dyadic field term,
the opposite parity and the proved character ranges have been bounded.
Theorem~\ref{tcr:global} proves both theta recursions.
Propositions~\ref{dg:identity} and~\ref{munf:span} give the doubled
kernel and its spanning consequence. Proposition~\ref{cmeig:degree}
computes the integral CM degree used in the final weights.
Section~\ref{scope:account} records the scope of each application
to the family. Source links refer to the fixed release~\cite{source}.

\section{Galois representations from simultaneous Hecke congruences}\label{tr:section}

The supply of arithmetic windows uses a local--global compatibility theorem for Galois representations attached to automorphic data. We isolate the passage from congruences of Hecke eigenvalues to a continuous representation. It illustrates how a conditional existence theorem is discharged and why approximating individual coordinates would be insufficient. The geometrically constructed seed representations, the automorphic congruence theorem for the requisite weights and types, and the crystalline closure result used at places above \(\ell\) are cited inputs.

\subsection{Extension of trace data}
Let \(k=\overline{\mathbb Q}_\ell\) with its \(\ell\)-adic norm, let \(\Gamma\) be the profinite Galois group in the construction and \(x:J\to\Gamma\) its Frobenius coordinates. The target values \(\xi_j\) lie in one finite extension \(L/\mathbb Q_\ell\) inside \(k\). The seed representations \(r_i:\Gamma\to\mathrm{GL}_a(k)\) are continuous. Their systems \(\eta_i:J\to k\) need agree with \(t_i(x_j)=\operatorname{tr}r_i(x_j)\) only on sets \(D_i\subseteq J\). For each finite set \(S\) of indices, \(x(\bigcap_{i\in S}D_i)\) is dense. In the application Chebotarev density after exclusion of finitely many primes is combined with the intersection of the finitely many additional dense open conditions attached to the seeds. The congruence theorem supplies
\begin{equation}\label{tr:approx}
 \forall e>0\quad \exists S\ {\rm finite},\ (c_i)_{i\in S}\subset k\quad
 \forall P\in\mathbb Z[X_j:j\in J],\qquad
 \left|P(\xi)-\sum_{i\in S}c_iP(\eta_i)\right|_\ell\le e.
\end{equation}
In particular the coefficients \(c_i\) are independent of \(P\). The source also supplies this property for specified subfamilies carrying one prescribed finite set of places. \src{src/Development/AutomorphicArithmetic/Foundations/GaloisRepresentations/TraceLimits/LimFinOddOf.lean}{69}{families and joint coordinates} and \src{src/Development/AutomorphicArithmetic/Foundations/GaloisRepresentations/TraceLimits/LimFinOddOf.lean}{91}{density for finite intersections}.

\paragraph{The common good index in the arithmetic application.}
The indices can be taken to be quadruples \((w,A,s,n)\), with \(w\nmid\ell\)
split and spherical for the target, \(A\) a prolongation of \(w\) to the
algebraic closure, \(s\) an arithmetic Frobenius lift at \(A\), and
\(n\ge1\). Their group coordinate is \(s^{-n}\). In particular all
prolongations and all lifts are allowed at exponent one.
\src{src/Development/AbelianVarietyTheory/Proofs/SatakeParameters/FrobOpsIdx.lean}{24}{The index set}
and \src{src/Development/AbelianVarietyTheory/Proofs/SatakeParameters/FrobOpsIdx.lean}{71}{the exponent-one witness}.

Here is the density argument. Fix a nonempty open set \(W\subset G_E\) and
a finite family \(I_0\) of seeds. For each such seed the attachment supplies
a dense open \(V_i\subset G_E\) on which to test arithmetic Frobenius
lifts and a finite set \(X_i\) of
exceptional places. The open set \(W^{-1}\cap\bigcap_i V_i\) is nonempty
and contains a coset of an open normal subgroup. In the corresponding
finite Galois quotient, Chebotarev gives primes of \(E\) with the conjugacy
class of that coset. We can take one unramified of absolute residue degree
one, away from \(\ell\), the bad places of the target and those of all
these seeds, also avoiding the primes below these finite sets in \(E^+\).
Indeed a Chebotarev class has positive density among primes of \(E\),
whereas primes of absolute degree at least two have density zero (their
norms up to \(X\) lie over rational primes up to \(\sqrt X\)).
Omit also the primes dividing \(2d\) and the rational primes
ramified in \(E\). Then \(E_w=\mathbb Q_p\). Its local degree
over \(E^+\) is one, so \(w\) splits over \(E^+\). The given
embedding of \(\mathbb Q(\sqrt{-d})\) into \(E\) makes \(-d\)
a square in \(\mathbb Q_p\), as the usable-place condition requires.
Conjugating a prime above \(w\) in the finite Galois extension selects
the desired element of the class. Prolong it to \(A\) and lift Frobenius
there. The resulting \(s\) belongs to the prescribed coset. At this
single \((w,A,s)\) each seed has its normalized characteristic polynomial
and the level supplies a spherical local vector. The exponent-one Hecke
operator therefore reads \(\eta_i(j)=\operatorname{tr}r_i(s^{-1})\)
for every \(i\in I_0\) at the same \(j=(w,A,s,1)\). Thus
\(x_j\in W\) and \(j\in\bigcap_iD_i\). This proves the asserted density
even in a nonabelian finite quotient.
\src{src/Development/AutomorphicArithmetic/Foundations/GaloisRepresentations/TraceLimits/LimFinOddOf.lean}{101}{Finite exclusions and inversion}
and \src{src/Development/AutomorphicArithmetic/Foundations/GaloisRepresentations/TraceLimits/LimFinOddOf.lean}{112}{simultaneous Hecke comparison}.

\begin{proposition}[Trace interpolation]\label{tr:interpolation}
Under the preceding hypotheses there is a continuous \(T:\Gamma\to L\). Whenever \((S,c,e)\) satisfies the inequality in \eqref{tr:approx}, it also satisfies
\[
  \left|Q(T)-\sum_{i\in S}c_iQ(t_i)\right|_\ell\le e
    \qquad(Q\in\mathbb Z[X_g:g\in\Gamma]).
\]
If a subfamily still gives approximations to every precision and all its indices satisfy \(j\in D_i\) at a particular \(j\), then \(T(x_j)=\xi_j\).
\end{proposition}
\begin{proof}
For each \(g\) use the filter on \(J\) whose basic sets consist of \(j\) such that \(x_j\) is in a prescribed neighborhood of \(g\) and \(j\) is good for a prescribed finite set of indices. Density makes this a proper filter. Given \(e>0\), take one approximant \(A=\sum_{i\in S}c_it_i\) from \eqref{tr:approx}. On the common good set for \(S\), substitution of \(P=X_j\) gives \(|\xi_j-A(x_j)|_\ell\le e\). Continuity of \(A\) gives \(|A(x_j)-A(g)|_\ell\le e\) on a smaller neighborhood. Thus the \(\xi_j\)'s are Cauchy along the filter for arbitrarily small \(e\). They take values in the complete \emph{finite extension} \(L\), giving a limit \(T(g)\). For any fixed approximant, continuity and passage to the limit give \(|T(g)-A(g)|_\ell\le e\) uniformly in \(g\). This proves continuity of \(T\) by uniform approximation. At the specified coordinate an approximant from the indicated subfamily also has \(|\xi_j-A(x_j)|_\ell\le e\). Sending \(e\) to zero proves the last assertion.

For a polynomial \(Q\), list its finitely many group-coordinate variables. Approximate each pair \((g,T(g))\) independently by \((x_j,\xi_j)\) with \(j\in\bigcap_{i\in S}D_i\). Substitute those indices for the variables of \(Q\). This gives a polynomial in the \(J\)-variables, to which the \emph{same} \(S,c,e\) apply. On that set \(\eta_i(j)=t_i(x_j)\). The inequality survives passage to the product limit since the polynomial evaluations and \(t_i\) are continuous and the closed norm ball is closed. \src{src/Development/AutomorphicArithmetic/Foundations/GaloisRepresentations/TraceLimits/ExistsTraceLimit.lean}{15}{filter construction} and \src{src/Development/AutomorphicArithmetic/Foundations/GaloisRepresentations/TraceLimits/ExistsTraceLimit.lean}{98}{transport of polynomial tests}.
\end{proof}

\subsection{Realization and recovery of local data}
\begin{proposition}[Trace identities, continuity and inertia]\label{tr:realize}
The trace in Proposition~\ref{tr:interpolation} is the trace of a continuous semisimple \(a\)-dimensional \(k\)-representation \(r\). If a subfamily still gives approximations at every precision for a specified place, and all its representations kill an inertia element \(s\), the resulting \(r\) kills \(s\) as well.
\end{proposition}
\begin{proof}
An integral polynomial vanishing on every seed trace vanishes on \(T\) by Proposition~\ref{tr:interpolation} and \(e\downarrow0\). In particular \(T(1)=a\) and \(T(gh)=T(hg)\). For \(g_1,\ldots,g_{a+1}\) consider the alternating sum over permutations of the product, over cycles, of the trace of the cyclic word. For a representation on \(V=k^a\), this sum is the trace of the tensor product of its \(a+1\) matrices times the antisymmetrizer on \(V^{\otimes(a+1)}\). Indeed, expansion in a tensor basis contracts precisely the indices in each cycle. The antisymmetrizer is zero since \(\bigwedge^{a+1}V=0\). Hence the integral polynomial defining this sum vanishes on \(T\). These are the characteristic-zero pseudocharacter identities. \src{src/Development/AbelianVarietyTheory/Proofs/InertiaEigenvalues/IsPseudoCharacterOfApproximatedByTraces.lean}{29}{passage of trace identities}.

Extend \(T\) by linearity to \(R=k[\Gamma]\). The alternating identity just proved is multilinear in its \(a+1\) arguments, since every argument occurs once in each permutation term. Expansion in the group basis therefore proves that identity throughout \(R\). Cyclicity extends by bilinearity. As \(k\) is a \(\mathbf Q\)-algebra, Chenevier's Proposition~1.27 gives a determinant \(D\) of dimension \(a\) whose trace is \(T\). Since \(k\) is algebraically closed, his Theorem~2.12 realizes \(D\) as the determinant of a unital semisimple \(R\)-module of dimension \(a\). Each \(g\in\Gamma\) acts invertibly with inverse the action of \(g^{-1}\), and the linear coefficient of \(D(1+tg)\) gives \(\operatorname{tr} r(g)=T(g)\). This proves algebraic realization over the specified field. \cite{chenevier}. The original construction uses the trace radical and the multiplicities of its matrix blocks. \src{src/Development/AbelianVarietyTheory/Proofs/InertiaEigenvalues/FiniteSemisimple.lean}{16}{Finite semisimplicity} and \src{src/Development/AbelianVarietyTheory/Proofs/InertiaEigenvalues/FiniteSemisimple.lean}{100}{block multiplicities}.

The image algebra \(A=\operatorname{span}_k r(\Gamma)\) is a product of the blocks occurring in this realization. The pairing \((u,v)\mapsto\operatorname{tr}(uv)\) is nondegenerate on \(A\), since each occurring block carries a positive multiplicity in characteristic zero and the ordinary matrix trace pairing is nondegenerate. Choose \(r(h_1),\ldots,r(h_b)\) spanning \(A\). The linear map \(u\mapsto(\operatorname{tr}(ur(h_j)))_j\) has a \(k\)-linear left inverse on its image, extendible to \(k^b\). Its matrix entries are fixed elements of \(k\), so it is continuous. The coordinates of \(r(g)\) under this map are \(T(gh_j)\), continuous functions of \(g\). Both the entries of \(r(g)\) and those of its inverse \(r(g^{-1})\) are continuous. Finally, for an inertia element as in the statement, uniform approximation of \(T(sg)\) and \(T(g)\) by the same weighted traces gives \(T(sg)=T(g)\). Hence \(\operatorname{tr}((r(s)-1)r(g))=0\) for all \(g\). Since the \(r(g)\)'s span \(A\), nondegeneracy gives \(r(s)=1\). \src{src/Development/AbelianVarietyTheory/Proofs/InertiaEigenvalues/ContinuousOfContinuousTrace.lean}{23}{trace coordinates} and \src{src/Development/AbelianVarietyTheory/Proofs/InertiaEigenvalues/KerOfLinApprox.lean}{10}{inertia kernel transfer}.
\end{proof}

At a split spherical place \(w\nmid\ell\), outside the exceptional levels stipulated in the theorem, let \(B_w\) be its Satake multiset and \(q_w\) the residue norm. The Hecke operator for the \(n\)th power sum, \(n\ge1\), has rational coefficients of the form \(m/q_w^e\). The normalization is \(B_w^{\mathrm{alg}}=q_w^{(a-1)/2}B_w\). For a chosen Frobenius representative \(\sigma_w\) the coordinates correspond to \((\sigma_w^{-1})^n\). With the place included in the prescribed subfamily, Proposition~\ref{tr:interpolation} gives the power sums for the same \(r\). The Newton relations
\[
 j e_j=\sum_{h=1}^j(-1)^{h-1}e_{j-h}p_h,\qquad e_0=1
\]
recover all elementary symmetric coefficients from \(p_1,\ldots,p_a\), so after the chosen \(k\simeq\mathbb C\) the characteristic polynomial of \(r(\sigma_w^{-1})\) is \(\prod_{\alpha\in B_w^{\mathrm{alg}}}(X-\alpha)\). The argument uses traces of powers of the matrix and works for a nondiagonalizable Frobenius. At \(w\mid\ell\), the actual proof also invokes a theorem on closedness of the crystalline and de Rham conditions and uses joint normalized Frobenius coordinates. Semisimplicity by itself does not give these conclusions. \src{src/Development/AbelianVarietyTheory/Proofs/SatakeParameters/FrobOpsIdx.lean}{20}{Frobenius indices}, \src{src/Development/AutomorphicArithmetic/Foundations/GaloisRepresentations/TraceLimits/LimFinOddOf.lean}{133}{recovery away from \(\ell\)} and \src{src/Development/AutomorphicArithmetic/Foundations/GaloisRepresentations/TraceLimits/LimFinOddOf.lean}{166}{separate conditions at \(\ell\)}.

In the active construction the odd-rank seed theorem and its weight-congruence theorem are themselves supplied by weight-one geometry and integral automorphic forms. The local--global row then uses the odd-rank limit with theta lifting, closedness of the Frobenius condition, descent off the selected set of places, and local compatibility at the remaining places. These last steps are cited for their original range of automorphic data. Proposition~\ref{tr:realize} explains the limit operation on this chain. The closed row at its end is what is supplied to the window and theta comparison branches. \src{src/Development/AutomorphicArithmetic/Foundations/GaloisRepresentations/IntegralStructures/GeometricFlagBases/GeometricRealizationOfWeightOne.lean}{25}{weight-one realization}, \src{src/Development/AutomorphicArithmetic/Foundations/GaloisRepresentations/IntegralStructures/Development.lean}{20}{assembly through the odd-rank limit} and \src{src/Development/AutomorphicArithmetic/Foundations/GaloisRepresentations/IntegralStructures/GaloisRepresentationsOfThetaLift.lean}{6}{instantiation of the geometric seed}.

\section{From weighted traces to crystalline lattices}\label{newgal:section}

Fix a CM field $E$ with maximal real subfield $E^+$, a prime $\ell$ and an
abstract field isomorphism $\iota:k=\overline{\mathbb Q}_\ell\simeq\mathbb C$.
All limits are taken on the $\ell$-adic side. We describe how the representation
of Section~\ref{tr:section} inherits a local condition at a place $w\mid\ell$.
The approximants are finite linear combinations of global characters. Their
coefficients can have arbitrarily large denominators. The local condition concerns
the restriction of the actual global representation, including any extensions in it.

\subsection{The geometric input and the rank of the seed}

In the odd-rank step there is a $d>0$ with $\sqrt{-d}\in E$ and $-d$ a square in
$\mathbb Q_\ell$. The degree $[E^+:\mathbb Q]$ is even, and each nonsplit finite
place of $E/E^+$ is unramified and has residue characteristic different from $\ell$.
Let $a$ be a positive odd integer. A usable place away from $\ell$ is a split
place $w$ whose rational residue characteristic $q$ satisfies
$-d\in(\mathbb Q_q^\times)^2$. Two places are apart when their restrictions
to $E^+$ differ. The auxiliary places $w_0,v'\mid\ell$ are apart. The definite
unitary automorphic representation $U$ is unramified at nonsplit places outside
a given finite set $T$ of nonsplit places. A seed has at $w_0$ a nontrivial
character of a cyclic $\ell$-power quotient of a torsion-free compact open
subgroup, chosen so that every irreducible smooth representation containing
that character is cuspidal. At $v'$ the seed is spherical and has a lacunary
set of $a$ jumps at one embedding. Lacunarity means that two disjoint nonempty
subsets have different arithmetic means.
\src{src/Development/AbelianVarietyTheory/Definitions/GaloisGeometry/LatticeDataThicknessBound.lean}{16}{Field and level assumptions},
\src{src/Development/AbelianVarietyTheory/Definitions/Divisors/GenOf.lean}{14}{local type}, and
\src{src/Development/AbelianVarietyTheory/Definitions/GaloisGeometry/DivisionOfDeterminantCharacter.lean}{15}{lacunarity}.

The geometric theorem has an intermediate output of variable dimension.
One takes a Hecke eigenspace in a cohomology tower, passes to a semisimplification
and twists by the Galois character attached to an algebraic Hecke character.
The resulting representation $V$ is continuous and semisimple. On the usable
spherical primes outside a finite set its geometric Frobenius is annihilated
by the normalized Hecke polynomial of degree $a$. At an auxiliary embedding
at least $a$ lacunary jumps occur. One can request additional geometric
information at one place by changing the tower or the Hecke eigenspace.
The separate requests yield a crystalline $V$ containing the prescribed jumps,
or a crystalline $V$ whose Frobenius polynomial at a degree $-1$ Weil element
divides a power of the predicted polynomial. The geometric statement requires
the specified places to restrict compatibly to $\mathbb Q(\sqrt{-d})$.
Above $\ell$ the two possibilities are exchanged by CM conjugation. Replacing
an auxiliary place by its conjugate preserves apartness at $E^+$, and the
argument uses conjugation of the local type, the spherical condition and
the lacunary condition to make these choices.
\src{src/Development/AutomorphicArithmetic/Foundations/GaloisRepresentations/IntegralStructures/GeometricFlagBases/HodgeTrAt.lean}{79}{Separate geometric outputs}
and \src{src/Development/AutomorphicArithmetic/Foundations/GaloisRepresentations/IntegralStructures/ZetaCharacterDivisibility/SeedFinOddOf.lean}{46}{conjugate-place choice}.

The rank-extraction input takes a continuous semisimple $k$-representation
$V$ of $G_E$ with a lacunary set of $a$ jumps at one labelled embedding
as above. Let $\psi:G_E\to k^\times$ be a continuous character. On a
dense set of $g$ suppose there is a monic polynomial $Q_g$ of degree $a$
annihilating $g$ on $V$ and satisfying $(-1)^a Q_g(0)=\psi(g)$.
The source supplies a continuous semisimple $W$ of dimension $a$ with
equivariant injections
\begin{equation}\label{newgal:injections}
                 W\hookrightarrow V\hookrightarrow W^{\oplus m}
\end{equation}
for some $m$, such that the elements acting with distinct eigenvalues on $W$
form a dense open set. The determinant character in the application comes
from the central character and the Satake normalization. Proposition~\ref{famgal:rank} below proves the rank extraction from the
Hodge cocharacter. The existence of the initial geometric witness is an
arithmetic input here.
\src{src/Development/AutomorphicArithmetic/Foundations/GaloisRepresentations/IntegralStructures/ZetaCharacterDivisibility/PhiOfDivOf.lean}{122}{Rank-extraction hypothesis and conclusion}
and \src{src/Development/AutomorphicArithmetic/Foundations/GaloisRepresentations/IntegralStructures/ZetaCharacterDivisibility/DivisionOneGeometricRepresentation.lean}{52}{division and cocharacter theorems}.

\begin{proposition}\label{newgal:seed}
These inputs identify the same $W$ for all geometric witnesses for a fixed
seed eigensystem. If the witness at $w\mid\ell$ is crystalline and contains
the $a$ prescribed distinct jumps at a labelled embedding, $W$ is crystalline
there and has precisely those jumps. If the crystalline Frobenius polynomial
of the witness divides a power of the predicted polynomial, the predicted
polynomial is the Frobenius polynomial of $W$ whenever the Satake multiset
satisfies the multiplicative genericity condition below.
\end{proposition}
\begin{proof}
At a good geometric Frobenius in the open set of distinct eigenvalues,
the degree $a$ Hecke polynomial annihilates its action on $W\subset V$.
The minimal and characteristic polynomials both have degree $a$, so both
equal the monic Hecke polynomial. For two geometric witnesses, intersect
their dense open sets and discard the union of their finite exceptional
sets. Usable geometric Frobenius elements remain dense there by the density
argument of Section~\ref{tr:section}. Equality of traces extends by continuity
to $G_E$, and semisimplicity in characteristic zero identifies the two $W$'s.
The same argument applied to the constant coefficient gives $\det W=\psi$.
Thus the $W$ already selected admits \eqref{newgal:injections} with each
new witness, up to isomorphism.

Since the global representations are semisimple, both injections split
equivariantly. The splittings persist after restriction to a decomposition
group. Crystalline periods commute with finite direct sums and direct
summands, as do the labelled graded pieces of the de Rham filtration.
Hence $W$ is crystalline when $V$ is, and every jump of $V$ occurs in
$W^{\oplus m}$, hence in $W$. At a fixed embedding the graded dimensions
of $W$ sum to $a$. The occurrence of $a$ distinct prescribed integers forces
multiplicity one at each of them.

Write $S=q_w^{(a-1)/2}B_w$. Multiplicative genericity says that a multiset
of size $a$ supported on $S$ and having product $\prod_{s\in S}s$ equals
$S$, with its multiplicities. Scaling all entries by the same nonzero
scalar preserves this condition. The Frobenius polynomial on $W$ divides
that on $V$ by the splitting. Its roots are supported on $S$, it has
degree $a$, and its root product agrees with that of $S$ by the local
compatibility of the rank-one determinant character. This proves equality
with $\prod_{s\in S}(X-s)$. At a crystalline place the associated Weil
action is unramified, so the characteristic polynomials for all Weil
elements are determined by powers of the degree $-1$ element. Away from
$\ell$ one uses the annihilating polynomial and the same product test.
\src{src/Development/AutomorphicArithmetic/Foundations/GaloisRepresentations/IntegralStructures/ZetaCharacterDivisibility/PhiOfDivOf.lean}{166}{Comparison across geometric realizations}
and \src{src/Development/AutomorphicArithmetic/Foundations/GaloisRepresentations/IntegralStructures/ZetaCharacterDivisibility/PRootPin.lean}{18}{root multiplicities}.
\end{proof}

\subsection{Extracting the rank from a lacunary Hodge type}

\begin{proposition}[Rank extraction in every component]\label{famgal:rank}
Let \(\Gamma\) be profinite and \(V\) a continuous semisimple representation
over \(k=\overline{\mathbf Q}_\ell\). Suppose that on a dense subset the
operator on \(V\) is annihilated by a monic polynomial of degree \(a\ge1\).
Suppose that the algebraic monodromy group of \(V\), after scalar extension
to \(\mathbf C_\ell\), contains a cocharacter for which \(a\) distinct
weights occur. Assume that no two disjoint nonempty subsets of these
\(a\) integers have the same arithmetic mean. There is a semisimple
\(k\)-representation \(W\) of dimension \(a\) with
\(W\hookrightarrow V\hookrightarrow W^{\oplus h}\) for some \(h\).
The elements of \(\Gamma\) having simple spectrum on \(W\) form a dense
open subset.
\end{proposition}
\begin{proof}
Write \(G\) for the Zariski closure of the image and \(H=G^\circ\).
Every element \(A\) of \(G\) has minimal polynomial of degree at most
\(a\). Indeed, linear dependence of \(1,A,\ldots,A^a\) is the
vanishing of minors, and it holds first on the given dense set.
These equations persist over \(\mathbf C_\ell\). The representation of
\(G\) is semisimple since its invariant subspaces are exactly those of
\(\Gamma\). The unipotent radical of \(H\) is normal in \(G\), and its
fixed space in every irreducible \(G\)-summand is nonzero by triangularization.
It thus acts trivially on the faithful representation \(V\), so \(H\)
is reductive.

Choose a maximal torus \(T\) and a Borel subgroup containing it over \(k\).
By conjugacy of maximal tori, we may evaluate the given cocharacter on
the \(T\)-weights after extending scalars. At a general point of \(T\)
all distinct weights give distinct eigenvalues, since their pairwise
quotients are nontrivial characters. There are at most \(a\) weights
and the cocharacter separates at least \(a\). Thus there are exactly
\(a\), all separated, and their cocharacter values are the given set.

In an irreducible \(H\)-module all weights differ by elements of the
root lattice. Averaging over the Weyl group annihilates each root,
since its reflection negates it. The average of each Weyl orbit in
this module is consequently the same rational character. Two different
orbits would give disjoint nonempty subsets of the cocharacter weights
with equal means. There is therefore just one orbit. Its highest weight
has multiplicity one, and Weyl invariance of multiplicities shows that
every weight in the orbit has multiplicity one. If two irreducibles
share a weight they share the orbit and its dominant highest weight,
so they are isomorphic. We use here the ordinary highest-weight theorem
for connected reductive groups in characteristic zero.

We check the action of the components on the multiplicity spaces.
Let \(n\in G(k)\) normalize \(T\) and the chosen Borel. It permutes
the \(a\) weights. For a cycle \(C\) of length \(d_C\), put
\(w_C=\prod_{\alpha\in C}\alpha\). On one weight space \(V_\alpha\)
in the cycle, \((tn)^{d_C}\) acts as \(w_C(t)n^{d_C}\), \(t\in T(k)\).
The spectrum of \(tn\) on the cycle contains all \(d_C\)-th roots of
\(w_C(t)c\) for each eigenvalue \(c\) of \(n^{d_C}|_{V_\alpha}\).
Indeed, propagate an eigenvector around the cycle, dividing by the
chosen root at each step. The resulting components close up on return.
Two different cycles \(C,D\) have disjoint spectra for general \(t\).
An overlap would impose
\(w_C(t)^{d_D}/w_D(t)^{d_C}=c_D^{d_C}/c_C^{d_D}\) for a pair of eigenvalues.
The character on the left is nontrivial by the distinct arithmetic
means on the cocharacter. There are only finitely many pairs.
Thus \(tn\) has at least \(\sum_C d_C=a\) distinct eigenvalues.
Its minimal polynomial has degree at most \(a\), so it is diagonalizable
with exactly this many eigenvalues. Each \(n^{d_C}|_{V_\alpha}\)
is diagonalizable and has only one eigenvalue, hence is scalar.
In particular \(n\) acts by a scalar on any weight space that it fixes.

In each orbit of \(G/H\) on the \(H\)-isotypic spaces choose one
representative and let \(\alpha\) be its highest weight.
The stabilizer of this isotypic space preserves the \(H\)-submodule
generated by any vector of \(V_\alpha\). Indeed, write a stabilizing
element as \(h n\) with \(h\in H\) and \(n\) normalizing the Borel
and torus, using conjugacy of Borel--torus pairs in \(H\).
The element \(n\) fixes \(\alpha\) and acts by a common scalar on
\(V_\alpha\). An \(H\)-intertwiner taking one highest vector to another
also commutes with the stabilizer: on \(h'v\) this follows by writing
\(n h'v=(n h'n^{-1})nv\). A basis of \(V_\alpha\) therefore gives
isomorphic stabilizer modules whose sum is the chosen isotypic space.
Transport each of these modules to the other isotypic spaces in its
orbit. Choice of coset representatives has no effect on the transported
subspaces, and the same intertwiners transport consistently.
Their orbit sums give isomorphic \(G\)-modules. In each orbit retain
one such sum and take their direct sum for \(W\). Every \(T\)-weight
now occurs exactly once. This gives \(\dim W=a\) and the two
injections, taking \(h\) to be the maximum multiplicity.
As a summand of \(V\), \(W\) is semisimple and continuous.

Every component has a representative \(n\) normalizing the Borel and
torus. The cycle calculation on its one-dimensional weight spaces
gives simple spectrum on \(W\) for general \(tn\). The nonvanishing
of the discriminant is thus a nonempty Zariski open set in every
component. The preimage of \(H\) in \(\Gamma\) is open and closed.
Its image is Zariski dense in \(H\), as is the image of each finite-index
subgroup of it, since an irreducible variety cannot be a finite union
of proper closed subsets. Each open subset of \(\Gamma\) contains a
coset of such an open subgroup. Its image meets the discriminant open
in its component. Continuity gives openness on \(\Gamma\) as well.
\end{proof}

For the seed above, use the labelled Hodge--Tate decomposition of the
geometric witness at its distinguished embedding. Its grading defines
a cocharacter in the local algebraic monodromy group, hence in \(G\).
To see the membership, the Hodge--Tate decomposition respects tensor
products, duals, morphisms and subrepresentations of Hodge--Tate
representations. The grading therefore preserves each stable line in
tensor constructions, including direct sums of such constructions.
The stabilizer description of an algebraic subgroup of a general linear
group (Chevalley's theorem) puts the grading in the scalar extension of
that subgroup. With finite coefficients, take the component of
\(\mathbf C_\ell\otimes_{\mathbf Q_\ell}L\) at the embedding giving the
specified label, extending the local embedding to \(\mathbf C_\ell\)
accordingly. This is the labelled cocharacter construction of
Dospinescu, Pa\v{s}k\=unas and Schraen~\cite{senCharacters}.
Changing the sign convention for all weights preserves lacunarity.
Proposition~\ref{famgal:rank} now supplies \eqref{newgal:injections}
for every such witness. The determinant identification is then the
dense Frobenius calculation in Proposition~\ref{newgal:seed}.
The construction still needs the geometric witness and the simultaneous
Hecke congruences at the claimed weights and places.
\src{src/Development/AutomorphicArithmetic/Foundations/GaloisRepresentations/IntegralStructures/ZetaCharacterDivisibility/HtCocharOf.lean}{20}{Hodge cocharacter in the closure},
\src{src/Development/EtaleCohomology/DivisorCycles/GoalHolds.lean}{20}{algebraic rank-extraction input}, and
\src{src/Development/AutomorphicArithmetic/Foundations/GaloisRepresentations/IntegralStructures/ZetaCharacterDivisibility/PDivSet.lean}{16}{application to the Galois seed}.


\subsection{Integral quotients of the regular image}

Let $\Gamma$ be a compact Hausdorff group, $r:\Gamma\to\mathrm{GL}_a(k)$
a continuous semisimple representation and $(r_i)$ a family of continuous
$a$-dimensional representations over $k$. Give $k$ the norm $|\ell|=\ell^{-1}$.
The approximation needed for crystalline closure is
\begin{equation}\label{newgal:lin}
 \inf_{I,c}\ \sup_{g\in\Gamma}
 \left|\operatorname{tr}r(g)-\sum_{i\in I}c_i\operatorname{tr}r_i(g)\right|=0 ,
\end{equation}
where $I$ is a finite subfamily and $c_i\in k$. Fix a finite extension
$K/\mathbb Q_\ell$ and a continuous homomorphism $G_K\to\Gamma$. In the
application $K=E_w$ and this homomorphism identifies $G_K$ with a decomposition
subgroup. Write $V_i,V$ for the restricted modules.

\begin{lemma}[Finite coefficient fields]\label{newgal:coeff}
Each continuous representation of $\Gamma$ over $k$ can be written, in its
given basis, over a finite extension of $\mathbb Q_\ell$.
\end{lemma}
\begin{proof}
There are countably many finite subextensions of $k/\mathbb Q_\ell$.
For each degree there are finitely many local extensions up to isomorphism
by the local-field finiteness theorem, and each has finitely many embeddings.
For each such $L$ the compact image group meets $\mathrm{GL}_a(L)$ in a
closed subgroup, since $L$ is complete. These subgroups cover the image.
Baire's theorem gives one with nonempty interior. It is open and of finite
index. Adjoin the entries of finitely many coset representatives to $L$.
All matrices are defined over the resulting finite extension.
\end{proof}

Choose a field of definition $L$ for $r$. The sum of the translates of any
$\mathcal O_L$-lattice is bounded by compactness and contains the original
lattice. It is therefore a stable lattice. Choose its basis, and do the same
for each $r_i$ over a field $L_i$. Define
\[
 R= r(\mathbb Z_\ell[\Gamma])\subset M_a(\mathcal O_L),\qquad
 R_I=(r_i)_{i\in I}(\mathbb Z_\ell[\Gamma])
       \subset\prod_{i\in I}M_a(\mathcal O_{L_i}).
\]
These are finite free $\mathbb Z_\ell$-modules, stable under left translation
by $\Gamma$. We use the ordinary group ring with finite supports. Both image
modules are already closed, since finite free $\mathbb Z_\ell$-modules are compact.

\begin{lemma}[Lattice quotients from translated traces]\label{newgal:quot}
For every positive integer $n$ there is a finite $I$ and a
$\Gamma$-equivariant surjection $R_I\twoheadrightarrow R/\ell^nR$,
given by $((r_i(f))_i)\mapsto r(f)\bmod\ell^nR$.
\end{lemma}
\begin{proof}
In the $k$-algebra spanned by $r(\Gamma)$ the matrix trace pairing is
nondegenerate by Proposition~\ref{tr:realize}. Choose finitely many $h_j$
so that $r(h_j)$ span this algebra over $k$. On the finite-dimensional
$\mathbb Q_\ell$-space $R[1/\ell]$ the map
\[
   x\longmapsto(\operatorname{tr}(xr(h_j)))_j\in L^b
\]
is therefore injective. The maximum of the norms of its coordinates is
a norm. Comparison with the lattice norm having unit ball $R$ gives an
integer $c\ge0$, depending only on $r$, such that
\[
 \max_j|\operatorname{tr}(xr(h_j))|\le\ell^{-n-c}
       \quad\Longrightarrow\quad x\in\ell^n R.
\]
Use \eqref{newgal:lin} with error at most $\ell^{-n-c}$.
If $f=\sum b_g[g]$ maps to zero in $R_I$, then
$\sum_g b_g\operatorname{tr} r_i(gh_j)=0$ for every $i,j$.
The ultrametric inequality and $|b_g|\le1$ give
$|\operatorname{tr}(r(f)r(h_j))|\le\ell^{-n-c}$ for all $j$.
This proves well-definedness. The formula gives both surjectivity and
equivariance. The kernel contains $\ell^nR_I$, and so is a full stable
lattice in $R_I[1/\ell]$. Global semisimplicity was used before restriction
to $G_K$.
\end{proof}

\begin{theorem}[Crystalline closure from weighted traces]\label{newgal:liu}
Suppose all $V_i$ are crystalline over $K$ and all their labelled filtration
jumps belong to a common finite interval $[u,v]$. Under \eqref{newgal:lin},
$V$ is crystalline over $K$ with jumps in the same interval.
The assertion also holds throughout upon replacing crystalline by semistable.
\end{theorem}
\begin{proof}
Use the usual Fontaine period rings for $K$. With left multiplication on
matrices, the ambient $\mathbb Q_\ell[G_K]$-module of $R_I[1/\ell]$
is a direct sum of $a$ copies of each $\operatorname{Res}_{L_i/\mathbb Q_\ell}V_i$,
using models over the finite extensions $L_i/\mathbb Q_\ell$. Forgetting coefficients collects the jumps
at the different embeddings. In particular the common interval remains valid.
Crystalline (respectively semistable) representations with jumps in that
interval are stable under direct sums and subquotients by exactness and
strictness of the filtered period functor on the admissible category.
Consequently $R_I[1/\ell]$ has this property. Lemma~\ref{newgal:quot}
exhibits $R/\ell^nR$ as a quotient of two full $G_K$-stable $\mathbb Z_\ell$-lattices
inside that single $\mathbb Q_\ell$-representation for each $n$. The ambient
dimension may grow with $n$.

Apply the bounded torsion theorem of T. Liu to $T=R$. In Liu's
contravariant convention the weights are those of the filtered module
of the dual, so our interval becomes $[-v,-u]$. A common integral Tate
twist translates it to $[0,v-u]$. Twisting the quotient sequences uses
the same rank-one $\mathbb Z_\ell$-lattice on both sides. Liu's result
yields the indicated property for $R[1/\ell]$ after undoing the twist.
It allows an arbitrary finite extension $K/\mathbb Q_\ell$, arbitrary
positive height, and $\ell=2$ with the author's correction. For $u=v$ apply
the positive-height statement to both $[u-1,u]$ and $[u,u+1]$ and intersect.
T. Liu,
\emph{Torsion $p$-adic Galois representations and a conjecture of Fontaine},
proved Conjecture 1.0.1, pp.~633--634~\cite{liuTorsion} and
Liu,
\emph{A note on lattices in semi-stable representations}, errata for the 2007 paper~\cite{liuLattices}.

For any $v_0\in L^a$, evaluation $x\mapsto xv_0$ on $R[1/\ell]$ is
$G_K$-equivariant. Taking $v_0$ through a $\mathbb Q_\ell$-basis of $L^a$
yields a surjection from a finite sum of copies of $R[1/\ell]$ onto $L^a$,
since $1\in R$. The quotient is the full restriction of $r$.
The $L$-action commutes with the period comparison map, giving the
coefficient version of the assertion.
\end{proof}

For the automorphic limit take the subfamily indexed by the singleton
$\{w\}$, where $w\mid\ell$ is spherical and apart from $w_0,v'$.
At every embedding $\tau:E\to k$ inducing $w$ the weight at $\iota\tau$
agrees throughout the subfamily. In the zero-based indexing of the highest
weight $\lambda_{\iota\tau}$ the prescribed jumps are
\begin{equation}\label{newgal:jumps}
             J_\tau=\{j-\lambda_{\iota\tau,j}:0\le j<a\}.
\end{equation}
Indeed the $j$th infinitesimal-character entry equals
$\lambda_{\iota\tau,j}+(a-1)/2-j$. Taking the minimum and maximum of
\eqref{newgal:jumps} over $j$ and the finitely many $\tau$ over $w$
gives the common interval. The singleton subfamily retains polynomial
approximation at every precision, hence \eqref{newgal:lin} on the whole
global group. Theorem~\ref{newgal:liu} therefore proves crystallinity
of the limit at this $w$.
\src{src/Development/AutomorphicArithmetic/Foundations/GaloisRepresentations/TraceLimits/LimFinOddOf.lean}{166}{Singleton restriction}
and \src{src/Development/AbelianVarietyTheory/Proofs/SatakeParameters/BlowUpFibreSemilocal.lean}{29}{jump convention}.

\subsection{Comparison of exact local invariants}

For exact jumps and Frobenius the source keeps more information than the
surjection of Lemma~\ref{newgal:quot}. Let $b$ be a $\mathbb Z_\ell$-basis
of $R$, let $e$ be the column of coordinates of $1$, and let $t$ be the
row of traces of the $b_j$. For precision $n$ choose the approximation of
Lemma~\ref{newgal:quot}. Enlarge its support $I_n$ to contain all previous
supports, assigning coefficient zero to the added members. Set
\[
 D_n=\prod_{i\in I_n}M_a(\mathcal O_{L_i}),\qquad
 B_n=R_{I_n}+\ell^{M_n}D_n .
\]
Choose increasing $M_n$ with $\ell^{-M_n}\max(1,\max_i|c_i|)
\le\ell^{-n-c}$. If $(r_i(f))_i\in\ell^{M_n}D_n$, then the
translated traces of $r(f)$ against all $r(h_j)$ have norm at most
$\ell^{-n-c}$ by \eqref{newgal:lin} and integrality of all $r_i(h_j)$.
The map of Lemma~\ref{newgal:quot} thus extends to $B_n$ by zero on
$\ell^{M_n}D_n$. Choose a basis of $B_n$ and any integral matrix $P_n$
lifting that map. Write $e_n$ for the coordinates of the tuple of identities,
and $t_{ni}$ for the row of traces after projection of the basis to factor $i$.
Then
\begin{equation}\label{newgal:marked}
 |P_ne_n-e|\le\ell^{-n},\qquad
 \left|tP_n-\sum_{i\in I_n}c_i t_{ni}\right|\le\ell^{-n}.
\end{equation}
Here matrix norms are maximum-entry norms. The first bound follows
by evaluating the quotient map at $1$. For the second, write each basis
vector as $(r_i(f))_i+\ell^{M_n}z$. Its chosen lift agrees with $r(f)$ modulo
$\ell^nR$, whose traces have norm at most $\ell^{-n}$. The approximation
applied to $f$, and the choice of $M_n$ applied to $z$, give the bound.
The coordinate projection $B_{n+1}\to B_n$ intertwines the actions exactly
and commutes with the maps to $R/\ell^n R$. Both claims follow on group
elements and on $\ell^{M_{n+1}}D_{n+1}$. These constructions work after
restriction to $G_K$ and are compatible with any additional scalar
approximated by the same coefficients.
\src{src/Development/AbelianVarietyTheory/Proofs/InertiaEigenvalues/CoverExists.lean}{23}{Integral model with trace and unit coordinates}.

\begin{proposition}[Reading a labelled Hodge multiset]\label{newgal:hodge}
Fix an embedding $\tau:K\to k$. Suppose the comparison on the preceding
lattices supplies a non-root of unity $z\in\mathbb C_\ell^\times$ and matrices
$U,U_n$ on $R,B_n$ over $\mathbb C_\ell$ with the following properties.
Their norms are bounded by a common $B\ge1$. For every $N$ there is $n\ge N$
with $|UP_n-P_nU_n|\le B\ell^{-N}$. After the same isometric coefficient
embedding on all trace rows and all $c_i$, for $j\ge0$ one has
\[
 tU^je=\sum_{h\in{\rm HT}_\tau(V)}z^{jh},\qquad
 t_{ni}U_n^je_n=\sum_{h\in J_\tau}z^{jh}\quad(i\in I_n).
\]
Here $\operatorname{HT}_\tau$ denotes the multiset of labelled filtration jumps.
Then $\operatorname{HT}_\tau(V)=J_\tau$.
\end{proposition}
\begin{proof}
Put $S_j=\sum_{h\in J_\tau}z^{jh}$. For fixed $j$, the telescoping formula
for $U^jP_n-P_nU_n^j$ bounds its norm by $B^j\ell^{-N}$ for $j\ge1$.
The maximum-entry norm is submultiplicative in all dimensions here.
In $tU^je$ replace $e$ by $P_ne_n$, commute $P_n$ past $U^j$, then
replace $tP_n$ by $\sum c_it_{ni}$. Equations~\eqref{newgal:marked}
bound the total error by $\max(1,|t|)B^j\ell^{-N}$. Thus the result is
$S_j\sum c_i$ to that precision. At $j=0$ both power sums equal $a$,
so $|1-\sum c_i|=O(\ell^{-N}/|a|)$. The same coefficients serve at
$j$ and at zero. Let $N$ tend to infinity. Newton's identities now
identify the multisets $\{z^h:h\in{\rm HT}_\tau(V)\}$ and
$\{z^h:h\in J_\tau\}$. The powers of $z$ distinguish the integers.
\end{proof}

The period-theoretic hypothesis in Proposition~\ref{newgal:hodge} is
quantitative. In the source the height is bounded by
$H=\max_{\tau,h\in J_\tau}|h|$. For its Breuil--Kisin period
$u\in A_{\rm inf}$ put $s_0=\theta(u^{2H})\ne0$, where
$\theta:A_{\rm inf}\to\mathcal O_{\mathbb C_\ell}$ is Fontaine's map.
The linked frames give integral evaluated change-of-basis matrices
$Y,Y'$ with integral matrices $Y_1,Y'_1$ satisfying
$YY_1=Y'Y'_1=s_0\,1$ in the respective dimensions.
For integral cocycle matrices $C,C'$ the operators satisfy $UY=YC$,
$U_nY'=Y'C'$. A link at level $N$ gives an integral rectangular matrix
$D$ with $|P_nY'-YD|,\ |CD-DC'|\le\ell^{-N}$. Thus
$|U|,|U_n|\le |s_0|^{-1}$. The identity
\[
 (UP_n-P_nU_n)Y'=U(P_nY'-YD)+Y(CD-DC')-(P_nY'-YD)C'
\]
and multiplication by $Y'_1/s_0$ give the claimed bound with
$B=|s_0|^{-2}$, independent of the growing dimensions.
The other period assertion identifies the operator on a Hodge-adapted
basis with a diagonal matrix of powers $z^h$ and transports its powers
through each column projection. The construction of those linked
bounded-height frames and of the adapted comparison, needed at every
precision and for the selected embedding, are cited here.
\src{src/Development/AbelianVarietyTheory/Proofs/InertiaEigenvalues/ExistsDeRhamFrameLimitLatticeBreuilKisinTower.lean}{16}{Linked frames},
\src{src/Development/AbelianVarietyTheory/Proofs/InertiaEigenvalues/ThetaConjBoundedCompatibleBreuilKisinFrames.lean}{117}{uniform estimate}, and
\src{src/Development/AbelianVarietyTheory/Proofs/InertiaEigenvalues/ClosedOfCover.lean}{159}{labelled comparison}.

At a crystalline place put $D_i=\iota\mathrm{WD}(V_i)$ and
$D=\iota\mathrm{WD}(V)$, and set $p_i=\iota^{-1}\operatorname{tr}D_i(g)$
and $p=\iota^{-1}\operatorname{tr}D(g)$ for $g$ in the Weil group of $K$.
Write $B_w$ for the spherical Satake
multiset. At $g^j$ the predicted value is
\[
 \vartheta_j=\iota^{-1}\left(\sum_{x\in q_w^{(a-1)/2}B_w}x^{-j\deg g}\right).
\]
The Hecke approximation at each fixed $j$ uses the same $c_i$ for the global
trace and this additional scalar, to arbitrary $\ell$-adic precision.
One removes a fixed nonzero normalization factor by increasing the required
precision. For the next proposition write $\vartheta$ for the target at $g$.

\begin{proposition}[Comparison of a crystalline Frobenius trace]\label{newgal:phi}
Assume the marked approximations \eqref{newgal:marked} satisfy also
$|\vartheta-\sum_{i\in I_n}c_i p_i|\le\ell^{-n}$. Put
$F=\mathbb Q_\ell^{\mathrm{nr}}\subset k$ and use the unramified embedding
$F\hookrightarrow B_{\rm cris}$ taking $\mathcal O_F$ into $A_{\rm cris}$.
Suppose there is a single $s\in A_{\rm cris}$ invertible in $B_{\rm cris}$,
and matrices $A,A_n$ over $A_{\rm cris}$ on the respective lattices, with
the following two properties. For every $N$ some $n\ge N$ has every entry
of $AP_n-P_nA_n$ in $\ell^NA_{\rm cris}$. For every $F$-linear
$\lambda:k\to F$, acting entrywise on the rows, one has
\[
 \lambda(t)Ae=s\lambda(p),\qquad
 \lambda(t_{ni})A_ne_n=s\lambda(p_i)\quad(i\in I_n)
\]
in $B_{\rm cris}$. Then $p=\vartheta$.
\end{proposition}
\begin{proof}
Suppose $\Delta=p-\vartheta\ne0$. Fix integers $c_1,c_2\ge0$ such that
$\ell^{c_1}t$ is integral and $\ell^{-c_2}\le|\ell^{c_1}\Delta|\le1$.
Choose $n$ as in the first hypothesis for $N$ and write
$\delta=tP_n-\sum_i c_i t_{ni}$ and $\eta=\vartheta-\sum_i c_i p_i$.
Both have norm at most $\ell^{-N}$. There is an $F$-linear $\lambda$
taking integral values on the entries of $\ell^{c_1}t,\ell^{-N}\delta,
\ell^{-N}\eta$ and on $z=\ell^{c_1}\Delta$, with $|\lambda(z)|\ge|z|$.
Indeed their $\mathcal O_F$-span is finite free over the discrete valuation
ring $\mathcal O_F$. All elements of that span are integral in $k$. In any
basis of the span one coordinate of $z$ has absolute value at least $|z|$.
Extend this coordinate functional $F$-linearly to $k$.

Apply the seed identities to $x\mapsto\lambda(c_i x)$, subtract and
insert $P_n$. This gives
\[
 \begin{split}
 s\lambda(\Delta)={}&\lambda(t)\big[(AP_n-P_nA_n)e_n+A(e-P_ne_n)\big]\\
                    &+\lambda(\delta) A_ne_n-s\lambda(\eta).
 \end{split}
\]
Multiplication by $\ell^{c_1}$ puts the right side in $\ell^NA_{\rm cris}$.
Multiplication next by $\ell^{c_2}/\lambda(z)\in\mathcal O_F$ gives
$\ell^{c_2}s\in\ell^NA_{\rm cris}$ for every $N$.
This contradicts $\ell$-adic separation and torsion-freeness of
$A_{\rm cris}$, since $s$ is invertible in $B_{\rm cris}$.
\end{proof}

Here is the integral calculation giving the first Frobenius comparison
from the period frames. Use a target precision \(N\) and perform the
calculation in \(A_{\rm cris}/\ell^N\). For the two Frobenius iterates the
linked period bases give square matrices \(W,Z\) and \(W',Z'\) on the
target and seed-product lattices, respectively. Their scaled inverses
satisfy \(Z^*Z=b1,\ Z'Z'^*=b1\) for the same scalar \(b\). The two
frame congruences read \(PZ'=ZD,\ PW'=WD\), where \(P=P_n\) and \(D\)
have size \(\operatorname{rank}R\) by \(\operatorname{rank}B_n\).
In this quotient,
\[
 bWZ^*P=WZ^*PZ'Z'^*=WZ^*ZDZ'^*=bWDZ'^*=bPW'Z'^*.
\]
Thus the integral operators \(A=bWZ^*,\ A_n=bW'Z'^*\) intertwine to
precision \(N\), without inverting \(P\) or cancelling \(b\).
In the construction, if \(H\) bounds the height, the integer \(k_0\)
for the period denominators is chosen before the level. With
\(\tau_1\) the image of the chosen period in \(A_{\rm cris}\), one has
\[
 b=\varphi^{n_2}(\tau_1^{2H}\ell^{2k_0}),\qquad
 s=\varphi^{n_2}(\tau_1)^{3H}\ell^{4k_0}\varphi^{n_1}(\tau_1)^H.
\]
Both lattice dimensions use these same scalars. To obtain the two frame
congruences modulo \(\ell^N\), first ask for the frame-link precision
\(N'\) required by the period comparison, then choose one level
\(n\ge\max(N',N)\). At that level both iterations compare with the
same \(D\). The unit error, weighted trace-row error and the error in
the additional scalar are already at most \(\ell^{-n}\). In particular
one never multiplies an error of size \(\ell^{-N}\) for an individual
seed by \(c_i\). The potentially unbounded \(\max_i|c_i|\) entered
earlier in the choice of \(M_n\) in \eqref{newgal:marked}.
\src{src/Development/AutomorphicArithmetic/Foundations/GaloisRepresentations/ClosedPhiFunctoriality/OpExHolds.lean}{75}{Common denominator before the level},
\src{src/Development/AutomorphicArithmetic/Foundations/GaloisRepresentations/ClosedPhiFunctoriality/OpExHolds.lean}{86}{one level for both iterates}
and \src{src/Development/AutomorphicArithmetic/Foundations/GaloisRepresentations/ClosedPhiFunctoriality/LinkOp.lean}{8}{the integral rectangular identity}.

One can also see exactly where the unit and trace marks enter the
second Frobenius comparison. For either lattice let \(b_i\) be its
matrix readings (on a product lattice, project to one seed factor);
thus \(\sum_i e_i b_i=1_a\), \(t_i=\operatorname{tr}(b_i)\).
Work after extending the Weil module to the common coefficient and
period ring \(k\otimes_F B_{\rm st}\), writing the two inclusions of
scalars implicitly. Write \(U_j=\varphi^{n_j}(U)\) for the invariant
period basis of this lattice, so \(AU_2=sU_1\). Denote the \(h\)th
standard vector in coefficient dimension \(a\) by \(u_h\).
The column comparison for the crystalline periods supplies vectors
\(v_{jh}\) in the Weil module, whose comparison coordinates are the
columns \(\sum_i (U_2)_{ij}b_i u_h\). Their transforms under \(g\)
have coordinates \(\sum_i (U_1)_{ij}b_i u_h\). This is the period
input identifying the sign of the Weil action. Now set \(z=U_2^{-1}e\).
For each \(h\) the vector \(\sum_j z_jv_{jh}\) has comparison
coordinates \(\sum_i e_i b_i u_h=u_h\).
Applying the scalar-extended Weil operator and summing its diagonal
coordinates gives
\[
 p\otimes 1=\sum_i \operatorname{tr}(b_i)\otimes(U_1z)_i,\qquad
 p\otimes s=\sum_i t_i\otimes(Ae)_i .
\]
Applying \(\lambda\otimes1\) for any \(F\)-linear \(\lambda:k\to F\)
gives precisely \(\lambda(t)Ae=s\lambda(p)\); the terms belong to
\(B_{\rm cris}\subset B_{\rm st}\). The same proof for each column
projection of the product lattice gives every seed identity before
weighting it. Only the column comparison and the invariant periods
are used here beyond linear algebra.
\src{src/Development/AutomorphicArithmetic/Foundations/GaloisRepresentations/ClosedPhiFunctoriality/TieRep.lean}{75}{The two systems of period columns},
\src{src/Development/AutomorphicArithmetic/Foundations/GaloisRepresentations/ClosedPhiFunctoriality/TieRep.lean}{94}{The unit frame and trace contraction}
and \src{src/Development/AutomorphicArithmetic/Foundations/GaloisRepresentations/ClosedPhiFunctoriality/ClosedPhiOfClosedTr.lean}{65}{Application to the target and each product projection}.

In the crystalline comparison the two nonnegative Frobenius iterates are
$n_1=\max(0,-\alpha)$ and $n_2=\max(0,\alpha)$, where
$\alpha=[k(w):\mathbb F_\ell]\deg g$ is the signed exponent on the
unramified coefficient field and $k(w)$ denotes the residue field of $K$.
For an invertible matrix $U$ of invariant periods one requires
$A\varphi^{n_2}(U)=s\varphi^{n_1}(U)$, and likewise at each level.
The common $s$ and congruences of these operators are assertions about
the linked integral frames. The second hypothesis of the proposition
identifies their values on the marked unit, paired with the marked trace
rows, with the Weil trace via the period comparison. Using $F$ and its
linear functionals avoids choosing embeddings of all ramified coefficient
fields into $B_{\rm cris}$. Given these two frame assertions, apply
Proposition~\ref{newgal:phi} at the first $a$ positive powers of a
geometric Frobenius. Newton's identities give the complete characteristic
polynomial. Crystallinity gives trivial inertia and zero monodromy,
so this determines the Frobenius-semisimplified Weil--Deligne parameter.
\src{src/Development/AutomorphicArithmetic/Foundations/GaloisRepresentations/ClosedPhiFunctoriality/IsOp2One.lean}{45}{Integral operators and common denominator},
\src{src/Development/AutomorphicArithmetic/Foundations/GaloisRepresentations/ClosedPhiFunctoriality/IsOp2One.lean}{81}{comparison with Weil traces}, and
\src{src/Development/AutomorphicArithmetic/Foundations/GaloisRepresentations/ClosedPhiFunctoriality/ClosedPhiOfClosedTr.lean}{13}{passage through the integral comparison}.

In particular, Theorem~\ref{newgal:liu} supplies an interval for the
weights and the existence of the crystalline period space. It gives no
formula pairing $t$ with a specific operator on that space. A finer
published congruence theorem for odd $\ell$ gives equality of the labelled
Hodge types of two semistable representations of the same rank over one
$\mathcal O_L$ if their matrices are congruent modulo $\ell^n$, with
$n>c(K,H)$ for a common positive height bound $H$ after twisting.
Here $P_n$ can be rectangular and the approximation of $t$ uses a
weighted sum. Applying that theorem to $P_n$ would require an additional
equal-rank congruence. Propositions~\ref{newgal:hodge} and~\ref{newgal:phi}
specify the uniform comparisons used for this more general family. An
independent construction of the linked integral frames with the two
period identifications is still required for this proof in that generality.
T. Liu,
\emph{Filtration associated to torsion semi-stable representations},
Theorems 1.1 and 4.16~\cite{liuFiltration}.

In particular, at \(\ell=5\) the trivial and the unramified quadratic
character have the same single Hodge jump \(0\), while their Frobenius
traces are \(1\) and \(-1\). The local approximation at a spherical
\(w\mid\ell\) therefore uses the extra Hecke coordinates explicitly.
Take \(Y=\{w\}\) for \(w\) apart from the two auxiliary places. The
congruence theorem here holds on the disjoint union of the global
indices and the local integer powers. For a fixed Weil element \(g\)
and power \(j\), choose the coordinate with index \(-j\deg g\).
Its nonzero Hecke normalization is fixed before requesting the
precision. Divide it out of the approximated scalar and request the
correspondingly finer error in the original approximation. Genericity
at this \(w\) in every selected seed identifies its local power sum
with its Weil trace. This gives exactly the simultaneous global and local
test of Proposition~\ref{newgal:phi}.
\src{src/Development/AutomorphicArithmetic/Foundations/GaloisRepresentations/TraceLimits/LimFinOddOf.lean}{82}{Joint coordinate set}
and \src{src/Development/AutomorphicArithmetic/Foundations/GaloisRepresentations/TraceLimits/LimFinOddOf.lean}{185}{normalization and the same coefficients}.

The labelled comparison counts multiplicities. At power zero the
augmentation has error at most the prescribed error by testing the
constant polynomial \(1\) in \eqref{tr:approx}. At each fixed power
the same-level rectangular estimate in Proposition~\ref{newgal:hodge}
then recovers the full power sum, including repeated summands. Newton's
identities require no separation of jumps. For a dominant integral
weight in \eqref{newgal:jumps}, even equal adjacent highest-weight
entries give distinct shifted jumps, since
\((j+1-\lambda_{j+1})-(j-\lambda_j)\ge1\). Neither this observation nor
the multiset argument produces an automorphic seed at a singular
intermediate stage. The input there is the weight and local-type
congruence together with the geometric seed and the theta-lift and
descent steps in their stated ranges. For the odd-rank limit the
singleton construction compares weights only away from \(w_0,v'\).
\src{src/Development/AutomorphicArithmetic/Foundations/GaloisRepresentations/TraceLimits/LimFinOddOf.lean}{166}{Range of the singleton comparison}
and \src{src/Development/AutomorphicArithmetic/Foundations/GaloisRepresentations/IntegralStructures/Development.lean}{27}{Assembly from geometry, congruences and lifting}.

\subsection{The classical polarized cuspidal case}

There is a direct proof of the local assertions for the final classical
$\mathrm{GL}_a$ input which also explains the normalization. Let $\pi$
be a cuspidal automorphic representation of $\mathrm{GL}_a(\mathbb A_E)$
whose archimedean infinitesimal character is that of an irreducible
algebraic representation. Assume a continuous Hecke character $\chi$
of $E^+$ satisfies
\[
 \pi^c\simeq\pi^\vee\otimes(\chi\circ N_{E/E^+}\circ\det),
 \qquad \chi_v(-1)=(-1)^a\quad(v\mid\infty).
\]
Here $c$ is CM conjugation. The invariant bilinear pairing in the
polarization gives precisely this duality when the smooth finite and
archimedean modules are those of $\pi$. Choose the classical local
Langlands correspondence with uniformizers mapping to geometric Frobenius.
\src{src/Development/AbelianVarietyTheory/Definitions/GaloisGeometry/DualIsConjByOne.lean}{53}{Polarization}
and \src{src/Development/AutomorphicArithmetic/Foundations/GaloisRepresentations/TypeConditions/WeightsConjAlt.lean}{85}{cuspidal input}.

\begin{proposition}
For this input\label{newgal:classical} there is a continuous semisimple $r_\pi:G_E\to
\mathrm{GL}_a(k)$ with labelled jumps \eqref{newgal:jumps} at all
$w\mid\ell$. It is de Rham there and crystalline when $\pi_w$ is
spherical. Write
\[
 Y_w=\operatorname{rec}_{E_w}\bigl(\pi_w\otimes|\det|_w^{(1-a)/2}\bigr).
\]
At every $w\nmid\ell$ one has $\iota\mathrm{WD}(r_\pi|_{G_{E_w}})^{F\text{-}\mathrm{ss}}
\simeq Y_w$. At $w\mid\ell$ this equality also holds if the local
representation has an Iwahori-fixed vector or if $Y_w$ has zero monodromy.
Any continuous semisimple $r$ with the same normalized traces on a dense
set of good Frobenius elements inherits these conclusions.
\end{proposition}
\begin{proof}
Apply Barnet-Lamb, Gee, Geraghty and Taylor's Theorem~2.1.1 to $(\pi,\chi)$.
Their weight $\mu$ uses the infinitesimal character of the dual of the
algebraic module of highest weight $\mu$. Thus, with their index $1\le i\le a$,
$\mu_{\iota\tau,i}=-\lambda_{\iota\tau,a-i}$. The theorem gives jumps
$\mu_{\iota\tau,i}+a-i$, precisely \eqref{newgal:jumps}. It gives an
integer $w_\pi$ such that the jumps at $\tau c$ are the complements
$w_\pi-h$ of those at $\tau$. Parts (2) and (4) give the asserted
comparisons away from $\ell$ and at Iwahori level respectively, including
crystallinity at spherical level. Twisting by the indicated power of the
norm multiplies spherical Satake eigenvalues at geometric Frobenius by
$q_w^{(a-1)/2}$.
Barnet-Lamb,
Gee, Geraghty and Taylor, \emph{Potential automorphy and change of weight},
\S2.1, Theorem 2.1.1(2)--(4)~\cite{blggtPotential}.

For $a\ge2$ and $w\mid\ell$, the same authors' monodromy bound gives
the semisimplified Weil representation of $Y_w$. In each class of
irreducible Weil constituents up to unramified twist, the decreasing
list of lengths of the special representations on the Galois side is
dominated in partial sums by that on the automorphic side. If the
automorphic monodromy is zero, every length is one. The bound on the
first partial sum forces every Galois length to be one. Hence its
monodromy also vanishes, proving the claimed equality. For $a=1$
the same equality follows from the local compatibility of the
$\ell$-adic character of an algebraic Hecke character. Finally the
two global traces agree by continuity on $G_E$ and the characteristic-zero
Brauer--Nesbitt theorem identifies $r$ with $r_\pi$.
Barnet-Lamb,
Gee, Geraghty and Taylor, \emph{Local-global compatibility for $l=p$, II},
\S1, Theorems 1.1 and 3.1~\cite{blggtLocal}.
\end{proof}

For the classical correspondence these conclusions cover the ramified
places tested in the final statement. Above $\ell$ its extra local
comparison explicitly assumes zero monodromy on the automorphic side.
Away from $\ell$ the full comparison gives equality of all the tested
inertial polynomials and monodromy conditions. Applying this external
argument to intermediate definite-unitary eigensystems would require
their identification with suitable classical $\mathrm{GL}_a$ data and
equality of the specified Hecke systems. Applying it with an abstract
reciprocity map would also require the tested local parameters of that
map to agree with the classical ones. The weighted trace construction
and integral comparisons above describe the source's route prior to
these identifications.

\subsection{The principal level of an attached classical system}

Let $F$ be an imaginary quadratic field in rank three, or $\mathbb Q$ in
rank two, and let $M\ge1$. Write $\varepsilon_\ell$ for the cyclotomic
character. A character $\chi:G_F\to k^\times$ trivial on
$G_{F(\mu_M)}$ has finite image. For $v\mid p$ write $d=v_v(M)$,
where $v_v$ gives a uniformizer value one. Put
$U_v^0=\mathcal O_v^\times$ and $U_v^d=1+\mathfrak p_v^d$ for $d>0$.
Under local class field theory $\chi_v$ is trivial on $U_v^d$.
Indeed $\operatorname{Gal}(F(\mu_M)/F)$ is a subgroup of
$(\mathbb Z/M)^\times$. Extend its character to this finite abelian
group, using divisibility of $k^\times$. The local character on $F_v$
is then the norm pullback of a character on $\mathbb Q_p^\times$ whose
restriction to $\mathbb Z_p^\times$ factors modulo $p^{v_p(M)}$.
When $d>0$, any element of $U_v^d$ acts by multiplication on the
$\mathbb Z_p$-lattice $\mathcal O_v$ through a matrix congruent to
one modulo $p^{v_p(M)}$, since $\mathfrak p_v^d=M\mathcal O_v$.
Taking its determinant proves the claim. When $d=0$ the character
on $\mathbb Q_p^\times$ is unramified.

\begin{lemma}[Fixed vector criterion]
Let $\pi_v$\label{newgal:fixed} be a generic irreducible smooth representation of
$\mathrm{GL}_n(F_v)$, and write $(D,r_W,N)$ for its
Frobenius-semisimple Weil--Deligne parameter, allowing an unramified
normalizing twist. Suppose its inertia representation is a sum of
restrictions of characters of $W_{F_v}$ of conductor exponent at most
$d$. If $d=0$, assume also $N=0$. Then $\pi_v$ has a nonzero vector
fixed by the full principal congruence group $K_v(d)$, where
$K_v(0)=\mathrm{GL}_n(\mathcal O_v)$.
\end{lemma}
\begin{proof}
Write the generic representation as the irreducible parabolic induction
of essentially square-integrable representations $\sigma_j$ of ranks
$n_j$. Their parameters $D_j$ are the corresponding direct summands of
$D$. For $d>0$, set $t_j=\dim D_j^{I_v}$. The Artin conductor of its
Weil part satisfies $a(r_{W,j})\le (n_j-t_j)d$: finite inertia acts
semisimply, and every nontrivial constituent has the same restriction
to the ramification groups as one of the specified characters. The
additivity of the Artin conductor on this restriction gives the bound.
The Weil--Deligne conductor adds
$t_j-\dim(\ker N_j)^{I_v}\le t_j$. Thus
\[
 a(\sigma_j)=a(r_{W,j},N_j)\le (n_j-t_j)d+t_j\le n_jd .
\]
Here equality of the automorphic and Weil--Deligne conductors follows
from equality of their standard epsilon factors with the same additive
character, and an unramified twist preserves the exponent.
Miyauchi and Yamauchi's Theorem~1.2 characterizes $K_v(d)$-invariants
by exactly these inequalities for every inducing block. If $d=0$, all
the conductors vanish since inertia and monodromy are trivial. The
same theorem then gives a spherical vector.
Miyauchi and Yamauchi,
\emph{A remark on conductor, depth and principal congruence subgroups},
Theorems 1.1--1.2~\cite{miyauchiYamauchi}.
\end{proof}

\begin{proposition}[Transfer from the prescribed Galois level]\label{newgal:level}
Suppose a continuous semisimple $r$ is attached on a cofinite set of
places to a regular algebraic cuspidal $\pi$ over $F$. Assume $\pi$
is conjugate self-dual in the rank-three case, and comes from the
normalized classical weight $2e+1$ newform, $e\ge1$, in the rank-two
case. In rank three use $r$ after the inverse cyclotomic twist occurring
in Lemma~\ref{rk3:purity}. Suppose $r$ satisfies the following level
conditions for one prime $\ell$ and a tuple $m$.
\begin{enumerate}
\item At $v\nmid\ell$ the restriction to inertia is diagonal, with entries
$\chi_i|_{I_v}$, where each $\chi_i$ is trivial on $G_{F(\mu_M)}$.
\item At $v\mid\ell$ there is a $G_{F_v}$-stable full flag with graded
characters $\kappa_i$. On inertia $\kappa_i=\varepsilon_\ell^{m_i}\chi_i$,
with $\chi_i$ as above. At this $v$ either $r$ is potentially crystalline
or $\kappa_i\ne\kappa_j\varepsilon_\ell$ for all $i,j$.
\end{enumerate}
Then $\pi_f^{K_F(M)}\ne0$. At each embedding over $v\mid\ell$ its
attached representation has filtration jumps $\{-m_i\}$. In the two
counting applications $m=(e-1,-1,-e-1)$ and $m=(e,-e)$ respectively.
\end{proposition}
\begin{proof}
Chebotarev and semisimplicity identify $r$ with the representation
of $\pi$; this uses attachment to the classical cuspidal system as
a hypothesis. Every $\pi_v$ is generic by global cuspidality for
$\mathrm{GL}_n$. At $v\nmid\ell$ the prescribed finite inertia action
implies zero monodromy, since it is trivial on an open subgroup of
inertia, and gives precisely the specified inertia characters for the
Weil parameter. Full local--global compatibility away from $\ell$
therefore gives Lemma~\ref{newgal:fixed} at such $v$.

At $v\mid\ell$ put $L=F_v$. The graded characters have the form
$\kappa_i=u_i\chi_i\varepsilon_\ell^{m_i}$, with $u_i$ continuous
and unramified. Its value on geometric Frobenius is an $\ell$-adic unit
by compactness of $G_L$. Each graded line has the labelled jump
$-m_i$ at every embedding of $L$, and the attached representation is
de Rham. Strict exactness of filtered de Rham comparison now gives
the stated jumps. Let $q_v$ be the residue cardinality. The Weil
character of $\varepsilon_\ell$ is the unramified character $\nu$
taking geometric Frobenius to $q_v^{-1}$. On a finite character and on
a continuous unramified character the Weil functor gives restriction
to $W_L$, with coefficients transported by $\iota$.
Exactness of potentially semistable comparison gives the
semisimplified Weil part as $\bigoplus_i\mathrm{WD}(\kappa_i)$.
In particular its inertia is the sum of the $\chi_i|_{I_v}$.

We check zero monodromy when $d=0$. Potential crystallinity gives this
immediately in the first case. Otherwise, a nonzero monodromy matrix
entry between the $j$th and $i$th character spaces would give
$\mathrm{WD}(\kappa_i)=\nu\,\mathrm{WD}(\kappa_j)$ by
$r_W(g)Nr_W(g)^{-1}=\nu(g)N$. The character
$\theta=\kappa_i/(\kappa_j\varepsilon_\ell)$ would then have trivial
Weil character. Write $\theta=u\chi\varepsilon_\ell^s$ as above. On
geometric Frobenius, after returning to $k$, the Weil character takes
the value $u(\mathrm{Frob})\chi(\mathrm{Frob})q_v^{-s}$. Its
$\ell$-adic valuation forces $s=0$. The finite and unramified factors
now give $\theta|_{W_L}=1$, hence $\theta=1$ by density. This contradicts
the second alternative in the flag hypothesis.

Finally the needed full comparison at a possibly ramified
$v\mid\ell$ follows in rank three because every odd rank is
Shin-regular. In rank two the difference of the two dominant weight
entries is $2e-1>0$, as obtained by subtracting the half-sum of
positive roots from the normalized infinitesimal character. This is
the Shin-regular condition in rank two. Apply Theorem~1.1 for the CM
case and Corollary~1.2 for the totally real case in
Barnet-Lamb,
Gee, Geraghty and Taylor, \emph{Local-global compatibility for $l=p$, II}~\cite{blggtLocal}.
Thus Lemma~\ref{newgal:fixed} applies to the actual parameter on the
automorphic side, including at the coefficient prime. Taking the restricted
tensor product of the resulting fixed vectors gives a nonzero vector
at principal level $M$. Use $M=N$ in rank three and
$M=N|\operatorname{disc}F'|$ in rank two, where $F'$ is the imaginary
quadratic field to which the rank-two representation is subsequently
restricted. These are precisely the levels used in the fixed-field
count of Proposition~\ref{rk3:count}.
\end{proof}

The exported rank-three counting argument calls both
\src{src/Development/AutomorphicArithmetic/AutomorphicVarieties/ModularityArgument/ThresholdArguments/DenseThresholdCases/KeepLevel.lean}{95}{the conjugate self-dual CM level theorem}
and \src{src/Development/AutomorphicArithmetic/AutomorphicVarieties/ModularityArgument/ThresholdArguments/DenseThresholdCases/KeepLevel.lean}{83}{the polarized CM level theorem},
as recorded at the
\src{src/Development/AutomorphicArithmetic/AutomorphicVarieties/ModularityArgument/ThresholdArguments/DenseThresholdCases/KeepNumU.lean}{79}{count with the rank-two alternative}.
Both theorems use an additional local comparison witness and give
normalized good-place Frobenius polynomials as well as fixed vectors.
Proposition~\ref{newgal:level} establishes the fixed-vector and weight
conclusions for the stated classical representations. Its rank-two
calculation is over $\mathbb Q$ before restriction to the CM field.

\subsection{The polarized adjoint Selmer group over \texorpdfstring{$\mathbb Q$}{Q}}

Let $K$ be imaginary quadratic, $n\ge2$, $k=\overline{\mathbb Q}_\ell$,
and fix $\iota:k\simeq\mathbb C$. Write $\varepsilon$ for the $\ell$-adic
cyclotomic character, $\delta$ for the quadratic character of $K/\mathbb Q$,
and $c\in G_{\mathbb Q}$ for complex conjugation at a chosen embedding.
Given $s:G_K\to\mathrm{GL}_n(k)$, a multiplier $\mu$ on $G_K$ and
$P\in\mathrm{GL}_n(k)$ with
\[
 s(g)^t P s(cgc^{-1})=\mu(g)P ,
\]
let $A$ be a $G_{\mathbb Q}$-action on $M_n(k)$ which acts by conjugation
by $s(g)$ for $g\in G_K$ and by $X\mapsto-P^{-1}X^tP$ for $c$. In
particular its scalar line affords $\delta$.

\begin{proposition}[The two adjoint vanishing assertions]\label{newgal:adselmer}
Assume $s$ is continuous and absolutely irreducible on
$G_{K(\mu_{\ell^\infty})}$. Suppose one of the following holds, with
attachment meaning equality of normalized Frobenius polynomials at all
but finitely many places for a classical regular algebraic cuspidal
representation $\Pi$ of $\mathrm{GL}_n(\mathbb A_K)$.
\begin{enumerate}
\item $s$ is attached to a conjugate self-dual $\Pi$ and
$\mu=\varepsilon^{1-n}|_{G_K}$.
\item $s=r|_{G_K}$ for a continuous $r:G_{\mathbb Q}\to\mathrm{GL}_n(k)$
and $s$ is attached to $\Pi$. There is a nondegenerate pairing on $r$
with multiplier $\theta=\chi\varepsilon^w$, where $w\in\mathbb Z$ and
$\chi$ factors through $\operatorname{Gal}(\mathbb Q(\mu_M)/\mathbb Q)$
for some $M\ge1$. Here $\mu=\theta|_{G_K}$.
\end{enumerate}
Then $H^1_f(\mathbb Q,A)=0$ for every such $A$.
\end{proposition}
\begin{proof}
In the first case $(\Pi,\delta^n)$ is polarized, since $\delta$ is
trivial on norms from $K$ and $\delta_\infty(-1)=-1$. In the second,
$s^c\simeq s$ by $r(c)$ and $s\simeq s^\vee\otimes\mu$ by the pairing.
Let $\xi$ be the finite order Hecke character of $\mathbb Q$ with
$\ell$-adic realization $\chi$, using geometric reciprocity. Put
$\eta=\xi|\cdot|_{\mathbb A}^{w+n-1}$. Its $\ell$-adic realization is
$\theta\varepsilon^{n-1}$. If the normalized eigenvalues of $\Pi_v$
are $b_i=q_v^{(n-1)/2}\alpha_i$ at a good place, the dual has normalized
eigenvalues $q_v^{n-1}/b_i$. Twisting the dual by $\eta\circ N_{K/\mathbb Q}$
multiplies these by $\iota(\theta(\mathrm{Frob}_v))q_v^{1-n}$.
The result is the multiset for $\Pi^c_v$, by the two isomorphisms above.
Strong multiplicity one gives
$\Pi^c\simeq\Pi^\vee\otimes(\eta\circ N_{K/\mathbb Q}\circ\det)$.
Multiply $\eta$ by $\delta$ if necessary to arrange
$\eta_\infty(-1)=(-1)^n$. The twist on $\mathrm{GL}_n(\mathbb A_K)$
is unchanged, and the Galois multiplier
$\varepsilon^{1-n}\eta_\ell$ restricts to $\mu$ on $G_K$.

In both cases Chebotarev and Brauer--Nesbitt identify $s$ with
$r_{\Pi,\iota}$. Fixing the same $c$, the automorphic polarization gives
a pairing with multiplier $\mu$ and the negative adjoint action above
(Newton--Thorne,
\emph{Adjoint Selmer groups of automorphic Galois representations of unitary
type}, \S1~\cite{newtonThorne}). By Schur's lemma the pairing matrices differ by a scalar,
so they give the same action of $c$. Now apply
Thorne, \emph{On the vanishing
of adjoint Bloch--Kato Selmer groups of irreducible automorphic Galois
representations}, Theorem~6.1~\cite{thorneAdjoint}, with $F=K$ and $F^+=\mathbb Q$.
The hypothesis on the cyclotomic kernel is precisely its irreducibility
hypothesis. This theorem covers $\ell=2$ and imposes no ordinarity or
residual image condition. The earlier Newton--Thorne Theorem~5.2
assumes enormous image on the cyclotomic kernel.

To match the local conditions, a continuous cocycle $z$ and the action
$A$ descend simultaneously to a finite extension $L/\mathbb Q_\ell$:
apply Lemma~\ref{newgal:coeff} to the representation with diagonal
blocks $A(g),1$ and upper right block $z(g)$. Continuity of $A$ follows
from its restriction to the open subgroup $G_K$. At $v\nmid\ell$
the prescribed equation $z(g)=A(g)x-x$ on inertia is the unramified
condition. At $\ell$ the equation is $z(g)\otimes1=gx-x$ in
$A\otimes_{\mathbb Q_\ell}B_{\mathrm{cris}}$, the Bloch--Kato finite
condition. Indeed, for an embedding
$j:\overline{\mathbb Q}\hookrightarrow\overline{\mathbb Q}_\ell$,
the subgroup of compatible pairs of local and global automorphisms
projects isomorphically to $G_{\mathbb Q_\ell}$: restriction to the
algebraic closure of $\mathbb Q$ inside $\overline{\mathbb Q}_\ell$
determines the global component. Other choices conjugate the local
condition. In either equation an $L$-linear retraction $k\to L$
fixing $L$, applied coefficientwise and tensored with the identity
on the period ring when needed, gives a witness over $L$. Outside a
finite set containing $\ell$ and the ramification of $A$, inertia
acts trivially on $A$, so the local equation gives $z|_{I_v}=0$.
The extension representation thus factors through the Galois group
unramified outside this set. Hence $z$ represents a Selmer class with
coefficients in $L$, and vanishing gives a global coboundary.
\end{proof}

In rank three the cyclotomic twist cancels from the adjoint, and its
square supplies the multiplier $\varepsilon^{-2}$. The rank-two block
has the alternating determinant pairing and $w=0$. In each case the
cyclotomic-kernel irreducibility used in Proposition~\ref{newgal:adselmer}
is a separate input to the classical argument. The exported rank-two
proof uses its own cyclotomic irreducibility result, the cuspidal
attachment of the modular block, and additional local witnesses in
\src{src/Development/AutomorphicArithmetic/AutomorphicVarieties/ModularityArgument/ThresholdArguments/DenseThresholdCases/RRRigII.lean}{25}{the second rigidity case}.
The exported rank-four automorphy predicate supplies both a cuspidal
attachment over each imaginary quadratic field on which the restriction
is irreducible and a further comparison witness.
\src{src/Development/AutomorphicArithmetic/AutomorphicVarieties/ModularityArgument/AutomorphyLifting/IsAutCsdA.lean}{16}{Quadratic restriction and its witness}.
At a place above \(\ell\) where the ordinary flag characters
\(\kappa_i\) satisfy \(\kappa_i\ne\kappa_j\varepsilon\) for
all \(i,j\), this witness requires a flag for which
\(\kappa_i^m\ne(\kappa_j\varepsilon)^m\) for every \(m\ge1\).
It also requires, for every regular algebraic cuspidal attachment
\(\Pi\), a single real \(t\) such that at every place away from
\(\ell\) and every place above \(\ell\) satisfying the latter
condition, its cuspidal support consists of characters with
uniformizer values of absolute value strictly between
\(q_v^{t-1/2}\) and \(q_v^{t+1/2}\). Thus the first assertion
involves finite-order ratios too, and the second uses the same
\(t\) across the indicated places.
\src{src/Development/AutomorphicArithmetic/AutomorphicVarieties/ModularityArgument/AutomorphyLifting/WeakOfAtt.lean}{27}{The two requirements on the attachment}
and \src{src/Development/AutomorphicArithmetic/AutomorphicVarieties/ModularityArgument/LeviSmoothLocalization/WinUpDown.lean}{38}{The local support and its absolute-value bounds}.
If $J$ is the alternating similitude form, the pairing
matrix is $P=Jr(c)^{-1}$. Write $X^*=J^{-1}X^tJ$. The map
$X\mapsto (X-X^*)/2$ projects $A$ equivariantly onto the usual
conjugation action on $\mathfrak{sp}(J)$. Indeed conjugation by a
similitude commutes with $*$, and $c$ acts on $A$ by
$X\mapsto-r(c)X^*r(c)^{-1}$. Thus a finite symplectic tangent cocycle
which is a coboundary in $A$ has a primitive in $\mathfrak{sp}(J)$.
This division by two takes place in $k$. On the exported rank-four path,
\src{src/Development/AutomorphicArithmetic/AutomorphicVarieties/ModularityArgument/ThresholdArguments/DenseThresholdCases/ExistsCmFieldTwistH1.lean}{28}{descent from a linked CM tower}
produces an imaginary quadratic field and polarized adjoint vanishing
from an ordinary adjoint lifting theorem. The
\src{src/Development/AutomorphicArithmetic/AutomorphicVarieties/ModularityArgument/ThresholdArguments/DenseThresholdCases/RigSp4.lean}{181}{rank-four rigidity theorem}
combines this with the finite local conditions for the tangent cocycle
and the symplectic projection above. Proposition~\ref{newgal:adselmer}
identifies that vanishing for the classical points satisfying its
hypotheses. The exported lifting theorem concerns points of an ordinary
determinant ring and also uses their local flags, genericity and weight
conditions; establishing the hypotheses for those points is a separate
step.

\subsection{Exact local data from the elementary spherical Hecke operators}

At a spherical place \(w\mid\ell\) apart from \(w_0,v'\), the
simultaneous Hecke congruence used in the odd-rank construction accepts
any collection of integral spherical operators away from \(\ell\).
Using the elementary symmetric operators in that collection gives a
second proof of the two exact local comparisons. The geometric seed
representations and the simultaneous congruence are the inputs to this
argument.

Write \(k=\overline{\mathbb Q}_\ell\) and \(\mathcal O_k=\{z\in k:|z|\le1\}\).
For a representation \(V\) of \(G_K\), with \(K/\mathbb Q_\ell\) finite,
write \(X(V)_{g,j}\) for the coefficient of \(t^j\) in
\(\det(1+tV(g))\), for \(g\in G_K\) and \(1\le j\le a=\dim V\).
These coefficients are integral on continuous representations, by a stable
lattice over a finite coefficient field as in Lemma~\ref{newgal:coeff}.

\begin{lemma}[Polynomial tests at a residual determinant]\label{addg:residual}
Let \((X_0,z_0)\) and \((X_i,z_i)\) consist of such coefficient systems
and an additional scalar in \(k\). Suppose for each \(\eta>0\) there
are finitely many \(i\) and coefficients \(c_i\in k\) such that
\[
 \left|P(X_0,z_0)-\sum_i c_iP(X_i,z_i)\right|\le\eta
 \qquad(P\text{ an integer polynomial in the coordinates}).
\]
Fix a finite coefficient DVR \(\mathcal O\subset k\) over which the
residual determinant \(\bar D\) of \(X_0\) is defined. Let \(R\) be
its universal local pseudodeformation ring over \(\mathcal O\) for
\(G_K\). For each \(\eta\), one can retain only the points with residual
determinant \(\bar D\), and choose weights \(d_i\) on these points,
such that their evaluation maps \(x_i\) and the target map \(x_0\) satisfy
\[
  \left|x_0(f)-\sum_i d_i x_i(f)\right|\le\eta\quad(f\in R),
  \qquad \left|z_0-\sum_i d_i z_i\right|\le\eta .
\]
\end{lemma}
\begin{proof}
The original tests hold for polynomials with coefficients in
\(\mathcal O_k\) at the same precision. Apply the integer test to each
monomial, multiply by its coefficient, and use the ultrametric inequality.
Choose a support and weights at precision \(\min(\eta,1/2)\).
For each point \(i\) with a
different residual determinant, some coordinate \(j_i\) differs from the
target modulo the maximal ideal of \(\mathcal O_k\). Indeed the
characteristic coefficients on group elements determine a determinant on
the group algebra by Amitsur's formula. Form the integral polynomial
\[
 H(X)=\prod_{i\,\mathrm{excluded}}
       \frac{X_{j_i}-(X_i)_{j_i}}{(X_0)_{j_i}-(X_i)_{j_i}} .
\]
The denominators are units. Testing \(HP\) gives the same inequality
for \(P\) with \(d_i=c_iH(X_i)\) on the retained points, since
\(H(X_0)=1\) and all excluded evaluations vanish. The constant test
ensures that some point remains.

The universal ring is topologically generated over \(\mathcal O\) by
the characteristic coefficients of group elements. The absolute Galois
group of a finite extension of \(\mathbb Q_\ell\) satisfies Mazur's
finiteness condition, so this ring is Noetherian and its topology is
the maximal-ideal topology. See Chenevier, Corollary 1.14,
Proposition 3.3, Remark 3.5, Example 3.6 and Proposition 3.7
of~\cite{chenevier}, with finite extension of the coefficient DVR.
The evaluations are continuous local maps into integer rings of finite
extensions. For an extension with larger residue field, the coefficient
values and \(\mathcal O\) lie in the inverse image of the chosen residue
field, so the same universal property applies. Approximate any \(f\in R\)
by integral coordinate polynomials in the maximal-ideal topology and
evaluate at the target and the retained points. There are finitely many
maps and the \(d_i\) are already fixed. Thus the precision at these maps
can be chosen to absorb \(\max(1,\max_i|d_i|)\). Passing to the limit
gives the first inequality. Testing the additional coordinate itself
gives the second.
\end{proof}

\begin{proposition}[Exact periods on the odd-rank spherical subfamily]
\label{addg:local}
Use the geometric seed and simultaneous congruence hypotheses of the
odd-rank construction in Sections~\ref{tr:section} and~\ref{newgal:section},
at a spherical \(w\mid\ell\) apart from both auxiliary places. For the
global representation \(r\) constructed from traces, \(V=r|_{G_{E_w}}\)
has the labelled jumps \(J_\tau\) of \eqref{newgal:jumps}. Its
Frobenius-semisimplified crystalline Weil parameter has the prescribed
Satake eigenvalues in the normalization of Proposition~\ref{newgal:phi}.
These conclusions hold for every coefficient prime \(\ell\).
\end{proposition}
\begin{proof}
Set \(K=E_w\). At a usable spherical prime \(v\nmid\ell\), put
\(q=\mathrm N v\) and let \(T_j\) be the hyperspecial double-coset
operator for \(j\) uniformizer entries. On spherical fixed vectors,
\(q^{\binom j2}T_j\) has eigenvalue the \(j\)-th elementary symmetric
function of \(q^{(a-1)/2}B_v\). The eigenvalue of \(T_j\) is
\(q^{j(a-j)/2}e_j(B_v)\), and the exponents
add to \(j(a-1)/2\). The representatives and coefficients are integral
at \(\ell\). A representation with zero fixed space satisfies the required
scalar-action assertion there as well.
\src{src/Development/AbelianVarietyTheory/Proofs/SatakeParameters/SphericalOperatorsExist.lean}{146}{the elementary spherical operators and their action}
and \src{src/Development/AbelianVarietyTheory/Statements/EllipticCurveLevels/HasLevel.lean}{36}{the admissible operator data}.

For a fixed Weil element \(s\) and power, let \(b_0\) denote the
predicted Weil trace and \(b_i\) its value at a seed. The congruence for
the singleton \(\{w\}\) includes the coordinate \(z=cb\), with one
nonzero scalar \(c\) fixed by the place, power and normalization. Take
as the away-from-\(\ell\) list all the preceding elementary operators on
usable spherical primes, indexed also by the prolongations and Frobenius
lifts as in Section~\ref{tr:section}. The congruence permits this list
and the local coordinate together, with the same coefficients for every
polynomial.
\src{src/Development/AbelianVarietyTheory/Definitions/GaloisGeometry/LimWU4OfLimWU3.lean}{46}{quantification over operator lists and the additional local coordinates}
and \src{src/Development/AutomorphicArithmetic/Foundations/GaloisRepresentations/TraceLimits/LimFinOddOf.lean}{185}{normalization of the Weil power}.

On a fixed finite support of this approximation the common good
geometric Frobenius elements are dense in the global Galois group.
The target identification away from \(\ell\) in Section~\ref{tr:section}
and the seed characteristic-polynomial identities identify all the
elementary coefficients at those elements. For a test in finitely many
group elements, approach each element by common good Frobenius elements,
keeping all its elementary degrees at the same prime. Hold \(z\) fixed.
Continuity of the representations passes the polynomial inequality to
the chosen elements, with the same finite support. Restriction to
\(G_K\) gives exactly the hypothesis of Lemma~\ref{addg:residual} for
\(X_0=X(V)\), \(X_i=X(V_i)\) and \(z_j=cb_j\). Taking the elementary
operators directly avoids division by factorials at the small primes.

Choose the fraction field of \(\mathcal O\) to contain \(c\) and the
images of all embeddings of \(K\) over \(\mathbb Q_\ell\).
Theorem~\ref{newgal:liu} first gives crystallinity of the actual
restriction \(V\), including its extension classes. All the local seeds
are crystalline of the common labelled type. In the ring \(R\) of the
lemma, take the semistable quotient \(R^{\mathbf v,1}\) with this Hodge
type and trivial inertial type, taking the semistability field to be
\(K\). The quotient and its characterization on field-valued points for
arbitrary residual multiplicities are given by A'Campo, Hevesi, Thorne
and Whitmore, Theorem 3.3.6 in
\href{https://arxiv.org/abs/2607.11763}{\emph{Local-global compatibility of automorphic Galois representations over CM fields at \(p\)}}.
The local pseudorepresentation records the semisimplification. Exactness
of crystalline periods and strictness of the de Rham filtration on
crystalline representations give the same labelled graded multiplicities
for a seed and its semisimplification, as also recorded for de Rham
representations in Lemma 3.3.19 of the same paper. In particular every retained
seed map kills the defining ideal. Apply Lemma~\ref{addg:residual} to
each element of this ideal and let \(\eta\) tend to zero. The target
factors through \(R^{\mathbf v,1}\). Applying the same exactness to the
already crystalline \(V\) gives its claimed labelled jumps.

The integers in each \(J_\tau\) are distinct by dominance, so this
type fits the highest-weight convention of \S3.3.9 of the same paper.
Use the spherical Bernstein block. By \S3.3.11 the invariant function
given by the trace at the prescribed Weil element and power is an element
\(F\in R^{\mathbf v,1}[1/\ell]\). Its specialization is the trace of
the semisimple Weil representation. This trace also agrees with the
trace on the original crystalline module by exactness. Choose \(h\ge0\)
and \(f\in R\) lifting \(\ell^h cF\). The choice is fixed
before \(\eta\). Each retained point has \(x_i(f)=\ell^h z_i\). The
two inequalities of the lemma then give
\[
       |x_0(f)-\ell^h z_0|\le\max(1,|\ell^h|)\eta .
\]
Letting \(\eta\) decrease to zero and cancelling \(\ell^h c\) gives
the predicted Weil trace on \(V\). Apply this to the first \(a\)
powers of geometric Frobenius and use Newton's identities in
characteristic zero. Crystallinity gives trivial inertia and zero
monodromy. This determines the asserted parameter.
\end{proof}

\subsection{Base change for the definite-unitary eigensystems}

\begin{proposition}[Attachment on the standard definite unitary group]
\label{addg:unitary}
Use the CM field, odd rank \(a\) and definite unitary group of
Section~\ref{newgal:section}, with the classical local Langlands
correspondence and Fontaine period functors. For each of its
automorphic representations \(U\) there is a continuous semisimple
\(r_U:G_E\to\mathrm{GL}_a(k)\). At every split spherical place
\(w\nmid\ell\) it is unramified with geometric Frobenius multiset
\(q_w^{(a-1)/2}B_w\). At every \(w\mid\ell\) its restriction is de Rham
with labelled jumps \(J_\tau\) of \eqref{newgal:jumps}. If \(U\) is
spherical there, it is crystalline and its crystalline Weil parameter
has geometric Frobenius multiset \(q_w^{(a-1)/2}B_w\).
This holds for all dominant integral weights and for any finite
exceptional set \(T\) of nonsplit places in the level.
\end{proposition}
\begin{proof}
The group in the construction is defined by \({}^{t}\!\bar g g=1\)
for the identity Hermitian form over \(E/E^+\). The field hypotheses
give an extension unramified at every finite place with every place
over \(\ell\) split. These are the standard assumptions of Newton and
Thorne, \S1, and their formula (1.0.1) is the displayed group.
Their map at a split place is projection to the chosen matrix factor,
which is the map defining the local component in \(B_w\). At infinity
the pullback of the algebraic module along the chosen embedding gives
the weight \(\lambda_{\iota\tau}\) used here. On the compact adelic
quotient the continuous automorphic functions of fixed finite level
and finite type at infinity are square integrable. Their finite-dimensional
translates at infinity are smooth by the representation theory of compact
Lie groups. An irreducible submodule of these automorphic forms therefore
occurs in the discrete automorphic spectrum. The projection to some
irreducible Hilbert summand is nonzero and maps this submodule into
its smooth, finite-type vectors.
\src{src/Development/AbelianVarietyTheory/Definitions/GaloisGeometry/LatticeDataThicknessBound.lean}{16}{field, automorphic and split local conditions}.
Theorem 1.2 and Corollary 1.3 of Newton and Thorne give the discrete
isobaric base change \(\Pi\) and \(r_U\), including compatibility
with \(\operatorname{rec}_{E_w}(\Pi_w|\det|^{(1-a)/2})\) at split places
and de Rhamness over \(\ell\). Their forward base change allows
ramification at nonsplit places. At a spherical place the parameter has
trivial inertia and zero monodromy. For a de Rham representation these
two conditions on its Weil--Deligne parameter give crystallinity by
semistable descent. The displayed norm twist gives the claimed factor
of \(q_w^{(a-1)/2}\).
Newton and Thorne,
\href{https://arxiv.org/abs/1912.11261}{\emph{Symmetric power functoriality for holomorphic modular forms}}, \S1.

For the labelled weights, let one discrete block of \(\Pi\) have
cuspidal constituent \(\mu\) on \(\mathrm{GL}_m\) and length \(b\).
Its archimedean constituents have twists \(\mu|\det|^{(b+1-2j)/2}\),
\(1\le j\le b\). In the construction of \(r_U\) this block contributes
\[
   \bigoplus_{j=1}^{b}
   r_{\mu|\det|^s}\otimes\varepsilon_\ell^{\,1-j},
   \qquad s=(m+b-1-a)/2.
\]
Here \(\mu|\det|^s\) is regular algebraic and polarizable. This is the
discrete-spectrum formula used in the proof of Newton and Thorne's
corollary. See Anastassiades and Thorne,
\href{https://arxiv.org/abs/1912.11267}{\emph{Raising the level of automorphic representations of
\(\mathrm{GL}_{2n}\) of unitary type}}, proof of Corollary 3.4.
For each labelled infinitesimal-character entry \(x\) of \(\mu\),
the cuspidal weight formula in Proposition~\ref{newgal:classical} and
the jump \(-1\) of \(\varepsilon_\ell\) give the entry
\[
  (m-1)/2-(x+s)+j-1
      =(a-1)/2-\bigl(x+(b+1-2j)/2\bigr).
\]
Taking all entries and blocks gives \((a-1)/2\) minus the infinitesimal
character of \(\Pi\), with multiplicity. By Theorem 1.2(3) this character
is that of the algebraic module of \(U\). Its entries in zero-based
indexing are \(\lambda_{\iota\tau,i}+(a-1)/2-i\). Subtraction gives
\(i-\lambda_{\iota\tau,i}\) for \(0\le i<a\), including equal adjacent
highest-weight entries.
\end{proof}

For example, the constant automorphic representation in odd rank
\(a\ge3\) has highest weight zero. Its algebraically normalized Satake multiset at a
split spherical place is \(\{1,q_w,\ldots,q_w^{a-1}\}\), so the attached
representation is \(\bigoplus_{i=0}^{a-1}\varepsilon_\ell^{-i}\) by
Chebotarev and semisimplicity. Its jumps \(0,\ldots,a-1\) fail
lacunarity since \(\{1\}\) and \(\{0,2\}\) have equal means.
For a general \(U\) the same uniqueness identifies \(r_U\) with any
semisimple trace limit matching its good Frobenius eigenvalues.
Proposition~\ref{addg:unitary} thus supplies the rank-\(a\) seed
and target Galois conclusions in the classical interpretation.
Realizing a particular seed in the prescribed cohomology tower,
proving the simultaneous automorphic congruence, and comparing any
separately specified period functors with Fontaine's are further assertions.

\subsection{Inserting the auxiliary cuspidal type at fixed weight}

\begin{proposition}\label{addg:auxiliary}
Fix any algebraic weight of the definite group and a level away from
the split auxiliary place \(w_0\mid\ell\). Let \(f\) be a simultaneous
eigenform of that weight and level, with system \(\theta\).
Take any list of integral spherical Hecke operators away from
\(\ell\) at the hyperspecial places of the level, or more generally
integral operators away from \(\ell\) acting by a scalar on each
relevant local fixed space. One may also include
fixed spherical double-coset operators or finite linear combinations
at hyperspecial places over \(\ell\) apart from \(w_0\), after a nonzero scalar
normalization for each operator. For every \(m\ge1\) there are finitely
many eigensystems \(\theta_i\) of the same weight and level away from
\(w_0\), all containing the auxiliary type of Section~\ref{newgal:section}
at \(w_0\), and coefficients \(c_i\in k\) such that
\[
   \left|P(\theta)-\sum_i c_iP(\theta_i)\right|
              \le |\ell|^m
      \qquad(P\text{ an integer polynomial in the listed operators}).
\]
\end{proposition}
\begin{proof}
In the finite-adelic model let \(A\) be the algebraic coefficient
representation, with \(f(\gamma g)=A(\gamma)f(g)\). Transport coefficients
by \(\iota^{-1}\), and measure the values of \(f\) using
\(A(g_\ell)^{-1}f(g)\). These values are left invariant. A lattice
preserved by the compact coefficient-place part of the original level
exists over a finite coefficient field, after shrinking the level at
\(w_0\) if needed. The resulting supremum norm is a maximum on the
finite double quotient and uses the same lattice at finer levels.
Integral Hecke operators away from \(\ell\) are contractive since their
finite sums of translations have trivial \(\ell\)-component. At each
marked coefficient place there are finitely many matrices \(A(t_\ell)\)
in a given operator. Multiplying its coefficients by a sufficiently
large power of \(\ell\) makes all the summands integral in the lattice.
These normalizations depend only on the operators and weight.
Rescale \(f\) so a coordinate evaluation \(L\) of norm at most one
satisfies \(L(f)=\|f\|=1\).

Inside a torsion-free compact open subgroup fixing \(f\) at \(w_0\),
choose an open \(J\) and a quotient
\(\theta_J:J\twoheadrightarrow C=\mathbb Z/\ell^m\mathbb Z\) such
that every nontrivial character of \(C\) inflated to \(J\) is a
supercuspidal type. Such subgroups exist for every prime by
Beuzart-Plessis, Theorem D.1 in the appendix to Fintzen and Shin,
\href{https://arxiv.org/abs/2009.08476}{\emph{Congruences of algebraic automorphic forms and supercuspidal representations}}.
For \(m=1\) one may compose a quotient of order \(\ell^2\)
with reduction modulo \(\ell\).
A sufficiently deep principal congruence subgroup is torsion-free by
the convergent matrix logarithm. For \(a>1\), surjectivity follows also
because a nontrivial character trivial on the image would occur in
the trivial representation. In rank one a coordinate of the logarithm
on the deep unit group gives a surjection to \(C\), and every smooth
irreducible representation is cuspidal. In level \(\ker\theta_J\), the right
action of \(C\) on the finite double quotient is free. Indeed the
rational stabilizer at level \(J\) is finite by compactness at infinity
and discreteness of the rational group in its adeles. Its projection
at \(w_0\), after conjugation, lies in the torsion-free group \(J\).
The faithful embedding of rational matrices into this completion
then makes the stabilizer trivial.

Choose one point of each \(C\)-orbit and let \(h\) be 1 minus
\(\ell^m\) times the characteristic function of the selected points.
Its sum on each orbit is zero and \(\|1-h\|=|\ell|^m\).
Since \(f\) is \(J\)-fixed, \(hf\) has only nontrivial
\(C\)-characters. Decomposing at the fixed weight and the finer
level gives a finite sum \(hf=\sum_i f_i\) of Hecke eigenvectors
in these character spaces. The decomposition follows from the
discrete automorphic decomposition and the scalar action of the
spherical operators on hyperspecial invariants. Each occurring
local representation at \(w_0\) contains the indicated type.
Every integral polynomial in the operators is contractive, so applying
\(L P(T)\) to \(f-hf\) proves the bound with \(c_i=L(f_i)\).
\end{proof}

This fixed-weight argument preserves any prescribed finite collection
of polynomial nonvanishing conditions already satisfied by the target.
Integral coefficients in \(k\) are allowed by the monomial argument of
Lemma~\ref{addg:residual}. After clearing denominators, let \(D\) be the product
of their integral test polynomials. Apply the estimate to \(DP\)
and divide by \(D(\theta)\ne0\). Terms with \(D(\theta_i)=0\)
drop out, and increasing \(m\) in advance absorbs the fixed loss
\(|D(\theta)|^{-1}\). For a polynomial required to be nonzero at
every ordering of a Satake multiset, take the product over permutations and express this
symmetric polynomial through the elementary spherical operators.
This proves the auxiliary-type step at arbitrary fixed weights,
including repeated highest-weight entries. Changing the weight to
produce lacunarity at \(v'\), or imposing a nonvanishing condition
that fails on the target, requires an additional argument.

\section{Supplying the arithmetic family}\label{sup:section}

We first select two primes for which every auxiliary twist in a fixed finite family is small. Theorem~\ref{pr:pair} proves the required uniform statement, using the density estimate below. The cohomological selection then requires an interval on which the strict Selmer groups have bounded total length for those very primes. Theorem~\ref{adds:dyadic-window} combines the low-layer and tower estimates on one interval for the whole field contribution at two and the prescribed fixed odd layers. Proposition~\ref{adds:order} gives three ranges for the character sums. The other varying-character inputs are described in Sections~\ref{famsel:complement} and~\ref{rk3:sec}.

\subsection{A density bound in the auxiliary box}

Write \(\sigma_0=399/400\). The object to be counted is a primitive character of conductor at most \(Q\) with at least one zero in \(\Re s>\sigma_0,\ |\Im s|\le Q\). We give the elementary mechanism of the density estimate, including the source of the large-sieve exponent. The classical analytic facts used here are nonvanishing for \(\Re s\ge1\) for nonprincipal characters, their Dirichlet series in \(\Re s>0\) by partial summation, and the divisor bound \(\tau(n)\ll_\alpha n^\alpha\). \src{src/Formalism/GL3MeanSquares/GenusMeanAsymptotics/JacobiCharacterSums/CountZAll.lean}{13}{density assembly}

\begin{lemma}[Auxiliary-box density]\label{sup:density}
For \(Q\ge3\), the number of such primitive characters is \(O(Q^{1/4})\), with an absolute implied constant.
\end{lemma}
\begin{proof}
It suffices to work with large \(Q\). The smaller range has only finitely many possible conductors. Put \(B=\lceil Q\rceil\), \(V=B^4\), \(J=8(\lfloor\log_2B\rfloor+1)\), and \(N=2^J\). Ignore conductor \(1\), which contributes at most one character. Choose one zero \(\rho_\chi=\beta+it\) for each remaining character. For the partial sums \(A_\chi(u)=\sum_{n\le u}\chi(n)\), periodicity and vanishing of the sum over a period give \(|A_\chi(u)|\le Q\). Partial summation at a zero gives
\[
 \left|\sum_{n\le N}\chi(n)n^{-\rho_\chi}\right|
 \le Q N^{-\beta}(1+|\rho_\chi|/\beta).
\]
The mollifier \(\sum_{a\le V}\mu(a)\chi(a)a^{-\rho_\chi}\) has modulus at most \(V\). Since \(N\ge B^8\) and \(\beta>399/400\), the product is \(O(B^{6-8\sigma_0})\) and thus at most \(1/2\) for large \(Q\). Multiplicativity expands this product into \(1+D\), where \(D\) is supported on \(V<n\le VN\), with coefficients
\[
 c_n=\sum_{\substack{ab=n\\a\le V,\ b\le N}}\mu(a).
\]
Indeed \(N\ge V\), so for \(1<n\le V\) all divisors are present and the coefficient vanishes. Consequently \(|D|\ge1/2\). Partition \((V,VN]\) into \(J\) intervals \((H,2H]\) with \(H=V2^j\), and assign the character to a block whose sum has modulus at least \(1/(2J)\).

For a set \(T\) assigned to such a block put \(b_n=c_n n^{-\sigma_0}\) and \(v_\chi(n)=\chi(n)n^{-(\rho_\chi-\sigma_0)}\). The diagonal Gram entries of the vectors \(v_\chi\) are at most \(H\). For distinct primitive characters \(\chi,\chi'\), their product with one conjugated, viewed at the common product modulus, is nonprincipal. Otherwise uniqueness of the primitive inducer would identify them. Its sums over intervals have absolute value at most \(2Q^2\). Let \(w=\rho_\chi+\overline{\rho_{\chi'}}-2\sigma_0\), so \(\Re w\ge0\) and \(|w|\le2Q+2\). Partial summation on \((H,2H]\) bounds the corresponding Gram entry by
\[
 2Q^2\left((2H)^{-\Re w}+|w|\int_H^{2H}u^{-1-\Re w}\,du\right)\ll Q^3 .
\]
This shows why one needs to count characters with a chosen zero, rather than impose separation among the heights of zeros for one character. The row-sum inequality for the Hermitian Gram matrix follows from \(2|z_i||z_j|\le |z_i|^2+|z_j|^2\). By duality,
\[
 \frac{|T|}{4J^2}\le \sum_{\chi\in T}\left|\sum_n b_n v_\chi(n)\right|^2
 \le\bigl(H+C|T|Q^3\bigr)\sum_{H<n\le2H}|b_n|^2 .
\]
Use \(|c_n|\le\tau(n)\ll n^{1/200}\). Then \(\sum|b_n|^2\ll H^{1-2\sigma_0+1/100}=H^{-197/200}\). Uniformly in all blocks,
\(J^2 Q^3 H^{-197/200}\ll(\log Q)^2 Q^{-47/50}\to0\).
Thus the term proportional to \(|T|\) can be absorbed, leaving \(|T|\ll J^2H^{3/200}\). Summing over the blocks and using \(VN\ll Q^{12}\) gives \(O(J^3Q^{9/50})=O(Q^{1/4})\).
\end{proof}
\noindent
The derivation uses the same block and exponents as the source, but the proof above uses an absolute \(O(Q^3)\) for the off-diagonal. The source gives the explicit row bound \(H+8|T|Q^3\) and handles the small \(Q\) range separately. \src{src/Formalism/GL3MeanSquares/GenusMeanAsymptotics/JacobiCharacterSums/CountZDefs.lean}{44}{coefficient exponent and block lengths}\quad
\src{src/Formalism/GL3MeanSquares/GenusMeanAsymptotics/JacobiCharacterSums/CountZDetect.lean}{73}{detection from a selected zero}\quad
\src{src/Formalism/GL3MeanSquares/GenusMeanAsymptotics/JacobiCharacterSums/CountZBlock.lean}{23}{explicit Gram bound}\quad
\src{src/Formalism/GL3MeanSquares/GenusMeanAsymptotics/JacobiCharacterSums/CountZFinal.lean}{10}{density conclusion}.

\subsection{Selecting auxiliary primes and retaining the same witnesses}

The common selection lemma takes fixed \(0<\varepsilon\le1/4,\tau>0,K\ge1,\eta>0\) and a fixed upper bound \(R\) on the number of base characters. At sufficiently large \(Y\), it applies uniformly to
\[
 \begin{gathered}\mathcal P=\{p:\sqrt U<p\le U,\ p\equiv a_0\bmod m_0,\ p\text{ prime}\},
 \\ 1\le F\le\sqrt U,\qquad Q=FU,\qquad Z=(Y^K)^{400}.\end{gathered}
\]
The base characters are primitive of conductor \(\le F\). The progression is reduced, \(8\mid m_0\), and
\[
 \tfrac14(\log Y)^2\log m_0\le\log U,\quad
 Y^{\varepsilon/6}\le\log U\le Y^{\varepsilon/5}(\log Y)^2 .
\]
Primes dividing \(m_0\) are \(\le Y^{\varepsilon/5}\) or \(>Z\). The proved relative prime number theorem in progressions, with its possible exceptional term and the source's exceptional-zero lower bound, gives \(|\mathcal P|\ge U^{9/10}\). This uses a relative estimate in the specified range, not a uniform assertion of an exceptional-zero-free progression formula. \src{src/Development/AutomorphicArithmetic/MeanSquareBounds/StagePow.lean}{41}{selection hypotheses}\quad
\src{src/Formalism/GL3MeanSquares/GenusMeanAsymptotics/JacobiCharacterSums/MSqCount.lean}{346}{candidate count}\quad
\src{src/Formalism/GL3MeanSquares/GenusMeanAsymptotics/JacobiCharacterSums/PNTFinal.lean}{22}{proved relative progression input}.

For each fixed base character of conductor \(f\le F\), its twist by the quadratic character of \(p>\sqrt U\) is primitive of conductor \(fp\). Even if one chooses a different bad twist for each bad prime, the conductor uniquely recovers \(p\). Indeed the relation \(fp=f'p'\), with \(f'\le F<p\), forces \(p=p'\). Lemma~\ref{sup:density} thus gives \(O(Q^{1/4})=O(U^{3/8})\) first-kind bad primes, simultaneously for the full base list. Their fraction tends to zero as \(O(U^{-21/40})\). \src{src/Development/AutomorphicArithmetic/MeanSquareBounds/CountBadCore.lean}{18}{injective conductor argument}

For the second kind of bad prime, let
\[
 S_p(y)=\sum_{n\le Z}\Lambda(n)\min(1,y/n)\kappa(n)(n/p),\qquad |\kappa(n)|\le1.
\]
The source proves \(\sum_{p\in\mathcal P}|S_p(y)|^2\le C_1 y\log y\,|\mathcal P|\) on \([Y,Y^K]\). To obtain this, expand the square, use reciprocity at \(p>Z^2\) in the fixed progression, and apply the relative progression estimate to the odd squarefree kernel of \(nn'\) after removing primes of \(m_0\). Each nontrivial kernel gives cancellation of size \(Y^{-K-1}|\mathcal P|\), except at most one from the exceptional character. For the diagonal and the exception, the prime-power support of \(\Lambda\) allows a separate bound. The grouped coefficients by prime base have square sum \(O(y\log y)\) and maximum \(O(\log y)\), and the coefficients whose residual odd squarefree kernel is one have total mass \(O(\sqrt y+Y^{\varepsilon/5}\log y)\). The explicit decomposition therefore bounds both the diagonal and the single exception by \(O(y\log y)\) after normalizing by \(|\mathcal P|\). \src{src/Development/AutomorphicArithmetic/MeanSquareBounds/twistK.lean}{148}{single exceptional reciprocity sum}\quad
\src{src/Development/AutomorphicArithmetic/MeanSquareBounds/MeanSqAbs.lean}{8}{diagonal and exceptional decomposition}\quad
\src{src/Development/AutomorphicArithmetic/MeanSquareBounds/meanSqK.lean}{51}{pointwise mean square}.

To pass to all \(y\) for a single prime, differentiate the piecewise linear ramp away from its finitely many breakpoints. The derivative is the sum over \(n>y\) with coefficients \(\Lambda(n)\kappa(n)(n/p)/n\). Apply the same mean square to \(\kappa(n)\mathbf1_{n>y}\) to bound the sum of squared derivatives by \(C_1|\mathcal P|\log y/y\). For each dyadic block \([a,b]\), \(b\le2a\), and arbitrary \(w_p\in[a,b]\), the fundamental theorem of calculus and Cauchy--Schwarz give
\[
 \sum_p |S_p(w_p)|^2
 \le 2\sum_p|S_p(a)|^2+2a\int_a^b\sum_p|S_p'(v)|^2\,dv
 \le 4C_1a\log(2a)|\mathcal P|.
\]
Choose a witnessing bad point \(w_p\) when \(|S_p(y)|>y^{1/2+\tau}/2\) somewhere in the block. Chebyshev's inequality, the bound \(\log(2a)\ll_\tau a^\tau\), and a geometric sum give a bad proportion \(O_\tau(Y^{-\tau})\) for one \(\kappa\). Budget at most \(\eta/(2(R+1))\) for each base character and \(\eta/2\) for the zero boxes. This proves simultaneous goodness across the base list and the entire interval. The source supplies the uniform constant \(32C_1/(\tau(1-2^{-\tau}))\) in the bound for the number of bad primes divided by \(Y^{-\tau}|\mathcal P|\). \src{src/Development/AutomorphicArithmetic/MeanSquareBounds/BadShare.lean}{146}{derivative bound}\quad
\src{src/Development/AutomorphicArithmetic/MeanSquareBounds/BadShare.lean}{251}{uniform bad proportion}\quad
\src{src/Development/AutomorphicArithmetic/MeanSquareBounds/StagePow.lean}{141}{combining the two bad sets}.

For a good auxiliary twist at this stage, the source uses the exact splitting \(G_\chi(y)=-S_p(y)-y\,\mathrm{tailL}_\chi(Z)\). It proves \(Y^K|\mathrm{tailL}_\chi(Z)|\le Y^\varepsilon\) by using the finite-height zero box to move a contour. Together with the half-sized threshold for \(S_p\), this gives \(|G_\chi(y)|\le y^{1/2+\tau}\). A separate short truncation and Mertens bound yield \(1000\log\log Q+O(1)\) for the logarithmic derivative. In the final application \(\log Q\le3Y^{\varepsilon/5}(\log Y)^2\). Since \(200\varepsilon<50/13<4\), this gives \(\left|L'/L(1,\chi)\right|\le4\log Y\le4\log y\) eventually. \src{src/Development/AutomorphicArithmetic/MeanSquareBounds/BoxOfS.lean}{28}{from the truncated ramp to the box}\quad
\src{src/Development/AutomorphicArithmetic/MeanSquareBounds/TailKExp.lean}{92}{finite-height contour}\quad
\src{src/Development/AutomorphicArithmetic/MeanSquareBounds/Kappa4Loglog.lean}{126}{active logarithmic derivative}.

The first progression is \(q\equiv3\bmod8\) with \(U=e^z,\ z=Y^{\varepsilon/5}\), and the base list consists of the primitive representatives of characters modulo \(N_0\). The second fixes this \(q\), uses \(U=e^{z'}\), \(z(\log Y)^2\le z'\le K^2z(\log Y)^2\), and restricts \(d\equiv7\bmod16\) and \(d\equiv-1\) modulo \(q\) and every odd prime \(\le z\). The modulus is \(16q\prod_{p\le z,\ p\ {\rm odd\ prime}}p\). Its logarithm is at most \(4z\) eventually. Its large prime factor \(q\) exceeds \((Y^K)^{400}\). The base list now includes the old characters and their \(q\)-twists, at most \(2\varphi(N_0)\), with conductor bound \(F=N_0e^z\). The common selection is used with \(\varepsilon'=6\varepsilon/5\) in the second progression. These choices give the hypotheses just stated, with the admissible \(z'\) range and \(F\le\sqrt U\) holding eventually, uniformly in the selected \(q\). Taking \(\eta=1/2\) successively produces both primes. At the top, \(N_0=8M,\tau=\varepsilon/2\) and \(z'=z(\log Y)^2\), with \(M\) the conductor of the target character. The bounds concern the three nonempty binary twist patterns \(J_q,J_d,J_qJ_d\) and their specified primitive base characters, with no bound asserted here for the untwisted target. \src{src/Development/AutomorphicArithmetic/MeanSquareBounds/Stage1Succ.lean}{36}{first progression}\quad
\src{src/Development/AutomorphicArithmetic/MeanSquareBounds/CharOf.lean}{61}{second progression}\quad
\src{src/Development/AutomorphicArithmetic/MeanSquareBounds/PairShare.lean}{26}{nested good sets}\quad
\src{src/Development/AutomorphicArithmetic/MeanSquareBounds/pairEx.lean}{112}{active parameters and final pair}.

\subsection{Cohomological windows and compatibility with the host}

We use the source's proved \(\mathrm{VWinAt}\) input in its actual quantifier order. For a fixed primitive target \(\psi\), \(0<\varepsilon<1/52\), and a page-index function \(e(v)\) with \(v^{2+8\varepsilon}\le 2^{e(v)}\le2v^{2+8\varepsilon}\) for \(v\ge1\), it provides \(J\ge1,X_0\). For every \(X\ge X_0\) and every pair meeting the prime conditions at \(Y=X^{1/J}\), it then gives \(W\ge Y\) with \(W^2\le X\). For all \(v\in[W,W^2]\) and every field \(E\) satisfying the corresponding \(\mathrm{IsChiField}\) relation,
\[
 \mathrm{VVle}\!\left(E,\lfloor v^{401}\rfloor,\exp(v^\varepsilon/2)\right)
\]
holds. Here, writing \(S_\ell\) for the places above \(\ell\), the kernel is the intersection of the kernels of restriction from \(H^1(G_{E,S_\ell},(\mathbb Q_\ell/\mathbb Z_\ell)(-j))\) to the local decomposition groups at \(w\mid\ell\). Its length is finite for every \(1\le j<t\). For every finite set of primes their lengths summed over twists and weighted by \(\log\ell\) are at most the \(\mathrm{VVle}\) budget. Equivalently, these nonnegative prime terms form a summable family within that budget. The assertion is uniform over the algebraic closures and local choices in the source. \src{src/Development/AutomorphicArithmetic/CharacterBoxes/ExistsPageN.lean}{12}{prime-length budget}\quad
\src{src/Development/AutomorphicArithmetic/CharacterBoxes/ExistsPageN.lean}{44}{cohomological window quantifiers}.

The ingredients of the closed input should be distinguished from the selection we prove here. The source first bounds a finite number of exceptional twist indices uniformly in the auxiliary primes. A local \(\ell\)-adic centre can capture at most one integer \(1\le j\le U_0\) at distance strictly smaller than \(\ell^{-\lfloor\log_\ell U_0\rfloor}\). The difference of two would be divisible by \(\ell^{\lfloor\log_\ell U_0\rfloor+1}>U_0\). The ranks controlling the number of centres are uniform in the quadratic auxiliary primes. With \(R\) the sum of these rank bounds, remove the at most \(R\) troublesome cutoffs on the logarithmic scale. More explicitly, to make every surviving \(\alpha=\log j/C_t\) satisfy \(\alpha\le u\Rightarrow B_0\alpha\le u\) on an interval \(u\in[v,2v]\subseteq(0,\log X]\), with \(B_0=2C_KC_t/\varepsilon\ge1\), divide the current upper endpoint by \(2B_0\) whenever a bad cutoff is at most the current endpoint and strictly greater than that endpoint divided by \(2B_0\). That cutoff can never trigger the antecedent in the smaller interval. After at most \(R\) such steps take \(v\) half the remaining endpoint. Thus \(v\ge\log X/[2(2B_0)^R]\). The source allows \(J=2^{R+2}B_0^R\). Returning to the original scale forces a bad index still below \(x^{C_t}\) to be at most \(x^{\varepsilon/(2C_K)}\). This permits the coarser finite-layer estimate on those few indices. For the remaining indices the quantitative tower bounds apply without that reduction. \src{src/Development/AutomorphicArithmetic/AutomorphicVarieties/LayeredIsomorphisms/Near1.lean}{6}{one exception per centre}\quad
\src{src/Development/AutomorphicArithmetic/AutomorphicVarieties/LayeredIsomorphisms/ThmBabOf.lean}{27}{uniform rank and interval exponent}\quad
\src{src/Development/AutomorphicArithmetic/AutomorphicVarieties/LayeredIsomorphisms/Window1.lean}{6}{logarithmic window argument}.

Section~\ref{famsel:complement} gives the resulting inequality, including its dyadic, odd-prime minus and odd-prime plus terms. For Case I the rank-three input produces a nonzero integer of polynomially bounded logarithmic size in the conductor and twist. Its valuations bound a whole orbit of character lengths away from the test prime. Using a second integer at the first test prime gives a bound uniform in the finite set of coefficient primes. The formalization passes the local--global input \(\mathrm{galoisLocalGlobalStrong\_recParameter}\) to the cohomology construction and combines these counts with duality. Proposition~\ref{addg:unitary} proves the classical attachment for the standard definite group. The nonconstant curves of the deformation problem and the complementary character counts retain the hypotheses specified in Section~\ref{famsel:complement}. \src{src/Development/AutomorphicArithmetic/AutomorphicVarieties/UnitaryThreeGroups/CorGSumLog.lean}{37}{two-test-prime valuation bound}\quad
\src{src/Development/AutomorphicArithmetic/AutomorphicVarieties/HeightAndVanishing/JointWinOfCounts.lean}{30}{vanishing cases to windows}\quad
\src{src/Development/AutomorphicArithmetic/AutomorphicVarieties/ModularityArgument/ThresholdArguments/DenseThresholdCases/VwinOfRow.lean}{4}{active row to window}\quad
\src{src/Development/AutomorphicArithmetic/CharacterBoxes/MainConclusion.lean}{7}{actual choice of the row}.

To use both selections together, first take \(X\) large enough that \(Y=X^{1/J}\) exceeds the pair threshold and any required lower scale. Choose \(q,d\) using the analytic pair theorem on \([Y,X]\), and only then use \(\mathrm{VWinAt}\) to get \(W\) for exactly those primes. The host frame has degree parameter \(c2^{e}\) for a fixed \(c\ge1\) depending on \(\psi\), and uses \(\kappa=c^{1/(2+8\varepsilon)}\). Taking \(W\ge\kappa^{400}\) suffices for
\[
 [\kappa W,(\kappa W)^{3/2}]\subseteq[W,W^2]\subseteq[Y,X].
\]
At \(x\) in that interval use the page of \(x/\kappa\). Multiplication by \(c\) gives \(x^{2+8\varepsilon}\le c2^{e(x/\kappa)}\le2x^{2+8\varepsilon}\). The scale identity \(x=\kappa v\) with \(v\ge W\ge\kappa^{400}\) gives \(x^{400}=\kappa^{400}v^{400}\le v^{401}\). The source first uses monotonicity to replace the cutoff \(\lfloor v^{401}\rfloor\) by \(\lfloor x^{400}\rfloor\), then applies the actual rank estimate \(t\le x^{400}\). On the rescaled window the stored budget is at most \(\sqrt b=\exp(x^\varepsilon/2)\). \src{src/Development/AutomorphicArithmetic/CharacterBoxes/WinPairs.lean}{34}{same pair and common window}\quad
\src{src/Development/AutomorphicArithmetic/CharacterBoxes/WinPairs.lean}{58}{page and cohomological rescaling}.

For each point, \(\mathrm{fieldChi}\) carries this page and the same primes into the host and yields \(n=c2^e\), fixed \(o=\Lambda.o\), and \(QS\mid\operatorname{lcm}(2^{e+2},M)\). Thus the only possible odd divisors of \(QS\) are primes dividing \(M\). The geometry retains that host when adding the ranks and vertices, so the cohomological budget is for the very field used at that point, and the discarded-character estimates use the same pair. Write \(\delta_1=x^{-\varepsilon/8}\). The actual rank construction chooses even \(m,k\) in
\[
 [\sqrt{x\delta_1},2\sqrt{x\delta_1}]\quad\hbox{and}\quad
 [x^{9/2+10\varepsilon},2x^{9/2+10\varepsilon}]
\]
by taking \(2\lceil A/2\rceil\) at the respective lower endpoints. It then defines \(t\) by the integer rank equation of \(\mathrm{RankVal}\). The eventual comparisons needed in Proposition~\ref{ar:four} follow directly: \(xm\le k\), \(m\in[x^{1/2-\varepsilon},x^{1/2+\varepsilon/2}]\), and \(k/(n^2x^{1/2-\varepsilon})\le 2x^{-5\varepsilon}\le1\). The source separately bounds \(t\le x^{400}\) eventually for the active small-\(\varepsilon\) range. The section on simultaneous transport then chooses the full translate map within the family at this \(x\). Choosing these ranks or local copies independently of the point's host would not suffice for that transport. \src{src/Development/AutomorphicArithmetic/CharacterBoxes/FieldChi.lean}{20}{fixed host field and ramification}\quad
\src{src/Development/AbelianVarietyTheory/Proofs/GlobalInterpolation/ExistsRanksRankValFrame.lean}{45}{rank construction}\quad
\src{src/Development/AutomorphicArithmetic/CharacterBoxes/EventuallyRankLeRpowOfLtOneDivFiftytwo.lean}{6}{rank bounded by the cohomology cutoff}\quad
\src{src/Development/AutomorphicArithmetic/ShortNorms/supplyOfAdmOf.lean}{93}{same-host supply and budget}.



\subsection{Low-complexity cohomology and the exceptional twists}\label{sup:length}

The exceptional indices in the window argument are handled by a bound on a fixed low cyclotomic layer. Let \(L\) be totally imaginary, \(n=[L:\mathbf Q]\), \(D=|\operatorname{disc} L|\), and \(\ell\) a prime. Assume that \(L\) contains \(\mu_\ell\) for odd \(\ell\) and \(\mu_4\) for \(\ell=2\). The bound called \code{LemmaK} asserts the existence of an absolute \(C_K>0\) such that, for every such \(L,\ell\), every choice of algebraic closure and local data, and every integer \(E\ge1\),
\begin{equation}\label{sup:lemmaK}
 H^2(G_{L,S},\mathbf Z_\ell(E+1))\ \text{is finite},\qquad
 \log\#H^2(G_{L,S},\mathbf Z_\ell(E+1))
 \le \exp\bigl((2En)^{C_K}\log(2D)\bigr).
\end{equation}
Theorem~\ref{newcoh:lemmaK} proves this bound with \(C_K=32\) using published results. The remainder of this subsection identifies the estimates in the alternative argument of the formalization. Here \(S\) consists of the finite places over \(\ell\), and the right-hand side bounds the logarithm of the order. The compact module has action \(\chi_\ell^{E+1}\). Its Tate dual is the discrete module \(A_E=(\mathbf Q_\ell/\mathbf Z_\ell)(-E)\), with action \(\chi_\ell^{-E}\). In particular the local invariant term below uses the same discrete module as the degree-one restriction kernel \(\Sha^1(A_E)\). The negative sign follows from \(1-(E+1)=-E\). \src{src/Development/AutomorphicArithmetic/AutomorphicVarieties/GeneratorCountTuples/LemmaK.lean}{9}{roots of unity and cardinality bound}
\src{src/Formalism/GaloisCohomology/Statement.lean}{67}{two coefficient conventions}.

\paragraph{Cyclotomic descent.}
Take the layer \(L'/L\) of degree \(d=\ell^m\) with \(m=\lceil\log_\ell(4E+4)\rceil\). Let \(n'=nd\), and \(\ell^c\) be the highest order of an \(\ell\)-power root of unity in \(L'\). Then \(2E+1<\ell^c,r_2(L')\), \(d\le(n+1)(4E+4)\) and \(\ell^c\le 2n'\). Indeed \(\ell^{m-1}<4E+4\), \(\ell-1\le n\), and \(r_2(L')=n'/2\). The cyclotomic degree formula gives the last bound. The class-module descent supplies one common module and identifies the lengths after taking the cokernels of \(g-1\) and \(g^d-1\) on it, where \(g\) already includes the twist. Now
\[
 g^d-1=(g-1)\sum_{i=0}^{d-1}g^i,\qquad
 \operatorname{coker}(g^d-1)\twoheadrightarrow\operatorname{coker}(g-1).
\]
Thus the higher layer bounds the original length, including when the latter is initially allowed to be infinite. This comparison uses the supplied class-module descent, and makes no injectivity assumption about restriction on global cohomology. For the discriminant, write \(X^N-1=P(X)Q(X)\) with \(P\) the minimal polynomial of a root \(\zeta\) generating \(L'/L\) and \(N\le\ell^c\). Differentiation gives \(N=\zeta P'(\zeta)Q(\zeta)\). The different contains \(P'(\zeta)\), hence contains \((N)\). Taking norms in the discriminant tower gives \(D'\le D^d(\ell^c)^{n'}\). These estimates keep the field growth polynomial in \(2En\) on the \(\log(2D)\) scale. \src{src/Development/EtaleCohomology/CyclotomicDiscriminants/SLayerCut.lean}{11}{relative different}
\src{src/Development/EtaleCohomology/CyclotomicDiscriminants/SLayerCut.lean}{40}{cokernel length comparison},
\src{src/Development/EtaleCohomology/CyclotomicDiscriminants/SLayerCut.lean}{55}{choice of layer}.

\paragraph{A count uniform in the remaining layer number.}
Work over the enlarged base \(L'\), set \(M=2E+1,\ b=2E+M(c+M)\), and let \(q=\ell^\nu\) be the root order used for the finite-layer test. The finite-layer class-field injection and the arithmetic homology detector give the following bounds for the count of equivariant homomorphisms \(\mathrm{EqHom}\). Write \(T_j=\#\operatorname{tors}H_j(\mathrm{GL}_{4E+6}(\mathcal O_{L'}),\mathbf Z)\), and \(g_0=2n'+\lfloor\log_2 h_{L'}\rfloor+2\). If \(I\) is the image of the twisted norm test, then
\begin{equation}\label{sup:lengthcount}
\begin{aligned}
 \#I&\le\#\operatorname{im}(\ell^{M-1}\Phi)\le q^{r_2(L')}T_{2E+1}T_{2E},\\
 \#\mathrm{EqHom}\ q^{r_2(L')}
 &\le \#I\,\ell^{(b+v_\ell(E))g_0+c+v_\ell(E)} .
\end{aligned}
\end{equation}
Here is the first injection. Write \(F/L'\) for the finite cyclotomic test layer,
\(\gamma\) for its generator acting on \(\mu_q\) by \(\kappa=1+\ell^c\),
and \(a=\kappa^E\). Put \(\Delta\beta=\gamma^{-1}(\beta)\beta^{-a}\).
Let \(U\subset F^\times\) consist of elements whose divisor has the form
\((\mathfrak d_{\gamma P}-a\mathfrak d_P)_P\) for an integer-valued
divisor \(\mathfrak d\). The test on the finite group
\(W=U/(\Delta F^\times\,\mu(F))\) is
\[
 [\alpha]\longmapsto [(N_{\gamma,a}\alpha)^{\ell^b}]
 \quad\text{in }F^\times/\big((F^\times)^q\mu(F)\big),\qquad
 N_{\gamma,a}\alpha=\prod_{j=0}^{\operatorname{ord}(\gamma)-1}
                       \gamma^j(\alpha)^{a^j}.
\]
For each \(z\in I\) choose a lift \(\alpha_z\) with a divisor
\(\mathfrak d_z\) satisfying
\(v_P(\alpha_z)=(\mathfrak d_z)_{\gamma P}-a(\mathfrak d_z)_P\).
The arithmetic cycle theorem supplies a mod-\(q\) cycle \(y_z\)
of degree \(M\) over \(\mathcal O_{L'}\) and a nonzero denominator in
\(\mathcal O_F\). Exclude from the primes of \(F\) those dividing
\(\ell\) or any of these finitely many denominators. At each remaining
\(P\) of relative residue degree one choose a multiplicative index
\(v_P^\times:k(P)^\times\to\mathbf Z/q\) with kernel \(k(P)^{\times q}\).
Such an index exists because \(k(P)^\times\) is cyclic and reduction
preserves the primitive \(q\)-th root. The residue evaluations
\(\Phi_P\) provided by the source then satisfy
\[
 \Phi_P(y_z)=
 \prod_{\substack{0\le i<M\\ i\ne E}}(\kappa^E-\kappa^i)\,
 v_P^\times\!\left(\overline{N_{\gamma,a}\alpha_z}\right).
\]
The bar denotes reduction at \(P\). Each factor of the product has
valuation \(c+v_\ell(|E-i|)\): removing a unit power of \(\kappa\)
reduces this to the binomial valuation of \(\kappa^{|E-i|}-1\),
with \(c\ge2\) at two. Since \(0<|E-i|<M\), the valuation of the
product is at most \((M-1)(c+M)\). Hence it divides
\(\ell^{M(c+M)}\) in \(\mathbf Z/q\).
Thus equality of \(\ell^{M-1}\Phi_P(y_z)\) and
\(\ell^{M-1}\Phi_P(y_{z'})\) forces the quotient of the two
\(\ell^b\)-powered norms to be a \(q\)-th power in \(k(P)^\times\).
If equality holds at every test place, the relative-degree-one power
criterion in Section~\ref{newcoh:ktheory} makes that quotient a
\(q\)-th power in \(F^\times\), so \(z=z'\).
Consequently \(z\mapsto(\ell^{M-1}\Phi_P(y_z))_P\) is an injection
of sets for any such choices. The product evaluation \(\Phi=(\Phi_P)_P\)
is primitive. The scaled homology count in Section~\ref{newcoh:ktheory}
bounds its image, with the exponent \(r_2(L')\) coming from the field
over which the cycles live.
\src{src/Development/AutomorphicArithmetic/AutomorphicVarieties/ModularityArgument/ClassCounting/SelmerOrderCounts/ActUPow.lean}{37}{representatives of the test image},
\src{src/Development/AutomorphicArithmetic/AutomorphicVarieties/ModularityArgument/ClassCounting/LemmaKOfSteps.lean}{86}{cycle evaluation on twisted norms},
\src{src/Development/EtaleCohomology/CyclotomicDiscriminants/SFewOfCut.lean}{67}{divisibility of the evaluation coefficient}, and
\src{src/Development/EtaleCohomology/CyclotomicDiscriminants/SFewOfCut.lean}{197}{product injection and scaled count}.

For the second bound one may see the length loss via the \(\ell\)-power filtration of the norm-test kernel. If \(W\) is the finite source and its kernel is killed by \(\ell^s u\), \(\ell\nmid u\), its \(\ell\)-primary part is killed by \(\ell^s\). For every \(i\ge1\), multiplication by \(\ell^{i-1}\) embeds \(W[\ell^i]/W[\ell^{i-1}]\) into \(W[\ell]\). Hence its kernel has \(\ell\)-primary order at most \(\ell^{s g_0}\) when \(\#W[\ell]\le\ell^{g_0}\). The image contributes at most \(\#I\) by exactness. The twisted powers telescope to \(N_{\gamma,a}(\Delta\beta)=\gamma^{-1}(\beta)^{1-a^{\operatorname{ord}\gamma}}\), with \(a=(1+\ell^c)^E\). The valuation of the scalar is \(\nu+v_\ell(E)\), giving \(s=b+v_\ell(E)\) after removing the modulus \(q\). The unit-lattice order calculation gives the factor \(q^{r_2(L')}\) and the remaining factor of \(\ell^{c+v_\ell(E)}\). Thus the same power of \(q\) cancels and the bound is uniform in \(\nu\).
\src{src/Development/EtaleCohomology/CyclotomicDiscriminants/SFewOfCut.lean}{187}{joint detector},
\src{src/Development/AutomorphicArithmetic/AutomorphicVarieties/ModularityArgument/ClassCounting/SelmerOrderCounts/PCount.lean}{12}{finite primary filtration},
\src{src/Development/AutomorphicArithmetic/AutomorphicVarieties/ModularityArgument/ClassCounting/SelmerOrderCounts/KerFDelNm.lean}{20}{norm and difference},
\src{src/Development/AutomorphicArithmetic/AutomorphicVarieties/ModularityArgument/ClassCounting/SelmerOrderCounts/CardEQ.lean}{36}{unit-lattice order},
\src{src/Development/AutomorphicArithmetic/AutomorphicVarieties/ModularityArgument/ClassCounting/NumThree.lean}{223}{cancellation and exponents}.

The class-field injection and layer realization apply to any finite subset of the image of \(\Sha^1(A_E)\) in the infinite tower by taking one layer beyond the finitely many thresholds for its elements. This subset is bounded in cardinality by the right side of the cancelled inequality, independently of which layer was chosen. A larger finite subset would contradict that bound, so the whole image is finite. The homology torsion bound, obtained from an arithmetic face complex and orbit and stabilizer estimates, now applies to the two factors \(T_j\). Writing \(u=2En'\), \(Y=\log(2D')\), the remaining contribution \(((b+v_\ell(E))g_0+c+v_\ell(E))\log\ell\) is at most \(u^9Y\). Taking logarithms of the count and using the preceding discriminant and length comparisons gives \(\operatorname{length}_{\mathbf Z_\ell}\Sha^1(A_E)\log\ell\le \exp((2En)^C\log(2D))\) for an absolute \(C\). The face complex and the class-field injection are used here at their supplied scope.
\src{src/Development/AutomorphicArithmetic/AutomorphicVarieties/ModularityArgument/ClassCounting/RingEquivalences/OrderOfOneAddPrimePowZmodPrimePow.lean}{454}{one finite layer for a finite subset},
\src{src/Development/AutomorphicArithmetic/AutomorphicVarieties/HermiteNormalForms/Assembly.lean}{14}{arithmetic torsion input},
\src{src/Development/AutomorphicArithmetic/AutomorphicVarieties/ModularityArgument/ClassCounting/NumThree.lean}{350}{uniform numerical absorption}.

\paragraph{Duality and a local length calculation.}
Put \(B=H^2(G_{L,S},\mathbf Z_\ell(E+1))\) and let \(\rho\) be its map into the product of the local \(H^2\)'s. Poitou--Tate identifies the length of \(\ker\rho\) with the length of \(\Sha^1(A_E)\) and the length of \(\operatorname{coker}\rho\) with that of \(H^0(G_{L,S},A_E)\). Local Tate duality gives the length of each local \(H^2\) as that of \(H^0(G_{L_w},A_E)\). The source obtains these claims from the finite-level dualities and finiteness through the cohomology limits. Additivity for the kernel, image and cokernel gives the equality with the global \(H^0\) term on the left, and therefore bounds \(\operatorname{length}B\) by \(\operatorname{length}\Sha^1(A_E)+\sum_{w\mid\ell}\operatorname{length}H^0(G_{L_w},A_E)\). Total imaginarity supplies the real-place hypothesis also at \(\ell=2\).
\src{src/Development/AutomorphicArithmetic/AutomorphicVarieties/GaloisLibrary/GaloisCohomologyLength/Formula.lean}{15}{finite-level dualities and limiting length formula}.

For an explicit bound, let \(U_w\subseteq\mathbf Z_\ell^\times\) be the local cyclotomic image and \(a_w=[\mathbf Z_\ell^\times:U_w]\). The invariants in \(A_E\) are killed by \((1+\ell^2)^{Ea_w}-1\), since \((1+\ell^2)^{a_w}\in U_w\) and replacing its action by the inverse changes the difference from \(1\) by a unit. For \(r\ge2\), raising a unit \(1+\ell^r v\), \(\ell\nmid v\), to the \(\ell\)-th power increases \(r\) by exactly one by the binomial expansion. Raising it to an exponent prime to \(\ell\) leaves \(r\) unchanged. Thus the annihilator has valuation \(2+v_\ell(Ea_w)\), including at \(\ell=2\). The submodule of \(\mathbf Q_\ell/\mathbf Z_\ell\) killed by \(\ell^k\) is \(\ell^{-k}\mathbf Z_\ell/\mathbf Z_\ell\), of length \(k\). As \(a_w\le n\) by the inertia-image index bound and there are at most \(n\) local terms, this adds at most
\[
 n\bigl(2\log\ell+\log(nE)\bigr)
\]
to \(\operatorname{length}B\log\ell\). This bound is polynomially bounded in \(2En\) since \(\ell\le n+1\). Finite length gives order \(\ell^{\operatorname{length}B}\) by the composition factors, each of size \(\ell\). Increasing the absolute exponent proves the form \eqref{sup:lemmaK}. \src{src/Development/AutomorphicArithmetic/AutomorphicVarieties/ModularityArgument/ClassCounting/SFinal.lean}{35}{local test unit and inertia bound}
\src{src/Development/AutomorphicArithmetic/AutomorphicVarieties/ModularityArgument/ClassCounting/SFinal.lean}{100}{finite order and final comparison}.

\paragraph{Return to the same window.}
The active assembly passes the odd arithmetic input, the Poitou--Tate statement, local finiteness and local Tate duality to the class-counting proof of \code{lemmaK}. The resulting bound is one input in \code{vwin\_of\_hypotheses\_H\_E\_alt}. For an exceptional index \(E\) still below \(x^{C_t}\) in a good window, \(E\le x^{\varepsilon/(2C_K)}\). The fixed low layer used in the tower comparison has degree at most \(n_0\) and log discriminant at most \(c_0 x^{\varepsilon/4}\) in the form \(\log(2D)\). Consequently \((2En)^{C_K}\log(2D)\le (2n_0)^{C_K}c_0 x^{3\varepsilon/4}\). The length bound uses the factor \(2\), since \(1/\log\ell\le2\). The tower comparison adds at most \(R(m-j_0)\) in natural subtraction to the fixed-layer length, while nonexceptional indices use the linear tower bound directly. This puts the exceptional contribution below the final \(\exp(x^\varepsilon/2)\) scale for large \(x\). Indeed any fixed power of \(x\), and any fixed power of \(\log x\), multiplying \(\exp(C x^{3\varepsilon/4})\) can be absorbed at that scale by taking logarithms. The other ramification and parity cases in the \(\mathrm{VWin}\) assembly are separate bounds. \src{src/Development/AutomorphicArithmetic/AutomorphicVarieties/ModularityArgument/ThresholdArguments/DenseThresholdCases/LemmaKOfLitN.lean}{12}{instantiated duality and odd input}
\src{src/Development/AutomorphicArithmetic/AutomorphicVarieties/ModularityArgument/ThresholdArguments/DenseThresholdCases/WindowBoundOfLiteratureLevel.lean}{25}{actual window assembly},
\src{src/Development/AutomorphicArithmetic/AutomorphicVarieties/LayeredIsomorphisms/ImportsAlt.lean}{9}{exceptional and ordinary tower bounds},
\src{src/Development/AutomorphicArithmetic/AutomorphicVarieties/LayeredIsomorphisms/Klen.lean}{23}{order to length}.

\subsection{Integral cycles detecting the twisted norm}\label{addc:section}

We construct the cycles and the common family of evaluations used in
\eqref{sup:lengthcount}. Write \(A=\mathbf Z/q\), \(q=\ell^\nu\),
\(w=\ell^c\), \(\kappa=1+w\), \(a=\kappa^E\), \(M=2E+1<w\).
Here \(\nu\ge c\), and \(c\ge2\) if \(\ell=2\). The layer \(F/L'\)
is cyclic with generator \(\gamma\), and contains a primitive
\(q\)-th root \(\zeta\) on which \(\gamma\) acts by \(\kappa\).
In this cyclotomic layer \(\gamma\) and \(\kappa\bmod q\) have the
same order. Write \(C_q=\langle h\rangle\) for the group of order \(q\).
All bar chains below have coefficients in \(A\).

\begin{lemma}[One integral test for each residue field]\label{addc:residue}
Let \(k\) be a finite field containing a primitive root
\(\xi\) of order \(q\), and let \(v:k^\times\to A\) be a
homomorphism. There are homology evaluations \(\theta_V\) of degree
\(M\) on \(\mathrm{GL}(V)\) for finite-dimensional \(k\)-spaces \(V\).
They are invariant under change of basis and additive under direct
sum of representations. Put
\[
 z=\sum_{j_1,\dots,j_E\bmod q}
       [h^{j_1}|h|\cdots|h^{j_E}|h]\in C_{2E}(C_q;A).
\]
If \(h\) acts on \(V=\bigoplus_x V_x\) by \(\xi^x\) on \(V_x\),
and \(g=\bigoplus_x g_x\) preserves the summands, then the prism
of \(g\) on the image of \(z\) has value
\(\sum_x x^E v(\det g_x)\).
\end{lemma}
\begin{proof}
The first face in the boundary of \(z\) occurs \(q\) times. The
other faces cancel in consecutive pairs after translating a summation
index by one, including the last face. Thus \(z\) is a cycle.
Write \(|k^\times|=q m\) and choose a generator with \(m\)-th power
\(\xi\). Such a generator exists since units modulo \(qm\) surject
onto units modulo \(q\), by prime-power reduction and the Chinese
remainder theorem. With \(L(b)\) its integer discrete logarithm in
\([0,qm)\), put
\[
   u(b,d)=\frac{L(b)+L(d)-L(bd)}{qm}\pmod q .
\]
This symmetric two-cocycle has zero coboundary by cancellation of
the four logarithmic differences. On \(k^\times\) take the cup
product \(u^{\smile E}\smile v\). The prism of \(g\) inserts \(g\)
in every position with alternating signs. Terms in the first \(2E\)
positions cancel in pairs by symmetry of \(u\). The final insertion
contributes \(v(g)\) times the product of the \(E\) carries. For
\(h\mapsto\xi^x\), summing each carry over \(j\bmod q\) gives
\(q L(\xi^x)/(qm)\equiv x\bmod q\) by telescoping. This gives
\(x^E v(g)\) in dimension one.

For general \(V\), sum this cocycle over the lines of \(V\).
More explicitly, lift a bar simplex through the action on
\(\mathbf P(V)(k)\), once for each starting line. Choose one nonzero
vector \(v_x\) on each line. For arrows \(g_1,\ldots,g_M\) and initial
line \(x_0\), put \(x_i=g_i^{-1}x_{i-1}\) and read the scalar \(b_i\)
from \(g_i v_{x_i}=b_i v_{x_{i-1}}\). Evaluate the cocycle on these
scalars. Composition multiplies the scalars and
the arrows permute the starting lines, so the sum commutes with the
bar faces. Changing the vectors gives naturally isomorphic functors
from the action groupoid to \(k^\times\). The simplicial prism of
that natural isomorphism gives the same value on cycles. This also
proves invariance under change of basis. For a block sum \(V\oplus W\),
lines with nonzero first projection lie over \(\mathbf P(V)\) with
\(|k|^{\dim W}\) lifts per line. Normalize their first projections
by the chosen vectors of \(V\). Their scalars are the scalars in \(V\).
The other lines are precisely \(\mathbf P(W)\).
Since \(|k|\equiv1\bmod q\), the evaluation is additive.

On a scalar block the prism value is a homomorphism in \(g\).
Indeed, composing the prisms for two successive arrows gives the
prism for their product modulo the boundary over their two-simplex.
Here both arrows centralize the cycle. This homomorphism factors
through the determinant. In dimension at least two the elementary
matrices generate \(\mathrm{SL}(V)\), and each is a commutator
with a diagonal matrix since \(|k|>2\). Its value on
\(\operatorname{diag}(b,1,\dots,1)\) follows from additivity and the
one-dimensional computation. Applying additivity to all character
blocks gives the formula. In this application each block prism is
the image of the same cycle on \(C_q\times\mathbf Z\).
\end{proof}

\begin{lemma}[Isolation of one cyclic direction]\label{addc:filter}
Set \(d_t(h)=(h,h^t)\) for \(0\le t\le M\) and
\(\chi_t(j_1,j_2)=j_1+t j_2\) on \(J=A^2\).
There are \(n_t\in A\), with \(n_{t_0}\) a unit, and a chain \(B\)
on \(C_q^2\) with \(\partial B=\sum_t n_t(d_t)_*z\).
For this \(t_0\) there is an \emph{integer}-valued function \(\eta\)
on \(J\) such that, in the integral group ring of the additive group
of \(A\),
\[
 (\chi_t)_*\eta=0\quad(t\ne t_0),\qquad
 (\chi_{t_0})_*\eta=\sum_{i=0}^M p_i[\kappa^i],\qquad
 \sum_i p_i X^i=P(X)=\prod_{j=0}^{M-1}(X-\kappa^j).
\]
Here pushforward adds coefficients in each fibre and \([x]\) is a
group-ring basis element.
\end{lemma}
\begin{proof}
Tensor the two cyclic periodic free resolutions over \(\mathbf Z\)
to resolve the trivial module for \(C_q^2\). After tensoring with
trivial coefficients \(A\), the maps \(h-1\) and the norm both
vanish. Thus \(H_{2E}(C_q^2;A)\simeq A^{2E+1}\). The \(M+1\)
classes give an \(M\)-by-\((M+1)\) matrix. Its Smith form over
\(\mathbf Z/\ell^\nu\) has a free zero column. Transporting its
coordinate vector back gives a relation with a unit entry, hence
the bounding chain \(B\).

Let \(U=X^w-1\) and let \(T\) act on \(\mathbf Z[X]\) by
\(Tf(X)=f(X^\kappa)\). Then \(Tf\equiv f\bmod U\) and
\(TU=(1+U)^\kappa-1\equiv\kappa U\bmod U^2\). Consequently
\((T-\kappa^j)(U^j f)\) is divisible by \(U^{j+1}\). Applying
the \(M\) factors to \(X\) shows that \(\sum_i p_i X^{\kappa^i}\)
is divisible by \(U^M\) over \(\mathbf Z\). Substitute \(X=[1]\)
in the group ring. Write the resulting element as
\(([w]-[0])^M H\). For each \(t\ne t_0\) choose \(b_t\in A\) with
\((t_0-t)b_t=w\). This is possible since \(0<|t_0-t|\le M<w\)
and \(q\) is an \(\ell\)-power. Lift \(H\) by \(x\mapsto(x,0)\)
and multiply it by
\(\prod_{t\ne t_0}([(-t b_t,b_t)]-[(0,0)])\).
Pushing the product forward gives precisely the two asserted
formulas. The argument also covers \(c=\nu\).
\end{proof}

\begin{theorem}[Simultaneous cycles for the arithmetic detector]\label{addc:cycle}
Suppose that \(\alpha\in F^\times\) and an integral-coefficient
divisor \(\mathfrak d\) satisfy
\(\operatorname{div}\alpha=(\gamma^{-1}-a)\mathfrak d\), with
\((\gamma^{-1}\mathfrak d)(P)=\mathfrak d(\gamma P)\).
There are a matrix size \(N\), a cycle \(y_\alpha\) of degree \(M\)
on \(\mathrm{GL}_N(\mathcal O_{L'})\), and a finite set of excluded
primes. At every remaining prime \(P\) of \(F\) of relative residue
degree one, reduction followed by Lemma~\ref{addc:residue}, with
\(\xi=\bar\zeta\), gives
\[
 \Phi_P(y_\alpha)=
  \prod_{\substack{0\le j<M\\j\ne E}}(a-\kappa^j)\,
       v_P^\times\!\left(\overline{N_{\gamma,a}\alpha}\right).
\]
The evaluation family can be fixed independently of \(\alpha\) once
the indices \(v_P^\times\) have been chosen.
\end{theorem}
\begin{proof}
Put \(Q(X)=P(X)/(X-a)\) and \(\beta=Q(\gamma^{-1})\alpha\),
where an integral group-ring operator acts multiplicatively on
\(F^\times\). Then
\(\operatorname{div}\beta=P(\gamma^{-1})\mathfrak d\).
Let \(\mathfrak a\) be the fractional ideal of divisor \(\mathfrak d\).
On \(F^J\) consider the two \(\mathcal O_{L'}\)-lattices
\[
 \Lambda_0=\bigoplus_{j\in J}\mathcal O_F,\qquad
 \Lambda_1=\bigoplus_{j\in J}\mathfrak a^{\eta(j)}.
\]
The diagonal representation \(\rho\) of \(C_q^2\), with characters
\(j\) evaluated at \(\zeta\), preserves both. We construct
\(L'\)-linear isomorphisms \(g_t:\Lambda_0\to\Lambda_1\) commuting
with \(\rho d_t\). For \(t\ne t_0\), use the blocks
\(\chi_t(j)=x\). The determinant ideal of each target block is the
unit ideal by Lemma~\ref{addc:filter}. Over the Dedekind domain
\(\mathcal O_F\), a sum of \(s>0\) fractional ideals is isomorphic
to \(\mathcal O_F^{s-1}\) plus their product. Indeed, after scaling
two ideals to coprime integral representatives their sum map onto
\(\mathcal O_F\) splits with kernel isomorphic to their product.
Iteration proves the assertion. Thus one may choose the map on
each block with determinant \(1\), correcting its determinant by
a unit on a source coordinate.

At \(t_0\), on each coordinate with \(\chi_{t_0}(j)=\kappa^i\)
modulo \(q\), change coordinates by the \(L'\)-automorphism
\(\gamma^{-i}\) of \(F\). This is unambiguous by the equality of
orders in the layer. The character on these coordinates becomes \(1\)
in additive exponent notation, so \(h\) acts as \(\zeta\).
Keep the other coordinates unchanged. On the combined exponent-one
block the determinant divisor of the transformed target is
\(\sum_i p_i\gamma^{-i}\mathfrak d=\operatorname{div}\beta\).
Every other block has determinant divisor zero. The same ideal
argument gives a map of determinant \(\beta\) on the exponent-one
block and \(1\) elsewhere. This block is nonempty since \(j=(1,0)\)
occurs. Undo the coordinate change to obtain \(g_{t_0}\).

Write \(H=\operatorname{Aut}_{\mathcal O_{L'}}(\Lambda_0)\),
\(g_*=g_{t_0}\), \(\rho'=g_*^{-1}\rho g_*\), \(k_t=g_*^{-1}g_t\).
Both \(\rho,\rho'\) and all \(k_t\) take values in \(H\).
In the bar construction write \(P_g\) for the homotopy from the
identity to conjugation by \(g^{-1}\). Form the chain in \(H\)
\[
 y=n_{t_0}^{-1}\left(\rho'_* B-\rho_* B+
           \sum_t n_t P_{k_t}(\rho'd_t)_*z\right).
\]
Indeed \(k_t^{-1}(\rho'd_t)k_t=\rho d_t\), so the prism boundary
formula and \(\partial B=\sum_t n_t(d_t)_*z\) give \(\partial y=0\).
In the group containing \(H\) and the \(g_t\), this cycle is
homologous to \(n_{t_0}^{-1}\sum_t n_t P_{g_t}(\rho d_t)_*z\).
To see this, apply the prism identity to \(\rho_*B\) at \(g_*\).
Modulo a boundary it replaces \(\rho'_*B-\rho_*B\) by
\(P_{g_*}\rho_*\partial B\). On each summand the two consecutive
prisms are homologous to the prism of \(g_*k_t=g_t\), using the
boundary of the triangle for the two composable conjugations.
The finite projective \(\mathcal O_{L'}\)-module \(\Lambda_0\)
is a summand of a finite free module. Extending the \(H\)-action
by the identity on a projective complement gives the asserted cycle
\(y_\alpha\) on an integral general linear group.

Exclude ramification, primes dividing \(q\), and all conjugates of
the supports of \(\operatorname{div}\alpha\) and \(\mathfrak d\).
Also discard the finitely many places needed to make the lattices
equal and the displayed maps and their inverses integral. These
exclusions make all conjugates of \(\alpha,\beta\) units before
the residue test. At a remaining relative-degree-one place the Galois extension
splits over the residue field \(k\). The reduction of \(\Lambda_0\)
therefore has coordinates \((\sigma,j)\), with
\(\sigma\in\operatorname{Gal}(F/L')\). In those coordinates conjugating
an \(F\)-coordinate by a power of \(\gamma\) merely permutes its
embeddings. Write \(\sigma\zeta=\zeta^{k_\sigma}\).
The trivial complement contributes zero in positive degree.
By Lemma~\ref{addc:residue}, every \(t\ne t_0\) contributes zero
because its block determinants are \(1\). The remaining prism
contributes \(\sum_\sigma k_\sigma^E v_P^\times(\overline{\sigma\beta})\).
The unit coefficient \(n_{t_0}\) has already been divided out.
Replacing \(\alpha\) by \(\gamma^{-i}\alpha\) in this last weighted
sum for \(\alpha\) multiplies it by \(a^i\), by reindexing
\(\sigma\mapsto\sigma\gamma^i\). For \(\alpha\) the sum is
\(v_P^\times(\overline{N_{\gamma,a}\alpha})\).
Applying \(Q\) gives the claimed factor \(Q(a)\). All residue
evaluations use the single construction of Lemma~\ref{addc:residue}.
\end{proof}

In particular one can use the union of the excluded sets for all
representatives of the finite image \(I\) in \eqref{sup:lengthcount}.
Stabilization preserves the value and repeating a block \(b'\) times
multiplies it by \(b'\), by direct-sum additivity. The exact valuation of \(Q(a)\) is
\((M-1)c+2v_\ell(E!)\): the distances \(|j-E|\) run twice through
\(1,\ldots,E\), and \(v_\ell(\kappa^n-1)=c+v_\ell(n)\) by the
binomial formula, with \(c\ge2\) at two. Thus the scaled test
\(\ell^{M-1}\Phi\) already detects the norm raised to the power
\(\ell^{M-1+(M-1)c+2v_\ell(E!)}\), which suffices for the stated
choice of \(b\). The finite-direction filter, the residue cocycle
and the assembly on the two lattices have corresponding constructions
in the pinned source.
\src{src/Development/AutomorphicArithmetic/AutomorphicVarieties/ModularityArgument/ClassCounting/ExpStep.lean}{116}{Integral direction filter},
\src{src/Development/AutomorphicArithmetic/AutomorphicVarieties/ModularityArgument/ClassCounting/FProd.lean}{101}{cyclic residue cocycle},
\src{src/Development/AutomorphicArithmetic/AutomorphicVarieties/ModularityArgument/ClassCounting/STrans.lean}{411}{transfer over the projective lines}, and
\src{src/Development/AutomorphicArithmetic/AutomorphicVarieties/ModularityArgument/ClassCounting/LatticeBlockClasses/CJoin.lean}{5}{arithmetic cycle on the common lattice}.

\section{Integral cohomology and Frobenius congruences}
\subsection{Integral vanishing, cyclotomic descent and the number of exceptions}
\label{newcoh:section}

For a number field \(F\), let \(D_v,I_v\) denote absolute decomposition
and inertia groups at finite places in a fixed algebraic closure. For a
discrete \(\ell\)-primary \(G_F\)-module \(A\), write \(S_\ell(F,A)\)
for the classes in continuous \(H^1(F,A)\) which restrict to zero on
every \(I_v\) and on every \(D_v\) above \(\ell\). By conjugacy the
definition is independent of the choices of prolongations. When \(A\)
is unramified away from \(\ell\), this is exactly the kernel on
\(H^1(G_{F,S},A)\) used in Section~\ref{sup:length}. Additional
decomposition conditions at the primes of \(N\) define the subgroup
used at level \(\ell N\).
\src{src/Development/AutomorphicArithmetic/AutomorphicVarieties/CharacterLevelBounds/ShaFamMono.lean}{17}{coefficient action and the two levels}.

\begin{proposition}[The trivial-character integral vanishing]\label{newcoh:zero}
Suppose \(\ell\ge5\) is prime, \(e\ge2\) is even and \(\ell-1\mid 2e\).
Let \(O\) be the ring of integers in a finite extension of
\(\mathbf Q_\ell\). For \(A=(O\otimes_{\mathbf Z_\ell}
\mathbf Q_\ell/\mathbf Z_\ell)(-e)\) one has \(S_\ell(\mathbf Q,A)=0\).
In particular its subgroup at level \(\ell N\) has length zero for
every \(N\).
\end{proposition}
\begin{proof}
Every continuous cocycle on an absolute Galois group takes only finitely
many values in \(A\), by compactness. Its class is therefore killed by
some power of \(\ell\). It suffices to show that a class \(x\) killed
by \(\ell\) is zero. Fix a place above \(\ell\). Subtract a coboundary
to obtain a representative \(z\) vanishing on \(D_\ell\). Write
\(\ell z(g)=g b-b\). Then \(b\in A^{D_\ell}\). Every extension
\(\mathbf Q_\ell(\mu_{\ell^n})/\mathbf Q_\ell\) is totally ramified
of degree \(\varphi(\ell^n)\), as follows from the Eisenstein
polynomial \(\Phi_{\ell^n}(1+X)\). Hence the local cyclotomic image is
the full global image, so \(b\in A^{G_{\mathbf Q}}\). It follows that
\(z\) actually has values in \(A[\ell]\).

On \(A[\ell]\) the action is via the scalar character
\(\bar\chi_\ell^{-e}\). Its order divides two. If it is trivial put
\(K=\mathbf Q\). Otherwise \(e\) is an odd multiple of
\((\ell-1)/2\). Since \(e\) is even, \(\ell\equiv1\bmod4\), and the
quadratic Gauss sum identifies the fixed field of its kernel as
\(K=\mathbf Q(\sqrt{\ell})\), of discriminant \(\ell\).
Now \(z|_{G_K}\) is a homomorphism. It vanishes on inertia away from
\(\ell\), where the action on all of \(A\) is trivial and the
cohomology restriction of \(x\) is zero. At \(\ell\) it vanishes on
\(D_\ell\cap G_K\) by normalization. For its conjugates use
\[
 z(hgh^{-1})=h z(g)+(1-hgh^{-1})z(h).
\]
The second term is zero for \(g\in G_K\), and the first term is zero
when \(z(g)=0\). Real inertia has order two and cannot contribute.
Global class field theory now makes this homomorphism factor through
the ordinary ideal class group of \(K\).

For \(K=\mathbf Q\) this group is trivial. For the quadratic field put
\(y=\sqrt{\ell}\). Minkowski's bound supplies each class with an
integral ideal of norm at most \(y\). There are at most \(\tau(n)\)
ideals of norm \(n\), by the three quadratic splitting possibilities
and unique ideal factorization. Hence
\[
 h_K\le\sum_{n\le y}\tau(n)=\sum_{a\le y}\lfloor y/a\rfloor
 \le y(1+\log y)<y^2=\ell .
\]
Thus \(\ell\nmid h_K\), and the restriction vanishes in both cases.
The cocycle factors through \(\operatorname{Gal}(K/\mathbf Q)\) with
values in \(A[\ell]\). This quotient has order dividing two, which
acts invertibly on \(A[\ell]\), so its first cohomology vanishes by
averaging. This proves \(x=0\).
\end{proof}
The local normalization is what makes reduction to \(A[\ell]\)
legitimate even in the residually trivial case. This gives a direct
proof of the source's \(\mathrm{PropZ}\).
\src{src/Development/AutomorphicArithmetic/AutomorphicVarieties/LevelCohomology/PropZ.lean}{11}{two residual characters},
\src{src/Development/AutomorphicArithmetic/AutomorphicVarieties/LevelCohomology/HQuad.lean}{11}{quadratic class bound}.

\paragraph{The split unramified module.}
Assume \(F\) contains \(\mu_\ell\) for odd \(\ell\) and \(\mu_4\) for
\(\ell=2\). Put \(F_\infty=F(\mu_{\ell^\infty})\),
\(\Gamma=\operatorname{Gal}(F_\infty/F)\) and \([F_m:F]=\ell^m\) in
this tower. Let \(X\) be the Galois group of the maximal abelian
pro-\(\ell\) extension of \(F_\infty\) unramified at every finite
place in which every place above \(\ell\) splits completely. Equivalently
it is the maximal abelian pro-\(\ell\) quotient of \(G_{F_\infty}\)
which kills all finite inertia groups and all those decomposition
groups. This definition also specifies what unramified means over the
infinite field. Conjugation induces an action of \(\Gamma\).
Choose a generator \(\gamma\) of \(\Gamma\simeq\mathbf Z_\ell\), write
\(g\) for its action on \(X\), and put \(u=\chi_\ell(\gamma)\).
With \(c=v_\ell(u-1)\) one has \(c\ge1\) and \(c\ge2\) at \(\ell=2\).
The group \(\Gamma\) is an open subgroup of the principal units
under the root-of-unity hypotheses.

\begin{lemma}[Descent with the local conditions]\label{newcoh:desc}
For \(e\ge1\), the Pontryagin dual of
\(S_\ell(F_m,(\mathbf Q_\ell/\mathbf Z_\ell)(-e))\) is
\[
 X\big/\big((u^eg)^{\ell^m}-1\big)X .
\]
\end{lemma}
\begin{proof}
Put \(A_e=(\mathbf Q_\ell/\mathbf Z_\ell)(-e)\).
The quotient \(\Gamma_m=\operatorname{Gal}(F_\infty/F_m)\) is
procyclic with generator \(\gamma^{\ell^m}\). Its first cohomology on
\(A_e\) is zero since \(u^{-e\ell^m}-1\ne0\) acts surjectively on
\(\mathbf Q_\ell/\mathbf Z_\ell\). Its second cohomology is zero since
\(\mathbf Z_\ell\) has \(\ell\)-cohomological dimension one.
Inflation--restriction identifies \(H^1(F_m,A_e)\) with the
\(\Gamma_m\)-invariant continuous homomorphisms from \(G_{F_\infty}\)
to \(A_e\), on which \(G_{F_\infty}\) acts trivially.
The tower is unramified away from \(\ell\), so the inertia subgroups
there are unchanged. At \(w\mid\ell\) the local cyclotomic image is an
open subgroup of \(\Gamma_m\), since \((F_m)_w/\mathbf Q_\ell\)
is finite. The same procyclic first-cohomology calculation gives
injectivity of local restriction. Thus the finite-level local
conditions are exactly vanishing on the defining subgroups of \(X\).
We obtain \(\operatorname{Hom}_{\rm cont}(X,\mathbf Q_\ell/
\mathbf Z_\ell)(-e)^{\Gamma_m}\). Duality turns invariants into
coinvariants, and the dual twist acts by \(u^e\). The image of the
continuous endomorphism of the compact group \(X\) is closed, so its
ordinary quotient gives those coinvariants.
\end{proof}
\noindent
\src{src/Development/AutomorphicArithmetic/AutomorphicVarieties/GaloisLibrary/KummerTheory/Mk0IdMapCongr.lean}{90}{local conditions over the tower},
\src{src/Development/AutomorphicArithmetic/AutomorphicVarieties/NumberFieldTowers/ExistsDescends.lean}{26}{descent with the positive dual twist}.

\begin{proposition}[Integral growth and centres]\label{newcoh:centres}
Suppose \(X\) has at most \(R\ge1\) generators over \(\mathbf Z_\ell\)
and \(g\) is unipotent on \(X/\ell X\). Set
\(j=\lfloor\log_\ell R\rfloor+3\) and
\(L_m(e)=\operatorname{length}_{\mathbf Z_\ell}
\big(X/((u^eg)^{\ell^m}-1)X\big)\), allowing infinite length.
There are \(R\) elements \(s_1,\ldots,s_R\in\mathbf Z_\ell\),
independent of \(e,m\), such that
\begin{align}
 L_m(e)&\le L_j(e)+R(m-j)_+, \label{newcoh:growth}\\
 L_m(e)&\le\sum_{i=1}^R
 \left(\max(m,j)+c+v_\ell(e-s_i)\right). \label{newcoh:centreformula}
\end{align}
For each integer \(Z\ge1\), at most \(R\) values of \(1\le e\le Z\)
need be removed to bound every \(v_\ell(e-s_i)\) by
\(\lfloor\log_\ell Z\rfloor\) for the remaining \(e\).
\end{proposition}
\begin{proof}
Write \(h=u^eg\) and \(k=\lfloor\log_\ell R\rfloor+1\).
On \(X/\ell X\) the dimension is at most \(R\) and
\((h-1)^R=0\). In characteristic \(\ell\),
\(h^{\ell^k}-1=(h-1)^{\ell^k}=0\) there. Thus
\((h^{\ell^k}-1)X\subseteq\ell X\). Raising this endomorphism
plus the identity to the \(\ell\)-th power shows
\((h^{\ell^{k+1}}-1)X\subseteq\ell^2X\) by binomial expansion.
This inclusion persists at higher layers.

For \(w=h^{\ell^n}\), \(n\ge j\), put \(P=1+w+\cdots+w^{\ell-1}\).
Then \(Px\equiv\ell x\bmod\ell^2X\), whence \(PX\subseteq\ell X\)
and \(PX+\ell^2X=\ell X\). Nakayama gives \(PX=\ell X\).
Consequently
\[
 (w^\ell-1)X=\ell(w-1)X,\qquad
 (w-1)X/(w^\ell-1)X=(w-1)X/\ell(w-1)X .
\]
The last module has dimension at most \(R\), since \((w-1)X\) is a
quotient of \(X\). This proves the claimed one-step loss even when
\(X\) has \(\ell\)-torsion. Addition of lengths proves
\eqref{newcoh:growth} for \(m\ge j\). For \(m<j\) use the
divisibility \(T^{\ell^m}-1\mid T^{\ell^j}-1\).

Choose a surjection \(\mathbf Z_\ell^R\to X\) and lift the images of
the generators under \(g^{\ell^j}\) to columns of an integral matrix
\(C\). Put \(a=u^{\ell^j}\). Then \(X/((u^eg)^{\ell^j}-1)X\)
is a quotient of \(\operatorname{coker}(a^eC-I)\). By Smith normal form
the length of the latter equals \(v_\ell(\det(a^eC-I))\), with
\(v_\ell(0)=+\infty\). Let \(z_i\) be the eigenvalues of \(C\)
in a splitting field, counted with multiplicity. For a unit
eigenvalue \(z\), choose \(s\in\mathbf Z_\ell\) minimizing
\(\lvert a^s z-1\rvert_\ell\). Such \(s\) exists by compactness and
the continuous extension of powers on the principal units. For
integers \(e\ge0\), the ultrametric inequality gives
\[
 |a^e-a^s|_\ell\le\max(|a^e z-1|_\ell,|a^s z-1|_\ell)
 = |a^e z-1|_\ell .
\]
The binomial valuation formula from Section~\ref{sup:length}, extended
by continuity from integers, gives
\(v_\ell(a^e-a^s)=j+c+v_\ell(e-s)\), including equality at infinity.
For a nonunit eigenvalue \(|z|_\ell<1\) the valuation of \(a^e z-1\)
is zero and any centre suffices. All eigenvalues are integral.
Factoring the determinant proves the sum bound at \(j\), and
\eqref{newcoh:growth} proves \eqref{newcoh:centreformula}. Finally
integers with \(v_\ell(e-s_i)>\lfloor\log_\ell Z\rfloor\) form at
most one residue class modulo \(\ell^{1+\lfloor\log_\ell Z\rfloor}>Z\);
this class meets \([1,Z]\) at most once.
\end{proof}
In Lemma~\ref{newcoh:desc} the action on the finite vector space
\(X/\ell X\) factors through a finite \(\ell\)-group, which gives
the required unipotence. The centres are attached to a chosen base
field. Only their number must be uniform as the auxiliary primes vary.
\src{src/Development/AutomorphicArithmetic/AutomorphicVarieties/AdditiveValuationNorms/ExistsCentresLevel.lean}{19}{integral tower inequalities},
\src{src/Development/AutomorphicArithmetic/AutomorphicVarieties/LayeredIsomorphisms/Exc.lean}{8}{exceptional set from centres}.

\paragraph{Uniformity for the moving quadratic fields.}
Let \(M/\mathbf Q\) be the fixed abelian target field, \(d_M=[M:\mathbf Q]\),
and let \(q\equiv3\bmod8\), \(d\equiv7\bmod16\) be distinct primes
coprime to \(2\) and to the chosen modulus of \(M\).
Put \(B=M(\mu_{2^\infty})\). The unramified elementary
two-characters of \(G_B\) have dimension at most a constant \(g_M\)
depending only on \(M\). Indeed, Ferrero--Washington gives
\(\mu=0\) for the fixed abelian base \(M(i)\). The cyclotomic
Iwasawa structure theorem then gives a bound for the two-ranks of the
class groups throughout this fixed tower. Class field theory gives
the same bound on unramified two-characters in its finite layers.
For the passage to the union, every finite collection of unramified
characters at infinite level descends unramified to one finite level.
To see this for one character, extend its locally constant function on
the closed subgroup \(G_B\) to a function factoring through a finite
quotient of \(G_{M(i)}\). On some finite layer its image subgroup in
this finite quotient is already that of \(G_B\), so the extension
there is a homomorphism. Only finitely many places of \(M(i)\) can
ramify in that quotient. For each of them choose one absolute inertia
group and its finitely many conjugates in the quotient. More precisely,
apply stabilization inside each conjugate of the absolute inertia
group to the descending compact intersections with the layer groups.
Their images eventually equal the image of the intersection with
\(G_B\). One representative per coset of the open normal kernel of
the finite quotient suffices, since the finite layer groups are normal
over \(M(i)\) and conjugation by that kernel acts trivially in the
quotient. The extension is therefore unramified on a sufficiently
high layer. Taking the maximum for a finite collection proves the
bound at infinite level.
See Ferrero--Washington, Ann. of Math. 109 (1979)~\cite{ferreroWashington}.
\src{src/Development/AutomorphicArithmetic/AutomorphicVarieties/PublishedTheorems/FerreroWashington.lean}{7}{finite-layer rank formulation},
\src{src/Development/AutomorphicArithmetic/AutomorphicVarieties/QuadraticExtensions/LevUAnti.lean}{9}{passage to the union}.

We compute explicitly how much this dimension can increase when
adjoining \(\sqrt{-q},\sqrt{-d}\). In the full two-power cyclotomic
tower over \(\mathbf Q\), the number of primes over an odd prime \(p\)
is the index of the closed subgroup generated by \(p\) in
\(\mathbf Z_2^\times\). For the two congruence classes at hand
\(v_2(q^2-1)=3\), \(v_2(d^2-1)=4\). The even powers of each
prime therefore generate \(1+2^3\mathbf Z_2\) and
\(1+2^4\mathbf Z_2\), respectively, by the binomial valuation formula.
Each prime itself is \(3\bmod4\), so adjoining its odd powers gives
indices \(2\) and \(4\) in \(\mathbf Z_2^\times\). Base change to
\(B\) multiplies the number of places by at most \(d_M\).
Allowing ramification at \(q,d\) thus raises the dimension of
two-characters by at most \(6d_M\). Indeed a character's restriction
to each new inertia group is specified by one bit, since wild inertia
has odd residue characteristic and the tame pro-two quotient is
procyclic. The kernel of these evaluations is unramified everywhere.

Each of the two quadratic adjunctions multiplies the resulting dimension
by at most two. To see this, quotient the absolute group by the closed
normal subgroup generated by inertia outside \(q,d\). Each adjunction
is unramified there, so defines a subgroup \(U\) of index at most two
in the resulting group \(V\). Normality of \(U\) and the use of all
conjugate inertia groups ensure that the same quotient inside \(U\)
imposes all the required inertia conditions. For index two let
\(W=H^1(U,\mathbf F_2)\). The inflation--restriction sequence gives
\(\dim W^{V/U}\le\dim H^1(V,\mathbf F_2)\), since the inflation kernel
in degree one contributes one dimension and the transgression has
target \(H^2(C_2,\mathbf F_2)\) of dimension one. On \(W\) the
operator \(\sigma-1\) for the nonidentity coset satisfies
\((\sigma-1)^2=0\), so \(\dim W\le2\dim W^{V/U}\).
It follows that the split unramified \(X\) for
\(M(i,\sqrt{-q},\sqrt{-d})\) has generator bound
\[
 R_2=4(g_M+6d_M).
\]
Indeed \((X/2X)^\vee\) is a subspace of the characters just counted,
and lifts of any basis of \(X/2X\) generate the compact abelian
pro-two group \(X\) over \(\mathbf Z_2\). The closed span of the
lifts has zero quotient by compact Nakayama; it equals their algebraic
\(\mathbf Z_2\)-span by compactness. The same fixed-field argument
before the quadratic adjunctions applies at each of the finitely many
fixed odd primes with its own abelian base containing \(\mu_\ell\).
Consequently the cardinality bound used to choose the window precedes
the choice of \(q,d\).
\src{src/Development/AutomorphicArithmetic/AutomorphicVarieties/QuadraticExtensions/IsAlgClosureRat.lean}{69}{moving two-rank and its exact constant},
\src{src/Development/AutomorphicArithmetic/AutomorphicVarieties/LayeredIsomorphisms/ThmBabOf.lean}{27}{sum of the uniform ranks}.

\subsection{An independent bound for the low layer, including the small primes}
\label{newcoh:ktheory}

\begin{theorem}[Uniform order of the low-layer group]\label{newcoh:lemmaK}
Let \(F\) be totally imaginary of degree \(n\), with absolute
discriminant \(D\). Let \(E\ge1\) and let \(\ell\) satisfy the
root-of-unity hypotheses of Section~\ref{sup:length}. With \(C=32\)
one has
\[
 H^2(G_{F,S_\ell},\mathbf Z_\ell(E+1))\ \text{is finite},\qquad
 \log\#H^2(G_{F,S_\ell},\mathbf Z_\ell(E+1))
 \le (2D)^{(2En)^C}.
\]
\end{theorem}
\begin{proof}
Put \(m=2E\), \(N=2m+1\ge5\), \(U=\operatorname{Spec}
\mathcal O_F[1/\ell]\), and \(A=K_m(\mathcal O_F)\{\ell\}\), where
braces indicate the primary subgroup. Quillen's finite generation and
Borel's rational rank formula imply that \(K_m(\mathcal O_F)\) is
finite. The Quillen--Lichtenbaum comparison for number rings gives
\[
 A\simeq H^2_{\mathrm{\acute et}}(U,\mathbf Z_\ell(E+1)).
\]
Here \(E+1\ge2\). At the prime two all correction terms arising from
real embeddings vanish. Inverting \(\ell\) in the \(K\)-group also
gives the same primary group, since the finite residue fields above
\(\ell\) have vanishing \(K_{2E}\) and \(K_{2E-1}\) of order
\((\#k_v)^E-1\). These comparisons hold before taking orders.
The rank calculation is Borel's Proposition~12.2~\cite{borelStable}.
For the comparison see the norm-residue theorem~\cite{voevodskyNorm},
Kahn's Theorem~1 at two~\cite{kahnTwo}, and Weibel, Chapter~VI,
Theorems~8.2 and 8.4~\cite{weibelKbook}.

For the desired upper bound an injection of Galois cohomology into
this group suffices. For each finite locally constant coefficient sheaf
on \(U\), the fundamental-group comparison in degree two is injective.
Indeed degree one on the universal pro-finite-\'etale cover vanishes,
because each torsor for a finite constant abelian group is trivialized
on a finite \'etale cover. The five-term sequence of that cover gives
the injection, and \(\pi_1(U)=G_{F,S_\ell}\). Apply this to
\(\mu_{\ell^a}^{\otimes(E+1)}\) and take inverse limits. The finite
degree-one groups on the number ring and its fundamental group satisfy
the Mittag--Leffler condition, so no derived limit in degree one enters.
Their finiteness also follows by trivializing the finite action in a
finite extension and applying class field theory to extensions of
bounded degree unramified outside \(S_\ell\).

\emph{The number of generators at the small primes.}
If \(r=\dim_{\mathbf F_\ell}A/\ell A\), reduction of the displayed
\'etale cohomology injects \(A/\ell A\) into
\(H^2_{\mathrm{\acute et}}(U,\mu_\ell^{\otimes(E+1)})\).
Choosing a primitive \(\ell\)-th root over \(F\) identifies the latter
coefficient sheaf with \(\mu_\ell\). By the Kummer sequence,
\[
 0\longrightarrow\operatorname{Pic}(U)/\ell
 \longrightarrow H^2_{\mathrm{\acute et}}(U,\mu_\ell)
 \longrightarrow\operatorname{Br}(U)[\ell]\longrightarrow0 .
\]
The class group surjects onto \(\operatorname{Pic}(U)\). The local
invariant sequence for the Brauer group embeds
\(\operatorname{Br}(U)[\ell]\) in the \(\ell\)-torsion of the
sum-zero subgroup of \(\bigoplus_{v\mid\ell}\mathbf Q/\mathbf Z\).
All archimedean places are complex. In particular
\[
 r\log\ell\le\log h_F+(n-1)\log\ell .                 \tag{*}
\]
Minkowski supplies each ideal class with an integral ideal of norm at
most \(2^n\sqrt D\). For \(y\ge1\) the number of such ideals of norm
at most \(y\) is at most \(y^2\zeta_F(2)\le y^2\zeta(2)^n\), since
each rational prime has at most \(n\) Euler factors of residue degree
at least one. Thus \(\log h_F\le\log D+O(n)\). Also
\(\ell\le n+1\). The constants introduced here are absolute.

\emph{The integral Hurewicz map and transfer.}
The identity component \(B\mathrm{GL}(\mathcal O_F)^+\) of the
algebraic \(K\)-theory infinite loop space is connected. Apply Mathew's
unstable Hurewicz theorem to this component in degree \(m\), with
connectivity parameter one. Its \(\ell\)-primary kernel is killed by
\(\ell^{m+3}\). Indeed his two bounds on the exponent of \(\ell\) are
\(\lceil(m-1)/2\rceil+3\) at two and
\(\lceil(m+2)/(2\ell-2)\rceil+1\) at odd primes. The plus
construction preserves integral homology. By homological stability
for the ring of integers (a Dedekind domain of stable rank at most
two), \(H_m(\mathrm{GL}(\mathcal O_F),\mathbf Z)\) identifies with
\(H_m(\Gamma,\mathbf Z)\) for
\(\Gamma=\mathrm{GL}_N(\mathcal O_F)\) and \(N=2m+1\).
See Mathew, Theorem~6.3~\cite{mathewHurewicz}, van der Kallen,
Theorem~4.11~\cite{vanderKallen}, and Bayer-Fluckiger, Emery and
Houriet, \S3.1~\cite{bayerEmeryHouriet}.

Let \(\Gamma'\) be the kernel of reduction modulo \(3\) in \(\Gamma\)
and \(I=[\Gamma:\Gamma']\le3^{nN^2}\). The group \(\Gamma'\) is
torsion-free. In the action on the underlying integral lattice, a
nonidentity element in it has the form \(1+3^tB\) for \(t\ge1\) and
an integral matrix \(B\) nonzero modulo three. Raising to a prime
exponent \(p\ne3\) leaves the first nonzero
three-adic term at \(t\); cubing puts it exactly at \(t+1\), by the
binomial formula. In particular there can be no element of prime order.
Compose the Hurewicz map with transfer to
\(H_m(\Gamma',\mathbf Z)\). If \(x\in A\) is in the kernel of this
composite, the pushforward of its transfer gives \(Ih(x)=0\), so
\(\ell^{m+3+v_\ell(I)}x=0\). A subgroup of a finite abelian
\(\ell\)-group on \(r\) generators killed by \(\ell^a\) has at most
\(\ell^{ar}\) elements, by the elementary divisors. We have proved
\begin{equation}\label{newcoh:transfer}
 \log\#A\le
 (m+3+v_\ell(I))\,r\log\ell+
 \log\# H_m(\Gamma',\mathbf Z)_{\mathrm{tors}} .
\end{equation}
This bounds the transfer loss also for primes dividing \(I\).

\emph{Bounding the torsion of the congruence subgroup.}
Use Ash's contractible well-rounded retract of the space of positive
archimedean Hermitian forms on the free lattice \(\mathcal O_F^N\).
Normalize the minimum to one. Each cell has finite stabilizer and is
determined by the set of its minimal vectors \(v\), those with
\(Q(v)=1\). For a specified minimal set the stratum is the relative
interior of the convex face imposed by its equalities in the region
\(Q(w)\ge1\) for every nonzero lattice vector \(w\). Its boundary
adds minimal vectors. This polyhedral structure gives the subdivision
by flags of faces, ordered by reverse inclusion of their minimal
sets. Put
\(b=nN^2\), which bounds the dimension of the real vector space of
Hermitian forms. By Bayer-Fluckiger--Emery--Houriet there exists a
finite \(\Phi\subset\mathcal O_F^N\) into which every one of these
minimal sets can be moved by \(\Gamma\), with
\[
 |\Phi|\le N^{3N^2n^2} n^{5N^3n^4} D^{9N^3n^3}.
\]
Their result applies to free lattices, \(N\ge5\) and \(F\ne\mathbf Q\).
These all hold, including for \(E=1\).
See Bayer-Fluckiger--Emery--Houriet, Thm.~2.2 and Cor.~5.2~\cite{bayerEmeryHouriet}.

Here is a count for the barycentric subdivision which allows redundant
minimal-vector equations. In any flag move the minimal-vector set of
its smallest face into \(\Phi\). Every other face in the flag then
has its minimal-vector set in \(\Phi\). A single such set \(T\) is
encoded by at most \(b+1\) elements of \(T\). Choose a basis for the
span of the affine functions \(Q\mapsto Q(v)-1\), \(v\in T\), on
the \(b\)-dimensional ambient space (or its actual smaller space).
The equalities for this basis recover all minimal vectors at that
face. In fact for the affine solution space of these equations, each
vector in \(T\) has \(Q(v)=1\) identically, whereas for \(v\notin T\)
the chosen interior point of the face has \(Q(v)>1\) when \(v\ne0\).
Thus distinct faces cannot share the same encoding. A flag has at
most \(b+1\) members. Encoding each choice as a padded list from
\(\Phi\) with an extra termination symbol gives the bound
\[
 C_{\rm cell}=(|\Phi|+2)^{(b+2)^2}
\]
for the number of \(\Gamma\)-orbits of simplices in the subdivision.
Consequently the quotient by \(\Gamma'\) has at most \(I C_{\rm cell}\)
simplices in every degree. The free action on the contractible
complex gives a classifying space for \(\Gamma'\).

For the boundary in degree \(m+1\), each column has Euclidean norm
at most \(m+2\), since there are \(m+2\) faces counted with sign;
possible identifications in the quotient respect this upper bound.
The torsion order of the cokernel of an integral matrix of rank \(s\)
is the gcd of its full-rank minors. Hadamard bounds every such minor
by the product of \(s\) column norms. The torsion in homology embeds
in the torsion of this cokernel, so
\[
 \log\#H_m(\Gamma',\mathbf Z)_{\mathrm{tors}}
 \le I C_{\rm cell}\log(m+2).
\]
Finally put \(P=2En\ge4\), \(Y=\log(2D)\ge\log2\).
Then \(N=4E+1\le2P\), \(n,m\le P\), \(b\le P^4\), and
\(\log I\le2P^4\). Using \(\log N\le2P,\ \log n\le P\),
the displayed bound for \(\Phi\) gives
\(\log|\Phi|\le24P^5+40P^8+72P^6Y\le P^{12}Y\).
Thus \(\log C_{\rm cell}\le(2P^4)^2(2P^{12}Y)\le P^{22}Y\).
By \((*)\), \(r\log\ell\le\log D+n\log(4\zeta(2))+n\log(n+1)
\le 6P^2Y\), while \(m+3+v_\ell(I)\le5P^4\).
Equation~\eqref{newcoh:transfer} now gives
\[
 \log\#A\le30P^6Y+\exp(2P^4+P^{22}Y)\log(m+2)
 \le\exp(P^{32}Y).
\]
For the last inequality each summand is at most
\(\exp(P^{24}Y)\), and \(2\exp(P^{24}Y)\le\exp(P^{32}Y)\)
throughout \(P\ge4\). This proves the estimate.
\end{proof}

Theorem~\ref{newcoh:lemmaK} applies directly over the field at the
fixed layer \(j\) in \eqref{newcoh:growth}. No enlargement of that
field or uniformity loss in a cyclotomic descent is needed in this
proof of the low-layer estimate. Poitou--Tate identifies the dual of
the strict degree-one kernel with the kernel of the degree-two
localization map, hence \(L_j(e)\le
\operatorname{length}H^2(G_{F_j,S_\ell},\mathbf Z_\ell(e+1))\).
Inserting the degree and discriminant bounds of Section~\ref{sup:length}
therefore gives the exceptional contribution calculated there. The
original proof reaches the same bound through the integral homology
cycles in \eqref{sup:lengthcount}; the proof here uses a different
set of classical theorems.

\paragraph{The primitive homology count in the original argument.}
One can also derive its factor \(q^{r_2}T_{2E+1}T_{2E}\) directly.
Let \(d=2E+1\), \(q=\ell^\nu\), and take the stable integral homology of
\(\mathrm{GL}(\mathcal O_{L'})\). A finite collection of mod-\(q\)
cycles of varying matrix sizes can be moved into one stable matrix
size, by stabilization followed by the natural map. Suppose their
joint evaluation \(\Phi\) is a homomorphism on stable mod-\(q\)
homology taking values in a \(\mathbf Z/q\)-module and satisfying
\(\Phi\psi_b=b\Phi\), where \(\psi_b\) denotes the \(b\)-fold
block-diagonal operation. These are precisely the boundary,
stabilization and direct-sum requirements on the primitive families.
Use integral stability in degrees \(d,d-1\) to write \(T_i\) for the
torsion order in either stable homology or matrix size \(2d+4\).

Consider \(P=\prod_{j=2}^{d}(\psi_\ell-\ell^j)\) on rational
homology in degree \(d\), induced by its integral counterpart. Rational
homology of the connected infinite loop space is the free graded
commutative algebra on its rational homotopy under Hurewicz, by the
rational Hurewicz theorem for connected homotopy-commutative
\(H\)-spaces. The product is block sum; its \(b\)-fold power operation
acts on a monomial of \(a\) positive-degree primitives by \(b^a\).
Since \(a\le d\), the image of \(P\) is contained in the primitive
subspace. For the totally imaginary field \(L'\), Borel's rank theorem
gives dimension \(r_2(L')\) in this odd degree \(d\ge3\).
Thus the integral image has rank at most \(r_2(L')\), and its torsion
has order at most \(T_d\). For any integral class \(y\),
\[
 \Phi(\overline{Py})=
 \ell^{d-1}\prod_{j=2}^{d}(1-\ell^{j-1})\,\Phi(\bar y).
\]
The product is a unit modulo \(q\). Hence the image of
\(\ell^{d-1}\Phi\) on reductions of integral classes has size at most
\(q^{r_2(L')}T_d\). The coefficient exact sequence identifies the
quotient of \(H_d(-,\mathbf Z/q)\) by integral reductions with
\(H_{d-1}(-,\mathbf Z)[q]\), of order at most \(T_{d-1}\).
Evaluating on representatives for this quotient multiplies the number
of possibilities by at most \(T_{d-1}\). Thus on \emph{all} mod-\(q\) cycles
\[
 \#\operatorname{im}\bigl(\ell^{d-1}\Phi\bigr)
 \le q^{r_2(L')}T_dT_{d-1}.
\]
This includes cycles with nonzero Bockstein and applies to the joint
product of residue evaluations used in \eqref{sup:lengthcount}. The
scalar is essential when \(q=\ell\) and \(d\ge2\), for then every
scaled evaluation is zero regardless of the unscaled image.
See Borel, \S\S10.6 and 12.1--12.2~\cite{borelStable}.
\src{src/Development/AutomorphicArithmetic/AutomorphicVarieties/ModularityArgument/ClassCounting/SUpperOf.lean}{41}{scaled map on all cycles and count by Bockstein fibres}.

The residual power test in that argument has a short field-theoretic
proof. Let \(F/L\) be finite, \(\mu_q\subset F\), and
\(\beta\in F^\times\). Suppose \(\beta\) is a \(q\)-th power in the
residue fields at all but finitely many primes of \(F\) of relative
residue degree one over \(L\). Then \(\beta\in F^{\times q}\).
Indeed, if not, \(F(\beta^{1/q})/F\) is a nontrivial cyclic Kummer
extension. Take the normal closure over \(L\) and an automorphism
fixing \(F\) which induces a nonidentity element of the Kummer
group. Chebotarev gives primes of \(L\) outside the exclusions and
prolongations with exactly this Frobenius. The primes of \(F\) below
these prolongations have residue degree one, since Frobenius fixes
\(F\) pointwise. It multiplies any \(q\)-th root of \(\beta\) by the
same nontrivial root of unity. Away from the primes of \(q\), \(\beta\)
and its denominator, these roots remain distinct on reduction, so
the residue polynomial has no root in the residue field of \(F\).
This contradiction proves the assertion.
\src{src/Development/EtaleCohomology/CyclotomicDiscriminants/SDetectCut.lean}{101}{residual power criterion}.

The original route additionally constructs cycles whose evaluations
encode twisted norms of actual divisors; the abstract evaluation-count
inequality just proved requires such an encoding to bound class-group
characters. Theorem~\ref{newcoh:lemmaK} yields the required cohomological
conclusion independently of that construction. The higher-rank
automorphic estimates enter elsewhere. For example, when \(N\) is
the host modulus, \(\psi(-1)=(-1)^e\), \(2\le e<t\), and the
prime-to-\(\ell\) part of \(\psi\) after deleting its local
\(\ell\)-component has even order, case II applies if either
\(\ell\) lies in the host's fixed exceptional prime set or
\(\ell-1\mid2e\). At odd \(\ell\), for the admitted hosts the
required estimate is finiteness of the strict length at level
\(\ell N\) together with
\[
 \frac{\log\ell}{e(O/\mathbf Z_\ell)}
 \operatorname{length}_O S_{\ell N}
   \big(\mathbf Q,(\operatorname{Frac}O/O)(\psi^{-1}\chi_\ell^{-e})\big)
 \le (\ell N e)^{C_\chi}.
\]
Here \(S_{\ell N}\) imposes zero on all inertia and on \(D_v\)
for \(v\mid\ell N\), \(\psi\) occurs in the host character set,
and \(C_\chi\) depends only on the fixed target. Case III requires the
same inequality away from the exceptional prime set when
\(\ell-1\mid2e\) and \(\psi\ne1\) has odd order, for the applicable
positive twists. The cited proof uses the adjoint lifting theorem,
the symplectic rank-four counting result, and the residual Bernoulli
test. The low-layer theorem has a different dependence on \(e\) and
\(N\). Propositions~\ref{famsel:hostinertia} and~\ref{famsel:splitray}
prove vanishing in further cases at every congruence depth. The
complementary budget in Section~\ref{famsel:complement} uses the
cited polynomial bounds on the remaining entries.
\src{src/Development/AutomorphicArithmetic/AutomorphicVarieties/HeightAndVanishing/Counts.lean}{60}{quantified normalized bounds},
\src{src/Development/AutomorphicArithmetic/AutomorphicVarieties/ModularityArgument/ThresholdArguments/DenseThresholdCases/AdjointLiftingOfCohomology.lean}{33}{automorphic and residual inputs for cases II and III}.

There is a parity obstruction to using the ordinary special value as
the nonzero integer in this comparison. For either primitive inducer
\(\varphi\) of \(\psi\) or \(\psi^{-1}\), let \(f\) be its conductor
and \(m=e+1\). If \(f>1\), in the expression
\[
 B_{m,\varphi}=f^{m-1}\sum_{a=1}^{f}\varphi(a)B_m(a/f)
\]
the substitution \(a\mapsto f-a\) on units multiplies the sum by
\(\varphi(-1)(-1)^m=-1\), since \(B_m(1-x)=(-1)^m B_m(x)\).
Thus \(B_{m,\varphi}=0\). At \(f=1\), \(e\) is even and positive,
so \(m\ge3\) is odd and \(B_m(1)=B_m(0)=0\), by the same reflection
and \(B_m(1)-B_m(0)=0\) for \(m>1\). Hence \(L(-e,\varphi)=0\)
in these branches. Valuations of this value supply no finite bound
for the strict lengths. The trace products below require an actual
nonzero integer.

\begin{proposition}[Finiteness and prime support before the effective estimate]
\label{newcoh:fullsupport}
Fix \(N,e\ge1\), let \(F=\mathbf Q(\mu_{4N})\), \(n=[F:\mathbf Q]\),
and put \(K_e=K_{2e}(\mathcal O_F)\). At each prime \(\ell\) choose
any finite coefficient extension \(E_\ell/\mathbf Q_\ell\), with
integers \(O_\ell\) and ramification index \(e_\ell\), and any
\(O_\ell^\times\)-valued Dirichlet character \(\psi_\ell\) of
modulus \(N\). Let \(S_\ell\) be the strict kernel with coefficients
\((E_\ell/O_\ell)(\psi_\ell^{-1}\chi_\ell^{-e})\), trivial on every
inertia group and locally trivial at the decomposition groups above
\(\ell N\), in the sense of restriction of cohomology classes.
Each \(S_\ell\) is finite. It vanishes when
\(\ell\nmid n\,\#K_e\). Moreover
\[
 \sum_{\ell}\frac{\log\ell}{e_\ell}\operatorname{length}_{O_\ell}S_\ell
 \ \le\ \log\#K_e+n\log(4Nne).
\]
If one sums over a family of at most \(\varphi(N)\) such characters
at each prime, the right side is multiplied by \(\varphi(N)\).
\end{proposition}
\begin{proof}
Over \(F\) each \(\psi_\ell\) is trivial. On inertia away from \(\ell\)
the remaining cyclotomic twist also acts trivially. The restriction
of a strict cocycle therefore vanishes on each such inertia group,
since a coboundary there is zero. It factors through the quotient
unramified outside \(\ell\), and its class restricts to zero on the
decomposition groups at \(\ell\). There are no real places of \(F\).
Choose a \(\mathbf Z_\ell\)-basis of \(O_\ell\) and put
\(d_\ell=[E_\ell:\mathbf Q_\ell]\). The restricted coefficient module
is a sum of \(d_\ell\) copies of \((\mathbf Q_\ell/\mathbf Z_\ell)(-e)\).
Poitou--Tate identifies the dual of the strict kernel for one copy
with the kernel of localization in
\(H^2(G_{F,S(\ell)},\mathbf Z_\ell(e+1))\), where \(S(\ell)\)
is the set of primes of \(F\) above \(\ell\).
The degree-two injection, localization and Quillen--Lichtenbaum
comparison at the start of the proof of Theorem~\ref{newcoh:lemmaK}
embed this kernel in \(K_e\{\ell\}\). Those steps only require that
\(F\) be totally imaginary and \(e\ge1\). Equivalently, the even
case of Weibel's Theorem~VI.8.4 gives
\(K_e\{\ell\}\simeq H^2_{\mathrm{et}}(\mathcal O_F[1/\ell],
\mathbf Z_\ell(e+1))\), also for \(\ell=2\)~\cite{weibelKbook}.
The root-of-unity assumption entered the subsequent generator count.
Thus the order of the restriction image is at most
\(\#(K_e\{\ell\})^{d_\ell}\).

On the restriction kernel, restriction followed by corestriction
gives multiplication by \(n\), so the kernel vanishes if
\(\ell\nmid n\). For \(\ell\mid n\) write
\(r=\max(1,v_\ell(4N))\); at two this is at least two.
The image of \(G_F\) under \(\chi_\ell\) contains \(1+\ell^r\).
Indeed a cyclotomic automorphism with that \(\ell\)-adic component
and identity at every other prime fixes \(\mu_{4N}\), by the Chinese
remainder theorem. The binomial valuation formula gives
\(v_\ell((1+\ell^r)^e-1)=r+v_\ell(e)\), including at two since
\(r\ge2\) there. The \(G_F\)-invariant subgroup of the original
coefficient module therefore has order at most
\(\ell^{d_\ell(r+v_\ell(e))}\).
By inflation--restriction the remaining kernel injects into the
first cohomology of \(\operatorname{Gal}(F/\mathbf Q)\) with these
invariants. Counting all functions on that group bounds its order
by the \(n\)-th power of the displayed bound. This proves the claimed
finite support since \(K_e\) is finite.

For any finite \(O_\ell\)-module \(M\),
\(\log\#M/d_\ell=(\log\ell/e_\ell)\operatorname{length}_{O_\ell} M\).
After this normalization the restriction image contributes at most
\(\sum_\ell\log\#(K_e\{\ell\})=\log\#K_e\).
The kernel contributes at most \(n\log(4Nne)\), since the product
of the distinct primes dividing \(n\) divides \(n\), the product
of their \(\ell^{v_\ell(4N)}\) divides \(4N\), and their
\(\ell^{v_\ell(e)}\) have product at most \(e\). This estimate
is independent of the character and the coefficient field.
\end{proof}

\begin{corollary}[Summing distinct characters before bounding their lengths]
\label{newcoh:tamepool}
Retain the notation of Proposition~\ref{newcoh:fullsupport}. Regard
the characters at \(\ell\) as distinct characters of
\((\mathbf Z/N\mathbf Z)^\times\) in one algebraic closure of
\(\mathbf Q_\ell\). Put \(b_\ell=\ell^{v_\ell(\varphi(N))}\),
\(a_\ell=\max(1,v_\ell(4N))+v_\ell(e)\), and let \(s_\ell\)
be the least number of generators of the Sylow \(\ell\)-subgroup of
\(G=\operatorname{Gal}(F/\mathbf Q)\), taking zero for the trivial
subgroup. For any subset of these characters one has
\[
 \sum_{\psi}
  \frac{\log\ell}{e_{\ell,\psi}}\,
     \operatorname{length}_{O_{\ell,\psi}} S_{\ell,\psi}
 \le b_\ell\bigl(\log\#K_e\{\ell\}
                     +s_\ell a_\ell\log\ell\bigr).
\]
In particular, at every \(\ell\nmid [F:\mathbf Q]\) the sum over
\emph{all} the characters costs at most \(\log\#K_e\{\ell\}\).
\end{corollary}
\begin{proof}
Enlarge all coefficients to a common finite extension with integers
\(O\), splitting the finite abelian group \(G\). A further extension
of discrete valuation rings is finite free. Continuous cochains and
the kernels of their restriction maps commute with this scalar
extension: after choosing a module basis the underlying operation
is a finite direct sum. The length in the larger integers is the
old length times the relative ramification index, by a composition
series. Thus the normalized lengths are unchanged.

Write \(G=P_\ell\times H\), with \(P_\ell\) its Sylow \(\ell\)-subgroup
and \(H\) of order prime to \(\ell\). Over \(F\) use the single strict
kernel \(M\) with coefficients \((\mathbf Q_\ell/\mathbf Z_\ell)(-e)\)
and localizations at \(\ell\), as in the preceding proof.
It is a finite \(G\)-module of order at most \(\#K_e\{\ell\}\).
Untwisting the restriction of \(S_{\ell,\psi}\) puts its image in
the \(\psi|_H\)-part of \(M\otimes_{\mathbf Z_\ell}O\): the restricted
class was invariant for the action twisted by \(\psi^{-1}\).
These parts for different characters of \(H\) form a direct sum
by the character projectors \( |H|^{-1}\sum_{h\in H}\eta(h)^{-1}h\).
At most \(b_\ell\) of the original characters have a given restriction
to \(H\), since \(G\) surjects onto \((\mathbf Z/N\mathbf Z)^\times\)
and the characters of this quotient trivial on the image of \(H\)
are its \(\ell\)-power order characters. Dividing log orders by
\(d_\ell=[\operatorname{Frac}O:\mathbf Q_\ell]\), the total image
contribution is at most \(b_\ell\log\#K_e\{\ell\}\).

For the kernel put \(U=((\mathbf Q_\ell/\mathbf Z_\ell)(-e))^{G_F}\).
It is cyclic of order at most \(\ell^{a_\ell}\) by the cyclotomic
calculation above. If it is nonzero, \(H\) acts on it through one
prime-to-\(\ell\) character. Indeed the kernel of reduction of the
units of \(\mathbf Z/\ell^j\mathbf Z\) modulo \(\ell\) is an
\(\ell\)-group, and the prime-to-\(\ell\) roots have unique lifts.
Thus \((U\otimes O)(\psi^{-1})^H\), with invariants taken after
twisting, is zero except for at most \(b_\ell\) of the characters.
Inflation--restriction identifies the cohomology in degree one of
\(G\) on this twisted module with that of \(P_\ell\) on its
\(H\)-invariants, since \(|H|\) is invertible. A cocycle on \(P_\ell\)
is determined by its \(s_\ell\) values on generators, by the cocycle
identity. The normalized log order for each of the remaining kernels
is therefore at most \(s_\ell a_\ell\log\ell\). Adding the two
contributions proves the assertion. If \(\ell\nmid n\), then
\(s_\ell=0\) and \(b_\ell=1\) since \(\varphi(N)\mid n\).
\end{proof}

This proves qualitative finiteness for every character under
discussion. To get the required size in the moving windows by this
route one would need a sufficiently sharp uniform bound as \(N,e\)
vary. Finiteness of the single group \(K_e\) for each pair \(N,e\)
alone provides no such rate. Theorem~\ref{famsel:rate} gives an explicit
all-prime bound, and Corollary~\ref{famsel:subwindow} identifies the
part of the required budget it covers.

\subsection{Two vanishing ranges in case III and the integral cotangent inequality}
\label{newcoh:oddbranch}

In case III the conductor of the odd-order character \(\psi\) divides
the fixed target conductor \(q\). Indeed, its component at \(2\) has
both odd order and \(2\)-power order. At either auxiliary prime away
from \(q\) the host condition says that the square of the component
is trivial. At a prime dividing \(q\) the host allows no higher
conductor there than the local factor of \(q\) (the auxiliary primes
occur to exponent one). The Chinese remainder theorem proves the
claim. Thus only finitely many primitive characters occur as the host
varies. The following criterion applies before any conditions at the
auxiliary level are imposed.
\src{src/Development/AutomorphicArithmetic/AutomorphicVarieties/HeightAndVanishing/HostOdd.lean}{70}{odd-order host characters}.

\begin{proposition}\label{newcoh:oddzero}
Let \(\psi\) have conductor \(f\), odd order \(m>1\), and fixed field
\(K=\overline{\mathbf Q}^{\ker\psi}\). Let \(\ell\ge5\) be prime
to \(mf\), let \(e>0\) be even with \(\ell-1\mid2e\), and set
\(s=\operatorname{ord}_m(\ell)\). Let \(O\) be the ring of integers
of a finite extension of \(\mathbf Q_\ell\) containing the values of
\(\psi\), and put \(A=(\operatorname{Frac}O/O)(\psi^{-1}\chi_\ell^{-e})\).
In either of the following situations \(S_\ell(\mathbf Q,A)=0\).
\begin{enumerate}
\item \(\ell-1\mid e\), \(\psi(\ell)\ne1\), and \(\ell^s\nmid h_K\).
\item \(\ell-1\nmid e\) and \(\ell^s\nmid h_{K(\sqrt{\ell})}\).
In particular this holds, for fixed \(K\), for all sufficiently large
\(\ell\) with \(s>m/2\).
\end{enumerate}
Here \(h_L\) denotes the ordinary class number of \(L\).
\end{proposition}
\begin{proof}
Write \(\theta=\psi^{-1}\chi_\ell^{-e}\), let \(\pi\) be a uniformizer,
and set \(k=O/(\pi)\), \(\alpha=\bar\theta\) and
\(L=\overline{\mathbf Q}^{\ker\alpha}\).
Any nonzero strict group contains a nonzero class \(x\) killed by
\(\pi\), by the finite-image argument in Proposition~\ref{newcoh:zero}.
For a representative \(z\) we have \(\pi z=\delta b\). Divide \(b\)
by \(\pi\) in the divisible module \(A\) and subtract the resulting
coboundary from \(z\). This gives a representative valued in
\(A[\pi]\simeq k(\alpha)\).

In the first case \(\alpha\) is nontrivial on \(D_\ell\) by the
Frobenius hypothesis. In the second it is nontrivial on inertia,
since \(\psi\) is unramified there and the residual Tate character
is quadratic. In both cases \(A^{D_\ell}=0\), since some element
acts by a scalar whose difference from one is a unit. On \(D_\ell\)
write \(z=\delta b_\ell\) by strictness. Then \(\pi b_\ell\) is
invariant and \(b_\ell\in A[\pi]\). Consequently \(z\) vanishes on
\(D_\ell\cap G_L\). For \(p\ne\ell\) the scalar on
\(I_p\cap G_L\) is already one in \(O\), since the cyclotomic factor
is unramified and reduction is injective on the \(m\)-th roots of
unity. The inertia condition again gives \(z=0\) there. These
conclusions hold at every prolongation, by the conjugation identity
in Proposition~\ref{newcoh:zero}. The restriction to \(G_L\)
therefore factors through \(\operatorname{Cl}(L)\), since real
inertia has order at most two. It is nonzero by inflation--restriction,
as \([L:\mathbf Q]\mid2m\) is prime to \(\ell\). Its image is an
additive subgroup of \(k\) stable under \(\alpha(G_{\mathbf Q})\)
by the cocycle identity. This cyclic scalar group contains the
\(m\)-th roots of unity, since the two factors have coprime orders.
The image is thus a nonzero vector space over
\(\mathbf F_\ell(\mu_m)\), giving \(\ell^s\mid h_L\).

In the first case \(L=K\). In the second, the residual quadratic
factor cuts out \(\mathbf Q(\sqrt{\ell})\), since even \(e\) forces
\(\ell\equiv1\bmod4\). Coprimality of the character orders gives
\(L=K(\sqrt{\ell})\), of degree \(2m\). The tensor product of
integral bases generates an order of discriminant
\(|D_K|^2\ell^m\) in \(L\). Minkowski supplies an ideal in each
class of norm at most \(y=2^m|D_K|\ell^{m/2}\).
There are at most \(\tau_{2m}(n)\) ideals of norm \(n\), by comparison
of each Euler factor with \((1-p^{-z})^{-2m}\) coefficientwise.
Summing the divisor expansion, with the final variable counted by
its upper bound, gives
\[
 h_L\le y(1+\log y)^{2m-1}.
\]
If \(s>m/2\), this is less than \(\ell^s\) for all sufficiently
large \(\ell\), uniformly in \(e\). For example, cubic characters
give the second range at primes \(\ell\equiv5\bmod12\).
\end{proof}

For the remaining prime and weight combinations the rank-four
argument compares the strict group with integral derivations.
We give the length inequality in the form needed for this comparison.
Let \(O\) be the ring of integers of a finite extension of \(\mathbf Q_\ell\),
\(B=\bigoplus_i B_i\) a graded \(O\)-algebra, and \(x:B\to O\)
an \(O\)-point, \(J=\ker x\), and
\(U=\{b\in B:\exists\,0\ne t\in O,\ tb\in J^2\}\).
Fix \(j\) with \(x(B_j)=0\) and equations \(F_i\in B_j\).
In \(B/U\) let \(C_j\) be the image of \(B_j\), and put
\(Q_j=C_j/\langle\bar F_i\rangle_O\).

\begin{lemma}\label{newcoh:cotangent}
Suppose \(H\) is an \(O\)-module, \(0\ne a\in O\), and its elements
index distinct \(O\)-linear maps \(D_h:B\to O/(a)\) killing all
\(F_i\). Suppose each \(D_h\) lifts to an \(O\)-valued derivation
at \(x\) vanishing on \(B_i\) for every \(i\ne j\). Then
\(\operatorname{length}_O H\le\operatorname{length}_O Q_j\).
\end{lemma}
\begin{proof}
The Leibniz rule kills \(J^2\) under an integral lift, and
torsion-freeness of \(O\) then kills \(U\). Thus each \(D_h\)
induces a functional \(Q_j\to O/(a)\). Its values on \(B_j\)
determine its values on all of \(B\) by the weight assumption,
so these functionals remain distinct. If \(Q_j\) has finite length,
write its elementary divisors as \(O/(\pi^{b_i})\). The dual with
values in \(O/(a)\) has length
\(\sum_i\min(b_i,v_\pi(a))\). Comparing cardinalities now proves
the inequality, since a finite \(O\)-module of length \(b\) has
\(|k|^b\) elements. Infinite right-hand length gives the assertion
automatically.
\end{proof}
\src{src/Development/AutomorphicArithmetic/AutomorphicVarieties/CommutativeAlgebra/GradedModulesLength/Mprime.lean}{38}{the saturated integral cotangent inequality}.

To apply this comparison let \(H\) be the strict group for the
scalar \(\theta\). Fix \(b\ne0\) and work on its subgroup \(H[b]\).
Suppose \(\theta(g_0)-1\) is a unit. Every class in this subgroup
then has a unique representative \(z\) with \(z(g_0)=0\), obtained
by subtracting \(\delta(z(g_0)/(\theta(g_0)-1))\).
Uniqueness gives additivity and \(bz=0\) as a cocycle. On each of
finitely many designated decomposition groups choose \(s_v\) with
\(\theta(s_v)\ne1\). If \(z|_{D_v}=\delta u_v\), the equality
\(bz(s_v)=0\) shows that \(b(\theta(s_v)-1)u_v=0\). Hence for
\[
 a=b\prod_v(\theta(s_v)-1)
\]
all representatives and all local witnesses have values in
\(A[a]\simeq O/(a)\). This yields an additive injection of \(H[b]\)
into finite-coefficient cocycles satisfying the prescribed local
coboundary conditions. The denominator \(a\) may grow with \(b\).
\src{src/Development/AutomorphicArithmetic/AutomorphicVarieties/SymplecticGSp4/ShaCoc.lean}{112}{common denominator for the strict local witnesses}.

Section~\ref{ct:section} constructs the two square-zero representations
and, from the specified relative presentations, sends their cocycles to
reductions of integral weight \(-1\) derivations killing the equations modulo \(a\).
In the lower left matrix entry with indices \((3,1)\) in one-based
notation the derivative recovers \(z(g)\) times a unit for each \(g\).
This makes the assignment injective, so Lemma~\ref{newcoh:cotangent}
bounds the length of \(H[b]\) by that of \(Q_{-1}\), independent
of \(b\). The subsequent geometric inequality is
\(\operatorname{length}_O Q_{-1}\le2\sum_\xi v_\pi(\xi(a_0))\),
where \(\xi\) runs through a finite cover of the relevant orbits
and \(a_0\) interpolates a trace difference. Applying it still
requires the integral deformation model, rigidity of the relevant
irreducible points, and simultaneous nonvanishing of that trace
difference. In particular the extension groups characterized by
Bella\"iche and Chenevier for a fixed Cayley--Hamilton quotient
are the extensions realized by that quotient. Their theorem alone
does not realize every strict class in this rank-four model.
See Bella\"iche and Chenevier, Theorems 1.5.5--1.5.6~\cite{bellaicheChenevier}.
\src{src/Development/AutomorphicArithmetic/AutomorphicVarieties/SymplecticGSp4/MPrime.lean}{102}{weight minus one and the cocycle entry},
\src{src/Development/AutomorphicArithmetic/AutomorphicVarieties/SymplecticGSp4/GaloisRepIneqOfAllIrreducibleNonvanishing.lean}{152}{rigidity and the orbit valuation inequality}.

\subsection{Strict cocycles as integral tangent functionals}
\label{ct:section}

Let \(O\) be the integers in a finite extension of \(\mathbb Q_\ell\),
\(0\ne a\in O\), \(M=O/aO\), and \(\Gamma=G_{\mathbb Q}\).
Write \(O\ltimes M=O\oplus\epsilon M\), with \(\epsilon^2=0\).
This is a complete local \(O\)-algebra. Its quotients
\((O/\pi^h)\ltimes M\), \(h\ge v_\pi(a)\), are Artin rings and their
inverse limit is \(O\ltimes M\). For a continuous augmented
\(O\)-algebra \(B\xrightarrow{x}O\), a continuous map to this
extension over \(x\) has the unique form \(x+\epsilon d\), with
\(d(fg)=x(f)d(g)+x(g)d(f)\).

Use one-based indices for matrix units. The two diagonal points
\(r_0\), their infinitesimal directions \(X\), and the weights \(j\)
of their nonzero entries are as follows:
\[
\begin{array}{c|c|c|c}
 n&\text{matrix of }r_0& X& j\\ \hline
 3&\operatorname{diag}(\Psi,1,\Psi^{-1})&E_{21}-E_{32}&-1\\
 4&\operatorname{diag}(\alpha\Psi,\Psi,\alpha,1)&E_{31}+E_{42}&(0,-1).
\end{array}
\]
In rank three one appends the sign character \(\eta\) to the matrix,
in the semidirect product in which the negative sign acts by
\(g\mapsto Jg^{-t}J\), \(J\) the antidiagonal matrix of ones. Its
differential \(Y\mapsto-JY^tJ\) fixes \(X\).
In rank four the symplectic adjoint sends \(X\) to \(-X\);
the diagonal point has similitude \(\alpha\Psi\).
In both ranks \(\operatorname{Ad}(r_0(g))X=\theta(g)X\) with
\(\theta=\Psi^{-1}\). The rank-four bidegrees of the four basis
vectors are \((1,1),(0,1),(1,0),(0,0)\).

\begin{proposition}[The square-zero representation]\label{ct:point}
Let \(z:\Gamma\to M\) be a continuous cocycle for \(\theta\).
Then \(r_z(g)=(1+\epsilon z(g)X)r_0(g)\) is a continuous
representation over \(O\ltimes M\) with the prescribed sign or
similitude. If on a subgroup \(D\) one has
\(z(g)=(\theta(g)-1)m\), the restriction is the reduction modulo \(a\)
in the second coordinate of an integral dual-number conjugate of
\(r_0|_D\).
\end{proposition}
\begin{proof}
Multiplication adds the two square-zero coefficients as
\(z(g)+\theta(g)z(h)=z(gh)\). In rank four the linear terms in the
similitude cancel because the symplectic adjoint of \(X\) is \(-X\).
Choose a lift \(\widetilde m\in O\) of \(m\). Conjugation by
\(1-\epsilon\widetilde mX\) changes the left infinitesimal coefficient
of \(r_0(g)\) to \(\widetilde m(\theta(g)-1)X\).
The matrix entries are continuous in the adic topology since \(r_0\)
is continuous and \(z\) is locally constant.
\end{proof}

For rank three take the normalized representative of each class in
\(H[b]\), where \(H\) is the strict group and \(0\ne b\in O\). At the
marked element \(g_1\) one has \(\Psi(g_1)^2-1\in O^\times\);
subtracting a unique coboundary makes \(z(g_1)=0\). The normalized
representative has values in \((\operatorname{Frac}(O)/O)[b]\simeq O/bO\).
Here the local group at \(\ell\) contains \(s\) with
\(\Psi(s)^2-1\in O^\times\), so its coboundary witness is also killed
by \(b\). Take \(a=b\). At complex conjugation \(c\), the identities
\(\theta(c)=1,\ c^2=1\) give \(2z(c)=0\), hence \(z(c)=0\) since
\(\ell\ne2\). The cocycle identity then gives \(z(cg_1^{-1}c)=0\).
Consequently the matrices at the three elements entering the two
normalization equations agree with those of the base point.
On the prescribed inertia elements away from \(\ell\) where the
residual representation is trivial, the datum requires \(r_0(g)=1\).
There \(\theta(g)=1\), so local strictness gives \(z(g)=0\).
\src{src/Development/AutomorphicArithmetic/AutomorphicVarieties/UnitaryThreeGroups/ShaInjGen.lean}{145}{normalized rank-three cocycles and their root entries}.

In rank three the local condition at \(\ell\) specifies particular
components of a closed locus. Its first-order verification can be made
directly. Let \(U\) represent unrestricted framed local lifts with the
prescribed sign. Put \(H_L=\ker\eta\cap\ker\tau\), where
\(\Psi|_{I_\ell}=\varepsilon_\ell^e\tau\), \(e\ge1\), with \(\tau\)
of finite order. Let \(I_{\rm ord}\subset U\) be the ideal of the
inertia characteristic-polynomial identities and the ordered products
of three shifted matrices on \(H_L\cap I_\ell\), with successive
characters \(\Psi,1,\Psi^{-1}\). Write \(\Delta=\ell d_1d_2\) for the
product of the two consecutive Frobenius-gap equations used to saturate
this ideal. Let \(I_t\) be the intersection of the minimal primes of
\(\sqrt{I_{\rm ord}:\Delta^\infty}\) contained in the kernel of the
diagonal point. The framed universal property makes this point unique.
\src{src/Development/AutomorphicArithmetic/AutomorphicVarieties/PadicHodgeTheory/IsInertiaInfiniteCyclotomic.lean}{133}{the selected ordinary components}.

\begin{lemma}\label{ct:local}
For \(Y\in M_3(O)\), the local framed representation
\((1+\epsilon Y)r_0(1+\epsilon Y)^{-1}\) over \(O[\epsilon]\),
\(\epsilon^2=0\), kills \(I_t\).
\end{lemma}
\begin{proof}
Let \(d:G_{\mathbb Q_\ell}\to\widehat{\mathbb Z}\to\mathbb Z_\ell\)
be unramified degree, normalized on arithmetic Frobenius. The character
\(u_T(g)=(1+T)^{d(g)}\) is continuous for the \((\pi,T)\)-adic
topology by the binomial expansion. Over the domain \(O[[T,V]]\)
use the diagonal point for \(\Psi u_T\), with the same sign as before,
conjugated by \(1+VY\). The antidiagonal inverse-transpose involution
fixes every matrix \(\operatorname{diag}(b,1,b^{-1})\), so this
does define a representation of the full local group. On
\(H_L\cap I_\ell\) the diagonal characters are unchanged. The ordered
product of the three shifted diagonal matrices is zero since each
diagonal entry has its own vanishing factor. Conjugation preserves this
identity and the three characteristic coefficients.

At the distinguished inertia element \(s_0\in H_L\) the three diagonal
values \(\lambda_i\) are distinct: \(\varepsilon_\ell(s_0)\) has
infinite order. Set \(c_i=\prod_{j\ne i}(\lambda_i-\lambda_j)\).
The product of the other two shifted matrices is \(c_iE_{ii}\) in
the diagonal basis. At the specified \(\phi\in H_L\) of nonzero
unramified degree, put \(b=\Psi(\phi)(1+T)^{d(\phi)}\). Its three
diagonal values \(b_i\) are \(b,1,b^{-1}\). The trace defining the
\(i\)-th projected Frobenius value is \(c_i b_i\), and hence
\[
 d_i=c_i c_{i+1}\bigl(b_i-\varepsilon_\ell(\phi)b_{i+1}\bigr)
 \qquad(i=1,2).
\]
These are invariant under the displayed conjugation since all the
elements in the traces have positive sign. Both are nonzero: \(b\)
has nonzero linear coefficient. Thus the map \(U\to O[[T,V]]\)
kills \(I_{\rm ord}\) and its saturation and radical. Its prime kernel
contains a minimal prime \(P\) of that radical. Evaluation at
\(T=V=0\) is the diagonal point, so \(P\) occurs in the intersection
defining \(I_t\). In particular the map kills \(I_t\). Specializing
\(T=0,\ V=\epsilon\) proves the assertion, including when a gap
vanishes at the original diagonal point.
\end{proof}

In rank four use the common denominator for the local witnesses as
in Lemma~\ref{newcoh:cotangent} and the paragraph following it.
At \(\ell\) the indicated ordinary equations prescribe the
characteristic polynomial at \(s_0\), the vanishing of the forbidden
entries between its polynomial projectors, and the prescribed diagonal
inertia characters between equal projectors. These hold on the
diagonal representation and survive symplectic conjugation, since a
polynomial in a matrix transforms by the same conjugation. At
\(p\in S_b=\{p\in S:p\ne\ell,\ \ell\mid\#\psi(I_p)\}\) the condition
is fixed similitude and triviality on \(I_p\cap\ker\psi\); the diagonal
datum has this property, and conjugation preserves it. At all of these
places Proposition~\ref{ct:point} supplies an integral conjugate using
the local witness. Away from \(T=\{\ell\}\cup S_b\), the datum has
\(r_0(g)=1\) on inertia elements with trivial residual action, and the
strict cocycle vanishes there. Thus the local lifts and global \(r_z\)
agree after reduction, with every comparison frame equal to the identity.
The representing framed global ring consequently gives a map into
\(O\ltimes M\); its restriction to the local coefficient ring lifts
to \(O[\epsilon]/(\epsilon^2)\). The same assertions hold for the
normalized rank-three ring by Lemma~\ref{ct:local} and the checks above.
\src{src/Development/AutomorphicArithmetic/AutomorphicVarieties/SymplecticGSp4/ExistsLocalFirstOrderLifts.lean}{188}{integral local conjugates}
and \src{src/Development/AutomorphicArithmetic/AutomorphicVarieties/SymplecticGSp4/ExistsQuotientDerivationOfShaCocycle.lean}{252}{the global framed point}.

\begin{lemma}[Integral lift on a relative presentation]\label{ct:lift}
Let \(L\) be an augmented complete local \(O\)-algebra with weighted
coordinates at the augmentation, and \(P=L[[t_1,\ldots,t_s]]\) with
homogeneous variables of augmentation zero. Suppose a map
\(P\to O\ltimes M\) over the augmentation restricts on \(L\) to the
reduction of an integral derivation of weight \(j\) in its second
coordinate. Suppose its second coordinate on \(t_i\) is zero when
\(\mathrm{wt}(t_i)\ne j\). Then its second coordinate lifts to an
integral derivation of weight \(j\) on \(P\).
\end{lemma}
\begin{proof}
Lift its value on each \(t_i\) to \(d_i\in O\), taking \(d_i=0\) in
all other weights. If \(D_L\) is the given local lift and \(x_L\) is
the augmentation of \(L\), the required map is
\[
 D\Bigl(\sum_\nu b_\nu t^\nu\Bigr)
     =D_L(b_0)+\sum_i x_L(b_{e_i})d_i .
\]
Comparing constant and linear coefficients proves Leibniz and
reduction. The augmentation is fixed by the grading, so this map
vanishes on the other homogeneous pieces.
\end{proof}

Apply the lemma on the weighted presentation before the global
relations. The presentation variables are shifted matrix entries or
comparison-frame entries. Derivatives of the frames vanish, and the
matrix derivatives are precisely the entries of \(z(g)Xr_0(g)\).
The rank-three presentation also uses the inverse-transpose block:
the derivative of inversion at a diagonal \(D\) is
\(Y\mapsto-D^{-1}YD^{-1}\); transposition and reversal of the indices
preserve the assigned block weight. Thus all these variables satisfy
the hypothesis on the weights. The local matrix-entry derivatives of
the integral conjugates have the same weight. More generally, suppose
maps from the local representing ring into integral dual numbers are
determined by the local matrix entries. For an integral torus unit
\(t\), precompose the local map by the weight action of \(t\), or
multiply its second coordinate by the character \(t^j\). The two maps
agree on all matrix entries and hence everywhere. On a homogeneous
element of weight \(i\) this gives \((t^i-t^j)d(f)=0\). For \(i\ne j\)
one can separate the two weights using \(1+\ell\) in one torus
coordinate. Since \(O\) is a domain this gives \(d(f)=0\). This
argument verifies the weight condition on the whole local derivation.
\src{src/Development/AutomorphicArithmetic/AutomorphicVarieties/UnitaryThreeGroups/ModelOfPres.lean}{302}{the relative rank-three presentation}.

In rank four the presentation is further cut by six homogeneous frame
or matrix entries per distinguished place, each with nonzero first
weight. The integral derivation of weight \((0,-1)\) kills them.
In the resulting quotient \(P'\), the model next divides by the kernel
of its characteristic-zero formal chart. The precise property needed
here is that for \(f\) in this kernel some \(0\ne c\in O\) satisfies
\(cf\in(\ker x_{P'})^2\). Then \(cD(f)=0\), so the integral derivation
descends through the chart kernel. This is the cotangent injectivity
condition of the chart after extension to \(\operatorname{Frac}(O)\).
The other normalization equations involve the frames at \(S_b\) and
off-diagonal entries of \(r(g_0)+\alpha(g_0)\Psi(g_0)r(g_0^{-1})\).
Their derivatives vanish modulo \(a\) because the frames are constant
and \(z(g_0)=z(g_0^{-1})=0\). All original global relations also vanish
modulo \(a\), since \(r_z\) defines a point of the framed ring.
\src{src/Development/AutomorphicArithmetic/AutomorphicVarieties/SymplecticGSp4/ExModel.lean}{496}{passage to the normalized weighted quotient}
and \src{src/Development/AutomorphicArithmetic/AutomorphicVarieties/SymplecticGSp4/ExistsSmoothQuotientByWeightedEquations.lean}{104}{the chart-kernel condition}.

On the sum of homogeneous pieces used in Lemma~\ref{newcoh:cotangent},
restrict the resulting derivation to weight \(-1\) (weight \((0,-1)\)
before taking the first-weight-zero algebra in rank four). It kills
the saturation of the square of the augmentation ideal because it
is integral, and its reduction kills the indicated equations.
The functional on \(Q_{-1}\) recovers \(z(g)\) from a homogeneous
lift of matrix entry \((2,1)\) or \((3,1)\), respectively, by dividing
by \(r_0(g)_{11}\in O^\times\). Distinct normalized cocycles therefore
give distinct functionals. In rank three take \(a=b\). In rank four
take \(a=b\prod_v(\theta(s_v)-1)\), as above, so that each local
witness is represented at this denominator. In either rank the same
\(Q_{-1}\) bounds the length of \(H[b]\) for every nonzero \(b\).
Every continuous cocycle into the discrete torsion module has finite
image by compactness, so \(H=\bigcup_hH[\pi^h]\). If \(Q_{-1}\)
has finite length, this increasing sequence stabilizes since a strict
inclusion would increase length. This proves finiteness and the
length bound for \(H\) in both ranks. The following sections construct the weighted presentations, count their
equations and prove cotangent injectivity at the rank-four cut.
Finiteness of the invariant quotient and rigidity at the remaining
global points are further requirements of Theorem~\ref{newcoh:orbitdegree}.

\subsection{The framed presentation and its real-place defect}
\label{dp:section}

Let \(O\) be the ring of integers of a finite extension \(E/\mathbf Q_\ell\),
with residue field \(k\) of odd characteristic. The two presentations used
here are presentations of \emph{framed} representations. In rank three
the form-preserving group in the two dual blocks has Lie algebra
\(\mathfrak {gl}_3\). In rank four, fixing the multiplier gives
\(\mathfrak {sp}_4\). Write \(M_j\) for the residual Lie algebra of
weight \(j\), with Galois action by the adjoint of the residual point.
The grading records the weight of an entry \(ab\) as \(w_a-w_b\).
In particular the tangent vector that reads that entry is indexed here
by the weight of the entry it reads.

Write \(G=G_{\mathbf Q,S}\) for the Galois group unramified outside the
finite primes \(S\) and infinity. Partition \(S=T\sqcup Q\), with
\(\ell=v_0\in T\), and use the full local decomposition groups \(G_v\)
with their maps to \(G\). At \(v\in T\) a complete local augmented
algebra \(B\) supplies a simultaneous family \(\rho_v\) of local
representations. In rank three \(T=\{\ell\}\) and \(B\) is the
prescribed quotient of the local ring. In rank four
\(T=\{\ell\}\cup S_b\). At \(q\in Q\) representations must kill
\(N_q=I_q\cap\ker\bar r\), where \(I_q\) is inertia. In both
applications \([I_q:N_q]\) is prime to \(\ell\). Let \(r_0\) be the
chosen integral lift. An object over a local \(O\)-algebra \(C\)
consists of \(i:B\to C\), a lift \(r\) with the prescribed form and
multiplier, and congruence-to-identity isometries \(\alpha_v\), with
\[
 \alpha_{v_0}=1,\qquad
 r(s)\alpha_v=\alpha_v\,i(\rho_v(s))\quad(s\in G_v).
\]
All matrix-valued maps are continuous. Representations over finite
Artin \(O\)-algebras use the discrete topology.

\begin{proposition}[Relative generators and relations]\label{dp:relative}
Suppose that \(B\) has a finite weighted power series presentation, the
entries of \(\rho_v\) have the indicated weights, and the base point
has weight zero. The framed ring \(U\) has a presentation over \(B\)
by \(g_j\) variables and at most \(o_j\) relations of weight \(j\).
The following formulas describe both integers. Put \(M=M_j\), write
\(dy(s)=s y-y\), and set
\[
 {\cal T}(M)=\{(z,(y_v)):\ z\in Z^1(G,M),\
 z|_{N_q}=0,\ z|_{G_v}=dy_v,\ y_{v_0}=0\}.
\]
Then \(g_j=\dim_k{\cal T}(M)\). An obstruction cocycle is a continuous
triple \((c,(b_v),(b_q))\), where \(c\in Z^2(G,M)\) and
\(db_w=c|_{G_w}\) for \(w\in S\). Divide by the triples
\[
 \bigl(dh,(h|_{G_v}-dy_v)_v,(h|_{G_q}-e_q)_q\bigr),
 \qquad e_q\in Z^1(G_q,M),\quad e_q|_{N_q}=0.
\]
The dimension of the resulting space \({\cal O}b(M)\) is \(o_j\).
The variables can be chosen among entries of \(r(g)-r_0(g)\) and
\(\alpha_v-1\). Zero equations can be appended in each weight to
give exactly \(o_j\) listed equations.
\end{proposition}

\begin{proof}
For clarity, the usual finite presentation argument starts with finitely
many matrices. The maximal pro-\(\ell\) quotient of \(\ker\bar r\) in
\(G\) is topologically finitely generated. Indeed its Frattini quotient
is dual to \(H^1(\ker\bar r,\mathbf F_\ell)\), which is finite by
finiteness of number-field cohomology with ramification in a finite
set. The kernel of this maximal quotient is characteristic in
\(\ker\bar r\). Every Artin lift factors through the corresponding
quotient of \(G\), since the congruence kernel of a matrix group over
a finite local \(O\)-algebra is an \(\ell\)-group. Choose topological
generators there, adjoining lifts of the finite residual quotient.
Start with \(B\) and matrix-entry variables shifted by their base
values for the generators and frames. Invertibility holds modulo the
maximal ideal and hence in the power series ring. The free profinite
group on the generators maps continuously to its matrix group by
passage through the finite coefficient quotients. Impose its relators,
the form, block, multiplier, inertia and frame equations. These are
closed ideal conditions and give precisely the stated universal
property, first over Artin rings and then over their inverse limits.
The noetherian ring already contains the limits of its ideals.
Every matrix-product summand has the weight of the resulting entry.
The base lift preserves the weight spaces and normalizes the form and
block conditions. Thus these equations respect the matrix weights,
including the relations obtained by profinite limits. This weight
argument works termwise modulo powers of the maximal ideal.

Hold \(B\) fixed at its residual point and evaluate on \(k[\epsilon]\),
\(\epsilon^2=0\). Write \(r=(1+\epsilon z)\bar r\) and
\(\alpha_v=1-\epsilon y_v\). Multiplication is exactly the cocycle
identity and the frame equations are the equations defining
\({\cal T}(M)\). In particular no quotient by global conjugation
occurs. Matrix-entry evaluations span the dual of this finite
dimensional space because they separate all the pairs. Choose a basis
among these evaluations in each weight. With the corresponding
variables \(t_i\), the map \(P=B[[t_1,\ldots,t_g]]\to U\) is
surjective by successively lifting modulo powers of the maximal ideal.
Its kernel \(I\) lies in \(\mathfrak m_B P+\mathfrak m_P^2\), by the
chosen basis. The kernel is closed under taking weight components:
substitution of the homogeneous images is componentwise modulo every
power, and both rings are separated. Therefore \(I/\mathfrak m_PI\)
has a homogeneous basis lifting to homogeneous generators. Finiteness
of this vector space permits only finitely many weights. Nakayama's
lemma gives generation of the full ideal.

For the bound, put \(C=I/\mathfrak m_PI\). In
\(P/\mathfrak m_PI\) the maximal ideal annihilates \(C\). Choose \(h\)
so large that \(C\) meets the \(h\)-th maximal-ideal power trivially.
This follows from Krull intersection and finite dimension of \(C\).
After truncating by that power, the extension onto \(P/(I+\mathfrak
m_P^h)\) therefore still has kernel \(C\). Lift the universal matrices
as functions across this finite extension, preserving their form,
and lift the frames. A form-preserving congruence matrix admits such
lifts: lift its entries and correct the form error to first order
by subtracting half its adjoint-symmetrization error in the kernel.
In the paired blocks one may equivalently lift the first block and
use its inverse transpose in the second. At \(v\in T\) the specified
lift is supplied by the structural map of \(B\). At \(q\in Q\)
there is a local lift killing \(N_q\): \(G_q/N_q\) is an extension
of \(\widehat{\mathbf Z}\) by the finite prime-to-\(\ell\) group
\(I_q/N_q\). Averaging on the finite subgroup and
\(\ell\)-cohomological dimension one of \(\widehat{\mathbf Z}\)
give \(H^2(G_q/N_q,M)=0\), which kills the local multiplication
error. Locally constant lifts of matrices exist since all the rings
in this step are finite.

For a set-theoretic lift \(\widetilde r\), the excess of
\(\widetilde r(s)\widetilde r(t)\widetilde r(st)^{-1}\) over 1
is a Lie-algebra-valued two-cocycle with coefficients in \(C\).
Comparison with the true local lifts, conjugated by the lifted
frames, gives the cochains \(b_w\). Changing the global lifts
adds \(dh\), changing frames adds a local coboundary, and changing
a local lift killing \(N_q\) adds a cocycle vanishing there.
Conversely these changes exhaust the displayed coboundaries. The
frame at \(v_0\) may be restored to 1 by common conjugation. Thus
the class after a linear pushout \(C\to k\) vanishes exactly when
the point lifts over \(B\) to that pushout. The component supported
on \(C_j\) has its matrix values in \(M_j\): choose homogeneous lifts
of the entries through the graded truncation and take the coefficient
of a degree-\(j\) kernel element in each error. For the auxiliary
local corrections one may project cochains onto the matching weight
since the residual action preserves every weight. Equivalently first
push out to a homogeneous kernel line, where only entries of its own
weight admit a nonzero homogeneous correction.

The resulting linear map \(C_j^\vee\to{\cal O}b(M_j)\) is injective.
Indeed a zero obstruction gives \(U\to A\) lifting the truncated
point in the pushed-out extension \(A\). Its composite with \(P\to U\)
and the canonical map \(P\to A\) differ by a \(B\)-derivation into
the kernel \(k\). This derivation kills
\(\mathfrak m_BP+\mathfrak m_P^2\), hence \(I\). On \(I\) the first
map is zero and the second map is exactly the chosen functional.
This proves the bound, including the case of a zero obstruction
space. The padding by zero relations makes no minimality assertion.
\src{src/Development/AutomorphicArithmetic/AutomorphicVarieties/CommutativeAlgebra/LemmaP.lean}{256}{The relative presentation argument}
and \src{src/Development/AutomorphicArithmetic/AutomorphicVarieties/CommutativeAlgebra/FramedTwoCoboundaries.lean}{249}{its precise functor and indexing}.
\end{proof}

\begin{proposition}[The equation count]\label{dp:count}
For a finite \(k[G]\)-module \(M\) on which every \(N_q\) acts
trivially, the above spaces satisfy
\[
 \dim{\cal O}b(M)-\dim{\cal T}(M)=
 \dim M^{c=1}+\dim H^0(G,M^\vee(1)),
\]
where \(c\) is complex conjugation and \(M^\vee=\operatorname{Hom}_k(M,k)\).
\end{proposition}

\begin{proof}
Let \(L_q\subset H^1(G_q,M)\) be the inflated classes from
\(G_q/N_q\) and let \(A^1\) be the product of \(H^1(G_v,M)\)
for \(v\in T\) and \(H^1(G_q,M)/L_q\) for \(q\in Q\).
Write \(r_1:H^1(G,M)\to A^1\) and
\(r_2:H^2(G,M)\to\prod_{w\in S}H^2(G_w,M)\).
Sending a pair in \({\cal T}\) to its global class has image
\(\ker r_1\). For any class there choose the local primitives and
subtract the global coboundary of the primitive at \(v_0\).
Over the zero class the choices are local invariant vectors at
\(T\) modulo the simultaneous global invariant vector, as is also
seen by starting with an arbitrary global primitive. Forgetting the
local cochains in \({\cal O}b\) has image \(\ker r_2\) by definition,
and kernel \(\operatorname{coker}r_1\), by subtracting a global
primitive of \(c\). Thus
\[
 \dim{\cal T}=\dim\ker r_1+\sum_{v\in T}h^0_v-h^0_G,\qquad
 \dim{\cal O}b=\dim\operatorname{coker}r_1+\dim\ker r_2.
\]
For \(q\in Q\), averaging over \(I_q/N_q\) expresses \(L_q\)
as the coinvariants of Frobenius on \(M^{I_q/N_q}\).
The kernel and cokernel of Frobenius minus 1 have the same dimension,
so \(\dim L_q=h_q^0\). The final segment of Poitou--Tate identifies
\(\dim\operatorname{coker}r_2=\dim H^0(G,M^\vee(1))\).
The real place has zero positive-degree cohomology since \(\ell\)
is odd. With \(\chi=h^0-h^1+h^2\), the remaining terms are
\(\chi_G-\sum_{w\in S}\chi_w\). The local Euler characteristics
are \(-\dim M\) at \(\ell\) and zero at all other finite primes.
The global Euler characteristic is \(-\dim M^{c=-1}\).
These formulas give the assertion. For the finite module formulas
see Milne~\cite{milneADT},
I, Theorems 2.8, 4.10 and 5.1. Passing from group orders to
\(k\)-dimensions divides all logarithms by \(\log |k|\).
\end{proof}

In rank three the auxiliary datum supplies \(g\) with
\(\eta(g)=\bar\Psi(g)=1\) and \(\bar\chi_\ell(g)=-1\).
Indeed its field-disjointness condition gives \(h\) in the joint
kernel of the residual cyclotomic and finite characters, outside
\(\ker\eta\). Hence \(\bar\Psi(h)=1\) and \(\eta(h)=-1\).
Since \(\eta(c)=-1\), \(\Psi(c)=1\), take \(g=ch\).
It acts as minus the identity on the twisted dual. On
\(\mathfrak {gl}_3(k)\) the action of \(c\) is negative transpose
across the antidiagonal. Its fixed part in weight zero is the line
\(\operatorname{diag}(a,0,-a)\). In either first off-diagonal the
two entries are negative partners, giving one line, and each outer
corner is its own negative, giving zero. The relative list therefore
has \(g_j+{\bf1}_{|j|\le1}\) equations. The eight normalization
entries have weights \(0,0,1,-1,2,-2,1,-1\). The resulting equation
counts in weights \((-1,-2,1,2,0)\) are
\[
 (g_{-1}+3,\ g_{-2}+1,\ g_1+3,\ g_2+1,\ g_0+3).
\]
\src{src/Development/AutomorphicArithmetic/AutomorphicVarieties/UnitaryThreeGroups/BoxPresOf.lean}{149}{The rank-three framed ring}
and \src{src/Development/AutomorphicArithmetic/AutomorphicVarieties/UnitaryThreeGroups/MopBase.lean}{105}{the normalization entries}.

In rank four put \(u=(1,0)\), \(v=(-1,1)\), \(p=(0,1)\),
\(t=(1,1)\). The positive roots are \(u,v,p,t\). The zero
weight space has dimension two, and each of the eight root spaces has
dimension one. On weight \(j=(a,b)\), \(G\) acts by
\(\bar\alpha^a\bar\Psi^b\) and \(c\) by \((-1)^a\).
For \(j=0,\pm p\) this character differs from \(\bar\chi_\ell\)
at \(c\). For the six roots \(\pm u,\pm v,\pm t\), the specified
\(s_\ell\) has cyclotomic residue 1 and root-character residue
different from 1. Thus the twisted dual has no invariants and
\[
 o_j=g_j+2{\bf1}_{j=0}+{\bf1}_{j=p\ \mathrm{or}\ j=-p}.
\]
There is a separate arithmetic premise for \(g_j=0\) at the six
roots. It is the strict residual class vanishing for the characters
\(\bar\alpha^{\pm1},\bar\beta^{\pm1},\bar\theta^{\pm1}\) imposed by
the good datum. It kills \(\ker r_1\); the local invariants at every
\(v\in T\) vanish because the root character is nontrivial at \(s_v\).
The first dimension formula of the proof now gives \(g_j=0\).
\src{src/Development/AutomorphicArithmetic/AutomorphicVarieties/SymplecticGSp4/NumsZero.lean}{17}{The residual vanishing characters and their use in the count}.

\subsection{The rank-four local chart at the diagonal point}
\label{dp:chart}

Write \(\chi\) for the \(\ell\)-adic cyclotomic character.
The base characters in order are
\[
 (\theta,\Psi,\alpha,1),\qquad
 \Psi=\chi^e\psi,\quad\alpha=\chi\omega^{-1}\xi,\quad
 \beta=\Psi/\alpha,\quad\theta=\alpha\Psi .
\]
The multiplier is \(\theta\). The good datum has \(e\ge2\), finite
characters \(\psi,\xi,\omega\), and \(\xi\) unramified at \(T\).
At \(\ell\) it gives \(s_0\in I_\ell\cap\ker\psi\cap\ker\omega\)
with \(\chi(s_0)\) of infinite order. It also gives
\(\chi(s_\ell)=1\), \(\bar\omega=\bar\chi\), and
\(\bar\xi(s_\ell)\bar\Psi(s_\ell)^i\ne1\) for every integer \(i\).

\begin{proposition}[Local charts of the framed symplectic ring]\label{dp:local}
At the diagonal point, the completed characteristic-zero local ring
for the ordinary projector equations at \(\ell\) is a weighted power
series ring over \(E\). It has two coordinates of each weight
\(0,u,v,p,t\), and one of each weight \(-u,-v,-p,-t\).
For \(q\in S_b\) the corresponding ring for the condition
\(r(I_q\cap\ker\psi)=1\) has two weight-zero coordinates and one
of each of the eight root weights.
\end{proposition}

\begin{proof}
Use finite local Artin \(E\)-algebras \(A\) with residue \(E\),
given their finite-dimensional \(\ell\)-adic vector space topology.
At \(\ell\) write \(d_i=\psi_i(s_0)\) for the four diagonal values.
They are distinct, since their exponents on \(\chi(s_0)\) are
\(e+1,e,1,0\). The source equations specify
\(\det(X-r(s_0))=\prod_i(X-d_i)\). Hence
\[
 P_i=\prod_{k\ne i}\frac{r(s_0)-d_k}{d_i-d_k}
\]
are orthogonal projectors summing to one. Each image is a direct
summand of rank one by reduction to \(E\). The other equations say
\(P_jr(g)P_i=0\) for \(j>i\), and
\(P_i(r(h)-\psi_i(h))P_i=0\) for \(h\in I_\ell\). They specify
stability of the ordered flag and the prescribed inertial action on
its lines. The multiplier pairs lines \(i\) and \(5-i\) perfectly:
pairing any two eigenlines with \(d_i d_j\ne\theta(s_0)\) gives zero,
and nonsingularity of the entire pairing gives the assertion.
Rescaling eigenvectors lifting the original basis now produces a
symplectic basis reducing to that basis. Conversely an upper
triangular lift in such a basis with the given inertia characters
obeys the projector equations, since the partial sums of the
eigenlines of \(r(s_0)\) are its prescribed flag. The flag is
uniquely recovered by the projectors, including over nonreduced \(A\).

In the open cell of the symplectic flag variety at this flag there
is a unique lower unipotent matrix carrying the standard flag to a
given flag. This is Gaussian elimination with diagonal pivots units:
factor a near-identity change of symplectic basis into a lower
unipotent and an upper triangular factor. The involution
\(g\mapsto J^{-1}g^{-t}J\) preserves the two factors, so uniqueness
in \(\mathrm{GL}_4\) makes the factors symplectic. The lower
unipotent Lie algebra has one line in each negative root. Exponential
and logarithm are polynomial inverse maps here in characteristic
zero, since strictly triangular matrices are nilpotent. Thus the
moving-flag parameters give four smooth negative-weight directions.
It remains to consider a lift in the upper triangular group with
standard flag and fixed multiplier.

The diagonal deformations of \(\alpha\) and 1 are independent
unramified twists, determining the other two by the multiplier.
They lift across every small extension \(A'\to A\): use logarithm
on \(1+\mathfrak m_A\), lift the value at Frobenius additively, and
use \(\widehat{\mathbf Z}\to\mathbf Z_\ell\) to define the lifted
continuous additive character. Here and below a small extension has
kernel \(V\) killed by the maximal ideal of \(A'\). With lifted
diagonal fixed, the obstruction to lifting the upper-unipotent part
is in
\[
 H^2(G_{\mathbf Q_\ell},\mathfrak n\otimes_E V),\qquad
 \mathfrak n=E(\alpha)\oplus E(\beta)\oplus E(\Psi)\oplus E(\theta).
\]
Indeed a continuous vector-space section on the nilpotent entries,
in upper-unipotent exponential coordinates, gives a lifted cochain.
The multiplication error takes values in the additive kernel.
Conjugation on that kernel uses only the residual diagonal, so the
four root spaces split as written. A coboundary error can be
removed by subtracting a continuous one-cochain. Local Tate duality
gives \(H^2(G_{\mathbf Q_\ell},E(\gamma))^\vee
=H^0(G_{\mathbf Q_\ell},E(\chi/\gamma))\).
For \(\gamma=\alpha,\beta,\theta\) the residual value at \(s_\ell\)
differs from 1 by the displayed nontriviality with \(i=0,-1,1\),
respectively, while \(\chi(s_\ell)=1\). For \(\gamma=\Psi\) use
\(\chi(s_0)^{e-1}\ne1\). The four obstructions vanish. Every
surjection of finite local Artin \(E\)-algebras factors into small
extensions, giving formal smoothness.

The framed tangent space for a positive root \(\gamma\) consists of
all \(Z^1(G_{\mathbf Q_\ell},E(\gamma))\). By local Euler
characteristic and the preceding \(H^2=0\), its dimension is
\((1+h^0)+(1-h^0)=2\). The two unramified characters give two
zero-weight parameters. Together with the moving flag this gives
fourteen. The \(\ell\)-adic vector-space duality and Euler formula
used here follow from the finite-module theorems by applying them
to the finite quotients of a stable lattice and inverting \(\ell\).
For this passage the inverse systems of finite cohomology groups
are Mittag--Leffler, and every continuous cochain on the compact
powers of the Galois group takes values in a bounded lattice.

At \(q\in S_b\) the effective local group is \(G_q/(I_q\cap\ker\psi)\).
Its inertia quotient \(\Delta\) is finite since \(\ker\psi\) is
open. Its order may be divisible by \(\ell\). Over \(E\), averaging
over \(\Delta\) is exact. Consequently for each weight space \(W\)
of the symplectic adjoint, \(H^2\) of the effective group is zero
and \(H^1\) is the space of Frobenius coinvariants of \(W^\Delta\).
Indeed a stable lattice in \(W^\Delta\) reduces this last calculation
to cohomological dimension one of \(\widehat{\mathbf Z}\) on its
finite \(\ell\)-power quotients. Thus \(h^1=h^0\).
The small-extension multiplication-error argument with smooth
symplectic matrices gives formal smoothness, and its framed tangent
space has dimension \(h^1+\dim W-h^0=\dim W\) weight by weight.

To pass from this calculation to the completion, the integral local
ring \(R\) represents the actual matrix equations: both the
undivided projector expressions and the inertia-kernel conditions
are closed matrix equations in the unrestricted framed ring.
Write \(x:R\to O\) for the diagonal point and \(J=\ker x\).
Then \(\widehat{(R[1/\ell])_{J[1/\ell]}}\) represents precisely the
above Artin \(E\)-valued lifts. For such a lift, its matrix entries
minus their diagonal base values form a bounded subset of the
nilpotent ideal of \(A\) by compactness. Put them into a finite
\(O\)-lattice \(L\) in that ideal. If \(\mathfrak m_A^d=0\), the
order \(O+L+\cdots+L^{d-1}\) is complete local with residue \(k\)
and contains the entire representation. Its finite \(O\)-module
topology is its maximal-ideal topology, since its nilpotent ideal
is killed by the \(d\)-th power. This gives a map from \(R\)
reducing to \(x\), hence a local map on the indicated completion.
Conversely such a map kills \(J^d\), and \(R/J^d\) is a finite
\(O\)-module by the filtration with quotients \(J^i/J^{i+1}\).
It thus induces a continuous lift with bounded entries. Two
choices of order embed in a common one obtained by enlarging \(L\),
so the universal property also gives uniqueness.

Choose homogeneous lifts of a basis of \( (J/J^2)\otimes_O E\)
in \(J[1/\ell]\), using shifted matrix entries. They give a
surjection from a power series ring onto the completion. Formal
smoothness forces it to be an isomorphism: lift the inverse tangent
map successively to the power series ring modulo its maximal-ideal
powers. The resulting composite on the power series ring has
invertible linear term and is an automorphism by recursive
substitution. The coordinates and inverse are homogeneous in each
finite order, by uniqueness of the inverse. This also supplies
cotangent injectivity for the original integral ring after tensoring
with \(E\). The kernel of the map \(R\) to this completion is the
kernel of \(R\to(R[1/\ell])_{J[1/\ell]}\) by Krull intersection.
It is a finitely generated ideal, so a single \(y\in R\) with
\(x(y)\ne0\) annihilates it, by clearing the finitely many
denominators separately and taking their product.
\src{src/Development/AutomorphicArithmetic/AutomorphicVarieties/PadicHodgeTheory/SpFlag.lean}{309}{The inertia projector equations},
\src{src/Development/AutomorphicArithmetic/AutomorphicVarieties/PadicHodgeTheory/SpLiftL.lean}{32}{the four character checks}
and \src{src/Development/AutomorphicArithmetic/AutomorphicVarieties/SymplecticGSp4/LocBase.lean}{142}{the integral chart requirements}.
\end{proof}

\subsection{Cotangent faithfulness after the symplectic cuts}
\label{dp:cut}

\begin{lemma}\label{dp:cutlemma}
Let \(P\) be an augmented noetherian complete local \(O\)-algebra,
with augmentation ideal \(J\). Suppose a weighted chart
\(\phi:P\to E[[Y_1,\ldots,Y_n]]\) induces an isomorphism from
\((J/J^2)\otimes_O E\) to the cotangent at \(Y=0\).
Let \(u_1,\ldots,u_h\in J\) be homogeneous with independent images in
this rational cotangent. There is a chart of \(P'=P/(u_i)\) with
one fewer variable of the given weight for each cut, again inducing
a cotangent isomorphism. If \(P'\) is further divided by the kernel
of this latter chart, every \(O\)-valued point derivation at its
augmentation descends to the resulting domain.
\end{lemma}

\begin{proof}
Complete the linearly independent \(d\phi(u_i)\) to a basis in each
weight. Replace the first coordinates by the full series
\(\phi(u_i)\), keeping the complementary homogeneous linear
coordinates. The substitution has zero constant terms and invertible
linear part. Its inverse is obtained recursively in total degree.
Each finite degree respects the weights even when negative and zero
weights occur. Set the first \(h\) new coordinates to zero. On
cotangents this kills precisely the span of \(du_i\), and the
cotangent of \(P'\) is \(J/(J^2+(u_i))\). Exact scalar extension
therefore gives the asserted isomorphism. For \(a\) in the new
chart kernel its augmentation vanishes and its differential is zero
after inverting \(\ell\). Some \(0\ne c\in O\) then satisfies
\(ca\in(\ker x_{P'})^2\). Every integral point derivation kills this
square, hence kills \(a\) by torsion-freeness of \(O\). The chart
target is a domain.
\end{proof}

For the application take \(P=B[[t_1,\ldots,t_g]]\) from
Proposition~\ref{dp:relative}. The simultaneous local ring \(B\)
represents the product of the finitely many local functors. Its
rational completion is the power series ring on the union of the
local coordinates by Proposition~\ref{dp:local}. The relative
variables add their own free summands to the cotangent, so \(P\)
has the required chart. The global relations cut out \(U\) inside
this ambient formal spectrum. The next cuts are lifted to \(P\)
from the following entries of \(U\), in one-based indices:
\[
 {\cal P}=\big((1,2),(1,4),(3,2),(2,1),(4,1),(2,3)\big).
\]
At \(\ell\) take these entries of \(r(\phi_0)\), where
\(\phi_0=\lim_n s_\ell^{\ell^{n!}}\) is the element specified by the
datum. At every other \(v\in T\) take these entries of \(\alpha_v\).
All have zero augmentation. Use the six Lie algebra elements,
in the same order,
\[
 E_{12}-E_{34},\quad E_{14},\quad E_{32},\quad
 E_{21}-E_{43},\quad E_{41},\quad E_{23}.
\]
Each lies in \(\mathfrak {sp}_4(O)\) for the antidiagonal pairing with
signs \(1,1,-1,-1\), and has just one nonzero entry among \({\cal P}\).
For a selected element \(X\) conjugate \(r_0\) and all the local lifts
by \(1+\epsilon X\) over \(O[\epsilon]/(\epsilon^2)\), leaving all
frames equal to 1. This preserves the local equations by conjugation.
At \(\ell\) the only nonzero selected derivative is
\(\psi_j(\phi_0)-\psi_i(\phi_0)\). Indeed the residual diagonal at
\(\phi_0\) is the same as at \(s_\ell\), since the values belong to
\(k^\times\) and the residue degree \([k:\mathbf F_\ell]\) divides
\(n!\) eventually. Each pair has different first weight, so
\(\bar\xi(s_\ell)\bar\Psi(s_\ell)^m\ne1\) for all \(m\) makes
the residual entries distinct. At \(v\ne\ell\) keep the global
lift fixed, move only \(\alpha_v\) to \(1+\epsilon X\), and
conjugate the corresponding local diagonal lift by its inverse.
This tests one entry with derivative 1. These point derivations of
\(U\) pull back to \(P\), proving independence for any homogeneous
lifts of all the indicated cuts there.

Lemma~\ref{dp:cutlemma} now gives the post-cut chart and the
integral descent of derivations needed in
Section~\ref{dp:section}'s framed construction. Let \(b=|S_b|\).
In each odd root the cut removes \(b+1\) variables. The remaining
dimensions there are 1 for \(u,v,t\) and 0 for their negatives,
using the residual vanishings above. In weights \((-p,p,0)\)
the dimensions are
\[
 (1+b+g_{-p},\quad 2+b+g_p,\quad 2+2b+g_0).
\]
The remaining frame equations at each bad prime have weights
\(p,-p,0,0\), and the two normalization entries of
\(r(g_0)+\theta(g_0)r(g_0^{-1})\) have weights \(p,-p\).
Adding their counts to Proposition~\ref{dp:count} gives final lists
of lengths \(2+b+g_{-p},\,2+b+g_p,\,2+2b+g_0\), respectively.
They define the quotient of the resulting domain used by the model.
\src{src/Development/AutomorphicArithmetic/AutomorphicVarieties/SymplecticGSp4/ExModel.lean}{155}{The entries and six root vectors}
and \src{src/Development/AutomorphicArithmetic/AutomorphicVarieties/SymplecticGSp4/ExModel.lean}{425}{the normalized model}.

\subsection{The selected rank-three branch at a resonant character}
\label{dp:ordinary3}

Here \(G=G_{\mathbf Q_\ell}\), \(I\) is its inertia, and \(\eta\)
is an unramified character of order dividing two. Put \(H=\ker\eta\)
and \(I_0=I\cap H\cap\ker\tau\), where \(\tau\) has finite image.
The source specifies \(\Psi|_I=\chi^e\tau|_I\) with \(e\ge1\).
The matrix group is \(\mathrm{GL}_3\rtimes\{1,-1\}\), where
\(-1\) sends \(A\) to \(J A^{-t}J\) and \(J\) reverses the basis.
Put \((d_1,d_2,d_3)=(\Psi,1,\Psi^{-1})\). The base point of
the unrestricted local framed ring \(R\) is
\(r_0(g)=(\operatorname{diag}(d_i(g)),\eta(g))\), and its
augmentation is \(x\). For \(s\in I_0\) put \(A_i(s)=r(s)-d_i(s)\)
in the matrix block. The ordinary ideal \(I_{\rm ord}\) fixes the
characteristic polynomial on \(I_0\) to \(\prod_i(X-d_i(s))\)
and kills all entries of \(A_1(s)A_2(s')A_3(s'')\) for independent
\(s,s',s''\in I_0\). The definition then saturates by a specific
element \(\delta\), described below. Write
\[
 I_Y=(I_{\rm ord}:\delta^\infty),\qquad
 I_t=\bigcap_{\substack{\mathfrak p\text{ minimal over }I_Y\\
                          \mathfrak p\subseteq\ker x}}\mathfrak p .
\]
Taking the radical of \(I_Y\) before selecting these primes gives
the same intersection. The framed universal property makes \(x\)
the unique map with this specified matrix point.
\src{src/Development/AutomorphicArithmetic/AutomorphicVarieties/PadicHodgeTheory/IsInertiaInfiniteCyclotomic.lean}{97}{The ordinary equations, saturation and selected primes}.

\begin{proposition}\label{dp:rank3chart}
There is a unique selected prime. The completion of
\((R/I_t)[1/\ell]\) at \(x\) is a power series ring over \(E\)
with coordinate counts \(1,2,3,4,2\) in weights \(-2,-1,0,1,2\).
The kernel of the resulting chart from \(R\) is \(I_t\).
These statements hold also when a ratio of two adjacent diagonal
characters is cyclotomic.
\end{proposition}

\begin{proof}
As in Proposition~\ref{dp:local}, use continuous lifts over Artin
\(E\)-algebras, now of \(r_0\). Choose the specified
\(s_0\in I_0\) with \(\chi(s_0)\) of infinite order.
The eigenvalues \(d_i(s_0)\) are distinct. The characteristic
polynomial gives three rank-one projectors. Write \(V_1\) for the
first eigenline and \(V_2\) for the sum of the first two.
In eigen-coordinates, \(A_2(s_0)A_3(s_0)\) is a unit multiple of
the first projector, and \(A_1(s_0)A_2(s_0)\) is a unit multiple
of the third. The triple products therefore say that \(r(s)\) acts by
\(d_1(s)\) on \(V_1\), preserves \(V_2\), and acts on the last
quotient by \(d_3(s)\). Subtracting these two diagonal terms from
the specified trace gives \(d_2(s)\) on \(V_2/V_1\). Conversely
these three conditions imply every triple product, by applying its
factors from right to left along this flag.

The flag is unique by its ordered characters at \(s_0\). Conjugation
by \(g\in G\) preserves \(I_0\) and all the \(d_i\). When
\(\eta(g)=-1\) the induced action on flags first takes the dual flag
in reversed order via \(J\); its ordered quotient characters on
inertia are again \(d_1,d_2,d_3\) since \(d_{4-i}^{-1}=d_i\).
Uniqueness therefore gives stability under the full semidirect
action. By lower-unipotent change of basis the moving flag contributes
three free coordinates, two of weight \(-1\) and one of weight
\(-2\). Once the flag is standard all matrix blocks are upper
triangular, since inverse transpose across the antidiagonal preserves
the standard flag. Conjugation here and below is in the semidirect
group, so a change of basis also changes the block of an element
outside \(H\) by the prescribed involution.

The diagonal part relative to \(r_0\) is an unramified cocycle in
the nilpotent diagonal torus: on inertia the action on the torus
is trivial, and a character trivial on \(I_0\) factors there
through a finite group with no nonzero additive image in an
\(E\)-algebra. Take logarithms. The three coordinates \(b_i\) of
the value at a chosen Frobenius are free, even if \(\eta\ne1\).
Indeed a cocycle on \(\widehat{\mathbf Z}\) with Frobenius action
\(D=1\) or \(D(b_1,b_2,b_3)=(-b_3,-b_2,-b_1)\) has value
\(\sum_{k=0}^{n-1}D^kb\) at its \(n\)-th power. This extends
continuously using parity and the projection to \(\mathbf Z_\ell\).
Thus the base for the upper unipotent part is
\(\Lambda=E[[b_1,b_2,b_3]]\).

First divide the upper unipotent group by its central line \(E_{13}\).
The two adjacent entries form an additive module. For fixed diagonal
they are precisely continuous one-cocycles, with no quotient by
coboundaries. If \(H=G\), they give two independent rank-one
problems with respective parameters \(t=b_1-b_2,b_2-b_3\).
Their characters are \(\Psi\exp(t\nu)\), where the continuous
unramified additive character \(\nu\) is 1 on Frobenius.
If \([G:H]=2\), the two adjacent lines are exchanged by the
antidiagonal inverse transpose on their Lie algebra. The module
is induced from the first line on \(H\): this line and its translate
by a Frobenius form a basis, also for the deformed diagonal over
every Artin algebra. Its character on \(H\) is
\(\Psi|_H\exp(t\nu_H)\), with \(t=b_1-b_3\) and \(\nu_H\) equal
to 1 on Frobenius of \(H\). Indeed the value of the unramified
diagonal logarithm on the square of Frobenius is \(b+Db\).
This description by induction works functorially in the diagonal
parameters. In particular at \(t=0\) the indicated module is
isomorphic over the whole diagonal base to the constant module.

Here is the local equation for either one of these problems. Let
\(G_*=G\) in the split case and \(G_*=H\) in the induced case.
Denote the corresponding \(G\)-module at the origin by \(W\) and
put \(m=\dim_E Z^1(G,W)\). For its cocycle functor over \(\Lambda\),
a minimal presentation uses \(m\) variables and at most
\(\dim_E H^2(G,W)\) relations. To justify this version over \(E\),
the functor preserves fibre products of Artin algebras since its
objects are functions subject to equalities, and its tangent is
finite-dimensional by local Euler characteristic. Schlessinger's
prorepresentability criterion~\cite[Theorem 2.11]{schlessingerArtin}
therefore applies (the functor at the residue field is
one point). One can include the three diagonal
parameters in the functor and then work relative to their power
series ring. The small-extension argument in the proof of
Proposition~\ref{dp:relative} applies verbatim using continuous
linear sections of surjections of finite-dimensional \(E\)-spaces;
here the single obstruction is the coboundary of a lifted one-cochain.
The relation bound concerns the actual ideal modulo its maximal-ideal
multiple. Local Tate duality, with Shapiro's lemma in the induced
case, gives
\[
 \dim_E H^2(G,W)=
 {\bf1}_{\,\Psi|_{G_*}=\chi|_{G_*}} .
\]
If this is zero the cocycle algebra is a power series ring.

If the dimension is one, its possible single equation is, up to a
change of variables, exactly \(tY=0\). We verify its order of
vanishing. Choose the relative variables \(X_1,\ldots,X_m\) among
linear evaluations of the cocycle; they have weight one for scalar
multiplication of the cocycle. Modulo \(t\) the module is constant
up to the functorial basis isomorphism just proved, so every cocycle
is an arbitrary linear combination of a basis of \(Z^1(G,W)\) with
coefficients in the maximal ideal of the test algebra. Its completed
algebra over \(\Lambda/(t)\) is free in the \(m\) variables. The
chosen evaluations still give coordinates there by the inverse
function theorem. It follows that every equation is divisible by
\(t\). No relation has an \(X\)-independent term, since zero is
a cocycle for every diagonal.

Vary \(t\) to \(h\) and the cocycle to \(kz\) over
\(E[h,k]/(h^2,k^2,hk)\), choosing any linear diagonal section of
the parameter \(t\). The obstruction to lifting to
\(E[h,k]/(h^2,k^2)\), in the \(hk\) coefficient, maps to
\(\nu_{G_*}\smile [z_*]\) under the preceding cohomology
identification. In the split case this follows by expanding the
cocycle identity with character \(\chi(1+h\nu)\): its error on
\((g,g')\) is \(hk\,\nu(g)\chi(g)z(g')\) up to the sign of the
coboundary convention. In the induced case the cochain map to
\(G_*\) restricts and projects to the inducing line; it commutes
with reduction for the induced basis and gives the identical
calculation. Every class of \(H^1(G_*,E(\chi))\) occurs by
Shapiro. The local Tate pairing of this space with \(H^1(G_*,E)\)
is perfect and \(\nu_{G_*}\ne0\). Hence for some \(z\) the
obstruction is nonzero.

In particular there is one relation, with nonzero mixed term
\(t\sum_i c_i X_i\), some \(c_i\ne0\). It may be chosen
homogeneous of weight one. Indeed the ideal is stable under scalar
multiplication of all \(X_i\), which fixes \(b\). Its weight
components belong to the ideal by polynomial interpolation in each
finite truncation over the infinite field \(E\). All relations lie
in \(t(X_1,\ldots,X_m)\), so a relation with this nonzero
quadratic term generates the one-dimensional space of relations
modulo the maximal-ideal multiple. Its weight-one component still
generates. Write it \(t\sum_i a_i(b)X_i\); some \(a_i\) is a
unit. Use \(Y=\sum_i a_i(b)X_i\) as one of the relative
coordinates. This proves the claimed equation including its
multiplicity.

Finally restore the central root. Its Galois action on \(I_0\)
is \(\Psi^2=\chi^{2e}\), hence its \(H^2\) vanishes by local
duality since \(2e\ne1\). Multiplication errors for a lifted
central entry in a small extension are two-cocycles for this root
module, as above. The map to the first-layer functor is therefore
formally smooth with relative tangent its space of one-cocycles.
This adds free variables also over the singular first-layer ring:
choose relative cotangent coordinates and split the corresponding
power series surjection order by order by formal smoothness.
Equivalently the small-extension relation bound with zero obstruction
gives zero relative relations. The central coordinate has weight two.

We compute the effect of the prescribed saturation in these
coordinates. The source chooses \(f\in H\cap\ker\tau\) mapping to
a positive power of Frobenius. For \(i=1,2,3\) put
\(c_i=\prod_{k\ne i}(d_i(s_0)-d_k(s_0))\in E^\times\).
Its trace expression is \(c_i\,\mathrm{tr}(P_i r(f))\), where
\(P_i\) is the corresponding eigenprojector of \(r(s_0)\). In an
upper triangular basis this equals \(c_i\kappa_i(f)\), with
\(\kappa_i\) the \(i\)-th diagonal character on \(H\). Indeed each
\(P_i\) is upper triangular with only its \(i\)-th diagonal entry
equal to 1. Thus \(\delta\) is, up to a unit in the characteristic-zero
completion, precisely
\[
  \prod_{i=1}^2\bigl(\kappa_i(f)/\kappa_{i+1}(f)-\chi(f)\bigr).
\]
For a resonant adjacent block the corresponding factor is
\(\chi(f)(\exp(a t)-1)\), with \(a=\nu_{G_*}(f)>0\).
It is a unit times \(t\). This formula also covers the two equal
ratios on \(H\) in the induced case. In the quotient by all the
coordinates \(Y\) belonging to resonant blocks, the ring is a
power series domain and the image of \(\delta\) is nonzero:
each adjacent ratio still varies as
\(\Psi(f)\exp(a t)\) for a nonzero diagonal parameter \(t\).
Since \(tY=0\) and \(t\) divides \(\delta\) in each resonant
block, the \(\delta\)-power torsion is exactly the ideal of these
\(Y\)'s. This establishes the calculation even if \(\delta(x)=0\).

For any \(G\)-module \(W\) in an adjacent block the local Euler
formula over \(\mathbf Q_\ell\) gives
\(\dim Z^1(G,W)=2\dim W+\dim H^2(G,W)\).
Saturation removes exactly one variable for each of the latter
dimensions. The adjacent roots thus leave four weight-one
coordinates in all; the central root adds two of weight two.
The three unramified diagonal coordinates and the three lower flag
coordinates give the other asserted counts.

Localization at the characteristic-zero point and completion are
flat and commute with \((I_{\rm ord}:\delta^n)\): take the kernel
of multiplication by \(\delta^n\) on the quotient. The chain
stabilizes by noetherianity. It follows that the completion of the
localization modulo \(I_Y\) is the power series domain just computed.
Faithful flatness makes the localized quotient a domain.
Consequently there is exactly one minimal prime over \(I_Y\)
contained in \(\ker x\), the contraction of the zero ideal of this
local quotient. Its contraction to \(R\) is also the kernel of the
completed chart. This proves the assertion about the selected
integral ideal \(I_t\).
\src{src/Development/AutomorphicArithmetic/AutomorphicVarieties/PadicHodgeTheory/RtChart.lean}{41}{The corresponding twelve-coordinate conclusion}.
\end{proof}

The preceding computation is a statement about the ordinary matrix
ideal in the completed local ring. Application to the integral
presentation uses a weight grading on \(R/I_t\). The unrestricted
matrix ring has homogeneous coordinates given by shifted entries.
The ordinary equations are weight-homogeneous, \(\delta\) has
weight zero, and the map to the completed chart preserves weights.
Its kernel \(I_t\) is thus closed under projection onto each weight.
Conjugation by a matrix congruent to 1 also preserves the selected
branch over arbitrary integral test algebras. Indeed conjugate the
universal \(R/I_t\)-valued family by the generic matrix \(1+X\)
in \((R/I_t)[[X_{ij}]]\). Embed this ring coefficientwise into
\(E[[Y_1,\ldots,Y_{12},X_{ij}]]\) by the chart. The ordinary
equations are conjugation invariant; \(\delta\) is unchanged because
its defining traces use elements of \(H\), and it is nonzero in this
domain. Hence the conjugate annihilates \(I_Y\). For \(a\in I_t\)
some \(s\) with \(x(s)\ne0\) satisfies \(sa\in I_Y\), by the
calculation of the localized prime. The conjugate image of \(s\) has
constant coefficient \(x(s)\), so \(a\) also maps to zero.
Injectivity then gives the same conclusion before passage to \(E\).
Substituting any matrix of elements of the maximal ideal for \(X\)
proves the integral assertion. This supplies the local presentation required by
Proposition~\ref{dp:relative} in rank three. Identifying potentially
crystalline points of the required type on the entire model is
a further arithmetic statement.

In rank four the local and relative arguments use the good-datum
hypotheses, including the separate residual vanishings.
In both ranks, finiteness of the integral invariant quotient,
exclusion of extraneous one-sided points and rigidity of the
nontrivial two-sided points still require arithmetic control of
representations away from the diagonal base point. The local charts
and equation counts alone yield none of these global-point assertions.

\subsection{The special fibre of the invariant algebra}\label{dp:finitecurves}

There is a separate arithmetic test for the finiteness hypothesis in
Theorem~\ref{newcoh:orbitdegree}. In rank three it leaves only
strongly irreducible three-dimensional specializations on \(G_K\).
In rank four it leaves an absolutely irreducible four-dimensional
specialization or an even absolutely irreducible two-dimensional
constituent. The fields of coefficients in these statements have
characteristic \(\ell\).

\begin{lemma}[Detecting a nonfinite invariant algebra]\label{dp:detect}
Let \(O\) be a complete DVR with finite residue field \(k\) and
uniformizer \(\pi\). Let \(S\) be a quotient of a finite formal
power series algebra over \(O\) by a closed ideal stable under
projection onto the weight components for integral vector weights.
Write \(C\subset S\) for the image of the weight-zero power series.
Assume the variables are in the maximal ideal. Then \(C\) is finite
over \(O\) if and only if every continuous local map
\[
               S\longrightarrow k'[[t]],\qquad \pi\longmapsto 0,
\]
for finite extensions \(k'/k\), takes \(C\) into the constants.
\end{lemma}
\begin{proof}
The nonzero minimal exponent vectors of weight zero are finite
by Dickson's lemma, and they generate the monoid of all balanced
exponents by subtraction. Denote the corresponding monomials by
\(m_j\). Substituting these into a formal power series over \(O\)
surjects onto \(C\): fix one factorization for each exponent and
group the terms accordingly. The bounded total degrees of the
\(m_j\) ensure convergence in both descriptions. Thus \(C\) is
complete noetherian local, with maximal ideal \((\pi,m_j)\).
Weight-zero projection gives \(C\cap\pi S=\pi C\).
If \(C/\pi C\) is artinian, a power of its maximal ideal is zero.
Consequently the \(\pi\)-adic and maximal-ideal topologies on \(C\)
are equivalent; lifting a \(k\)-basis of \(C/\pi C\) and expanding
successively in powers of \(\pi\) gives finitely many \(O\)-module
generators. Conversely a finite algebra has artinian special fibre,
whose image in the domain \(k'[[t]]\) is a finite field.

Otherwise some \(m_j\) is not nilpotent modulo \(\pi C\). By the
projection identity it survives on a prime quotient of \(S/\pi S\).
Cut this complete local domain to dimension one while preserving
the image \(b\ne0\) of \(m_j\). At each step of dimension \(>1\)
choose an element of the maximal ideal outside the finitely many
height-one primes dividing \(b\), and take a minimal prime over
that nonzero element. This prime has height one, avoids \(b\), and
is strictly below the maximal ideal. The resulting one-dimensional
complete domain has finite normalization since it is excellent.
This normalization is local by completeness and the domain property,
and is an equicharacteristic complete DVR with finite residue
extension. Its coefficient field is the unique lift of this finite
field, so a uniformizer identifies it with \(k'[[t]]\) over \(k\).
Here \(b\) has strictly positive valuation. This proves the criterion.
\end{proof}

In the two presentations one may use homogeneous coordinates taken
from matrix entries minus their residual diagonal lifts. Indeed these
entries generate the representing ring topologically, by the framed
universal property, and homogeneous cotangent generators suffice by
successive approximation. Equivalently, on the post-cut rank-four
model the entries generate \(S\) in the sense that they distinguish
its maps to complete local algebras. They then span its relative
cotangent space by testing the dual numbers of the residue field.
For these coordinates \(C\) is the closure of the algebra generated
by weight-zero products of entries. Closure does not change module
finiteness: a finite \(O\)-submodule is compact, hence closed, and
an \(O\)-submodule of a finite \(O\)-module is finite. Thus
Lemma~\ref{dp:detect} applies to the actual weight-zero entry algebra.
The image of the invariant part \(B\) of the presentation in the
post-cut ring is \(B/(I\cap B)\). In rank four one first takes
weight zero for the first coordinate and subsequently for the second;
this is simultaneous weight \((0,0)\).
\src{src/Development/AutomorphicArithmetic/AutomorphicVarieties/CurveDVRSpecialization/FinCrit.lean}{122}{Curve detection in rank three};
\src{src/Development/AutomorphicArithmetic/AutomorphicVarieties/SymplecticGSp4/FgMapMk.lean}{24}{From the rank-four weight-zero entry algebra to the invariant quotient}.

Here are two arithmetic facts about such curves. We may enlarge
\(k'((t))\) by a finite extension and normalize its valuation ring
to define constituent spaces, their flags and stable lattices.
Take intersections and saturated quotients of the original lattice
to get stable lattices on subquotients. Their residual constituents
together give the residual semisimplification of the original
representation, by reduction of these exact lattice sequences.

First, a continuous character of a number-field Galois group into
the units of an equal-characteristic DVR with finite residue field
is constant if its inertia characters at every finite place have
finite image, with ramification at only finitely many places.
Indeed a prime-to-\(\ell\) root of unity is the unique lift of its
residue and the principal units have no torsion in characteristic
\(\ell\). Dividing by the lift of the residual character leaves an
everywhere unramified character into principal units. It factors
through the finite narrow class group by global class field theory,
and is therefore trivial. This argument applies on every finite
number-field extension as well. If a matrix representation becomes
triangular there, and its inertial eigencharacters have finite
image, all its diagonal characters there are constant. On the full
group every eigenvalue is then a root of unity, by raising a group
element to a positive power in the subgroup. Its characteristic
polynomial is the constant lift of its reduction, even if the
subgroup and power depend on the curve.

\begin{lemma}[Odd two-dimensional subquotients on a curve]\label{dp:pan}
Let \(\ell>2\). Suppose a continuous two-dimensional representation
of \(G_{\mathbb Q}\) over a finite extension of \(k'((t))\) has a
stable lattice, is unramified outside finitely many primes, has
finite inertial eigencharacters everywhere, and is triangular on
\(G_{\mathbb Q_\ell}\) over an algebraic closure. Suppose its
residual semisimplification is a sum of two characters of opposite
parity. Then its characteristic polynomials agree with the constant
lifts of their reductions; in particular it is geometrically
reducible.
\end{lemma}
\begin{proof}
After a finite coefficient extension twist by the constant inverse
of the first residual character. Write \(1,\bar\chi\) for the
resulting residual characters. By the preceding character argument
the determinant is the constant lift of \(\bar\chi\). Use its
Teichm\"uller lift \(\chi\) to \(W(k')\), enlarging \(k'\) to the
residue field of the chosen extension. It is an odd finite-order
character and is de Rham locally at \(\ell\). We use the number
field \(\mathbf Q\), an abelian totally real field unramified at
\(\ell\), so the required extension of the residual determinant
to \(G_{\mathbf Q}\) is already given. Choose a finite set \(S\)
of finite primes containing \(\ell\) and every ramified prime
of the twisted representation. Pan's ordinary
pseudodeformation ring for \(G_{\mathbb Q,S}\) at \(1+\bar\chi\),
with determinant \(\chi\), is finite over
\(\Lambda=W(k')[[\mathbb Z_\ell^\times(\ell)]]\), where the group
in brackets is the pro-\(\ell\) part of the local units.
Use the distinguished theorem when
\(\bar\chi|_{G_{\mathbb Q_\ell}}\ne1\). If its restriction is
trivial, use the character-labelled theorem and the actual
subcharacter \(\lambda_1\) of the local triangular representation.
Both local diagonal characters are integral units, since the
inverses are integral too. In a triangular basis
\((r(\sigma)-\lambda_1(\sigma))(r(\tau)-\lambda_2(\tau))=0\)
for local \(\sigma,\tau\); taking the trace after multiplying by
any \(r(g)\) gives the extra ordinary pseudocharacter relation.
The coefficient homomorphism from \(W(k')\) to the normalized
curve DVR is reduction to \(k'\) followed by its constant inclusion.
Thus the trace and, in the second case, the chosen character
define the asserted map of universal rings on passing to the
inverse limit of the Artin quotients. The finiteness assertions
used here hold for every odd prime, including three.
The map on \(\Lambda\) comes from an inertia
character which has finite image in the residue-field units, so
it kills the augmentation of the pro-\(\ell\) group and also
\(\ell\). The image of the finite algebra in this DVR is a finite
field. The trace and determinant therefore equal their residual
values as constants. The characteristic-polynomial criterion for
semisimplification (Brauer--Nesbitt) finishes the proof.
\cite[Theorems~5.1.1(1) and~6.1.1(1) in the published
version, Theorems~5.1.2(1) and~6.1.2(1) on pp.~71 and~95--96
of the author manuscript]{panFM}
\end{proof}

We explain how traces control the marked entries required here. Put
\(\mathcal E=\operatorname{span}\{r(g)\}\) over the algebraic
closure of the fraction field of the curve, using \(G_K\) in rank
three. It is an algebra. It contains the individual coordinate
projectors \(E_{ii}\) in both models. In rank three the marked
difference \(Z=r(g_1)-Jr(g_1)^tJ\) belongs to this algebra since
the second matrix is \(r(cg_1^{-1}c)\). The normalization kills
entries \(01,10\), and antisymmetry under reverse transpose makes
\(Z=\operatorname{diag}(z,0,-z)\). The reduction of \(z\) is
\(\bar\Psi(g_1)-\bar\Psi(g_1)^{-1}\ne0\). Lagrange interpolation
gives all three projectors. In rank four the marked block matrix
\(r(\phi)\) has coprime characteristic polynomials on the blocks
\(\{0,2\}\) and \(\{1,3\}\), by the distinct residual values
between blocks. Polynomial interpolation gives the first block
projector. In that block \(r(g_0)+\theta(g_0)r(g_0)^{-1}\) is
diagonal by normalization, with diagonal difference reducing to
\((\bar\alpha(g_0)-1)(\bar\Psi(g_0)-1)\ne0\). We obtain
\(E_{00},E_{22}\) and then their symplectic adjoints \(E_{33},E_{11}\).
\src{src/Development/AutomorphicArithmetic/AutomorphicVarieties/SymplecticGSp4/EiiUnit.lean}{169}{The four marked projectors}.

If the semisimplification is a sum of characters, a full invariant
flag is consequently a coordinate flag: every stable subspace is
preserved by all \(E_{ii}\). Order the coordinates by this flag,
assigning numbers \(b_i\) so \(b_i>b_j\) for every nonzero
off-diagonal entry in position \(ij\). The reverse position
\((j^*,i^*)\), with \(i^*=n-1-i\), also occurs. This follows by
reverse transpose in rank three and symplectic adjunction in rank
four, both preserving \(\mathcal E\). Therefore \(h_i=b_i-b_{i^*}\)
is strictly decreasing on these directed entries. The functional
\(ij\mapsto h_i-h_j\) factors through their weights. For rank three
this follows from \((h_0,h_1,h_2)=(h_0,0,-h_0)\). For rank four it
is the linear form \((h_0-h_1)\lambda+(h_0-h_2)\mu\) on differences
of the coordinates \((\lambda_i,\mu_i)=(1,1),(0,1),(1,0),(0,0)\).
Thus a nonzero balanced monomial uses diagonal entries only. These
are the diagonal characters on the flag and are constant by class
field theory. In rank three all negative-sign elements have the
matrix of a positive-sign element, by multiplication by \(c\).
Conversely, an absolutely irreducible constituent of dimension \(>1\)
forces a nonconstant characteristic-polynomial coefficient if the
residual semisimplification consists of characters. Such coefficients
are themselves sums of weight-zero monomials.

\begin{proposition}[Remaining characteristic-\(\ell\) curves]\label{dp:curvebranches}
For finiteness of the invariant quotient in rank three it suffices
to exclude curves on which \(r|_{G_K}\) is strongly absolutely
irreducible. In rank four it suffices to exclude curves on which
\(r\) is absolutely irreducible, and those with an absolutely
irreducible two-dimensional constituent on which complex conjugation
is scalar.
\end{proposition}
\begin{proof}
The needed local flags specialize even where the marked inertia
eigenvalues coalesce modulo \(\pi\). In rank three take the closed
incidence scheme in the projective full flag variety requiring
stability under \(H\) and the ordered characters \(d_i\) on \(I_0\).
It has a point over the fraction field of the chart constructed
in Proposition~\ref{dp:rank3chart}. Its proper image contains the
generic point of \(R/I_t\) by the kernel assertion there, so every
geometric fibre has a flag. In rank four use the symplectic flag
variety, stability under the full decomposition group and the four
prescribed inertia characters. The condition \(\mathrm{OnComp}\)
gives a domain carrying the local equations and an augmentation to
\(O\) at the diagonal point, through which the test factors. Thus
the differences \(d_i(s_0)-d_j(s_0)\) are nonzero in its fraction
field. There the projectors give a flag by the calculation in
Proposition~\ref{dp:local}. Properness specializes it to the given
geometric point. The closed
conditions in both assertions are intersections of the closed
stability and action equations, so no finite presentation of the
inertia group is needed. Away from \(\ell\) the fixed residual
or finite character kernel bounds inertia. At \(\ell\) the
inertia characters on these flags have finite image in equal
characteristic, also on the full inertia in rank three since
\([I:I_0]\) is finite. These properties pass to subquotients.

In rank three the algebra \(\mathcal E\) contains \(E_{ij}\)
precisely when some matrix entry in that position is nonzero, by
sandwiching between \(E_{ii}\) and \(E_{jj}\). Its entry relation
is thus transitive. Its equivalence classes for mutual
reachability are the simple diagonal blocks after the classes
are topologically ordered, each supporting a full matrix algebra.
Reverse transpose permutes the classes and reverses the order.
The only possible two-element class is \(\{0,2\}\): its reverse
must be the unique class of size two. The middle singleton is
incomparable to that class, since both classes are fixed by
reversal. Thus a two-dimensional constituent forces the actual
outer block decomposition. The middle character and the determinant
on \(G_K\) have finite inertia and trivial reduction and hence are
one by the character argument. The outer block \(U\) has determinant
one. It extends to \(W\) on \(G_{\mathbb Q}\) with
\(W(c)=D=\operatorname{diag}(-1,1)\), by the matrix identity
\(J_2 U(h)^{-t}J_2=D U(h)D^{-1}\) for \(h\in G_K\), with
\(J_2\) the two-by-two antidiagonal unit matrix. The marked
outer coordinates retain the original integral lattice, which
is also preserved by \(D\). The formula on the two open cosets
is continuous, and \(D^2=1\) gives the group law
as in Proposition~\ref{spreal:end}.
Its residual characters are \(\eta\bar\Psi,\bar\Psi^{-1}\).

On \(I_\ell\subset G_K\) the two inertia constituents of \(W\)
are finite and hence constant, equal to \(\bar\Psi,\bar\Psi^{-1}\)
by reduction. They are distinct since \(\psi_\ell^2=1\) and
\(\ell-1\nmid2e\). The restriction to inertia has a common eigenline
by the local flag. The eigenspace for its character is exactly one
line by distinctness. Its character is invariant under the full
decomposition group because it is the restriction of one of the
global residual characters. The line is therefore stable under
that group. Away from \(\ell\) the extension adds only the finite
quadratic character. Lemma~\ref{dp:pan} now excludes an absolutely
irreducible outer block. If an irreducible three-dimensional
representation becomes reducible over a finite number-field
extension, pass to its normal core. Its restriction there is
semisimple: the sum of the simple submodules is nonzero and is
stable under the full group. The isotypic components are permuted
transitively, with equal multiplicities. Primality of dimension
three makes every simple constituent on this subgroup a character.
The preceding potential triangularity argument gives constant
characteristic polynomials, contradicting absolute irreducibility.

In rank four any reducible symplectic similitude representation
has simple constituents of dimension at most two. A stable line
lies in its stable three-dimensional perpendicular and a stable
hyperplane gives such a line by duality. A stable two-plane gives
two pieces of dimension two directly. Every two-dimensional
constituent inherits a local flag by intersection and quotient,
and its residual semisimplification is a pair from
\((\bar\theta,\bar\Psi,\bar\alpha,1)\). Lemma~\ref{dp:pan} excludes
all opposite-parity pairs. The equal-parity pairs are precisely
\(\{\bar\theta,\bar\alpha\}\) and \(\{\bar\Psi,1\}\), with parity
read on the two coordinate pairs \(\{0,2\}\), \(\{1,3\}\).
The preceding marked-entry argument and Lemma~\ref{dp:detect}
complete both reductions.
\end{proof}

The rank-four source separates this issue from its characteristic-zero
purity and automorphy assertions. Its first deduction of finiteness
uses the positive-characteristic exclusion \(\mathrm{Con}(D)\), whose
test includes the nonsplit extension with irreducible even block on
indices \(1,3\). In the tower the ideal \(J\) is required to satisfy
\(\mathfrak m_{X.B}^n\subset J+\mathfrak m_{O'}X.B+
(\psi_{v,i}(s)-1:s\in I_v)\). This gives finiteness on the
fixed-inertia special fibre \emph{of the quotient by \(J\)}.
Using it for a given curve also requires showing that the curve
annihilates \(J\) on the tower. Its localization clause for
irreducible, locally distinguished curves already assumes this
annihilation. The characteristic-zero seed and purity clauses
alone impose no equation on a purely characteristic-\(\ell\)
component. In rank three the source uses a common finite-index
subgroup for constancy of all the special-fibre characteristic
polynomials. The curve criterion shows why it is enough to
settle the remaining representations of Proposition~\ref{dp:curvebranches}
one curve at a time.
\src{src/Development/AutomorphicArithmetic/AutomorphicVarieties/CurveDVRSpecialization/GfResZpow.lean}{179}{Rank-three finiteness from special-fibre characteristic polynomials};
\src{src/Development/AutomorphicArithmetic/AutomorphicVarieties/ModularityArgument/ThresholdArguments/DenseThresholdCases/VWinPm.lean}{75}{The characteristic used in each of the three conclusions};
\src{src/Development/AutomorphicArithmetic/AutomorphicVarieties/ModularityArgument/ThresholdArguments/DenseThresholdCases/PotentialModularityTower.lean}{38}{Finiteness and localization conditions on the tower ideal}.


\subsection{Intersection of the weighted equations}
\label{newcoh:orbitsection}

We prove the geometric inequality used above. Let \(O\) be a complete
discrete valuation ring with uniformizer \(\pi\) and fraction field \(E\)
of characteristic zero. Let \(K\) be an algebraically closed extension
of \(E\) with a real valuation extending the valuation \(v(\pi)=1\).
Write \(A=\bigoplus_{j\in\mathbf Z}A_j\)
and \(B=A_0\). Assume \(B\) is a complete Noetherian local
\(O\)-algebra with the same residue field and local structure map,
and \(A\) is generated over \(B\) by finitely many elements of
weights \(-2,-1,1,2\). An \(O\)-point \(x\) kills all nonzero
weights. Write \(J_x=\ker x\) and use the saturated cotangent module
of Lemma~\ref{newcoh:cotangent}, taking in weight zero the image of
\(J_x\cap B\).

Suppose there is a weight-preserving formal chart
\(A\to E[[X_1,\ldots,X_n]]\) with constant term \(x\), whose linear
terms on \(J_x\) span all coordinate directions over \(E\) and have
kernel the saturation of \(J_x^2\) there. The table gives the number
of chart variables and of homogeneous equations defining an ideal \(I\).
The neutral equations vanish at \(x\).
\[
\begin{array}{c|rrrrr}
 \text{weight}&-1&-2&1&2&0\\ \hline
 \text{variables}&r&s&r'+1&s'+1&m\\
 \text{equations}&r+1&s&r'&s'&m\\
 \text{names}&F&G&P&H&Z
\end{array}
\qquad r\ge1 .
\]
Assume the equations in weights \(-1,-2,0\) span the corresponding
rational cotangent spaces. Set \(X=\operatorname{Spec}(A/I)\).
A geometric point is called two-sided if some positive and some
negative weight both have nonzero value. A positive or negative
one-sided point has nonzero values on precisely one of the two signs.
An invariant value of a point means its restriction to \(B\).

\begin{theorem}[Integral degree of the two-sided orbits]
\label{newcoh:orbitdegree}
Suppose \(B/(I\cap B)\) is finite as an \(O\)-module. Assume every
fixed \(K\)-point of \(X\) other than \(x_K\) has at most one sign of
one-sided points with its invariant value, and the geometric tangent
space of \(X_K\) at each two-sided point has dimension at most one.
Let \(a\in B\) vanish at \(x\) and be nonzero at every two-sided point.
There are finitely many two-sided \(K^\times\)-orbits, and for any finite
set \(T\) of two-sided points meeting every such orbit,
\[
 \operatorname{length}_O Q_{-1}\ \le\ 2\sum_{\xi\in T}v(\xi(a)).
\]
\end{theorem}

\begin{proof}
\emph{The chart and the three gradings.}
Every \(A_j\) is finite over \(B\). Indeed, in a monomial in the
chosen generators one may remove weight-zero submonomials until
none remains. Among exponent vectors of fixed weight there are
only finitely many minimal elements in the componentwise order
by Dickson's lemma, and the removed factors belong to \(B\).
The weight \(-1\) cotangent lattice is thus finite over \(O\)
since \(B\) acts through \(x\) on it. It is torsion-free by
saturation and has rank \(r\) by the chart. Apply Smith normal
form to the map from the \(r+1\) equations to this lattice.
After invertible \(O\)-linear changes of equations and lattice basis
there are \(y_i\in A_{-1}\) and \(c_i\ge0\) such that, modulo
the saturation of \(J_x^2\),
\[
 F_i=\pi^{c_i}y_i\quad(1\le i\le r),\qquad F_*=0,\qquad
 \operatorname{length}_O Q_{-1}=\sum_i c_i .
\]
Here \(F_*\) denotes the last equation. In the formal chart replace
the \(m\) neutral coordinates by the \(Z\)'s using their invertible
linear part and the formal inverse function theorem. Setting those
coordinates to zero gives a chart \(\phi\) with the four nonzero
weights and the same linear coordinates in these weights. The inverse
is weight-preserving by uniqueness of the formal inverse. We may
arrange \(\phi(a)\ne0\). If its image is zero, add \(uF_i\) to \(a\),
choosing \(u\in A_1\) and \(i\le r\) whose images have nonzero linear
parts. The product has nonzero image in a power series domain and
changes no value on \(X\).

Form the following standard \(\mathbf N^3\)-graded subalgebra
of \(A[t_0,t_1,t_2]\), with the ordinary weight in \(A\) kept
separate from the three exponents of \(t\):
\[
 U=A_{-2},\qquad V=A_2,\qquad W=BF_*^2+U^2V,\qquad
 R=B[Ut_0,Vt_1,Wt_2].
\]
Let \(R'\) be its image under the coefficient map \(\phi\), so
\(a\) is a nonzerodivisor on \(R'\). Denote by \(\mathcal S\) the
following list in \(R\) or any of its quotients:
\[
 (F_i^2t_0)_{i\le r},\quad (G_i t_0)_{i\le s},\quad
 (P_i^2t_1)_{i\le r'},\quad (H_i t_1)_{i\le s'},\quad F_*^2t_2 .
\]
Put \(M=R'/(\mathcal S)\). We first determine its horizontal
support, including at the boundary where one sign vanishes.
\src{src/Development/AutomorphicArithmetic/AutomorphicVarieties/CommutativeAlgebra/GradedModulesLength/ToDegInjective.lean}{187}{multigraded algebra and marked equations}.

\emph{Support of high graded pieces.}
Call a prime of \(R\) relevant if it misses some generator in each
of the three standard directions. Let \(\mathfrak p\) be such a prime
above the relations of \(M\), and \(\mathfrak b=\mathfrak p\cap B\).
Choose surviving symbols \(ut_0,vt_1,wt_2\). Since \(W\subset U\),
the identity \((F_*^2t_0)(wt_2)=(wt_0)(F_*^2t_2)\) puts
\(F_*^2t_0\) in \(\mathfrak p\). We have \(I\cap B\subset\mathfrak b\).
Indeed, for a marked coefficient \(h\) of weight \(-2\) and \(k\)
of weight \(2\), multiply \(kh\) by the surviving product
\((ut_0)(vt_1)\). The result is \((ht_0)(kuv t_1)\) in
\(\mathfrak p\). The positive case interchanges the signs. For an
equation of weight \(1\) or \(-1\) apply this argument to its square
and to the square of its complementary coefficient. Neutral equations
are already zero in \(R'\). Taking weight-zero parts of all
multiples of the generators of \(I\) proves the inclusion.

Put \(D=UV\subset B\). We claim that if \(D\subset\mathfrak b\)
then \(\mathfrak b\) is the maximal ideal of \(B\). Otherwise,
since \(B/(I\cap B)\) is finite over \(O\), the point is horizontal
and \(B/\mathfrak b\) embeds in \(K\). We construct one-sided points
of both signs with this invariant value. For \(n\ge1\),
\(A_{2n}=V^n\) and \(A_{-2n}=U^n\). To see this, remove a positive
generator of weight two or a pair of positive generators of weight
one from a monomial of weight \(2n\) until the remaining weight is
two. Negating the weights gives the other equality.
On the even-weight subalgebra define a map to \(\kappa(\mathfrak p)\)
by \(h\mapsto ht_1^n\bmod\mathfrak p\) for \(h\in A_{2n}\), \(n\ge0\),
and send negative weights to zero. This is multiplicative. A mixed
product with nonnegative weight and a negative factor has its
coefficient in \(D V^n\) at weight \(2n\), by the two power
identities. The map kills the appropriate marked even equations and
squares, and takes \(v\) to a nonzero element. The full algebra is
integral over the even-weight subalgebra, because the square of
each odd homogeneous element belongs to it. Lying over extends the
map to a field extension, where all original equations vanish.
Over the invariant value under consideration these equations, the
vanishing of the negative generators and \(v\ne0\) describe a
nonempty finite-type locus over \(\kappa(\mathfrak b)\), hence they
have a \(K\)-solution. The same construction with \(t_0\) gives the
negative point. Taking the scaling limit of either point toward
zero gives the same fixed point. Away from \(x_K\) this contradicts
the hypothesis on the signs. Above \(x_K\) a negative point would
give a curve through \(x_K\), taking \(h\in A_{-j}\) to
\(\xi(h)T^j\), \(j>0\). Choose the least \(j\) occurring with a
nonzero coefficient. Since \(A\) is generated in the stated weights,
\(j=1\) or \(2\). All of \(J_x^2\), and hence its saturation, vanish
to order at least \(2j\) on the curve. The \(F\)'s or \(G\)'s span
the rational cotangent in weight \(-j\) and vanish identically on
the curve, contradicting minimality. This proves the claim.

Here and below \(d\gg0\) means all three components are sufficiently
large. The two prime inclusions just proved give
\[
 (I\cap B)^n M_d=0,\qquad \pi^{n'}M_d\subset D M_d
 \qquad(d\gg0)
\]
for some \(n,n'\). Indeed, multiply the respective ideals \(I\cap B\)
and \((\pi)\) by the irrelevant ideal (the product of the three
direction ideals). The products lie in the radicals of the defining
ideals of \(M\) and \(M/DM\), respectively. Noetherianity gives
powers, and sufficiently high graded pieces are generated by the
corresponding powers of the irrelevant ideal times lower pieces.
The first inclusion makes \(M_d\) finite over \(O\), by the finite
filtration of \(B/(I\cap B)^n\) over \(B/(I\cap B)\). The second
and Nakayama's lemma put its generic support where some
\(u v\ne0\), \(u\in U,\ v\in V\).
\src{src/Development/AutomorphicArithmetic/AutomorphicVarieties/CommutativeAlgebra/GradedModulesLength/SameDegreeZeroOfRestriction.lean}{17}{exclusion of the one-sided boundary}.

\emph{The generic multiplicities.}
After inverting \(uv\), \(u\) is a homogeneous unit in \(A\), and
the three coefficient modules \(U,V,W\) are freely generated over
the localized \(B\) by \(u,u^{-1},u\), respectively. If
\(d=(i,j,k)\) with \(i,j,k\ge1\), taking coefficients identifies
\(R_d\) there with \(A_{2(j-i-k)}\). Multiplication by a marked
section is multiplication by its coefficient between the corresponding
weight spaces. Hence \(M_d\) is a quotient of this weight space
modulo the ideal \(I_2\) obtained from \(I\) by squaring the
\(q=r+1+r'\) equations \(F,P\), and retaining \(G,H,Z\).
Since any product of \(q+1\) generators of \(I\) contains either
a retained equation or a repeated squared generator, \(I^{q+1}\subset I_2\).
In particular all weight spaces occurring in the successive
thickenings are modules over \(B/(I\cap B)^{q+1}\), which is finite
over \(O\) by the same ideal-power filtration. After base change
their local multiplicities may thus be computed over the local
factors of a finite-dimensional \(K\)-algebra.

Before squaring, each weight has geometric multiplicity at most
one at an invariant point where \(uv\ne0\). In fact
\(X_K[(uv)^{-1}]\) is of finite type by finiteness of the invariant
algebra. More precisely, taking the degree-zero summand commutes
with extension of scalars and with localization at the weight-zero
element \(uv\), so its invariant ring is
\(\big((B/(I\cap B))\otimes_O K\big)[1/(uv)]\).
Each closed point of \(X_K[(uv)^{-1}]\) belongs to a one-dimensional
torus orbit, and its tangent space has dimension at most one.
Thus this open scheme is smooth of dimension one there. Its orbits
are open in their components and every closed point lies in such an
orbit, so each connected component is a single orbit. Equivalently,
each orbit is closed here since the invertible coordinate \(u\)
scales by weight \(-2\). The component idempotents are invariant,
so they distinguish the components on the indicated invariant base.
Each weight space on one orbit embeds by pullback into
the corresponding weight space on \(K[t,t^{-1}]\), of dimension at
most one. A finite collection of opens \(uv\ne0\) covers the
two-sided locus, since nonzero weights can be witnessed by
generators and weight-one generators can be squared. This also proves
finiteness of its orbits.

Replacing a homogeneous equation \(f\) of weight \(h\) by \(f^2\)
at most doubles a bound valid for all weights. Indeed, the kernel of
\((L/f^2L)_w\to(L/fL)_w\) is a quotient of \((L/fL)_{w-h}\)
under multiplication by \(f\). Iterate after localizing at the given
geometric invariant point. The multiplicity of \(M_d\otimes_O K\)
at that point is at most \(2^{r+1+r'}\).
On the finite \(O\)-module \(M_d\) put
\(\theta(a,M_d)=\operatorname{length}_O(M_d/aM_d)
-\operatorname{length}_O(M_d[a])\).
Its torsion submodule contributes zero to this difference.
On its free quotient the difference equals the valuation of the
determinant of \(a\), by Smith normal form. Over \(K\) this
determinant is the product of the values of \(a\) at invariant
points counted with the above multiplicities, by decomposition of
the finite-dimensional module into its generalized eigenspaces.
These values are nonzero and integral over \(O\). It follows that
\begin{equation}\label{newcoh:upperorbit}
 \theta(a,M_d)\le 2^{r+1+r'}\sum_{\xi\in T} v(\xi(a)).
\end{equation}
\src{src/Development/AutomorphicArithmetic/AutomorphicVarieties/CommutativeAlgebra/GradedModulesLength/TwoSigned.lean}{17}{multiplicity upper bound for the three gradings}.

\emph{The component through the origin.}
In a chart monomial let \(b_-,b_+\) count its negative and
positive weights with absolute values, so its weight is \(b_+-b_-\).
In degree \(d=(i,j,k)\) of \(R'\) keep the terms with
\[
 (b_-,b_+)=(2i+4k,\,2j+2k).
\]
These are the terms of least possible negative weight. The bounds
for \(U,V\) follow from their weights. For \(W\), the term of
bidegree \((1,0)\) in \(F_*\) is zero since \(F_*\) has zero linear
part, giving \(b_-\ge4\) for its square. The same bound holds on
\(U^2V\). Multiplication of lowest terms therefore defines a
graded quotient \(C\) of \(R'\), retaining the three exponents to
distinguish components. On \(C\) the ring \(B\) acts through
\(x:B\to O\), so \(C\) is a quotient of \(R'/aR'\).
Each \(C_d\) is a finite torsion-free \(O\)-module inside the
polynomials of the indicated bidegree.

Over \(E\) these polynomials all occur. The leading terms in
weight \(-1\) contain all linear coordinates of that weight after
extending scalars. Those from weight \(-2\) span its linear
coordinates modulo quadratics in the weight \(-1\) coordinates.
The quadratics themselves are products of the preceding lifts,
so \(U\) supplies all pure monomials of negative weight two. The
same argument applies to \(V\) on the positive side. A monomial
with even total weight on either side splits into blocks of weight
two, by pairing the weight-one variables. The required blocks form
\(i\) factors from \(U\), \(j\) from \(V\), and \(k\) from
\(U^2V\subset W\). Thus the length of \((C/\pi C)_d\) is
\[
 h(d)=N_{r,s}(i+2k)\,N_{r'+1,s'+1}(j+k),\qquad
 N_{b,c}(z)=[t^{2z}](1-t)^{-b}(1-t^2)^{-c}.
\]
In \(C\) the first \(r\) marked sections are
\(\pi^{2c_i}Y_i\), where \(Y_i\) is the marked leading term of
\(y_i^2\). The Smith relations give this equality because any
weight \(-1\) element in the saturated square has zero linear term.
\src{src/Development/AutomorphicArithmetic/AutomorphicVarieties/CommutativeAlgebra/GradedModulesLength/ReesGradedPieceSurjection.lean}{199}{Hilbert function of the leading terms}.

\emph{Euler characteristics.}
For a finite graded module \(L\) over \(R\), write
\(\chi(\mathbf{s};L)_d\) for the alternating sum of the \(B\)-lengths
of the degree \(d\) Koszul homology for a list of homogeneous
sections. We use this when the high pieces of \(L/(\mathbf{s})L\)
have finite length. Then the high pieces of the homology have finite
length as well, since these finite graded modules are supported on
the same closed locus. Equivalently, off the irrelevant locus all
primes of this support lie over the closed point of \(B\). All
lengths of finite-length \(B\)-modules here agree with their
\(O\)-lengths because the residue fields agree.

For completeness, \(\chi(\mathbf{s};L)_d\ge0\) for \(d\gg0\).
Use induction on the list and a homogeneous prime filtration of \(L\),
using additivity of Koszul Euler characteristics in short exact
sequences. Irrelevant factors vanish in high degree. In a relevant
domain factor, a nonzero first section is regular and the calculation
passes to its quotient with one fewer section. If the first section
vanishes, its two-term complex gives the difference of the Euler
characteristics of the shorter list in degrees \(d\) and
\(d-\deg s_1\). This difference is also the Euler characteristic
of the shorter list on the quotient by any nonzero section of that
degree, by the multiplication exact sequence. Such a section exists
in a relevant domain by taking a product of surviving standard
generators, or \(1\) in degree zero. The shorter list already has
the required support in this case. This proves positivity and,
by additivity, monotonicity on passing to a quotient whenever the
original list has the stated support.

The support condition holds for \((\mathcal S,a)\) on \(R'\) by
the preceding description of \(M_d\) and the nonvanishing of \(a\)
on its generic support. Koszul complexes are invariant up to
isomorphism under permutation. Since \(a\) is regular on \(R'\),
monotonicity and the one-section kernel--cokernel formula give
\[
 \chi(\mathcal S;C)_d
 \le\chi(\mathcal S,a;R')_d
 \le\chi(a;M)_d=\theta(a,M_d)\qquad(d\gg0).
\]
For the second inequality successively replace each of the first
sections by its ordinary quotient. At each step the omitted term
is the Euler characteristic of the remaining list on the shifted
kernel of that multiplication and is nonnegative by the preceding
argument.

In \(C\), split each \(\pi^{2c_i}Y_i\) into its two factors using
additivity of the one-section Koszul class for a product. In fact,
writing \(e_0=(1,0,0)\), use the degree-preserving composite
\[
 C(-e_0)\xrightarrow{\,Y_i\,}C\xrightarrow{\,\pi^{2c_i}\,} C .
\]
The three mapping cones for the first map, the composite, and the
second map form an exact triangle in that order. Tensor with the
remaining Koszul complex, whose terms are graded free, and take
degree \(d\). Taking Euler characteristics expresses the contribution
of the composite as the sum of those of the two factors, with
both summands taken in degree \(d\). The same triangle splits the
degree-zero power into \(2c_i\) copies of the section \(\pi\).
Each factor list has support contained in
the original intersection. In each term with a single \(\pi\),
regularity of \(\pi\) on \(C\) passes to \(C/\pi C\).
Its terms already have finite length degree by degree, so the Euler
characteristic is simply the alternating sum of the lengths of its
terms, independent of the actual sections of the remaining degrees.
Peeling off the powers successively leaves one final nonnegative
Euler characteristic. With \(\Delta_i,\Delta_j,\Delta_k\) denoting
unit backward differences in the indicated variable this gives
\[
 \chi(\mathcal S;C)_d\ge
 2\left(\sum_i c_i\right)
 \Delta_i^{\,r+s-1}\Delta_j^{\,r'+s'}\Delta_k h(d).
\]
For \(b\ge1\), partial fractions at \(t=1,-1\) show that
\(N_{b,c}(z)\) is eventually polynomial of degree \(b+c-1\)
with leading coefficient \(2^{b-1}/(b+c-1)!\). Indeed the pole
at \(1\) has order \(b+c\) and leading factor \(2^{-c}\);
the pole at \(-1\) has order at most \(c\), of lower degree.
The indicated \(i\)-difference therefore gives \(2^{r-1}\) on
the first factor, and the \(j\)-difference gives a linear polynomial
of slope \(2^{r'}\) in \(j+k\) on the second. Its \(k\)-difference
is \(2^{r'}\). We have proved
\[
 2^{r+r'}\sum_i c_i\le \theta(a,M_d)
 \le 2^{r+1+r'}\sum_{\xi\in T}v(\xi(a))
 \qquad(d\gg0).
\]
Dividing gives the theorem.
\end{proof}
\src{src/Development/AutomorphicArithmetic/AutomorphicVarieties/CommutativeAlgebra/GradedModulesLength/AlgMemCoef.lean}{92}{comparison of the two Koszul calculations}.

For the arithmetic application the valuation on a finite coefficient
extension has \(v(z)=-e(O/\mathbf Z_\ell)\log|z|_\ell/\log\ell\).
Thus Lemma~\ref{newcoh:cotangent} and Theorem~\ref{newcoh:orbitdegree}
give the factor \(2\) in the normalized trace inequality once the
rank-four deformation datum supplies the chart, equations, one-sided
exclusion, infinitesimal rigidity and trace function with the stated
nonvanishing. Establishing these hypotheses for all cases II and III
and bounding the resulting automorphic family remain the global steps.
 

\subsection{One integer for all coefficient primes in rank three}\label{rk3:sec}

Fix an imaginary quadratic field \(K\), a level \(N\ge1\), a twist
\(e\ge1\), a prime \(p\nmid30N|\operatorname{disc}K|\), and a place
\(w\mid p\). Write \(q=|\mathcal O_K/w|=p^h\), \(h=1\) or \(2\).
In this subsection \(q\) is a residue norm, independent of the notation
for the first auxiliary prime elsewhere. For \(H\le(\mathbf Z/N)^\times\),
let \(J\) be the complex characters modulo \(N\) with kernel exactly \(H\).
If the quotient is cyclic then \(|J|=\varphi([(\mathbf Z/N)^\times:H])\);
otherwise \(J\) is empty. This follows by choosing the image of a generator
under a faithful embedding in \(\mathbf C^\times\). Put \(z_j=j(q)^{-1}\).
The multiset of the \(z_j\) is invariant under every automorphism of
\(\overline{\mathbf Q}\).

The two spectral lists used at \(w\) are the irreducible rank-three,
polarized automorphic representations with normalized weight sequence
\((e,0,-e)\), and the sums \(W|_{G_K}\oplus1\) from a rank-two newform
of weight \(2e+1\), with the normalization giving weights \((e,-e)\).
The conductor bound for the second list involves \(N|\operatorname{disc}K|\).
The following lemma isolates the exact use of classification by finite
lists. The same trace can belong to several distinct entries.

\begin{proposition}[Counting the classical spectral lists]\label{rk3:count}
Fix \(F=\mathbf Q\) or an imaginary quadratic field and \(n\ge2\).
Let \(K_F(M)\) be the kernel of reduction modulo \(M\) in
\(\mathrm{GL}_n(\widehat{\mathcal O}_F)\).
Consider regular algebraic cuspidal representations of
\(\mathrm{GL}_n(\mathbb A_F)\) with \(K_F(M)\)-fixed vectors and
one prescribed infinitesimal character at each embedding. Suppose the
pairwise differences of its Harish--Chandra parameters are at most \(H\)
in absolute value. The number of isomorphism classes is at most
\[
 C_{F,n}(3M)^{n^2[F:\mathbf Q]}
       (1+H)^{[F:\mathbf Q]n(n-1)/2}.
\]
On the split real centre take the prescribed character and the inverse
character on the coefficient system; the product descends to the quotient.
\end{proposition}
\begin{proof}
The regular algebraic cuspidal condition gives nonzero relative Lie
algebra cohomology with the irreducible coefficient module of the
prescribed infinitesimal character, using the dual if necessary.
Take the connected maximal compact at infinity, together with the
split central factor. Distinct cuspidal representations give distinct
summands in cuspidal cohomology. This embeds in ordinary cohomology by
the cuspidal summand in Franke's theorem. At finite level take compact-open
invariants, an exact functor over \(\mathbf C\) by averaging. Since
\(K_F(3M)\subseteq K_F(M)\), every representation counted contributes
at least one dimension to the sum over all cohomological degrees at
level \(3M\).
Grobner--Raghuram,
Thms.~6.3 and 7.6, Cor.~7.7 and Prop.~7.16~\cite{grobnerRaghuram}.

At level three there are finitely many connected components. Indeed the
determinants give finitely many ray classes and strong approximation
for \(\mathrm{SL}_n\) gives the component description, up to the
finite group of real signs. The stabilizers are torsion-free: in each
adelic lattice their action on its underlying integral lattice is
congruent to the identity modulo three. The elementary binomial
argument in Theorem~\ref{newcoh:lemmaK} applies. The Borel--Serre
compactifications of these components have the homotopy type of the
interior and are compact manifolds with corners. Fix finite
triangulations with \(C_{F,n}\) simplices in total. The finite cover
at \(3M\) has at most \([K_F(3):K_F(3M)]\) lifts per simplex, at most
\((3M)^{n^2[F:\mathbf Q]}\) by counting all residue matrices modulo
\(3M\). A cellular cochain group with local coefficients has one
copy of the fibre per cell. Thus the sum of the Betti numbers is
bounded by the number of lifted simplices times the fibre dimension.
Borel--Serre,
\emph{Corners and arithmetic groups}~\cite{borelSerre}.

At one embedding write the ordered infinitesimal character as
\(\mu+\rho\), with \(\mu\) the highest weight. Weyl's formula for
the coefficient dimension multiplies its pairwise positive differences
divided by \(j-i\) over \(i<j\). There are \(n(n-1)/2\) factors,
each at most \(1+H\); changing a determinant twist leaves the
differences fixed. Taking the product over all embeddings proves the
bound. For the present rank-three triples the infinitesimal character
is \((e,0,-e)\). Subtracting \(\rho=(1,0,-1)\) gives dimension \(e^3\)
at each of the two embeddings, hence \(e^6\) in all. In rank two
\((e+1/2,1/2-e)-\rho=(e,1-e)\) has dimension \(2e\).
\end{proof}

\begin{corollary}[Discriminant and level in the rank-three count]\label{rk3:volume}
Let \(D=|\operatorname{disc}K|\) for an imaginary quadratic field \(K\).
Let \(S\) consist of the two classical lists of
Proposition~\ref{rk3:count}, at principal levels \(N\) over \(K\)
and \(ND\) over \(\mathbf Q\), respectively, with one entry per
system. For integers \(N,e\ge1\),
\[
 |S|\ll h_K D^4 N^{18}e^6+(ND)^4e
       \ll D^5N^{18}e^6 .
\]
The constants are absolute. Consequently the integer of
Lemma~\ref{rk3:integer}, under its trace hypotheses, satisfies
\[
 \log|C_p|\ll |J|D^5N^{18}e^7\log p .
\]
\end{corollary}
\begin{proof}
Put \(L=3N\) and \(X=\mathrm{SL}_3(\mathbf C)/\mathrm{SU}(3)\)
with a fixed invariant Riemannian metric. As in
Proposition~\ref{rk3:count}, pass to principal level \(L\) and
divide out the real split centre. At infinity
\(\mathrm{GL}_3(\mathbf C)/(\mathbf R_{>0}\mathrm U(3))\simeq X\).
Determinant and strong approximation for \(\mathrm{SL}_3\) identify
the connected components with the finite ray class group modulo
\(L\mathcal O_K\). In fact the determinant image of the principal
compact open is the subgroup of finite integral idelic units
congruent to one at that modulus, as is seen on diagonal matrices.
The map to the ordinary class group has kernel a quotient of
\((\mathcal O_K/L\mathcal O_K)^\times\). There are thus at
most \(h_K L^2\) components.

For a finite adelic representative \(g\) of a component put
\(\Lambda_v=g_v\mathcal O_{K_v}^3\). A member of its arithmetic
stabilizer preserves these lattices and acts trivially modulo \(L\).
Its determinant is a unit \(\epsilon\in\mathcal O_K^\times\) with
\(\epsilon\equiv1\bmod L\). If \(\epsilon\ne1\), then
\(0<|\operatorname{N}_{K/\mathbf Q}(\epsilon-1)|\le4\), whereas
divisibility by \(3\) makes that norm divisible by \(9\).
Hence the stabilizer is the level-\(L\) subgroup of
\[
 \Gamma_g^1=\mathrm{SL}_3(K)\cap
          \prod_{v<\infty}\operatorname{Stab}_{\mathrm{SL}_3(K_v)}
                                     (\Lambda_v).
\]
Its index in \(\Gamma_g^1\) is at most \(L^{16}\). Indeed, in a
basis of each \(\Lambda_v\) the image of reduction at
\(\mathfrak p_v^b\Vert L\mathcal O_K\) has size at most
\[
 |\mathrm{SL}_3(\mathcal O_{K_v}/\mathfrak p_v^b)|
 =q_v^{8(b-1)}|\mathrm{SL}_3(\mathbf F_{q_v})|\le q_v^{8b}.
\]
Reduction is surjective, since a lift with unit determinant can be
corrected by scaling one row. Each successive kernel is the additive
space of trace-zero matrices over the residue field. Multiplication
over \(v\) gives \(\operatorname{N}(L\mathcal O_K)^8=L^{16}\).

Prasad's covolume formula applies to the split simply connected
\(K\)-group \(\mathrm{SL}_3\). For every \(g\) each local stabilizer
is hyperspecial, with integral model \(\mathrm{SL}(\Lambda_v)\)
and reductive special fibre \(\mathrm{SL}_3\). They form a coherent
collection, equal to the standard one almost everywhere. Strong
approximation at the noncompact complex place gives the principal
arithmetic group appearing in the formula. Its Haar covolume in
\(\mathrm{SL}_3(\mathbf C)\) is
\[
 c_\infty D^{8/2}\prod_{v<\infty}
   \frac{q_v^8}{|\mathrm{SL}_3(\mathbf F_{q_v})|}
 =c_\infty D^4\zeta_K(2)\zeta_K(3).
\]
Here the Tamagawa number is one and the archimedean constant is
fixed for the single complex place and the chosen Haar measure.
In particular there is no dependence on the ideal class of \(g\).
The local zeta factors give \(\zeta_K(j)\le\zeta(j)^2\) for \(j>1\).
Prasad,
\emph{Volumes of \(S\)-arithmetic quotients of semi-simple groups},
\S\S3.3, 3.7 and 3.11~\cite{prasadVolume}.

The level-\(L\) subgroup acts faithfully and without torsion, by
its action on the underlying integral lattice modulo three. Its
quotient of \(X\) is noncompact since \(\mathrm{SL}_3\) is \(K\)-split
of positive rank. Its Riemannian volume is at most a fixed constant
times \(D^4L^{16}\), by the covolume and index bounds. Gelander's
theorem gives a homotopy-equivalent simplicial complex with at most
a fixed multiple of this volume in vertices and uniformly bounded
vertex degree. A vertex of degree \(b\) belongs to at most \(2^b\)
simplices, so the same linear bound holds for all simplices.
Local-coefficient cochains contribute at most \(e^6\) dimensions per
simplex. Summing over the \(h_K L^2\) possible components and
using cuspidal cohomology as in Proposition~\ref{rk3:count} proves
the first term. The second follows from the same proposition over
the fixed field \(\mathbf Q\) in rank two with \(M=ND\) and
coefficient dimension \(2e\).
Gelander,
\emph{Homotopy type and volume of locally symmetric manifolds},
Theorem 1.5(1)~\cite{gelanderVolume}.

Minkowski supplies each ideal class with an integral ideal of norm
at most \(\sqrt D\). There are at most \(\tau(m)\) ideals of norm
\(m\), by the quadratic prime splitting possibilities. Consequently
\(h_K\le\sqrt D(1+\log\sqrt D)\le D\).
Finally \(q\le p^2\), so the per-factor logarithmic cost
\(2e\log q+1\) of Lemma~\ref{rk3:integer} is \(O(e\log p)\).
\end{proof}


\begin{lemma}[Size and denominators of the classical traces]\label{rk3:purity}
For the conjugate self-dual rank-three cuspidal systems just described,
take principal level \(N\) and the weights for the inverse cyclotomic
twist \(m=(e-1,-1,-e-1)\) in the convention that the labelled jumps
are \(-m_i\). For the rank-two summands take principal level
\(N|\operatorname{disc}K|\) and weights \((e,-e)\). Assume the usual
classical polarized cuspidal realizations with these levels and weights.
The resulting traces at inverse arithmetic Frobenius, indexed with
multiplicity by the systems, are algebraic, stable as a list under
conjugation when all such systems are taken, and satisfy
\(|t_s|\le3\), \(q^et_s\) an algebraic integer.
\end{lemma}
\begin{proof}
Use geometric Frobenius, the inverse of arithmetic Frobenius. In rank
three the spherical parameters \(\alpha_i\) in unitary normalization
have absolute value one. Indeed the complex representation is cuspidal,
cohomological and conjugate self-dual over a CM field, so Caraiani's
temperedness theorem applies. By the normalization of
Proposition~\ref{newgal:classical}, the geometric eigenvalues of the
inverse cyclotomic twist are \(b_i=q\alpha_i\). The cyclotomic
character at geometric Frobenius is \(q^{-1}\), so restoring the
original representation gives \(t_s=q^{-1}\sum_i b_i\).
The cohomological conjugation theorem permutes the systems: all
embeddings have the same weight triple, the level is defined
arithmetically, and conjugation commutes with the contragredient and
field-automorphism relations on the finite parts. Strong multiplicity
one identifies the conjugate cuspidal representation by its finite
part. The spherical parameters in algebraic normalization are algebraic
and conjugate accordingly. Here the factor \(q\) is rational.
Caraiani,
\emph{Local-global compatibility and the action of monodromy on
nearby cycles}, Thm.~1.2~\cite{caraianiNearby};
Grobner--Raghuram,
Thm.~7.17 and Remark~7.18~\cite{grobnerRaghuram}.

At a coefficient prime \(v\nmid p\), a stable integral lattice for
the corresponding member of the compatible system shows that both the
eigenvalues of Frobenius and their inverses are integral: its matrix
and its inverse preserve the lattice. At the coefficient prime \(p\),
the local representation is crystalline since \(p\) divides neither
the level nor the discriminant. Let \(q=p^h\). The normalized Newton
slopes of its filtered isocrystal are \(v_p(b_i)/h\), with
\(v_p(p)=1\). They lie between the smallest and largest labelled
filtration jumps. To see the lower bound, take the isoclinic subobject
of minimum slope with its induced filtration, after extending
coefficients as necessary. Weak admissibility gives \(t_N\ge t_H\)
there, and each jump is at least the minimum of the ambient jumps.
Divide by its rank, with the usual average over embeddings. Apply the
same argument to the dual for the upper bound. In rank three the
endpoints are \(1-e,1+e\). Hence \(q^{e-1}b_i\) is \(p\)-integral.
The argument applies to every embedding of the coefficient field.
Integrality at every finite place gives the assertion for \(q^et_s\).

For a normalized rank-two member \(W\) over \(\mathbf Q\) the jumps
are \(\{-e,e\}\). The purity weight in the polarized cuspidal theorem
is zero, since the reflected pair of jumps about \(w/2\) must be that
same pair. Its two geometric Frobenius eigenvalues \(b_1,b_2\) at a
good \(p\) therefore have absolute value one under every embedding.
The \(p\)-adic slope calculation gives \(v_p(b_i)\ge-e\). Restriction
to \(K_w\) raises Frobenius to the \(h\)th power modulo inertia, which
acts trivially here. Thus \(t_s=1+b_1^h+b_2^h\), and integrality at
every finite place follows as above. Equivalently, if the newform is
in motivic weight \(2e\), its eigenvalues are divided by \(p^e\);
this is an integral power and commutes with conjugation. This proves
the assertions in both ranks. For the comparison between spherical
parameters, labelled jumps and crystalline Frobenius see
Barnet-Lamb--Gee--Geraghty--Taylor,
Thm.~2.1.1~\cite{blggtPotential}.
\end{proof}

\begin{lemma}[The integral trace product]\label{rk3:integer}
Suppose the finite list \((t_s)_{s\in S}\) of algebraic numbers is
Galois invariant with multiplicities, each \(|t_s|\le3\), and each
\(q^et_s\) is an algebraic integer. Set \(Q=q^e\) and
\[
 t^0_j=Q^{-1}z_j+1+Qz_j^{-1},\qquad
 C_p=\prod_{j\in J}\prod_{s\in S}Q(t_s-t^0_j).
\]
Then \(C_p\) is a nonzero rational integer with
\[
 \log|C_p|\le |J|\,|S|\,(2e\log q+1).
\]
For any \(\ell\ne p\), fix an embedding of the numbers involved into
\(\overline{\mathbf Q}_\ell\), with \(v_\ell(\ell)=1\). If at every
\(j\) a sublist of indices \(S_j\subseteq S\) has been selected, then
\[
      \sum_{j\in J}\sum_{s\in S_j}v_\ell(t_s-t^0_j)
       \ \le\ v_\ell(C_p).
\]
\end{lemma}
\begin{proof}
In a factor \(Qt_s-z_j-Q-Q^2z_j^{-1}\), all terms are algebraic
integers. The last has absolute value \(Q^2\), the sum of the others
has magnitude at most \(4Q+1<Q^2\) since \(Q\ge7\), and the magnitude
of the factor is at most \(2Q^2\). The action of any Galois automorphism
permutes the full product. A rational algebraic integer is an integer.
Sum \(\log(2Q^2)\le1+2e\log q\) over the factors, with the empty product
interpreted as one. At \(\ell\ne p\), \(Q\) is a unit. Every omitted
factor has nonnegative valuation by integrality. Removing those factors
proves the inequality. The embedding may vary with \(\ell\); the integer
in the bound remains fixed.
\end{proof}

Here are the prime and character conditions at the rank-three
application. The modulus \(N\) is even, \(e\ge1\),
\(\psi(-1)=(-1)^e\), and \(\ell\) is odd with
\(\ell\nmid\operatorname{ord}(\psi)|\operatorname{disc}K|h_K\) and
\(\ell-1\nmid 2e\). The \(\ell\)-local component of \(\psi\) has
square one and \((N,|\operatorname{disc}K|)=1\). If \(\psi^{(\ell)}\)
is obtained by deleting that local component, the quadratic field \(K\)
must not be contained in the compositum of \(\mathbf Q(\mu_\ell)\)
and the field cut out by \(\psi^{(\ell)}\). These conditions have the
same truth value for every faithful character of the cyclic image:
raising its values to a power coprime to the image order preserves
the kernels, the order of every component, and the sign at \(-1\).
\src{src/Development/AutomorphicArithmetic/AutomorphicVarieties/HeightAndVanishing/Counts.lean}{9}{Prime and character hypotheses}.

\begin{lemma}[Independence of the integral coefficient ring]\label{rk3:coeff}
Let \(L_0\) be the extension of \(\mathbf Q_\ell\) generated by the
values of the character on units, and let \(\mathcal O_0\) be its ring of integers.
For a finite extension \(L/L_0\), put \(A_L=(L/\mathcal O_L)
(\psi^{-1}\chi_\ell^{-e})\). The strict groups at level \(\ell N\)
satisfy
\[
 S_{\ell N}(\mathbf Q,A_L)\simeq
 S_{\ell N}(\mathbf Q,A_{L_0})\otimes_{\mathcal O_0}\mathcal O_L .
\]
Finiteness and length divided by the absolute ramification index
are independent of \(L\).
\end{lemma}
\begin{proof}
The character values on units have finite order prime to \(\ell\), so lie in
an unramified extension by Hensel lifting the roots of \(X^m-1\)
over a residue field containing \(\mu_m\), \(m=\operatorname{ord}\psi\).
The remaining Tate character takes values in \(\mathbf Z_\ell^\times\).
The finite free extension of valuation rings is flat, giving
\((L_0/\mathcal O_0)\otimes\mathcal O_L=L/\mathcal O_L\).
In a free \(\mathcal O_0\)-basis of \(\mathcal O_L\), each group of
continuous cochains becomes the finite direct sum of its counterparts
over \(\mathcal O_0\), with coordinatewise coboundaries and restrictions.
The common kernel of the restriction maps is obtained coordinatewise,
even for the conditions over all inertia subgroups. This proves the
displayed isomorphism. Finite length is preserved and reflected since
the extension is finite free and faithfully flat. Each elementary factor
\(\mathcal O_0/(\pi_0^k)\) gives length \(k\,e(L/L_0)\) after extension.
Dividing by \(e(L/\mathbf Q_\ell)=e(L/L_0)e(L_0/\mathbf Q_\ell)\)
proves the length assertion.
\end{proof}

To see precisely what must be supplied by the deformation argument, let
\(\mathcal O\) be a coefficient discrete valuation ring of ramification
index \(a\) over \(\mathbf Z_\ell\) and write \(\lambda_j\) for the
finite length over \(\mathcal O\) of the relevant strict Selmer group,
divided by \(a\). This normalization makes length and \(v_\ell\)
compatible under extension of coefficients. In the rank-three case the
needed assertion gives, for each \(j\), a finite set of pairwise
nonisomorphic admissible representations \(\rho\) with
\[
       \lambda_j\le
          2\sum_\rho v_\ell\bigl(\operatorname{tr}
                       \rho(\mathrm{Frob}_w^{-1})-t^0_j\bigr).
\]
It includes finiteness of the Selmer group. The trace-list theorem must
classify these representations by distinct indices within each such set,
identify their normalized inverse-Frobenius traces with the \(t_s\),
and give \(|S|\le C_K(N(e+1))^{A_K}\) with constants independent of
\(p,\ell,e,N,H\). When one character with kernel \(H\) exists over
\(\mathcal O\), every other embedding of the cyclic image is a coprime
power of that character, so the orbit used in the length sum is exactly
the complex list \(J\). The lemma now gives
\(\sum_j\lambda_j\le2v_\ell(C_p)\). It explains why one cannot replace
the spectral list by an unlabelled set of different trace values.

Choose another allowable prime \(p'\ne p\). For any finite set of
allowable coefficient primes, sum the first estimate away from \(p\),
and use \(p'\) only at \(p\). Prime factorization gives
\[
 \sum_{\ell}\log\ell\sum_j\lambda_j
       \le 2\log|C_p|+2\log|C_{p'}|
       \ll \varphi([(\mathbf Z/N)^\times:H])\,
          |\operatorname{disc}K|^5 N^{18}e^7\log(pp').
\]
This bound for finite subsums also bounds the full nonnegative sum
whenever the stated local hypotheses hold throughout it. For the
classical lists Corollary~\ref{rk3:volume} and Lemma~\ref{rk3:purity}
give the displayed exponents and the integral factors. The implied
constant is absolute. For each fixed \(K\), this gives the
polynomial dependence on \(N\) and \(e+1\) used in the rank-three
window construction, with explicit dependence here on the discriminant.
\src{src/Development/AutomorphicArithmetic/AutomorphicVarieties/UnitaryThreeGroups/KatoUse.lean}{52}{The fixed-field bound used on the dependency path}.
If two semisimple Galois representations are assigned the same full
Hecke system, equality of traces at almost all Frobenius elements,
Chebotarev and Brauer--Nesbitt make them isomorphic. For the reducible
list, assigning the same rank-two system gives isomorphic restrictions
and hence the same rank-three class. This supplies the injection on
each pairwise nonisomorphic sublist. For classical cuspidal attachment
and the stated local Galois conditions, Proposition~\ref{newgal:level}
transfers the principal level, including at the coefficient prime.
The outstanding requirement is a construction from every relevant strict
Selmer group which gives the length inequality and produces such
automorphic systems.
It is this construction which makes the integer relevant to cohomology.

\subsection{Quadratic restriction and the rank-four trace product}

In the rank-four problem the automorphy hypothesis over \(\mathbf Q\)
refers to continuous representations after restriction to imaginary
quadratic fields where the restriction is irreducible. Here is how a
bounded list over those fields can be used without choosing a
cuspidal representation over \(\mathbf Q\).
The following finite-list argument is an independent way to use this
part of the stronger automorphy hypothesis on the active path.
\src{src/Development/AutomorphicArithmetic/AutomorphicVarieties/ModularityArgument/ThresholdArguments/DenseThresholdCases/AutQ0.lean}{20}{Forgetting the additional automorphy witnesses};
\src{src/Development/AutomorphicArithmetic/AutomorphicVarieties/ModularityArgument/ThresholdArguments/DenseThresholdCases/VWLevel.lean}{18}{Use of level, automorphy and polarization in rank four}.

\begin{lemma}[Descent of the finite list]\label{rk4:quadratic}
Fix \(n\ge2\) and \(n^2+3\) distinct imaginary quadratic fields of
pairwise coprime discriminants. At a prime \(p\) discard the fields
ramified at \(p\); in each remaining field \(K_i\) choose a place
above \(p\), of residue degree \(f_i\in\{1,2\}\). Suppose finite lists
\(S_i\), independent of the coefficient prime, are given, with each
index \(s\) carrying an algebraic eigenvalue multiset \(B_{i,s}\)
of size \(n\), stable as a list under Galois conjugation.
For every \(\ell\ne p\), fix the coefficient embedding. Consider a
family of absolutely irreducible \(n\)-dimensional representations of
\(G_{\mathbf Q}\) unramified at \(p\). Whenever \(\ell\) is unramified
in \(K_i\) and the restriction is irreducible, suppose \(S_i\)
classifies the restrictions up to isomorphism and records their
arithmetic Frobenius eigenvalues at the chosen place. Then there is a
Galois-invariant list of at most \(2^{n+1}\sum_i|S_i|\) multiset
entries classifying the original representations with their arithmetic
Frobenius at \(p\). Here classification means an injective labelling
of isomorphism classes; different labels may have the same eigenvalues.
\end{lemma}
\begin{proof}
The prescribed fields can, for example, be obtained from distinct
primes \(q_i\equiv3\bmod4\) by adjoining \(\sqrt{-q_i}\).
For an irreducible representation \(r\), every quadratic self-twist
\(\eta\) gives a nonzero operator in \(\operatorname{End}(r)\) on
which conjugation acts by \(\eta\). Vectors belonging to distinct such
characters are independent. Indeed, apply \(g-\eta(g)\) to a shortest
nonzero relation, choosing \(g\) at which some other character differs
from \(\eta\); a shorter nonzero relation would result. Thus there are
at most \(n^2\) self-twists. For \(H=\ker\eta\) of index two, if
\(r|_H\) is reducible, take a simple \(H\)-submodule \(U\).
For \(g\notin H\), \(U+gU\) is stable under the whole group since
\(g^2\in H\). Irreducibility and properness of \(U\) imply it is
the direct sum \(U\oplus gU\). The operator acting by \(+1,-1\)
on the summands intertwines \(r\) with \(r\otimes\eta\).
At most one of the fields ramifies at \(p\), and at most one at
\(\ell\), so at least one of the remaining restrictions is irreducible.

If two extensions of an irreducible \(H\)-representation agree on \(H\)
in a chosen basis, the quotient of their matrices at \(g\notin H\)
commutes with \(r(H)\). Schur's lemma makes it scalar, and its square
is one since their matrices at \(g^2\) agree. Thus there are at most
two isomorphism classes of extensions. At the chosen place, arithmetic
Frobenius restricts as the \(f_i\)-th power of Frobenius at \(p\)
modulo inertia, up to conjugacy. Include for each \(i,s\) all multisets
\(A\) whose entrywise \(f_i\)-th powers form \(B_{i,s}\).
Ordering \(B_{i,s}\) temporarily shows there are at most
\(f_i^n\le2^n\) possible multisets, including when some entries
coincide. Give each two labels for the possible extensions. This is a
conjugation-stable finite list. Assign each representation to the first
eligible field in a fixed order, its restriction label, its actual
multiset and one of the two extension labels. This assignment is
injective for each fixed coefficient embedding. The whole list is
independent of that embedding and of \(\ell\).
\end{proof}

\begin{corollary}[The integer in rank four]\label{rk4:integer}
Take \(n=4\), \(p\ge5\), \(e\ge2\). Suppose, in the preceding lemma,
that every member of \(B_{i,s}\) is an algebraic integer of absolute
value \(p^{f_i(e+1)/2}\) in every complex embedding. Let \(S\) be
the resulting list and \(t_s\) its entrywise traces. Fix
\(z_1,z_2,z_3\in\mu_k\), \(k\ge1\). There is a nonzero rational
integer \(I_p\) with
\[
 \log|I_p|\le\varphi(k)|S|\bigl((e+1)\log p+\log 2\bigr)
\]
which bounds the sum of the \(v_\ell\)'s of the trace differences for
any pairwise nonisomorphic sublist at \(\ell\ne p\), for the comparison
trace \(p^{e+1}z_1+p^e z_2+pz_3+1\).
\end{corollary}
\begin{proof}
Every lifted eigenvalue is integral by the monic root equation,
and its absolute value is \(p^{(e+1)/2}\). For each
\(j\in(\mathbf Z/k)^\times\) put \(T_j=p^{e+1}z_1^j+p^ez_2^j+pz_3^j+1\);
at \(k=1\) use the unique automorphism of the trivial root-of-unity
group. Set \(P=p^{e+1}\). Since \((p^e+p+1)/P\le1/5+1/25+1/125<1/2\)
and \(P^{1/2}\ge5^{3/2}>8\), we have
\[
 |T_j|>P/2,\qquad |t_s|\le4\sqrt P<P/2,\qquad |t_s-T_j|<2P.
\]
The product \(I_p=\prod_{s,j}(t_s-T_j)\) is a nonzero algebraic
integer fixed by every Galois automorphism: the latter permutes the
three powers \emph{simultaneously} through the indicated indices.
This proves the size assertion. Every difference is integral at
\(\ell\); keeping just the factors belonging to a chosen sublist at
the identity automorphism gives the valuation bound.
\end{proof}
This gives the trace product in the arithmetic-Frobenius convention
of rank four. If the quadratic lists have size polynomial in \(M(e+1)\),
the bound has the required polynomial dependence, since only finitely
many fields were fixed for the rank. To apply it to a particular
Selmer group one needs the classical quadratic attachments of its
deformation points with the stated levels, weights and purity.
Section~\ref{spreal:section} identifies the branch on which an
attachment follows from a known modularity theorem. \src{src/Development/AutomorphicArithmetic/AutomorphicVarieties/ModularityArgument/ThresholdArguments/DenseThresholdCases/AfStmt.lean}{34}{Rank-four comparison required of those deformations}.

\subsection{Which ordinary points give classical spectra?}\label{spreal:section}

Write \(k=\overline{\mathbb Q}_\ell\) and
\(\varepsilon=\varepsilon_\ell\). A representation is called strongly
irreducible if its restriction to every open subgroup is absolutely
irreducible. We work with continuous representations at actual
characteristic-zero points. They and their invariant lines can be defined
over a finite coefficient extension, by Lemma~\ref{newgal:coeff} and
extension of the coefficients of a finite basis.

\begin{lemma}[Finite self-twists at an ordinary place]\label{spreal:twists}
Let \(V\) be an irreducible representation of \(G_F\) of dimension \(n\).
Suppose the semisimplification on inertia at one place \(v\mid\ell\)
is a sum of characters \(\lambda_i\) whose restrictions to an open
subgroup of inertia are \(\varepsilon^{m_i}\) for distinct integers \(m_i\).
Every finite-order self-twist of \(V\) is unramified at \(v\).
If \(\varepsilon(I_v)=\varepsilon(G_F)\), then
\(V|_{G_{F(\mu_{\ell^\infty})}}\) is irreducible.
\end{lemma}
\begin{proof}
If \(V\simeq V\otimes\nu\), then multiplication by \(\nu|_{I_v}\)
permutes the \(\lambda_i\). On an open subgroup killing \(\nu\) the
characters remain distinct, since the cyclotomic image there is
infinite. The permutation is the identity, so \(\nu|_{I_v}=1\).
Put \(H=\ker\varepsilon\). The restriction to \(H\) is semisimple:
the sum of its simple submodules is nonzero and \(G_F\)-stable.
If it were reducible, its commutant would contain a nonzero traceless
endomorphism, obtained by subtracting a scalar from an equivariant
projection. The abelian quotient \(G_F/H\) acts on the space of such
endomorphisms by commuting operators. A common eigenvector \(T\ne0\)
gives an intertwiner with a character twist. Its kernel is
\(G_F\)-stable, so \(T\) is invertible. Schur's lemma and
\(\operatorname{tr}T=0\) exclude the trivial character. Taking
determinants shows that the character has order dividing \(n\).
It is continuous from its action on the line \(kT\). Triviality on
\(I_v\) and surjectivity onto the cyclotomic quotient give a contradiction.
\end{proof}

For the rank-three datum let \(K/\mathbb Q\) be the imaginary quadratic
field, \(\eta\) its character, and \(\Psi=\psi\varepsilon^e\).
Use the prime and character hypotheses in Section~\ref{rk3:sec}.
In particular \(\ell\nmid D_K h_K\operatorname{ord}(\psi)\),
\(\psi_\ell^2=1\), and \(\ell-1\nmid2e\). Thus \(\ell\ge5\).
The normalized matrix at complex conjugation is the identity with
negative sign. On \(G_K\) the residual matrices are diagonal with
characters \(\bar\Psi,1,\bar\Psi^{-1}\). Away from \(\ell\) the
prescribed inertia kernel kills elements with trivial residual action.
At \(\ell\) put \(J=I_\ell\cap\ker\psi_\ell\); here \(I_\ell\subset G_K\).
The ordinary equations of Section~\ref{ct:section} apply on \(J\).
\src{src/Development/AutomorphicArithmetic/AutomorphicVarieties/UnitaryThreeGroups/RhoKApply.lean}{76}{Defining equations for a characteristic-zero point}.

\begin{lemma}[The inertia flag]\label{spreal:flag}
On \(J\) the ordinary equations give an upper triangular flag with
ordered diagonal \((\varepsilon^e,1,\varepsilon^{-e})\). This holds
over \(k\) and over \(k[\delta]/(\delta^2)\) at deformations of such
a point. The flag is stable under every positive-sign element of the
decomposition group. Its graded inertia characters on \(I_\ell\)
are the indicated powers times characters of \(I_\ell/J\).
\end{lemma}
\begin{proof}
Choose \(s\in J\) with infinite cyclotomic value. The prescribed
characteristic polynomial has three distinct constant roots at \(s\).
Its Lagrange projectors split the free module into three free lines,
also over the dual numbers. In the ordered product of three shifted
matrices, put \(g\in J\) in position \(a\) and \(s\) in the other
two. In the eigenbasis the \((i,j)\) entry for \(i\ge a\ge j\)
is the corresponding entry of the \(a\)-th shifted matrix times
a unit. Taking \(a=j\) annihilates all lower entries and imposes
the stated diagonal. The common \(\varepsilon^e\)-eigenline is
uniquely specified by its eigenvalue at \(s\). The same description
in the quotient gives the next line. As \(J\) is normal and these
characters are conjugation invariant, the flag is preserved.
The extra inertia factors on the graded lines are trivial on \(J\).
\end{proof}

\begin{proposition}[The endoscopic block]\label{spreal:end}
Suppose the normalized rank-three point is of type II: on \(G_K\) it
preserves the middle coordinate and the outer coordinate plane,
and its action on the latter is irreducible. There is an odd
representation \(W\) of \(G_{\mathbb Q}\), with determinant \(\eta\),
such that the point restricts to \(1\oplus W|_{G_K}\). This \(W\)
is the normalized representation of a classical cuspidal newform
of weight \(2e+1\). Its automorphic representation has invariants
at principal level \(N|D_K|\).
\end{proposition}
\begin{proof}
On \(G_K\) denote by \(m\) the middle character and by \(d\) the
determinant of the whole representation. Both reduce to one. On \(J\)
the determinant is one by Lemma~\ref{spreal:flag}, and \(m\) is one
of its three constituents. The reduction of \(\varepsilon^e\) is
nontrivial on \(J\): otherwise the image of
\(\bar\varepsilon^e(I_\ell)\) has order at most two, contrary to
\(\ell-1\nmid2e\). Here \(\bar\varepsilon(I_\ell)=\mathbb F_\ell^\times\).
It follows that \(m|_J=1\). Since \([I_\ell:J]\le2\), reduction
detects both characters on \(I_\ell\). The argument holds at every
place of \(K\) over \(\ell\): conjugating from the marked place
either conjugates the matrices or takes their conjugate inverse
transposes, reversing and negating the three exponents together.
At primes away from \(\ell\), the inertia action on \(G_K\) factors
through its residual image by the kernel condition. This image has
order dividing \(\operatorname{ord}(\psi)\), prime to \(\ell\).
Thus \(m,d\) are unramified globally, and their images in the group
of principal units factor through \(\operatorname{Cl}(K)\). Since
\(\ell\nmid h_K\), both are trivial.

Let \(B(h)\in\mathrm{SL}_2(k)\) be the outer block for \(h\in G_K\).
Take \(W(h)=B(h)\) and \(W(c)=D=\operatorname{diag}(-1,1)\) at
complex conjugation. With \(J_2\) the antidiagonal unit matrix,
\(J_2B^{-t}J_2=DBD^{-1}\) whenever \(\det B=1\).
This is the action of \(c\) on the corner, and \(c^2=1\), giving the
claimed extension, determinant and oddness. It is continuous on the
two cosets and irreducible already on \(G_K\).

On \(J\) the highest inertia line of Lemma~\ref{spreal:flag} belongs
to the outer plane, since the middle character is one. In \(W\) it
is the unique common \(\varepsilon^e\)-eigenline, so normality of
\(J\) makes it stable under the full decomposition group. The two
graded characters on inertia have powers \(e,-e\) with finite parts
trivial on \(J\). Over the finite local extension killing the finite
parts the two characters are crystalline. Their ratio \(U\), upper
line divided by quotient, is an unramified twist of \(\varepsilon^{2e}\).
Over that local field \(L\) both \(H^0(L,U)\) and \(H^0(L,U^*(1))\)
vanish since \(2e\ge2\). The local Euler characteristic formula gives
\(\dim H^1(L,U)=[L:\mathbb Q_\ell]\) over the coefficient field.
The Bloch--Kato dimension formula gives the same dimension for
\(H^1_f(L,U)\), since \(\operatorname{Fil}^0 D_{\rm dR}(U)=0\).
The extension is therefore crystalline over \(L\) and de Rham over
\(\mathbb Q_\ell\), with distinct jumps \(\{-e,e\}\).

The representation is unramified outside \(\ell N D_K\).
Pan's regular Fontaine--Mazur theorem now gives a cuspidal eigenform
up to twist. Its weight is one plus the difference of the two jumps.
In the geometric Frobenius convention of Section~\ref{rk3:sec}, the
cohomological representation is \(W\varepsilon^{-e}\), with jumps
\(0,2e\). A finite character twist is absorbed into the newform.
Pan's only extra regular-case restriction for odd primes concerns
\(\ell=3\), already excluded here.
See \cite[Theorem 1.0.4]{panFM}.

For \(v\nmid\ell\), \(I_v\cap G_{K(\mu_N)}\) acts trivially by the
residual kernel condition. Thus inertia acts through a subgroup of
the finite abelian group \(\operatorname{Gal}(K(\mu_N)/\mathbb Q)\),
and its characters extend to that group. At \(\ell\), the finite
inertia characters are trivial on \(J\), hence are either trivial
or the prescribed quadratic local component, which likewise extends
with modulus dividing \(N\). Since \(K\subset\mathbb Q(\mu_{|D_K|})\),
Proposition~\ref{newgal:level} gives the claimed level. In particular
the classical trace and its nonvanishing comparison in
Lemmas~\ref{rk3:purity} and~\ref{rk3:integer} apply to every such block.
\end{proof}

\begin{proposition}[Strong irreducibility at the rank-three points]\label{spreal:strong}
The representation on \(G_K\) at every type I point and the outer
two-dimensional constituent at every type II point are strongly
irreducible. Both remain irreducible on \(G_{K(\mu_{\ell^\infty})}\).
\end{proposition}
\begin{proof}
Call the relevant irreducible representation \(V\) and put
\(d=\dim V\in\{2,3\}\). The inertia characters of
Lemma~\ref{spreal:twists} occur by Lemma~\ref{spreal:flag} and
Proposition~\ref{spreal:end}. As \(K\) is unramified at \(\ell\),
its local inertia already surjects onto \(\mathbb Z_\ell^\times\),
by total ramification of the local cyclotomic tower. This proves
the second assertion.

Suppose restriction to an open subgroup is reducible, and replace
that subgroup by its normal core \(H\). As in Lemma~\ref{spreal:twists},
the restriction to \(H\) is semisimple. Conjugation permutes its
isotypic spaces transitively, so their number times their common
multiplicity times the dimension of a simple constituent is \(d\).
Primality of \(d\) leaves two possibilities. A scalar restriction
to \(H\) contradicts the distinct characters on \(H\cap J\).
Otherwise there are \(d\) distinct lines permuted transitively.
Choose a stable lattice and replace it by the direct sum of its
intersections with these lines, enlarging coefficients if necessary.
In generators of these rank-one lattices the matrices are monomial
with unit coefficients. Their reduction has order prime to \(\ell\):
its diagonal subgroup lies in a product of finite residue-field
multiplicative groups and its permutation quotient has order
dividing \(d!\). Thus the reduction is semisimple. Independence of
residual semisimplification from the lattice now identifies it with
a sum of the given global characters. Its image is abelian, and
its permutation quotient, unchanged by reduction of the units, is
transitive abelian of prime degree. It is cyclic of order \(d\).

Let \(\nu\) be a faithful character of that quotient. The operator
which multiplies the successive lines in cyclic order by successive
powers of a primitive \(d\)-th root intertwines \(V\) with the
corresponding twist (up to replacing \(\nu\) by its inverse).
Lemma~\ref{spreal:twists} gives \(\nu|_{I_\ell}=1\). Its reduction
has order \(d\), and twisting must permute the residual constituents.
In dimension three those constituents are
\(1,\bar\Psi,\bar\Psi^{-1}\), so the orbit of \(1\) consists of all
three powers of \(\bar\nu\). In dimension two the two constituents
are \(\bar\Psi,\bar\Psi^{-1}\), so \(\bar\nu=\bar\Psi^2\).
In either case \(\bar\Psi^2|_{I_\ell}=1\), contrary to
\(\bar\Psi^2|_{I_\ell}=\bar\varepsilon^{2e}|_{I_\ell}\ne1\).
\end{proof}

\begin{corollary}[The polarized adjoint at an endoscopic point]\label{spreal:endrig}
For the characteristic-zero coefficient space \(A=M_3(k)\) with
the polarized adjoint action of a type II point,
\(H^1_f(\mathbb Q,A)=0\).
\end{corollary}
\begin{proof}
Decompose matrices into the outer endomorphism block, the middle
scalar, and the two off-diagonal rectangles taken together. All
three spaces are stable, since \(c\) acts by negative reversed
transpose. The outer block is the polarized adjoint of \(W|_{G_K}\).
Indeed its polarization matrix for this action is \(J_2\), the
product of the alternating determinant pairing
\(\left(\begin{smallmatrix}0&1\\-1&0\end{smallmatrix}\right)\)
and \(W(c)^{-1}\). Cyclic base change of the modular form to \(K\)
is cuspidal: a noncuspidal quadratic base change is induced from a
character of \(K\), which would give reducible restriction to
\(G_K\). The normalized Tate twist does not affect the adjoint.
Propositions~\ref{spreal:strong} and~\ref{newgal:adselmer} give its
Selmer vanishing. For the other summands restriction to \(K\) is
injective on rational cohomology by restriction and corestriction
and takes finite classes to finite classes. Over \(K\) the scalar
is trivial. A finite local class in \(H^1(K_v,k)\) at \(v\mid\ell\)
is unramified: the unramified line is crystalline and exhausts the
finite subspace by its dimension. A global everywhere unramified
additive character of \(G_K\) vanishes by class field theory.

The off-diagonal part over \(K\) is \(W|_{G_K}\oplus W^*|_{G_K}\).
These two representations are isomorphic since \(\det W|_{G_K}=1\).
For the weight \(k_0=2e+1\) newform just obtained, denote its
cohomological representation by \(V_f\). In Kato's convention
\(W=V_f(e)\). Here \(1\le e\le k_0-1\) and \(2e\ne k_0\).
Kato's Theorem~14.2(1) over the finite abelian extension \(K/\mathbb Q\)
gives finiteness of the discrete Bloch--Kato group for any stable
lattice in this twist. To pass to rational coefficients, the image of
the vector-space Selmer group under reduction in the discrete group
is divisible and hence zero. Its kernel is contained in the image
of integral \(H^1(G_{K,S},T)\), a finitely generated lattice module
for a finite \(S\) containing the ramification and \(\ell\).
A vector space contained in such an image is zero. This argument
may be made over the original finite coefficient field and then
extended to \(k\) as in Proposition~\ref{newgal:adselmer}.
See \cite[\S14.1 and Theorem 14.2(1)]{katoZeta}.
\end{proof}

The fixed ordinary equations place every continuous tangent cocycle
at this point in \(H^1_f\).
At a prime away from \(\ell\) its restriction to inertia factors
through a finite group, both for the point and the lift, by the
fixed kernel condition. Finite-group cohomology over \(k\) vanishes
in positive degree. At \(\ell\) pass to a finite local extension
inside the positive-sign subgroup killing the finite graded inertia
characters. By Lemma~\ref{spreal:flag} the lift over the dual numbers
has a stable ordered flag. Each rank-one graded module, regarded as
a two-dimensional representation over the coefficient field, is an
unramified representation tensored with its fixed cyclotomic power
and is crystalline. Successive extensions, taken in the order of the
flag, are crystalline as well. Indeed the Hom from a later graded
piece to the already constructed part has a filtration with factors
unramified characters times \(\varepsilon^s\) for \(s\ge e\ge1\).
It is crystalline and its de Rham filtration in degree zero vanishes.
When \(e\ge2\), distinct exponents alone imply that its dual twist
has no invariants. For \(e=1\) we use the modular block at the
coefficient prime. Its normalized Weil--Deligne parameter has zero
monodromy by potential crystallinity in Proposition~\ref{spreal:end}.
The full comparison at \(\ell\) in Proposition~\ref{newgal:level}
identifies it in the weight-zero normalization with the local
parameter of the unitary newform. This parameter is tempered by
Blasius's Theorem~1 for holomorphic forms of weight at least two.
With zero monodromy its Weil Frobenius eigenvalues have magnitude
one, also after finite local restriction. The middle character is one.
Every ratio of two graded characters after setting \(\delta=0\)
thus has Weil Frobenius eigenvalue of magnitude one, whereas
\(\varepsilon\) gives \(q_L^{-1}\) on geometric Frobenius.
No such ratio is \(\varepsilon\), so the dual twist has no invariants
in this case either.
See \cite[Theorem 1]{blasiusRam}.
The same local dimension
calculation as in Proposition~\ref{spreal:end} gives \(H_f^1=H^1\).
The lifted representation and its reduction being crystalline there,
their extension class becomes a coboundary over \(B_{\rm cris}\).
Restriction on cohomology with these period coefficients is injective
for finite local extensions by corestriction. The global tangent
class is therefore finite. Corollary~\ref{spreal:endrig} makes it a
coboundary, giving infinitesimal rigidity modulo conjugation at
these endoscopic points.

\begin{proposition}[Finite tangent classes on the rank-four component]\label{spreal:four}
For the good rank-four datum, at every actual characteristic-zero
point every tangent arising from the specified component gives a
finite local extension class in the adjoint at every finite prime.
\end{proposition}
\begin{proof}
Put \((\theta_i)=(\alpha\Psi,\Psi,\alpha,1)\). The ordinary
equations in Proposition~\ref{dp:local} give a flag with graded inertia
characters \(\theta_i\), also over the dual numbers. The proof with
projectors there applies at any point: the distinguished inertia
matrix has the same four distinct eigenvalues \(d_i=\theta_i(s_0)\).
After killing the finite inertia factors the four powers of
\(\varepsilon\) are \(e+1,e,1,0\). Denote the characteristic-zero
graded characters on the full decomposition group by \(\kappa_i\).

Let \(\sigma=s_\ell\) be the specified Frobenius lift with
\(\varepsilon(\sigma)=1\). We claim
\(\overline{\kappa_i(\sigma)}=\overline{\theta_i(\sigma)}\)
with the indicated labels, even if some residual eigenvalues coincide.
The component hypothesis provides a complete local domain \(B\)
with an ordinary matrix family \(R\), specializing to the diagonal
point and to the given point. Invert the nonzero constant
\(a=\prod_{i<j}(d_i-d_j)\). On \(B[1/a]\) the projectors
\(P_i\) of \(R(s_0)\) have rank one. Indeed their ranks are locally
constant on this connected spectrum and are one at the diagonal
specialization. The triangular equations make
\(b_i=\operatorname{tr}(P_iR(\sigma))\) the scalar on the
corresponding flag quotient, so that \(b_i\) is a root of the
characteristic polynomial of \(R(\sigma)\).
Over the Henselian ring \(B\) that polynomial splits into monic
factors \(F_u\), one for each distinct residual eigenvalue \(u\),
with \(\bar F_u=(X-u)^{n_u}\). In the domain \(B[1/a]\)
exactly one \(F_u(b_i)\) vanishes, by coprimality. The diagonal
specialization shows that its label is
\(u=\overline{\theta_i(\sigma)}\). At the other specialization
the same factor vanishes, proving the claim. In this argument any
\(O\)-algebra specialization \(B\to\overline{\mathbb Q}_\ell\)
is integral and reduces to the given residual map. Indeed an
element of negative valuation has a polynomial relation over
\(O\) with constant term one and all higher coefficients in
\(\mathfrak m_O\), by using the minimal polynomial of its inverse
over \(E\). Such a relation gives a unit in \(B\). For an element
of \(\mathfrak m_B\) mapping to a unit, its image has monic
minimal polynomial with unit constant term, again impossible.

Let \(I_0=I_\ell\cap\ker\psi\cap\ker\omega\). In the quotient of
the local group by \(I_0\) take the closed subgroup generated by
\(\sigma\) and call its preimage \(G_L\). The generated procyclic
subgroup maps isomorphically to \(\widehat{\mathbb Z}\), since its
generator maps to a Frobenius generator there. Thus \(L\) is finite,
\(G_L\cap I_\ell=I_0\), and \(\sigma\in G_L\).
For any higher and lower graded characters, their ratio differs from
\(\varepsilon\) on this group. If their exponents do not differ
by one this follows on \(I_0\). The ratios with exponent difference
one are those with underlying diagonal ratio \(\alpha\), or
\(\Psi/\alpha\) if \(e=2\). Since \(\bar\alpha(\sigma)=
\bar\xi(\sigma)\), the good-datum condition
\(\bar\xi(\sigma)\bar\Psi(\sigma)^i\ne1\) for \(i=0,-1\)
proves the assertion by the labeled congruence.

The dual-number flag reduces modulo \(\delta\) to this same flag,
by the distinct labeled eigenvalues at \(s_0\). On \(G_L\) its
graded pieces are unramified two-dimensional representations times
their cyclotomic powers. For successive extensions from a later
piece to the already constructed part, the Hom has strictly positive
cyclotomic exponents and its dual twist has no invariants by the
preceding test. As in the rank-three calculation the local Euler
and Bloch--Kato formulas give \(H^1=H^1_f\), inductively giving
crystallinity over \(L\) and finiteness of the extension class over
\(\mathbb Q_\ell\). At every other prime the point and its tangent
kill a common finite-index subgroup of inertia by the prescribed
kernel conditions. Averaging over the finite quotient gives the
unramified cohomology condition.
\src{src/Development/AutomorphicArithmetic/AutomorphicVarieties/SymplecticGSp4/DecompSigmaOrd.lean}{27}{The domain and the two specializations},
\src{src/Development/AutomorphicArithmetic/AutomorphicVarieties/PadicHodgeTheory/SplAdjUnit.lean}{264}{the labeled Frobenius congruence}
and \src{src/Development/AutomorphicArithmetic/AutomorphicVarieties/SymplecticGSp4/TbTriDiag.lean}{310}{the local nonresonance test}.
\end{proof}

For each irreducible point choose an imaginary quadratic \(K\)
unramified at \(\ell\) with irreducible restriction as in
Lemma~\ref{rk4:quadratic}. Its local inertia has full cyclotomic
image. Thus Lemma~\ref{spreal:twists} verifies the cyclotomic
irreducibility needed by Proposition~\ref{newgal:adselmer}.
A regular polarized cuspidal attachment over this \(K\) therefore
suffices for rank-four rigidity by Proposition~\ref{spreal:four}
and the symplectic projection of
Section~\ref{newgal:section}.
\src{src/Development/AutomorphicArithmetic/AutomorphicVarieties/SymplecticGSp4/MapRepApply.lean}{99}{Polynomial projectors in the ordinary local deformation condition}.

\begin{lemma}[The excluded rank-four block]\label{spreal:even}
Suppose a rank-four point preserves the two coordinate planes of
the block condition. If its representation \(V\) on coordinates
\(1,3\) (one-based indices) is irreducible, it is strongly
irreducible, potentially crystalline with distinct weights at
\(\ell\), and even. Its residual semisimplification is
\(\bar\alpha\bar\Psi\oplus\bar\alpha\).
\end{lemma}
\begin{proof}
The integral specialization gives the asserted residual characters.
At complex conjugation their common value is \(-1\), since
\(\Psi(c)=1\), \(\xi(c)=-1\) and \(\bar\omega=\bar\varepsilon\).
An involution reducing to \(-I\) is itself \(-I\) in odd residual
characteristic, by the idempotent \((1+c)/2\). The block inherits
potential crystallinity from Proposition~\ref{spreal:four}'s
local argument. Intersecting its plane with the full flag gives
two distinct exponents on open inertia, since all four are distinct.
The image of inertia is finite away from \(\ell\).

If \(V\) becomes reducible on an open subgroup, pass to its normal core.
The scalar case is excluded by the exponents. The other possibility,
by the isotypic decomposition in Proposition~\ref{spreal:strong},
is induction of a character \(\lambda\) from a quadratic field \(L\).
This field is real since \(c\) fixes both lines. If \(\ell\) is
nonsplit in \(L\), its decomposition group exchanges the two lines.
Their characters on a common open inertia subgroup then have the
same cyclotomic exponent, as conjugation preserves the cyclotomic
character. If \(\ell\) splits, write \(m_1,m_2\) for the exponents
of \(\lambda\) at the two places. Take a nontrivial totally positive
unit \(u\) of norm one in \(L\). Replacing it by a power kills its
image under the finite inertia characters at every other ramified
place, and puts it in the open unit subgroups with the specified
cyclotomic powers at the two places over \(\ell\). Global reciprocity
on the principal idele of \(u\) gives
\(\iota_1(u)^{m_1}\iota_2(u)^{m_2}=1\), up to simultaneous
inversion according to the reciprocity convention. As
\(\iota_1(u)\iota_2(u)=1\) and \(u\) has infinite order, this gives
\(m_1=m_2\), again impossible. This proves strong irreducibility.
\end{proof}

The rank-four block exclusion would therefore follow from the even
two-dimensional Fontaine--Mazur assertion for these ordinary, residually
split representations. The oddness in Pan's theorem fails precisely
on this block. Calegari's ordinary theorem assumes
\(\ell>7\) and residual image containing \(\mathrm{SL}_2(\mathbb F_\ell)\).
The displayed residual semisimplification rules out this image for
every choice of lattice. See \cite[Theorem 1.2]{calegariEven}.

The missing higher-rank assertion is cuspidal attachment for every
relevant characteristic-zero irreducible point. For example, the
residually reducible lifting theorem of Allen, Newton and Thorne
requires in its Theorem~1.1(4) that the distinct residual constituents
be individually conjugate self-dual with the same multiplier.
After a common twist the actual rank-three constituents have ratios
\(\bar\Psi,\bar\Psi^{-1}\) to the middle character, both pulled back
from \(\mathbb Q\). Individual polarization would make each ratio
self-inverse, thereby identifying the two outer characters, contrary
to the distinctness hypothesis. This obstruction persists under
restriction to a CM extension and any common character twist:
either the square remains nontrivial or the outer characters coincide.
In rank four the residual characters up to twist are
\(1,a,b,ab\), with \(a=\bar\alpha\) and \(b=\bar\Psi\), again
restricted from \(\mathbb Q\). If condition (4) holds, \(a,b\)
are independent quadratic characters. For \(H=\ker a\) the sum
of these four characters is \(\operatorname{Ind}_{H}^{G_F}(1\oplus b|_H)\),
contradicting the residual primitivity in (7). This is also invariant
under common twisting and applies after any CM restriction.
Conditions (5)--(6) of the same theorem independently require a place
away from \(\ell\) with consecutive unramified cyclotomic twists in
the characteristic-zero local semisimplification, and an ordinary
congruent cuspidal seed which is a Steinberg twist there.
See \cite[Theorem 1.1]{antLifting}.

The precise rank-three statement used by the source gives a conjugate
self-dual cuspidal attachment of \(r|_{G_K}\varepsilon^{-1}\) for
every ordinary point over a finite coefficient DVR whose restriction
to every open subgroup of \(G_K\) is irreducible. Proposition
\ref{spreal:strong} supplies this last hypothesis. The proposed passage
uses an auxiliary CM frame, arbitrarily large weights in a prescribed
residue class supporting adjoints of rank-two seeds, propagation in the
ordinary determinant space to the specified point, and descent to
\(K\). To reach the given point one must show its restriction lies
on the automorphic quotient and then descend the attachment. Existence
of a seed at each high weight supplies only points on the seed locus.
\src{src/Development/AutomorphicArithmetic/AutomorphicVarieties/ModularityArgument/ThresholdArguments/DenseThresholdCases/AutEnd3AltAlt.lean}{20}{The ordinary-point assertion},
\src{src/Development/AutomorphicArithmetic/AutomorphicVarieties/ModularityArgument/ThresholdArguments/DenseThresholdCases/SymUse.lean}{27}{the cofinal-weight seed assertion}
and \src{src/Development/AutomorphicArithmetic/AutomorphicVarieties/ModularityArgument/ThresholdArguments/DenseThresholdCases/GdEndOf.lean}{22}{the linked-point and descent inputs}.

For rank four the asserted output gives, for every imaginary quadratic
\(K\) on which the restriction stays irreducible, the cuspidal
representation and the prescribed comparison witness. On a sufficiently
deep auxiliary tower the proposed proof constructs an ideal \(J\)
in a universal determinant algebra. Specializations killing \(J\)
at the prescribed weights must be pure and, in the irreducible
case, automorphic after descent. To use these implications at the
given point one also needs a specialization killing \(J\) which
realizes its restriction, up to conjugacy, with the labeled inertia
characters. In the case with four seed characters the chain assertion
upgrades membership of some component over a residual specialization
to membership of all components after the stated dimension bound
\(30(C+1)^2\le d\). Establishing the existence of \(J\) with those
arithmetic properties and the coverage of the specific point remains
necessary to convert the finite orbit estimate into a comparison for
the entire Selmer module.
\src{src/Development/AutomorphicArithmetic/AutomorphicVarieties/ModularityArgument/AutomorphyLifting/IsAutCsdA.lean}{16}{The quantifier over quadratic restrictions},
\src{src/Development/AutomorphicArithmetic/AutomorphicVarieties/ModularityArgument/ThresholdArguments/DenseThresholdCases/PotentialModularityTower.lean}{63}{the ideal and its required specialization},
\src{src/Development/AutomorphicArithmetic/AutomorphicVarieties/ModularityArgument/ThresholdArguments/DenseThresholdCases/PromodularityCaseThree.lean}{26}{the component propagation}
and \src{src/Development/AutomorphicArithmetic/AutomorphicVarieties/ModularityArgument/ThresholdArguments/DenseThresholdCases/ModularityRankFourOfPromodularity.lean}{20}{the implication to finiteness, block exclusion, and automorphy}.

\subsection{Orthogonal ternary points and quadratic induction in degree four}
\label{hr:section}

Work with the marked representations in Propositions~\ref{spreal:strong},
\ref{spreal:four} and~\ref{dp:curvebranches}. The coefficient field \(E\)
below is a finite extension of \(\mathbb Q_\ell\) for an integral
arithmetic point, or the fraction field of a normalized test curve in
characteristic \(\ell\). Its residue field is finite. We allow finite
extensions of \(E\). Statements about classical attachment use
geometric Frobenius and the normalization
\(\operatorname{rec}(\pi_v|\det|^{(1-n)/2})\) in degree \(n\).

\begin{proposition}[The odd orthogonal ternary branch]\label{hr:orth}
Suppose that \(\rho\) on \(G_K\) is absolutely irreducible and is
isomorphic to its own dual. Let \(S\) be its invariant bilinear form
and \(J_3\) the antidiagonal unit matrix in the marked basis. The
projective class of \(J_3S\) has a unique representative \(C\) in
\(\mathrm{SO}(S)\). Suppose \(C\ne1\). In equal characteristic
such a point cannot occur. In characteristic zero there is a
regular algebraic, conjugate self-dual cuspidal representation
\(\Pi\) of \(\mathrm{GL}_3(\mathbb A_K)\) whose associated
representation is \(\rho\varepsilon^{-1}\).
\end{proposition}
\begin{proof}
The form is unique up to scalar, nondegenerate and symmetric by
irreducibility and odd dimension. The determinant of \(\rho\)
squares to one and reduces to one, so is one. The marked polarization
is \(\rho(chc^{-1})=J_3\rho(h)^{-t}J_3\). Consequently \(D=J_3S\)
implements conjugation by \(c\), \(D^2\) is scalar and
\(D^t S D=\mu S\). Since \((\det D)^2=\mu^3\), the required
representative is \(C=(\mu/\det D)D\), which is independent of the
scale of \(S\). Its square is a scalar in \(\mathrm{SO}_3\), hence
is one. This extends \(\rho\) continuously to \(G_{\mathbb Q}\)
in \(\mathrm{SO}(S)\) by sending \(c\) to \(C\).

Split \(S\) over a finite extension. Identify its special orthogonal
group with \(\mathrm{PGL}_2\) on trace-zero matrices, with quadratic
form their determinant. Indeed its isotropic conic consists of the
nilpotent lines indexed by their kernels in \(\mathbb P^1\).
Conjugation by \(\mathrm{PGL}_2\) induces every automorphism of this
conic, and a linear isometry fixing it pointwise is projectively
scalar, hence is the identity in \(\mathrm{SO}_3\). Write \(P\) for
the resulting projective representation.

We construct a linear lift over a DVR. After adjoining square roots
of representatives of \(E^\times/(E^\times)^2\), the compact image
of \(P\) lifts setwise to \(\mathrm{SL}_2\). The square-class
group is finite by Hensel's lemma and odd residue characteristic.
The inverse image of the compact group under this finite morphism
is a compact group: local sections are obtained by fixing one
nonzero entry of a matrix representative and taking a continuous
square root of its determinant near a given lift. A compact group
of matrices preserves a lattice, by taking the integral span of
the translates of one lattice. In particular we can reduce \(P\)
projectively on this lattice.

The normal subgroup \(H=\ker(\bar\Psi|_{G_K})\) has pro-\(\ell\)
image in the original ternary matrices, since all its matrices
reduce to the identity. Its projective reduction is a normal
\(\ell\)-subgroup. If nontrivial, it fixes a unique line over the
algebraic closure of the residue field (use the upper unitriangular
Sylow subgroup of \(\mathrm{PGL}_2\) over a finite field).
If trivial, the residual projective image is abelian, since
\(G_{\mathbb Q}/H\) embeds in the product of the images of \(\eta\)
and \(\bar\Psi\). Some inertia element has residual eigenvalue
ratio of order \(>2\): its adjoint characteristic polynomial on
the original lattice has roots \(1,z,z^{-1}\) with
\(z=\bar\Psi(s)\), and \(\bar\Psi^2|_{I_\ell}\ne1\).
Characteristic polynomials of the adjoint are independent of
lattice. The centralizer of this projective residual element is a
torus. In both cases the projective reduction is triangular after
a finite residue extension.

It can be made diagonal. In an upper triangular residual basis,
extend ramification by adjoining a square root \(u\) of the
uniformizer and conjugate by \(\operatorname{diag}(u,1)\).
The lower left entries were divisible by \(u^2\). Write
\(\bar\lambda\) for the residual diagonal ratio. The nontrivial
involution \(P(c)\) cannot be in the pro-\(\ell\) congruence
kernel, so \(\bar\lambda(c)=-1\). In equal characteristic take the
constant character \(\lambda=\bar\lambda\). In characteristic zero
take
\[
 \lambda=\widetilde{\bar\lambda}
                 (\varepsilon/\widetilde{\bar\varepsilon})^e,
\]
where the tilde denotes Teichm\"uller lift. For each \(g\) there is
exactly one representative \(W(g)\) with determinant \(\lambda(g)\)
and reduction \(\operatorname{diag}(\bar\lambda(g),1)\). Indeed
first arrange the specified reduction by scalar multiplication,
then correct the determinant by the unique square root in the
principal units. Uniqueness gives multiplicativity, and the same
construction on coordinate neighborhoods gives continuity. The lift
is odd and unramified outside finitely many places: where the
projectivization, determinant and residual characters are
unramified, its inertia is scalar, squares to one and reduces to
one. We have \(\operatorname{Ad}^0W|_{G_K}\simeq\rho\).

In equal characteristic the flag of Proposition~\ref{dp:curvebranches}
triangularizes inertia on \(\rho\). The inertia subgroup with all
diagonal entries one consists precisely of its unipotent elements.
It is therefore normal under the whole decomposition group, acting
through the extended ternary representation. Its image under \(P\)
consists of projective unipotents. If nontrivial they fix a unique
common line. Indeed two with distinct fixed lines have lifts
\(\left(\begin{smallmatrix}1&b\\0&1\end{smallmatrix}\right)\) and
\(\left(\begin{smallmatrix}1&0\\d&1\end{smallmatrix}\right)\)
with \(bd\ne0\) in the basis of those lines. The product of the
first and the \(j\)-th power of the second has trace \(2+jbd\),
which for some \(j\) differs from both \(2\) and \(-2\).
Normality then gives a decomposition-group line. If the subgroup
has trivial image, inertia is finite abelian by the finite
diagonal characters on the flag and faithfulness of the adjoint
action on \(\mathrm{PGL}_2\). It preserves the two eigenspaces of
the inertia element with residual ratio of order \(>2\). They
reduce to distinct lines in the diagonal residual basis, so the
decomposition group, which normalizes inertia and preserves each
residual line, cannot interchange them. The inertia eigenvalue
ratios have finite order at every place, as do the determinants,
so both inertia eigencharacters have finite image. Lemma~\ref{dp:pan}
now makes \(W\) geometrically reducible; its adjoint on \(G_K\)
would preserve the nilpotent line associated to an invariant line.

In characteristic zero the eigenline of the adjoint with inertia
character \(\varepsilon^e\) on \(I_\ell\cap\ker\psi_\ell\) is
unique by Lemma~\ref{spreal:flag}. It is preserved by the full
decomposition group since the character is conjugation invariant.
Its elements are nilpotent: their determinant is fixed by
conjugation and is multiplied by \(\varepsilon^{2e}\) under the
specified scaling. The kernel of any nonzero element therefore
gives a decomposition-group line of \(W\). The ratio of the
subcharacter to the quotient character on this open inertia
subgroup is \(\varepsilon^e\). After further finite restriction
the determinant there is \(\varepsilon^e\), so the ordered
characters there are \((\varepsilon^e,1)\). Both local characters
are de Rham, the upper one of strictly smaller Hodge--Tate weight
in the convention \(\mathrm{HT}(\varepsilon)=-1\). The residual
ratio is nontrivial on inertia. The distinguished ordinary
modularity theorem of Pan applies over \(\mathbb Q\) with the
specified odd determinant, to the maximal ideal defined by \(W\)
in the pseudocharacter ring with \(\ell\) inverted. It is maximal
because the integral image is a submodule of the finite coefficient
DVR over \(\mathbb Z_\ell\), and its localization is a field.
Irreducibility follows from that of the adjoint on \(G_K\).
If Pan writes \(W=W_f\otimes\theta\), its upper character and
the character of the same stable line in the de Rham representation
\(W_f\) are de Rham, so their ratio \(\theta_\ell\) is de Rham.
It is a cyclotomic integer power on open inertia. Away from
\(\ell\), the determinant identity shows that \(\theta^2\), hence
\(\theta\), has finite inertia. Dividing by the cyclotomic power
leaves a finite-order character by global class field theory.
Thus \(W\) is classical with its algebraic normalization.
\cite[Theorem~5.1.1(2)]{panFM}

Let \(\pi\) be its cuspidal representation of \(\mathrm{GL}_2\).
Its quadratic base change to \(K\) is cuspidal, since the
alternative gives a quadratic self-twist, hence reducibility of
\(W|_{G_K}\) by the two eigenspaces of the twist intertwiner.
The base change itself has no nontrivial self-twist: one would be
quadratic by the central character and split \(W\) on its open
kernel, contrary to strong irreducibility of the adjoint
(Proposition~\ref{spreal:strong}). These implications use equality
of unramified Frobenius polynomials and Chebotarev. Apply
Gelbart--Jacquet to the unitary normalization of this base change;
the adjoint parameter is unaffected by any common scalar twist.
The resulting cuspidal
adjoint lift is self-dual and invariant under
\(\operatorname{Gal}(K/\mathbb Q)\). Its archimedean parameter has
three distinct exponents \(e,0,-e\). At unramified places its
Satake parameters are the eigenvalues of \(\operatorname{Ad}^0W\):
the common scalar in the degree-two normalization cancels. The
degree-three arithmetic normalization multiplies them by
\(|\cdot|^{-1}\), giving \(\rho\varepsilon^{-1}\) as claimed.
\cite[Chapter~III, Theorems~4.2 and~5.1]{arthurClozelBC}
\cite[Theorem~9.3]{gelbartJacquetAdj}
\end{proof}

\begin{corollary}[Essential self-duality]\label{hr:ess}
For a strongly irreducible ternary point, an isomorphism
\(\rho^\vee\simeq\rho\otimes\gamma\) on \(G_K\) already implies
\(\gamma=1\). If its orthogonal extension has \(C=1\), it too has
a linear adjoint lift over \(\mathbb Q\), this time with scalar
complex conjugation, and with the same ordinary local description
as in Proposition~\ref{hr:orth}.
\end{corollary}
\begin{proof}
In characteristic zero the determinant of \(\rho\) is one by the
determinant calculation in Proposition~\ref{spreal:end}, which
uses the local flag, the residual inertia kernels and
\(\ell\nmid h_K\) without any use of the block decomposition.
In equal characteristic it follows directly from the finite inertia
character argument and its trivial reduction. Taking determinants
therefore gives \(\gamma^3=1\). Applying \(c\) to
\(\rho^\vee\simeq\rho\otimes\gamma\) and using
\(\rho^c\simeq\rho^\vee\) shows that \(\gamma\gamma^c\)
is a self-twist. Every finite-order self-twist of a strongly
irreducible representation is trivial: its intertwiner is scalar
on the open kernel by Schur's lemma. Hence \(\gamma^c=\gamma^{-1}\).
The residual characters of \(\rho\) are
\(1,\bar\Psi,\bar\Psi^{-1}\); comparison of the two residual
multisets forces \(\bar\gamma\) to be one of them. All three are
fixed by \(c\). Thus \(\bar\gamma^2=1\) and \(\gamma=1\), since
reduction is injective on the cube roots for \(\ell\ge5\).
If \(C=1\), the construction of \(P,W,\lambda\) and their local
flags is unchanged, with \(\bar\lambda(c)=1\) and \(W(c)=1\).
This reduces that case to an even two-dimensional problem;
the odd-determinant theorem of Pan used above does not settle it.
\end{proof}

\begin{proposition}[A real quadratic self-twist in degree four]\label{hr:ind}
Let \(R\) be an absolutely irreducible point of the good rank-four
problem admitting a nontrivial even quadratic self-twist \(\delta\).
Put \(a=\bar\alpha=\bar\xi\) and \(b=\bar\Psi\). Then \(b=\bar\delta\)
and the inducing field \(M\) is the real quadratic field cut out by
\(b\). In equal characteristic such a point cannot occur if
\(\ell\) is unramified in \(M\). In characteristic zero \(\ell\)
necessarily splits in \(M\), and \(R\) has a regular cuspidal
attachment on \(\mathrm{GL}_4/\mathbb Q\). Its base change to any
imaginary quadratic field on which \(R\) remains irreducible is
cuspidal and has the restricted Frobenius parameters.
\end{proposition}
\begin{proof}
A self-twist intertwiner has scalar square. After scaling it to
square one, its two eigenspaces are preserved by \(G_M\) and
interchanged by the other coset, so both have dimension two. Thus
\(R=\operatorname{Ind}_{G_M}^{G_{\mathbb Q}}U\), with \(U\)
irreducible. Conversely any such quadratic induction admits this
self-twist. On residual semisimplifications multiplication by
\(\bar\delta\) preserves the multiset \(\{1,a,b,ab\}\).
The only possibilities for the image of the trivial character are
the even entries \(1,b\); the quadratic character survives
reduction since \(\ell\) is odd. Over \(G_M\) the residual multiset
is consequently \(1,1,a,a\). The two conjugate blocks have
identical residual semisimplifications: their characters are
chosen from \(1,a|_{G_M}\), both fixed by conjugation from
\(G_{\mathbb Q}\). This follows by reducing lattices in their
direct sum and using independence of the residual semisimplification.
Thus each block has residual characters \(1,a|_{G_M}\).
They have opposite signs at both real places of \(M\).

Intersections of the ordinary local flag with the plane of \(U\)
give a full flag after omitting repeated subspaces. This applies
at either place of \(M\) over \(\ell\) by using both conjugate
planes. In equal characteristic the inertial eigencharacters are
finite at every finite place. The determinant of \(U\) is therefore
the constant character \(a|_{G_M}\), by the character argument
preceding Lemma~\ref{dp:pan}. Apply the proof of that lemma over
\(M\). This is an abelian totally real field unramified at
\(\ell\), the residual ratio extends to \(G_{\mathbb Q}\), and
the set of places may contain every ramified place and those over
\(\ell\). At the places over \(\ell\) the residual ratio is either
nontrivial at all of them or trivial at all, since it is the
restriction of a character of \(G_{\mathbb Q}\) and \(M\) is
Galois. In the latter case use the actual upper character at each
place in the character-labelled ring. The map from the ordinary
weight algebra kills \(\ell\) and each local pro-\(\ell\) inertia
augmentation, including when a place has residue degree two.
Finiteness over this algebra therefore makes every
characteristic-polynomial coefficient constant. Brauer--Nesbitt
contradicts irreducibility of \(U\).
\cite[Theorems~5.1.1(1) and~6.1.1(1)]{panFM}

In characteristic zero the four distinct exponents of \(R\) on
an open subgroup of inertia are \(e+1,e,1,0\), in descending
order. On a common open subgroup of \(I_\ell\cap G_M\) the two
planes have disjoint pairs of exponents. A decomposition element
interchanging them would preserve these exponents, since
\(\varepsilon\) is a global character. Thus \(\ell\) splits in
\(M\). The exponents of \(\det U\) at its two places above \(\ell\)
are the sums of the two respective pairs. Global reciprocity on a
power of a nontrivial totally positive norm-one unit of \(M\)
equates these sums, by the calculation in Lemma~\ref{spreal:even}.
Their common value is \(e+1\), and the pairs are
\(\{0,e+1\},\{1,e\}\). The character \(\det U/\varepsilon^{e+1}\)
has finite inertia at all places and is unramified outside finitely
many. A common power is everywhere unramified and factors through
the finite narrow class group; hence this character has finite
image.

The local upper character of \(U\) has the larger cyclotomic
exponent in each pair and hence strictly smaller Hodge--Tate
weight than the quotient. Both characters are de Rham, since
they have the indicated powers on open inertia. The residual
ratio \(a\) is distinguished there: the chosen Frobenius lift
\(\sigma_\ell\) has \(a(\sigma_\ell)=\bar\xi(\sigma_\ell)\ne1\)
by the good-datum hypothesis and the splitting of \(\ell\).
The ordinary theorem now makes \(U\) Hilbert modular over \(M\),
up to the standard algebraic twist. The coefficient map gives a
maximal ideal after inverting \(\ell\), as in
Proposition~\ref{hr:orth}. To absorb the twist in \(U=W_f\otimes\theta\),
compare the upper character of \(U\) at each place above \(\ell\)
with that of the same line in the de Rham \(W_f\). Their ratio
\(\theta_v\) is de Rham and hence a cyclotomic integer power on
open inertia at the split place. The determinant identity gives
finite inertia for \(\theta\) at all other finite places.
The same norm-one unit argument equates the two exponents; dividing
by their common cyclotomic power leaves a finite-order character.
Both twists thus have algebraic automorphic counterparts.
\cite[Theorem~5.1.1(2)]{panFM}

Write \(\pi_M\) for the corresponding cuspidal representation in
the normalization giving \(U\). Its Galois conjugate has the other
pair of exponents at the first place above \(\ell\), so the two
are inequivalent. Cyclic automorphic induction gives a cuspidal
representation, and the arithmetic normalization in degree four
is obtained by taking
\(\Pi=\operatorname{AI}_{M/\mathbb Q}(\pi_M)|\det|\).
Indeed \(1-3/2=-1/2\), the shift used for \(U\) in degree two.
The unramified parameters therefore give \(R\). At the real place
the parameter is the sum of those at the two real places of
\(M\), with the same shift of normalization; its four arithmetic
exponents are \(0,1,e,e+1\), giving regular algebraicity.
If quadratic base change to \(K\) failed to be cuspidal,
\(\Pi\) would admit its quadratic self-twist. The unramified
polynomials would give the same self-twist of \(R\) and split
\(R|_{G_K}\) by the intertwiner argument above.
\cite[Chapter~III, Theorems~4.2, 5.1 and~6.2, Corollary~6.5]{arthurClozelBC}
\end{proof}

\begin{corollary}[Imprimitive quartic points]\label{hr:imprimitive}
Every absolutely irreducible characteristic-zero point of the good
quartic problem which fails strong irreducibility is induced from a
plane over the real quadratic field cut out by \(b=\bar\Psi\).
It has the cuspidal attachment of Proposition~\ref{hr:ind} and
its finite symplectic adjoint Selmer group over \(\mathbb Q\)
vanishes. No such point permutes a system of four lines.
\end{corollary}
\begin{proof}
First consider any self-twist \(R\simeq R\otimes\delta\).
Determinants give \(\delta^4=1\). Over the decomposition group
at \(\ell\) the flag of Proposition~\ref{spreal:four} has
characters \(\kappa_i\). On an open subgroup of inertia in
\(\ker\delta\) they are four distinct cyclotomic powers.
Thus the permutation matching \(\{\kappa_i\delta\}\) with
\(\{\kappa_i\}\) fixes each character, and
\(\delta|_{G_{\mathbb Q_\ell}}=1\).
The image of the trivial residual constituent under twisting gives
\(\bar\delta\in\{1,a,b,ab\}\). The specified element \(s_\ell\)
of this decomposition group satisfies
\(a(s_\ell)b(s_\ell)^i\ne1\) for all integers \(i\). This
eliminates \(a,ab\). Reduction is injective on the fourth roots of
unity, so the self-twist group has order at most two, and its sole
possible nonidentity element is even quadratic with reduction \(b\).
\src{src/Development/AutomorphicArithmetic/AutomorphicVarieties/SymplecticGSp4/DecompSigmaOrd.lean}{114}{The local character condition of the good datum}.

For four permuted lines choose the direct sum of their intersections
with a stable lattice. The nonzero entries of the resulting monomial
matrices are units, so their permutation quotient \(P\) survives
reduction. Independence of residual semisimplification gives an
invariant residual flag with one-dimensional quotients. The derived
group of the finite residual image is therefore an \(\ell\)-group,
as it is contained in the upper unitriangular matrices. Hence
\(P'\) is a normal \(\ell\)-group in the transitive group \(P\).
Its orbits on the four lines all have the same size, a divisor
of four and a power of \(\ell\). Thus it fixes every line and is
trivial as a permutation group. This holds also for \(\ell=3\).
The faithful transitive abelian group \(P\) is regular of order
four. Index the lines by \(P\). Multiplying the line indexed by
\(g\) by \(\chi(g)\) gives a self-twist for each character
\(\chi\) of \(P\). These four distinct twists contradict the
bound of two.

Now restrict to an open normal subgroup
\(H\) on which the representation is reducible. The translates
of any simple \(H\)-subspace span \(R\), so the restriction is
semisimple and its isotypic subspaces are transitively permuted.
Every simple constituent has multiplicity one, because some
element of \(I_\ell\cap H\) has four distinct eigenvalues by the
four distinct cyclotomic exponents on open inertia. Thus the
remaining alternative is two inequivalent planes. Their permutation
gives the quadratic self-twist just identified, so
Proposition~\ref{hr:ind} applies. Choose an imaginary quadratic
\(K\) unramified at \(\ell\) with irreducible restriction by
Lemma~\ref{rk4:quadratic}. Its cyclotomic-kernel irreducibility
follows from Lemma~\ref{spreal:twists}, since the ordinary characters
remain distinct on open inertia and inertia surjects onto the
cyclotomic quotient. The cuspidal base change supplied by
Proposition~\ref{hr:ind} and the given similitude pairing now
satisfy Proposition~\ref{newgal:adselmer}(2). Its symplectic
direct summand has zero finite Selmer group. In particular every
tangent of Proposition~\ref{spreal:four} is a coboundary there.
\end{proof}

\begin{proposition}[Strong irreducibility on the quartic test curves]
\label{hr:case3strong}
For the original good quartic datum in its third host character case,
every absolutely irreducible point in either characteristic remains
absolutely irreducible over every finite number-field extension.
In the second host case an absolutely irreducible equal-characteristic
curve failing this condition requires \(b\) to cut out a real
quadratic field ramified at \(\ell\).
\end{proposition}
\begin{proof}
Write \(N\) for the host level and \(a=\bar\xi\),
\(b=\bar\varepsilon^e\bar\psi\). For a self-twist, determinants
give \(\delta^4=1\) in either characteristic. Its reduction is one
of \(1,a,b,ab\). We check the orders of these characters on the
original group.
\src{src/Development/AutomorphicArithmetic/AutomorphicVarieties/AutomorphicCases/SymplecticGSp4.lean}{15}{The conductor factors of the auxiliary character}
and \src{src/Development/AutomorphicArithmetic/AutomorphicVarieties/HeightAndVanishing/DirichletPrimeToLPartMulLPowerPart.lean}{221}{The exponents in the two host cases}.

In case three, \(m=\operatorname{ord}(\psi)>1\) is odd and prime to
\(\ell\). The auxiliary odd prime \(p_1\) divides neither \(N\)
nor \(\ell\), and \(\ell\nmid N\). The character \(\xi\) is the
product of a character of two-power conductor and a character of
order \(m\) modulo \(p_1\). Thus on inertia at \(p_1\) the
characters \(a,ab\) both have order \(m\): inertia surjects
onto \(\operatorname{Gal}(\mathbb Q(\mu_{p_1})/\mathbb Q)\),
and reduction preserves its prime-to-\(\ell\) character values.
Also the image of \(b\) on \(\ker\bar\varepsilon\) contains the
image of \(\bar\psi\), by the Chinese remainder theorem at
\(N,\ell\). None of these three characters has fourth power one.
\src{src/Development/AutomorphicArithmetic/AutomorphicVarieties/HeightAndVanishing/BchiFacts.lean}{8}{The coprime conductors in case three}
and \src{src/Development/AutomorphicArithmetic/AutomorphicVarieties/WeakShimuraTheory/DirichletGroupPowers/DirGPowOne.lean}{77}{The order prime to the coefficient prime}.

In case two \(\xi\) is primitive of conductor \(2^n\). Put
\(u=v_2(N)\ge2\). By construction \(2^n\ge\ell N\);
the odd part of \(N\) contains an odd host prime, so \(n\ge u+3\).
The subgroup \(1+2^u\mathbb Z_2\) modulo \(2^n\) is cyclic
of order \(2^{n-u}\). Its order-two subgroup is the kernel of
reduction modulo \(2^{n-1}\), where primitivity makes \(\xi\)
nontrivial. Thus \(a\) is faithful on that cyclic subgroup.
On the corresponding inertia at two, \(b=1\), since \(\psi\)
has level \(N\) and \(\ell\) is odd. Consequently neither
\(a\) nor \(ab\) has fourth power one. Reduction is injective
on the fourth roots of unity, so the self-twist group has size at
most two, and any nontrivial twist is the even quadratic lift of
\(b\). In case three it is trivial. Corollary~\ref{hr:imprimitive}
now gives the characteristic-zero assertion.

In equal characteristic let \(H\) be an open normal subgroup of
\(G_{\mathbb Q}\) on which the representation becomes reducible
over an algebraic closure of the coefficient field. Its restriction
is semisimple by the sum-of-translates argument above. Thus the
commutant of \(H\) has dimension greater than one, and its trace-zero
part \(W\) is nonzero. Conjugation on \(W\) has finite image since
\(H\) acts trivially. In any stable lattice the residual
semisimplification is a sum of characters. Triangularize its
residual action. The preimage of the upper unitriangular subgroup
in the integral image is pro-\(\ell\), by the successive congruence
quotients, and its quotient is diagonal and abelian. Hence
the finite group acting on \(W\) has a normal \(\ell\)-subgroup
with abelian quotient. A finite \(\ell\)-group acting in characteristic
\(\ell\) has nonzero invariants: use a central element of order
\(\ell\), take its fixed space, and induct on the quotient.
The abelian quotient has a common eigenline in that invariant space.
A nonzero endomorphism on this line intertwines the original
representation with a character twist and is invertible by
irreducibility. The character is nontrivial, since a scalar
endomorphism of trace zero is zero for odd \(\ell\) in dimension
four. This excludes the curve in case three. In case two it forces
the quadratic self-twist of \(b\); Proposition~\ref{hr:ind}
excludes it whenever \(\ell\) is unramified in the inducing field.
\end{proof}

There is a matrix criterion for this induction argument. Suppose
the absolutely irreducible \(R\), already symplectic with multiplier
\(\mu\), also preserves a symmetric form up to a character \(\tau\).
Comparing the two identifications with the dual gives a self-twist
\(\delta=\tau/\mu\). It is nontrivial by Schur's lemma since one
form is symmetric and the other alternating. For the two original
host cases the character calculation in
Proposition~\ref{hr:case3strong} forces case two and the real
quadratic character lifting \(b\). Thus the characteristic-zero
point is covered by Proposition~\ref{hr:ind}, and an
equal-characteristic point requires ramification of \(\ell\) in
that quadratic field.
In particular a genuine tensor product \(R=X\otimes Y\) of two
degree-two representations has a symmetric form with
\(\tau=(\det X)(\det Y)\), by tensoring their alternating determinant
pairings. Here \(\tau^2=\det R=\mu^2\), so the self-twist is
quadratic. These criteria concern the representation on the
original group; the existence of an auxiliary tensor-factor
tower alone supplies no such form there.

The automorphy assertions here concern classical cuspidal objects
and their unramified arithmetic parameters. The rank-four source
predicate also asks for a witness in its specified automorphic
model for every qualifying quadratic CM restriction.
\src{src/Development/AutomorphicArithmetic/AutomorphicVarieties/ModularityArgument/AutomorphyLifting/IsAutCsdA.lean}{16}{The two clauses of the quadratic-restriction predicate}.

\begin{lemma}[An unramified test on the residual field]\label{hr:class}
In either rank let \(F_0\) be the field of the kernel of the actual
diagonal residual representation (on \(G_K\) in rank three).
Suppose the \(\ell\)-part of the narrow class group of \(F_0\) is
cyclic. Every equal-characteristic test curve unramified at all finite
places on \(G_{F_0}\) has constant characteristic polynomials.
It suffices that the curve be tame above \(\ell\) over \(F_0\)
and have inertia image of order prime to \(\ell\) at the remaining
finite places over the allowed ramified primes.
\end{lemma}
\begin{proof}
The image of \(G_{F_0}\) lies in the first congruence group and is
pro-\(\ell\). Tame inertia alone can still have an \(\ell\)-part
away from \(\ell\); at those places the defining inertia kernels
must also kill that part. In rank three they kill the kernel of the
residual action. In rank four, at a prime away from \(\ell\) where
the prescribed condition only kills the kernel of \(\psi\), require
in addition that the order of its allowed inertia quotient be
prime to \(\ell\), or directly that the image on \(G_{F_0}\) be
trivial. At places over \(\ell\) tame inertia has pro-prime-to-\(\ell\)
order. Under these local conditions the pro-\(\ell\) image is
everywhere unramified. Any finite quotient is an \(\ell\)-group
whose abelianization is a quotient of the cyclic \(\ell\)-class
group, so is cyclic by the Frattini criterion. Its order is bounded
by the same class number. Passing to the inverse limit gives
finite image, also over the original number field since \(F_0\)
has finite degree. Every eigenvalue of every matrix has finite
order and is thus constant over the algebraic closure of the finite
residue field. Reduction recovers each characteristic polynomial.
\end{proof}

\subsection{Which ideal the tower controls}\label{hr:tower}

Write \(F\) for the real field of a fit, deep rank-four tower and
\(L/F\) for its quadratic CM extension. Let \(X.A\) be its framed
ring and \(X.B\hookrightarrow X.A\) the coefficient ring of the
characteristic polynomials and the ordered local characters. The
following finite-fibre calculation isolates the use of an ordinary
Hecke algebra in the source.

\begin{lemma}\label{hr:towerfin}
Suppose an auxiliary link supplies an ordinary determinant ring
\(D_1\) over its weight algebra \(\Lambda_1\), an ideal \(J_1\)
such that \(D_1/J_1\) is finite over \(\Lambda_1\), and the link
map \(f_1:D_1\to X.B\). This map identifies the determinant on
\(G_L\) and its local characters with the restricted ones twisted
by the link character \(\kappa\). Put \(O=c'.O\) and
\[
 U_{\mathrm{in}}=(\psi_{v,i}(s)-1\mid v\mid\ell,\ 
                   1\le i\le4,\ s\in I_v)\subset X.B .
\]
Then for some \(m\),
\[
 \mathfrak m_{X.B}^{m}
       \subset f_1(J_1)X.B+\mathfrak m_O X.B+U_{\mathrm{in}}.
\]
\end{lemma}
\begin{proof}
Pass to the quotient by any prime containing the right hand side.
It is a local domain of characteristic \(\ell\) with the fixed
finite residue field. The inertia weight generators of \(\Lambda_1\)
map to constants: after the fixed link twist the remaining ordered
inertia characters are one. The local units of \(O\), including
the values of the link twist, map to their residues. By the
universal property of the weight algebra its map is the residual
one, since that map has the same labelled characters. Finiteness
therefore makes the image of \(D_1/J_1\) integral over the finite
residue field. In a local domain with residue field \(k\), any
element integral over the coefficient field \(k\) belongs to \(k\):
the subalgebra it generates is a field and reduction embeds that
field into \(k\) over \(k\).

Hence the characteristic coefficients on \(G_L\) are constant
after removal of the scalar twist. For \(g\in G_F\), use
\(g^2\in G_L\). Each elementary symmetric polynomial in four
variables \(x_i\) is integral over the symmetric polynomials in
\(x_i^2\), since each \(x_i\) satisfies
\(\prod_j(T^2-x_j^2)=0\). The resulting monic identities with
integral symmetric coefficients hold on every matrix by the
characteristic polynomial, hence in \(X.B\) by its injection
into the framed ring. They survive the prime quotient without
any lifting of its prime to \(X.A\). Thus the characteristic
coefficients for \(g\) are constant there. Each marked decomposition
group lifts to the one on \(L\) by the splitting at these places;
the corresponding labelled characters, after division by the link
twist, are constant as well. These data topologically generate
\(X.B\), so the prime is its maximal ideal. In the noetherian ring
\(X.B\), the radical inclusion gives the asserted power.
\src{src/Development/AutomorphicArithmetic/AutomorphicVarieties/ModularityArgument/ThresholdArguments/DenseThresholdCases/TPmFin.lean}{126}{The prime-by-prime calculation for one link}.
\end{proof}

The ideal actually used is
\[
 J=f_1(J_1)X.B+f_2(J_2)X.B .
\]
The two links have the same twisting character \(\kappa\) and
distinct auxiliary primes. Lemma~\ref{hr:towerfin} needs only
the first link. The two links give separate local vanishing
statements at a seed and at the prescribed curve points. The
characteristic-zero automorphy assertions away from the separate
auxiliary primes are combined using their distinctness and the
common twist. The actual quantifier chooses a lower coefficient
extension for the tower, then provides such an ideal for every
larger coefficient DVR and every representing space \(X\).
\src{src/Development/AutomorphicArithmetic/AutomorphicVarieties/ModularityArgument/ThresholdArguments/DenseThresholdCases/LinkTwo.lean}{13}{Two links with a common twist};
\src{src/Development/AutomorphicArithmetic/AutomorphicVarieties/ModularityArgument/ThresholdArguments/DenseThresholdCases/PmTwo.lean}{21}{Sum of the two extended ideals}.

On the unitary side \(J_i\) is the unit ideal when the specified
ordinary module vanishes. Otherwise it is the kernel of a
surjection of the universal ordinary determinant ring onto the
localized Hecke algebra: the interpolated determinant gives the
map, the unramified Hecke polynomials and the ordinary uniformizer
eigencharacters recover the Hecke generators, and the map agrees
on the weight algebra. Finiteness of this Hecke algebra over the
weight algebra then proves finiteness of the quotient. Construction
of the interpolated determinant uses separation by arithmetic
points with ordinary Galois attachments. The additional hypotheses
in the modular ideal include density, the point theorem,
purity and classicality at the indicated points, twist compatibility
and the two component assertions.
\src{src/Development/AutomorphicArithmetic/AutomorphicVarieties/ModularityArgument/ThresholdArguments/DenseThresholdCases/ExistsHeckeIdeal.lean}{28}{The zero-module case and the Hecke kernel};
\src{src/Development/AutomorphicArithmetic/AutomorphicVarieties/ModularityArgument/ThresholdArguments/DenseThresholdCases/TJRing.lean}{172}{Transfer of ordinary Hecke finiteness};
\src{src/Development/AutomorphicArithmetic/AutomorphicVarieties/ModularityArgument/ThresholdArguments/DenseThresholdCases/THidaU.lean}{38}{Assembly of the unitary ideal properties}.

For a given curve the required extra assertion is that after
the specified finite base changes and an integral choice of
lattice there is \(y:X.A\to\mathcal O_{\rm curve}\) cutting out
its restricted representation, killing \(J\), and giving the
original labelled inertia characters. On a fit tower those
characters reduce to one on the marked inertia subgroups. The
lemma kills \(\mathfrak m_{X.B}\) under \(y\) since the target is
a domain. All characteristic polynomials on \(G_F\) are then the
residual ones. For every \(g\in G_{\mathbb Q}\) some positive
power belongs to \(G_F\). Every eigenvalue of \(g\) is thus a
root of unity, hence constant in equal characteristic; reduction
again gives the specified polynomial. This excludes every such
curve with a higher-dimensional simple constituent by
Brauer--Nesbitt. On the marked model the projectors used in
Proposition~\ref{dp:curvebranches} then give constancy of the
weight-zero entry monomials themselves.
\src{src/Development/AutomorphicArithmetic/AutomorphicVarieties/ModularityArgument/ThresholdArguments/DenseThresholdCases/PotentialModularityTower.lean}{39}{The finite-fibre, seed and component requirements};
\src{src/Development/AutomorphicArithmetic/AutomorphicVarieties/ModularityArgument/ThresholdArguments/DenseThresholdCases/ContradictionStarPromodular.lean}{135}{Recovering the residual characteristic polynomial over the tower}.

In the rank-four case-three construction, local vanishing at a
characteristic-zero seed gives inclusion on some component at its
residual point. A chain of linked residual frames carries this
to the residual point of the proposed test. At each good residual
point a separate component theorem promotes inclusion on some
component to inclusion on all of them. Only the latter conclusion
gives \(y(J)=0\) for the chosen lift. The source chooses a layer
with threshold \(30(C+1)^2\), where \(C\) is supplied by its
cohomology bound, and requires a seed over that layer and a
descending chain for the specific residual lattice. Its seed
local-vanishing proof also uses adjoint lifting. Merely producing
the ideal, or one characteristic-zero point of its zero set,
does not yield this quantifier over curves.
\src{src/Development/AutomorphicArithmetic/AutomorphicVarieties/ModularityArgument/ThresholdArguments/DenseThresholdCases/PromodularityCaseThreeOfLemmas.lean}{131}{The choice of layer, seed and chain};
\src{src/Development/AutomorphicArithmetic/AutomorphicVarieties/ModularityArgument/ThresholdArguments/DenseThresholdCases/BirthPm.lean}{28}{Adjoint lifting at a seed}.

In the rank-three construction there are separate ordinary
determinant ideals in degrees two and three. For each degree
the input supplies realizations with an auxiliary prime avoiding
any prescribed nonzero integer, and a seed specializing into
each ideal satisfying the point theorem. Component propagation
first treats degree two; degree three also uses the splitting
condition from the third component property and the degree-two
conclusion. It produces the degree-three finite-point assertion.
A second choice of auxiliary prime in each degree removes the
exceptional place from the two automorphy assertions.
\src{src/Development/AutomorphicArithmetic/AutomorphicVarieties/ModularityArgument/ThresholdArguments/DenseThresholdCases/PrLinkedEnd.lean}{22}{Specialization and the two modular ideals};
\src{src/Development/AutomorphicArithmetic/AutomorphicVarieties/CurveDVRSpecialization/AbsMain.lean}{12}{Propagation in the two degrees};
\src{src/Development/AutomorphicArithmetic/AutomorphicVarieties/ModularityArgument/ThresholdArguments/DenseThresholdCases/TwoRuns.lean}{80}{The two realizations and change of auxiliary prime}.

\paragraph{The ternary seed on the linked CM field.}
Here \(L\) is the CM field already chosen for the test. Write
\(\varepsilon\) for the cyclotomic character,
\(\Phi=\varepsilon^j\chi'\) for the prescribed lift of the residual
character and \(\eta\) for the auxiliary quadratic character. The
seed theorem requires \(\chi'\ne\eta\), \(M\ne0,\ 2\ell\mid M\),
the extension \(L/L^+\) unramified at every finite place and split
at all places dividing \(M\),
and triviality of \(\chi'|_{G_L}\) on inertia away from \(M\).
The linked field and character conditions include \(K\subseteq L\)
for the imaginary quadratic field of \(\eta\), \(L/\mathbb Q\) abelian,
and triviality of the auxiliary character \(\theta\) on \(G_L\).
In the seed datum the positive auxiliary conductor \(N'\) is prime to
the odd prime \(\ell\). The prime \(w\nmid N'\) acts trivially through
both \(\eta\) and \(\chi'\) on \(\mathbb Q(\mu_{N'})\), through which
they factor. The character \(\theta\) has order \(\ell\) and factors
through \(\mathbb Q(\mu_w)\). Also \(\chi'(c)=(-1)^j\) for complex
conjugation and \(\bar\Phi^2|_{G_L}\ne1\). The link and the congruence with
the prescribed level character give the two local hypotheses at
the point of use. For every lower bound \(b\) the theorem produces
\(m\ge b,\ m\equiv j\bmod(\ell-1)\), and a characteristic-zero
representation of \(G_L\) irreducible on open subgroups, with residual
characters \((\bar\Phi,1,\bar\Phi^{-1})|_{G_L}\), away-from-\(\ell\)
inertial types \((\chi',1,(\chi')^{-1})|_{G_L}\), ordered ordinary
exponents \((m,0,-m)\), a conjugate-dual pairing and automorphic
attachment after twisting by \(\varepsilon^{-1}|_{G_L}\).
Here one takes \(b=2\).
\src{src/Development/AutomorphicArithmetic/AutomorphicVarieties/IdeleGenerators/SlopeTestRemainingChoices.lean}{136}{Field and character conditions},
\src{src/Development/AutomorphicArithmetic/AutomorphicVarieties/ModularityArgument/ThresholdArguments/DenseThresholdCases/SymUse.lean}{27}{Ternary seed theorem}
and \src{src/Development/AutomorphicArithmetic/AutomorphicVarieties/ModularityArgument/ThresholdArguments/DenseThresholdCases/LinkedEndE.lean}{51}{Its use on the linked field}.

In the construction this is \(\operatorname{Ad}^{0}(\rho)|_{G_L}\)
for an auxiliary degree-two Eisenstein congruence. An upper triangular
action on a plane with diagonal \(a,b'\) induces on its traceless
endomorphisms a flag with diagonal \(a/b',1,b'/a\): take first the
strictly upper triangular endomorphisms, then all upper triangular
traceless endomorphisms. This proves the claimed exponents from
the difference \(m\) of the two ordinary exponents. The auxiliary
\(\eta,\theta\) disappear on \(G_L\). The trace pairing on traceless
endomorphisms is symmetric and invariant in characteristic zero.
Open irreducibility, existence of \(\rho\) and automorphic attachment
of its twisted adjoint are arithmetic inputs to this calculation.
\src{src/Development/AutomorphicArithmetic/AutomorphicVarieties/ModularityArgument/ThresholdArguments/DenseThresholdCases/AdjointLiftFromEisensteinSeedOfHyps.lean}{184}{Inputs and construction}
and \src{src/Development/AutomorphicArithmetic/AutomorphicVarieties/ModularityArgument/ThresholdArguments/DenseThresholdCases/AdjointLiftFromEisensteinSeedOfHyps.lean}{250}{Restriction of the adjoint and its local conditions}.

For either degree \(n=2,3\) the seed specialization is taken after
the link twist \(\kappa\). The integral comparison characters are
\(\Psi^{w_i}|_{G_L}\kappa\), with \(w_i=(1,-1)\) or \((1,0,-1)\)
and \(\Psi\) the original integral character. Their reductions are
the labelled residual characters in the determinant problem.
Away from \(\ell\) an inertial eigenvalue in the required kernel is
a root of unity \(u\) of order dividing a number \(h\) prime to \(\ell\),
with \(u\equiv1\) in the residue field. Since
\(1+u+\cdots+u^{h-1}\equiv h\ne0\), the identity \(u^h=1\)
gives \(u=1\). Diagonality on that inertia subgroup kills the
whole action there. Above \(\ell\) scalar twisting multiplies each
graded character by \(\kappa\) in the same flag, preserving their
order. Thus the representation and its labelled characters give a
point of the ordinary determinant ring. The cited arithmetic
point theorem puts this particular point in \(V(J)\).
\src{src/Development/AutomorphicArithmetic/AutomorphicVarieties/ModularityArgument/ThresholdArguments/DenseThresholdCases/PrSeedPtOf.lean}{132}{Integral comparison characters},
\src{src/Development/AutomorphicArithmetic/AutomorphicVarieties/ModularityArgument/ThresholdArguments/DenseThresholdCases/PrSeedPtOf.lean}{198}{Inertia and flags}
and \src{src/Development/AutomorphicArithmetic/AutomorphicVarieties/ModularityArgument/ThresholdArguments/DenseThresholdCases/PrSeedPtOf.lean}{219}{Application of the point theorem}.

Write \(B\) for this determinant ring and choose a minimal prime \(Z\)
of \(J\) below the kernel of the seed point. Put \(C=B/Z\). If the
ordinary quotient is finite torsion-free over
\(\Lambda\simeq\mathcal O[[X_1,\ldots,X_{nd}]]\), every nonzero
element of \(\Lambda\) is a nonzerodivisor on \(B/J\) and therefore
misses its minimal primes. Hence \(\Lambda\hookrightarrow C\) is
integral, and lying over, going up and incomparability give
\(\dim C=1+nd\). Absolute irreducibility persists at the generic
point. Choose \(n^2\) group matrices spanning the full matrix algebra
at the seed. Their trace Gram determinant is nonzero there, since
the matrix trace pairing is nondegenerate. Its value for the
universal determinant is consequently nonzero in \(C\).
Over the algebraic closure of \(\operatorname{Frac} C\), the
determinant realization and untwisting theorem gives the framed
point. It uses the condition that the determinant kills the specified
auxiliary inertia modulo \(J\). The Gram determinant shows its
image algebra has dimension \(n^2\).
The reverse map from the framed coefficient ring to \(C\) composed
with the determinant link is \(B\twoheadrightarrow C\), by uniqueness
on the determinant and ordered local characters. It is surjective,
so its kernel has quotient dimension \(1+nd\) and lies under the
framed generic point. This is the seed dimension needed for curve
selection. In the actual frame \(d=\sum_v d_v>3C_0+7\) and every
\(d_v>5\), with \(C_0\) the cohomological bound in its hypotheses.
The arithmetic transfers at separated curves are invoked in degree
two first and then in degree three, where the additional exceptional-chain
transfer is used.
\src{src/Development/AutomorphicArithmetic/AutomorphicVarieties/CurveDVRSpecialization/SeedOf.lean}{501}{The chosen minimal prime and untwisting},
\src{src/Development/AutomorphicArithmetic/AutomorphicVarieties/CurveDVRSpecialization/SeedOf.lean}{525}{The framed point and quotient dimension},
\src{src/Development/AutomorphicArithmetic/AutomorphicVarieties/FrobeniusSmoothVectors/ContinuousCoboundaries2.lean}{212}{The strict degree inequalities}
and \src{src/Development/AutomorphicArithmetic/AutomorphicVarieties/CurveDVRSpecialization/AbsMain.lean}{21}{The arithmetic transfers}.

\paragraph{The quartic seed and the four links.}
In the rank-four case-three tower, the input is a good host datum
with \(0<e<t\) in its third character case. The integral test
\((r,u)\) has the prescribed residual matrices and similitude,
the stated global ramification bound and kernel on inertia away
from \(\ell\), and a symplectic flag matrix \(u\) above \(\ell\).
Its ordered diagonal characters agree with the prescribed
characters on inertia and reduce to them on the decomposition
group. The test is defined over a small coefficient ring.
On the fraction field the case under consideration is either
irreducible or the specified nonsplit block extension with full
two-dimensional block algebra. These are the hypotheses on the
test to which the propagation is applied.
\src{src/Development/AutomorphicArithmetic/AutomorphicVarieties/WeakShimuraTheory/NearSplit.lean}{59}{Host and integral test}
and \src{src/Development/AutomorphicArithmetic/AutomorphicVarieties/WeakShimuraTheory/RedIdeal.lean}{65}{The block-extension alternative}.

Start with its actual lattice shape on the first tensor-factor tower \(T_0\).
The seed theorem supplies a rational prime \(w\) split in the
specified tower sense and a residual framed point \(s_0\), initially
allowing \(w\) in its ramification set. Its assertion at higher
levels is universal. For every allowable \(T\ge T_0\) at which
all continuous order-\(\ell\) characters of \(G_{\mathbb Q}\)
unramified away from \(w\) become trivial on \(G_{F_T}\), it
supplies an integral pair \(Q_T\) over a finite coefficient DVR.
Its reduction lies over \(s_0\), has the specified root-entry
nonvanishings, and obeys the frame with the auxiliary set now empty.
Its generic fibre is irreducible on the designated subgroup.
After restriction to the CM field of \(T\) and a specified scalar
twist it is attached to a regular algebraic cuspidal degree-four
representation. Its consecutive weight differences are \(2n_i\),
\(n_i\ge2\). The source uses a tensor product of two planes. It
chooses \(1\le j\le9\), a positive \(m\), and \(k\ge2(9m)+3\).
The three half-gaps in this construction are
\(2mj,\ k-1-2mj,\ 2mj\), all at least two. The restriction
condition at \(w\) kills the auxiliary character on \(G_{F_T}\),
which is where the integral pairing without the extra allowed prime
is checked.
\src{src/Development/AutomorphicArithmetic/AutomorphicVarieties/WeakShimuraTheory/NearSplit.lean}{175}{Condition at the auxiliary prime},
\src{src/Development/AutomorphicArithmetic/AutomorphicVarieties/ModularityArgument/ThresholdArguments/DenseThresholdCases/BirthSeedIV.lean}{34}{Tensor seed and size of the gaps},
\src{src/Development/AutomorphicArithmetic/AutomorphicVarieties/ModularityArgument/ThresholdArguments/DenseThresholdCases/BirthSeedIV.lean}{48}{Residual and integral pair on each allowed layer}
and \src{src/Development/AutomorphicArithmetic/AutomorphicVarieties/ModularityArgument/ThresholdArguments/DenseThresholdCases/SeedOf4P.lean}{50}{Weights and regular cuspidal attachment}.

For the given test lattice the cohomological construction chooses a
finite set of split auxiliary primes containing \(w\) and then
\(T_1\ge T_0\) killing all the indicated order-\(\ell\) characters
for each auxiliary prime separately. This condition
persists under further restriction of Galois groups. The cohomology
bound supplies \(C\) over the resulting tame tower, before the choice
of a deep layer \(T\ge T_1\) with \([F_T:\mathbb Q]\ge30(C+1)^2\).
The seed theorem consequently applies at this very layer. Local
conditions on the restricted test can be checked in coordinates.
For a selected \(v\mid\ell\) choose \(g_v\) carrying the two chosen
prolongations into agreement. In the transported flag the \(i\)th
character at \(s\in I_v\) is the earlier \(i\)th character at
\(g_v^{-1}s g_v\) (using restriction when viewing \(s\) over
\(\mathbb Q\)). This equals the prescribed global character
\(\psi_i(s)\), since a character is invariant under conjugation.
To see why the flag has an integral normalized basis, put
\(Y=r(g_v)n_0\) in the small local coefficient ring, with \(n_0\)
the original lower unipotent flag matrix. The prescribed residual
support gives \(\bar Y_{11}=\bar r(g_v)_{11}\),
\(\bar Y_{12}=0\) and \(\bar Y_{22}=\bar r(g_v)_{22}\) in one-based
indices. Both leading principal minors of orders one and two are
therefore units. The similitude multiplier \(\mu\) is a unit.
The third leading minor is \(\det(Y)(Y^{-1})_{44}\), a unit because
\((Y^{-1})_{44}=Y_{11}/\mu\) in the chosen antidiagonal
symplectic basis. Gaussian elimination over the local ring now
gives \(Y=NU\) with \(N\) lower unipotent and \(U\) upper triangular
invertible. Let \(X^\dagger\) denote the symplectic adjoint.
The automorphism \(X\mapsto(X^\dagger)^{-1}\) preserves lower
unipotent and upper triangular matrices and takes \(Y\) to
\(\mu^{-1}Y\). Uniqueness of triangular factorization gives
\((N^\dagger)^{-1}=N\), so the new flag basis is symplectic as
required. Restricting the multiplier identity and the
specified inertia kernels then gives the framed pair on \(T\).
\src{src/Development/AutomorphicArithmetic/AutomorphicVarieties/WeakShimuraTheory/RootsOfUnityShapes/ShapeQ.lean}{6}{The residual pivots},
\src{src/Development/AutomorphicArithmetic/AutomorphicVarieties/WeakShimuraTheory/CompactIntegralFields/RestrQ.lean}{77}{Transport and integral normalization of the flag}
and \src{src/Development/AutomorphicArithmetic/AutomorphicVarieties/ModularityArgument/ThresholdArguments/DenseThresholdCases/PromodularityCaseThreeOfLemmas.lean}{138}{Choice of lattice, cohomology bound, layer and seed}.

At that layer the explicit residual chain, in the notation of the
construction, is
\[
 q_0\ \longleftrightarrow\ q_3\ \longleftrightarrow\ q_4\
 \longleftrightarrow\ q_2\ \longleftrightarrow\ \overline Q_T .
\]
The first vertex is the restricted test, identified by uniqueness
over its chosen lattice shape. The four edges alternate between
two splittings into pairs of coordinate planes, with indices
\((1,3)\mid(2,4)\) on the first and third edges and
\((1,2)\mid(3,4)\) on the second and fourth. Here is the
reason they give actual lines in the representation scheme.
In the indicated order the two endpoint matrices are block lower
triangular with identical diagonal blocks. The same property
holds for their local flag matrices. Put
\(X(t)=X_0+t(X_1-X_0)\), with \(t\) an indeterminate. For two
such pairs, \((X_1-X_0)(Y_1-Y_0)=0\), since each difference has
only a lower left block. Expansion gives
\[
 X(t)Y(t)=X_0Y_0+t(X_1Y_1-X_0Y_0).
\]
The product pair again has the same diagonal blocks. The
symplectic adjoint is linear and preserves this block condition
in both orders. The formula consequently interpolates the group
law, the similitude identity and the conjugated flag matrices.
Equal diagonal entries, prescribed zeros and the inertia kernels
persist along the line. Evaluating at \(t=0,1\) recovers the
two framed endpoints in one irreducible image.
\src{src/Development/AutomorphicArithmetic/AutomorphicVarieties/WeakShimuraTheory/CompactIntegralFields/AffSmulOne.lean}{43}{First block decomposition}
and \src{src/Development/AutomorphicArithmetic/AutomorphicVarieties/WeakShimuraTheory/CompactIntegralFields/GlobalTheory.lean}{131}{Second block decomposition and flag matrices}.

In fact each nonconstant link is a closed affine line in the framed
special fibre. Write \(A=X.A\), \(B=X.B\) and \(k\) for the
residue field of the common coefficient DVR. The representing map
\(A\to k[t]\) kills \(\mathfrak m_B A\) by the line clause of
the representing space. If the framed endpoints differ, some
matrix or flag entry maps to \(c+dt\), \(d\ne0\). The coefficient
DVR maps onto \(k\), so the image contains \(t\). Thus the map is
surjective and its kernel lies below both point ideals. Every proper
coefficient ideal, including \(J\), vanishes scheme-theoretically
on all four links and on their formal arcs. A minimal prime of
\(A\) below the kernel supplies one ambient component through both
endpoints. This gives the transfer from \emph{all} components at
one end to \emph{some} component at the other.
\src{src/Development/AutomorphicArithmetic/AutomorphicVarieties/WeakShimuraTheory/RedIdeal.lean}{233}{The representing map for the residual line}
and \src{src/Development/AutomorphicArithmetic/AutomorphicVarieties/WeakShimuraTheory/CompactIntegralFields/AlgHomResUnique.lean}{56}{The two directions of component transfer}.

Even seed local vanishing and a finite weight fibre do not by
themselves extend this conclusion to a chosen test curve. For an
explicit calculation over a complete DVR \(O\) of uniformizer
\(\pi\) and residue field \(k\), take
\[
 B=O[[u,v]]/(uv),\qquad A=B[t]/(ut),\qquad J=(u).
\]
With weight parameter \(v\), the fibre \(B/(J,\pi,v)=k\).
The integral seed \((u,v,t)=(0,\pi,1)\) satisfies the stated
coefficient local vanishing at its generic point: every prime
below its kernel misses \(v\), hence contains \(u\). Its reduction
and \((0,0,0)\) lie on the closed line \(A/\mathfrak m_B A=k[t]\).
However the equal-characteristic test
\((u,v,t)=(z,0,0)\) in \(k[[z]]\) specializes to the latter
endpoint and pulls \(J\) back to \((z)\). The two minimal primes
of \(A\) there are \((u)\) and \((v,t)\); only the first contains
\(J\). This identifies the precise additional conclusion required
at a vertex with more than one component.

The root-entry and support tests at the five vertices verify the
good-point hypotheses of the component theorem. Enlarge the coefficient
DVR once to contain the seed coefficients and realize the residual
chain, a ring of definition of the test, and the lower coefficient
extension for the tower ideal over a common coefficient base.
The universal lifts now give points of one representing space. Local
vanishing at the generic seed puts some component at
\(\overline Q_T\) in \(V(J)\). The good-point component theorem at this
vertex gives all components there. The last line then gives some
component at \(q_2\), and the same theorem gives all components
there. Repeating backwards along the three remaining lines gives all
components at \(q_0\). Thus the specific generization representing
the test kills \(J\). This argument uses the cited arithmetic
construction of the lines, the seed local-vanishing theorem and
the good-point component theorem for the case-three hypotheses.
\src{src/Development/AutomorphicArithmetic/AutomorphicVarieties/ModularityArgument/ThresholdArguments/DenseThresholdCases/CuPChainIV.lean}{25}{Enlargement killing the auxiliary primes},
\src{src/Development/AutomorphicArithmetic/AutomorphicVarieties/ModularityArgument/ThresholdArguments/DenseThresholdCases/CuPChainIV.lean}{64}{The five good vertices and the four lines},
\src{src/Development/AutomorphicArithmetic/AutomorphicVarieties/ModularityArgument/ThresholdArguments/DenseThresholdCases/PromodularityCaseThreeOfLemmas.lean}{59}{Common coefficient extension}
and \src{src/Development/AutomorphicArithmetic/AutomorphicVarieties/ModularityArgument/ThresholdArguments/DenseThresholdCases/PromodularityCaseThreeOfLemmas.lean}{117}{Application to the chosen generization}.

The marked ternary polarization alone gives
\(\rho^c\simeq\rho^\vee\); it does not supply the invariant form
or the nontrivial projective involution in
Proposition~\ref{hr:orth}. In degree four
Proposition~\ref{hr:ind} requires a real quadratic self-twist,
and in equal characteristic it also requires the unramified
condition there. In characteristic zero Corollary~\ref{hr:imprimitive}
leaves the strongly irreducible quartic points. In equal
characteristic Proposition~\ref{hr:case3strong} also leaves those
points, together with the possible ramified real quadratic induction
in the second host case. The even two-dimensional block of
Lemma~\ref{spreal:even} and the complementary ternary cases also
remain in the general result.

\subsection{Propagation between components in degrees two and three}

In the ordinary deformation argument the closed set to be propagated
is \(V(J)\), the inverse image of the specified Hecke ideal.
Here are the component and curve arguments with their dimension and
point hypotheses visible.

\begin{lemma}\label{cprop:curve}
Let \((R,\mathfrak n)\) be a Noetherian local ring, \(\pi\in\mathfrak n\),
\(b\ge1\), and \(I,H_1,\ldots,H_t\) ideals. Suppose
\(\dim R/I\ge b+1\) and \(\dim R/H_i<b\) for each \(i\).
There is a prime \(\mathfrak p\supset I+(\pi)\) with
\(\dim R/\mathfrak p=1\) containing none of the \(H_i\).
\end{lemma}
\begin{proof}
In \(R/I\) a system of parameters modulo \(\pi\), together with
\(\pi\), generates an ideal primary to the maximal ideal.
Krull's height theorem gives \(\dim R/(I+(\pi))\ge b\). Take a
minimal prime \(Q\) of this quotient of dimension at least \(b\),
viewed as a prime of \(R\). Each \(H_i\) contains an element
\(a_i\notin Q\), by the strict inequality of dimensions.
Choose \(u\in\mathfrak n\setminus Q\) and put \(z=u\prod_i a_i\).
Choose a prime \(\mathfrak p\) maximal among primes containing \(Q\)
and avoiding \(z\). In the local domain \(D=R/\mathfrak p\), the
image of \(z\) is a nonzero nonunit belonging to every nonzero prime.
If \(\dim D\ge2\), the height-one primes containing \(z\) are the
finitely many minimal primes over \((z)\), by the principal ideal
theorem. Prime avoidance gives a nonunit in the maximal ideal
outside all of them. A minimal prime over that nonzero element
has height one and also contains \(z\), a contradiction. Thus
\(\dim D=1\), and the choice of \(z\) excludes every \(H_i\).
\end{proof}
For the source chart of rank \(n=2,3\) and layer degree parameter
\(d\), take \(b=nd\). The forbidden ideals are the opposite triangular
ideal and, for \(i\ne j\) at each marked \(v\), the sum of the Hecke
ideal, the coefficient maximal ideal and the ideal of all differences
\(\psi_{vi}(g)-\psi_{vj}(g)\) on the specified inertia subgroup.
If the first triangular ideal already generates the unit ideal on
the chart and each forbidden quotient has dimension \(<nd\), the
lemma produces a characteristic-\(\ell\) curve off both triangular
loci, with all pairs of marked characters distinct there. For
the difference ideals this uses that the initial prime contains
the Hecke ideal and that the curve contains \(\pi\).
\src{src/Development/AutomorphicArithmetic/AutomorphicVarieties/CurveDVRSpecialization/NicePrimeEx.lean}{87}{The one-dimensional prime} and
\src{src/Development/AutomorphicArithmetic/AutomorphicVarieties/CurveDVRSpecialization/ExistsNicePrimeAboveOfLargeDim.lean}{7}{the exceptional ideals in the chart}.

\paragraph{The path through a local chart.}
Write \(U\) for the open locus off both triangular ideals and, in
rank three, off the ideal of matrix entries of odd weight difference.
Write \(E_0\) for the part off both triangular ideals where that
third ideal does vanish. In degree two \(E_0\) is empty. Call a prime
\emph{selected} when it is a one-dimensional-quotient prime as just
described (including the separation and characteristic conditions).
It has type \(U\) or \(E_0\). In the next lemma a designated subset
of component primes of a chart ideal is allowed to carry the
exceptional links of the rank-three construction. Set \(D=nd+1\)
and \(d(P)=\dim R/P\).

\begin{lemma}\label{cprop:local}
Let \(J\) define the closed Hecke locus in a Noetherian chart.
Let \(\mathcal P\) be the minimal primes of a specified chart ideal,
including every minimal prime of \(R\) lying in \(U\). Assume each
vertex of \(\mathcal P\) is either in \(U\) or designated exceptional.
Join \(P,Q\in\mathcal P\) when \(\dim R/(P+Q)\ge D\), and assume
this graph is connected. Suppose the following point assertions hold.
\begin{enumerate}
\item Every prime over \(J\) of quotient dimension at least \(D\)
has a selected specialization. Above any designated exceptional
vertex, every prime of that dimension has a selected specialization
of type \(E_0\).
\item At a selected \(p\in U\) with \(J\subset p\), every prime
below \(p\) contains \(J\).
\item Link selected \(E_0\)-primes when they share a designated
exceptional vertex below them. If \(p_0\) has a \(U\)-prime
containing \(J\) below it, every \(U\)-prime below every selected
\(E_0\)-prime linked by a finite path to \(p_0\) contains \(J\).
Paths of length zero are included.
\end{enumerate}
If one \(U\)-prime \(I\) contains \(J\) and \(d(I)\ge D\), every
minimal prime of \(R\) in \(U\) contains \(J\).
\end{lemma}
\begin{proof}
At any selected \(p\) above a \(U\)-prime containing \(J\), all
\(U\)-primes below \(p\) contain \(J\): use assertion 2 or the
length-zero case of assertion 3. Choose a minimal prime \(P_0\)
below \(I\); openness gives \(P_0\in U\). Applying assertion 1
at \(I\) and the preceding observation gives \(J\subset P_0\).

Travel along a path in \(\mathcal P\) from \(P_0\). At a \(U\)-vertex
retain the conclusion that it contains \(J\). At any vertex outside
\(U\) retain a selected \(E_0\)-prime above it reachable from a
selected \(p_0\in E_0\) above some \(U\)-prime containing \(J\).
For an edge \(P,Q\) take a prime \(M\supset P+Q\) with \(d(M)\ge D\),
from a maximal-dimensional component of \(R/(P+Q)\).
If \(Q\in U\) and \(P\in U\), assertion 1 at \(M\supset J\)
and the observation give \(J\subset Q\).
If \(Q\in U\) and \(P\notin U\), select an \(E_0\)-prime above
\(M\) by assertion 1 at the designated vertex \(P\). The retained
path extends to it through \(P\). Assertion 3 gives \(J\subset Q\).
If \(Q\notin U\), select an \(E_0\)-prime above \(M\) using \(Q\).
For \(P\in U\) it starts a path with witness \(P\). For \(P\notin U\)
it extends the retained path through \(P\). This proves the claimed
induction. Every minimal prime in \(U\) is a vertex of the connected graph.
\end{proof}
\src{src/Development/AutomorphicArithmetic/AutomorphicVarieties/CurveDVRSpecialization/LocProp.lean}{17}{Local propagation}
and \src{src/Development/AutomorphicArithmetic/AutomorphicVarieties/CurveDVRSpecialization/PropagationTheoremGL2.lean}{32}{the two prime-transfer hypotheses}.

\begin{lemma}\label{cprop:global}
Let \(X\) be Noetherian, \(U\subset X\) open and \(V(J)\) closed.
Call an irreducible component eligible if it meets \(U\).
Let \(B\subset X\) be connected and let \(T\subset B\) be test points
such that each closed set meeting \(B\) meets \(T\).
Suppose every eligible component meets \(B\), their union contains
a dense subset of \(B\), and at any \(t\in T\), containment in
\(V(J)\) passes between eligible components through \(t\).
If one eligible component lies in \(V(J)\), then \(U\subset V(J)\).
\end{lemma}
\begin{proof}
Form the unions of the eligible components contained in \(V(J)\)
and of the others. Both unions are closed by finiteness. Together
they cover \(B\) by the density assumption. If the second group
exists, both groups meet \(B\), so connectedness gives an intersection
over \(B\). The intersection is closed, hence contains a point of
\(T\). The asserted transfer there contradicts membership of a
component in the second group. Every point of \(U\) lies on an
eligible component, proving the result.
\end{proof}

For clarity, transfer at \(t\) follows from Lemma~\ref{cprop:local}
on a flat local chart if each eligible component through \(t\)
has a chart minimal prime mapping to its generic point, of quotient
dimension at least \(D\). Indeed openness pulls the generic point
into \(U\). The first component supplies a \(J\)-containing prime,
the local lemma puts the other preimage prime over \(J\), and closure
gives the other whole component. For an affine flat local chart
centered at \(t\) the required lifting of generic points follows by
going down from its closed point to each minimal prime under \(t\),
and then taking a minimal prime below the lift. The contraction is
still the same minimal prime. The dimension inequality is additional.
In the global argument \(B\) is the connected locus used to test the
deformations, \(T\) its specified maximal points, and the density
required here is obtained in the source away from the closed
odd-weight-entry locus. When \(n=3\) the exceptional vertices also
use a finite ideal obtained from degree two. If the Hecke ideal \(J\)
works on \(U\) and the exceptional ideal \(I_3\) on \(E_0\),
then \(J\cap I_3\) works on their union. Finiteness of its fibre follows
from the radical test: any prime over a sum of the weight maximal ideal
and \(J\cap I_3\) contains \(J\) or \(I_3\), since \(JI_3\subset J\cap I_3\).
Thus the closed fibre in the Noetherian local ring is zero-dimensional,
hence Artinian. Its composition factors are its residue field, finite
over the coefficient residue field in the ordinary deformation chart.
This gives finite dimension over that coefficient field.
This makes explicit the use of both ranks. It still requires the
dimension and connectedness inputs, a seed on the actual tower and
the two arithmetic point-transfer assertions in the local lemma.
\src{src/Development/AutomorphicArithmetic/AutomorphicVarieties/CurveDVRSpecialization/GlobProp.lean}{3}{Global propagation},
\src{src/Development/AutomorphicArithmetic/AutomorphicVarieties/CurveDVRSpecialization/GlobOfCharts.lean}{4}{passage through flat charts} and
\src{src/Development/AutomorphicArithmetic/AutomorphicVarieties/CurveDVRSpecialization/AbsMain.lean}{12}{the order of the two rank arguments}.

\subsection{The residual Bernoulli test in the higher-rank argument}

The auxiliary character used in the deformation argument is odd. Its
first Bernoulli number can be a unit even when the target value
\(L(-e,\psi)\) vanishes by parity. We give the auxiliary unit with
the precise conductor threshold, and the resulting cohomological
implication in both character directions.

\begin{theorem}[Unit for the auxiliary two-power twist]\label{famsel:unit}
Let \(\ell\) be prime, \(d\ge1\), \(Q=\ell^d\), and
\(a=v_2(Q-1)\ge2\). Let \(f\) be odd, \(\phi\) primitive of conductor
\(f\) and order dividing \(Q-1\), and \(\xi\) primitive of conductor
\(M=2^n\). Take the character values in \(\overline{\mathbf Q}_\ell\).
If \(M\ge f\,4^a\) and \(\phi(-1)\xi(-1)=-1\), then
\(B_{1,\phi\xi}\) is an \(\ell\)-adic unit. This includes \(\ell\mid f\).
\end{theorem}
\begin{proof}
Put \(m=n-a\) and \(s=f2^a\). Then \(m\ge a\ge2\) and
\(1<s\le2^m\). Extend characters by zero at nonunits of their modulus.
For \(u\in\mathbf Z/M\mathbf Z\) set
\[
 A(u)=\sum_{\substack{0\le b<fM\\ b\equiv u\bmod M}} b\phi(b),
 \qquad S=\sum_{b\bmod f}\phi(b),\qquad
 g(u)=\frac{A(u)-A(0)}{fM}.
\]
The sum of the unweighted \(\phi(b)\) in each progression is \(S\),
since \(M\) is coprime to \(f\). If \(t\ge0\), shifting the letters
of this progression by \(ft\) gives
\[
 g(u)-g(u-ft)=tS/M-
       \sum_{\substack{0\le b<ft\\b\equiv u\bmod M}}\phi(b)
       \qquad(0\le ft\le fM).
\]
Indeed the terms which wrap past \(fM\) lose \(fM\) in the weighted
sum, and \(\phi\) has period \(f\). Choose \(0\le t<M\) with
\(ft\equiv u\bmod M\). This formula proves \(g(u)\) integral at \(\ell\)
and \(\overline{g(u)}\in\mathbf F_Q\), since \(g(0)=0\), \(\ell\)
is odd and the nonzero values of \(\phi\) reduce to \(\mathbf F_Q\).
In particular no division by \(f\) is needed on reduction.
Write \([u]\) for the least nonnegative representative. At \(t=2^a\)
we obtain
\[
 \bar g(u)-\bar g(u-s)=c-{\bf1}_{[u]<s}\,\bar\phi([u]),
 \qquad c=\overline{2^a S/M}.
\]
The character \(\xi\) has sum zero. Grouping the defining first
Bernoulli sum by progressions gives
\(B_{1,\phi\xi}=\sum_u\xi(u)g(u)\), as \(\phi\xi\) is nontrivial.

Suppose the reduction of this last sum vanishes. For every unit \(y\)
we also have \(\sum_u\bar\xi(u)\bar g(uy)=0\), by the substitution
\(uy=v\). Apply normalized finite-field trace from
\(\mathbf F_Q(\mu_{2^n})\) to \(\mathbf F_Q\). Its degree is a power
of two, invertible in characteristic \(\ell\). The trace fixes roots
of order dividing \(2^a\) and kills every root of order \(2^j\)
with \(j>a\). Indeed \(Q^{2^{j-a-1}}\equiv1+2^{j-1}\pmod{2^j}\)
by successive squaring of \(Q\equiv1\bmod4\). The corresponding
Frobenius automorphism takes such a root to its negative. Thus
\[
 \sum_{u\in D}\bar\xi(u)\bar g(uy)=0,\qquad
 D=\{u\in(\mathbf Z/M\mathbf Z)^\times:\bar\xi(u)^{2^a}=1\}.
\]
Every \(u\in D\) is congruent to \(1\) or \(-1\) modulo \(2^m\).
In fact the units are \(\{\pm1\}\times\langle5\rangle\) with
\(5\) of order \(2^{n-2}\), by \(v_2(5^{2^j}-1)=j+2\).
Primitivity gives \(\bar\xi(1+M/2)=-1\), so \(\bar\xi\) is faithful
on \(\langle5\rangle\). Writing \(u=\pm5^i\), we have
\(2^{m-2}\mid i\), which gives the claimed congruence.
Let \(\epsilon(u)=1\) on the first congruence class and \(-1\)
on the second. Multiplying by \(1+M/2\) preserves the two classes
and \(D\) and negates \(\bar\xi\). Hence
\(\sum_D\epsilon(u)\bar\xi(u)=0\).

Subtract the identities for \(y=1-s\) and \(y=1\), taking the latter
minus the former. Since \(2^a\mid s\), we have \(u(1-s)=u-s\)
in the positive class and \(u(1-s)=u+s\) in the negative class.
In the first class the only boundary contribution \([u]<s\) occurs
at \(u=1\), where it contributes \(-1\). In the second class use the
shift formula at \(u+s\). The inequality \([u+s]<s\) forces
\([u+s]=s-1\), because \(s\le2^m\) and \(u+s\equiv s-1\bmod 2^m\).
Here \(u=-1\), and its contribution is
\(\bar\xi(-1)\bar\phi(s-1)=-1\). Both exceptional units belong to
\(D\). The resulting equality is
\(0=c\sum_D\epsilon(u)\bar\xi(u)-2=-2\), impossible at odd \(\ell\).
\end{proof}
\src{src/Development/AutomorphicArithmetic/AutomorphicVarieties/LevelCohomology/LxPartial.lean}{15}{progression difference},
\src{src/Development/AutomorphicArithmetic/AutomorphicVarieties/LevelCohomology/LxCore.lean}{65}{two boundary contributions}, and
\src{src/Development/AutomorphicArithmetic/AutomorphicVarieties/LevelCohomology/LemmaX.lean}{78}{exact conductor threshold}.

At the threshold itself take \(\ell=5,\ d=f=1,\ M=16\), and the
primitive character with \(\xi(5)=i,\ \xi(-1)=-1\).
Pairing the residues \(b=1,5,9,13\) with \(16-b\) gives
\[
 16B_{1,\xi}=(2-16)+(10-16)i-(18-16)-(26-16)i=-16(1+i).
\]
The norm of \(B_{1,\xi}\) is \(2\), a unit at either place over five.

\begin{proposition}[The residual classes annihilated by the unit]\label{famsel:stick}
Let \(\eta\ne1\) be an integral Dirichlet character modulo \(D_0\),
of order prime to an odd prime \(\ell\). Set \(\gamma=\bar\eta\)
as a Galois character and suppose it differs from the mod-\(\ell\)
cyclotomic character. If \(B_{1,\eta^{-1}}\) is a unit at \(\ell\),
every continuous cocycle \(z:G_{\mathbf Q}\to k(\gamma)\), for a
discrete residue field \(k\) containing its values, is a coboundary
under the following local condition. On the intersection of each
finite inertia group with \(\ker\gamma\), the cocycle vanishes on a
subgroup of finite index prime to \(\ell\).
In Theorem~\ref{famsel:unit} this applies to both
\(\eta=\phi\xi\) and \(\eta=(\phi\xi)^{-1}\).
\end{proposition}
\begin{proof}
On \(\ker\gamma\) the cocycle is an additive homomorphism, so the
local vanishing extends to the whole inertia intersection: a nonzero
finite image would have order divisible by \(\ell\), dividing the
given index. On \(G_L\), \(L=\mathbf Q(\mu_{D_0})\), it therefore
factors through \(\mathrm{Cl}(L)\) by global class field theory.
Real inertia has order at most two and is killed as well. If
\(\sigma_b(\zeta)=\zeta^b\), conjugating a group element in \(G_L\)
by a lift of \(\sigma_b\) multiplies the restricted homomorphism by
\(\bar\eta(b)\), by the cocycle identity and equivariance of reciprocity.

Choose a positive odd integer \(c\) prime to \(D_0\) with
\(c-\eta(c)\) a unit at \(\ell\). Indeed choose a Galois element on
which \(\gamma\) and the residual cyclotomic character differ and
lift its unit class modulo \(\operatorname{lcm}(2,\ell,D_0)\).
The classical Stickelberger theorem says that
\[
 (c-\sigma_c)\Theta\in\mathbf Z[\operatorname{Gal}(L/\mathbf Q)]
 \quad\text{annihilates }\mathrm{Cl}(L),\qquad
 \Theta=D_0^{-1}\!\!\sum_{b\in(\mathbf Z/D_0\mathbf Z)^\times}
                  [b]\,\sigma_b^{-1}.
\]
This is the integral annihilator for the full cyclotomic field
\cite{stickelberger1890}. Evaluating the integral group-ring element
at \(\eta\) in characteristic zero gives
\((c-\eta(c))B_{1,\eta^{-1}}\), since \(\eta^{-1}\ne1\).
Reduction of this scalar is nonzero even if \(\ell\mid D_0\).
Thus the restricted homomorphism is zero. The cocycle descends to
the abelian group \(\operatorname{Gal}(L/\mathbf Q)\). Choose \(h\)
with \(\gamma(h)\ne1\), possible since reduction is faithful on the
prime-to-\(\ell\) roots. Comparing \(z(hg)\) and \(z(gh)\) gives
\(z(g)=(\gamma(g)-1)z(h)/(\gamma(h)-1)\).

In the application the character and its inverse both have odd parity
and order prime to \(\ell\). Primitivity gives \(\xi(1+M/2)=-1\).
Choose a unit congruent to \(1+M/2\) modulo \(M\) and to \(1\) modulo
\(\ell f\), using coprimality. On a Galois lift both residual characters
equal \(-1\), whereas the cyclotomic character is \(1\).
Theorem~\ref{famsel:unit} applies to the inverse pair of characters
as well. This proves both instances.
\end{proof}
\src{src/Development/AutomorphicArithmetic/AutomorphicVarieties/LevelCohomology/ResidualBernoulliNonvanishing.lean}{174}{residual cocycle criterion},
\src{src/Development/AutomorphicArithmetic/AutomorphicVarieties/LevelCohomology/ChrInvNeg.lean}{197}{both auxiliary characters}, and
\src{src/Development/AutomorphicArithmetic/AutomorphicVarieties/ModularityArgument/ThresholdArguments/DenseThresholdCases/ConductorAtMapInjectiveCoeff.lean}{132}{use in the two higher-rank counts}.
These two statements discharge the residual Bernoulli premise of the
rank-four counting theorem. Its other premises include integral
deformation finiteness, rigidity and an automorphic count. The cocycles
in Proposition~\ref{famsel:stick} have residual coefficients in an
auxiliary odd character, whereas the required lengths count all
congruence depths in the original coefficient module.

\subsection{Integral vanishing for varying nontrivial characters}

\begin{proposition}\label{famsel:class}
Let \(\ell\) be odd, \(O\) a coefficient DVR finite over
\(\mathbf Z_\ell\), with uniformizer \(\pi\), and \(\psi\) a Dirichlet
character of order prime to \(\ell\) valued in \(O^\times\).
Put \(\alpha=\psi^{-1}\chi_\ell^{-e}\),
\(A=(\operatorname{Frac}O/O)(\alpha)\) and
\(K=\overline{\mathbf Q}^{\ker\bar\alpha}\), where the bar denotes
reduction modulo \(\pi\). Suppose \(\bar\alpha(D_\ell)\ne1\).
Put \(s=[\mathbf F_\ell(\bar\alpha(G_{\mathbf Q})):\mathbf F_\ell]\).
If the order of the group
\[
        \operatorname{Cl}(K)/\langle[\mathfrak l]:\mathfrak l\mid\ell\rangle
\]
is not divisible by \(\ell^s\), the entire strict group
\(S_\ell(\mathbf Q,A)\) is zero. Its normalized length at every host
level \(\ell N\) is then zero.
\end{proposition}
\begin{proof}
Each continuous cocycle into the discrete module \(A\) has finite
image, so each class is killed by a power of \(\pi\). Take a strict
class killed by \(\pi\), and subtract a coboundary so that its
representative \(z\) vanishes on \(D_\ell\).
We have \(\pi z(g)=g b-b\) for some \(b\in A\). Thus
\(b\in A^{D_\ell}=0\), since some element acts by a scalar whose
difference from one is a unit. The values of \(z\) lie in \(A[\pi]\),
so its restriction to \(G_K\) is a homomorphism. It vanishes on
every decomposition group above \(\ell\) in \(G_K\), by the cocycle
conjugation formula and the trivial residual action of \(G_K\).
For \(p\ne\ell\), inertia acts trivially through \(\chi_\ell\) and
reduction is injective on the prime-to-\(\ell\) image of \(\psi\).
Thus \(I_p\cap G_K\) acts trivially on the full \(A\). On that
subgroup the restriction of \(z\) is a coboundary by strictness,
hence is zero. Inertia at infinity has order at most two and has no
image in \(A[\pi]\). By global reciprocity this homomorphism
factors through the displayed quotient of the ordinary ideal class
group. Its image in \(A[\pi]\) is stable under every residual scalar,
by the cocycle conjugation formula. It is an additive subgroup of
the residue field, hence a vector space over the field generated by
these scalars. A nonzero image would have order divisible by \(\ell^s\).
Thus it vanishes. Inflation--restriction leaves a class of
\(\operatorname{Gal}(K/\mathbf Q)\) in \(A[\pi]\). This Galois
group has order dividing \(\#(O/\pi)^\times\), so averaging kills
the class. Vanishing of the \(\pi\)-torsion proves the claim at every depth.
\end{proof}

In particular, let the primitive conductor \(f\) be prime to \(\ell\),
let \(\psi\) be even and nontrivial, and take any positive multiple
\(e\) of \(\ell-1\). If \(\psi(\ell)\ne1\), the local hypothesis holds.
The field \(K\) is the totally real field cut out by \(\psi\), independent
of \(e\). Writing \(d=\operatorname{ord}(\psi)\), one explicit sufficient
condition for vanishing, uniform in the character and the twist, is
\[
        \ell^{\operatorname{ord}_d(\ell)}>
              f^{d/2}\bigl(1+\tfrac d2\log f\bigr)^{d-1}.
\]
Indeed \(|\Delta_K|\le f^d\), by the discriminant formula in the tower
\(\mathbf Q(\mu_f)/K/\mathbf Q\) and the elementary cyclotomic bound
\(|\Delta_{\mathbf Q(\mu_f)}|\le f^{\varphi(f)}\).
Minkowski puts a representative of every ideal class at norm at most
\(B=\sqrt{|\Delta_K|}\). At a rational prime \(p\) the coefficients of
the Euler factor for counting ideals are bounded by those of
\((1-p^{-s})^{-d}\), since there are at most \(d\) prime ideals and
their residue degrees are positive. Counting ordered products of \(d\)
positive integers with product at most \(B\) therefore bounds \(h_K\)
by \(B(1+\log B)^{d-1}\), using the harmonic sum \(d-1\) times.
Here the residual character still has order \(d\), so the stated
inequality gives \(h_K<\ell^s\). This extends the odd-order test of
Proposition~\ref{newcoh:oddzero} to either parity of the character
order and accounts for the required splitting above \(\ell\).

Apply the class-group criterion of Proposition~\ref{famsel:class} directly. The quadratic character modulo \(5\) at \(\ell=7\)
gives zero for every \(e=6j\), and either order-three character
modulo \(7\) at \(\ell=5\) gives zero for every \(e=4j\), \(j\ge1\).
The Frobenius residues are \(2\bmod5\) and \(5\not\equiv\pm1\bmod7\).
The first field is \(\mathbf Q(\sqrt5)\), with Minkowski bound
\(\sqrt5/2<2\). The second is \(\mathbf Q(\mu_7)^+\).
The polynomial \(X^3+X^2-2X-1\) of \(\zeta_7+\zeta_7^{-1}\)
has discriminant \(49\), so its Minkowski bound is at most
\(6\cdot7/27<2\). Both class groups are trivial.
Whenever these criteria hold for an admitted host character in case II
or III, they give the required polynomial inequality at all the indicated
twists. Characters whose residual action is locally trivial at \(\ell\),
and those for which the class-group test fails, still enter the counts.
\src{src/Development/AutomorphicArithmetic/AutomorphicVarieties/HeightAndVanishing/Counts.lean}{65}{Case II}
and \src{src/Development/AutomorphicArithmetic/AutomorphicVarieties/HeightAndVanishing/Counts.lean}{67}{case III}.

\begin{proposition}[A host-level class test allowing \(\ell\)-power inertia]
\label{famsel:hostinertia}
Let \(\ell\) be odd, \(N\) a host modulus, and \(\psi\) any character
modulo \(N\) with values in \(O^\times\). Use \(A,\bar\alpha,K,s\)
as in Proposition~\ref{famsel:class}. Let
\[
 T=\{p\mid\ell N:\bar\alpha(D_p)\ne1\}.
\]
Suppose \(\ell\in T\), and every \(p\ne\ell\) with
\(\ell\mid\#\psi(I_p)\) belongs to \(T\).
If \(\ell^s\) does not divide the order of
\(\operatorname{Cl}(K)/\langle[\mathfrak p]:\mathfrak p\mid p,\ p\in T\rangle\),
the full strict group at level \(\ell N\) vanishes.
\end{proposition}
\begin{proof}
Normalize a \(\pi\)-torsion class to a cocycle \(z\) vanishing on
\(D_\ell\) as in Proposition~\ref{famsel:class}. Then
\(z\) takes values in \(A[\pi]\). For \(p\in T\), the host condition
gives \(z|_{D_p}=\delta b_p\) in \(A\). Since \(\pi z=0\),
\(\pi b_p\in A^{D_p}=0\). Thus \(b_p\in A[\pi]\) and
\(z|_{D_p\cap G_K}=0\). At every other \(p\ne\ell\), the order of
\(\psi(I_p)\) is prime to \(\ell\), so reduction is injective on
that group, and \(I_p\cap G_K\) acts trivially on \(A\). Strictness
on inertia again makes \(z\) zero there. The homomorphism
\(z|_{G_K}\) factors through the displayed quotient by reciprocity.
The image-size and inflation--restriction arguments of
Proposition~\ref{famsel:class} apply verbatim, since
\(\operatorname{Gal}(K/\mathbf Q)\) has order prime to \(\ell\)
regardless of the order of the original \(\psi\).
\end{proof}

\src{src/Development/AutomorphicArithmetic/AutomorphicVarieties/CharacterLevelBounds/ShaFamMono.lean}{30}{strict length at level \(\ell N\)} and
\src{src/Development/AutomorphicArithmetic/AutomorphicVarieties/GaloisLibrary/SelmerTransport/ConjStableShaFam.lean}{94}{decomposition conditions at the primes of the level}.

For a fixed exceptional prime, take \(\ell=3\) and either primitive
order-six character \(\psi\) modulo \(7\). Let \(7\mid N\) and inflate
\(\psi\) to \(N\). For every positive odd \(e\) its reduction after
the inverse Tate twist is the product of the quadratic characters of
discriminants \(-7\) and \(-3\): reduction kills the cubic part,
and inversion fixes both quadratic factors. Hence
\(K=\mathbf Q(\sqrt{21})\); its residual character is nontrivial on
inertia at both \(3\) and \(7\). The only prime away from \(3\)
where \(\psi\) has inertia of order divisible by \(3\) is \(7\).
Minkowski puts each ideal class at norm at most \(\sqrt{21}/2<3\).
The integer \((1+\sqrt{21})/2\) has polynomial \(X^2-X-5\),
irreducible modulo \(2\), so \(2\) is inert and no ideal has norm
two. Thus \(h_K=1\), and Proposition~\ref{famsel:hostinertia} gives
zero at every congruence depth. With either sextic character as the
fixed target, its inflation occurs on each of its hosts at all odd
twists, by the defining kernel of the host. For odd \(e\ge3\) these
are case II: \(3\) is exceptional and the prime-to-\(3\) part away
from \(3\) has order two. This includes arbitrarily high twists with
\(3\mid\operatorname{ord}(\psi)\).
\src{src/Development/AutomorphicArithmetic/AutomorphicVarieties/HeightAndVanishing/DirichletPrimeToLPartMulLPowerPart.lean}{179}{host kernel and parity test},
\src{src/Development/AutomorphicArithmetic/AutomorphicVarieties/HeightAndVanishing/DirichletPrimeToLPartMulLPowerPart.lean}{212}{case II}, and
\src{src/Development/AutomorphicArithmetic/AutomorphicVarieties/StmtFields.lean}{66}{the fixed exceptional set}.

\begin{proposition}[Characters of prime-power conductor]\label{famsel:primary}
Let \(\ell\) be odd and suppose \(\ell\nmid h_{\mathbf Q(\zeta_\ell)^+}\).
For every Dirichlet character \(\psi\) of conductor dividing a power
of \(\ell\) and every \(e\ge1\) with \(\psi(-1)=(-1)^e\), the strict
group \(S_\ell(\mathbf Q,A)\) vanishes, also on every host to which
\(\psi\) inflates.
\end{proposition}
\begin{proof}
The full scalar \(\alpha\) factors through the \(\ell\)-power
cyclotomic tower. This tower is totally ramified at \(\ell\), so
\(\alpha(D_\ell)=\alpha(G_{\mathbf Q})\). For a class killed by
\(\pi\), normalize \(z|_{D_\ell}=0\) and write \(\pi z=\delta b\).
Then \(b\in A^{D_\ell}=A^{G_{\mathbf Q}}\), hence \(\pi z=0\);
this works even when \(\bar\alpha=1\).
The prime-to-\(\ell\) character \(\bar\alpha\) factors through
\((\mathbf Z/\ell)^\times\), and parity makes its field \(K\)
a subfield of \(L=\mathbf Q(\zeta_\ell)^+\).
Extension of ideals embeds the \(\ell\)-part of \(\mathrm{Cl}(K)\)
in \(\mathrm{Cl}(L)\), since its composite with norm is raising to
\([L:K]\), a unit at \(\ell\).
The homomorphism \(z|_{G_K}\) is zero on inertia away from \(\ell\)
by strictness and is zero on decomposition groups over \(\ell\)
by normalization and conjugation. Ordinary class field theory now
kills it, since \(\ell\nmid h_K\). Inflation--restriction and
\(\ell\nmid[K:\mathbf Q]\) kill the original class. As before,
vanishing of the \(\pi\)-torsion suffices.
\end{proof}

In fact, for the sextic target modulo \(7\), every power of the
target has zero strict group at both exceptional primes \(3\) and
\(7\), at all twists of its admitted parity. At \(7\) this follows
from Proposition~\ref{famsel:primary} and \(h_{\mathbf Q(\zeta_7)^+}=1\).
At \(3\) the odd powers, including the quadratic character, all
have the residual field \(\mathbf Q(\sqrt{21})\) at odd twists, so
Proposition~\ref{famsel:hostinertia} applies as above. For an even
power and an even twist, \(\bar\alpha=1\). Inertia at \(3\) maps
onto the full \(3\)-adic cyclotomic image, and a Frobenius at \(3\)
generates the image of the unramified character \(\psi\) of order
one or three, since \(3\not\equiv\pm1\bmod7\). Thus
\(\alpha(D_3)=\alpha(G_{\mathbf Q})\), and the normalization just
proved makes \(z\) a homomorphism of exponent dividing three, zero
on \(D_3\). It is unramified away from \(7\) by strictness and
is tame at \(7\) by its exponent. By reciprocity it factors through
the ordinary ray group modulo \(7\) of \(\mathbf Q\), which is
\((\mathbf Z/7)^\times/\{\pm1\}\). The class of \(3\) generates
this group and maps to zero, proving the claim. This treats the
inflated target powers on each host. Other host characters can have
additional conductor factors.

\begin{proposition}[A unit test at a completely split prime]\label{famsel:splitray}
Let \(\psi\) be a nontrivial even character of order prime to an odd
prime \(\ell\). Suppose \(\ell\) splits completely in its character
field \(K=\overline{\mathbf Q}^{\ker\psi}\), and \(\ell\nmid h_K\).
At the embeddings \(w:K\hookrightarrow\mathbf Q_\ell\) put
\[
 q_w(u)=\frac{w(u)^{\ell-1}-1}{\ell}\pmod\ell
 \qquad(u\in\mathcal O_K^\times).
\]
If the vectors \((q_w(u))_w\) span the coordinate-sum-zero hyperplane
over \(\mathbf F_\ell\), then \(S_\ell(\mathbf Q,A)=0\) for every
positive multiple \(e\) of \(\ell-1\), also after imposing any host
conditions. Here \(\bar\alpha\) is trivial on \(D_\ell\).
\end{proposition}
\begin{proof}
A strict class killed by \(\pi\) lifts to a cocycle \(z\) in
\(A[\pi]\) by the exact sequence for multiplication by \(\pi\);
no local condition on this lift at \(\ell\) is needed. The residual
action is \(\bar\psi^{-1}\). On \(G_K\) the lift is a homomorphism
unramified away from \(\ell\), since every such inertia group there
acts trivially on \(A\). Its image has exponent \(\ell\). At each
\(w\mid\ell\), local reciprocity therefore kills
\(1+\ell^2\mathbf Z_\ell=(1+\ell\mathbf Z_\ell)^\ell\), by the
\(\ell\)-adic logarithm. Real inertia has order two. Thus the
homomorphism factors through the ray class group of modulus
\(\ell^2\mathcal O_K\) modulo \(\ell\)-th powers.

The ray class exact sequence and \(\ell\nmid h_K\) identify this
quotient with a quotient of
\((\mathcal O_K/\ell^2)^\times/
((\mathcal O_K/\ell^2)^\times)^\ell\simeq\bigoplus_w\mathbf F_\ell\)
by the unit vectors above. The displayed isomorphism is the
coordinatewise Fermat quotient. Its kernel is the subgroup of
\(\ell\)-th powers, and it is surjective on each factor, as is seen
on \(1+\ell a\). The vectors sum to zero since the norm of a unit is
\(\pm1\). Under the span hypothesis the ray class quotient has trivial
Galois action: the coordinate-sum map is invariant under permutation.
The restriction of \(z\) is equivariant with target of nontrivial
scalar character \(\bar\psi^{-1}\), hence must vanish. Finally
\([K:\mathbf Q]=\operatorname{ord}(\psi)\) is invertible on \(A[\pi]\),
so inflation--restriction kills the class there and in \(A\).
Every cohomology class with coefficients in \(A\) is killed by a power of \(\pi\), which proves the assertion.
For the ray class sequence see
\href{https://math.mit.edu/classes/18.785/2021fa/LectureNotes21.pdf\#page=5}{Sutherland, Theorem 21.8}.
\end{proof}

\paragraph{A full case III progression, including the split local action.}
For either cubic character modulo \(7\) and \(\ell=29\), the strict
group vanishes for \emph{every} positive \(e\) with \(28\mid2e\).
First take \(28\mid e\). The field \(K=\mathbf Q(\zeta_7)^+\) has
class number one by the Minkowski computation above, and \(29\equiv1\bmod7\)
splits completely. The unit \(u=\zeta_7+\zeta_7^{-1}\) has polynomial
\(X^3+X^2-2X-1\); \(u^2-2\) is a conjugate. Modulo \(29^2=841\)
the three roots, ordered cyclically under \(r\mapsto r^2-2\), are
\(409,761,511\). Direct exponentiation gives
\[
 (q_w(u))_w=(20,16,22),\qquad (q_w(u^2-2))_w=(16,22,20).
\]
Both vectors have coordinate sum \(0\bmod29\), and their first
two columns have determinant \(20\cdot22-16^2\equiv10\bmod29\).
Proposition~\ref{famsel:splitray} applies at every multiple of \(28\).

For \(e\equiv14\bmod28\), Proposition~\ref{newcoh:oddzero}(2)
reduces the question to \(29^2\nmid h_{K(\sqrt{29})}\), since
\(\operatorname{ord}_3(29)=2\). The compositum is totally real of
degree six and discriminant at most \(49^2\,29^3\), by the order
obtained by tensoring integral bases. Its Minkowski bound is
at most \((6!/6^6)49\,29^{3/2}<120\).
The primes \(2,3,5\) all have degree three in \(K\), by their
classes modulo \(7\) up to sign. Also \(2\) is inert in
\(\mathbf Q(\sqrt{29})\), so its degree in the compositum is six.
Consequently no prime ideal has norm below seven and the only ones
of norm at most ten lie over \(7\). There are at most two of these,
since \(7\) is totally ramified in \(K\).
Every ideal of norm \(<120\) is the unit ideal, a prime ideal, or a
product of two prime ideals. There are at most \(6\cdot60\) prime
ideals of norm \(<120\), since at most 60 rational primes are below
120. In a product the smaller norm is at most ten and the other is
at most \(120/7<18\), giving at most \(2\cdot6\cdot17\) products.
Hence \(h_{K(\sqrt{29})}\le1+360+204<29^2\), as required.

Taking a cubic character modulo \(7\) as the target gives
\(p_K=11,\ p_1=13\) in the exceptional set defined below. The class
number of \(\mathbf Q(\sqrt{-11})\) is one: its Minkowski bound
\((2/\pi)\sqrt{11}<3\) and \(2\) is inert. Thus \(29\notin B_\chi\).
The characters inflate to every host for this target. The two
computations exhaust the twists in case III at this prime, uniformly
in their size and in the host.

\begin{corollary}[Case III on two families of targets]\label{famsel:hostdensity}
For every target of two-power order, including the zeta target,
case III is empty on every one of its hosts. For either primitive
cubic target modulo \(7\), there is an explicit set of coefficient
primes of relative natural density \(1/2\) among primes on which
every case III strict group vanishes for all sufficiently large
primes in the set, simultaneously in the host and the twist.
\end{corollary}
\begin{proof}
In the host group of Proposition~\ref{host:art}, the kernel of
the projection to the cyclic target image is a two-group. Every
odd-order character is therefore inflated from that cyclic image.
If the target order is a power of two, only the trivial character
is possible, whereas case III requires a nontrivial one. This
applies in particular to \(\psi=1\). There are still case II tests
on the zeta hosts: for \(\ell=3\) and each odd \(e\ge3\) in the
window take \(J_d\), where \(d>3\) is the auxiliary prime.
It is a host character, is odd since \(d\equiv3\bmod4\), and its
prime-to-three part away from three has order two.

For a cubic target the only nontrivial odd-order host characters
are the two cubic characters modulo \(7\). Their common field is
\(K=\mathbf Q(\zeta_7)^+\), of class number one. Take primes
\[
    \ell\bmod 7\in\{2,3,4,5\},\qquad
    \ell\bmod 12\in\{5,7,11\}.
\]
Both characters are nontrivial at \(\ell\). If \(\ell-1\mid e\),
Proposition~\ref{newcoh:oddzero}(1) gives vanishing. For
\(\ell\equiv3\bmod4\), admissible \(e\) is even, so
\(\ell-1\mid 2e\) already gives \(\ell-1\mid e\).
In the only remaining residue class, \(\ell\equiv5\bmod12\),
the order of \(\ell\) modulo three is \(2>3/2\).
Proposition~\ref{newcoh:oddzero}(2) gives vanishing for the other
twists for all sufficiently large such \(\ell\), with a threshold
depending only on \(K\). Imposing the further host local conditions
preserves zero. There are \(4\cdot3=12\) such reduced classes
out of \(\varphi(84)=24\) by the Chinese remainder theorem.
The prime number theorem for this fixed modulus gives their
relative natural density. Omitting the finitely many exceptional
or small primes preserves that density.
\end{proof}

\subsection{An effective bound simultaneous in the conductor, twist and prime}

\begin{theorem}\label{famsel:rate}
For \(N,e\ge1\), sum the normalized strict lengths of
Proposition~\ref{newcoh:fullsupport} over every prime and any set of
distinct characters of modulus \(N\) at each prime. The coefficient
fields may vary. With \(P=16Ne\) one has
\[
 \sum_{\ell,\psi}\frac{\log\ell}{e_{\ell,\psi}}
       \operatorname{length}_{O_{\ell,\psi}}S_{\ell,\psi}
       \ \le\ \exp(P^{102}).
\]
\end{theorem}
\begin{proof}
Put \(M=4N,\ F=\mathbf Q(\mu_M),\ d=[F:\mathbf Q],\ j=2e\)
and \(H=\max(d+1,2j+2)\). The group \(K_j(\mathcal O_F)\) is finite.
For its primary parts with \(\ell>H\), Bayer-Fluckiger, Emery
and Houriet's Theorem~1.1~\cite{bayerEmeryHouriet} bounds the sum of
their log orders by
\[
 B=(2j+1)^{71j^4d^3}\, d^{293j^5d^5}\,
             |\operatorname{disc}F|^{528j^5d^4}.
\]
Indeed the subgroup generated by elements of order at most \(H\)
has zero projection to these primary parts. We only use the
large-prime consequence of their theorem.

For \(\ell\le H\), let \(F'=F(\mu_\ell)\).
Its degree over \(F\) divides \(\ell-1\). The extension of the two
integer schemes after inverting \(\ell\) is finite \'{e}tale:
adjoining \(\mu_\ell\) is unramified outside \(\ell\).
On \'{e}tale \(H^2\) with the Tate twist \(e+1\), trace composed
with pullback is multiplication by \([F':F]\). It is invertible
at \(\ell\), by the finite-level trace identity and passage to the
inverse limit. Via the even Quillen--Lichtenbaum comparison in
Theorem~\ref{newcoh:lemmaK}, the primary group over \(F\) thus embeds
in the one over \(F'\). The proof of that theorem bounds the order of
the entire primary \(K\)-group, by \eqref{newcoh:transfer} and its
subsequent estimates. The hypotheses on roots of unity hold over
\(F'\) also at two, since \(4\mid M\). We have \(d'=[F':\mathbf Q]\le dH\)
and \(|\operatorname{disc}F'|\le(MH)^{d'}\).
Indeed \(F'\) has conductor dividing \(\operatorname{lcm}(M,\ell)\);
the derivative of \(X^c-1\) at a primitive \(c\)-th root is divisible
in the integers by \(\Phi_c'\) there. Taking norms gives the usual
discriminant bound \(c^{\varphi(c)}\) for the cyclotomic field,
using \(\mathbf Z[\zeta_c]=\mathcal O_{\mathbf Q(\mu_c)}\).
It follows that
\[
 \log\#K_j(\mathcal O_F)
 \le B+H\exp\!\left((jdH)^{32}(1+dH)\log(2MH)\right).
\]
Here \(d,j,H,M\le P\), \(2j+1\le P\), and
\(\log|\operatorname{disc}F|\le d\log M\). Consequently
\(\log B\le 892 P^{10}\log P\). The exponent in the second term
is at most \(6P^{98}\log P\), using \(P\ge16\).
Each term on the right is at most \(\exp(P^{100})\), including
the factor \(H\). Thus \(\log\#K_j(\mathcal O_F)\le\exp(P^{101})\).
Proposition~\ref{newcoh:fullsupport} now bounds the required sum by
\(\varphi(N)(\exp(P^{101})+d\log(4Nde))\le\exp(P^{102})\).
\end{proof}

\begin{corollary}[Part of the window budget]\label{famsel:subwindow}
Fix \(\varepsilon>0\). At a common host level \(N\), retain, for each
positive integer twist \(e\), any distinct characters of primitive
conductor \(f\) satisfying \(16fe\le x^{\varepsilon/204}\).
At all sufficiently large \(x\), their total normalized strict length
at level \(\ell N\), summed over the primes, twists and characters,
is at most \(\exp(x^\varepsilon/2)\). The threshold depends only on
\(\varepsilon\), independently of the host and coefficient fields.
\end{corollary}
\begin{proof}
At fixed \(N\) an inducer modulo \(f\mid N\) has a unique inflation.
A strict class at level \(\ell N\) satisfies in particular all the
conditions at level \(\ell f\). Group by \(f,e\) and apply
Theorem~\ref{famsel:rate} at \(f\). Each pair contributes at most
\(\exp(x^{\varepsilon/2})\). Bounding each coordinate separately gives
at most \(x^{\varepsilon/102}\) pairs. The sum is therefore at most
\(x^{\varepsilon/102}\exp(x^{\varepsilon/2})\), whose logarithm
\(x^{\varepsilon/2}+(\varepsilon/102)\log x\) is at most
\(x^\varepsilon/2\) for sufficiently large \(x\).
\end{proof}

\subsection{The complementary twists and the field budget}\label{famsel:complement}
Use \(v\) for the page scale, \(t=\lfloor v^{401}\rfloor\), and \(N\)
for its host conductor. The scale of the final comparison is
\(x=\kappa v\) with \(\kappa\) fixed. Already \(f=1,\ e=\lfloor v\rfloor\)
eventually violates \(16fe\le v^{\varepsilon/204}\) in the active range
\(0<\varepsilon<1/52\). We record which additional estimates enter
at that same host.

Fix the primitive target \(\chi\), of order \(o\) and conductor \(m\).
Let \(p_K\) be the least prime \(p\ge7\), \(p\equiv3\bmod4\), prime
to \(m\), let \(h_* = h_{\mathbf Q(\sqrt{-p_K})}\), and let \(p_1\)
be the least prime congruent to \(1\) modulo the odd part of \(o\)
and prime to \(2mp_K\). The fixed exceptional set is
\[
 B_\chi=\{3\}\cup\{\ell\text{ odd prime}:\ell\mid om p_K h_*p_1\}.
\]
At odd \(\ell\) use the host characters \(\psi\) with
\(\psi(-1)=(-1)^e\), \(1\le e<t\). Case II requires \(e\ge2\), even
order of the prime-to-\(\ell\) part of \(\psi\) with its local
\(\ell\)-component removed, and
\(\ell\in B_\chi\) or \(\ell-1\mid2e\). Case III requires
\(\ell\notin B_\chi,\ \ell-1\mid2e\), and nontrivial \(\psi\) of
odd order. On every host whose two auxiliary primes exceed \(2t+1\)
and \(m,o,p_K,h_*,p_1\), the imported counts give a finite normalized
length \(w_{\ell,e,\psi}=(\log\ell)\operatorname{length}_O
S_{\ell N}(\mathbf Q,A)/e(O/\mathbf Z_\ell)\) at most
\((\ell N e)^{C}\) in both cases. One exponent \(C\), the larger of
the two source exponents, works for every such host and choice of
coefficient rings for this target.
\src{src/Development/AutomorphicArithmetic/AutomorphicVarieties/StmtFields.lean}{59}{auxiliaries and exceptional primes},
\src{src/Development/AutomorphicArithmetic/AutomorphicVarieties/HeightAndVanishing/DirichletPrimeToLPartMulLPowerPart.lean}{185}{host size and six cases}, and
\src{src/Development/AutomorphicArithmetic/AutomorphicVarieties/HeightAndVanishing/Counts.lean}{60}{uniform character bounds}.

Let \(\mathcal S\) consist of all character entries with
\(16fe\le v^{\varepsilon/204}\), and \(\mathcal R\) of the cases II
and III outside \(\mathcal S\), deleting any entries proved zero
above. For \(L=\max(2t,\max B_\chi)\), every entry of \(\mathcal R\)
has \(\ell\le L\), since \(\ell-1\mid2e\) implies \(\ell\le2t-1\).
There are at most \(N\) characters per prime and twist. Applying the
explicit estimate in the proof of Corollary~\ref{famsel:subwindow}
only to the first sum therefore gives
\[
 \sum_{\mathcal S}w_{\ell,e,\psi}+\sum_{\mathcal R}w_{\ell,e,\psi}
 \le v^{\varepsilon/102}\exp(v^{\varepsilon/2})
       + LNt\,(LNt)^{C}.
\]
On the admitted page \(\log(Nt)\le 2v^{\varepsilon/4}\).
Since \(t\le v^{401}\), the logarithm of each summand is
\(o(v^\varepsilon)\), uniformly on these hosts. The sum fits in
\(\exp(v^\varepsilon/2)\) for sufficiently large \(v\).
Only the second summand uses the imported bounds in cases II and III.
\src{src/Development/AutomorphicArithmetic/AutomorphicVarieties/HeightAndVanishing/SumOut.lean}{167}{related finite-prime count outside the exceptional set} and
\src{src/Development/AutomorphicArithmetic/AutomorphicVarieties/CMFieldExponentialBounds/LogConductorMulHeightLeOfPageN.lean}{6}{uniform conductor bound}.

For clarity, the full sum over the CM field \(E\) has further terms.
Away from \(B_\chi\) and two the remaining cases are I
(\(\ell-1\nmid2e\)) and Z (the trivial character with
\(\ell-1\mid2e\)). For I the imported estimate sums over any finite
set of its primes \emph{before} bounding by \((Nt)^{C_1}\) for each
\(e\), once both auxiliary primes exceed two additional test primes
depending only on \(\chi\). Proposition~\ref{newcoh:zero} gives Z.
Inside \(B_\chi\), case IV handles \(e\ge2\) with the indicated
prime-to-\(\ell\) part of odd order via the fixed cyclotomic layer,
and the two estimates for case V handle \(e=1\).
Using the field-to-character comparison and the fixed-layer bounds,
the source obtains the odd-prime plus contribution
\((Nt)^{C'}+C'B\), where \(B=\exp(v^{4\varepsilon/5})\).
The same window bounds the \(2\)-part by \(B\log2\); the odd-prime
minus estimate is at most \(K\exp(8v^{\varepsilon/4})\).
Consequently for every finite prime set \(S\), on that field,
\[
 \begin{aligned}
 &\sum_{\ell\in S}\log\ell\sum_{1\le e<t}
      \operatorname{length}_{\mathbf Z_\ell}
      S_\ell(E,(\mathbf Q_\ell/\mathbf Z_\ell)(-e))\\
 &\qquad\le B\log2+K\exp(8v^{\varepsilon/4})+(Nt)^{C'}+C'B
       \le \exp(v^\varepsilon/2).
 \end{aligned}
\]
Here \(C',K\) are independent of \(v\) and of the admitted host for
fixed \(\chi\). The last inequality follows by taking logarithms
and using the bound on \(\log(Nt)\). The comparisons on \(E\), the
case I sum and the remaining nontrivial character bounds are used at
the indicated source theorems. The preceding character inequality
identifies the part of their input supplied by the independent estimates.
\src{src/Development/AutomorphicArithmetic/AutomorphicVarieties/HeightAndVanishing/Counts.lean}{53}{case I with both test primes},
\src{src/Development/AutomorphicArithmetic/AutomorphicVarieties/HeightAndVanishing/LayerLe.lean}{16}{fixed layer and field comparisons},
\src{src/Development/AutomorphicArithmetic/AutomorphicVarieties/HeightAndVanishing/JointWinOfCounts.lean}{16}{assembly of the counts}, and
\src{src/Development/AutomorphicArithmetic/AutomorphicVarieties/CMFieldExponentialBounds/AdmissibleOfPageN.lean}{110}{dyadic, minus and plus contributions on one field}.

\subsection{The dyadic field contribution and characters of small order}
\label{adds:field-budget}

Fix a primitive target \(\psi\), write \(K\) for its abelian
field, and put \(t(v)=\lfloor v^{401}\rfloor\). For a prime \(p\)
and a number field \(L\) write
\[
 h_p(L,e)=\operatorname{length}_{\mathbf Z_p}
 S_p(L,(\mathbf Q_p/\mathbf Z_p)(-e)).
\]
The strict conditions here are those of Section~\ref{newcoh:section}.
They are invariant under changes of algebraic closure and prolongation.

\begin{theorem}[One window for the dyadic field and the fixed odd layers]
\label{adds:dyadic-window}
Fix \(0<\varepsilon<1/52\) and a nonnegative integer \(s\).
At finitely many fixed odd primes \(p\), prescribe abelian base fields
\(B_p/\mathbf Q\) containing \(\mu_p\), and fixed nonnegative layer
indices \(m_p\). Write \(B_{p,m_p}\) for the corresponding degree
\(p^{m_p}\) cyclotomic layer. There are \(J\ge1\) and \(X_0\),
depending only on these fixed choices and the target, with the
following property. For \(X\ge X_0\) put \(H=X^{1/J}\). Take any
distinct primes \(q\equiv3\bmod8,\ d\equiv7\bmod16\), coprime to
the target conductor, with \(\log(qd)\le H^{\varepsilon/4}\).
Allow any additional set \(\mathcal E_*\) of at most \(s\) positive
integers, chosen even after \(q,d\).

There is a \(W\ge H\), \(W^2\le X\), such that for every
\(v\in[W,W^2]\) and every page index \(r\) with
\(v^{2+8\varepsilon}\le2^r\le2v^{2+8\varepsilon}\), the actual field
\[
 E=K\mathbf Q(\mu_{2^{r+2}})^+(\sqrt{-q},\sqrt{-d})
\]
of Proposition~\ref{host:art} satisfies
\[
 \sum_{1\le e<t(v)} h_2(E,e)\le \exp(v^{4\varepsilon/5}),
 \qquad
 \sum_{1\le e<t(v)}h_p(B_{p,m_p},e)\le\exp(v^{4\varepsilon/5})
 \quad\text{for each prescribed }p .
\]
All these lengths are finite. On the same interval every
\(e\in\mathcal E_*\) with \(e\le v^{401}\) satisfies
\(e\le v^{\varepsilon/64}\).
\end{theorem}
\begin{proof}
Put \(F_0=K(i,\sqrt{-q},\sqrt{-d})\) and
\(F_\infty=F_0(\mu_{2^\infty})\). The moving-field argument in
Section~\ref{newcoh:section} bounds the number of \(\mathbf Z_2\)
generators of the split unramified module by
\(R_2=4(g_K+6[K:\mathbf Q])\), uniformly in \(q,d\).
Indeed the two primes have respectively \(2\) and \(4\) places in
\(\mathbf Q(\mu_{2^\infty})\), since
\(v_2(q^2-1)=3,\ v_2(d^2-1)=4\).
The fixed base change multiplies these numbers by at most
\([K:\mathbf Q]\). Each newly allowed inertia group contributes
at most one mod-two character, and passing to each quadratic
adjunction multiplies the dimension bound by at most two.
At each prescribed odd prime choose a generator bound \(R_p\ge1\)
for the split unramified module over \(B_p(\mu_{p^\infty})\).
These bounds follow from the fixed abelian \(\mu=0\) theorem and the
class-field argument in the same section~\cite{ferreroWashington}.

At any of these primes let \(j_p=\lfloor\log_p R_p\rfloor+3\).
For \(p=2\), the cyclotomic image over \(F_0\) is
\(1+2^c\mathbf Z_2\) for some \(c\ge2\), since \(i\in F_0\).
It follows that \(\mathbf Q(\mu_{2^c})\subset F_0\) and
\(2^{c-1}\le[F_0:\mathbf Q]\le8[K:\mathbf Q]\).
The layer \(F_m=F_0(\mu_{2^{c+m}})\) has degree \(2^m\) over \(F_0\).
In particular \(E(i)=F_m\) with \(m=\max(0,r+2-c)=O(1+\log v)\).
The field \(F_{j_2}\) has bounded degree and lies in the cyclotomic
field of modulus \(2^{c+j_2}qd\operatorname{cond}(K)\). The
cyclotomic discriminant bound and the tower discriminant formula give
\[
 \log(2|\operatorname{disc}F_{j_2}|)\ll_K 1+\log(qd)\ll_K v^{\varepsilon/4}
 \qquad(H\le v\le X).
\]
At the odd primes the layers \(B_{p,j_p}\) are fixed.

Set \(R=R_2+\sum_p R_p+s,\ \eta=\varepsilon/64\), where the sum
here ranges over the prescribed odd primes. Take \(J=4(802/\eta)^R\).
Apply Proposition~\ref{newcoh:centres} in each tower with
\(Z=\lfloor X^{401}\rfloor\). For each prime it gives at most \(R_p\)
exceptional positive indices in \([1,Z]\). Apply the finite-window
argument of Section~\ref{sup:section} to the union of all these
sets and \(\mathcal E_*\). With this \(J\) it gives one \(W\) as in
the statement for which every member of the union at most \(v^{401}\)
is at most \(v^\eta\).

By Lemma~\ref{newcoh:desc} the tower lengths are exactly the
coinvariant lengths in Proposition~\ref{newcoh:centres}.
A nonexceptional \(e<t(v)\) at its prime consequently costs
\(O(1+\log v)\) at the required layer, since
\(\log Z\le401J\log v\). The constants are uniform in the pair.
For an exceptional \(e<t(v)\), apply Theorem~\ref{newcoh:lemmaK}
at the low layer \(j_p\). It is totally imaginary and contains the
required roots of unity. By Poitou--Tate the dual of its strict
group embeds in the indicated \(H^2\). Its length is therefore at
most \(2\exp(Cv^{3\varepsilon/4})\): the degree is bounded,
\(e^{32}\le v^{\varepsilon/2}\), and the log discriminant costs
\(O(v^{\varepsilon/4})\). The growth inequality
\eqref{newcoh:growth} adds at most \(R_p(m-j_p)_+\) on passing to
the required layer, with \(m=m_p\) at the odd primes.

For descent from \(E(i)\) to the actual field \(E\), restriction
preserves strictness. If the extension is quadratic, its
restriction kernel on \(H^1\) is
\(H^1(C_2,A^{G_{E(i)}})\), where \(A=(\mathbf Q_2/\mathbf Z_2)(-e)\).
The invariants are finite cyclic since the cyclotomic image is open
and \(e>0\). Evaluation at the nonidentity element makes this
cohomology a subquotient of a cyclic group, and restriction followed
by corestriction shows it is killed by two. Thus the loss in length
is at most one per twist. There is zero loss if \(E(i)=E\).
Summing the preceding estimates for fewer than \(v^{401}\) twists
gives the claimed bounds, since \(3\varepsilon/4<4\varepsilon/5\).
\end{proof}

For the prime selection in Theorem~\ref{pr:pair}, first fix this
\(J\), then use initial scale \(H\) and interval exponent \(J\).
The selected pair has
\(\log(qd)\le H^{\varepsilon/5}(1+\log^2 H)\le H^{\varepsilon/4}\)
eventually. Its discarded-character bounds hold on \([H,X]\),
which contains the entire selected interval. The rescaling in
Section~\ref{sup:section} uses the same pair. In particular
the dyadic term in Section~\ref{famsel:complement} costs
\((\log 2)\exp(v^{4\varepsilon/5})\), and the fixed odd layers
may be chosen to be those for the exceptional primes of that section.
Corollary~\ref{adds:fixed-char} uses them for every Case IV character.
Proposition~\ref{addo:compare} controls passage to the field at odd primes.

\begin{corollary}[The dyadic restriction on a zeta host]
\label{adds:zeta-injection}
For the trivial target one can take \(R_2=24,\ j_2=7\).
Restriction of strict groups from \(E\) to \(E(i)\) is injective
at every positive twist.
\end{corollary}
\begin{proof}
Every \(\mathbf Q(\mu_{2^a})\), \(a\ge2\), has odd class number.
Start with the Gaussian integers. In a quadratic step \(L/K'\) of
this tower the unique dyadic prime is totally ramified. If \(L\) has
a nonzero unramified quadratic character then its space of such
characters has a nonzero invariant under the involution, since
\((\sigma-1)^2=0\) in characteristic two. Its fixed character defines
an unramified quadratic extension of \(L\) which is Galois of degree
four over \(K'\). This group is abelian and the inertia at the unique
dyadic prime has order two. Its fixed field gives an everywhere
unramified quadratic extension of \(K'\), contradicting the induction
hypothesis. The finite-layer descent in Section~\ref{newcoh:section}
then gives \(g_{\mathbf Q}=0\).

Globally \(\mathbf Q(\mu_{2^\infty})\cap
\mathbf Q(i,\sqrt{-q},\sqrt{-d})=\mathbf Q(i)\):
inertia at \(q\) and \(d\) flips the corresponding square root
independently, and the cyclotomic two-tower is unramified there.
Locally \(-q\equiv5\bmod8\) gives the unramified quadratic extension
of \(\mathbf Q_2\), while \(-d\equiv1\bmod8\) is a square.
Thus, at any dyadic place \(w\), the completion of \(E(i)\) is the
compositum of \(\mathbf Q_2(\mu_{2^{r+2}})\) and this unramified
quadratic field. The completion of \(E\) uses the subfield fixed
by inversion in the same cyclotomic field. Their local extension
has degree two. Since each higher cyclotomic layer is totally
ramified, the cyclotomic images of \(G_{E(i)}\) and
\(G_{E(i)_w}\) are both \(1+2^{r+2}\mathbf Z_2\).
The restriction kernels in global and local degree one are
therefore inflations from the very same \(C_2\) acting on the
same invariant subgroup of \(A\). Local inflation is injective.
A strict global class restricts to zero at the dyadic place,
so its restriction to \(E(i)\) detects the class.
\end{proof}

\begin{proposition}[An all-prime budget in terms of character order]
\label{adds:order}
On any sufficiently late page of Section~\ref{famsel:complement},
sum over distinct host characters \(\xi\) at each coefficient
prime \(\ell\), with \(1\le e<t(v)\). Write
\(a=\operatorname{ord}(\xi)\) and \(U=ea\). Include every entry
satisfying at least one of
\[
 \begin{gathered}
 U\le v^{\varepsilon/128},\qquad \ell U\le v^{\varepsilon/64},
 \qquad
 U\le v^{\varepsilon/10}\ \text{ and }\
                  \ell>\max(2a+1,4e+2).
 \end{gathered}
\]
For every finite choice of coefficient primes the sum of their
normalized strict weights \(w_{\ell,e,\xi}\) is at most
\(\exp(v^{4\varepsilon/5})\). The estimate is uniform in the
primitive conductors of the characters on that page.
\end{proposition}
\begin{proof}
For a fixed kernel of a host character let \(K_\xi\) be its cyclic
field of degree \(a=\operatorname{ord}(\xi)\). Set \(F=K_\xi(i)\),
\(n=[F:\mathbf Q]\le2a\), and \(M=4N\), so that
\(F\subset\mathbf Q(\mu_M)\). Grouping characters by kernel is
independent of the coefficient prime. For each of them, restriction
to \(F\), Poitou--Tate, and the even Quillen--Lichtenbaum comparison
in Proposition~\ref{newcoh:fullsupport} bound the normalized log
order of the image by \(\log\# K_{2e}(\mathcal O_F)\{\ell\}\).
The restriction kernel vanishes for \(\ell\nmid n\). Otherwise
put \(u_\ell=\max(1,v_\ell(M))+v_\ell(e)\).
Since \(F\subset\mathbf Q(\mu_M)\), its cyclotomic image contains
\(1+\ell^{\max(1,v_\ell(M))}\) as a unit (at two \(4\mid M\)).
The coefficient invariants over \(F\) have normalized log order
at most \(u_\ell\log\ell\). Inflation--restriction and the bound
by all functions on a group of order \(n\) cost at most
\(n u_\ell\log\ell\) for the kernel. Summing this cost uses
\(\sum_{\ell\mid n}u_\ell\log\ell\le\log(Mne)\).
At fixed character kernel there are
at most \(\varphi(a)\) characters at any \(\ell\). Summing yields
\[
 \sum_{\ell,\ \ker\xi\ {\rm fixed}} w_{\ell,e,\xi}
 \le\varphi(a)
       \bigl(\log\#K_{2e}(\mathcal O_F)+n\log(Mne)\bigr).
\]

For clarity, the degree-sensitive estimate for the \(K\)-group
follows from the two parts of the proof of Theorem~\ref{famsel:rate}.
Put \(j=2e,\ H_0=\max(n+1,2j+2)\). For \(\ell>H_0\),
\cite[Theorem 1.1]{bayerEmeryHouriet} bounds the sum of the
primary log orders by
\[
 B=(2j+1)^{71j^4 n^3}\,n^{293j^5 n^5}
          |\operatorname{disc}F|^{528j^5 n^4}.
\]
For \(\ell\le H_0\), adjoining \(\mu_\ell\) costs degree at
most \(H_0\) prime to \(\ell\) at that individual prime. Pullback
on the primary \(K\)-group is injective by the trace identity in
the same proof. The new field has degree at most \(nH_0\),
conductor at most \(MH_0\), and contains the required roots of
unity. The bound for the full primary \(K\)-group in the proof
of Theorem~\ref{newcoh:lemmaK} gives
\[
 \log\#K_{2e}(\mathcal O_F)
 \le B+H_0\exp\bigl((jnH_0)^{32}(1+nH_0)\log(2MH_0)\bigr).
\]
These bounds apply to the totally imaginary field \(F\) of
degree at least two. Write \(U=ea\ge1\). Then \(jn\le4U\),
\(H_0\le6U\), \(nH_0\le12U^2\) and
\(\log|\operatorname{disc}F|\le n\log M\).
Taking logarithms of the positive bounds on the right gives
\[
 \log\bigl(1+\log\#K_{2e}(\mathcal O_F)\bigr)
       \ll U^{66}\log(2MU)
\]
with an absolute constant. On a page, \(\log(Nt)\le2v^{\varepsilon/4}\).
Thus for \(U\le v^{\varepsilon/128}\) the last display is
\(O(v^{49\varepsilon/64})\). The sum of the multiplicities
\(\varphi(a)\) over the cyclic quotients that occur is at most
the number of host characters, at most
\(4[K:\mathbf Q]\,2^r\ll_\psi v^4\).
In the second range apply the same restriction bound separately
at \(\ell\) with \(F'=F(\mu_\ell)\). The exponent in the
small-prime estimate is now at most
\(C(\ell U)^{33}\log(2M\ell U)=O(v^{49\varepsilon/64})\):
here \(j[F':\mathbf Q]\le4\ell U\) and
\([F':\mathbf Q]\le2\ell U\). There are at most
\(v^{\varepsilon/64}\) primes to sum in this range. The restriction
kernel cost \(n\log(Mne)\) per character is also absorbed.
In the third range \(\ell>H_0\) and \(\ell\nmid n\), so the sum over
these primes costs at most \(\varphi(a)B\) for the kernel orbit.
In every factor defining \(B\), including the discriminant term,
\(\log B\ll U^5\log(2MU)=O(v^{3\varepsilon/4})\).
Multiplication by the polynomial counts of characters, twists and,
in the second range, primes is absorbed in \(\exp(v^{4\varepsilon/5})\),
since \(49/64<4/5\).
\end{proof}

At the zeta target this includes every quadratic host character
at every coefficient prime for \(2e\le v^{\varepsilon/128}\).
In particular the bound applies to \(J_d\) at three with the
moving discriminant \(-d\) independently of the exponent of
\(\operatorname{Cl}(\mathbf Q(\sqrt{3d}))[3^\infty]\). For
Case~V it includes every character of order at most
\(v^{\varepsilon/128}\), and also those of order at most
\(v^{\varepsilon/64}/\ell\), since \(e=1\) there. Entries outside
the three ranges still require the other estimates of
Section~\ref{famsel:complement}.

\begin{corollary}[Inducing fields on the quartic zeta tests]
\label{adds:quartic}
On a sufficiently late zeta host let \((\ell,e,\xi)\) be in Case II
and consider an absolutely irreducible point of the good quartic
problem in equal characteristic. If it becomes reducible over a
finite extension, then \(\xi\) is quadratic,
\(\ell-1\mid 2e\), and \(\ell-1\nmid e\). Writing \(\Delta_\xi\)
for the fundamental discriminant of \(\xi\), the point is induced
from a strongly irreducible plane over the real quadratic field
of discriminant
\[
            D=\Delta_\xi\,(-1)^{(\ell-1)/2}\ell .
\]
This same field is the base of the strict Selmer calculation:
Proposition~\ref{famsel:quadraticlayers} computes the initial layers
there, and Proposition~\ref{famsel:alltower} computes deeper layers
over its indicated cyclotomic extensions.
The number of possible input triples with \(2\le e<t\) is
\(O(t\log\log(3t))\), uniformly in the host. Among them the total
weight of those with
\(\exp\operatorname{Cl}(\mathbf Q(\sqrt D))[\ell^\infty]\mid
\ell^{v_\ell(e)}\) is \(O(t\log\log(3t)\log(2Nt))\).
\end{corollary}
\begin{proof}
The Case II primes are odd. The fixed exceptional ones are \(3,7\)
and the others satisfy \(\ell-1\mid2e\). Hence none equals the two
odd level primes \(q,d>\max(2t+1,7)\). The factor of the residual
character \(b=\bar\varepsilon_\ell^{\,e}\bar\xi\) at \(\ell\)
and its factor of level \(N\) thus have independent images by the
Chinese remainder theorem. Proposition~\ref{hr:case3strong} forces
\(b\) to be quadratic and ramified at \(\ell\). Its cyclotomic factor
therefore has order exactly two, giving both divisibility claims.
The host character has two-power order, which reduction at the odd
prime preserves. Case II requires its order to be even. It must
therefore be quadratic as well. The two primitive discriminants
are coprime, and their product is positive by the parity of the
input. This gives the stated \(D\) and identifies the field with
the first field in Proposition~\ref{famsel:quadraticlayers}.

The quadratic intertwiner splits into two planes over this field,
as in Proposition~\ref{hr:ind}. Each has residual semisimplification
a sum of characters. Suppose one plane \(U\) becomes reducible on
an open subgroup. On an open normal subgroup of \(G_{\mathbf Q(\sqrt D)}\)
contained in it, the restriction is semisimple: the translates of
an invariant line span the irreducible plane. It is thus a sum of
two characters. At every finite place their inertial images are
finite by the equal-characteristic inertia conditions in the proof
of Proposition~\ref{hr:ind}. Only finitely many places ramify.
A common power of each character is everywhere unramified and
factors through a finite narrow class group. Both characters have
finite image. Passing to powers in the finite coset quotient now
makes all eigenvalues of \(U(g)\) roots of unity for every \(g\).
In equal characteristic they are constants and reduce injectively.
The characteristic polynomials consequently agree with those of
the constant lifts of the two residual characters. Brauer--Nesbitt
contradicts the irreducibility of \(U\).

There are at most four quadratic host characters of each parity,
as in Corollary~\ref{famsel:zetasubsum}. For a fixed \(\ell\) at
most \(2t/(\ell-1)\) positive twists qualify.
For \(y=2t+1\) set \(s=1+1/\log y\). The Euler product gives
\(\sum_{p\le y}p^{-1}\le\mathrm e\log\zeta(s)
\le\mathrm e\log(1+\log y)\). This proves the claimed count.
Under the exponent condition Proposition~\ref{famsel:quadraticlayers}
gives weight \(\log|\operatorname{Cl}(\mathbf Q(\sqrt D))[\ell^\infty]|\),
at most \(\log(2Nt)\) by the ideal bound of
Corollary~\ref{famsel:zetasubsum}. Summing proves the last assertion.
\end{proof}

\begin{corollary}[The full Case IV character sum]\label{adds:fixed-char}
Let the primitive target have conductor \(M\), and fix any finite set
\(P\) of odd coefficient primes containing its exceptional primes
\(B_\psi\). The interval in Theorem~\ref{adds:dyadic-window} can
be chosen so that the sum of \(w_{\ell,e,\xi}\) over
\(\ell\in P\), \(1\le e<t(v)\), and all host characters of primitive
conductor dividing \(M\), is \(O_{\psi,P}(\exp(v^{4\varepsilon/5}))\).
In particular this bounds the entire Case IV contribution.
\end{corollary}
\begin{proof}
For each \(\ell\in P\), prescribe the fixed field
\(B_\ell=\mathbf Q(\mu_{4M\ell})\) and its base layer \(m_\ell=0\)
in Theorem~\ref{adds:dyadic-window}. A character of conductor dividing
\(M\) is trivial over \(B_\ell\). Restriction takes its strict group
into the strict group there after scalar extension. At the places of
the host level the class restricts to zero, and unramified classes
at the other primes away from \(\ell\) stay unramified. The kernel
bound in the proof of Proposition~\ref{newcoh:fullsupport}, applied
to this fixed field, gives
\[
 w_{\ell,e,\xi}\ \le\ (\log\ell)h_\ell(B_\ell,e)
                   +O_{\psi,P}(1+\log e).
\]
Indeed the group of the extension has bounded order, and the
cyclotomic image contains \(1+\ell^{v_\ell(4M\ell)}\). Its action
kills the coefficient invariants by
\(\ell^{v_\ell(4M\ell)+v_\ell(e)}\), so counting cocycles of the
finite quotient gives the asserted normalized kernel loss.
There are at most \(\varphi(M)\) distinct lifts at each prime.
Summing on the jointly selected interval proves the bound.

On a late host the primes \(q,d\) lie outside \(B_\psi\) and are
prime to \(2M\). Their local character factors have order at most
two by Proposition~\ref{host:art}; the local factor at two has
two-power order. Taking the prime-to-\(\ell\) part for odd \(\ell\)
preserves each such factor. Removing the local factor at
\(\ell\in B_\psi\) leaves them all in \(\chi_P\), which by Case IV
has odd order. They are consequently trivial. At each remaining
prime the conductor exponent is at most that of \(M\), since
the host conductor divides \(\operatorname{lcm}(2^{r+2},M)qd\).
Every Case IV character therefore occurs in the preceding sum.
\src{src/Development/AutomorphicArithmetic/AutomorphicVarieties/HeightAndVanishing/DirichletPrimeToLPartMulLPowerPart.lean}{86}{removing the local factor}
and \src{src/Development/AutomorphicArithmetic/AutomorphicVarieties/HeightAndVanishing/DirichletPrimeToLPartMulLPowerPart.lean}{217}{the Case IV order condition}.
\end{proof}

\subsection{The odd-prime field budgets for a general target}\label{addo:field}

Fix a primitive target \(\psi\) of order \(o\) and a host of
Proposition~\ref{host:art}. Write \(E\subset\mathbf Q(\mu_N)\) for
its CM field, \(G=\operatorname{Gal}(E/\mathbf Q)\), \(D=|G|\), and
\(c\in G\) for complex conjugation. The exact sequence for \(H_e\) in
that proposition shows that \(D/o\) is a power of two. For an odd
prime \(\ell\), a twist \(e\ge1\), and \(s\in\{1,-1\}\), put
\[
 M_{\ell,e}=S_\ell(E,(\mathbf Q_\ell/\mathbf Z_\ell)(-e)),\qquad
 b_s=(1+sc)/2,\qquad
 h_{\ell,e}^s=\operatorname{length}_{\mathbf Z_\ell} b_s M_{\ell,e}.
\]
Here \(c\) has its natural action on cohomology with the indicated
twist. Use the character weights \(w_{\ell,e,\xi}\) of
Section~\ref{famsel:complement}, with the same definition for either
sign. Write \(F/\mathbf Q_\ell\) for a splitting field for \(G\),
\(O\) for its ring of integers, and \(a_O=e(O/\mathbf Z_\ell)\).

\begin{proposition}[Comparison on each sign]\label{addo:compare}
For either sign and every odd prime,
\[
 (\log\ell)h_{\ell,e}^s
 \ \le\ \sum_{\xi(c)=s} w_{\ell,e,\xi}
       + C_\psi v_\ell(o)(\log\ell)D(1+\log(2N)).
\]
If \(\ell\nmid o\), restriction gives equality with the character sum.
The total error on any finite set of odd primes and all \(1\le e<t\)
is \(O_\psi(tD(1+\log(2N)))\).
\end{proposition}

\begin{proof}
Put \(n=e+1\), \(U=\operatorname{Spec}\mathcal O_E[1/(N\ell)]\), and
\(H=H^2_{\rm et}(U,\mathbf Z_\ell(n))\). At every finite
\(w\nmid\ell\), an unramified local class with the field coefficients
is zero. Indeed its group is the quotient of the divisible coefficient
module by the image of \((\mathrm Nw)^{-e}-1\), a nonzero scalar.
At \(w\mid\ell\) local vanishing is imposed. Thus the strict field
group is the kernel of localization for the restricted Galois group
with the primes over \(N\ell\) allowed. Poitou--Tate duality and the
degree-two comparison in the proof of Theorem~\ref{newcoh:lemmaK}
give a \(G\)-equivariant injection \(M_{\ell,e}^{\vee}\hookrightarrow H\).
The group \(H\) is finite. Over \(\mathcal O_E[1/\ell]\) this follows
from the even Quillen--Lichtenbaum comparison and finiteness of
\(K_{2e}(\mathcal O_E)\) as in Proposition~\ref{newcoh:fullsupport}.
Further localization adds subquotients of
\(H^1(k_w,\mathbf Z_\ell(e))\), finite since \((\mathrm Nw)^e\ne1\).
The character groups are finite by Proposition~\ref{newcoh:fullsupport}.

For \(\ell\mid o\) we need a bound for the number of generators,
uniform in the twist. The coefficient exact sequence injects
\(H/\ell H\) into \(H^2_{\rm et}(U,\mu_\ell^{\otimes n})\). Restrict
to \(L=E(\mu_\ell)\), writing \(d_L=[L:\mathbf Q]\) and
\(U_L=\operatorname{Spec}\mathcal O_L[1/(N\ell)]\). This restriction
is injective since the cover is finite \'etale of degree dividing
\(\ell-1\) and trace composes with restriction to multiplication by its
degree. A primitive \(\ell\)-th root identifies the coefficient sheaf
on \(U_L\) with \(\mu_\ell\). The Kummer sequence gives
\[
 \dim_{\mathbf F_\ell} H/\ell H
 \le \dim_{\mathbf F_\ell}\operatorname{Pic}(U_L)/\ell
       +\dim_{\mathbf F_\ell}\operatorname{Br}(U_L)[\ell].
\]
The first term is at most \(\log h_L/\log\ell\), where \(h_L\) is
the class number. Minkowski gives each ideal class an integral ideal
of norm at most \(2^{d_L}|\Delta_L|^{1/2}\). The number of ideals
of norm at most \(y\) is at most \(y^2\zeta_L(2)\le y^2\zeta(2)^{d_L}\).
Hence \(\log h_L\ll d_L+\log|\Delta_L|\). The cyclotomic discriminant
bound and the discriminant formula in towers give
\(\log|\Delta_L|\le d_L\log(\ell N)\), since \(L\subset\mathbf Q(\mu_{\ell N})\).
By local invariants the Brauer term is at most the number of places
of \(L\) over \(N\ell\), which is at most \(d_L(1+\omega(N))\).
Here \(L\) is totally imaginary, and \(\omega(N)\) counts the distinct
prime divisors of \(N\). Since \(d_L\le(\ell-1)D\) and \(\ell\mid o\),
we obtain
\[
 r_\ell:=\dim_{\mathbf F_\ell} H/\ell H
       \ll_\psi D(1+\log(2N)).
\]
Taking \(\ell\)-torsion in the injection of \(M_{\ell,e}^{\vee}\)
and using the invariant factors of the finite group \(H\) shows that
\(M_{\ell,e}\), and hence each direct summand \(b_s M_{\ell,e}\),
has at most \(r_\ell\) generators.

Let \(M_O=M_{\ell,e}\otimes_{\mathbf Z_\ell}O\). For a host character
\(\xi\) put \(T_\xi=\sum_{g\in G}\xi(g)^{-1}g\in O[G]\).
A class in \(M_O\) can be viewed over \(E\) with coefficients
\(A_\xi=(F/O)(\xi^{-1}\chi_\ell^{-e})\), since \(\xi\) is trivial on
\(G_E\). Finite free extension of scalars commutes with the cochains
and the local kernels. The class vanishes locally at every finite
place of \(E\), so its corestriction vanishes locally at every finite
place of \(\mathbf Q\) by the local restriction--corestriction formula.
In particular this corestriction lies in \(S_{\ell N}(\mathbf Q,A_\xi)\).
After untwisting over \(E\), its restriction is \(T_\xi\) applied to
the original class. Thus \(T_\xi M_O\) lies in the restriction image
of that strict character group. Character orthogonality gives the
integral identity \(\sum_{\xi(c)=s}T_\xi=D b_s\). Consequently the
length of \(D b_s M_O\) is bounded by the sum of the character lengths
of sign \(s\). At \(\ell\mid o\) the quotient \(b_s M_O/D b_s M_O\)
has length at most \(a_Ov_\ell(D)r_\ell\). Dividing by \(a_O\) and
multiplying by \(\log\ell\) proves the stated inequality. At
\(\ell\nmid o\) the exact restriction descent of
Proposition~\ref{zetafield:descent} applies verbatim, since its
averaging and local restriction arguments require only \(\ell\nmid D\).
Finally \(\sum_{\ell\mid o}v_\ell(o)\log\ell\le\log o\).
\end{proof}

\begin{theorem}[Uniform opposite-parity field sum]\label{addo:opposite}
For every finite set \(P\) of odd primes and every integer \(t\ge2\),
\[
 \sum_{\ell\in P}(\log\ell)\sum_{1\le e<t} h_{\ell,e}^{(-1)^{e+1}}
     \ \ll\ D t^2\bigl(D+\log(2Nt)\bigr)
\]
with an absolute constant.
\end{theorem}

\begin{proof}
Retain \(n=e+1\), \(U\), and \(H\) from the preceding proof. Set
\(b=b_{(-1)^n}\), \(R=b\mathbf Z_\ell[G]\), and let \(\Sigma\) contain
the infinite place and precisely the finite primes dividing \(N\).
By the refined Coates--Sinnott theorem
\cite[Theorem~6.11 and Lemma~6.12]{greitherPopescuCS},
the \(b\)-part of the zeroth Fitting ideal of \(H\) is the
annihilator of \(H^1_{\rm et}(U,\mathbf Z_\ell(n))_{\rm tors}\)
times the equivariant value \(\Theta_\Sigma(1-n)\). Here \(E/\mathbf Q\)
is abelian CM, \(\ell\) is odd, and \(\Sigma\) contains all ramified
primes. The required cyclotomic \(\mu\)-invariant vanishes by
Ferrero--Washington~\cite{ferreroWashington} for the abelian field
\(E(\mu_\ell)\). The cited lemma permits \(\ell\notin\Sigma\).
As in Theorem~\ref{zetafield:minus},
\[
 W_n=(3^{Dn}-1)(5^{Dn}-1),\qquad C_n=nN(n+1)!.
\]
The first integer annihilates the torsion in that \(H^1\). By the
coefficient exact sequence the torsion classes come from the
invariants of \((\mathbf Q_\ell/\mathbf Z_\ell)(n)\),
on which Frobenius over \(p\in\{3,5\}\), \(p\ne\ell\), acts by
\(p^{f_wn}\) with \(f_w\mid D\). In particular
\[
 x=C_nW_n b\Theta_\Sigma(1-n)\ \in\ \operatorname{Fitt}_R(bH).
\]
For every character of sign \((-1)^n\) its value is
\(x_\xi=C_nW_n L_\Sigma(1-n,\xi^{-1})\). These values are nonzero
by the functional equation and the nonvanishing of the removed Euler
factors. The generalized Bernoulli formula at modulus \(N\) gives
\(C_n L_\Sigma(1-n,\xi^{-1})\) integral and
\(\lvert L_\Sigma(1-n,\xi^{-1})\rvert\le 4n!N^n\) at every complex
embedding, exactly as in the proof of Theorem~\ref{zetafield:minus}.
Consequently their Galois-stable product satisfies
\[
 P_n=\prod_{\xi(c)=(-1)^n} x_\xi\in\mathbf Z\setminus\{0\},
 \qquad \log|P_n|\ll Dn\bigl(D+\log(2Nn)\bigr).
\]

We show directly that \(\operatorname{length}_{\mathbf Z_\ell} bH
\le v_\ell(P_n)\), including at \(\ell\mid D\). Take an \(R\)-presentation
\(R^m\xrightarrow{A}R^k\to bH\to0\). We may suppose \(k>0\).
Since \(bH\) is finite, each character matrix \(A_\xi\) of this sign
has full row rank in \(F\), so \(m\ge k\). With \(v_O\) the normalized
valuation and \(J\) ranging over the \(k\)-element column sets, put
\(u_\xi=\min_J v_O(\det A_{\xi,J})\). The Fitting ideal is generated
by the maximal minors, so \(u_\xi\le v_O(x_\xi)\).

For an \(m\)-by-\(k\) matrix of indeterminates \(Y\), Cauchy--Binet
shows that \(\pi_O^{-u_\xi}\det(A_\xi Y)\) is integral and has nonzero
reduction. Indeed the monomials in \(\det Y_J\) use exactly the rows
in \(J\), and at least one such minor has a unit coefficient after
normalization. The product of these finitely many reduced polynomials
is nonzero over an algebraic closure of the residue field. Choose a
point at which it is nonzero and lift its coordinates to a finite
unramified extension \(O'/O\). The resulting single scalar matrix
\(Y\) satisfies \(v_{O'}(\det(A_\xi Y))=u_\xi\) for every \(\xi\).
Let \(R'=R\otimes_{\mathbf Z_\ell}O'\). The cokernel of
\(AY:(R')^k\to(R')^k\) surjects onto \(bH\otimes_{\mathbf Z_\ell}O'\).
The module \(R'\) is finite free over \(O'\), and over its fraction
field the sign algebra splits as one copy for every \(\xi(c)=(-1)^n\).
The determinant of \(AY\) as an \(O'\)-linear endomorphism is therefore
\(\prod_\xi\det(A_\xi Y)\). Taking lengths gives
\[
 a_O\operatorname{length}_{\mathbf Z_\ell} bH
 \le \sum_\xi u_\xi\le \sum_\xi v_O(x_\xi)=a_Ov_\ell(P_n).
\]
Duality preserves the sign of the involution \(c\), so the injection
\(M_{\ell,e}^{\vee}\hookrightarrow H\) gives the same upper bound for
\(h_{\ell,e}^{(-1)^n}\). Summing the prime valuations costs at most
\(\log|P_n|\), and summing over \(2\le n\le t\) proves the estimate.
\end{proof}

\begin{corollary}[The residual Selmer sum]\label{addo:remainder}
Fix the target and \(\varepsilon\) as in Theorem~\ref{adds:dyadic-window}.
On its jointly selected late interval, take \(N,E,t=t(v)\) at the
same page. For any finite set \(S\) of primes, start with the triples
\((\ell,e,\xi)\) with \(\ell\in S\) odd, \(1\le e<t\), and
\(\xi(c)=(-1)^e\). Delete Cases IV and Z, the conductor range of
Corollary~\ref{famsel:subwindow}, and the three order ranges of
Proposition~\ref{adds:order}. Denote the rest by \(\mathcal R(S,v)\).
Then
\[
 \sum_{\ell\in S}(\log\ell)\sum_{1\le e<t} h_\ell(E,e)
 \ \le\ \sum_{(\ell,e,\xi)\in\mathcal R(S,v)}w_{\ell,e,\xi}
        +O_{\psi,\varepsilon}\!\left(\exp(v^{4\varepsilon/5})\right).
\]
One may also delete entries proved zero by
Propositions~\ref{famsel:class}, \ref{famsel:primary} and~\ref{famsel:splitray}.
A bound \(o(v^\varepsilon)\) for the logarithm
of two plus the residual sum, uniform in \(S,v\) on these windows, suffices
for the entire field-length budget of Formal input~\ref{red:design}.
\end{corollary}
\begin{proof}
Prescribe the fixed fields of Corollary~\ref{adds:fixed-char} in
Theorem~\ref{adds:dyadic-window} when choosing the common interval.
The latter theorem bounds the term at two. Apply Theorem~\ref{addo:opposite}
to the opposite sign and Proposition~\ref{addo:compare} to the matching
sign. Their field bound and total comparison error are polynomial in
\(D,t,\log(2Nt)\), with \(D\le N\) and \(\log(Nt)\le 2v^{\varepsilon/4}\).
The fixed fields bound Case IV. Case Z vanishes by
Proposition~\ref{newcoh:zero} since \(3\in B_\psi\).
In the small-conductor range the proof of Corollary~\ref{famsel:subwindow}
gives \(v^{\varepsilon/102}\exp(v^{\varepsilon/2})\).
Proposition~\ref{adds:order} bounds the union of the three order
ranges. The weights are nonnegative, so overlaps in these upper bounds
cost no more than their sum. The asserted logarithmic condition now
makes the full sum at most \(\exp(v^\varepsilon/2)\) eventually.
The rescaling in Section~\ref{sup:section} keeps these same hosts.
\end{proof}

\section{A paper proof of the auxiliary-prime selection}\label{pr:sec}

Fix \(0<\alpha\le1/4,\ K\ge1\). Write \(L=\log Y,\ Z=Y^{400K}\).
Consider the primes \(\mathcal P\) between \(\sqrt U\) (excluded) and \(U\)
in a reduced class \(a\bmod m\), with \(8\mid m\). Assume
\[
 L^2\log m\le4\log U,\qquad
 Y^{\alpha/6}\le\log U\le Y^{\alpha/5}L^2 .
\]
Every prime divisor of \(m\) is at most \(Y^{\alpha/5}\) or greater than \(Z\).
All thresholds below are uniform in \(m,a,U\) subject to these inequalities.
We prove the estimates on a whole interval used in Section~\ref{sup:section}.

\begin{lemma}[Primes and quadratic signs in a progression]\label{pr:signs}
For large \(Y\), \(|\mathcal P|\ge U^{9/10}\). For each fixed \(A,B>0\)
there is at most one exceptional squarefree kernel \(s_*\) to be omitted
from the following assertion. For odd squarefree \(1<s\le Y^A\) coprime to \(m\),
\[
 \left|\sum_{p\in\mathcal P}(p/s)\right|\le Y^{-B}|\mathcal P| .
\]
\end{lemma}
\begin{proof}
For the count, \(m\le U^{1/32}\) eventually. Let \(D=2^j m\) be the largest
such multiple at most \(U^{1/32}\). Then \(U^{1/32}/2<D\le U^{1/32}\);
\(a\) remains coprime to \(D\). Maynard's unconditional Theorem~3.2
\cite{maynard} gives
\[
 \pi(U;D,a)\gg \frac{\log D}{\log U}\frac{U}{\sqrt D\,\varphi(D)}
              \gg U^{61/64}.
\]
The theorem requires \(D\) sufficiently large and \(U\ge D^8\); both follow
from the construction. At most \(\sqrt U\) primes have been counted below the
desired interval. The fixed positive constant is absorbed by the strict
inequality \(61/64>9/10\).

For the second assertion put \(D_*=mY^A\). Uniformly for \(t\in[\sqrt U,U]\)
and \(M=ms\), one has \(t\ge M^{12}\) eventually and
\(\log t/\log D_*\gg L^2\). Indeed \(\log D_*=\log m+AL\),
and \(Y^{\alpha/6}/L^3\to\infty\). Apply Thorner--Zaman, Corollary~1.4
and Remark~1.5 \cite{tz}, to each unit class \(b\bmod M\). In terms of
\(\Theta_b(t)=\sum_{p\le t,\ p\equiv b\bmod M}\log p\), this gives
\[
 \Theta_b(t)=H_b(t)(1+O(\Delta)),\qquad
 H_b(t)=\frac{t-\chi(b)t^\beta/\beta}{\varphi(M)} .
\]
The second term is omitted when the exceptional real zero of modulus \(M\)
is absent. The main terms are positive in this range, by (1.7) of
\cite{tz}. The relative error is
\(\Delta=\exp(-c\log t/\log M)+
 \exp(-c(\log t)^{3/5}/(\log\log t)^{1/5})\).
For any fixed \(B'>0\), this is \(O(Y^{-B'})\) uniformly here.

There is an absolute \(c_0>0\) such that at most one primitive character
of conductor at most \(D_*\) can have a real zero in
\((1-c_0/\log D_*,1)\). Indeed two would induce distinct characters
modulo the least common multiple of their conductors, at most \(D_*^2\).
The classical uniqueness of the exceptional zero for that one modulus
(the zero-free region recalled in \cite[\S1.1]{tz}) excludes this when
\(c_0\) is sufficiently small. The extra Euler factors of an induced
character do not vanish on this interval. If this character exists, take
\(s_*\) to be the product of the primes of its conductor not dividing \(m\).

Sum the displayed formula over classes \(b\equiv a\bmod m\) in
\((\mathbf Z/M)^\times\), first without weights, then with weights
\(\lambda(b)=(b/s)\). The Chinese remainder theorem makes the second
sum of the primary terms zero. The exceptional term in the weighted sum
also vanishes unless the component of \(\chi\) modulo \(s\) equals
\(\lambda\): the product of two distinct real characters has sum zero.
In this last case that component is nontrivial at every prime of \(s\).
Its primitive conductor outside \(m\) has radical \(s\). For \(s\ne s_*\)
we must then have \(1-\beta\ge c_0/\log D_*\). The \emph{unweighted}
exceptional term sums to zero in this case, and the magnitude of the
weighted main term divided by the unweighted main sum is
\(t^{\beta-1}/\beta=O(\exp(-c_0\log t/\log D_*))\);
the exceptional zeros under consideration have \(\beta\ge1/2\).
In every case the absolute error summed over classes is
\(O(\Delta\sum_b H_b(t))\), since the individual estimates are relative
to positive main terms. This remains true when the base progression
itself has a small main term.

It follows that the signed prefix sum of \(\log p\), on these unit classes,
is \(O(Y^{-B'})\) times the unsigned prefix sum for \(a\bmod m\).
The latter is at most \((|\mathcal P|+\sqrt U)\log U\) throughout
\([\sqrt U,U]\). Partial summation with \(1/\log t\) on this interval
costs \(O(1/\log U)\), including both endpoints. Primes dividing \(s\)
are below \(\sqrt U\). Use the first assertion to absorb \(\sqrt U\)
and take \(B'>B\).
\end{proof}

\begin{proposition}[A uniform mean square]\label{pr:ms}
For \(|\kappa(n)|\le1\), put
\[
 S_p(y)=\sum_{n\le Z}\Lambda(n)\kappa(n)(n/p)\min(1,y/n).
\]
For \(Y\le y\le Y^K\) one has
\(\sum_{p\in\mathcal P}|S_p(y)|^2\ll_K |\mathcal P|y\log y\).
For each fixed \(\tau>0\), the proportion of primes for which
\(|S_p(y)|>\tfrac12y^{1/2+\tau}\) somewhere in this interval is
\(O_{K,\tau}(Y^{-\tau})\), uniformly in \(\kappa\).
\end{proposition}
\begin{proof}
Eventually \(\sqrt U>Z^2\). For \(n=\ell^j\), let \(r(n)=\ell\) if
\(j\) is odd and \(\ell\nmid m\), and \(r(n)=1\) otherwise. Quadratic
reciprocity expresses \((n/p)\) as a fixed sign times \((p/r(n))\):
the signs at \(2\) and at primes of \(m\) are fixed by the progression,
and \(r(n)\) is odd. Group the absolute coefficients by their square
class, writing \(b_g=\sum_{r(n)=g,\ n\le Z}\Lambda(n)\min(1,y/n)\)
in the finite elementary two-group generated by the primes up to \(Z\)
outside \(m\).

The elementary Chebyshev bound \(\sum_{n\le u}\Lambda(n)\ll u\)
gives \(\sum_g b_g\ll y(1+\log Z)\) by partial summation.
Even powers have cumulative weight at most \(C\sqrt u\) up to \(u\);
integration against \(y\,du/u^2\) for \(u\ge y\) bounds their ramp
weight by \(2C\sqrt y\). For each prime \(\ell\le y\), all its powers
together contribute at most \(\log y+2\log\ell\): after the last power
at most \(y\) the geometric series has total at most \(2\).
Thus \(b_1\ll\sqrt y+Y^{\alpha/5}\log y\ll\sqrt y\).
Among other groups with \(\ell\le y\), at most \(O(y/\log y)\) have
positive weight, each of size \(O(\log y)\). For \(\ell>y\) the bound
is \(2y\log\ell/\ell\). On each dyadic block above \(y\) Chebyshev's
bound gives \(\sum_{\ell>y}(\log\ell)^2/\ell^2\ll(\log y)/y\).
Consequently \(\sum_g b_g^2\ll y\log y\).

In the expansion of the second moment the square class \(gh\) is odd,
squarefree, coprime to \(m\) and at most \(Z^2\) whenever the coefficient
is nonzero. Its normalized prime sum costs at most \(Y^{-K-1}\), except
for \(gh=1\) and possibly \(gh=s_*\), by Lemma~\ref{pr:signs}.
For either exceptional value \(d\), multiplication by \(d\) permutes
the classes; Cauchy--Schwarz gives \(\sum_g b_g b_{gd}\le\sum_g b_g^2\).
The remaining total is \(O(Y^{-K-1}y^2(1+\log Z)^2)\ll_K y\log y\).
This proves the pointwise estimate.

Except at integer breakpoints, \(y S'_p(y)\) is the same sum with
\(\kappa(n)\) replaced by \(\kappa(n)\mathbf1_{n>y}\). Therefore
\(\sum_p|S'_p(y)|^2\ll_K|\mathcal P|(\log y)/y\).
On a dyadic block \([u,v]\subseteq[Y,Y^K]\), \(v\le2u\), integration
of the derivative and Cauchy--Schwarz bound
\(\sum_p\sup_{[u,v]}|S_p|^2\ll_K|\mathcal P|u\log(2u)\).
The proportion violating the asserted threshold on the block is
\(O_K(u^{-2\tau}\log(2u))=O_{K,\tau}(u^{-\tau})\). Sum the geometric
series over blocks, truncating the last block at \(Y^K\). For \(K=1\)
the pointwise estimate suffices.
\end{proof}

\begin{lemma}[Truncation using a finite-height box]\label{pr:perron}
Let \(\chi\) be a primitive nonprincipal character of conductor at most
\(Q\), with \(L(s,\chi)\ne0\) for \(\Re s>399/400,\ |\Im s|\le Q\).
For \(Q\) and \(z\) sufficiently large, \(z\) a half-integer and
\(2\le\log z\le Q^{1/8}\), one has
\[
 \left|\frac{L'}{L}(1,\chi)+\sum_{n\le z}\frac{\Lambda(n)\chi(n)}n\right|
 \ll z^{-1/400}(\log z)^2(\log Q)^2+
        \frac{(\log z)^2(1+\log Q)}{Q}.
\]
\end{lemma}
\begin{proof}
Put \(\sigma_0=399/400,\ \eta=1/\log z,\ T=Q/2\). Assume
\(\eta\le1/800\), as is sufficient for the assertion. For \(|t|\le T\)
the open disk of radius \(2-\sigma_0\) about \(2+it\) is zero-free.
There is a branch of \(\log L\) there with bounded central value, by
the absolutely convergent Euler product at real part two. Periodicity
and partial summation give \(|L(s,\chi)|\le Q|s|/\Re s\) in the disk,
so \(\Re\log L\ll\log Q\). For clarity, if \(F\) is the branch minus
its central value, take \(H\gg\log Q\) with \(\Re F\le H\).
The map \(F/(2H-F)\) takes the disk into the unit disk and vanishes at
the centre. Schwarz's lemma gives \(|F|\ll H/\eta\) on the disk
shortened by \(\eta/2\); Cauchy's derivative estimate on a circle of
radius \(\eta/2\) then gives
\(|L'/L(\sigma+it,\chi)|\ll(\log Q)/\eta^2\) for
\(\sigma_0+\eta\le\sigma\le1+\eta\).

Integrate \(-\frac{L'}L(1+w,\chi)z^w/w\) up the segment
\(\Re w=\eta,\ |\Im w|\le T\), with factor \(1/(2\pi i)\).
Its difference from the prime-power sum is \(O((\log z)^2/T)\).
Indeed the scalar Perron kernel for \(v=z/n\ne1\) has error
\(O(v^\eta/(T|\log v|))\), by integration by parts on the tails of
the vertical inversion integral for \(\mathbf1_{v>1}\).
Outside \(z/2<n<2z\) use \(\sum(\log n)n^{-1-\eta}\ll\eta^{-2}\).
Inside use \(|\log(z/n)|\ge|z-n|/(2z)\), \(|z-n|\ge1/2\), and sum
the resulting harmonic series. Absolute convergence on the initial
line justifies interchanging the sum and the finite integral.
Shift to \(\Re w=-1/400+\eta\); the only residue is
\(-L'(1,\chi)/L(1,\chi)\) at zero. The two horizontal integrals cost
\(O((\log Q)\eta^{-2}/T)\). On the new vertical line
\(|z^w|\le e z^{-1/400}\) and \(\int_{-T}^T|w|^{-1}dt\ll\log(2T)\),
since \(1/400-\eta\ge1/800\). The derivative bound finishes the proof.
\end{proof}

\begin{theorem}[Uniform selection of the two primes]\label{pr:pair}
Fix \(M\ge1,\ 0<\varepsilon<1/52,\ K\ge1\). For every sufficiently
large \(Y\), put \(z_0=Y^{\varepsilon/5}\), \(z_1=z_0(\log Y)^2\).
There are primes \(q\in(e^{z_0/2},e^{z_0}],\
d\in(e^{z_1/2},e^{z_1}]\) with \(q\equiv3\bmod8,\ d\equiv7\bmod16\)
and \(d\equiv-1\) modulo \(q\) and every odd prime at most \(z_0\).
They are coprime to \(8M\). For every primitive inducer \(\xi\) of a
character modulo \(8M\), and every \((i,j)\in\{0,1\}^2\setminus\{(0,0)\}\),
write \(\chi\) for the primitive character associated to
\(\xi(\cdot)(\cdot/q)^i(\cdot/d)^j\). Simultaneously for these choices,
\[
 |G_\chi(y)|\le y^{1/2+\varepsilon/2},\qquad
 |L'(1,\chi)/L(1,\chi)|\le4\log y\qquad(Y\le y\le Y^K).
\]
Every selected twist is nonprincipal and has the zero-free box in
Lemma~\ref{pr:perron} up to the conductor bound used at its selection.
\end{theorem}
\begin{proof}
For one progression take a list of at most \(R\) primitive base characters
of conductor at most \(F\le\sqrt U\), where \(R\) is fixed. For \(p\in
\mathcal P\) twisting a base character of conductor \(f\) by \((\cdot/p)\)
gives primitive conductor \(fp\), by the prime-by-prime conductor factorization.
For \(Q=FU\) the density estimate of Lemma~\ref{sup:density} bounds the
number of bad primes by \(O(Q^{1/4})\): a chosen bad character determines
its prime uniquely because its conductor has only one prime factor \(>\sqrt U\).
Divide by \(|\mathcal P|\ge U^{9/10}\) and use \(Q\le U^{3/2}\).
By Proposition~\ref{pr:ms} a further \(O_{K,\tau}(RY^{-\tau})\) proportion
fails the ramp estimate for some base character. Both proportions tend to zero.

For any survivor choose \(z=\lfloor Y^{400K}\rfloor+1/2\) in
Lemma~\ref{pr:perron}. Since \(\log Q\ll Y^{\alpha/5}(\log Y)^2\),
the error is \(O_K(Y^{-K+2\alpha/5}(\log Y)^6)\). The term divided by \(Q\)
is smaller, since \(\log Q\ge Y^{\alpha/6}\). The elementary identity
\[
 G_\chi(y)=-S_p(y)+y\left(\frac{L'}L(1,\chi)
                           +\sum_{n\le z}\frac{\Lambda(n)\chi(n)}n\right)
\]
follows from \(y\ge1,\ y\le z\) by splitting the ramp at \(n=y\).
The half-integer retains exactly the terms \(n\le Y^{400K}\), and
\(y\le Y^K<z\) for large \(Y\).
For \(\tau=\varepsilon/2\) and either \(\alpha=\varepsilon\) or
\(6\varepsilon/5\), the tail after multiplication by \(y\) is
\(o(Y^{1/2+\tau})\), giving the asserted ramp bound.
For the logarithmic derivative take instead
\(z=\lfloor(\log Q)^{1000}\rfloor+1/2\).
The error in Lemma~\ref{pr:perron} is \(o(1)\). The elementary estimate
\(\sum_{n\le u}\Lambda(n)/n=\log u+O(1)\) follows from
\(\sum_{n\le u}\Lambda(n)\lfloor u/n\rfloor=\log(\lfloor u\rfloor!)\),
Chebyshev's bound and Stirling's formula. Hence the derivative is bounded
by \(1000\log\log Q+O(1)\).

First use \(m=8,\ a=3,\ U=e^{z_0},\ \alpha=\varepsilon,\ F=8M\)
and all primitive inducers from \(8M\). Fix a good \(q\). Next take
\[
 m=16q\prod_{\substack{\ell\le z_0\\\ell\ {\rm odd\ prime}}}\ell,\quad
 U=e^{z_1},\quad \alpha=6\varepsilon/5,\quad F=(8M)e^{z_0}.
\]
The Chinese remainder theorem gives the required reduced class \(a\).
For large \(Y\), \(q>Z,\ F\le\sqrt U,\ Y^{\alpha/6}\le z_1
\le Y^{\alpha/5}L^2\). Also \(\log m\le4z_0\): the elementary bound
\(\sum_{\ell\le u}\log\ell\le(2\log2+o(1))u\) follows by bounding
the primes between \(n\) and \(2n\) by the middle binomial coefficient
and summing dyadically. Thus \(L^2\log m\le4z_1\).
Use both the old characters and their \(q\)-twists, at most \(2\varphi(8M)\)
characters. The resulting \(d\) provides the two remaining twist patterns.
At both stages \(\log Q\le3Y^{\varepsilon/5}L^2\) eventually, so
\(1000\log\log Q\le200\varepsilon L+O(\log L)\). The strict inequality
\(200\varepsilon<4\) gives the promised bound. All choices of \(q,d\)
precede the choice of \(y\).
\end{proof}

\section{The common field, its integral lattice and its translates}\label{host:sec}

Fix a primitive Dirichlet character \(\psi\) of conductor \(M\), of order
\(o\). The arithmetic field itself and the finite sets used in its
translates admit explicit constructions. Both matter in applying the
prime bounds: changing the field after the pair is chosen could change
the list of twists that was tested.

\begin{proposition}[Cyclotomic realization]\label{host:art}
Put \(N_e=\operatorname{lcm}(2^{e+2},M)\), \(U_e=(\mathbf Z/N_e)^\times\),
\(C_e=(\mathbf Z/2^{e+2})^\times/\{\pm1\}\), \(B=\operatorname{im}\psi\).
Let \(\eta_e:U_e\to C_e\times B\) send \(u\) to \(([u],\psi(u))\),
and set \(H_e=\operatorname{im}\eta_e,\ A_0=\ker(H_e\to B)\).
Then \(A_0\) is cyclic, \([H_e:A_0]=o\), and \(|H_e|=c\,2^e\) for all
sufficiently large \(e\), where \(1\le c\le o\) is independent of \(e\).

Take \(e\ge2\) and a pair \(q,d\) from Theorem~\ref{pr:pair}, with
\(Y^{\varepsilon/5}\ge4M\). There is an abelian CM field \(E\) with
Galois group \(\{\pm1\}^2\times H_e\). The characters trivial on
\(\{(1,1,h):h\in A_0\}\), pulled back to integers prime to \(N_e qd\),
are precisely
\[
      \psi^a J_q^{\,i}J_d^{\,j},
 \qquad 0\le a<o,\quad i,j\in\{0,1\},\qquad J_p(u)=(u/p).
\]
The retained characters are the \(o\) primitive inducers of the
\(\psi^a\). Each remaining one is among the twists tested on the same pair.
\end{proposition}
\begin{proof}
Reduction of unit groups modulo a divisor is onto, by the Chinese
remainder theorem prime by prime. It follows that both projections of
\(H_e\) are onto. The kernel of the second injects into \(C_e\), cyclic
of order \(2^e\): for \(e\ge1\) the class of \(5\) generates, since
\(v_2(5^{2^j}-1)=j+2\) by successive squaring and all its powers are
\(1\bmod4\). For \(e=0\) the quotient is trivial. The kernel of the first
projection identifies with \(B_e=\psi(\{u\in U_e:u\equiv\pm1\bmod
2^{e+2}\})\). Reduction from level \(e+1\) shows \(B_{e+1}\subseteq B_e\).
This descending chain in the fixed finite group \(B\) stabilizes. Take
\(c\) equal to the eventual order and use the exact sequence with quotient
\(C_e\).

Define \(w:\{\pm1\}\to H_e\) by \(w(h)=\eta_e(h)\).
By the Chinese remainder theorem the homomorphism
\[
    u\longmapsto
    \bigl(J_q(u),J_d(u),\,\eta_e(u)w(J_q(u))\bigr)
       \quad\text{on }(\mathbf Z/N_eqd)^\times
\]
is onto. Indeed the signs can be prescribed at \(q,d\), and one then
chooses the \(\eta_e\)-image to attain the specified third coordinate.
Take \(E\) to be the subfield of \(\mathbf Q(\zeta_{N_e qd})\) fixed
by the kernel. Its kernel condition is exactly
\[
 u\equiv\pm1\bmod 2^{e+2},\qquad \psi(u)=J_q(u)=J_d(u)=1.
\]
Thus \(E\) is the compositum of the maximal real \(2^{e+2}\)-cyclotomic
field, the field cut out by \(\psi\), and the quadratic fields of
discriminants \(-q,-d\). In fact \(q,d\equiv3\bmod4\).
Complex conjugation has coordinates \((-1,-1,1)\), since \(w(-1)^2=1\).
Its fixed field is totally real by abelianness, and \(E\) contains the
imaginary quadratic field \(K=\mathbf Q(\sqrt{-d})\).
The integer \(N_e\) is also the least divisor of \(N_e\) through which
\(\psi\) factors and which is divisible by \(2^{e+2}\): a modulus for
the primitive \(\psi\) must be divisible by \(M\).

A character of \(H_e\) killing \(A_0\) is a character of \(B\). Since
\(B\subset\mathbf C^\times\) is cyclic of order \(o\), it raises values
of \(\psi\) to a power. Its value on the \(w(J_q)\) term is a sign
character of \(J_q\), absorbed into \(i\). Conversely choose the character
on the first sign to undo this additional sign. This proves equality of
the two character lists. They consist of \(4o\) distinct characters:
one can vary \(J_q\) and \(J_d\) independently while \(u\bmod N_e\) is
fixed, and conversely can set both signs to one. Passing to primitive
inducers at a fixed common modulus preserves distinctness. In particular
there are exactly three nonempty binary choices tested per power of \(\psi\).
\end{proof}

\begin{lemma}[The half-different]\label{host:different}
There is an integral ideal \(\mathfrak b\) of \(E\), invariant under
\(\operatorname{Gal}(E/K)\), such that
\(\mathfrak D_{E/K}=\mathfrak b\,\overline{\mathfrak b}\) and
\(\mathfrak b+\overline{\mathfrak b}=\mathcal O_E\).
The \(\mathcal O_K\)-lattice \(I=\mathfrak b^{-1}\) is self-dual for
\((u,v)\mapsto\operatorname{Tr}_{E/K}(u\bar v)\).
\end{lemma}
\begin{proof}
All finite ramification in \(E/K\) lies above primes dividing \(N_eq\).
Indeed the other three factors in the compositum with \(K\) are
unramified outside \(N_eq\), and unramified local extensions stay
unramified after base change. Each such rational prime splits in \(K\).
For odd primes dividing \(M\), and for \(q\), this follows from
\(-d\equiv1\bmod p\). For \(p=2\) use the odd discriminant
\(-d\equiv1\bmod8\). The different is invariant under every
\(\mathbf Q\)-automorphism of \(E\), by the defining trace dual.
For each rational prime in its support choose one of the two places of
\(K\) and place in \(\mathfrak b\) every factor of the different above
that place with its full multiplicity. Conjugation supplies exactly
the disjoint remaining factors; automorphisms fixing \(K\) preserve
the choice. For any fractional ideal \(J\), the integrality of
\(\operatorname{Tr}_{E/K}(u\bar J)\) is equivalent to
\(u\bar J\subseteq\mathfrak D_{E/K}^{-1}\), since \(\bar J\) is closed
under multiplication by \(\mathcal O_E\). Hence
\(I^\#=(\bar I)^{-1}\mathfrak D_{E/K}^{-1}=\mathfrak b^{-1}\).
\end{proof}

\begin{lemma}[The rank fits the cohomology bound]\label{host:rank}
Suppose \(x^{2+8\varepsilon}\le n=|H_e|\le2x^{2+8\varepsilon}\) and
\(0<\varepsilon<1/52\). Put \(f=2n,\ L=x^\varepsilon\) and choose \(m,k\)
to be the least even integers at least \(x^{1/2-\varepsilon/16}\) and
\(x^{9/2+10\varepsilon}\), respectively. Define
\[
 \begin{split}
 r&=k+m-1,\qquad P=\lceil x^{377}\rceil,\\
 D&=8\left\lceil \frac{10r}{8}
       \left(f+\frac{f^2(k-1)^3L^2}{m}\right)\right\rceil,\qquad
 t=8fr+4+PD-r .
 \end{split}
\]
Then \(t\) is a positive odd integer and \(t\le x^{400}\) for all
sufficiently large \(x\), uniformly in \(0<\varepsilon<1/52\).
\end{lemma}
\begin{proof}
Here \(r\) is odd, \(f\ge1\), and \(D\) is divisible by eight, proving
positivity and parity. With \(u=x^{1/100}\ge2\), the stated ranges give
\[
 f\le4u^{216},\quad k\le2u^{470},\quad r\le4u^{470},\quad
 m\ge u^{49},\quad L^2\le u^4.
\]
The rounding bound \(2\lceil v/2\rceil\le2v\) holds for \(v\ge2\).
The exponent of \(f^2(k-1)^3L^2/m\) in \(u\) is at most
\(432+1410+4-49=1797\), with coefficient \(128\). Consequently
\(D\le5288u^{2267}\), while \(P\le2u^{37700}\). Thus
\[
 t\le 10708u^{39967}\le u^{40000}
\]
because \(10708\le2^{33}\le u^{33}\). This strict margin is needed
when the cohomology window is read at the rescaled value in
Section~\ref{sup:section}.
\end{proof}

\begin{proposition}[A dense set with bounded orbit intersections]\label{host:behrend}
For \(n\ge1,\ L\ge2\) integers there is \(S\subset[0,n/(4L))\cap\mathbf Z\)
with
\[
 |S|\ge (n/L)\exp\{-C\sqrt{\log n\,\log L}\}
\]
for an absolute \(C\). For distinct \(s_1,s_2,s_3\in S\) and distinct
nonzero \(i,j\in[-L,L]\cap\mathbf Z\) one has
\(j(s_2-s_1)\ne i(s_3-s_1)\). In \(\mathbf Z/n\), every translate of
\(S\) meets \(L\) consecutive terms of an orbit of step \(h\) in at most
two terms if \(\operatorname{ord}(h)>4L\). For
\(\operatorname{ord}(h)\le4L\) it meets the whole orbit in at most one
distinct point.
\end{proposition}
\begin{proof}
Set \(v=\log(4L)\). If \(\log n\le64v\), the singleton \(\{0\}\)
suffices on increasing the absolute constant, since \(v\le3\log L\).
Otherwise put \(S_0=\log(n/(4L))\),
\(d_0=\lfloor\sqrt{S_0/v}\rfloor,\ b=\lfloor e^{S_0/d_0}/(4L)\rfloor\)
and \(B=4Lb\). These choices give \(b\ge2,\ b\ge e^{S_0/d_0}/(8L)\).
Encode \(\{0,\ldots,b-1\}^{d_0}\) as base-\(B\) integers. All encoded
values are \(<B^{d_0}\le n/(4L)\). One level of the squared Euclidean norm
has at least \(b^{d_0}/(2d_0 b^2)\) vectors. Its log cardinality is at least
\[
 S_0-d_0\log(8L)-2S_0/d_0-\log(2d_0)
 \ \ge\ S_0-C_1\sqrt{S_0 v}.
\]
Here \(d_0\) lies between \(\tfrac12\sqrt{S_0/v}\) and \(\sqrt{S_0/v}\).
Absorbing \(\log4\) and using \(v\le3\log L\) gives the claimed size.
If the asserted integer relation held, each digit of the signed
difference would have absolute value at most \(2L(b-1)<B\).
Reducing successively modulo \(B\) forces every such digit to be zero.
This would give three distinct collinear vectors on the same Euclidean
sphere, which a line meets in at most two points.

For three hits at times \(a<b'<c\) in a length-\(L\) segment with step
of order \(>4L\), the representatives \(s_1,s_2,s_3\) are distinct.
The translates cancel, so
\((c-a)(s_2-s_1)-(b'-a)(s_3-s_1)\equiv0\bmod n\).
Its absolute value is \(<n\), and its vanishing over \(\mathbf Z\)
contradicts the preceding property. For an orbit of order \(m'\le4L\),
any two hits have \(n\mid m'(s_2-s_1)\) with absolute value \(<n\).
They therefore coincide.
\end{proof}

In the application \(n_0=|A_0|=n/o\), \(o\) is fixed and
\(L_w=10\lceil405\log x\rceil+2\lceil30/\varepsilon\rceil+22\).
Use \(S\) in the cyclic group \(A_0\) with \(n_0,L_w\) in place of
\(n,L\). For fixed \(\varepsilon\) its lower bound divided by \(n\)
is at least \(x^{-\varepsilon}\) eventually: the logarithm of its
reciprocal is bounded by
\(\log(oL_w)+C\sqrt{\log n_0\log L_w}=o(\log x)\).
Thus the set giving the lower density also supplies the orbit restriction
at every translate.

\begin{lemma}[One choice of the top set]\label{host:shares}
In a finite group \(H\) of order \(n\asymp x^{2+8\varepsilon}\) with
subgroup \(A_0\) of fixed index \(o\), put \(\delta=x^{-\varepsilon/8}\).
Include points independently with probability \(\delta\) on the identity
coset and \(3\delta/4\) on the others. Let \(D(T)\ge0\) be any weighted
cost with \(\mathbf E D(T)\le B_0\). For sufficiently large \(x\) there is
one \(T\) with \(D(T)\le4 B_0,\ |T|\ge n\delta/2\), and
\[
  \big||T\cap C|-\lambda_C|T|\big|\le |T|x^{-1+\varepsilon}
  \quad\text{at every coset},\qquad
  \lambda_C=\frac{3+\mathbf1_{C=A_0}}{3o+1}.
\]
\end{lemma}
\begin{proof}
Let \(\mu=\mathbf E|T|\ge3n\delta/4\). Each coset count has mean
\(\lambda_C\mu\) and variance at most \(\mu\), by independence. For
\(a=\mu x^{-1+\varepsilon}/(4(o+1))\), Chebyshev and a union bound
give failure probability at most \(o\mu/a^2=O_o(x^{2-2\varepsilon}/\mu)
=o(1)\). On success, \(\big||T|-\mu\big|\le oa\), hence
\(|T|\ge n\delta/2\) and \((o+1)a\le |T|x^{-1+\varepsilon}\).
Apply the triangle inequality to each coset. Markov's inequality gives
\(\Pr(D>4B_0)\le1/4\) for \(B_0>0\); for \(B_0=0\) the cost is zero
almost surely. The two good events intersect.
\end{proof}

For example, if the cost counts chosen points \(t\) for which at least
\(k+1\) points of some specified window \(W\) of size at most \(L\) are
chosen, its expectation is bounded by a union over the windows and their
subsets. For \(t\in W\), the event requires \(t\) and at least \(k\)
other points, costing at most \(\delta(L\delta)^k\); for \(t\notin W\),
the bound is \(\delta(L\delta)^{k+1}\). This supplies a valid input
\(B_0\) whenever the windows and their weights are specified.
Lemma~\ref{host:shares} does not give a simultaneous automorphic realization
of the resulting translates. That additional assertion identifies the
centre \(a_t\) and the weights \(w_{u,j,\rho}\) for all the local periods
with the \emph{same} values in the Fourier discrepancy. The quantitative
form still needed is given in Section~\ref{scope:account}.

\subsection{Bounds for all characters at the thin coordinates}

\begin{proposition}[Uniform thin-coordinate estimate]\label{host:thin}
Fix \(M,\varepsilon\) as in Theorem~\ref{pr:pair}. Suppose
\(Y\le x\), \(x^{2+8\varepsilon}\le |H_e|\le2x^{2+8\varepsilon}\),
and \(q,d\) are the pair selected at \(Y\). For all sufficiently large
\(x\) (with \(Y\) above the pair threshold), the primitive Dirichlet inducer
\(\chi\) of each character of \(\operatorname{Gal}(E/\mathbf Q)\), for the field \(E\) of Proposition~\ref{host:art}, satisfies
\[
        |G_\chi(y)|\le y e^{-x^{\varepsilon/8}}
        \qquad \bigl(y\ge \exp(x^{\varepsilon/3})/3\bigr).
\]
\end{proposition}
\begin{proof}
Write \(N\) for the primitive conductor. From \(N\mid N_eqd\),
\(2^e\le|H_e|\), and the ranges of the pair, we obtain
\(H=\log(2N)\ll_{M,\varepsilon}x^{\varepsilon/5}(\log x)^2\).
There is a fixed \(\nu>0\) such that any exceptional real zero of \(\chi\),
if present, is at most \(1-\nu\). A primitive real character
on the cyclotomic modulus \(N_eqd\) factors as a product of local
quadratic or trivial characters. At an odd prime the principal units
are a group of odd order and are killed by a quadratic character.
At 2 every odd square is \(1\bmod8\), and every unit \(1\bmod8\) is
a square modulo a power of 2 (lift an odd square one bit at a time).
The part away from \(q,d\) therefore has conductor dividing \(8M\).
If a \(q\)- or \(d\)-factor is present, Theorem~\ref{pr:pair} supplies
the gap \(1/400\) at real height zero. The remaining real primitive
characters form a fixed finite set. Classical nonvanishing at one and
analyticity give them a common positive gap near one. At conductor one
use analyticity and nonvanishing of \((s-1)\zeta(s)\) at one.

Let \(D_\chi(t)=\sum_{n\le t}\Lambda(n)\chi(n)
-\mathbf1_{\chi=1}t\).
For \(t\ge(2N)^{12}\) sufficiently large, Corollary~1.4 of \cite{tz}
applied at modulus \(N\ge2\) gives
\[
 |D_\chi(t)|/t\ll
 e^{-c\log t/H}+e^{-c\sqrt{\log t}}+t^{-\nu}+t^{-1/4}.
\]
Indeed, multiply the formula for each unit class by \(\chi\) and sum.
Its primary terms cancel for a nonprincipal character. Its exceptional
term cancels unless \(\chi\) is the corresponding real character, in
which case the real gap bounds it. The absolute errors add to
\(O(t\Delta)\) because \(H_b(t)\le 3t/\varphi(N)\) in the notation
of Lemma~\ref{pr:signs}. Prime powers of exponent at least two contribute
\(O(\sqrt t\log^2t)\), by summing the trivial bound for their prime bases.
The second exponential weakens the corresponding bound in Corollary~1.4.
For \(N=1\), use the same formula for the odd class modulo two, add the
prime \(2\) and then the higher powers.

Summation of the ramp gives
\(P_\chi(y)=y\int_1^y(\sum_{n\le t}\Lambda(n)\chi(n))t^{-2}\,dt\).
In fact, for both kinds of character,
\[
            G_\chi(y)=-y\int_y^\infty D_\chi(t)t^{-2}\,dt .
\]
To identify the constant, use
\(-L'/L(s,\chi)=s\int_1^\infty
(\sum_{n\le t}\Lambda(n)\chi(n))t^{-1-s}\,dt\) for real \(s>1\).
After subtraction, the displayed error is integrable against \(dt/t\),
so the limit \(s\downarrow1\) gives the nonprincipal constant. In the
principal case subtract \(s/(s-1)\) before taking the limit; the
result is \(-1-\gamma\), as required by \eqref{eq:centred}.

Put \(A=\log y\ge x^{\varepsilon/3}-\log3\). Integrating the first two
errors from \(A\) to infinity in the variable \(\log t\) gives
\(O(H e^{-c A/H}+(1+\sqrt A)e^{-c\sqrt A})\);
the power errors give \(O_\nu(e^{-\nu A}+e^{-A/4})\).
All are \(o(e^{-x^{\varepsilon/8}})\) uniformly up to \(y=\infty\).
Indeed \(A/H\gg x^{2\varepsilon/15}/(\log x)^2\) and
\(\sqrt A\gg x^{\varepsilon/6}\), with both exponents strictly larger
than \(\varepsilon/8\). In particular \(y\) is eventually above
\((2N)^{12}\).
\end{proof}

\begin{lemma}[Variation of the ramp]\label{host:variation}
For \(a,w,b>0\), extending the same ramp formula to \(0<y<1\) where
necessary, one has
\[
 |wG_\chi(a/w)-bG_\chi(a/b)|\le 3a\,|w-b|/\min(w,b).
\]
For \(y,y'\ge1\) the change in \(G_1\) is \(O(|y-y'|)\); that in
\(G_\chi\) for \(\chi\ne1\) is at most
\(|y-y'|(\log\max(y,y')+C+|L'(1,\chi)/L(1,\chi)|)\).
\end{lemma}
\begin{proof}
Write \(\Psi(v)=\sum_{n\le v}\Lambda(n)\). One has
\(\Psi(v)\le 4(\log2)v<3v\) also for small \(v>0\).
Indeed each prime power between \(2^{j-1}\) (excluded) and \(2^j\)
contributes its prime to \(\binom{2^j}{2^{j-1}}\) by the factorial
valuation formula. Bound the binomial coefficient by \(2^{2^j}\), sum
over \(j\), and round \(v\ge1\) up to a power of two. Below one
\(\Psi(v)=0\). Off the finitely many breakpoints on a compact interval
of \(u>0\), the derivative of
\(uP_\chi(a/u)=\sum_n\Lambda(n)\chi(n)(a/n-u)_+\)
has absolute value at most \(\Psi(a/u)\).
The nonprincipal compensation is constant after multiplication by \(u\).
In the principal case the total derivative is \(a/u-\Psi(a/u)\),
also of magnitude at most \(3a/u\). Integrate between \(w\) and \(b\).
For variation in \(y\), differentiate the ramp itself and use
\(\sum_{n\le y}\Lambda(n)/n=\log y+O(1)\), proved in
Theorem~\ref{pr:pair}. For the principal term subtraction of \(\log y\)
leaves a bounded derivative. The functions are continuous at breakpoints,
so integration gives the asserted bounds there as well.
\end{proof}

\begin{proposition}[The analytic CM scalar on the selected hosts]\label{host:period}
Fix a primitive character \(\psi\)
and a positive exponent \(\varepsilon\) as in
Theorem~\ref{pr:pair}. At a selected scale \(x\) let \(E\) be the
field in Proposition~\ref{host:art}, let \(d\) be the second
auxiliary prime and put \(L=x^\varepsilon\).
In the normalization of the determinant comparison in Section~9 the
real scalar is
\[
 A_{J_d}=\gamma+\log(2\pi)-\log\operatorname{rd}(E)
                         -\frac{L'(1,J_d)}{L(1,J_d)} .
\]
Uniformly in all allowed primes and cyclotomic levels at \(x\),
\[
 |A_{J_d}|+\log d+\log\operatorname{rd}(E)
          \ll_\psi (1+x^{\varepsilon/5})\log^2(2x)=o(L).
\]
Here and in the little-oh estimate the threshold may depend on
\(\psi,\varepsilon\). In particular the analytic size condition
\(|A_{J_d}|\le3L\) holds throughout every sufficiently late selected
window.
\end{proposition}
\begin{proof}
Write \(M\) for the primitive conductor of \(\psi\).
At the page used in the construction \(2^e\le2x^{2+8\varepsilon}\).
The field \(E\) is contained in the cyclotomic field of conductor
dividing \(C=\operatorname{lcm}(2^{e+2},M)qd\).
The discriminant tower gives \(\operatorname{rd}(E)\le
\operatorname{rd}(\mathbf Q(\zeta_C))\le C\). For the second
inequality, factor \(X^C-1=\Phi_C(X)Q(X)\), differentiate at
\(\zeta_C\), and take norms. Then
\(|N\Phi_C'(\zeta_C)|\le C^{\varphi(C)}\), since \(Q(\zeta_C)\)
is a nonzero algebraic integer and \(\mathbf Z[\zeta_C]\) is the
integer ring. On the initial interval starting at \(Y_*\le x\),
the prime bounds are
\(\log q\le Y_*^{\varepsilon/5}\) and
\(\log d\le Y_*^{\varepsilon/5}(\log Y_*)^2\).
They give the stated bound for \(\log C\), including the dyadic
level and the fixed conductor. The quadratic character of the
imaginary field of discriminant \(-d\) is \(J_d\), since \(d\)
is prime and \(d\equiv7\bmod16\). The formula for \(A_{J_d}\)
is the definition of the analytic scalar, and all its terms are
real. Apply Theorem~\ref{pr:pair} to the trivial base character
and its \(d\)-twist to obtain
\(|L'(1,J_d)/L(1,J_d)|\le4\log x\). Finally
\(x^{-4\varepsilon/5}\log^2(2x)\to0\).
\end{proof}
Consequently, wherever the determinant-degree comparison of
Section~9 identifies the corrected constant as
\[
 c_1=-\tfrac12(p-q')A_{J_d}
                +q'\big(\log\operatorname{rd}(E)-\log d\big),
\]
one has \(|c_1|=o(rL)\) uniformly even for varying signature
\((p,q')\), \(r=p+q'\ge1\). This removes a numerical period bound
from the list of hypotheses to be checked on the actual hosts.
Identifying an integral determinant degree with that scalar is a
geometric assertion.
\src{src/Development/AbelianVarietyTheory/Definitions/CMLattices/ExistsIsThetaDegOfIsLocalModel.lean}{23}{normalization of the analytic CM scalar} and
\src{src/Development/AutomorphicArithmetic/CharacterBoxes/ExistsPageN.lean}{39}{dyadic level at the selected scale}.

\section{Arithmetic geometry of the theta operator}\label{geo:engine}

The arithmetic-geometric argument uses two estimates on the same column pairing. At each finite place a local valuation gives an integral map after a divisor shift. At infinity the chain and its quotient metric give an upper bound on the height of that map. A separate Hermitian-frame argument turns the resulting maps into a lower bound on a target slope. The independent covolume estimate supplies a numerical input to the pairing construction. In the active proof \code{closed\_family\_admissible\_twist\_link\_alt} supplies \code{hA2}; \code{closed\_linkDIVt\_of\_closed\_bridge\_kernel} then supplies \code{hDIV}, which is passed to \code{topDesWin}. These are assembled in \src{src/Development/AutomorphicArithmetic/ThetaCorrespondence/ThetaPairLifts/ThirdPlaceSeries/TopShape.lean}{40}{local estimate and division bridge}.

\noindent For the signed translate mean one may form a labelled direct sum before composing the first edge. Section~\ref{geo:translate} proves the resulting single slope comparison, and Section~\ref{red:translate} substitutes it into the window estimate. \subsection{The operators in the length-\(m\) theta chain}\label{geo:chain}

We use the column-side realization \code{kerRec}, which calls \code{chKerC}, then \code{chKer}. For ordinary edges \code{chKer} takes \code{matK}. For the last two edges it takes the coupled pair \code{cp1K}, \code{cp2K}. The chain is indexed by \(m_2\) ordinary edges followed by those two, with \(m_2+2=m\) by \code{ChIdx.Ok.hm}. \src{src/Formalism/LocalToGlobal/Columns/FiniteSumNonzeroLargeDepth.lean}{993}{exact length} The distribution for the final edge is \code{top}, with the weight at level \(m\). See \src{src/Formalism/LocalToGlobal/Divisibility/DivisionHeightBoundStatement.lean}{943}{active kernel record} and \src{src/Formalism/LocalToGlobal/KernelLattices/ThDistC.lean}{427}{ordinary and coupled kernels}.

Fix an edge \(i\to i+1\), source and target decorated points \(p,p'\) with lattices \(N,N'\), frames \(f,f'\), and indices \(\alpha,\beta\) in their respective level families. Put \(M=\mathrm{homAt}_S(N,N')\), the maps preserving the localized lattices at the specified places, and abbreviate the positive scales of \code{cAof}, \code{cBof} to \(c_A,c_B\). The actual matrix entry is
\begin{equation}\label{geo:chain-entry}
 K_i(p,p')=\sum_{x\in(c_A^{-1}M)/(c_BM)}
              \Xi_i(x)\,\mathcal T_i(N,f,N',f';c_B,x).
\end{equation}
The sum uses classes in the indicated quotient with chosen lifts. The local tests and distribution are descended on the admissible range. Here \(c_A=\prod_P q_P^{A}\), \(c_B=T_u\prod_P q_P^{B}\), and \(S\) consists of the level places and the places used for the row. In the ordinary case \(\Xi_i\) is the level indicator times a product of \(\Xi_{i,P}(f_P'^{-1}xf_P)\). The indicator tests preservation of both the local decorated lattices and their prescribed maps at the level places. Write \(H_P=C_{i+1,P}\cap N_1\), where \(C_{i+1,P}\) is the local \code{Cbe} and \(N_1\) the specified unipotent subgroup. With \(U_P=N_1\cap\operatorname{stab}(\Lambda_{i+1,P})\), the volume ratio used by \code{XiAt} is \(V_{i+1,P}=\mathrm{relIndex}(U_P,H_P)/\mathrm{relIndex}(H_P,U_P)\), with both indices interpreted as in \code{volN}. Let \(d_{i,P}\) denote the number of auxiliary coordinates \code{FC}. \src{src/Formalism/LocalToGlobal/Columns/FiniteSumNonzeroLargeDepth.lean}{542}{relative-index volume} Then the local factor is
\[
 \Xi_{i,P}(z)=V_{i+1,P}^{-1}q_P^{-B d_{i,P}}
   \sum_{\mathbf t\in\mathrm{valsQ}^{d_{i,P}}}
     \Phi_{i,P}(x_1(z,\mathbf t))
       \psi_P(-\mathrm{chArg}(x_1(z,\mathbf t))).
\]
This is a finite weighted Fourier sum. The negative sign and the power of \(q_P^{-B}\) follow by substituting \code{E1at} into \code{XiAt} and expanding \code{E1}. The distribution is zero unless \(c_B\mid c_Ac_B\) and the map belongs to \((c_Ac_B)^{-1}M\). On its support, with \(M_i(N,f)\) the \code{muInvDet} factor for the edge, it is
\[
 \mathcal T_i=
 \frac{M_i(N,f)}{M_i(N',f')}\,
          \mathrm{jetOp}_{i}(J_i).
\]
The coefficients of \(J_i\) are the descents of the complex jet coefficients of the induced theta distribution after multiplication by \(\chi_F^{-1}\). Applied to a polynomial, \(\mathrm{jetOp}_{i}\) first multiplies by the prescribed determinant powers, then substitutes the block variables, contracts their coefficients with \(J_i\), and projects onto the target weight space. This follows by composing the maps in \code{jetMap} and \code{jetOp} in their given order. See \src{src/Formalism/LocalToGlobal/Columns/FiniteSumNonzeroLargeDepth.lean}{636}{lattice quotient sum}, \src{src/Formalism/LocalToGlobal/Columns/FiniteSumNonzeroLargeDepth.lean}{943}{local Fourier normalization}, \src{src/Formalism/LocalToGlobal/Columns/FiniteSumNonzeroLargeDepth.lean}{1110}{scales and local factors}, \src{src/Formalism/LocalToGlobal/KernelLattices/ThDistC.lean}{31}{jet contraction}, and \src{src/Formalism/LocalToGlobal/KernelLattices/ThDistC.lean}{137}{determinant ratio}.

\begin{proposition}[Composition on class representatives]\label{geo:chain-compose}
Let \(F_i\) denote the equivariant forms on level \(i\), with direct summands indexed by its finite label set. Ordinary steps act on these forms by
\[
 (C_i F_i)_\beta(p')=
       \sum_{\alpha}\sum_{c}
          \frac{K_i(p_c,p')\,F_{i,\alpha}(p_c)}{w_c},
\]
where \(p_c\) are representatives of the classes in label \(\alpha\) and \(w_c\) are the source class weights. For the last two steps, the kernel between a point \(p_0\) at level \(m_2\) and a row point \(d\) is
\[
 L(p_0,d)=\sum_c \frac1{w_c}
                \sum_{Y\in\mathcal Y(p_c)}
                   K_{\mathrm{top}}(p_c,Y,d)
                   K_{\mathrm{pre}}(p_0,p_c,Y).
\]
The terminal period on the resulting column \(C(f)=
L\,C_{m_2-1}\cdots C_0\,\mathrm{start}(f)\) and a row \(G\) is
\[
 \mathcal P(C(f),G)
       =\sum_e \frac1{w_e}
            \left\langle \mathrm{br}_u C(f)(d_e),G(d_e)\right\rangle_{K}.
\]
Here \(L\) in the formula for \(C\) denotes the operator of the displayed coupled kernel with its source-class sum, \(d_e\) are row class representatives, and the pairing is the bilinear \code{pairBK}. The map \(\mathrm{br}_u\) is coefficient extraction in the last-column variables followed by projection to the row weight and the scalar
\[
 \prod_\tau t_\tau^{-n_\tau}
        \prod_\tau c_{\mathrm{Br},\tau}\kappa_{\mathrm{Br},\tau},
 \qquad
 t_\tau=\iota_\tau\bigl(\overline{\pi_{T_u}}\bigr).
\]
\end{proposition}
\begin{proof}
For an ordinary step, substitute the class sum \code{ClsKer.appOf} in \code{Mat.opOf}. The normalization is \(w_c^{-1}\) on the source class, and equivariance makes each term independent of the chosen representative. Replacing \(p_c\) by \(g p_c\) contributes \(\rho(g^{-1})\) on the right of the kernel and \(\rho(g)\) on the form. For the coupled pair the two kernel entries use the \code{XiCpl1}, \code{XiCpl2} tests respectively, on the same admissible coset \(Y\). They are zero off the admissible set. Upon changing the intermediate representative, \(Y\) is transported bijectively and the intervening factors \(\rho(g^{-1})\rho(g)\) cancel in the product. Thus there is one intermediate class weight and one sum over \(Y\), rather than two unrelated coset sums. The recursion for \code{colTo} starts at the identity and composes each ordinary step on the left. The \code{colCpl} formula then applies the coupled pair. Finally \code{PtilFF} evaluates the branching map and pairs at the same row representative, with the row class weight. Substituting \code{brAt} and \code{brStdP} gives the displayed scalar. The formula is a dual pairing against the row input. It requires no identification of either of the last two maps with the Hermitian adjoint of the other.
\end{proof}
\noindent
\src{src/Formalism/LocalToGlobal/Columns/RhoEqOfSmulEq.lean}{372}{class action and normalization},
\src{src/Formalism/LocalToGlobal/Columns/RhoEqOfSmulEq.lean}{679}{matrix action},
\src{src/Formalism/LocalToGlobal/Columns/FiniteSumNonzeroLargeDepth.lean}{263}{chain recursion},
\src{src/Formalism/LocalToGlobal/Columns/FiniteSumNonzeroLargeDepth.lean}{298}{coupled contraction},
\src{src/Formalism/LocalToGlobal/Columns/FiniteSumNonzeroLargeDepth.lean}{128}{branching coefficient},
\src{src/Formalism/LocalToGlobal/Columns/FiniteSumNonzeroLargeDepth.lean}{215}{terminal row pairing}.

The active coordinate comparison \code{pairing\_eq\_colPairing} identifies the rational map \(B_k(u)\) with the column pairing. It uses the period identity supplied by the realized column data and its vanishing on the right coordinate kernel. It then evaluates on a right inverse of the surjective quotient-coordinate matrix to prove the equality on every target coordinate. The finite-place comparison \code{perToVal2F} carries the local period estimate to this rational pairing. Evaluating individual lifts would give neither this descended functional nor its accumulated finite-place constraint. In the ordinary local model the valuation is taken on the actual chain end, using the recursion \(C_{[0]}=1,\ C_{[j+1]}=C_j C_{[j]}\) and the source's compensated generator estimate. At the endpoint its scalar budget contains \(e_w\sum_{j=1}^{m}(g_S(p_w,j)+F(j,w)-E(u,j,w))\), while the period cost subtracts the stabilizer contribution once, as in \eqref{geo:budget} and Lemma~\ref{geo:generators}. These are quantitative bounds on the composition, with the model and comparison hypotheses stated there. At \(p_w=T_u\) the per-layer budget instead aggregates to \(-m\,\mathrm{EatT}(u,w)\) before multiplying by \(e_w\). See \src{src/Development/AbelianVarietyTheory/Proofs/IwasawaDecomposition/RationalPairingEqColumnPairing.lean}{60}{active endpoint comparison}, \src{src/Development/AbelianVarietyTheory/Proofs/IwasawaDecomposition/RationalPairingEqColumnPairing.lean}{23}{rational pairing on the quotient}, \src{src/Formalism/LocalToGlobal/GroupLattices/ServJointF.lean}{42}{finite-place comparison}, \src{src/Formalism/SharpRestriction/GrowthBounds/DataGChain.lean}{85}{local chain recursion}, and \src{src/Formalism/SharpRestriction/Divisibility/WeightDegreeSizes.lean}{887}{level-prime length factor}.


\subsection{Finite-place integrality and divisor translation}

Fix a word \(u\) and a finite place \(w\) of the coefficient field \(K\), with residue characteristic \(p_w\) and ramification index \(e_w=v_w(p_w)\). Use the convention \(v_w(0)=+\infty\) and normalize the other values in \(\mathbb Z\). Let \(m\) be the chain length, \(T_u\) the level prime, and \(\Gamma_u\) the positive integer \code{DGam} divisible by every class stabilizer order in each layer \(0\le i\le m\). \src{src/Formalism/LocalRingCardinality/DataKernels.lean}{25}{common multiple of stabilizer orders} Put
\begin{equation}\label{geo:budget}
 d_w=v_{p_w}(\Gamma_u),\qquad
 s_{u,w}=\sum_{j=1}^{m}\bigl(g_D(j,w)+F(j,w)-E(u,j,w)\bigr).
\end{equation}
Here \(g_D(j,w)\) is the recorded gain with the gain at every dropped prime set to zero. The model gain \(g_S(p_w,j)\) uses the positive parts of the orbital gains. In particular an empty set of selected orbits can leave a nonzero contribution from the complementary orbits. The comparison \(g_D\le g_S\) uses the filled parameters and the inclusion of the exceptional \code{IsCs} primes among the dropped primes. At a dropped prime or a place outside the gain support the recorded gain vanishes. Failed inertia tests and the two generator primes give zero by the local gain equations. In the remaining branch the filled parameters identify the recorded gain with the model before the exceptional-prime mask, and the exceptional branch is ruled out by the drop condition. The nonnegativity of \(g_S\) then completes the comparison. See \src{src/Formalism/SharpRestriction/GrowthBounds/WeilJacquetShalikaLaw.lean}{538}{comparison of recorded and model gains}.

\begin{lemma}[Generator and period budgets]\label{geo:generators}
Let \(A\) be a discrete valuation ring with fraction field \(L\), normalized valuation \(v\), and let \(M\) be an \(A\)-submodule of an \(L\)-vector space. Let \(C\) be an \(L\)-linear chain map. Suppose an input lattice is generated over \(A\) by \(a_i x_i\), with \(C x_i\in \pi^{b_i} M\), \(b_i\in\mathbb Z\), and \(v(a_i)+b_i\ge b\) for every \(i\), where \(b\in\mathbb Z\). Then every input \(f\) satisfies \(C f\in\pi^b M\). If a terminal \(L\)-linear period \(\mathcal P\) satisfies \(\mathcal P(M)\subset\pi^{-d}A\) for an integer \(d\), then \(v(\mathcal P(Cf))\ge b-d\).
\end{lemma}
\begin{proof}
For each generator, \(a_i\pi^{b_i-b}\in A\). Hence \(\pi^{-b}C(a_i x_i)\in M\). The coefficients of any finite linear combination of generators belong to \(A\). Applying \(\pi^{-b}C\) to such a combination therefore gives an element of \(M\), including for the zero input or a cancelling combination. Write \(C f=\pi^b z\) with \(z\in M\). Then \(\pi^{d-b}\mathcal P(C f)=\pi^d\mathcal P(z)\in A\). Negative values of \(b\) and \(d\) are allowed, since the powers are taken in \(L\).
\end{proof}

In the active ordinary-place branch the twisted lattice is generated by products \(a_{\mathcal V'}F_0\). Here \(\mathcal V'\) ranges over subsets of the selected orbits \(\mathcal V\). The complement contributes the \(p_w\)-power budget
\[
 \sum_{t\in\mathcal V\setminus\mathcal V'}
       \max\bigl(0,\lceil \mathrm{Sv}(t)\rceil\bigr).
\]
The contract \code{unitBoundAtTwistedPlaces} bounds the chain on \(F_0\). The orbit inequality \code{bg\_gain\_le} converts the generator-specific bound after compensation into the uniform bound. The Lean proof writes membership using the uniformizer and applies span induction, precisely the mechanism of Lemma~\ref{geo:generators}. It starts from a chain-end estimate on specified generators. The terminal column comparison includes the ordinary and twisted coordinates. When the selected-model branch is unavailable, the comparison instead uses the supplied unselected-orbit lattice bound. See \src{src/Formalism/SharpRestriction/GrowthBounds/AdmissibleIntegers/StatementsOfRemainingBound.lean}{111}{generator compensation and span induction} and \src{src/Formalism/SharpRestriction/GrowthBounds/AdmissibleIntegers/StatementsOfRemainingBound.lean}{162}{ordinary-place model alternatives}.

The resulting column has scalar valuation at least
\[
 e_w\left(d_w+\sum_{j=1}^m(g_S(p_w,j)+F(j,w)-E(u,j,w))\right).
\]
The terminal period on the remaining vector costs at most \(e_w d_w\). Indeed its two lattice-coordinate versions use the integrality at the ordinary or twisted points, respectively, and the loss from the same common multiple. The valuation of the product is therefore at least \(e_w s_{u,w}\), using \(g_D\le g_S\). The loss \(e_w d_w\) occurs once. This is the multiplication in \src{src/Formalism/SharpRestriction/GrowthBounds/WeilJacquetShalikaLaw.lean}{401}{ordinary column and terminal period}. At \(p_w=T_u\) the level-lattice calculation yields the budget \(-m\,\mathrm{EatT}(u,w)\) before multiplication by \(e_w\), and \(d_w=0\) by coprimality. At a generator prime distinct from \(T_u\), one multiplies the test vector by a local unit to bring it into the lattice on which the local bound holds. Dividing by that unit preserves its valuation. These branches and the final integer bound \(\lfloor e_ws_{u,w}\rfloor\) occur in \src{src/Formalism/LocalToGlobal/GroupLattices/AdmissibleAtGeneratorPlaces.lean}{214}{three finite-place cases}. The family version takes a common threshold for the inputs before supplying \code{adm2}. See \src{src/Development/AutomorphicArithmetic/ThetaCorrespondence/ThetaPairLifts/ThirdPlaceSeries/ClosedFamilyAdmissibleTwistLink.lean}{24}{family estimate}. The local generator estimates and the ordinary and twisted comparison maps are the geometric premises in this step.

\begin{proposition}[Integral shifted dual map]\label{geo:shift}
Let \(B:N_K\times R_K\to K\) be bilinear on the scalar extensions of full \(\mathcal O_K\)-lattices \(N,R\). At every finite \(w\), suppose the available period coordinates \(A_w\)-span \(R_w\), where \(A_w=\mathcal O_{K,w}\), and the period equals the evaluation of \(B\) at those coordinates up to a unit. Assume the valuation of the period for every element of \(N_w\) and every such coordinate is at least \(\sigma_w\in\mathbb R\). If \(\tau_w\le\lceil\sigma_w\rceil\) are integers of finite support, the shifted source lattice with local model \(\pi_w^{-\tau_w}N_w\) maps into \(R^\vee\), the integral linear dual of \(R\).
\end{proposition}
\begin{proof}
An integral valuation at least \(\sigma_w\) is at least \(\lceil\sigma_w\rceil\). The ideal \(\pi_w^{\lceil\sigma_w\rceil}A_w\) is an \(A_w\)-module, so comparison by a unit and extension from the spanning coordinates preserve this conclusion. For \(f\in\pi_w^{-\tau_w}N_w\) and \(r\in R_w\), bilinearity now gives \(v_w(B(f,r))\ge\lceil\sigma_w\rceil-\tau_w\ge0\). For a global element of the shifted source and a global element of \(R\), the value thus belongs to every \(A_w\). The intersection of the \(A_w\) inside \(K\) is \(\mathcal O_K\), since a principal fractional ideal with no negative prime exponent is integral. This proves the dual-lattice inclusion.
\end{proof}

The coordinate comparison and spanning premises are supplied as part of the division package. The actual construction introduces them in \src{src/Formalism/LocalToGlobal/Divisibility/RecBridge.lean}{633}{feeding the period estimate into the division bridge}. See also \src{src/Formalism/LocalToGlobal/GroupLattices/ServJointF.lean}{84}{local column comparison} and \src{src/Formalism/QuotientMetricDivisors/HermitianTwistDivisors/TwistedLatticeNaturality.lean}{720}{integrality after the shift}.

\begin{corollary}[Determinant cost of the shift]\label{geo:shift-index}
Suppose the lattices in Proposition~\ref{geo:shift} have a fixed nonsingular infinite-place metric, with \(K\) totally imaginary and \(N\) of rank \(d\). Write \(N'\) for the shifted source and \(l_w=\log|\mathcal O_K/\mathfrak w|/[K:\mathbb Q]\). Then
\[
 \mathrm{lcv}(N')-\mathrm{lcv}(N)=-2d\sum_w l_w\tau_w,\qquad
 \mu_{\rm def}(N')-\mu_{\rm def}(N)=2\sum_w l_w\tau_w\quad(d>0).
\]
For a rational prime \(p\), \(\sum_w l_w v_w(p)=\log p\).
\end{corollary}
\begin{proof}
At a finite place, take an \(A_w\)-basis of \(N_w\). The shifted basis is obtained by multiplying each of its \(d\) vectors by \(\pi_w^{-\tau_w}\), so the determinant has valuation \(-d\tau_w\), including the zero-rank determinant \(1\). The fixed-metric comparison of log-covolumes is twice the sum of determinant valuations with weights \(l_w\). Its factor two is part of the normalization in \src{src/Formalism/QuotientMetricDivisors/HermitianTwistDivisors/TwistedLatticeNaturality.lean}{145}{determinant comparison and slope normalization}. For \(d>0\), \(\mu_{\rm def}=-\mathrm{lcv}/d\), giving the stated sign. Finally \(v_w(p)=0\) away from the places over \(p\), where \(v_w(p)=e_w\) and \(\log|\mathcal O_K/\mathfrak w|=f_w\log p\). The degree formula \(\sum_{w\mid p}e_w f_w=[K:\mathbb Q]\) yields the rational-prime identity. Equivalently, specialize the active valuative product formula to \(p\). All complex embeddings have modulus \(p\), and their squared normalization uses \(2\,\#S_\infty=[K:\mathbb Q]\). \src{src/Formalism/QuotientMetricDivisors/HermitianTwistDivisors/ProdFormula.lean}{39}{valuative product formula}
\end{proof}


\subsection{Archimedean coercivity and slope transfer}

\begin{lemma}[Height of a descended pairing]\label{geo:descent}
Suppose a Hermitian quotient used by the right input of \(B_u\) has a linear section taking that input to a lift of norm at most \(\beta^{-1/2}\) times its norm, with \(\beta>0\). Suppose the chain maps \(C_j\) have operator norms squared at most \(e^{k_j}\), and the terminal pairing has operator norm squared at most \(e^{h_P}\). Then
\[
 \|B_u\|^2\le e^{h_P+\sum_j k_j}/\beta.
\]
In particular, if \(\beta=\kappa\min_a\prod_j\lambda_{j,a}\) with \(\kappa>0\) and all the \(\lambda_{j,a}>0\), the logarithm for a nonzero \(B_u\) is bounded by \(h_P+\sum_j k_j-\log\kappa-\min_a\sum_j\log\lambda_{j,a}\).
\end{lemma}
\begin{proof}
Apply the quotient section to the right input and the chain to the left input. Successive operator-norm bounds give \(\|C_m\cdots C_1 x\|^2\le e^{\sum_j k_j}\|x\|^2\). Applying the terminal pairing and then the section bound proves the claimed norm inequality for all input pairs. If a Hermitian Gram operator \(G\) controls the section norm through \(\|y\|^2=\langle G\widetilde y,\widetilde y\rangle\) on the lifting complement and all its eigenvalues are at least \(\beta\), an orthonormal eigenbasis gives \(\beta\|\widetilde y\|^2\le\|y\|^2\), including \(\widetilde y=0\). In the application the product lower bound for \(\beta\) comes from the common block description of the quotient Gram operator. The zero map satisfies the norm inequality and needs no logarithm.
\end{proof}

Averaging the logarithmic estimate over infinite places yields \(\mathrm{htMap}(B_u)\le h(u)\) for the nonzero maps with the chosen word height budget. The normalization uses the mean of \(\log\|B_u\|_v^2\). See \src{src/Formalism/QuotientMetricDivisors/HermitianTwistDivisors/DescendedNorm.lean}{823}{chain in the quotient metric} and \src{src/Formalism/LocalToGlobal/Divisibility/ConstantGeTwentyEight.lean}{123}{average height}. The companion estimate on target slope plus word height also needs the border-minor bound on the base side and the row-to-column comparison in \code{P.extra.passage}. These inputs are separate from the bound on the operator norm. They are combined in \src{src/Formalism/LocalToGlobal/Divisibility/ConstantGeTwentyEight.lean}{210}{target-slope budget}. The actual bridge packages the integral inclusion, the height bound for nonzero maps, and this combined estimate as \code{WordFacts}. See \src{src/Formalism/LocalToGlobal/Divisibility/RecBridge.lean}{128}{wordwise division data}.

\begin{lemma}[Coercivity of a nearly tight frame]\label{geo:coercivity}
Let \(A_i:V_i\to \mathcal H\) be a finite family of maps of finite-dimensional complex Hilbert spaces. Put \(r=\sum_i\dim V_i>0\), \(D=\dim\mathcal H\), and \(Q=\sum_i A_iA_i^*\). Suppose
\[
 \sum_i\|A_i^*A_i-I_{V_i}\|_{\mathrm{HS}}^2\le er,\quad
 0\le e\le\tfrac14,\qquad
 \operatorname{tr}(Q^2)\le (1+\eta)(\operatorname{tr}Q)^2/D,\quad \eta\ge0.
\]
The second inequality is imposed on the nonzero ambient space. For subspaces \(U_i\subseteq V_i\), write \(q=\sum_i\operatorname{codim}_{V_i}U_i/r\). Then
\[
 \frac{\operatorname{codim}_{\mathcal H}\sum_i A_iU_i}{D}
       \le\eta+4q+4\sqrt e.
\]
Positive frame weights are allowed by putting their square roots into \(A_i\).
\end{lemma}
\begin{proof}
Write \(E_i=A_i^*A_i-I\). Then
\[
 \operatorname{tr}Q=r+\sum_i\operatorname{tr}E_i
 \ge r-\sqrt{r\sum_i\|E_i\|_{\mathrm{HS}}^2}
 \ge r(1-\sqrt e).
\]
This also rules out a zero ambient space using just the first hypothesis. Let \(P\) project onto \(Z=(\sum_i A_iU_i)^\perp\) and set \(d=\dim Z\). Since \(P A_i\) kills \(U_i\), if \(S_i\) projects onto \(U_i^\perp\) then
\[
 \operatorname{tr}(PQ)\le\sum_i\operatorname{tr}(S_iA_i^*A_i)
 \le rq+\sqrt{rq\sum_i\|E_i\|_{\mathrm{HS}}^2}
 \le r(q+\sqrt{qe}).
\]
For the first inequality, take an orthonormal basis adapted to \(U_i\) and bound the squared norms of the projected images by the squared norms of the images. With \(M=\operatorname{tr}Q>0\) the second frame hypothesis gives \(\|Q-(M/D)I\|_{\mathrm{HS}}^2\le\eta M^2/D\). Pairing with \(P\) yields \(\operatorname{tr}(PQ)/M\ge z-\sqrt{\eta z}\) for \(z=d/D\). Since \(2\sqrt{\eta z}\le\eta+z\), this lower bound is at least \((z-\eta)/2\). The upper bound is at most \(2(q+\sqrt e)\), using \(q\le1\) and \(1-\sqrt e\ge 1/2\). The conclusion follows. The proof allows \(\dim V_i=0\), \(Z=0\), \(q=0\), and \(e=0\).
\end{proof}

The active linear-algebra estimate uses the trace of the complementary projection in \src{src/Formalism/HarderNarasimhan/HarderNarasimhanProofs/FrameLinearAlgebraFirst.lean}{149}{complementary trace bounds}. In its compressed version \(\mathcal H=\mathcal H_0\oplus\mathcal H_1\), where the target is \(\mathcal H_0\) of dimension \(D_0\ge(1-\delta)D\). If \(M_0\) is the projection to \(\mathcal H_0\) of \(M=\sum_i A_iU_i\), then \(\mathcal H/M\) surjects onto \(\mathcal H_0/M_0\). Division by \(D_0\) gives the preceding relative deficit divided by \(1-\delta\), with \(D_0>0\) since \(0\le\delta<1\). The main application supplies the two frame parameters by compressing the full matrices along isometric embeddings. It substitutes \(\eta=5\sigma+3\sqrt{d_{\rm rk}}\) and \(e=2e_0\) after the rank and Hilbert--Schmidt estimates. See \src{src/Formalism/SizeWindowBrackets/ChainSlopeWindows/JointDModelAppliesUnitLattice.lean}{164}{frame compression estimates} and \src{src/Formalism/SharpRestriction/GlobalSchurBounds/TopDensityNGeSixteen.lean}{538}{actual frame and rank data}.

\begin{theorem}[Finite-layer HN estimate]\label{geo:hn}
Use normalized Hermitian slopes \(\lambda=2\lambda_{\rm Euc}-\log 2+\log|\operatorname{disc}K|/[K:\mathbb Q]\) for a totally imaginary \(K\). Suppose a finite family of lattice maps \(B_i:N_i\to R_*\) of height at most \(h_i\) for nonzero maps has positive total source rank. Fix \(0\le\delta<1\). Assume the target lattice has positive rank \(D_0\) over \(K\), and that at a selected infinite place its Hermitian space identifies with a subspace \(\mathcal H_0\subseteq\mathcal H\) of dimension \(D_0\ge(1-\delta)D\) in the frame of Lemma~\ref{geo:coercivity}. Assume the matrices of the lattice maps at that place are the orthogonal projections of the frame maps to \(\mathcal H_0\), up to multiplication of each matrix by a nonzero scalar. Let \(s\) be the rank-weighted mean of all source slopes after subtracting \(h_i\), and assume these slopes are at most \(s+u\) with \(u\ge0\). Assume every target slope is at least \(s-H\), \(H\ge0\). Then
\begin{equation}\label{geo:hn-result}
 s-\mu(R_*)\le
 \frac{4u\log(1+H/u)+(\eta+4\sqrt e)H}{1-\delta}.
\end{equation}
The first term in the numerator is defined as zero when \(u=0\).
\end{theorem}
\begin{proof}
For \(t\in[0,H]\), truncate each source at adjusted slope \(s-t\). The slope inequality for a morphism of Hermitian lattices gives, for a nonzero image of the corresponding source truncation, minimum slope at least \(s-t+h_i-\mathrm{htMap}(B_i)\ge s-t\). It uses the image lattice with the induced target metric. The image is therefore contained in the target high-slope truncation by HN functoriality at the chosen threshold. For the zero image containment holds directly. The normalized slopes and height account for the factor of two in the Euclidean scale. The slope inequality used here is \src{src/Formalism/HarderNarasimhan/HarderNarasimhanProofs/SlopeIneqMean.lean}{249}{image slope under a height bound}. The truncation containment, including the zero branch, occurs in \src{src/Formalism/HarderNarasimhan/HarderNarasimhanProofs/TheoremFSharp.lean}{277}{images of high-slope lattices}.

Let \(q_t\) be the source fraction below \(s-t\). Since the adjusted source mean equals \(s\),
\[
 s\le q_t(s-t)+(1-q_t)(s+u).
\]
For \(u+t>0\) this gives \(q_t\le u/(u+t)\). When \(u=0\), the mean condition and the common upper bound imply every source slope equals \(s\), so \(q_t=0\) even at \(t=0\). The target fraction strictly below \(s-t\) is no greater than the codimension of the sum of the projected truncated images divided by \(D_0\). The preceding frame bounds thus majorize this fraction by \((\eta+4q_t+4\sqrt e)/(1-\delta)\).

Each target slope \(\lambda\) satisfies \((s-\lambda)_+=\int_0^H\mathbf1_{\lambda<s-t}\,dt\). Equality at a threshold changes no integral. Summing this finite identity gives the layer-cake bound on \(s-\mu(R_*)\). For \(u>0\), the integral of \(u/(u+t)\) is \(u\log(1+H/u)\). For \(u=0\) its replacement \(q_t\) is identically zero. At \(H=0\) the target minimum hypothesis already gives \(s-\mu(R_*)\le0\). These are the rank-count and integration steps in \src{src/Formalism/HarderNarasimhan/HarderNarasimhanProofs/TheoremFSharp.lean}{308}{frame rank deficit} and \src{src/Formalism/HarderNarasimhan/HarderNarasimhanProofs/TheoremFSharp.lean}{338}{finite-slope layer cake}.
\end{proof}

The application takes the weakened bound with an added \(\log|I|\) for the nonempty index set of words. The division estimate bounds \(s\) below by the target slope plus the gain, minus the numerical losses. Applying Theorem~\ref{geo:hn} bounds the same \(s\) above by the target slope plus the frame and counting losses. Subtracting eliminates the target slope, giving the one-sided gain estimate entering \(\mathrm{ASK}\). The height and slope transfer is \src{src/Formalism/SharpRestriction/GlobalSchurBounds/QuadraticLogDerivCongr.lean}{699}{conversion to the actual frame data}. See \src{src/Formalism/SizeWindowBrackets/ChainSlopeWindows/LevelOfCongruence.lean}{50}{lower bound on the source mean} and \src{src/Formalism/SizeWindowBrackets/ChainSlopeWindows/LevelOfCongruence.lean}{113}{slope elimination}.

\subsection{Matching theta integrals and constructing the evaluation lattice}

In the tensor Weil model for an admissible Hermitian triple \(V,V',W\), write
\[
 \Theta_\omega(g,h;\Phi)=\operatorname{thetaSum}_F(\omega(g,h)\Phi),
 \qquad I_V(\Phi)(h)=
 \int_{\mathcal F_V}\Theta_\omega(g,h;\Phi)\,d\mu_V(g).
\]
Here \(h\in U(W)(\mathbb A_F)\) and \(\mathcal F_V\) is a fundamental domain for \(U(V)(F)\backslash U(V)(\mathbb A_F)\). The corresponding domain and measure on the other side are those of its Weil model. In particular the global comparison constant depends on the fixed matching maps and model measures. The volume-one normalization below applies to the unipotent constant term. The definitions are in \src{src/Development/AbelianVarietyTheory/BlockTheory/DivisorBase/WeilTensorModels.lean}{286}{theta series and integrals} and \src{src/Development/AbelianVarietyTheory/Definitions/Divisors/CTSmallCone.lean}{82}{normalized constant term}.

At good finite places the two spherical vectors have the same nonzero image in the common quotient. For a finite set \(S\) containing the bad places, polynomial matching of pure tensors identifies the tensor of the real image with the images in the finite-place quotients:
\[
 R_\infty(\phi_\infty)\otimes\bigotimes_{v\in S}R_v(\phi_v)
 =
 R'_\infty(\phi'_\infty)\otimes\bigotimes_{v\in S}R'_v(\phi'_v).
\]
The condition extends linearly to finite sums. At good places the relevant kernel is the span of the differences \(\omega(g,1)\phi-\phi\), so a linear functional invariant under \(g\) factors through the quotient because the functional vanishes on each difference and hence on the kernel. The general matching condition is a limit of polynomial matching, with a support condition on the right in the split matrix model that forces vanishing on rank-deficient matrices. See \src{src/Development/AbelianVarietyTheory/BlockTheory/DivisorBase/WeilTensorModels.lean}{459}{polynomial matching and matched limits}.

\begin{proposition}[Proportionality from the matching inputs]\label{geo:theta}
In the tensor Weil models, fix real and finite matching maps satisfying \code{IsMatched}. Assume the anisotropic bottom comparison, the constant-term descent and reconstruction identities, nonvanishing of the descent factors, a nonzero test at the bottom, and continuity along the matched approximations. Assume also that the right-hand polynomial test obtained by nonvanishing admits a left-hand match. Then there is a unique \(c\in\mathbb C^\times\) such that \(I_V(\Phi)(h)=cI_{V'}(\Phi')(h)\) for every matched pair and every \(h\), and some matched pair has nonzero right-hand integral.
\end{proposition}
\begin{proof}
First use polynomial vectors and induct on the sum \(p+q\) of the two Witt indices in the comparison data. The bottom comparison supplies the step \(p=q=0\). When \(q>0\), remove an adapted hyperbolic plane in the primed tower and the corresponding plane on the target. If \(N\) is the resulting unipotent group, take the constant term \(\int_{\mathcal F_N} f(nh)\,d\mu_N(n)\), where \(\mathcal F_N\) represents \(N(F)\backslash N(\mathbb A_F)\) and has measure one. Its two summands are identified by the specified Levi transformation rules. For \(a'=b+2\), set \(k=\lfloor(a-(b+2))/2\rfloor\) using natural subtraction, write \(c_1\) for the comparison constant with Witt index sum lower by one, and set \(F_0=\mathrm{stepFac}(-k)\), \(L_0=L^{S_0}(\nu,k+1)\) in the supplied descent identities. With the nonzero local constant \(C\), the compatibility identity to feed the reconstruction is \(cL_0=c_1/(CF_0)\). Indeed \(c=c_1/(CF_0L_0)\) exists and is nonzero, and multiplying by the nonzero product verifies the identity. In the branch \(q=0<p\), the unprimed descent gives the comparison in the other direction. The specified uniqueness of the constant-term splitting identifies the first component with the lower-tower component. The nonzero descent factors transport the lower comparison to the upper reconstruction. This is the induction implemented in \src{src/Development/WeilBound/ThetaLifts/ConeHypotheses/CtSmallCone.lean}{88}{two constant-term descents}. The explicit primed coefficient is constructed at \src{src/Development/WeilBound/ThetaLifts/ConeHypotheses/CtSmallCone.lean}{115}{nonzero denominator in the primed tower}.

In the positive primed Witt index case descent gives a nonzero constant-term component. If its theta integral were zero, the zero pair would also split the constant term, contradicting uniqueness of splitting. For \(q=0\) use the supplied bottom nonvanishing. Choose the resulting right-hand polynomial test \(\Phi'_0\) and a point \(h_0\) with \(I_{V'}(\Phi'_0)(h_0)\ne0\). By the matching-existence hypothesis choose a left-hand match \(\Phi_0\) for this particular test. Continuity extends the polynomial comparison to every matched limit with the same \(c\). If \(c'\) is another comparison constant, subtract the two identities for \((\Phi_0,\Phi'_0,h_0)\). This gives \((c-c')I_{V'}(\Phi'_0)(h_0)=0\), hence \(c=c'\). \src{src/Development/WeilBound/ThetaLifts/ConeHypotheses/CtSmallCone.lean}{138}{Nonzero test integral}, \src{src/Development/WeilBound/ThetaLifts/ConeHypotheses/CtSmallCone.lean}{179}{passage to the matched limit} and \src{src/Development/WeilBound/ThetaLifts/ConeHypotheses/CtSmallCone.lean}{192}{matching the nonzero right-hand test}.
\end{proof}

The real and finite matching maps come from the spherical closure with real Lie-algebra lifting and the split, inert, and ramified local models. The strong identities used in the preceding proof are filled, in particular, by \code{eisBL} and the finite standard Siegel--Weil formula at the end of the spherical reduction. The proof above retains those identities and their nonvanishing hypotheses at their supplied scope, including the specific descent value of the \(L\)-factor. The actual assembly is in \src{src/Development/WeilBound/ThetaLifts/Reductions/FiniteClosure/TheoremKOfMod.lean}{14}{local models in the spherical closure}, \src{src/Development/WeilBound/ThetaLifts/SphericalLattices/SphericalClosure/CloseSph.lean}{188}{terminal comparison}, and \src{src/Development/AutomorphicArithmetic/ThetaCorrespondence/ThetaPairLifts/ThirdPlaceSeries/EveryCharacterOfLiteratureOfVWin.lean}{561}{supply of the two identities}. At the main consumer the integral comparison and local matching define the evaluation grid. Their transport to the square model is used there. See \src{src/Development/AutomorphicArithmetic/ThetaCorrespondence/ThetaPairLifts/ThirdPlaceSeries/EveryCharacterOfLiteratureOfVWin.lean}{810}{active theta consumer} and \src{src/Development/AutomorphicArithmetic/ThetaCorrespondence/ThetaPairLifts/ThirdPlaceSeries/SiegelEisensteinSecondTermHypotheses.lean}{50}{grid and associated bundles}.

\subsection{Component Hilbert--Samuel estimates and the final covolume}

For an admissible weight on a good deep cover write \(h_Z\) for the height of a component \(Z\), \(d_Z>0\) for its geometric degree, \(T\) for the section degree parameter and \(D_f\ge1\) for the flag dimension. The component Hilbert--Samuel input is \code{HSClaimS} with constant \(26\). It provides a lower degree bound through the auxiliary section parameters \(n,\Lambda,\beta,\Lambda_D,\Lambda_T,J\). The active numerical theorem arranges this as a lower bound on the degree in the associated bundle, with
\[
 B_Z=\binom{T-1}{D_f}d_Z,\quad
 z=7(D_f^2/T+D_f\rho),\quad
 A=26(26D_f^2+2zT),\quad
 R=4000b+2c.
\]
Here \(c\ge0\) is the cost per rank of the bundle transfer. Uniform choices are taken for the same point and stage across the components. See \src{src/Development/AutomorphicArithmetic/ThetaCorrespondence/ThetaPairLifts/ThirdPlaceSeries/HilbertSamuelLowerBoundOfComponents.lean}{89}{uniform choices and component input}.

\begin{lemma}[Rank conversion in the component degree estimate]\label{geo:component}
Let \(h\in\mathbb R\), \(b,B,d,c\) be nonnegative, \(0\le z\le1/2\), and \((1-z)d\le B\). Suppose the component estimate, after rearranging the Hilbert--Samuel ratio, gives
\[
 \deg\ge B\bigl(Th-3(D_f+1)|h|-A\log(2+|h|)-J_0\bigr),
 \qquad J_0\le4000b,
\]
and \(\deg_M\ge\deg-dc\). Then
\[
 B\bigl(Th-3(D_f+1)|h|-A\log(2+|h|)-(4000b+2c)\bigr)
 \le\deg_M .
\]
\end{lemma}
\begin{proof}
The rank comparison and \(z\le1/2\) give \(d\le2B\). Since \(c\ge0\), the cost \(dc\) is at most \(2Bc\). Since \(B\ge0\), replacing \(J_0\) by \(4000b\) only decreases the degree lower bound. Combining the two comparisons proves the result, without multiplying any inequality by a possibly negative height. The argument includes \(B=0\), in which case \(d=0\).
\end{proof}

The numerical discharge in \src{src/Development/AbelianVarietyTheory/Proofs/DivisorProofs/DegBoundOfClaimNumbers.lean}{9}{exact numerical hypotheses} requires \(1\le P_0\), \(1\le b\), \(TP_0\le\sqrt b\), \(\log b\le P_0\), and \(0\le\rho\) in addition to the rank comparison. Squaring the nonnegative sides yields \((TP_0)^2\le b\), while the formula for \(z\) gives \(D_f^2\le T/14\). The theorem combines these with the logarithmic bounds for the section parameters to obtain \(J_0\le4000b\). It applies the rank conversion of Lemma~\ref{geo:component} at \src{src/Development/AbelianVarietyTheory/Proofs/DivisorProofs/DegBoundOfClaimNumbers.lean}{121}{last remainder and transfer bounds}. The analytic estimate producing the component \code{HSClaimS} and its auxiliary section parameters is a premise at this stage.

An independent height comparison supplies \(c_s,c_c,X\ge0\) and real component slopes \(\ell_Z\) satisfying
\[
 0\le h_Z+c_s,\qquad
 (1-\rho)\max(0,\ell_Z+c_s)\le h_Z+c_s+c_c,\qquad
 -X\le \frac{\sum_Z\ell_Z}{\#\{Z\}}.
\]
The quantities \(TX\) and \(Tc_c\) and the prescribed multiple
\(c_s(6D_f+10+T(z+2\rho))\) are each at most \(C_1 E_P\).
The short-section Bergman kernel produces the first two comparisons, and
\code{MeanAll} supplies the average comparison. The remaining hypotheses
are positive component degrees comparable to a common degree, the rank
shares, and the inequality
\[
 2\frac{\sum_Z\deg_C(Z)}{\sum_Z\mathrm{rk}_C(Z)}
 \le \mu_{\rm cover}.
\]
The component bound above contributes the uniform budget
\(A\log\max(A,1)+R\le C_2E_P\).
With these inputs, the covering-slope result used here is
\(\mu_{\rm cover}\ge-2(3C_1+C_2+4)E_P\). The subsequent descent moves this
bound to the bundles used on the grid. The geometric comparison and its
numerical aggregation are used at the stated scope. See \src{src/Development/AutomorphicArithmetic/ThetaCorrespondence/ThetaPairLifts/ThirdPlaceSeries/HeightLowerBoundGeneral.lean}{21}{separate height comparison}, \src{src/Development/AutomorphicArithmetic/ThetaCorrespondence/ThetaPairLifts/ThirdPlaceSeries/ULowOf.lean}{16}{covering slope estimate}, and \src{src/Development/AutomorphicArithmetic/ThetaCorrespondence/ThetaPairLifts/ThirdPlaceSeries/TopCovolG.lean}{49}{deep cover and descent}.

\begin{proposition}[Uniform exponential absorption]\label{geo:covolume}
Fix \(a>0\), \(A_0,B_0\in\mathbb R\), and \(C_1,C_2\ge0\). Suppose the pointwise Gram, recurrence, cut, bad-set, lattice and size bounds yield
\[
 \mu_{\rm type}\le\mu_{\rm good}
 +12e^{-x^a}(x^{A_0}+x^{B_0}+C_1+C_2+1)E_P
\]
with \(E_P\ge0\). Suppose the core comparison and low-slope bound give \(\mu_{\rm good}\le \mathrm{mainT}+(C_1+C_2)E_P\), and \(\mu_{\rm type}\) identifies with the normalized log-covolume of the selected lattice. Then for all sufficiently large \(x\) the normalized log-covolume is at most \(\mathrm{mainT}+(C_1+C_2+1)E_P\), uniformly over the points and stages satisfying the hypotheses with these constants.
\end{proposition}
\begin{proof}
Set \(y=x^a\). For each exponent \(A_0/a,B_0/a\) choose a positive integer \(n\) larger than both exponents and zero. For \(y\ge1\), \(e^y\ge y^n/n!\). Each \(e^{-y}y^{A_0/a}\), \(e^{-y}y^{B_0/a}\), and \(e^{-y}\) tends to zero. The coefficient of \(E_P\) in the first inequality is therefore at most one after a threshold depending only on the fixed constants. Multiply this estimate by \(E_P\ge0\) and substitute the upper bound on \(\mu_{\rm good}\). The case \(E_P=0\) has the stated conclusion at every point satisfying the premises.
\end{proof}

In the actual construction \(a=e_z/7\) and the pointwise inequalities are those proved and combined in \src{src/Development/AbelianVarietyTheory/Proofs/DivisorProofs/CovolumeFromGramAndBounds.lean}{39}{covolume from the cut and slope data}. The family argument takes the maximum of this numerical threshold and the geometric thresholds, with \(C_1+C_2+1\) as the minimum error constant. The \code{hX} used at the main entry is obtained through the Chow-height and deep-cover inputs in \src{src/Development/AutomorphicArithmetic/ThetaCorrespondence/ThetaPairLifts/ThirdPlaceSeries/TopCovolSlim2Use.lean}{16}{covolume input on the actual family}. It is applied to the associated bundles on the evaluation grid to give \code{L3}, the inequality \(\mathrm{lcv}(X_{W,i},M_i)/D(P)\le N.l3_i\). See \src{src/Development/AutomorphicArithmetic/ThetaCorrespondence/ThetaPairLifts/ThirdPlaceSeries/LinkedBundlesL3OfOpen.lean}{31}{putting the covolume bound into L3}. Thus theta matching is used to realize the grid, while the Hilbert--Samuel estimates supply the quantitative slope information for the lattices selected on it.

\subsection{One slope comparison over the translate group}\label{geo:translate}

Fix one good scale and one choice of the data of Formal input~\ref{red:design}.
Write \(\mathbb E a=|A_0|^{-1}\sum_{u\in A_0}a_u\). The translate \(u\)
and a word \(i\) at \(u\) are different indices. In a presentation with
several points above \(u\), select one lift for each \(u\) and use it
at every stage. We make this choice throughout. A probability mixture
of lifts gives the same proof if each lift's two rank measures are
multiplied by its conditional probability given \(u\).
The tied weights, the integers \(a_t\) and the centre are the ones on
these data. An identification of the points with translates in the
character group supplies no identification of their lattices.

The direct sum can be formed at the first edge. Retain the label \(u\)
in every factor and compose maps of the same word. For composable
layers \(C_{u,j}\), the identity
\[
 (\bigoplus_u C_{u,m})\cdots(\bigoplus_u C_{u,1})
     =\bigoplus_u(C_{u,m}\cdots C_{u,1})
\]
follows by evaluation on each summand. Over different coefficient fields
this is a tuple of maps over the product of those fields. The last two
kernels of the theta chain use the same intermediate mark \(Y\) in
each path. Treat their coupled sum as a single map whenever only that
sum descends to the quotient. If the fibres are identified, a diagonal
inclusion and a mean projection at the two \emph{ends} give the mean
of the composites. Averaging the edges separately can change the
answer even for integral isometries. On two rank-one blocks take \(+1\)
and \(-1\) at each of two successive edges in the same order. The
average composite is \(1\), whereas both average edges are zero.

For integral matrices in common bases, \(h^{-1}\sum_{u}B_u\),
\(h=|A_0|\), guarantees only the minimum entry valuation among the
summands minus \(v_w(h)\). On the present host \(h\) is a power of
two. Swapping the coordinates of \(\mathbb Z_2^2\) gives the
invariant projector \(e_1\mapsto(1/2,1/2)\). If a terminal integral
quotient \(q\) satisfies \(q\rho(u)=q\) for an integral group action
\(\rho\), then \(q\sum_u\rho(u)=hq\), so the division cancels there.
The construction below needs only the block maps. It divides real
degrees and traces by \(h\). It also keeps the weights inside the
local prime gains. For example, averaging \((2-w)_+\) at \(w=1,3\)
gives \(1/2\), while substitution of the average weight gives zero.

For word \(i\) write \(B_{ui}:S_{ui,K_u}\to R_{u,K_u}\) for the
rational descended column in the prescribed divisor shifts.
Here \(S_{ui},R_u\) are Hermitian lattices over a totally imaginary
field \(K_u\). Omit zero-rank words and put \(r_{ui}=\operatorname{rk}S_{ui}\),
\(r_u=\sum_i r_{ui}>0\), \(d_u=\operatorname{rk}R_u>0\). On the
tied data \(R_u\) is the dual target lattice of the column pairing. Slopes
are normalized as in Theorem~\ref{geo:hn}. A slope in word \(i\) has
mass \(1/(h r_u)\), and a target slope has mass \(1/(h d_u)\).
The ordinary slope of a direct sum over one field weights blocks by
rank. Repeating both sides of each block \(m_u\) times would require
both \(m_u r_u\) and \(m_u d_u\) to be constant to recover these
two means. The ratio \(r_u/d_u\) need not be constant. We use the
two normalized traces explicitly.

\begin{lemma}[Cost of making the descended column integral]\label{trmean:graph}
For one rational map \(B:S_K\to R_K\) fix real archimedean
allowances \(h_v\) with normalized mean \(h_B\). Write
\(\langle a_v\rangle_\infty=(2/[K:\mathbb Q])\sum_{v\mid\infty}a_v\)
with one representative of each conjugate pair. Define
\[
\begin{aligned}
 b_w(B)&=\operatorname{length}_{\mathcal O_{K,w}}
              ((R_w+BS_w)/R_w),\\
 p_v(B;h_v)&=\sum_a\log\max(1,e^{-h_v}\sigma_a(B_v)^2),\\
 g(B)&=2\sum_{w\nmid\infty}\frac{\log Nw}{[K:\mathbb Q]}b_w(B)
                    +\langle p_v(B;h_v)\rangle_\infty ,
\end{aligned}
\]
where \(\sigma_a\) are singular values for the prescribed metrics
and \(Nw\) is the residue field cardinality. There is a source
\(S^\sharp\) with integral contractive inclusion in \(S\) such that
\(B:S^\sharp\to R\) is integral and has height at most \(h_B\)
when it is nonzero. Its mean slope satisfies
\(\operatorname{rk}(S)(\mu(S)-\mu(S^\sharp))=g(B)\).
For composable maps, each \(b_w\) is subadditive and each \(p_v\)
is subadditive with additive allowances.
\end{lemma}
\begin{proof}
Take \(S^\sharp=S\cap B^{-1}R\) at the finite places. The map \(B\)
identifies \(S_w/S_w^\sharp\) with \((R_w+BS_w)/R_w\).
In Smith coordinates of the actual shifted lattices,
\(b_w=\sum_a(-\nu_{w,a})_+\) over the nonzero elementary divisors.
Only finitely many places contribute. The matrix
\(\operatorname{diag}(\pi^{-a},\pi^a)\) for \(a>0\) has determinant
one and \(b_w=a\). At \(v\) give the source the Gram matrix
\(\max(I,e^{-h_v}B_v^*B_v)\) in original orthonormal coordinates,
using the spectral decomposition. It dominates \(I\) and
\(e^{-h_v}B_v^*B_v\), proving the two norm bounds. Its log determinant
is \(p_v\). The finite index changes the squared normalized
covolume logarithm by \(2\sum_w b_w\log Nw/[K:\mathbb Q]\), and
the metric contributes \(\langle p_v\rangle_\infty\). Subtracting
the two normalized degrees proves the slope identity. Both changes
are zero for the zero map.

For \(X\xrightarrow{B}Y\xrightarrow{A}Z\) at a finite place,
\((Z+AY+ABX)/(Z+AY)\) is a quotient of \((Y+BX)/Y\).
The inclusion \(Z+ABX\subset Z+AY+ABX\) now gives
\(b_w(AB)\le b_w(A)+b_w(B)\). At infinity,
\[
 p_v(B;h_v)=\log\max_{k\ge0}
                 \|\textstyle\bigwedge^k(e^{-h_v/2}B_v)\|^2 .
\]
The term \(k=0\) equals one. The identity follows by ordering singular
values. Submultiplicativity of each exterior norm proves the assertion.
Thus positive costs may be accumulated from the first layer, before any
HN comparison, using the same intermediate lattices and metrics on the
two sides of each composition. In particular, if
\(B_{ui}=C_{ui,m}\cdots C_{ui,1}\) and the local allowances add to
those of the column, then
\[
 \mathbb E\,\frac{\sum_i g(B_{ui})}{r_u}
 \le \sum_{j=1}^m\mathbb E\,\frac{\sum_i g(C_{ui,j})}{r_u}.
\]
Each layer on the right uses the \emph{initial} total source rank
of its translate even when intermediate ranks differ. A coupled
terminal map counts as one factor in this formula.
\end{proof}

At one chosen infinite place in each block, suppose the original
source is \(V_{ui}\), with frame maps \(A_{ui}:V_{ui}\to\mathcal H_u\),
\(n_u=\dim\mathcal H_u\). The target identifies isometrically with
a subspace of dimension \(d_u\). Orthogonal projection of \(A_{ui}\)
to it gives \(B_{ui}\) up to a nonzero scalar, as in
Theorem~\ref{geo:hn}. Put \(Q_u=\sum_i A_{ui}A_{ui}^*\) and define
\[
\begin{aligned}
 e&=\mathbb E\,\frac{\sum_i\|A_{ui}^*A_{ui}-I\|_{\rm HS}^2}{r_u},
 &v&=\mathbb E\,\frac{\|(n_u/r_u)Q_u-I\|_{\rm HS}^2}{n_u},\\
 \delta&=\mathbb E(1-d_u/n_u),
 &\theta&=\min(1,\delta+\sqrt e+\sqrt v).
\end{aligned}
\]
Write \(h_{ui}\) for the mean archimedean allowance of \(B_{ui}\).
Let \(s_u=\sum_i r_{ui}(\mu(S_{ui})-h_{ui})/r_u\) and
\(\mu_u=\mu(R_u)\). Choose \(u\) uniformly, then choose a source
slope uniformly by rank in that block. Let \(S\) be its value
minus \(h_{ui}+s_u\). Let \(S'\) use the repaired source with the
\emph{same} centring constants. For the target let \(T\) be a
rank-uniform target slope minus \(s_u\). These specify three finite
probability laws. For such a law \(L\), \(L^\uparrow(p)\) denotes
its increasing quantile, so integrating on \([0,1]\) gives its mean.

\begin{theorem}[A single normalized trace comparison]\label{trmean:trace}
With these data, put
\(\mathcal J=\mathbb E\big(\sum_i g(B_{ui})/r_u\big)\). Then
\begin{equation}\label{trmean:gap}
 \mathbb E(s_u-\mu_u)\ \le\ \mathcal F,\qquad
 \mathcal F=\mathcal J+\int_{1-\theta}^1 S^\uparrow(p)\,dp
                         -\int_0^\theta T^\uparrow(p)\,dp .
\end{equation}
The two integrals use only the highest source fraction and the lowest
target fraction, in their original metrics. All ranks and frame
defects in this comparison may vary with the translate.
\end{theorem}
\begin{proof}
The HN cutoff of an integral map of height \(\le h\) sends source
slopes \(\ge a+h\) into the target cutoff at \(a\). Indeed, a nonzero
map from this source part into the quotient by that target part
would have an image which is a quotient of the former and whose
saturation is a sublattice of the latter. The quotient with quotient
metric has mean slope at least \(a+h\). The saturated image with
induced metric has mean slope \(<a\). The determinant of the induced
rank-\(k\) map lowers degree by at most \(kh\), since it is integral
and its archimedean squared norm is bounded by the \(k\)-th power of
the squared operator norm. This is a contradiction. This argument
uses the defining extremal properties of the HN filtration with its
subspace and quotient metrics. For the slope inequality and its
filtration form see the proof of Proposition~2.2 in
\href{https://arxiv.org/pdf/0803.2583}{H.~Chen, \emph{Positive degree and arithmetic bigness},
\S2.5}. Applied to \(S^\sharp\hookrightarrow S\)
at height zero it gives \(S'^\uparrow\le S^\uparrow\).
Lemma~\ref{trmean:graph} gives the three means \(0,-\mathcal J\),
and \(\mathbb E(\mu_u-s_u)\), respectively.

Fix a real threshold \(t\). Let \(U_{ui}\) be the subspace of the
repaired source above or at \(s_u+t+h_{ui}\), viewed in the original
complex vector space. The change of metric preserves its dimension
and its image as a subspace. Put \(q=\Pr(S'<t)\), let \(P_u\) project
orthogonally onto \((\sum_i A_{ui}U_{ui})^\perp\), and put
\(z=\mathbb E(\operatorname{tr}P_u/n_u)\). Then
\[
 z-\sqrt{zv}\ \le\ \mathbb E\,\frac{\operatorname{tr}(P_uQ_u)}{r_u}
                     \ \le\ q+\sqrt{qe}.
\]
For the left inequality pair \(P_u\) with \((n_u/r_u)Q_u-I\) in
the normalized Hilbert--Schmidt inner product, then apply
Cauchy--Schwarz over the blocks. For the right inequality
\(P_u A_{ui}\) vanishes on \(U_{ui}\). Sum the squared norms over
an orthonormal basis of its complement in the original source metric.
Pair its projection with \(A_{ui}^*A_{ui}-I\) and use the same
inequality with weights \(1/(h r_u)\). The mean relative dimension
of these complements is exactly \(q\).

If the relative codimension of the projected images in the target
is \(z_{0u}\), projection induces a surjection of the ambient
quotient onto the target quotient. Hence
\((d_u/n_u)z_{0u}\le\operatorname{tr}P_u/n_u\), so
\(\mathbb E z_{0u}\le z+\delta\). By the cutoff property, all projected
images lie in the target above or at \(s_u+t\).
Since \(z,q\le1\), we have
\(\Pr(T<t)\le\Pr(S'<t)+\theta\), including the case \(\theta=1\).
For finite increasing quantiles this implies
\(S'^\uparrow(p)\le T^\uparrow(p+\theta)\) for almost every
\(0<p<1-\theta\). Indeed, a strict violation yields a threshold
between the two values that contradicts the preceding count.
Integrate on that interval and use \(S'^\uparrow\le S^\uparrow\)
on the upper fraction. Substituting the three means gives
\eqref{trmean:gap}, also at the endpoints \(\theta=0,1\).
\end{proof}

A slope tail is essential. Take one block with \(D-1\) orthogonal unit
columns into the first \(D-1\) coordinates of a rank-\(D\) target,
with \(D\ge2\). All sources have slope zero on a common shifted scale,
and the target slopes are \(0,\ldots,0,-L\), obtained by changing
only the last line metric. Here \(\mathcal J=e=\delta=0\) and
\(v=1/(D-1)\), independent of \(L>0\). The mean gap is \(L/D\).
Consequently even small trace defects leave a quantitative tail
estimate to prove in the actual arithmetic family.

\begin{corollary}[Comparison with the error budget of the chain]\label{trmean:budget}
Suppose \(M_u=\operatorname{tr}Q_u\ge r_u/2\) and
\(d_u/n_u\ge1-\delta_*>0\) for the chosen blocks. Set
\[
 \eta=\mathbb E\left(\frac{n_u\operatorname{tr}(Q_u^2)}{M_u^2}-1\right),
 \qquad U=\mathbb E\max_{i,a}(\lambda_{ui,a}-h_{ui}-s_u),
\]
where \(\lambda_{ui,a}\) run over the original source slopes. For
\(H>0\) put \(T_{\rm out}=\int_0^1(-H-T^\uparrow(p))_+\,dp\).
Then \(U\ge0\) and
\begin{equation}\label{trmean:budget-eq}
 \mathbb E(s_u-\mu_u)\le T_{\rm out}+
 \frac{(\eta+4\sqrt e)H+4\mathcal J+4U\log(1+H/U)}{1-\delta_*}.
\end{equation}
The last product in the numerator is zero when \(U=0\). Integral
columns satisfying the prescribed global height bounds admit allowances
with \(\mathcal J=0\).
\end{corollary}
\begin{proof}
In the preceding trace proof normalize \(Q_u\) by \(M_u/n_u\). The
mean squared normalized defect is exactly \(\eta\), by expansion of
the square. The same projection test gives
\[
 z-\sqrt{z\eta}\le\mathbb E\,\operatorname{tr}(P_uQ_u)/M_u
                       \le 2(q+\sqrt{qe}).
\]
Here the upper bound uses \(M_u\ge r_u/2\) and positivity.
Since \(2\sqrt{z\eta}\le z+\eta\), the target fraction below the
threshold is at most \((\eta+4q+4\sqrt e)/(1-\delta_*)\).
At threshold \(-t\) the discarded source fraction is
\(q=\Pr(S'<-t)\). Its integral on \(0<t<H\) exceeds the corresponding
integral for \(S\) by at most \(\mathcal J\). Indeed the integrals
give the expectations of \(\min(H,(-S')_+)\) and \(\min(H,(-S)_+)\).
This function of the slope is decreasing and 1-Lipschitz, while the
ordered quantiles decrease by total integral \(\mathcal J\).

Write \(U_u=\max_{i,a}(\lambda_{ui,a}-h_{ui}-s_u)\). The centred
source in block \(u\) has mean zero and maximum \(U_u\ge0\).
For \(t>0\) its fraction below \(-t\) is at most \(U_u/(U_u+t)\).
The function \(a\mapsto a/(a+t)\) is concave on \(a\ge0\), so
\(\Pr(S<-t)\le U/(U+t)\). Integrating the three terms of the
target fraction on \([0,H]\) and adding \(T_{\rm out}\) bounds
the negative part of \(T\), which dominates \(-\mathbb E T\).
This proves \eqref{trmean:budget-eq}. For a nonzero integral column
of height \(\le h_{ui}\), choose \(h_v=\log\|B_v\|^2\) plus the
nonnegative constant making its normalized mean \(h_{ui}\).
Both terms of \(g(B)\) vanish. The zero column costs zero for any
allowances with the prescribed mean.
\end{proof}

\section{Theta matching and the arithmetic Hilbert--Samuel estimate}
Fix the tower, measures and matching data of \S\ref{newtheta:global}.
Theorem~\ref{tcr:global} constructs a supported test with nonzero
primed period and proves \(I(\Phi)=cI'(\Phi')\) for every matched
pair, at every allowed real signature. The scalar \(c\ne0\) belongs
to the chosen data. In its two rank recursions the comparison scalar
\(C\) is uniform in the admissible finite data once the adapted
upper and lower geometries, real matching maps and measures are
fixed. The proof includes the inert finite and real local comparisons.
The evaluation lattice used in Section~\ref{main:assembly} requires
more than proportionality of these two linear functionals. On
marked components the doubled kernel gives spanning over the
coefficient field (Propositions~\ref{dg:identity} and~\ref{munf:span}).
Theorem~\ref{bfam:oncharts} controls the finite denominators and
Petersson norms of spanning sections under its chart and cusp
hypotheses. The final change of weight also involves the integral
determinants of two conjugate Hodge eigenspaces. On the indefinite
selected family their common local orders and the CM evaluation
in Theorem~\ref{cmeig:splitmodel} give this change at every weight.

\subsection{Local matching in split matrix coordinates}\label{newtheta:local}

We compute the two local functionals used at a selected split place. This
also gives test vectors on which the functions agree at \emph{every} group
element. Let \(k\) be a nonarchimedean local field, \(\mathcal O\) its
valuation ring, \(\pi\) a uniformizer, and \(|\pi|=|\mathcal O/\pi|^{-1}\).
Put \(d=a+a'\), with \(a,a'>0\), \(G_j=\mathrm{GL}_j(k)\),
\(K_j=\mathrm{GL}_j(\mathcal O)\), and give \(G_j\) Haar measure
such that \(K_j\) has volume one. Use the row matrices
\(X_0=[I_a\;0]\) and \(Y_0=[0\;I_{a'}]\).
The characters \(\alpha,\beta:k^\times\to\mathbb C^\times\) are smooth
and unitary, and \(\nu=\alpha/\beta\). In the two polarizations the
orbital functions, on locally constant functions compactly supported in
the respective full-row-rank sets, are
\begin{align}
 B_\Phi(h)&=\beta(\det h)|\det h|^{-a/2}
 \int_{G_a}\Phi(gX_0h^{-t})\,\nu(\det g)|\det g|^{d/2}\,dg ,
 \label{newtheta:big}\\
 S_\Psi(h)&=\alpha(\det h)|\det h|^{a'/2}
 \int_{G_{a'}}\Psi(gY_0h)\,\nu(\det g)|\det g|^{d/2}\,dg .
 \label{newtheta:small}
\end{align}
Here \(h^{-t}=(h^{-1})^t\). For any fixed full-rank \(X\) the map
\(g\mapsto gX\) is a homeomorphism onto the full-rank matrices with
the same row space as \(X\). The image is closed in the full-rank set,
since having rows in a fixed subspace is closed. The inverse is right
multiplication by any right inverse of \(X\). Thus the integrals here
have compact support in \(g\). The large-side definition in the original
model is a sum of integrals over Cartan double cosets. A compact subset of
\(G_a\) meets only finitely many of its disjoint open double cosets, so
the two definitions agree in this case.
\src{src/Development/AbelianVarietyTheory/Definitions/Divisors/RankEqCardIffExistsFrameMatrixFirstColumnBlock.lean}{703}{The two orbital functions and their determinant factors}.

\begin{proposition}[The common image of full-rank split tests]\label{newtheta:split}
Let \(P\subset G_d\) be the block upper triangular parabolic with
diagonal blocks of sizes \(a,a'\). The image of each of
\eqref{newtheta:big} and \eqref{newtheta:small} on these test spaces is
the full space of smooth functions \(f:G_d\to\mathbb C\) satisfying
\[
 f\left(\begin{pmatrix}A&*\\0&D'\end{pmatrix}h\right)
 =\alpha(\det A)\beta(\det D')|\det A|^{a'/2}
             |\det D'|^{-a/2}f(h).
\]
In particular there are matched tests with value one at the identity.
\end{proposition}
\begin{proof}
For \(p\in P\), \(X_0p^{-t}=A^{-t}X_0\) and \(Y_0p=D'Y_0\).
In the first integral substitute \(gA^{-t}\) for \(g\), and in the
second \(gD'\) for \(g\). Right translations preserve Haar measure on
\(G_j\), as is also seen from its density
\(|\det g|^{-j}\prod_{i,l}dg_{il}\) in additive matrix coordinates.
Using \(\alpha=\nu\beta\) gives the displayed character in both cases.
Right translation of the output is obtained by the corresponding linear
substitution in the test times its determinant prefactor. Compact support
and local constancy give invariance under some open subgroup on the right.

Let \(K_j(e)\) be the kernel of reduction modulo \(\pi^e\). Take
\(e\ge1\) so large that both characters are trivial on
\(1+\pi^e\mathcal O\). For \(\Phi_e,\Psi_e\) take the indicators of
\(X_0+\pi^eM_{a,d}(\mathcal O)\), \(Y_0+\pi^eM_{a',d}(\mathcal O)\)
divided by the volumes of \(K_a(e)\), \(K_{a'}(e)\), respectively.
Both supports are compact and have full row rank. For \(h\in K_d\),
the integrality and full rank modulo \(\pi\) of \(gX_0h^{-t}\) on
the first support force \(g\in K_a\), by multiplying on the right by
\(h^t\). The same argument with \(h^{-1}\) gives \(g\in K_{a'}\)
on the second support. Modulo \(\pi^e\), the equations become
\([g\;0]=X_0h^t\) and \([0\;g]h=Y_0\), respectively.
Each has solutions precisely when the lower left block of \(h\)
is zero modulo \(\pi^e\). In this case the diagonal blocks \(A,D'\)
are invertible modulo \(\pi^e\). The solutions are one congruence
coset, given by \(g\equiv A^t\) or \(g\equiv (D')^{-1}\).
The two normalized values are consequently
\(\beta(\det h)\nu(\det A)\) and
\(\alpha(\det h)\nu(\det D')^{-1}\). They agree because
\(\det h\equiv\det A\det D'\pmod{\pi^e}\).
In particular the common function is one at \(h=1\).

Write \(K=K_d\) and \(U=(P\cap K)K_d(e)\). This is the inverse image
of the block upper triangular subgroup modulo \(\pi^e\): one lifts
the two invertible diagonal blocks and the upper right block to obtain
an element of \(P\cap K\). Our common function on \(K\) is supported
on \(U\) and is right \(K_d(e)\)-invariant. Also \(G_d=PK\). Indeed
the intersection of any \(a'\)-dimensional row space with \(\mathcal O^d\)
is a saturated free summand of rank \(a'\) by the elementary divisor
theorem. Extending an integral basis of that summand gives a matrix of
\(K\) with the same last-row space as the given matrix of \(G_d\).
Their quotient belongs to \(P\).

For any smooth \(f\) with the stated covariance choose \(e\) still
larger so that \(f|_K\) is right \(K_d(e)\)-invariant, possible by
compactness. On the finitely many cosets \(U k_j\) in \(K\), the
restrictions of \(f\) are determined by \(f(k_j)\). In fact \(K_d(e)\)
is normal in \(K\), so \(f(uk_j)\) transforms by the prescribed
character on \(P\cap K\) and trivially on the congruence kernel.
The right translates \(h\mapsto B_{\Phi_e}(hk_j^{-1})
=S_{\Psi_e}(hk_j^{-1})\) have exactly these transformation laws,
disjoint supports on \(K\), and value one at \(k_j\). Their linear
combination with coefficients \(f(k_j)\) agrees with \(f\) first on
\(K\) and then on \(PK\). Each translate comes from a test of the
same allowed kind. This proves surjectivity on both sides.
\end{proof}

The ramified split condition asks for a large-side value even on
tests meeting the singular locus. There is a direct justification for
the finite shell prescription in this case.
\begin{lemma}[Cancellation near the singular strata]\label{newtheta:shell}
Suppose \(\nu\) is nontrivial on \(\mathcal O^\times\). For every
locally constant compactly supported \(\Phi\) on \(M_{a,d}(k)\) and
every fixed full-rank \(X\), only finitely many double cosets
\(K_a\operatorname{diag}(\pi^{m_i})K_a\) have a nonzero integral of
\(\Phi(gX)\nu(\det g)|\det g|^{d/2}\).
\end{lemma}
\begin{proof}
Choose a right inverse \(X^\dagger\). The identity
\(g=(gX)X^\dagger\) and compact support give a common lower bound for
the valuations of the entries of all contributing \(g\). Multiplication
by \(K_a\) preserves the supremum norm of matrix entries in either
direction, so the bound applies to each \(m_i\). Choose \(u\in
\mathcal O^\times\) with \(\nu(u)\ne1\). Compact support and local
constancy give a common additive period lattice
\(\pi^tM_{a,d}(\mathcal O)\) for \(\Phi\), by covering its compact
support by finitely many balls of constancy. If \(m_i\) is large,
replacing \(k_1\) by \(k_1\operatorname{diag}(1,\ldots,u,\ldots,1)\)
in \(g=k_1\operatorname{diag}(\pi^{m_j})k_2\), with \(u\) at
position \(i\), changes \(gX\) by an element of this period lattice
uniformly in \(k_1,k_2\in K_a\). The determinant character is
multiplied by \(\nu(u)\), while its absolute value is unchanged.
The normalized measure on a double coset is the pushforward of product
probability Haar measure on \(K_a^2\): both measures are invariant
under its transitive left and right \(K_a\) actions. Averaging over
\(k_1\) thus makes this shell integral zero. With all \(m_i\) bounded
above and below, only finitely many integer tuples remain.
\end{proof}

\begin{proposition}[The whole kernel on the smaller split side]\label{split:kernel}
Put \(n=a'<a\) and \(s=d/2>n\), and set
\(\lambda(g)=\nu(\det g)|\det g|^s\) on \(G=\mathrm{GL}_n(k)\).
On all locally constant compactly supported functions on \(M_{n,d}(k)\),
the integral \eqref{newtheta:small} converges absolutely and its kernel
is precisely the span of
\[
        \Phi(g\,\cdot)-\lambda(g)^{-1}\Phi,\qquad g\in G .
\]
For unramified \(\nu\), the value at the identity of the standard integral
matrix indicator is \(\prod_{j=0}^{n-1}(1-\nu(\pi)q^{j-s})^{-1}\ne0\),
where \(q=|\mathcal O/\pi|\).
\end{proposition}
\begin{proof}
For any full-rank \(X\), boundedness of \(gX\) bounds \(g\) by a right
inverse of \(X\). On that bounded set the absolute integration density
is a constant times \(|\det g|^{s-n}\) in additive matrix coordinates.
It is bounded, so the integral converges. Haar substitution shows that
the displayed differences have zero integral.

On the full-rank open set the map taking the row span is a principal
\(G\)-bundle over a Grassmannian, trivial on each pivot chart by
row reduction. Subdivide the projection of a compact support into
finitely many disjoint compact open subsets of such charts. In chart
coordinates each test is a finite sum of products of a base test and
a test on \(G\). This follows by covering its compact support by
compact open rectangles of constancy and taking a common refinement.
On the fibre \(G\) the quotient by the displayed differences is
one-dimensional, detected by integration against \(\lambda(g)\,dg\).
In fact, for any finite list of tests take a compact open subgroup
\(J\subseteq\ker\lambda\) fixing all of them on the right. Modulo
differences \(1_{tJ}\) equals \(\lambda(t)1_J\); its integral is
\(\lambda(t)\operatorname{vol}(J)\). This proves the assertion on one
fibre. A base refinement where the base factors are constant gives the
same kernel description throughout the full-rank open set. All its
row spaces are tested by \(Y_0h\) as \(h\) varies.

We show that rank \(r<n\) adds no classes to this quotient. Over
a pivot chart of the row-span Grassmannian the rank \(r\) stratum has
fibre the space of injective \(n\)-by-\(r\) matrices. Here a basis of
the row space identifies the original matrix with its matrix of row
coefficients. The source group acts transitively on this fibre by left
multiplication. Suppose a linear functional \(D\) there satisfies
\(D(f(g\,\cdot))=\lambda(g)^{-1}D(f)\).
Around \(X_r=(I_r,0)^t\), take
\(V=X_r+\pi^t M_{n,r}(\mathcal O)\) for \(t\ge1\) sufficiently large
that \(\nu\) is trivial on \(1+\pi^t\mathcal O\). The congruence
subgroup \(K_n(t)\) acts transitively on \(V\), by completing its
first \(r\) columns, and preserves additive volume on \(V\).
Its invariant functionals on tests supported in \(V\) are multiples of
that volume: averaging any such test over the compact group gives the
constant average on \(V\), and is a finite average by smoothness.
Take \(t_0=\operatorname{diag}(I_r,\pi I_{n-r})\), which fixes \(X_r\)
and sends \(V\) into \(V\). On the one hand
\(D(1_{t_0V})=|\det t_0|^r D(1_V)\), by additive volume columnwise.
On the other hand equivariance gives \(\lambda(t_0)D(1_V)\).
Since \(|\lambda(t_0)|=|\det t_0|^s\) and \(r<s\), we have \(D(1_V)=0\).
The translates of \(V\) cover the fibre, so \(D=0\). For \(r=0\)
the fibre is a point and the same character test gives zero directly.
Testing products with each base indicator proves vanishing of every
equivariant functional on the stratum. Its quotient by differences
is zero since linear functionals separate points of any vector space.

Filter the closed rank-deficient locus by the finitely many closed
sets \(\{\operatorname{rank}\le r\}\). Restriction of tests to a closed
subset is onto: a compact open subset of it can be extended to a finite
union of compact open neighborhoods in the ambient space, chosen to
avoid its closed complement within the subset. The kernel consists of
tests compactly supported on the open complement. Taking quotients by
differences is right exact. Induction therefore gives zero quotient for
the whole rank-deficient locus. Consequently each test on the matrix
space is congruent modulo differences to one supported on the full-rank
open set. Their orbital integrals agree, so the full-rank calculation
proves the asserted equality of kernels.

Finally the integral at the identity for the spherical indicator is over
the integral nonsingular matrices. Its right \(K_n\)-cosets correspond
to the finite-index sublattices of \(\mathcal O^n\) and have volume one.
For unramified \(\nu\), the weight of an index \(q^b\) lattice is
\((\nu(\pi)q^{-s})^b\). In upper triangular column Hermite form with
diagonal \(\pi^{b_i}\), the entries in row \(i\) above the diagonal
are taken modulo \(\pi^{b_i}\). There are
\(\prod_i q^{(n-i)b_i}\) choices and \(b=\sum b_i\).
Summing the \(n\) convergent geometric series gives the displayed product.
\end{proof}
\begin{corollary}[Both orbital maps on ramified split tests]\label{split:ram}
If \(\nu|_{\mathcal O^\times}\ne1\), the larger-side shell prescription
has kernel the corresponding span of source-action differences on
\emph{all} tests. Both full test spaces have the common image of
Proposition~\ref{newtheta:split}.
\end{corollary}
\begin{proof}
In the larger integral first replace \(d/2\) in the integration weight
by \(z\) with \(\Re z>a\). Then it is absolutely convergent by the
bounded-density argument above, and change of variables under every
source translation is valid. For fixed \(\Phi,X\), the additive-period
argument of Lemma~\ref{newtheta:shell} kills all but finitely many
Cartan shells, independently of \(z\), since the determinant norm is
constant on each shell. Their sum is a Laurent polynomial in \(q^{-z}\).
The identity for a given translate, including its factor
\(\nu(\det g)^{-1}|\det g|^{-z}\), consequently holds at \(z=d/2\)
as well. The same change of variables gives parabolic covariance at
that parameter. A compact open subgroup fixes the test on the right
in the target action, which gives smoothness of its orbital function.

The rank argument in Proposition~\ref{split:kernel} applies with
\(n=a\); whenever \(r\ne d/2\), the absolute values in the test
still differ. For \(r=d/2<a\) fix \(u\in\mathcal O^\times\) with
\(\nu(u)\ne1\) and use \(\operatorname{diag}(I_r,u,1,\ldots,1)\)
in the stabilizer. This preserves the compact open ball \(V\) and
changes the functional on its indicator by the factor \(\nu(u)\).
It must therefore vanish there and on all translates. Thus every
singular stratum has zero quotient by differences. The full-rank
kernel calculation and the shell transform on the entire space finish
the proof exactly as on the smaller side.
\end{proof}
In particular every \(\lambda^{-1}\)-equivariant functional for the written
translations on the smaller split model factors through its orbital integral; the same statement applies on
the larger side at the indicated ramified places. Full-rank tests span
both indicated images. For the larger unramified side when \(\nu=1\)
and \(r=d/2\), the stabilizer test gives equality for every determinant.
The next lemma shows that the infinite type excludes this resonance
at the good places of the global construction.

\begin{lemma}[The character at a totally split rational prime]\label{split:ray}
Let \(\mathfrak w\) lie above a rational prime completely split in the
Galois CM field \(E\). If its restriction to \(F\) is outside the
ramification set \(S\) of the triple, then
\(t=\nu_{\mathfrak w}(\pi)\) has infinite order, for either choice of
the place above \(F\) and every uniformizer \(\pi\).
\end{lemma}
\begin{proof}
Fix \(E\subset\mathbb C\). The chosen embeddings above the real places
form a CM type \(\Phi\), and we can regard them as automorphisms of
\(E\) by normality. Put \(m_\tau=m_\alpha(\tau)-m_\beta(\tau)>d\).
The infinite component of \(\nu\) at \(x\in E^\times\) is
\(\prod_{\tau\in\Phi}(\tau x/|\tau x|)^{m_\tau}\).
Take a finite ray modulus for \(\nu\), prime to \(\mathfrak w\).
By finiteness of the ray class group there are \(h>0\) and \(x\in
E^\times\) such that \((x)=\mathfrak w^h\) and \(x\equiv1\) modulo
this modulus. All ramified finite components are one at \(x\);
at every other finite place except \(\mathfrak w\), \(x\) is a unit.
The continuous Hecke character is trivial on diagonal \(E^\times\).
Its value on \(x\), squared, therefore gives
\[
 t^{-2h}=\prod_{\tau\in\Phi}
            \bigl(\tau x/\overline{\tau x}\bigr)^{m_\tau}.
\]
The principal fractional ideal on the right has exponents \(hm_\tau\)
and \(-hm_\tau\) at \(\tau\mathfrak w\) and
\(\bar\tau\mathfrak w\), respectively. These primes are all distinct
because the decomposition group is trivial. The fractional ideal is
nontrivial, whereas every root of unity generates the unit ideal.
This proves the assertion. Unramifiedness gives independence of \(\pi\).
\end{proof}
\src{src/Development/AbelianVarietyTheory/Definitions/Foundations/NotIsSplitOfIsInertAt.lean}{266}{Continuity and triviality on the diagonal};
\src{src/Development/AbelianVarietyTheory/BlockTheory/DivisorBase/WeilTensorModels.lean}{114}{Infinite character on ideles with trivial finite part};
\src{src/Development/AbelianVarietyTheory/BlockTheory/DivisorBase/WeilTensorModels.lean}{374}{Local ramification and infinite exponents}.

\begin{proposition}[The larger orbital transform at a good split place]
\label{split:good}
Suppose \(\nu\) is unramified and \(t=\nu(\pi)\) has infinite order.
Set \(n=a\), \(s=d/2\). On every matrix test \(\Phi\), the larger
transform in \eqref{newtheta:big}, initially integrated with exponent
\(z\) for \(\Re z>n\), has a regular continuation at \(z=s\).
This transform and the smaller transform \eqref{newtheta:small}
have the same image on their whole test spaces, namely the induced
representation in Proposition~\ref{newtheta:split}. The kernel on
each side is the span of its source-action differences. For
unramified \(\alpha,\beta\) their integral-matrix indicators can be
normalized to have the same nonzero image.
\end{proposition}
\begin{proof}
Fix a full-rank \(X\). The function \(H(g)=\Phi(gX)\) on \(M_n(k)\)
is locally constant and compactly supported, since multiplying by a
right inverse of \(X\) bounds its support. In additive matrix
coordinates the integral is a constant times
\[
 Z_n(H,z)=\int_{M_n(k)} H(g)\nu(\det g)|\det g|^{z-n}\,dg .
\]
Assign the integrand value zero on singular matrices, a null set
by successive one-variable polynomial testing and Fubini.
For \(0\le j\le n\) use the same exponent \(z-n\) in \(Z_j\) on
\(M_j(k)\), with \(Z_0\) evaluation at its single point.
We show by induction that \(Z_j\) is rational in \(q^{-z}\) and
its only possible poles for finite \(z\) satisfy
\[
           1-t^i q^{-i(z-n+i)}=0,\qquad 1\le i\le j .
\]
Away from a neighborhood of the zero matrix, a compact support has
a finite disjoint compact open subdivision on each part of which a
chosen nonzero pivot \(b\) has constant valuation. Permute rows and
columns to make it the top left entry. For the other three blocks use
coordinates \((B,C,D-Cb^{-1}B)\) instead of \((B,C,D)\). This
preserves additive measure by translation in \(D\) and factors the
determinant as \(b\det(D-Cb^{-1}B)\), up to a unit sign.
In these coordinates a compactly supported locally constant test is
a finite sum of products on compact open rectangles, by a finite
refinement of a cover by rectangles of constancy. The pivot weight
is a Laurent monomial times a constant. The remaining integrals thus
have the asserted form by induction.

Near zero subtract the constant value of the test times the
indicator of a sufficiently small origin ball. Scalar dilation
reduces that ball to \(M_j(\mathcal O)\) up to a Laurent monomial.
In this unit ball the complement of \(\pi M_j(\mathcal O)\) is
covered by the preceding pivot argument. The integral on the
inner ball is \(t^j q^{-j^2-j(z-n)}\) times the unit-ball integral.
Solving this identity in the half-plane of convergence introduces
exactly the indicated possible denominator with \(i=j\). This
proves rational continuation. For real \(z=s\), a zero of a
denominator would force first \(s-n+i=0\) by \(|t|=1\), and then
\(t^i=1\). Hence every test is regular there.

Source equivariance holds first in the convergent half-plane by
Haar substitution and hence at \(s\) by continuation. The same
reasoning gives the displayed parabolic covariance of
Proposition~\ref{newtheta:split} at \(s\); smoothness follows from
the compact open stabilizer of the test. On full-rank tests the
integrals are entire and the image and kernel were computed above.
On the rank-\(r<n\) stratum the stabilizer test in
Proposition~\ref{split:kernel} would require
\[
          |\det t_0|^r=\nu(\det t_0)|\det t_0|^s,\qquad
          t_0=\operatorname{diag}(I_r,\pi I_{n-r}).
\]
The moduli differ for \(r\ne s\), and for \(r=s\) equality would
give \(t^{n-r}=1\). Thus every singular stratum has zero quotient by
source differences. The rank filtration and the full-rank kernel
calculation prove the asserted kernel on all tests. The smaller
side is Proposition~\ref{split:kernel}.

In the convergent region the Hermite count in
Proposition~\ref{split:kernel} gives the larger spherical value
\(\prod_{j=0}^{n-1}(1-tq^{j-z})^{-1}\). Its continuation at \(s\)
is finite and nonzero; the smaller value has the same formula
with \(n=a'\). If \(\alpha,\beta\) are unramified both output
functions are right \(K_d\)-invariant. Such an induced function is
determined by its value at one since \(G_d=PK_d\), so dividing each
map by its indicated value identifies the two spherical outputs.
\end{proof}

The covariance is \(\delta_P^{1/2}
((\alpha\circ\det)\boxtimes(\beta\circ\det))\):
\(\delta_P(A,D')=|\det A|^{a'}|\det D'|^{-a}\), from conjugation on
the upper right block. The two determinant characters give an
irreducible unitary representation of the Levi, since their external
product is one-dimensional and unitary. Bernstein's irreducibility
theorem for normalized induction of irreducible unitary representations
of a Levi in \(\mathrm{GL}_d(k)\) therefore proves that this common
image is irreducible, both here and in Corollary~\ref{split:ram}.
Bernstein,
\emph{P-invariant distributions on \(\mathrm{GL}(N)\) and the classification
of unitary representations of \(\mathrm{GL}(N)\) (non-Archimedean case)},
\S0.2, corollary, proved in \S8.2~\cite{bernsteinInvariant}.

Lemma~\ref{split:ray} supplies the character hypothesis for every
good totally split place of the triple, at either splitting.
The construction here evaluates a rational continuation on unramified
tests. For the integral-matrix indicator infinitely many Cartan shells
have nonzero integral, so the finite-support sum defining the prescribed
shell map cannot replace this continuation. The matching condition
that specifies that shell map is imposed at \(S\), where the earlier
ramified calculation applies. The results here concern the displayed
matrix models. Passing to the given adelic Weil actions requires the
identifications with these polarizations and their lattice vacua.
\src{src/Development/AbelianVarietyTheory/Definitions/Divisors/RankEqCardIffExistsFrameMatrixFirstColumnBlock.lean}{707}{Finite-support Cartan sum in the prescribed map};
\src{src/Development/AbelianVarietyTheory/BlockTheory/DivisorBase/WeilTensorModels.lean}{470}{Places at which equality with that map is required}.



\begin{lemma}[The finite epsilon factors in a descent]\label{newtheta:gauss}
Let \(q=|\mathcal O/\pi|\) and let the additive character \(\psi\)
have conductor \(\mathcal O\). Suppose the unitary character \(\nu\)
has conductor exponent \(c\ge1\). Put
\[
 G=\sum_{y\in(\mathcal O/\pi^c)^\times}
          \nu(y)^{-1}\psi(y/\pi^c).
\]
Then \(|G|^2=q^c\). In particular
\(q^{-cs}G\,\nu(\pi)^c\ne0\) for every integer \(s\).
\end{lemma}
\begin{proof}
The two characters on units in the sum do not depend on the chosen
lifts. Write \(R=\mathcal O/\pi^c\) and \(\Psi(x)=\psi(x/\pi^c)\).
The additive pairing \((t,x)\mapsto\Psi(tx)\) on \(R\) is perfect
by the conductor assumption. Extend \(\nu^{-1}\) by zero to the
nonunits. Its unnormalized Fourier transform at \(t\in R^\times\)
is \(\nu(t)G\) by changing \(y\) to \(ty\). For \(t\in\pi R\)
choose a unit \(v\equiv1\bmod\pi^{c-1}\) with \(\nu(v)\ne1\),
using primitivity (for \(c=1\) choose any such unit). Then \(tv=t\)
in \(R\), and multiplication of the summation variable by \(v\)
forces that Fourier coefficient to vanish. Summing the squared
moduli of all Fourier coefficients gives \(q^c|R^\times|\) by
expanding and applying additive orthogonality. The same sum is
\(|R^\times|\,|G|^2\), proving the claim.
\end{proof}

\subsection{The global comparison required of the local maps}\label{newtheta:global}

In the global triple, \(E/F\) is CM quadratic, \(E/\mathbb Q\) is
Galois, and \(E/F\) is unramified at finite primes. The additive
character on \(\mathbb A_F/F\) has infinite component
\(\exp(-2\pi i\operatorname{Tr}(x))\). It has conductor \(\mathcal O_v\)
at the designated split places and at each inert place, where the residue
characteristic is odd and the residue cardinality is greater than three.
The three spaces have dimensions \(a,a',d=a+a'\) and signatures
\((a-p,p),(q,a'-q),(d-p-q,p+q)\) at each real place, using \(i\)
times the skew-Hermitian form on the third space. The integral structures
are self-dual for the specified pairings (the skew lattice uses the
trace dual). Each comes with a Cartan involution and the selected local
bases. The rank and weight restrictions are
\[
\begin{gathered}
a'<a,\qquad d\equiv0\pmod2,\qquad 2(p+q)+1<\min\{(a-a')/2,a'\},\\
m_\alpha(w)-m_\beta(w)>d\qquad(w\mid\infty).
\end{gathered}
\]
The unitary Hecke characters \(\alpha,\beta\) of \(E\) have these
infinite exponents in the chosen embeddings, and restrict to
\(\eta_{E/F}^{\,a'},\eta_{E/F}^{\,a}\) on \(\mathbb A_F^\times\).
They are unramified outside the primes above \(S=S_0\sqcup Q\), a
subset of the designated finite places with a chosen splitting
\(E\hookrightarrow F_v\). Their quotient \(\nu\) is ramified at
every prime above \(S\). A further place \(v_1\notin S\) has
underlying rational prime totally split in \(E\), with a choice of
splitting at \(v_1\). The Weil actions use respectively the pairs
\((\beta,\beta/\alpha)\) and \((\alpha,\alpha/\beta)\) in their
prescribed polarizations. The local vacua belong to the prescribed
lattice eigenlines and are standard characteristic functions almost
everywhere.
\src{src/Development/AbelianVarietyTheory/BlockTheory/DivisorBase/WeilTensorModels.lean}{388}{Base, signatures, characters and Weil actions}.

In these models the missing global assertion has three parts.
The first is the simultaneous construction of nonzero irreducible
quotients at infinity and each finite place, as in
Proposition~\ref{geo:theta}. The maps kill the infinitesimal and compact
source actions at infinity (derivatives map to zero and compact actions
preserve the image of each vector) and kill the span of source-action
differences at finite places. Target Lie and compact actions, or target
group actions at finite places, preserve matching. The two vacua have
the same nonzero image at every \(v\notin S\). At \(v\in S\) equality
of images must be equivalent, for \emph{all} Schwartz--Bruhat inputs,
to equality of the two orbital functions, in the coordinates taking the
vacuum to the integral-matrix indicator. At good places (outside \(S\)
and above a totally split rational prime), the kernel on Schwartz tests
must be exactly the span of the source-action differences. Full-rank
primed tests in the conjugate polarization at either splitting there
must still surject. Proposition~\ref{newtheta:split} proves the
full-rank image calculation in split coordinates. Corollary~\ref{split:ram}
and Proposition~\ref{split:good} determine the whole split kernels at
the specified ramified and good places. Section~\ref{threal:main}
establishes convergence for the two adelic kernels and constructs the
remaining local quotients, including all real signatures.

The second assertion is global convergence and proportionality. On the
unprimed side every adelic Schwartz input must give an absolutely
integrable theta kernel in \(g\) for each \(h\). On the primed side
use the span of factorizable inputs whose factor at \(v_1\) vanishes
on rank-deficient matrices in the indicated polarization. This space
must be stable under the Weil action. For each such input and each
integer \(N\ge0\), there must be \(C,k\) such that
\[
 |\Theta_{\omega'}(g,h;\Phi')|\,H_{[U(V')]}(g)^N
                 \le C H_{U(W)}(h)^k .
\]
Here the quotient height is the infimum of the matrix height over
rational left translates. The primed integral must also converge and
have moderate growth. For \emph{every} choice of matching maps, one
nonzero scalar must compare the two integrals for all matched pairs,
including the limits specified in Proposition~\ref{geo:theta}, and
at least one matched pair must have a nonzero value.

The third assertion controls how these scalars vary when passing to
smaller spaces. Write \(k=(a-a')/2\). In the primed descent \(q>0\)
the dimensions drop from \((a,a',d)\) to \((a,a'-2,d-2)\).
In the unprimed descent \(q=0<p\) they drop to \((a-2,a',d-2)\).
The upper and lower models use adapted integral hyperbolic planes
in both changing spaces. If their fractional ideals along the first
isotropic lines are \(\mathfrak b,\mathfrak c\), set
\[
 F_X(j)=\nu(\mathfrak b)^{e_1}\nu(\mathfrak c)^{e_2}
    \prod_{v\in S}\nu_v(s)\,
       q_v^{-c_v j}G(\nu_v,\psi_v)\nu_v(\pi_v)^{c_v}.
\]
The ideal evaluations use the fixed local uniformizer values. Each
branch has one triple of signs \(e_1,e_2,s\in\{1,-1\}\) valid across
the fields and all admissible descents in that branch. For fixed
base, spaces, planes and real matching maps a further \(C\ne0\) must
work for \emph{all} allowed characters, Weil actions and finite
matching maps. For upper and lower constants \(c,c_1\), the relations are
\[
 cL^{S_0}(\nu,k+1)=\frac{c_1}{CF_X(-k)}
       \quad(q>0),\qquad
 c=CF_X(1-k)L^{S_0}(\nu,k)c_1\quad(q=0<p).
\]
Lemma~\ref{newtheta:gauss} proves nonvanishing of the finite factors.
For \(u=k,k+1\) one has \(u>1\), and the corresponding partial
Euler product over \(E\) is nonzero. Indeed each logarithmic factor
is bounded in absolute value by
\(\sum_{j\ge1}N\mathfrak p^{-ju}/j\).
There are at most \([E:\mathbb Q]\) primes above each rational
prime \(l\), all of norm at least \(l\), and
\(\sum_l l^{-u}<\infty\). This proves absolute convergence of the
sum of logarithms. These arguments justify division by the displayed
factors. Theorem~\ref{tcr:global} proves the common-constant assertion and the
identities comparing the automorphic functions through the descent.
In particular equality of a constant term detects a difference only up
to the kernel of that constant-term operator.
\src{src/Development/AbelianVarietyTheory/BlockTheory/DivisorBase/WeilTensorModels.lean}{552}{Quantifiers and the two scalar recursions}.

\paragraph{A published comparison at the anisotropic base.}
Let \(A\) have dimension \(m\ge1\) and be definite at a real place.
Any \(F\)-rational isotropic vector would stay isotropic under that
embedding, so \(A\) has Witt index zero over \(F\). For the doubled
space \(Y\oplus(-Y)\) with \(\dim_EY=n\ge1\), the diagonal copy of
\(Y\) is isotropic of dimension \(n\), half the total dimension.
In the standard Schr\"odinger model let \(\Phi\) be a Schwartz test,
\(\ell\) evaluation at zero, \(\omega\) the split Weil action, and
\(f_{\Phi,s}\) the standard flat Siegel section through
\(x\mapsto\ell(\omega(x)\Phi)\) at \(s_0=(m-n)/2\).
The restriction of its splitting character to \(\mathbb A_F^\times\)
is \(\eta_{E/F}^{\,m}\). For this section the Eisenstein series is
holomorphic at \(s_0\), and its value is
\[
 E(x,s_0,f_{\Phi,s_0})=
 \begin{cases}1&m>n,\\2&m\le n\end{cases}
 \int_{[U(A)]}\Theta(g,x;\Phi)\,d\bar g ,
 \qquad \operatorname{vol}_{d\bar g}([U(A)])=1.
\]
This is the unitary specialization of Gan--Qiu--Takeda, Theorem 7.1.
In their notation \(r=0\), \(d(n)=n\) and
\(\tau(U(A))/[E:F]=1\). Their integral in \S3.1 has probability
measure, which agrees with the average over the fundamental domain in
the doubled formula for any choice of Haar measure. These facts fix
both the sign of the critical parameter and the value at \(m=n\).
Gan, Qiu and Takeda, \S\S3.1 and 7.1~\cite{ganQiuTakeda}.

To apply this identity to an arbitrary doubled presentation one must
intertwine its Darboux coordinates, actions and functional with this
standard model, preserving the rational lattice sum. The defining
identities require that the functional be invariant under the
diagonal Lagrangian translations and the source group, and that its
translates along the opposite rational coordinates sum to the theta
distribution. The separate meromorphic continuation input for a
continuous \(K\)-finite section of the doubled group uses the Siegel
law \(\chi(\det t)|\det t|_{\mathbb A_E}^{\,s+n/2}\).
At each \(s_0\) it requires one pole order and a common neighborhood
on which the pole-cleared family is weakly holomorphic for every
functional bounded on each polynomial-growth unit ball. An independent
application of continuation must identify this topology as well as the
section. The next two subsections construct a supported nonzero
undoubled primed period when \(q=0\), compare both periods when
\(p=q=0\), and calculate the constant terms needed to pass between
ranks. The normalization of the smaller matching maps enters at
that last step.
\src{src/Development/WeilBound/ThetaLifts/LatticeBasics/Hypotheses2.lean}{63}{Doubled actions and summation identity};
\src{src/Development/WeilBound/ThetaLifts/LatticeBasics/HypothesesE.lean}{46}{Growth-functional continuation};
\src{src/Development/WeilBound/ThetaLifts/LatticeBasics/HypothesesS.lean}{11}{Doubled normalized value}.

\subsection{Absolute convergence, support at one split place, and the inert quotient}
\label{threal:main}

Write \(d=a+a'\), \(f=[F:\mathbb Q]\), and use the height \(H\) of
\S\ref{newtheta:global}. Put \([G]=G(F)\backslash G(\mathbb A_F)\).
There is a difference between the estimates needed to define the two
periods and the identity comparing their values. We prove the estimates
for the undoubled groups.

\begin{proposition}\label{threal:decay}
For the signatures and ranks of \S\ref{newtheta:global}, the unprimed
theta kernel is absolutely integrable over \([U(V)]\) at every \(h\)
for every adelic Schwartz function. For a primed test \(\Phi'\) in the
specified full-rank space at \(v_1\) and every \(N\geq0\), there exist
\(C_{\Phi',N}\) and an integer \(k\geq0\) such that
\[
 |\Theta_{\Phi'}(g,h)|\,H_{[U(V')]}(g)^N
                      \leq C_{\Phi',N} H(h)^k .
\]
The primed kernel is absolutely integrable in \(g\), its integral has
moderate growth in \(h\), and the full-rank test space is invariant
under the two unitary actions.
\end{proposition}
\begin{proof}
Use \(Y\) for either source, \(m=\dim_EY\), and let \(r\) be its
\(F\)-rational Witt index. Fix rational hyperbolic pairs
\((e_i,e_i^*)_{1\le i\le r}\) and their orthogonal complement \(Y_0\).
In \(Y\otimes_EW\), take one whole half of each hyperbolic pair
tensored with \(W\) as position coordinates and complete by a rational
Lagrangian in the underlying symplectic space of \(Y_0\otimes_E W\).
The theta sum now has positions
\[
        (x_1,\ldots,x_r,y)\in W(E)^r\times F^b,\qquad
             b=(m-2r)d ,
\]
where \(W(E)\) denotes the underlying \(E\)-vector space.
Choose the half for which the torus in the positive chamber acts on
functions by
\[
 \phi(x,y)\longmapsto
       \left(\prod_i T_i^{d/2}\right)
       \phi(t_1x_1,\ldots,t_rx_r,y),\qquad
          T_i=t_i^{2f},\quad t_i>0 ,
 \tag{*}
\]
up to a scalar of absolute value one. Indeed each of the \(d\)
\(E\)-coordinates contributes \(2f\) real dimensions and the
prefactor is the square root of the real change of volume. The torus
has value \(t_i\) at all real places and unit finite components.
With the parabolic stabilizing the expanding isotropic flag one uses
the opposite position half: the action on arguments is inverse to
the action on translations.

The change from the given Darboux coordinates is rational. Indeed the
images of a rational basis are rational by assumption, and their
determinant is nonzero. Rational linear changes, rational quadratic
phases and partial Fourier transforms preserve the absolute value of
the rational summation functional, by the product formula, triviality
of \(\psi\) on \(F\), and Poisson summation with self-dual measures.
These operations generate the rational symplectic group by the Siegel
Bruhat decomposition. They also show that the absolute value of the
kernel is invariant under rational left translates in either group.
These arguments apply to each allowed unitary implementation: two
unitary operators implementing the same action on the Heisenberg
operators differ by a scalar of modulus one. In Schr\"odinger
coordinates their quotient commutes with every modulation and hence
is multiplication on \(L^2\), and commuting with the translations
makes that multiplier constant. This argument also permits the
use of the usual tensor-product implementations for estimates at a
fixed adelic group element.

Apply adelic reduction theory to the connected reductive group \(U(Y)\).
Its center is the norm-one torus and has no \(F\)-split part.
A minimal \(F\)-parabolic is the stabilizer of the chosen flag.
The quotient is covered by a Siegel set with representatives
\(n m_0 t k_0\), with \(n,m_0,k_0\) in fixed compact sets of the
unipotent radical, norm-one part of the centralizing Levi, and
maximal compact subgroup, respectively. Here \(t_i/t_{i+1}\) and
\(t_r\) are bounded below by a positive constant. When \(r=0\) the
set is compact. These are the compactness and fundamental-set theorems
in adelic reduction theory, with inversion to pass from the right
quotient to the left quotient \cite[\S\S2.1--2.2]{springerReduction}.
The elements \(t^{-1}nt\) stay in a compact set, since the inverse
of every positive root is bounded in this chamber. Thus the
representatives have the form \(tc\), with \(c\) in a fixed compact
set. In the integration coordinates above the Haar density in \(t\)
is a constant times
\[
       \prod_{i=1}^{r} T_i^{-(m-2i+1)}\,\frac{dT_i}{T_i}.
\]
Indeed for each \(j>i\) the two relative roots with weights
\(t_i/t_j,t_it_j\) have parameter in \(E\); the root \(t_i\) has
parameters in the \(m-2r\) dimensional space \(Y_0\) over \(E\);
the root \(t_i^2\) has one parameter over \(F\). Their determinant
on the unipotent radical is
\(\prod_i T_i^{2(r-i)+(m-2r)+1}\).
Integration of a nonnegative rationally invariant function over a
covering Siegel set bounds its quotient integral from above.

Here are the bounds on the test in (*) as \(c,h\) vary. Fix coordinates
over \(\mathbb Q\) using bases of the number fields and put \(D=fmd\).
For a pure archimedean and finite tensor, the finite part of
\(\omega(c,h)\Phi\) is supported on \(M^{-1}\widehat{\mathbb Z}^{D}\)
and has supremum at most \(C M^{D/2}\), for an integer
\[
             1\le M\le C_\Phi H(h)^A .
\]
To see this, start with an integer \(M_0\) such that translations and
modulations by \(M_0\widehat{\mathbb Z}^{2D}\) fix the finite test,
using the standard \(\mathbb Q\)-adelic character and its dual
coordinates. Increasing \(M_0\) also fixes coordinate denominators.
After conjugation it suffices to have
\((c,h)^{-1}M\widehat{\mathbb Z}^{2D}\subset
 M_0\widehat{\mathbb Z}^{2D}\) in the symplectic representation.
At each place \(v\), put \(n_v\) equal to the largest denominator
order of the entries of \(h,h^{-1}\), or zero if all are integral,
using all factors of \(E\otimes F_v\). The extension \(E/F\) is
unramified at finite places. In fixed bases the entries on tensors
involve only matrix multiplication and fixed changes of coordinates.
All the \(c\)'s carry bounded denominators, supported at finitely many
primes. Thus \(M=C_0\prod_{p}p^{\max_{v\mid p}n_v}\) suffices for
a fixed positive integer \(C_0\). Since \(p\le q_v\) and the finite
height factor is at least \(q_v^{n_v}\), this proves the bound.
Invariance by modulations forces the stated support; invariance by
translations makes the function constant on cosets of
\(M\widehat{\mathbb Z}^{D}\), each of volume \(M^{-D}\) in product
measure of integral volume one. Unitarity gives the supremum bound.

At infinity, for every \(B\ge0\), the absolute value of the transformed
test at \(z\) is at most \(C_B H(h)^{A_B}(1+\|z\|)^{-B}\).
Indeed commute multiplication and differentiation through the unitary
operators using Heisenberg covariance. Each is a linear combination of
the original position and momentum operators, with coefficients
bounded by a fixed power of the maximum of the archimedean matrix
norms of \(h,h^{-1}\) and a constant for \(c\). Apply this identity to
words of length at most \(B+D+2\) (round \(B\) up to an integer).
Euclidean Sobolev embedding applied to each monomial times the
transformed function then gives the assertion. Finite sums of pure
tensors satisfy the same majorant. In these estimates \(A_B\) may
depend on the fixed spaces, \(B\), and the test. The specified height
uses the local spectral absolute values over \(F_v\), extending
\(|\pi_v|=q_v^{-1}\), and takes the maximum of \(1\) and the norms of
both matrices.
\src{src/Development/AbelianVarietyTheory/Definitions/Foundations/EventuallyNormFinLeOne.lean}{405}{Matrix height across the places}.

For the unprimed integral fix \(h\); then \(M\) is bounded and
\(t_i\gg1\). Summing the Schwartz majorant with \(B>D+1\) on
\(M^{-1}\mathbb Z^D\) gives \(O_{\Phi,h}(\prod_i T_i^{d/2})\).
For \(r\le p\) every exponent \(d/2-(a-2i+1)\) is strictly negative,
since \(a-d/2=(a-a')/2>2(p+q)+1\). Integration even over the larger
rectangle \(T_i\ge c_1>0\) converges.

For the primed integral we verify the support after the change of
polarization. At \(v_1\) take the specified conjugate split coordinates,
so the original position variable is an \(m\)-by-\(d\) matrix over
\(F_{v_1}\). The two unitary actions there are invertible row and
column changes times a scalar, since they preserve each of the two
split Lagrangians. Therefore they preserve full row rank. This also
proves invariance of the adelic span by factorization, since all other
factors stay Schwartz and a restricted adelic element fixes the
standard lattice factors up to scalar almost everywhere. For clarity,
write the conjugate split matrix first in the rational Witt row basis.
In a hyperbolic pair the two position rows are indexed by \(e_i,e_i^*\).
In the new model the \(E\otimes F_{v_1}\)-valued position block \(x_i\)
contains \emph{one unchanged row} and the Fourier dual of its paired
row. Indeed the tensor symplectic pairing pairs the \(e_i\)-row on
one split factor with the \(e_i^*\)-row on the other and conversely.
One exchanges the position and momentum in one of these two row
pairs, using an invertible linear identification with \(W\) from
its form. The remaining change of polarization takes place in the
\(Y_0\)-block. If the unchanged row vanishes, the original matrix has
deficient rank for every value integrated over in these partial
Fourier transforms. The transformed function therefore vanishes when
\(x_i=0\). This holds separately for each \(i\) and survives finite
sums. The coordinates called for in the support definition are unique
up to scalar on Schwartz--Bruhat functions: an intertwiner with the
stated translation and modulation laws is unique by the local
Schr\"odinger commutant. Algebraically this follows by averaging
modulations over compact open subgroups to project onto support in
any ball, and then translating the indicators of balls.

For a surviving rational summand \(\|x_i\|\ge c_2/M\): the rational
points in the indicated finite support lie in \(M^{-1}\mathbb Z^D\),
and the coordinates were chosen block by block. Given \(L\ge0\),
spend \(rL\) powers in the archimedean majorant to obtain
\(\prod_i (1+t_i\|x_i\|)^{-L}\ll M^{rL}\prod_i t_i^{-L}\).
Another \(D+2\) powers suffice to sum the remaining majorant over
\(M^{-1}\mathbb Z^D\), with bound \(O(M^D)\) by unit-cube counting.
Thus, for every \(J\ge0\),
\[
             |\Theta_{\Phi'}(tc,h)|
                    \ll_{\Phi',J} H(h)^{A_J}\prod_i T_i^{-J}.
\]
Here \(L\) was chosen to absorb both \(2fJ\) and the prefactor in (*).
The quotient height at any translate is at most \(H(tc)\), which on
this set is bounded by \(C\prod_i T_i^{A_0}\), by the matrix formula
for \(H\). Taking \(J\ge NA_0\) gives the claimed inequality.
Integrating the bound for any sufficiently large \(J\) proves
absolute convergence and moderate growth, including the case \(r=0\)
with empty products. All bounds use absolute values, so they apply
to the prescribed Haar fundamental domains.
\end{proof}

\begin{corollary}\label{threal:limits}
An identity with one fixed scalar for all polynomially matched pairs
whose primed input satisfies the support condition implies the same
identity for every matched pair in \S\ref{newtheta:global}.
\end{corollary}
\begin{proof}
Fix \(h\) and a finite Schwartz factor. The unprimed proof bounded the
absolute integral by finitely many archimedean Schwartz seminorms.
The same is true on the primed side for a fixed admissible finite
factor: the integer \(M\), the size of the finite factor and the
compact set used in reduction are fixed, and the Sobolev argument
uses finitely many seminorms of the varying archimedean factor.
For the primed approximations the individual finite factors in the
definition of matching need not themselves satisfy the support
condition. Let \(B\) be their finite-dimensional span and \(B^\circ\)
the subspace of the full finite Schwartz space spanned by admissible
finite tensors. For \(S=\mathcal S(\mathbb R^D)\) one has
\[
      (S\otimes B)\cap(S\otimes B^\circ)=S\otimes(B\cap B^\circ)
\]
inside the algebraic tensor product with all finite Schwartz tests,
as follows by extending a basis of \(B\cap B^\circ\) in each factor
space \(B,B^\circ\). Products of functions realize this tensor product
faithfully: evaluations span the dual of every finite-dimensional
space of functions. Express the approximants and their limit in a
basis of \(B\) adapted to \(B\cap B^\circ\). The coefficients outside
the intersection vanish and the others converge in \(S\), being fixed
finite linear combinations of the original coefficients. Continuity
on each resulting finite slice permits passage to the limit at \(h\).
\end{proof}

\begin{proposition}[Even source ranks at an inert place]\label{threal:inert}
Suppose \(v\) is inert and \(a=2u,\ a'=2w\). The two local
Schwartz--Bruhat spaces have a common irreducible \(U(W)(F_v)\)
quotient, invariant under their respective source groups. The maps
can be chosen to send the prescribed vacua to the same nonzero
spherical vector. Its unramified parameter, in the coordinates
\(q^{2\sigma_i}\) up to permutation and separate inversions, is
\[
      (q,q^3,\ldots,q^{2u-1})\ \sqcup\
      (q,q^3,\ldots,q^{2w-1}),\qquad q=|\mathcal O_F/v|.
\]
\end{proposition}
\begin{proof}
Put \(k=F_v,\ l=E_{\mathfrak w}\), and \(R=u+w\). This is an
unramified quadratic extension of odd residual characteristic, \(q>3\),
and \(\psi_v\) has conductor \(\mathcal O_k\). Since all ramification
of \(\alpha,\beta\) is split, their local components are unramified.
Both are trivial on \(k^\times\) by evenness of the ranks. A
uniformizer of \(k\) is also a uniformizer of \(l\), so both local
characters, and both splitting pairs, are trivial.

The two source lattices are perfect Hermitian lattices. The absolute
trace-dual condition on the target also gives a perfect skew-Hermitian
lattice here. Indeed the standard global additive character with the
specified infinite component has local annihilator of \(\mathcal O_k\)
equal to the absolute codifferent; this follows from the trace
pairing and the uniqueness of the character on \(\mathbb A_F/F\)
with that infinite component (the dual is \(F\)). The conductor
hypothesis makes this codifferent the unit ideal, and \(l/k\) adds
no different. A trace-zero unit converts a perfect skew-Hermitian
form into a Hermitian one. Such a unit is obtained by subtracting
the conjugate of a lift of a residue element outside the base field.
Every even-rank perfect Hermitian lattice here is integrally
hyperbolic. Over the residue fields diagonalize and use surjectivity
of the norm to obtain a nonzero isotropic vector in dimension at
least two. The derivative of the norm form there is the surjective
residue trace pairing with this vector. Hensel lifting gives a
primitive isotropic \(x\). Choose \(y\) pairing to one with \(x\)
and replace \(y\) by \(y-h(y,y)x/2\). The resulting hyperbolic plane
splits integrally, giving induction. In particular the target admits
the integral hyperbolic polarization and both source groups have the
quasi-split sign \(+\).

Use the ordinary Schr\"odinger model on \(Y^R\), for \(Y\) one of the
two Hermitian sources. The lattice vector is, up to scalar,
\(\mathbf1_{\Lambda_Y^R}\). To verify the identification with the
prescribed vacuum, apply its Heisenberg eigenline condition on each
of the two tensor lattice halves of the hyperbolic polarization of
\(W\). The quadratic phase is one on each half, which is isotropic
over \(l\) and stable under the trace-one element. The two conditions
are invariance under the translations by \(\Lambda_Y^R\) and the
modulations by its dual. They force support in \(\Lambda_Y^R\) and
constancy there. The change to these coordinates is a unitary
Heisenberg intertwiner and the original vacuum is nonzero.

We recall why the splitting, including its scalar character, agrees
with the one in the spherical calculation. The prescribed splitting
on each factor extends to its hyperbolic doubling. In the doubled
diagonal polarization evaluation at zero has on its Siegel stabilizer
the law \(\chi(\det t)|\det t|_l^{\dim_EY/2}\), with the oscillator
partner's dimension in the exponent if it is the source that is
doubled. This is precisely the standard doubled splitting with
character \(\chi\). To use the adelic law at a single place, take
evaluation in the diagonal polarization at all places. This is an
admissible functional: rational changes to that polarization are
compositions of the linear changes, phases and Fourier transforms
above, and evaluation is continuous on each fixed finite Schwartz
slice. The functional is nonzero. After unitary identification of the
doubled Heisenberg implementations, their quotient on elements
supported at \(v\) is a scalar character by the commutant argument.
The common evaluation law makes it trivial on the Siegel stabilizer.
Its conjugates generate the hyperbolic unitary group: the Siegel
Bruhat double cosets are represented by interchanging some of the
hyperbolic pairs. In each hyperbolic plane the interchange, up to
a diagonal element, is the product of three alternating upper and
lower root unipotents by the \(2\)-by-\(2\) identity in
\(\mathrm{SL}_2(k)\simeq SU(1,1)(k)\). Upper unipotents are in the
stabilizer and lower ones in its opposite conjugate. Thus the
character is trivial. Restrict on block elements \((g,1)\) of the
doubling and apply the defining product rule on pure tensors to
recover the given single-factor action. This argument works on both
factors. In particular it excludes a scalar character twist that
Heisenberg covariance alone would permit.
\src{src/Development/AbelianVarietyTheory/Definitions/Divisors/HilbertNinetyQuotientHomEqIff.lean}{435}{Doubled Siegel and block restriction laws}.

In Liu's notation take the larger skew-Hermitian group \(G_R\) and
the Hermitian space of dimension \(2s\), first \(s=u\), then \(s=w\).
His unramified additive character and trivial splitting character are
the ones just checked (a unit scaling in the trace pairing changes
only the choice of conductor-zero additive character or hyperbolic
basis). The trivial representation of \(U(2s)\) has parameters
\((q,q^3,\ldots,q^{2s-1})\), since the exponents of the normalized
Borel modulus on its \(l^\times\)-coordinates are
\(s-i+1/2\), for \(1\le i\le s\). In fact the root computation in
the preceding proof gives the modulus before taking its square root.
Liu's Theorem~6.2(2), Proposition~3.4 and the parameter maps of
\S2 say that, for \(R\ge s\) and sign \(+\), the small lift of the
spherical representation on \(G_R\) with the \(s\) parameters inverted
and \(q,q^3,\ldots,q^{2(R-s)-1}\) appended is the spherical
representation with those \(s\) parameters on \(U(2s)\).
The small-lift quotient sends \(\mathbf1_{\Lambda_Y^R}\) to a
nonzero vector, also explicitly in the proof of his Theorem~1.1.
The two applications give the displayed identical Weyl orbit.
Accordingly the oscillator has in each case a quotient onto
\(\pi\boxtimes{\bf1}\) with the same irreducible spherical \(\pi\)
on \(G_R\). The image of each vacuum belongs to the one-dimensional
hyperspecial fixed line; rescale one map to give equality. If maps
on the full space of all functions are requested, extend linearly
from Schwartz--Bruhat functions using any vector-space complement:
the conditions concern only Schwartz--Bruhat vectors.
For these choices of spherical vector, parameter and normalized
induction see \cite[\S\S2, 3 and 6]{liuThetaInert}.
\end{proof}

\begin{proposition}[All inert places, including the odd row ranks]
\label{threal:allinert}
At every inert place, in either parity, the two local Schwartz spaces
have a common irreducible target quotient annihilating the source
differences. The prescribed vacua have the same nonzero image after
one scalar normalization.
\end{proposition}
\begin{proof}
Write \(d=2R\), \(Y=V\) or \(V'\), \(m=\dim_EY\), and
\(\xi=\beta_v\) or \(\alpha_v\), respectively. These last two
characters of \(E_v^\times\) are the same: both are unramified
and their values at a uniformizer of \(F_v\), also a uniformizer
of \(E_v\), are \((-1)^m\). Thus \(\xi^2=1\). On the source the
splitting character is trivial in both cases. Use the integral
hyperbolic polarization of \(W\) in the preceding proof. In this
model \(U(Y)\) acts by linear substitution on \(Y^R\): in the
Siegel Levi implementing its action the determinant modulus is
one, with trivial splitting character. Evaluation at zero is
source invariant. Consequently
\[
        \phi\longmapsto
          \bigl(h\longmapsto(\omega(1,h)\phi)(0)\bigr)
\]
gives a nonzero equivariant map for the target into the normalized
Siegel induction \(I(s_m,\xi)\) of \(U(2R)(F_v)\) from
\((\xi\circ\det)|\det|_{E_v}^{s_m}\),
where \(s_m=(m-R)/2\). Indeed the Siegel unipotent acts by a
quadratic character taking value one at zero. The Levi acts at
zero by \(\xi(\det b)|\det b|_{E_v}^{m/2}\), and the half modulus
of this Siegel parabolic is \(|\det b|_{E_v}^{R/2}\).
The image is nonzero already on the lattice indicator at \(h=1\).
These substitutions are the standard Siegel splitting of the
polarized Weil representation. The doubled uniqueness argument
above identifies them with the actions specified here.

The image has finite length as a submodule of parabolic induction
of an admissible finite-length representation. Choose any
irreducible quotient of this nonzero image. The composite is a
nonzero map from the oscillator onto the trivial source times
this quotient. The unramified vacuum theorem says that such a map
takes the lattice vacuum to a nonzero vector in the hyperspecial
fixed line when the source representation is spherical. Its
hypotheses are precisely odd residue characteristic, an unramified
quadratic extension and additive character of conductor \(\mathcal
O_v\), self-dual lattices and unramified splitting characters. To
interchange the Hermitian and skew-Hermitian spaces one can use a
trace-zero unit and its inverse. The source vacuum in the given
coordinates is a nonzero multiple of this lattice vector by the
Heisenberg characterization above.
\cite[Appendix, Theorem~2.A.1]{harrisSpherical}

In particular this irreducible quotient has the spherical Hecke
character of \(I(s_m,\xi)\): hyperspecial invariants are exact by
compact averaging, and the induction has a single fixed line by
Iwasawa decomposition. For \(m'=2R-m\) one has \(s_{m'}=-s_m\).
The two inducing data are Weyl conjugate on their unramified
parameters: the long Siegel Weyl element takes \(b\) to its
inverse conjugate transpose, \(\xi\) is conjugation invariant
and self-inverse, and the \(R\) exponents of the trivial
\(\mathrm{GL}_R\) representation are symmetric about zero.
The Satake classification thus identifies the two spherical
irreducibles. Scaling the maps to their common fixed line gives
the desired normalization.
\end{proof}

\begin{proposition}[Ordinary relatively split places]
\label{threal:ordinarysplit}
At every place split in \(E/F\) outside \(S=S_0\sqcup Q\), there
is a common irreducible quotient with matching nonzero vacua,
including when the underlying rational prime does not split
completely in \(E\).
\end{proposition}
\begin{proof}
Use the two row and dual-row models of
Proposition~\ref{newtheta:split}. The relative splitting
\(E\otimes F_v=k\times k\) suffices for these models. In either
polarization the two idempotent summands of the tensor lattice give
translations by a full lattice and modulations by its character dual.
Indeed the source lattice is perfect for its \(E\)-valued form,
and the partner lattice is perfect for the absolute trace pairing.
Thus, in dual bases, the modulation coefficients range over the
inverse different of \(k/\mathbb Q_\ell\) times the ordinary dual
of the position lattice. This gives exactly the character dual for
the trace additive character. The latter is forced by its infinite
component and triviality on \(F\), since \(F\) is dense in its
finite adeles. On each idempotent summand separately the quadratic
phase in the vacuum condition is one by isotropy, also in residue
characteristic two. Translation invariance and modulation invariance
force any nonzero vacuum to be a constant times the indicator of
the position lattice. Choose position lattice bases and the
conjugate dual bases so its indicator is the integral-matrix test
on each side and the two target compact subgroups agree. Scaling
these local coordinate intertwiners identifies the given vacua
with the two indicators. The unitary linear action is the prescribed
one by the block and Siegel laws used above; conjugating to the dual
polarization sends the target action to inverse transpose as in
\eqref{newtheta:big}.

In the notation of Proposition~\ref{split:good} put \(n=a\),
\(t=\nu(\pi)\), \(s=d/2\) and
\(D_n(z)=\prod_{j=0}^{n-1}(1-tq^{j-z})\).
For any matrix test \(H\), the cleared determinant integral
\(D_n(z)Z_n(H,z)\) is a Laurent polynomial in \(q^{-z}\).
Here is the refinement needed when \(t\) has finite order. Keep the
integrand exponent \(z-n\) fixed as in that proposition. In size
\(j\) a ball about zero scales to \(M_j(\mathcal O_v)\). Its
denominator divides
\(\prod_{i=0}^{j-1}(1-tq^{n-j+i-z})\) by the same Hermite
count with parameter \(z-n+j\) and the conversion from additive
to group Haar measure. Its numerator is a constant times a
Laurent monomial. Subtract the value at zero times such a ball.
The remaining compact support is covered by pivot charts; the
Schur-complement argument reduces to size \(j-1\), whose denominator
is a subproduct of the one just displayed. Finite disjoint
subdivision and induction prove the claim starting with size zero.

Multiply the larger transform by \(D_n(z)\) and evaluate at \(s\).
The source and parabolic covariance pass from the right half-plane
to this value by the meromorphic identity principle, including if
the value is a residue. The character and modulus factors in that
covariance are holomorphic in \(z\). Smoothness in the target
follows from the compact open stabilizer of the test. The larger
vacuum has value one at the identity by the exact Hermite count,
with \(\operatorname{vol}(\mathrm{GL}_n(\mathcal O_v))=1\).
The smaller vacuum also has nonzero value there, since \(a'<s\),
and can be normalized to one. The two spherical vectors are now
equal. The common induced target of Proposition~\ref{newtheta:split}
is irreducible by Bernstein's theorem, so both nonzero maps are
onto. The stronger assertion on the entire source-difference
kernel at totally split good primes was proved separately in
Proposition~\ref{split:good}.
\end{proof}

\begin{proposition}[The real quotient for \(p=q=0\)]\label{threal:compact}
At each infinite place write \(A=m_\alpha,\ B=m_\beta\) and
\[
        j=(A-B-d)/2,\qquad l=(B+a)/2 .
\]
The two real polynomial modules have the same source-coinvariant
representation of the compact target \(U(d)\), with dominant highest
weight \(((l+j)^a,l^{a'})\). Superscripts indicate repetitions.
The tensor product of these representations over the real places
realizes the real matching maps for this branch.
\end{proposition}
\begin{proof}
The restrictions to \(\mathbb R^\times\) give \(A\equiv a'\),
\(B\equiv a\pmod2\). Thus \(j,l\) are integers, and \(j>0\) by the
weight condition. Let \(D_0=\mathbb C^d\), \(U_0=\mathbb C^a\),
\(U_1=\mathbb C^{a'}\) with the standard actions. In orthonormal
column frames the operators on the polynomial Fock modules are the
following two representations of source times target:
\[
\begin{split}
 &(\det_{U_0}^{-j}\boxtimes\det_{D_0}^{\,l})
               \otimes\operatorname{Sym}(U_0\otimes D_0),\\
 &(\det_{U_1}^{\,j}\boxtimes\det_{D_0}^{\,l+j})
               \otimes\operatorname{Sym}(U_1^*\otimes D_0^*).
\end{split}
\]
We verify the weights to fix the sign convention. In one unprimed
complex matrix entry the tensor form is
\(-\operatorname{Tr}_{\mathbb C/\mathbb R}(-i z\bar z')
=2\,\operatorname{Im}(\bar z z')\); it changes sign on the primed
side. The additive character is \(x\mapsto e^{-2\pi i x}\). In the
positive plane the Heisenberg derivatives for \(1,i\) can be taken
as \(\partial_x,-4\pi i x\). Acting on
\(\exp(-2\pi x^2)\), the second derivative operator is \(i\) times
the first. Thus a rotation acts on the creation operator
\(d\rho(1)-i\,d\rho(i)\) by its ordinary character. Reversing the
form gives the conjugate character. The infinitesimal metaplectic
rotation on the positive plane has eigenvalues \(i(n+1/2)\):
in the creation and annihilation operators \(b^\dagger,b\) with
\([b,b^\dagger]=1\) its generator is the symmetrized quadratic
\(i(b^\dagger b+bb^\dagger)/2\); the negative plane reverses this
sign. Therefore the metaplectic vacuum has half the sum of the signed
entry weights.

The doubled Siegel normalization contributes the additional half of
the splitting-character exponent on the determinant of each single
unitary factor. Indeed on the doubled Siegel Levi \(t\in
\mathrm{GL}_n(\mathbb C)\) the action of the metaplectic Lie algebra
on evaluation gives the positive modulus; its infinitesimal phase
is zero, as one also sees from the real change of coordinates.
The prescribed additional phase on this Levi is
\((\det t/|\det t|)^m\). The determinant of the doubled unitary
element on the same Levi is \(\det t/\overline{\det t}\); hence the
difference of the two infinitesimal implementations is
\((m/2)\operatorname{tr}\). They differ by a Lie algebra character,
and every such character on \(\mathfrak u(n,n)\) is a multiple of
trace since its derived algebra is \(\mathfrak{su}(n,n)\).
Restriction to the first block preserves the trace. Consequently the
source vacuum powers are \((d+B-A)/2=-j\) and
\((-d+A-B)/2=j\); the target powers are \((a+B)/2=l\) and
\((-a'+A)/2=l+j\). The unitary groups are connected, so the
infinitesimal action determines these integrated representations.
The Gaussian times the Hermite polynomials form the polynomial
module here: the annihilation equations for the decaying Gaussian
separate into the one-coordinate first-order equations above, with
the opposite choice of sign giving an increasing exponential, and
the Heisenberg derivatives generate all the creation monomials.

Use the Cauchy decomposition
\(\operatorname{Sym}(U\otimes D_0)=
\bigoplus_\lambda S_\lambda U\otimes S_\lambda D_0\);
this follows by applying Schur--Weyl duality to each homogeneous
degree and taking simultaneous permutation invariants.
The only partition for which \(S_\lambda U_0=\det_{U_0}^j\) is the
rectangle \((j^a)\), with multiplicity one. On the dual side the
rectangle is \((j^{a'})\). Thus the first invariant space is
\(\det^l S_{(j^a)}D_0\), and the second is
\(\det^{l+j} S_{(j^{a'})} D_0^*\). Their highest weights agree,
since duality reverses and negates the list of weights. Both spaces
are irreducible. Averaging over the source compact group projects
onto the invariants and kills precisely the span of the differences
on the polynomial module, by complete reducibility on each finite
degree. Composing with an identification of the two highest-weight
modules gives the claimed surjections. Their tensor product is
irreducible for the product of the compact targets. The maps may
be extended linearly off the polynomial spaces, since the
specified infinitesimal and compact transformation laws concern
polynomial inputs and these actions preserve that subspace.
For comparison with the source conventions see
\src{src/Development/WeilBound/ThetaLifts/LocalTheory/ToralCohomology/PiInfUnitNe.lean}{116}{The unprimed target highest weight}.
\end{proof}

\begin{proposition}[The real quotients for all signatures]\label{threal:allreal}
At each infinite place let \(A=m_\alpha,\ B=m_\beta\).
In Paul--Trapa's normalization the common real representation
is the cohomologically induced module on the ordered Levi
\[
 U(a-p,p)\times U(a'-q,q),\qquad
 A_{\mathfrak q}(\det^{A/2}\boxtimes\det^{B/2}).
\]
The half-determinant notation uses the covers in their normalization.
The output after twisting is an honest module for the target Lie
algebra and its maximal compact group.
\end{proposition}
\begin{proof}
Put \(r=d-p-q,\ t=p+q,\ k=A-B>d\). Work first with the
metaplectic oscillator, before the two splitting twists. With the
additive character and tensor form in Proposition~\ref{threal:compact},
a positive source row has vacuum determinant exponent \((r-t)/2\);
a negative source row has the opposite exponent. This follows by
summing the signed one-entry weights \(1/2\) there. It is the
convention of Paul--Trapa's joint harmonics, their formula (3.3).
In passing to the prescribed actions the source determinant exponents
are shifted by \((B-A)/2\) and \((A-B)/2\), respectively.
The doubled trace argument for that shift in
Proposition~\ref{threal:compact} applies to each signature, since
the doubled group and its Siegel Levi have the same trace character.
Thus the prescribed trivial source corresponds in the metaplectic
oscillators to \(\det^{k/2}\) and \(\det^{-k/2}\).

Since \(|r-t|<k\), subtracting the vacuum weights from
\((k/2,\ldots,k/2)\) gives positive entries on both compact
blocks of \(U(a-p,p)\). In the notation of Paul--Trapa's Lemma~3.2
the two joint-harmonic bounds are \(a-p\le r,\ p\le t\).
For the primed determinant all the entries after subtraction
are negative. The bounds for \(U(q,a'-q)\) become
\(a'-q\le r,\ q\le t\). They hold since \(r=(a-p)+(a'-q)\)
and \(t=p+q\). Thus each character occurs in the joint harmonics.
Part (1) of their Theorem~4.1 gives on the unprimed side
\(A_{\mathfrak q}(\det^{k/2}\boxtimes 1)\) with first block
\(U(a-p,p)\). Part (2) gives
\(A_{\mathfrak q}(1\boxtimes\det^{-k/2})\) with second block
\(U(a'-q,q)\) on the primed side. These are the same ordered
two-block parabolic before assigning its two inducing characters.
Finally the target splitting adds \(B/2\) to both block weights
on the first side and \(A/2\) on the second. Tensoring by this
determinant character commutes with cohomological induction, by
pulling the character through the induction and derived compact
invariant functors. Both modules therefore have the displayed
description. The shifts inherent in the authors' normalization
of \(A_{\mathfrak q}\) are identical on the two sides.
\cite[\S\S2.1, 2.3, 2.8, Lemma~3.2, Theorem~4.1 and Corollary~4.3]{paulTrapaHowe}

These are small theta quotients of the polynomial Fock modules,
hence irreducible nonzero targets of equivariant surjections.
Taking their algebraic tensor products over the real places gives
the quotient for the product Lie algebra and maximal compact.
Indeed the source polynomial module is spanned by finite products
of creation monomials from the separate places. A submodule of the
external product of simple Harish-Chandra modules is either zero
or the whole product, by the density theorem applied to finite
tuples in each factor (their commuting endomorphisms are scalars).
The quotient maps kill source derivatives and compact differences;
target derivatives and compact actions preserve matching. Extending
the maps linearly from polynomial vectors to their ambient
Schwartz spaces satisfies the stated real-map conditions, all of
which test polynomial inputs. For simultaneous archimedean
one-parameter groups one differentiates the finite product of
the individual actions, giving the sum of the placewise derivatives.
\end{proof}

\begin{theorem}[A nonzero supported undoubled period]\label{threal:nonzero}
For \(q=0\) and every allowed \(p\) the primed period is nonzero for some polynomial
Schwartz test satisfying the imposed full-rank condition at \(v_1\).
Both parities of \(a'\) are allowed.
\end{theorem}
\begin{proof}
Put \(Y=V'\), \(m=a'\), \(H=U(Y)\), \(J=U(W)\). The source quotient
is compact by the rank-zero case of reduction theory. Take probability
quotient measures for now; other Haar normalizations multiply the
periods by positive constants. We first prove the identity
\begin{equation}\label{threal:norm}
 \int_{[J]}\left|\int_{[H]}\Theta_\phi(g,h)\,dg\right|^2 dh
  = c_H\int_{H(\mathbb A_F)}
                \langle\omega(g,1)\phi,\phi\rangle\,dg,
                 \qquad 0<c_H<\infty .
\end{equation}
Here the scalar simply accommodates the choice of group Haar measure
on the right; it is one for quotient volume one with rational counting
measure. We use the usual self-dual Hilbert pairing, linear in the
first variable. For the proof it suffices to take finite pure tensors
with polynomial Fock factors at infinity.

On the doubled space \(Y\oplus(-Y)\), change from the product
polarization for \(\phi\otimes\bar\phi\) to the diagonal polarization.
The value at zero in the latter coordinates is
\(\int\phi(x)\overline{\phi(x)}\,dx\): the change is a difference and
average of the two position coordinates followed by Fourier transform
in the average variable. The Jacobian scalars can be fixed using
rational coordinates and globally self-dual measures, by the product
formula. At \((g_1,g_2)\) the analogous functional gives
\(\langle\omega(g_1,1)\phi,\omega(g_2,1)\phi\rangle\) times the
second-block determinant character of the doubled splitting. In the
product with the conjugate kernel that character is divided out.
Its exact conventional form is \(\chi_W(\det g_2)\) for the splitting
character on the \(H\)-action in the notation of the doubling
construction. Its modulus is one and it is trivial on rational points.
This block identity is also the change-of-model identity in
Gan--Qiu--Takeda, \S11.3, equations preceding (11.2)~\cite{ganQiuTakeda}.
The rational summation on the two product factors changes to the
doubled rational summation by Poisson summation.

For the given presentations, transport the standard implementations
using their rational Darboux coordinates. On each doubled factor the
possible discrepancy with the prescribed implementation on elements
supported at a single place is a scalar character trivial on the
Siegel stabilizer, by the argument of
Proposition~\ref{threal:inert}. This argument applies over the
real fields as well, using hyperbolic planes. At split places the
doubled group is a general linear group, generated by the block
parabolic and its conjugates by elementary elimination. Hence every
single-place discrepancy is trivial on each single factor by the
block rule. The global discrepancy is a continuous character:
compare the two implementations with a nonzero matrix coefficient
in a neighborhood of the identity, using unitarity and weak
continuity. It is trivial on the dense subgroup of elements supported
on finitely many places, hence everywhere. As a result the standard
automorphic theta and these periods agree after the rational change
of coordinates, up to the single scalar inherent in choosing the
unitary coordinate transform. Taking its modulus to be one suffices
in \eqref{threal:norm}.

To interchange the three integrations, use the chamber estimate of
Proposition~\ref{threal:decay} with \(W\) as integrating space of
dimension \(d\) and \(Y\) as partner of dimension \(m\). Its proof
applies to a skew-Hermitian integrating space as well. If a
Hermitian integrating space is used in the statement of the
Siegel--Weil formula, multiply its given skew-Hermitian form by
any \(0\ne\delta\in E\) with \(\bar\delta=-\delta\) and multiply
the other form by \(\delta^{-1}\). The product of the forms, the
underlying isometry groups, the flags and their moduli are preserved
in these conventions.
For \(g\) in compact source representatives the single kernel is
\(O(\prod T_i^{m/2})\), uniformly in \(g\). The squared bound is
integrable since the \(F\)-Witt index of \(W\) is at most \(p\)
and \(d-2p+1>m\). Now in the doubled Siegel--Weil formula integrate first over \([J]\),
with integrating space of dimension \(d\) and doubled group \(U(2m)\).
The parameters are \(s_0=(d-m)/2>m/2\), as \(a>a'\).
The coefficient of the probability theta integral is one by Weil's
convergent-range formula, applied with \(d>2m\); in the anisotropic
case this is also Theorem~7.1 of Gan--Qiu--Takeda. The Siegel Eisenstein
series with normalized induction is absolutely and locally uniformly
convergent in \(\Re s>m/2\); absolute convergence holds uniformly
on compact representatives of \([H]^2\) by majorization by the
positive spherical section at the real parameter. Thus integration
over \([H]^2\) permits absolute unfolding.
We use the probability normalization for the theta integral
\cite[introduction, Theorem~2 and \S3]{yamanaUnitary}.

Every maximal rational isotropic subspace of \(Y\oplus(-Y)\) projects
injectively onto each copy of \(Y\) since \(Y\) is anisotropic. Its
dimension is \(m\), so it is the graph of a rational isometry.
The set indexing the Siegel series is therefore one
\((H\times H)(F)\)-orbit with stabilizer \(\operatorname{diag} H(F)\).
After unfolding we integrate the matrix coefficient above, with its
second-block character cancelled, over
\(\operatorname{diag}H(F)\backslash H(\mathbb A_F)^2\).
The change \(g=g_2^{-1}g_1\) gives \eqref{threal:norm} and proves
absolute convergence of its right-hand integral. The same calculation
with any finite replacements of the test factors is valid.

We choose the factors so that the last integral is nonzero. At infinity
the primed source is negative definite. In the entry calculation of
Proposition~\ref{threal:compact} there are \(d-p\) negative planes
and \(p\) positive planes in each source row. Its vacuum determinant
power as a source representation is
\((A-B-d+2p)/2=j+p>0\). The \(d-p\) columns from the negative
planes are copies of the dual standard source representation.
Choose \(m\) of those columns, possible since \(d-p\ge m\), and
take their determinant to the power \(j+p\) times the vacuum.
This is a nonzero source-invariant polynomial vector. Its local matrix coefficient on
the compact source is the positive constant \(\|\phi_v\|^2\).
At a designated finite split place the conjugate polarization has
positions \(X\in M_{m,d}(F_v)\) and the source acts by left row
operations times a scalar character and the square root of the
modulus. Take the indicator of a sufficiently small compact open
ball about \((I_m\mid0)\). If \(g\) moves this ball to intersect
itself, the invertible leading block implies
\(g\in 1+\pi^b M_m(\mathcal O_v)\) for a prescribed \(b\);
this follows by writing \(g=X_1X_2^{-1}\) on the first blocks, or
its inverse if the action is inverse. Choose \(b\) so the character
is trivial there. The coefficient is nonnegative, vanishes off a
compact set and is positive on a neighborhood of \(1\).
This test satisfies the full-rank condition at \(v_1\).

At any specified inert place \(W\) has an integral hyperbolic
polarization by the lattice argument above, irrespective of the
parity of \(m\). The source acts on \(Y^{d/2}\) by linear changes
on all vectors times a scalar character. Now \(d/2>m\). Pick a small
ball about a tuple whose first \(m\) vectors are a basis of \(Y\).
The same leading-block argument inside \(\mathrm{GL}_m(E_{\mathfrak w})\)
forces overlap only in any prescribed neighborhood of \(1\) in
\(H(F_v)\). The scalar on this compact subgroup is trivial after
shrinking: a continuous character of a profinite group has open
kernel. The corresponding local integral again is strictly positive.

It remains to ensure that the unrestricted product of local
integrals is nonzero. Start with the factors just chosen at infinity
and \(v_1\), and the usual unit-norm lattice indicators at almost
all other places. There the hyperspecial group \(K_v\) fixes the
indicator: in the split linear model the unramified character is
trivial on units; in the inert hyperbolic polarization the source
acts by pullback with the unramified splitting. More explicitly the
character on the source for this polarization is trivial when the
partner is hyperbolic and its splitting character is trivial, as
is the case for the source splitting \(\alpha/\beta\) at inert places:
the two unramified characters have equal value on a uniformizer
by their common parity. The doubled Schr\"odinger model makes this
pullback implementation explicit, and the doubled uniqueness already
proved identifies it with the prescribed one. Take the local Haar
measures to give \(K_v\) mass one at almost all places. Write
\(b_v(g)=\langle\omega_v(g,1)\phi_v,\phi_v\rangle\) and
\(B_v=\int|b_v|\). Outside a finite set \(S\) one has \(b_v|_{K_v}=1\)
and \(B_v\ge1\). The integral of \(\prod_v|b_v|\) is finite by
the absolute unfolding, and is positive on some cylinder about
the identity. Exhaustion by finite restricted products and Tonelli
give \(\prod_{v\notin S} B_v<\infty\), hence
\(\sum_{v\notin S}(B_v-1)<\infty\). Thus
\[
        \left|\int_{H(F_v)} b_v(g)\,dg-1\right|
              \le B_v-1
\]
at these places. Only finitely many local integrals can vanish;
replace those and any exceptional finite factors by the compact
open tests just constructed. Absolute factorization now gives a
nonzero product. Formula~\eqref{threal:norm} yields a nonzero
undoubled primed period. The finitely modified tensor is a valid
supported test since the prescribed vacuum is standard almost
everywhere.
\end{proof}

\begin{theorem}[The scalar at the definite base]\label{threal:base}
Suppose \(p=q=0\). Both undoubled periods are nonzero, the primed
one on the required supported subspace. Common matching quotient
maps exist. For every choice of such maps there is a unique
\(c\in\mathbb C^\times\) such that \(I_\Phi=cI'_{\Phi'}\) for
all matched tests. Both parities are allowed. The inert and ordinary
split maps are furnished by Propositions~\ref{threal:allinert} and
\ref{threal:ordinarysplit}; the specified orbital maps at \(S\)
and the stronger kernel and support conditions at the totally split
good primes are those of Corollary~\ref{split:ram} and
Proposition~\ref{split:good} in the row polarizations. The
identification of those polarizations and their vacua is the one
in Proposition~\ref{threal:ordinarysplit}; at \(S\) use the
specified lattice bases. Proposition~\ref{threal:compact} gives
the archimedean maps.
\end{theorem}
\begin{proof}
First establish unprimed nonvanishing. Write \(n=a\),
\(H=U(V)\), \(J=U(W)\), \(\chi=\beta/\alpha=\nu^{-1}\),
\(T=S=S_0\sqcup Q\) the exact set above which \(\nu\) is ramified,
and \(s_0=(d-n)/2=a'/2\). Both adelic quotients are compact.
Form the Siegel section of \(U(V\oplus -V)\) through the
doubled Schwartz vector \(\phi\otimes\bar\phi\) at \(s_0\),
and extend it flat in normalized induction using the compact
subgroups specified below. Pair the resulting Eisenstein series
on the two copies of \([H]\) against the constants, cancelling the
second-block determinant character as above. Write this scalar
as \(Z(s)\). The graph orbit and the change of variables in
Theorem~\ref{threal:nonzero} express \(Z(s)\), for \(\Re s\)
large, as a positive measure constant times the product of the
local integrals of \(f_{v,s}(g,1)\) over \(H(F_v)\). Constants
are fixed by giving the hyperspecials volume one outside \(T\)
and choosing product Haar measure up to one global positive
factor. The meromorphically continued Eisenstein series can be
paired over the compact quotient: its meromorphic continuation
as smooth automorphic forms is locally uniform on compact sets
after multiplication by a local scalar denominator. At \(s_0\),
the anisotropic Siegel--Weil formula with integrating space \(W\)
and \(d>n\) gives \(Z(s_0)=C\|I_\phi\|^2\), \(C>0\), by Fubini.
No initial absolute convergence at \(s_0\) is used.
\cite[Theorem~7.1 and \S11.3]{ganQiuTakeda}

At each infinite place take a nonzero source-invariant polynomial
from Proposition~\ref{threal:compact}. Choose the maximal compact
in the doubled group containing \(H(F_v)^2\); hence
\(f_{v,s}(g,1)=\|\phi_v\|^2\) there. The doubled section is
compact-finite, since the doubled maximal compact at this place
is the product of the two compact groups. At \(v\in T\), which is
split because \(T\subset P\), take the small-ball test about a
full row frame in \(M_{n,d}(F_v)\) constructed above. On
\((g,1)\) the flat section equals the compactly supported matrix
coefficient at \(s_0\) times a positive Iwasawa height raised to
\(s-s_0\). Its integral is entire and strictly positive at \(s_0\).

At every other finite place the doubled section can be taken to
be the basic spherical section of value one at the identity.
At split places use the characteristic function of integral
matrices in the linear model. On the doubled row space the two
blocks together give the integral matrices with \(2n\) rows.
The determinant character is unramified and the integral general
linear group fixes their indicator. At inert places the extension
and characters are unramified, the additive character has conductor
\(\mathcal O_v\), and the residue characteristic is odd. The
lattices are perfect, \(W\) has even dimension and is integrally
hyperbolic, as proved above. In its polarization the doubled row
group acts linearly on \((V\oplus -V)^{d/2}\) with unramified
splitting, and the product of the two lattice indicators is
fixed by the stabilizer of the doubled self-dual lattice. This
works also for \(n\) odd. Indeed an integral hyperbolic plane
splits off repeatedly from a perfect Hermitian lattice of rank at
least two over an unramified quadratic extension of odd residue
characteristic by reduction and Hensel lifting, leaving a unit
line in odd rank. Thus the original row group too is unramified.
Normalize the indicators to have \(L^2\)-norm one. Evaluation on
the doubled diagonal then has value one. Transforming back to the
given presentation uses only finitely many departures from its
prescribed vacuum, exactly as above.

Let \(\eta=\eta_{E/F}\), and use a superscript \(T\) to omit finite
Euler factors over \(T\). The unramified doubling computation
with trivial representation of \(U(n)\) gives for the product of
the factors just described outside \(T\) and infinity
\[
 \frac{\displaystyle\prod_{i=0}^{n-1}
                 L^{T}_{E,\mathrm{fin}}(s+n/2-i,\chi)}
      {\displaystyle\prod_{i=0}^{n-1}
                 L^{T}_{F,\mathrm{fin}}(2s+n-i,\eta^{\,i})}.
\]
Here the base change of the spherical trivial representation has
the \(n\) exponents \((n-1)/2-i\): it is the spherical
constituent with trivial representation's Satake parameter,
the half sum of positive roots, restricted to the split group
over \(E\). In the inert odd case this list includes zero. The
doubling \(L\)-function uses \(s+1/2\), and \(\chi|_{\mathbb A_F^\times}=1\)
because the two ranks have equal parity. These give exactly the
shifts and the denominator above. See the unramified formula
\cite[\S4.2.1, (53)]{ehlsDoubling}. The product equals the
meromorphic continuation from its region of absolute convergence.
In particular no place with an unramified numerator has been
removed arbitrarily: at every prime of \(E\) over \(T\) the
character is ramified by the hypotheses, so its local factor is one.

At \(s_0\) the denominator arguments are at least \(a'+1>1\).
The numerator arguments are integers \(k=d/2-i\).
At each infinite place the angular exponent \(m\) of \(\chi\)
has \(|m|>d\); the gamma factors
\(\Gamma_{\mathbb C}(k+|m|/2)\) are finite and nonzero throughout
this interval. The Euler product gives nonvanishing when \(k>1\).
At \(1\), one elementary proof uses that both \(\chi\) and
\(\chi^2\) have nonzero infinity type. They have no Hecke pole.
For \(u>1\),
\(\zeta_E(u)^3|L_{E,\mathrm{fin}}(u,\chi)|^4
 |L_{E,\mathrm{fin}}(u,\chi^2)|\ge1\) off the ramified factors,
by \(3+4\cos t+\cos(2t)\ge0\) in the logarithmic Euler products.
The finite omitted factors of the three functions are bounded
above and below near \(1\). A zero would have order at least four
in this product, while \(\zeta_E^3\) has pole order three. For
\(k\le0\), apply the Hecke functional equation to the completed
function and use nonvanishing at \(1-k\) for \(\chi^{-1}\);
the conductor power and root number are nonzero
\cite{tateThesis}. This proves
\(Z(s_0)\ne0\). The primed supported period is already nonzero
by Theorem~\ref{threal:nonzero}.

We next compare the actual automorphic copies. On the definite
source the primed integral exists also for tests without the
support restriction. Use that extension in the following argument.
All three quotients are compact, with compact groups at infinity.
For any specified finite open compact level and infinite compact
type the space of automorphic vectors is finite-dimensional:
there are finitely many finite-adelic double cosets at that
level, and matrix coefficients of the prescribed compact type
span a finite-dimensional space on each. Consequently the unitary
spectrum is discrete with finite multiplicities. More explicitly
the orthogonal projection to any such level and type has finite
rank; these projections exhaust a dense subspace, and a continuous
direct integral of inequivalent unitary representations would have
infinite rank on any projection on which it is nonzero.

Let \(P_v,P'_v\) be the given local quotient maps and \(\pi\)
the restricted tensor product of the common irreducibles, using
their normalized vacuum images. At infinity take the entire
product of the compact factors; the polynomial calculation of
Proposition~\ref{threal:compact} gives its source-invariant
quotient with multiplicity one. For each of the two periods,
project to any irreducible automorphic summand of the target.
A nonzero projection, with all but one test factor fixed and
the other target factors paired with linear functionals, gives
an intertwiner from the local oscillator to the source trivial
representation times the local target constituent. At finite
places local Howe duality gives a unique such irreducible quotient
and a one-dimensional space of maps to it. This applies to
unitary dual pairs in characteristic zero and to the split
general linear dual pair; twisting an oscillator splitting by
characters preserves the assertion
\cite{ganTakedaHowe,minguezHowe}. At infinity the identical
conclusion follows from the compact calculation. Hence every
nonzero projection takes values in a copy of \(\pi\) and
annihilates each local kernel. Pairing is legitimate on algebraic
tensors: a smooth finite-level, compact-finite vector of an
irreducible automorphic summand belongs to its algebraic restricted
tensor product, and the pairings on the remaining factors separate
these vectors. At each finite tensor stage the kernel of the tensor
quotient is the sum of the subspaces having a kernel in one
factor, by exactness of tensor products over \(\mathbb C\).
This proves factorization of both periods. Smoothness then gives
pointwise equality whenever two resulting automorphic vectors are
equal in \(L^2\).

In fact \(\pi\) has multiplicity one here. Let \(i_0\) be the
unitary embedding defined by the nonzero primed image, using on
\(\pi\) the norm transported from that image. If another copy
exists, choose a unitary embedding \(i_1\) with image orthogonal
to \(i_0\). For any polynomial primed test \(u\) and compact-finite
smooth \(x\in\pi\), form the reverse lift
\[
  T_u(i_jx)(g)=\int_{[J]}(i_jx)(h)\,
                \overline{\Theta'_u(g,h)}\,dh
         \qquad (g\in U(V')(\mathbb A_F)).
\]
This function is constant in \(g\). In fact pairing against any
irreducible summand of the compact-quotient source spectrum and
varying \(x,u\) gives the same local intertwiners as above,
now with target constituent \(\pi_v\) fixed. Injectivity of
local Howe correspondence forces the source constituent to
be trivial at every place. At infinity the equivariant Hom in
the reverse direction also has this property: the Cauchy
decomposition in the compact calculation pairs distinct partitions
on the two compact groups. The globally trivial representation
occurs precisely on constant functions. Therefore \(T_u(i_1x)=0\):
its average on \([U(V')]\) is a positive normalization times
\(\langle i_1x,I'_u\rangle=0\).

On the other hand
\[
                 \|T_u(i_0x)\|_2^2=\|T_u(i_1x)\|_2^2.
\]
To see this form the flat doubled section \(f_{u,s}\) on
\(U(W\oplus -W)\) from \(u\otimes\bar u\), with parameter
\(s_*=(a'-d)/2=-a/2\). Pair its Eisenstein series, after the
second-block determinant correction, over \([J]^2\) against the
two copies of the vector \(i_jx\) and its conjugate. For \(\Re s\)
large there is just the graph orbit, since \(W\) is anisotropic.
Unfolding gives the integral over \(J(\mathbb A_F)\) of the
restricted section against the matrix coefficient of the right
regular action on \(i_jx\). In the convention with the section at
\((1,g)\) its scalar weight is \(f_{u,s}(1,g)\chi_*(\det g)^{-1}\),
where \(\chi_*=\alpha\) is the primed target splitting character.
The matrix coefficient is \(\langle R(g)i_jx,i_jx\rangle\), or
its conjugate upon conjugating the whole pairing. Either is
independent of \(j\) by isometry and equivariance. The two
meromorphic integrals thus agree identically. At \(s_*\) the
anisotropic Siegel--Weil formula is holomorphic, and its coefficient
is two for the probability theta integral because \(a'<d\).
The doubled product identity and Fubini over the compact quotients
give the two asserted squared norms times the same positive constant
\cite[Theorem~7.1 and \S11.3]{ganQiuTakeda}. Choose \(u\) with
\(I'_u\ne0\) and \(i_0x=I'_u\); then the average of
\(T_u(i_0x)\) is a positive multiple of \(\|I'_u\|^2\).
This contradicts \(T_u(i_1x)=0\). The two nonzero maps from the
common quotient now lie in the same one-dimensional intertwining
space by Schur's lemma. Restriction back to supported primed tests
and Corollary~\ref{threal:limits} give the required nonzero scalar
and its uniqueness.
\end{proof}

\begin{proposition}[The two periods against the cuspidal spectrum]\label{threal:orthA}
The unprimed period pairs to zero with every irreducible cuspidal
representation of \(J=U(W)\) whenever \(p>0\), for arbitrary \(q\).
If \(q=0<p\), the primed period also has this property, and on
removing a target hyperbolic plane the zero-vector orbit of the
unprimed line constant term vanishes identically.
\end{proposition}
\begin{proof}
Choose a rational prime completely split in \(E\) outside the finitely
many exceptional primes and the prime under \(v_1\), using the
splitting-prime theorem; let \(v\)
be a place above it. Write \(k=F_v\). By Lemma~\ref{split:ray}, the
unramified characters \(\alpha,\beta\) there have distinct values
on a uniformizer. The source-invariant quotient of the unprimed local
oscillator is the entire irreducible normalized induction
\[
       \pi_v=(\alpha\circ\det_a)\times(\beta\circ\det_{a'})
\]
by Proposition~\ref{split:good}. Thus every cuspidal constituent
pairing nontrivially with the unprimed period must have this local
factor. In detail the pairing, with other Schwartz coordinates and
other cuspidal coordinates fixed, gives an equivariant map into
the smooth dual of the conjugate of the local cusp representation.
Smoothness follows from invariance of the test under an open
compact subgroup, and the smooth dual in question is the original
representation by unitarity. Source invariance follows by changing
variables in the source integral.

For \(p>0\) the source has a rational hyperbolic plane: its form is
isotropic at every real place, it is split at the split finite places,
and at the inert places its perfect lattice has Witt index \(\lfloor
a/2\rfloor\). For the underlying \(F\)-quadratic norm form local
isotropy therefore gives global isotropy by Hasse--Minkowski, and
one completes the isotropic vector to a hyperbolic pair.
Write \(\Pi\) for such a cuspidal representation. Its adjoint lift
back to \(U(V)\), with conjugated oscillator, has a vector of
nonzero average, by Fubini. Indeed the absolute source theta sum in
Proposition~\ref{threal:decay}, integrated in \(g\), is bounded by
a power of \(H(h)\): in its unprimed estimate retain the factors
\(O(M^D H(h)^{A_B})\) before integrating the strictly negative
torus exponents. A cuspidal automorphic vector of finite compact
type decays faster than any power on a target Siegel set.
These estimates justify both the pairing and the interchange.

At any earlier source in the Witt tower, of dimension \(a-2b\)
with \(b\ge1\), no irreducible local representation can lift to
\(\pi_v\) in the forward oscillator. Indeed M\'inguez's split
correspondence adds the characters
\[
   \beta|\cdot|^{(1-l)/2},\
   \beta|\cdot|^{(3-l)/2},\ldots,\beta|\cdot|^{(l-1)/2},
               \qquad l=d-(a-2b)=a'+2b ,
\]
to the cuspidal support of the transformed source parameter.
The bare matrix oscillator adds the segment of the trivial
\(\mathrm{GL}_l\) representation. The scalar in the target row
action in \eqref{newtheta:big} is \(\beta(\det h)\) after removing
the unitary change-of-variables multiplier. The inverse transpose
on matrices interchanges the segment with its reversed dual, which
has the same multiset of real powers. This gives exactly the
displayed twist. The actual support of \(\pi_v\) has only \(a'\)
entries with unitary part \(\beta\). The other \(a\) entries have
unitary part \(\alpha\). This excludes the displayed segment.
\cite[Corollary~6.3(2)]{minguezHowe}
The first nonzero global adjoint lift of a cusp representation in a
Witt tower is cuspidal, by the tower constant-term formula
\cite[Proposition~3.1]{wuThetaTower}. If a lower lift were nonzero,
the first such lift would therefore have an irreducible cuspidal
constituent. Pairing against it and conjugating gives the excluded
forward local intertwiner. Thus the lift on \(U(V)\) is itself
cuspidal. The integral of a cusp representation against the constant
function on this positive-rank group is zero: a nonzero invariant
functional on a constituent would make it the trivial representation,
whose constant term along the given rational parabolic is nonzero.
This contradicts the preceding nonzero average.

For the primed period suppose \(q=0<p\) and let \(Y=V'\). A target
plane exists by the same norm-form argument, after multiplying its
skew form by a nonzero trace-zero element of \(E\).
The group \(U(Y)\) is anisotropic.
If the pairing with a cusp representation \(\Pi\) is nonzero, its
adjoint lift on \(U(Y)\) has nonzero projection to constants.
The entire adjoint lift is cuspidal by anisotropy, and is irreducible
by \cite[Theorem~5.3]{wuThetaTower}. Consequently it is the constant
representation. Its occurrence as functions consists of actual
constants, since right-invariant functions on \([U(Y)]\) are
constant. Wu's involution theorem now says that the forward lift of
these constants to \(W_1\) vanishes, where \(W_1\) is obtained from
\(W\) by deleting the rational hyperbolic plane
\cite[Theorem~5.1(3), Remark~5.2]{wuThetaTower}.
To check the twists in this use of the theorem, let \(d_X,d_Y\)
be its two splitting characters restricted to norm-one determinants
on the two groups \(X=W,Y\), respectively. Complex conjugation
of the oscillator replaces \(\psi,\chi_X,\chi_Y\) by their
inverses. Relative to Wu's oscillator for \(\psi^{-1}\) and the
original splitting characters it thus twists by
\(d_X^{-1}\boxtimes d_Y^{-1}\): on the doubled Siegel Levi this
follows from \(\chi(z/\bar z)=\chi(z)^2\); the doubled uniqueness
argument identifies the restrictions to the two factors. Wu's
reverse input is therefore \(d_X^{-1}\Pi\) and its output is
\(d_Y\) times the constants. The factor \(d_Y^{-1}\) in his
Theorem~5.1 cancels it in the forward lift.

This lower forward lift is nonzero by the calculation in
Theorem~\ref{threal:nonzero}. Here \(Y\) is unchanged, and the
new target has dimension \(d-2>2a'\), signature
\((d-p-1,p-1)\) with \(d-p-1\ge a'\), and is still hyperbolic
at the inert finite places, by Witt cancellation. The convergence
inequality there is \((d-2)-2(p-1)+1=d-2p+1>a'\).
These give respectively the absolute doubling identity, its real
source-invariant polynomial, and the required finite tests in that
proof. The splitting characters have unchanged parity on deleting
a plane. This contradiction proves the assertion for the primed
period as well.

Finally the zero-vector term on the unprimed side uses the
oscillator for the same \(V\) and the smaller \(W_1\). In the
polarization with position variable \(x\in V\) for the removed
target plane, setting \(x=0\) leaves the theta sum for
\(V\otimes W_1\); the extra source determinant character is
trivial for this added hyperbolic plane, whose source action on
the \(x\)-block is ordinary pullback. The source integral converges
by the proof of Proposition~\ref{threal:decay} with \(d-2\).
At a real place, the positive-character joint-harmonic condition in
Proposition~\ref{threal:allreal} now requires \(p\le p-1\).
Hence this local oscillator has zero trivial-source coinvariant
module: if its big theta module were nonzero, its unique irreducible
quotient would occur in the local correspondence. The source
integral at each \(h\) kills all the Lie algebra and compact-group
differences and so vanishes on the polynomial module. Continuity
and Hermite approximation give the same conclusion for all Schwartz
inputs. For the primed source the only rational isotropic vector is
zero, since \(Y\) is definite. Thus it has no nonzero-vector
orbit in this constant term.
\end{proof}

\begin{proposition}[Unfolding the larger nonzero orbit]\label{threal:coneA}
Let \(p>0\). Choose rational hyperbolic pairs \(x,y\) in \(V\)
and \(x',y'\) in \(W\), with orthogonal complements \(V_1,W_1\).
Write \(m(t,h_1)\) for the target line Levi and \(k=(a-a')/2>1\).
Use rational mixed Schr\"odinger coordinates and probability
measure on the target unipotent quotient. Let \(e_A\) be evaluation
in the coordinates dual to the modulation space
\[
        S_A= V\otimes_E Ex'\ \oplus\ Ex\otimes_E W_1.
\]
With \(K\) the source compact stabilizer and \(\rho\) the Heisenberg
action, put
\[
 T_A\Phi=\int_K\int_{\mathbb A_E^\times}
  e_A\rho(l\,x\otimes y')\omega(k_0,1)\Phi\
             \nu(l)|l|_E^k\,d^\times l\,dk_0 .
\]
The entire contribution of the rational nonzero isotropic vectors
to the unnormalized line constant term, at \(m(t,h_1)h\), is
\[
   c_A\,\alpha(t)|t|_E^{a'/2}\,
      I_{V_1,W_1}\big(T_A(\omega(1,h)\Phi)\big)(h_1).
\]
Here \(c_A\ne0\) depends on the rational models, the evaluations
and Haar measures; it is independent of the splitting characters.
At \(q=0<p\) this is the entire unprimed line constant term.
The entire primed constant term there is
\(\alpha(t)|t|_E^{a'/2}\) times its smaller zero-evaluation
period, with coefficient one for coherent mixed models and the same
source Haar measure.
\end{proposition}
\begin{proof}
First use \(V\otimes Ex'\) as modulation space. In its dual
coordinate the sum runs over \(z\in V(E)\). For fixed source and
target translates the sums are absolutely uniform on the compact
unipotent quotient. The center acts by an additive character of a
rational nonzero multiple of \(\langle z,z\rangle_V\); averaging
it leaves exactly the null vectors. The remaining target unipotent
coordinates have quotient \(W_1(\mathbb A_F)/W_1(F)\), writing
\(W_1\) here for its underlying \(F\)-vector group. They act on
the residual oscillator by the Heisenberg translations or
modulations \(z\otimes w\). This follows by applying the unipotent
matrix to the target coordinates \(x',W_1,y'\): the \(W_1\) shift
of \(y'\) is \(w\), and the quadratic correction in the \(x'\)
coordinate multiplies \(\langle z,z\rangle_V\), hence gives no
phase on this cone.

All nonzero rational null vectors form one \(U(V)(F)\)-orbit by
Witt extension. At the representative \(z=x\), use \(Ex\otimes W_1\)
as a further modulation space in the residual oscillator.
The last compact additive average kills every nonzero dual rational
coordinate by character orthogonality. Its surviving value is
therefore
\[
    \sum_{u\ {\rm rational}}
      (e_A\rho(x\otimes y')\,\omega(g,h)\Phi)(u),
\]
where \(u\) belongs to a position space for \(V_1\otimes W_1\).
The changes of rational polarization here preserve the theta
summation by Poisson summation with self-dual measures. Dual hyperbolic
bases may introduce a nonzero rational \(E\)-multiple of \(x\)
in this expression. Replacing the representative by that multiple
does not change the orbit, and its idele norm and Hecke character
are one.

Unfold the orbit to \(H_x(F)\backslash H(\mathbb A_F)\),
\(H=U(V)\), where the stabilizer of the vector is
\(H_x=N_x\rtimes U(V_1)\); the line stabilizer has the additional
Levi factor \(E^\times\). Write \(g=n\,m(b,g_1)k_0\) in source
Iwasawa coordinates. The adelic compact in this decomposition can
be taken to be the prescribed \(K\). At split finite places this
is integral row reduction. At an inert finite place one scales the
line generator to a primitive lattice vector and pairs it to one
by self-duality. The second vector can be made isotropic by the
surjectivity of the integral trace in the unramified extension.
Their unimodular hyperbolic plane splits off. Any two unitary
lattice complements here are isometric: a nonsingular finite-field
Hermitian space has a vector of nonzero norm, so integral
orthogonalization gives a unit diagonal, and every base-field unit
is an integral norm. The skew self-dual case is reduced to the
Hermitian case by the local different generator. These facts give
transitivity on isotropic lines. At infinity the two compact unitary
factors carry any isotropic line to the chosen line, by moving its
two components of equal length. The resulting local decompositions
are restricted products.

On the evaluated residual oscillator \(N_x\) acts trivially:
it fixes \(x\), acts with determinant one on \(S_A\), and acts
trivially on \(S_A^\perp/S_A=V_1\otimes W_1\). In mixed coordinates
the off-diagonal entries produce Heisenberg operators with parameters
linear in the dual \(S_A\)-coordinate, and quadratic phases in that
coordinate; both disappear at zero. It also fixes the displacement
\(x\otimes y'\). The source line Levi contributes
\(\nu(b)^{-1}|b|_E^{(d-2)/2}\) and the residual action of \(g_1\):
split off the whole target plane first, on which the source acts
by pullback on \(V\); on the remaining oscillator the source
Siegel character is \(\beta/\alpha\) and its modulation block
\(Ex\otimes W_1\) has dimension \(d-2\). Commuting the Levi left
past the displacement replaces \(x\otimes y'\) by
\(b^{-1}x\otimes y'\). The source Haar density on the line Levi
contains \(|b|_E^{1-a}\), because \(N_x\) has \(a-2\) ordinary
\(E\)-coordinates and one central \(F\)-coordinate of weight
\(b\bar b\). Integration over \([N_x]\) contributes its volume.
After the substitution \(l=b^{-1}\) the total character and modulus
are \(\nu(l)|l|_E^k\), since \(a-1-(d-2)/2=k\).
Integrating over \([U(V_1)]\) gives the smaller period.
The quotient integration constant and the unipotent volume are
positive for positive compatible Haar measures.

For justification of this unfolding one can integrate the
\emph{absolute value of the single orbit summand} over
\(H_x(F)\backslash H(\mathbb A_F)\) first. In the coordinates just
used its integral is bounded by the \(K,\mathbb A_E^\times\)
integral with weight \(|l|_E^k\) of the absolute smaller theta
sum integrated over \([U(V_1)]\). Choose the rational
source-adapted polarization of Proposition~\ref{threal:decay} for
this last sum. Its absolute integral is bounded by finitely many
Schwartz seminorms for each fixed finite test. The gap in dimension
is now \(k-1\) and the rational source Witt index is at most \(p-1\),
so the same negative-exponent estimates apply. In mixed coordinates
the new test is obtained by restriction to the linear coordinates
\((l,u)\); the partial Fourier transforms for changing the smaller
polarization act on \(u\) alone. At infinity these seminorms are
bounded by arbitrarily rapidly decreasing functions of \(l\),
uniformly on the compact \(K\)-orbit. At finitely many modified
finite places the restrictions are finite sums of products of
compactly supported tests in \(l,u\). At every other place they
are the indicator of the integral \(l\) times the smaller lattice
test. The absolute product is dominated by the Dedekind zeta
Euler product at \(k>1\); the finitely many local integrals at zero
also converge since \(k>0\). Tonelli applied to the translates of
the absolute orbit summand now justifies its global unfolding.
The same bounds prove that \(T_A\Phi\) is Schwartz and justify
interchanging its evaluation with the smaller period.
With multiplicative unit volume one at the basic places, the actual
scalar there is the absolutely convergent product of
\((1-\nu_w(\pi_w)(Nw)^{-k})^{-1}\), by the integral of the
lattice indicator along the displacement line.

On the target line Levi the evaluation scalar is
\(\beta(t)|t|_E^{a/2}\) because the whole \(V\otimes Ex'\)
block was removed first. Moving the Levi to the left past the
displacement changes \(l\) to \(l\bar t\); hence substitution in
the Mellin integral multiplies this scalar by
\(\nu(\bar t)^{-1}|t|_E^{-k}=\nu(t)|t|_E^{-k}\).
The target character is thus \(\alpha(t)|t|_E^{a'/2}\).
An evaluation with the same specified Heisenberg laws changes the
formula by a fixed scalar. Locally, invariance under the modulations
factors evaluation through restriction to zero: at a finite place
average on balls; at infinity differentiate the modulation law
to annihilate the coordinate multiples, which generate Schwartz
functions vanishing on that subspace (use Taylor's formula near it
and a cutoff away from it). The remaining intertwiner commutes with
all smaller translations and modulations, so is scalar. The
prescribed vacuum images fix the local normalizations at almost
every place. Thus changing the evaluations or the product measures
leaves the asserted independence from the splitting characters.
Finally Proposition~\ref{threal:orthA} removes the other unprimed
orbit when \(q=0\). On the primed side definiteness leaves only
zero, and the same first mixed evaluation gives the stated formula
with exponent \(a'/2\) and character \(\alpha\).
\end{proof}

\begin{proposition}[Detection by the entire line constant term]\label{threal:detectline}
Let \(f\) be a linear combination of the two periods of polynomial
tests, with the required support on the primed tests. In all
signatures, vanishing of the constant term of \(f\) along one
rational isotropic line, as a function on \(J(\mathbb A_F)\),
implies that \(f\) is cuspidal. For \(p+q>0\) this implies \(f=0\).
\end{proposition}
\begin{proof}
Witt extension conjugates any two rational isotropic subspaces of the
same dimension. If \(f\) is not cuspidal, choose the smallest \(r\)
for which its constant term along the stabilizer \(P_r\) of a rational
isotropic \(r\)-plane is nonzero. These are all the conjugacy classes
of maximal proper rational parabolics. By assumption \(r\ge2\).
On the Levi of \(P_r\), fix the other factor and vary the first
factor \(\mathrm{GL}_r(\mathbb A_E)\). Some resulting function,
allowing a compact right translate of \(f\) by Iwasawa decomposition,
is nonzero. It is cuspidal on this general linear group: the further
constant term along a \(j\)-plane in this block is the constant term
along the \((j,r)\)-flag, computed by unipotent Fubini, and this
also factors through the vanishing \(j\)-plane term for \(0<j<r\).

Here is the spectrum to which this cuspidal function belongs.
The original polynomial periods are compact-finite at infinity
and finite under the center of the target enveloping algebra:
source invariance factors the polynomial oscillator at each place
through its big theta module of the trivial representation, which
is of finite length by real Howe duality \cite{howeTranscending}.
Moderate growth follows from Proposition~\ref{threal:decay}.
The constant term restricted to its Levi is again center-finite.
Indeed on left unipotent-invariant functions restriction of a
central right differential operator gives its relative
Harish-Chandra projection (up to the modulus twist for the
unnormalized constant term). The Levi center is finite over this
image, since the invariants of the Levi Weyl group in the Cartan
polynomials are finite over the full Weyl invariants. This also
gives center-finiteness on the first direct factor of the Levi.

One can thus take a nonzero eigenfunction for the positive real
split center in the finite-dimensional span of its translates
(equivalently, extract an exponential coefficient on its Lie
algebra). These operations commute with finite-place translations.
Project also to a character on the compact norm-one idele class
center, and twist by a common real power of the idele determinant
to make the resulting central character unitary. Some such projection
is nonzero by compact Fourier decomposition. The resulting smooth,
compact-finite, center-finite cusp form is square-integrable modulo
the center by rapid decay on Siegel sets. It has a nonzero projection
to an irreducible unitary cuspidal representation of
\(\mathrm{GL}_r(\mathbb A_E)\). Its local factors are generic by
the cuspidal Whittaker expansion \cite{shalikaGeneric}. Thus at
each place above the \(v\) used in Proposition~\ref{threal:orthA}
the constant-term space must admit a generic irreducible subquotient
unitarizable after a determinant twist. Equivalently one can carry
out the center projections on the finite-place submodule generated
by the original slice: its split-center translates span a finite
number of copies of that module, and restriction to a central
eigenspace preserves its possible finite local subquotients.

At \(v\) each period factors through \(\pi_v\) by the full split
kernel computations, whether or not its coefficients at the other
places factor through their small theta quotients. The local
\(P_r\) has block sizes \((r,d-2r,r)\). Evaluation on the Levi after
taking the constant term factors through its unipotent coinvariants,
since a right unipotent translation becomes a translation of the
compact integration quotient. Normalizing the Jacquet functor only
twists each general linear block by a common real determinant power.
The geometric lemma computes the first \(\mathrm{GL}_r(k)\) action
in the resulting filtration, with possible subquotients of
\[
 \left(\alpha|\cdot|^{-(a-i)/2}\right)\circ\det_i
       \ \times\
 \left(\beta|\cdot|^{-(a'-j)/2}\right)\circ\det_j,\qquad i+j=r .
\]
Indeed a shuffle selects an initial block of length \(i\) of the
first determinant representation and of length \(j\) of the
second. The inverse square root of the modulus of their internal
parabolics gives the displayed exponents. The other end of the
isotropic plane accounts for the last, separate block of size \(r\).
\cite[geometric lemma, \S2.12]{bzgLemma}
The first block contributes a nondegenerate Whittaker functional
only if \(i,j\le1\): the highest derivative of the induced
representation is the tensor product of the highest derivatives
of its factors, and a determinant character of dimension \(>1\)
vanishes under that functor. The Whittaker coinvariant functor is
exact, so the same exclusion holds for every subquotient
\cite[\S4]{bzgLemma}.

It remains to examine \(r=2,\ i=j=1\). This normalized principal
series of \(\mathrm{GL}_2(k)\) is irreducible since its character
ratio has unitary part \(\alpha/\beta\) of infinite order and
therefore is neither \(|\cdot|\) nor \(|\cdot|^{-1}\).
The two Satake eigenvalues have distinct arguments and their
absolute values differ by the factor
\(|\pi_k|^{-(a-a')/2}\ne1\), where \(\pi_k\) is a uniformizer.
For unitarity the multiset must be invariant under
\(z\mapsto1/\bar z\), by the isomorphism to the conjugate dual.
That involution preserves the argument of each eigenvalue.
Consequently their distinct arguments force both absolute values
to be one. No common change of modulus can achieve this, giving the
required contradiction. Thus \(f\) is cuspidal. For \(p+q>0\)
Propositions~\ref{threal:orthA}, \ref{tpd:zero} and \ref{tpd:orth}
cover the two periods in every signature. Thus \(f\) is orthogonal
to the cuspidal spectrum and hence to itself; continuity gives
pointwise zero.
\end{proof}

\begin{corollary}[The detector in the \(a\)-descent]\label{threal:endA}
For \(q=0<p\), suppose that for one fixed \(c\) the
\(\alpha|\cdot|_E^{a'/2}\) components of the two periods agree up to
\(c\) on \(m(1,h_1)\) for every \(h_1\) and every polynomially
matched pair with supported primed test. Then those periods agree
up to \(c\) on all of \(J(\mathbb A_F)\) for every such pair.
\end{corollary}
\begin{proof}
By Proposition~\ref{threal:coneA} each entire line constant term
is just its indicated component. The matching condition survives
right translation by the target compact group: its real compact
action preserves the polynomial module, and its finite action
preserves full rank at \(v_1\) and intertwines the matching maps.
The hypothesis therefore applies also to \(m(1,h_1)k_0\) for each
target compact \(k_0\). The common Levi character extends this to
\(m(t,h_1)k_0\); left unipotent invariance and Iwasawa decomposition
cover the whole adelic group. Proposition~\ref{threal:detectline}
now applies to the difference.
\end{proof}

\begin{proposition}[The detector in the \(a'\)-descent]\label{threal:endB}
Let \(q>0\). Suppose a difference \(f\) of the two polynomial
periods, with supported primed test, has its full line constant
term transforming by \(\alpha(t)|t|_E^{a'/2}\) alone. Then \(f=0\).
In particular equality of the \(\beta(t)|t|_E^{a/2}\) components
of two matched periods, with one scalar for all the tests on
\(m(1,h_1)\), implies equality with that scalar on \(J(\mathbb A_F)\).
\end{proposition}
\begin{proof}
Induct on \(p\), keeping \(q>0\). For \(p=0\) the constant terms
of both periods have only the \(\beta\)-character by
Propositions~\ref{tpd:zero} and~\ref{tpd:unfold}. The two line
characters have distinct moduli since \(a>a'\), so the constant
term vanishes and Proposition~\ref{threal:detectline} applies.

For \(p>0\), split off planes in \(V,W\). At the real places one
can take compact groups adapted to the chosen planes. In detail,
choose on each hyperbolic plane a Cartan involution exchanging its
null lines and choose one on its orthogonal complement. The
resulting positive structure is conjugate to the original by a
real isometry. Translating the test by this isometry transports its
polynomial Fock module to the new one. Source translation leaves
the period fixed and simultaneous target translation replaces it
by a right translate, preserving purity and the zero assertion
to be proved. The lattices for the unused smaller finite models
can be chosen split along the rational decomposition at almost all
places, and compatible with the specified split polarizations at
\(v_1\). The induction uses only these rational theta models, their
finite tests, the prescribed signatures and the gap inequalities.
These hold in the smaller ranks \(a-2,a',d-2\), with indices
\(p-1,q\).

For each fixed target compact translate, restriction of this
pure constant term to the smaller unitary group is a difference
\(f_1\) of its two polynomial periods, with supported primed test.
Indeed the \(\alpha\)-component of the unprimed form is precisely
the cone in Proposition~\ref{threal:coneA}; the zero evaluation
has character \(\beta(t)|t|^{a/2}_E\). The smaller unprimed test
is polynomial at infinity by the orthogonal Fock factorization
of Lemma~\ref{tpd:mellin}, now with \(S_A\).
On the primed side the supported zero evaluation of
Lemma~\ref{tpd:so} gives precisely the \(\alpha\)-component.
The adjustment by source differences in that lemma may be made
anew for each fixed translate. Each resulting \(f_1\) in turn
has a two-character line constant term in its own ranks, by the
same calculations (and their definite-source versions). For an
arbitrary upper target translate, source Iwasawa is not needed
again: upper target Iwasawa and the pure character express its
slice as a scalar times a smaller target translate of a
compact-translate slice. Thus the same two line laws hold for all
slices used below.

We prove that the \(\beta\)-component of every such smaller
slice is zero. Choose any rational isotropic line in \(W_1\);
together with the first line it spans a rational isotropic
two-plane \(D\) in \(W\). For any fixed upper right translate
take the \(D\)-parabolic constant term of \(f\), and restrict its
first Levi block to \(\mathrm{SL}_2(\mathbb A_E)\), writing \(F\)
for the resulting automorphic function. Its ordinary upper
triangular constant term on this group is the average for the
isotropic \((1,2)\)-flag, by Fubini over the compact unipotent
quotients. Compute the same flag average in the other order by
first taking the outer line constant term and then the smaller
line term. These are averages along the two filtrations of the
\emph{same flag radical}. On \(d(t)=\operatorname{diag}(t,t^{-1})\)
the outcome, for each fixed right translate \(h\), is
\[
              F_N(d(t)h)=A(h)+\nu(t)|t|_E^{1-k} B(h),
                 \qquad k=(a-a')/2.
\]
The first outer factor is always \(\alpha(t)|t|_E^{a'/2}\).
The inner factors are \(\alpha(t^{-1})|t|_E^{-a'/2}\) and
\(\beta(t^{-1})|t|_E^{-(a-2)/2}\), respectively. These give
the two terms displayed. Their coefficients are continuous and
bounded on compact sets: they can be recovered by evaluating at
two fixed \(t\)'s with different idele norms.

In particular \(F\) belongs to
\(L^2(\mathrm{SL}_2(E)\backslash\mathrm{SL}_2(\mathbb A_E))\).
To check this, the displayed constant term is bounded on a Siegel
set since \(1-k<0\). Its difference from \(F\) decays faster than
any power of the split chamber parameter there. Indeed the
height exponent in Proposition~\ref{threal:decay} can be held fixed
for all right derivatives of the original period: choose once
the seminorm orders needed for integration; differentiating first
changes the test in these seminorms, not the powers of height in
the conjugation bounds. Integrating over the compact \(D\)-radical
therefore gives uniform moderate growth of derivatives of \(F\).
At fixed finite level the upper additive unipotent quotient is a
compact real torus. Subtract the mean of \(n\mapsto F(n d(t)c)\)
and use Fourier series there to bound its value by a Sobolev norm
of its positive-order derivatives of any sufficiently high order.
Along the archimedean positive chamber each unipotent derivative
gains the inverse positive root of \(d(t)\), as a right derivative
on \(F\). Compact factors \(c\) only change the fixed bounds. This
is the estimate in Corollary~\ref{tpd:detect} on the rank-one
group, with the bounded mean retained. Boundedness and finite
quotient volume give the assertion. This applies also to every
fixed finite local translate of the vector giving \(F\).

Use the good split place \(v\) of Proposition~\ref{threal:orthA}.
At one of its two \(E\)-places \(w\) the right
\(\mathrm{SL}_2(E_w)\) action on \(F\) factors through restriction
of the first block in the \(r=2\) Jacquet filtration computed in
Proposition~\ref{threal:detectline}. Its possible factors before
restriction are determinant characters and the one principal series
with \(i=j=1\). The latter has no unitarizable restriction
constituent. Here is a verification. Write \(G=\mathrm{GL}_2(E_w)\)
and \(G^+=Z(G)\mathrm{SL}_2(E_w)\). The quotient \(G/G^+\) is
the finite group of square classes. The irreducible principal
series \(\tau\) is finitely generated over \(G^+\) by coset
representatives applied to a generating vector. Choose a maximal
proper \(G^+\)-submodule; its finitely many \(G\)-conjugates
intersect trivially by \(G\)-irreducibility. The restriction
embeds in a finite sum of simple quotients, so is semisimple, and
its simple factors are \(G\)-conjugate. The center acts by a scalar
character. After twisting \(\tau\) by a real determinant power this
central character is unitary. If any factor on \(\mathrm{SL}_2(E_w)\)
had a positive invariant Hermitian form, so would every conjugate,
and their direct sum would give a \(G^+\)-invariant form.
Averaging it over \(G/G^+\) would unitarize the twist of \(\tau\),
contrary to the Satake calculation in Proposition~\ref{threal:detectline}.

The local subspace under consideration is a quotient of a
finite-length admissible \(\mathrm{SL}_2(E_w)\)-module: in the
Jacquet module take invariants of compact open subgroups fixing
the vector in all the other Levi blocks. Their invariant spaces
in the finite filtration are finite-dimensional on those blocks.
The first block restricts as just described. A finite-length
admissible submodule of a unitary Hilbert representation has
unitarizable simple constituents. In fact the closure of any of
its submodules has exactly that submodule's smooth vectors within
it, by compact-open averaging and finite-dimensionality at each
level; orthogonal projection then splits the submodule. Hence
only the trivial \(\mathrm{SL}_2(E_w)\) factors remain in the
space of translates of \(F\), and its local action is trivial.
In particular \(F_N(d(t))=F_N(1)\) if \(t\) is a uniformizer
at \(w\) and one elsewhere. Its modulus makes
\(\nu(t)|t|_E^{1-k}\ne1\). The flag formula now implies \(B(1)=0\).
The upper right translate was arbitrary, so the entire smaller
\(\beta\)-component vanishes.

The smaller slice \(f_1\) is thus pure for the \(\alpha\)-character.
Induction gives \(f_1=0\) for each upper compact translate.
The character law, upper unipotent invariance and Iwasawa
decomposition give vanishing of the full upper line constant term.
Proposition~\ref{threal:detectline} gives \(f=0\).
Finally equality on \(m(1,h_1)\) for all matched pairs kills
their entire \(\beta\)-component by compact stability of matching
and Iwasawa, as in Corollary~\ref{threal:endA}. The preceding
assertion applies to the remaining pure difference.
\end{proof}

\begin{proposition}[Matching all tests at an integral good split place]
\label{threal:localcone}
Let \(v\) lie above a completely split rational prime outside the
ramification set. Use unramified characters and a pair of integral
hyperbolic splittings in which the source, target and smaller basic
lattices decompose as direct sums and the displacement line is
primitive. Normalize the unit groups and \(K_v\) to volume one.
Write \(P_{1,A},P_{1,B}\) for the two smaller quotient maps,
normalized on their common vacuum. For every pair of locally
matched tests \(\phi,\phi'\), in the \(a\)-step one has
\[
   P_{1,A}(T_{A,v}\phi)=
    \prod_{w\mid v}L_w(\nu,k)\, P_{1,B}(e'_{0,v}\phi').
\]
In the \(a'\)-step the analogous equality is
\[
   P_{1,B}(T_{B,v}\phi')=
    \prod_{w\mid v}L_w(\nu,k+1)\, P_{1,A}(e_{0,v}\phi).
\]
Here the evaluations and basic vectors in the mixed coordinates
are normalized compatibly: evaluation of a basic vector is the
smaller basic vector, and its displacement evaluates to zero
outside the integral line. \(T_{B,v}\) is the local operator in
Lemma~\ref{tpd:mellin}. Both products range over the two places
of \(E\) above \(v\).
At a designated ramified split place the same pairs of comparison
maps span a space of dimension at most one in each direction.
\end{proposition}
\begin{proof}
Set \(K_0=F_v\). By Propositions~\ref{split:kernel} and
\ref{split:good}, or by Corollary~\ref{split:ram} at a designated
ramified place, the two source coinvariants are the same
irreducible \(\mathrm{GL}_d(K_0)\) representation \(\pi_v\).
In particular these are the precise kernels of any of the
surjective local maps in the matching data: a map killing the
source differences induces a nonzero map out of this irreducible
module. We show that the two maps to the smaller output after
applying a cone or zero evaluation factor through these
coinvariants.

For the zero evaluation, equivariance for the unchanged source
and its action on the smaller quotient give this immediately.
For \(A\), project
\(e_A\rho(x\otimes y')\omega(g,1)\phi\) to the smaller quotient.
It is left invariant by the full \(K_0\)-point vector stabilizer
\(H_x\), since its unipotent radical acts trivially as in
Proposition~\ref{threal:coneA}, and its remaining source factor
acts trivially after projection. This stabilizer is unimodular:
the action of the residual general linear block on the row and
column root groups has opposite determinant moduli, and is trivial
on the center. Integration over \(H_x\backslash H(K_0)\) therefore
gives a right \(H(K_0)\)-invariant map. Its Iwasawa coordinates
give precisely \(P_{1,A}T_{A,v}\), up to a nonzero Haar factor,
by the source Levi calculation above. The integral is absolute
since \(k>0\); the locally constant compactly supported restrictions
in the displacement and smaller position variables span a finite
number of smaller tests. For \(B\) use the vector stabilizer of
\(x\) in the primed source. It fixes the \(S_B\)-reduced quotient
and the class of \(y\otimes x'\) modulo \(S_B\), with no central
Heisenberg phase: the difference lies in
\((V'_1+Ex)\otimes Ex'\), isotropic together with
\(y\otimes x'\). The residual \(U(V'_1)\) again disappears after
projection. The Iwasawa calculation in Lemma~\ref{tpd:mellin}
gives \(P_{1,B}T_{B,v}\) by this quotient integral, now with
positive exponent \(k+1\).

Both resulting maps from \(\pi_v\) are invariant under the
unipotent radical of the target line parabolic. For \(A\), the
displacement changes there by an element of
\(Ex\otimes(W_1+Ex')\) in \(S_A\), with no Heisenberg phase
because \(x\) is isotropic. For \(B\), the displacement is fixed;
the target unipotent preserves \(S_B\), has determinant one there,
and acts trivially on the smaller tensor quotient. For the zero
evaluation the assertion holds already before the smaller quotient.
The actions on the reduced spaces in these checks are the mixed
ones at zero, as in Proposition~\ref{threal:coneA}. Moreover the
two maps in either comparison are equivariant for the smaller
target group and have the same line character, by the two Levi
computations.

At this split place the line parabolic has block sizes
\((1,d-2,1)\). Write its outer entries as \(u,z\); the element of
\(E_v^\times\) acting on the isotropic \(E_v\)-line is
\((u,z^{-1})\). The unnormalized line character in the \(a\)-step
is therefore \(\alpha(u)\alpha(z)|u/z|^{a'/2}\). Dividing by
\(\delta_P^{1/2}=|u/z|^{(d-1)/2}\) prescribes the first and last
characters of the normalized Jacquet module to be
\[
       \alpha|\cdot|^{-(a-1)/2},\qquad
       \alpha|\cdot|^{(a-1)/2}.
\]
The geometric lemma applied to the two determinant blocks of
\(\pi_v\) has exactly one such subquotient: the first and last
rows must both come from the \(\alpha\)-block, since
\(\alpha/\beta\) has infinite order on a uniformizer in the good
case, and is nontrivial on units in the ramified case. They leave
the normalized induction of the determinant characters of lengths
\(a-2,a'\) on the middle block. Its \(\alpha\)-exponent has no
shift because one row was removed at either end. It is irreducible
by the unitary induction theorem used in Proposition~\ref{split:good}.
In the \(a'\)-step the line character is \(\beta(t)|t|_E^{a/2}\);
the prescribed ends are \(\beta|\cdot|^{\mp(a'-1)/2}\). Both
come from the other block, leaving lengths \(a,a'-2\). Thus the
equivariant Hom into the stipulated smaller representation and
the two given characters has dimension at most one in each case,
by the Jacquet filtration and Schur's lemma. \cite[\S2.12]{bzgLemma}

At the stated good places evaluate the two maps on the upper
common spherical vector.
The zero evaluation gives the nonzero smaller vacuum. In the
cone integral the compact group fixes the test, and the shifted
evaluation gives exactly the smaller vacuum for integral displacement.
At each of the two places \(w\) the resulting geometric series
is \(L_w(\nu,k)\) or \(L_w(\nu,k+1)\), respectively. The product
is nonzero, fixing the scalar on every matched test.
\end{proof}

\begin{corollary}[The ramified unprimed comparison on all tests]\label{threal:ramA}
At a designated ramified split place, use the local character
\(\chi=\nu_w\), conductor \(\pi^c\), residue cardinality \(Q\)
and measures of Proposition~\ref{tpd:ram}. For all upper matched
tests the scalar in the \(a\)-step, multiplying the smaller primed
zero evaluation, is
\[
       Q^{c(k-1)}\chi(\pi)^c
       \sum_{r\in(\mathcal O/\pi^c)^\times}
                         \chi(r)^{-1}\psi_v(r/\pi^c).
\]
\end{corollary}
\begin{proof}
Use the pivot-determinant tests of Proposition~\ref{tpd:ram},
now removing the first and last of the \(a\) target pivot columns
on the unprimed side. The upper and lower pairs are matched by the
same integral row computation there, and both orbital functions
have value one at the identity before applying the reductions.
The primed zero evaluation integrates one nonpivot column with
\(a'\) entries in \(\pi^c\mathcal O\), giving
\(Q^{-a'c}\phi'_1\). On the unprimed side write
\(l=(s,t)\) in the two split components. The shift \(l x\otimes y'\)
has position entry \(s\) at the first row and last pivot column
and momentum entry \(-t\) at the last row and first pivot column:
in the chosen idempotent the coordinates of \(x,y'\) are the
first and last basis vectors; pairing their conjugates into the
dual basis gives the last and first vectors, with the minus sign
from the skew target pair. The two entries have zero mutual
Heisenberg half-pairing.
\src{src/Development/WeilBound/ThetaLifts/LocalTheory/StandardValuations/AdelicInclusionMapPlaceSection.lean}{218}{Paired-coordinate identities for the unprimed displacement}.

In the \(S_A\)-evaluation one integrates the first pivot column
and the first source row except for the last pivot column. This
follows by using the translations in \(V\otimes Ex'\) first
and then those in \(Ex\otimes W_1\); the dual modulation
coordinates are set equal to zero. Thus the pivot matrix has
the form
\[
      \begin{pmatrix} *&*&s\\ *&D_1&0\\ u&0&0\end{pmatrix}
\]
with Fourier factor \(\psi_v(-tu)\). Its determinant is
\(-su\det D_1\). Apart from \(u\), the integrated pivot entries
have unit volume, and the \(a'\) integrated nonpivot entries of
the first row give \(Q^{-a'c}\). Replacing \(u\) by \(-u\)
removes the two minus signs. With
\(D_\chi=\mathbf1_{\mathcal O^\times}\chi^{-1}\), this yields
\(Q^{-a'c}D_\chi(s)\widehat D_\chi(t)\phi_1\).
The compact source average fixes the test. The Mellin weight is
\(\chi(s)\chi(t)^{-1}|st|^k\). The \(s\)-integral is one; the
\(t\)-integral is supported on \(t=\pi^{-c}\mathcal O^\times\),
where \(\widehat D_\chi(\pi^{-c}u)=
 Q^{-c}\chi(u)G(\chi,\psi_v)\) by the unit residue sum.
Consequently this integral is \(Q^{c(k-1)}\chi(\pi)^c
G(\chi,\psi_v)\). Cancel the common factor \(Q^{-a'c}\) and use
Proposition~\ref{threal:localcone} to pass from this nonzero
smaller vector to all tests.
\end{proof}

For an adapted integral plane let \(\mathfrak b x\) and
\(\mathfrak c x'\) be its two first-line lattices. The duality
condition on the skew lattice, using the absolute trace, gives
the second target line lattice
\(\mathfrak D_{E/\mathbb Q}^{-1}\bar{\mathfrak c}^{-1}y'\).
Indeed a coefficient must pair with every element of
\(\mathfrak c x'\) into the codifferent, and the other integral
summands are orthogonal. Thus the finite displacement
\(l x\otimes y'\) is integral precisely for \(l\) in
\(\mathfrak b\mathfrak D_{E/\mathbb Q}^{-1}
  \bar{\mathfrak c}^{-1}\widehat{\mathcal O}_E\).
These ideals are units at the specified ramified places by the
prescribed paired integral bases and total rational splitting.
The relative different for \(E/F\) is trivial and
\(\nu|_{\mathbb A_F^\times}=1\), by the common parity of the
two splitting characters. Hence the character of the codifferent
is one, and that of the displacement ideal is
\(\nu(\mathfrak b)\nu(\mathfrak c)\); the norm factor with
exponent \(-k\) depends only on the chosen lattices. This gives
the two positive ideal exponents for the \(a\)-step. In particular
the splitting-prime calculation of the two comparison factors also
allows nonprimitive line generators by the corresponding change of
variable in the local Mellin test of the lattice vacuum.
\src{src/Development/WeilBound/ThetaLifts/LocalTheory/ConeShifts/IdealId0.lean}{49}{The two displacement ideals and their character values}.

\paragraph{The scalar comparison for \(p+q>0\).}
The common local maps give an algebraic test for the
global discrepancy. Fix polynomial real and finite tests and take the
restricted tensor products with the specified vacuum lines. Let
\(P,P'\) be the tensor maps to the common quotient, using the
full-rank restriction on the primed domain. Suppose the quotient
and full-rank surjectivities in \S\ref{newtheta:global} have been
established at the places in question. Then \(P,P'\) are onto:
each pure tensor has finitely many non-vacuum factors; choose lifts
factor by factor and use the prescribed vacua elsewhere. Work with
the rational summation functional averaged over the source, denoted
by \(I_h,I'_h\). For fixed \(h\) its action on source differences
vanishes by change of variables wherever the summation descends to
the source quotient. The first additional global requirement is
\[
            I_h(\ker P)=0,\qquad I'_h(\ker P')=0 .
\]
In particular a local quotient that kills the source differences
does not identify the annihilator of the global integral unless its
extra kernel is killed too. Once these equalities hold, the two
periods give well-defined linear maps \(A,A'\) from the common tensor
module to functions of \(h\). Proportionality is equivalent, when
\(A'\ne0\), to the following directly testable equalities for all
module vectors \(x,y\) and adelic \(h,h_0\):
\[
                   A(x)(h)\,A'(y)(h_0)
                    =A'(x)(h)\,A(y)(h_0).
\]
Choose \(y,h_0\) with \(A'(y)(h_0)\ne0\) to prove this equivalence;
the scalar is nonzero exactly when \(A\ne0\). Corollary~\ref{threal:limits}
then gives the assertion on the entire matching space. Conversely a
nonzero primed period on some supported Schwartz test already gives
one on a polynomial test: keep its finite factors fixed, approximate
each archimedean Schwartz coefficient by its finite Hermite
expansions, and use the continuity proved above. Hermite expansions
converge in all Schwartz seminorms, by iterating the harmonic
oscillator on the coefficients and using the creation and annihilation
operators to control the seminorms of its eigenfunctions.

Theorem~\ref{threal:base} establishes the scalar comparison at the
definite base. Theorem~\ref{threal:nonzero} proves \(A'\ne0\) also
for \(q=0<p\) whenever the factorization exists, and gives a nonzero
supported period there before factorization. The graph-orbit proof
of the base comparison changes once \(W\) has isotropic lines:
maximal isotropic subspaces of \(W\oplus -W\) can project with
nonzero kernels, and their terms survive in the unfolding unless
one proves the necessary vanishing. Proposition~\ref{threal:allreal}
gives the real quotient in these signatures, and
Proposition~\ref{threal:detectline} reduces the global comparison to
one entire line constant term for all target translates.

The rank restriction on other available multiplicity formulas
matters here. The determinant characters on the two inducing blocks
give the two Satake strings of lengths \(a,a'>1\) at the completely
split good primes. The natural Arthur parameter with these strings
has two character blocks with nontrivial Arthur \(\mathrm{SL}_2\);
the generic multiplicity-free formula does not apply to it.
Chen--Zou's Theorem~7.7 treats elliptic parameters at \(F\)-rank
one and also those for which each nontrivial Arthur block has
cuspidal degree greater than the rational Witt index. Here those
degrees are one. Moreover this theorem describes a discrete
spectrum in terms of packets, defined as multisets; using it on a
specific period would require proving discrete realization and
identifying its packet member. Even a multiplicity of one does not
give the ratio of functionals in different ranks.
\cite[\S5.3, Assumption~7.2 and Theorem~7.7]{chenZouMultiplicity}

In the descent in \(a\), the actual comparison of tests must send
the larger cone integral to the smaller zero evaluation times
\(C_A\,\mathrm{stepFac}_A\,L^{S_0}(\nu,(a-a')/2)\).
Proposition~\ref{threal:coneA} identifies the two global constant
terms with those operations and Corollary~\ref{threal:endA} proves
the required detection. Identifying their images in the smaller
matching quotients with the uniform scalar is a separate assertion.
In the descent in \(a'\), the smaller cone integral must match the
larger zero evaluation times
\(C_B\,\mathrm{stepFac}_B\,L^{S_0}(\nu,(a-a')/2+1)\), and must
land in the smaller supported space. Lemma~\ref{tpd:mellin}
proves that support assertion, Proposition~\ref{tpd:unfold}
identifies the cone component, and Proposition~\ref{threal:endB}
proves detection by it in every signature. The signs in the two
finite products and the conductors are as in \S\ref{newtheta:global}.
The constants \(C_A,C_B\) precede the quantifiers over the characters,
both Weil systems, and both sets of finite quotient maps once the
evaluations, real maps and measures are fixed.
Proposition~\ref{threal:localcone} determines both ratios on all
tests at the integral good split places in its hypotheses.
Proposition~\ref{tpd:ordinary} extends these identities to every
unramified place split in \(E/F\), including resonant characters
and the factors from the adapted line lattices.
Corollary~\ref{threal:ramA} and Proposition~\ref{tpd:ram} determine
the ratios at the designated ramified split places.
At inert places Proposition~\ref{threal:allinert} gives the
irreducible quotient and its common vacuum. Propositions~\ref{tcr:descend} and \ref{tcr:unique} prove that the
cone and zero operations factor through the upper quotient on every
local test and compute their ratio after projecting to the smaller quotient. At the real places
Proposition~\ref{threal:allreal} identifies the two modules in
each fixed rank. Proposition~\ref{tcr:unique} compares the compact-averaged
Mellin operators and evaluations between the two ranks on the
polynomial vectors for the fixed real quotient maps. Theorem~\ref{tcr:global}
assembles the comparisons for the prescribed polarizations, lattice vacua
and product measures, giving the uniform constants and nonvanishing in
positive signature.
\src{src/Development/WeilBound/ThetaLifts/ConeHypotheses/PolDense.lean}{46}{The two normalized cone comparisons};
\src{src/Development/WeilBound/ThetaLifts/ConeHypotheses/PolDense.lean}{32}{Detection by the selected constant-term component}.

\subsection{The positive primed Witt index}
\label{tpd:main}

Put \(k=(a-a')/2\), and let \(H\) denote a hyperbolic plane of the
appropriate sign. In the step \(q\mapsto q-1\), the unchanged source
is \(V\), the other source changes from \(V'_1\oplus H\) to \(V'_1\),
and the target from \(W=W_1\oplus H\) to \(W_1\). We use the
splitting characters of \S\ref{newtheta:global} throughout this tower.
Write \(P=MN\) for the target parabolic stabilizing the first
isotropic line of \(H\), and normalize \([N]\) to have volume one.
Write \(m(t,h_1)=\operatorname{diag}(t,h_1,\bar t^{-1})\).

\begin{proposition}[The unprimed coefficient when \(p=0\)]
\label{tpd:zero}
Suppose \(p=0\) and \(q>0\). In rational mixed coordinates, let
\(e_0\) be evaluation at zero in the new \(V(\mathbb A_F)\)
variable. For the coherent mixed realizations and the same Haar
measure on the unchanged source, the full, unnormalized constant term is
\[
 (I_\Phi)_P(m(t,h_1)h)=
 \beta(t)|t|_{\mathbb A_E}^{a/2}\,
 I^{W_1}_{e_0(\omega(1,h)\Phi)}(h_1).
\]
Some polynomial archimedean test with an allowable finite factor has
nonzero coefficient. Every unprimed lift with polynomial archimedean
test is orthogonal to the cuspidal spectrum of \(U(W)\).
\end{proposition}
\begin{proof}
In the mixed model the rational summation variables are \((x,y)\),
where \(x\in V(E)\) and \(y\) belongs to a rational position space
for \(V\otimes_E W_1\). The center of \(N\) acts through
\(\psi(bc\langle x,x\rangle_V)\), with \(b\in\mathbb A_F\) and
a fixed nonzero rational factor \(c\in F\) determined by the paired
basis of \(H\). Its integral over \(F\backslash\mathbb A_F\) selects
\(\langle x,x\rangle_V=0\). Definiteness gives \(x=0\).
The other coordinates of \(N\) act on the smaller oscillator through
Heisenberg operators whose translation and modulation parameters are
linear in \(x\), with quadratic central phase zero at \(x=0\).
They therefore act trivially on this value. Both unipotent and source
quotients are compact, and the rational Schwartz sum is absolutely
and uniformly convergent on compact sets of representatives; this
justifies both integrations. The line Levi acts at zero by
\(\beta(t)|t|_{\mathbb A_E}^{a/2}\), the square root of the modulus
of scalar substitution on the \(a\) coordinates times the prescribed
Siegel character. The smaller target uses the same splitting; on the
source the added plane acts on this variable by pullback. This is
also the mixed-model tower formula for a cuspidal source, applied to
the trivial representation of anisotropic \(U(V)\).
\cite[Proposition~3.1, (3.2)--(3.4)]{wuThetaTower}

There are \(q\) rational hyperbolic planes in \(W\). At split finite
places it is split, at inert places it is hyperbolic by the perfect
lattice calculation of Proposition~\ref{threal:allinert}, and at
every real place its Witt index is \(q\). To split off one plane,
multiply the skew form by a nonzero trace-zero element of \(E\) and
apply Hasse--Minkowski to the resulting quadratic norm form on the
underlying \(F\)-space. Complete an isotropic vector to a hyperbolic
pair. Witt cancellation preserves the remaining local indices, so
the argument can be iterated. The bottom \(X\) is definite of
dimension \(d_0=d-2q\).

The unprimed nonvanishing calculation in Theorem~\ref{threal:base}
applies to \((V,X)\): its argument up to the evaluation of
\(Z((d_0-a)/2)\) uses only these two spaces. Indeed
\(a'-2q>1\), \(d_0\) is even and the infinite angular exponent still
exceeds \(d_0\). At inert places \(X\) is hyperbolic and admits a
self-dual lattice, and at split places one may take dual bases.
These lattices agree almost everywhere with the lattice obtained by
a rational orthogonal decomposition. The denominator arguments
are at least \(a'-2q+1>1\) and the numerator integers are
\(d_0/2-i\) for \(0\leq i<a\); the same functional equation and
nonvanishing at one apply. At infinity the unprimed invariant
determinant used there has exponent \((m_\alpha-m_\beta-d_0)/2>0\)
and there are \(d_0\geq a\) columns. This gives a nonzero bottom
period.

Evaluation in each new variable is onto the smaller Schwartz
space by tensoring with a test of value one at zero. These extensions
are restricted tensors, since away from finitely many places the
chosen rational hyperbolic basis and its inverse are integral. The
constant-term formula gives nonvanishing successively in the tower.
With the finite test fixed, the resulting nonzero coefficient at
a fixed Levi element is a continuous functional of the infinite
Schwartz test, by compactness of the integrating quotients.
Polynomial times Gaussian tests in the original coordinates are
dense in the Schwartz topology, giving the asserted polynomial test.

For orthogonality we keep the unconjugated convention of Wu, in which
the two directions use inverse additive characters. Set
\(Y=V\), and let \(\chi_1,\chi_2\) be the splitting characters on
the target \(U(X)\) and the source \(U(Y)\), respectively, in
Wu's labeling (the twists below mean composition with determinant).
The bottom lift \(\sigma=\theta^X_{Y,\psi}(\mathbf1)\) is nonzero
and cuspidal since \(X\) is anisotropic. The input is itself an
irreducible cuspidal representation. Wu's Theorem~5.5 gives
irreducibility of \(\sigma\) and Theorem~5.4(1), with the two
additive characters in their indicated order, gives
\[
 \theta^Y_{X,\psi^{-1}}(\chi_1^{-1}\sigma)=\chi_2\mathbf1.
\]
The input and output in this last equality are again nonzero,
irreducible and cuspidal. Applying Theorem~5.1(2) with
\(\pi=\chi_1^{-1}\sigma\) therefore shows that
\(\theta^{X\oplus H^q}_{Y,\psi}
 (\chi_2^{-1}(\chi_2\mathbf1))\) pairs to zero with every cusp
form for \(q>0\). This is the asserted unprimed lift.
\cite[Theorems~5.1, 5.4 and 5.5]{wuThetaTower}
\end{proof}

\begin{proposition}[Cuspidal orthogonality of the supported primed lift]
\label{tpd:orth}
For every \(q>0\), including \(p>0\), every supported primed
theta integral pairs to zero with all cuspidal forms on \(U(W)\).
\end{proposition}
\begin{proof}
First take a compact-finite automorphic vector \(u\) in a single
irreducible cuspidal constituent on \(J=U(W)\). Such vectors span
a dense subspace of the cuspidal spectrum. Set
\[
 B(g)=\int_{[J]}\Theta_{\Phi'}(g,h)\overline{u(h)}\,dh,
        \qquad g\in U(V')(\mathbb A_F).
\]
Use Siegel representatives for \(h\) in this integral. Rapid decay of
\(u\) there dominates every fixed positive power of \(H(h)\), by
the rapid decrease theorem for cusp forms. Proposition~\ref{threal:decay}
then permits Fubini, and \(B\) is a rapidly decreasing, square-integrable
reverse theta lift. More explicitly its proof gives a bound
\(C_{L,u}\prod_i T_i^{-L}\) on each source Siegel set, with the
same bound when the right compact set of representatives is enlarged.
The finite supported condition is preserved by translations in either
group, so the estimate also applies to every fixed right translate.

Consider any rational maximal parabolic of \(U(V')\), stabilizing a
totally isotropic \(j\)-plane. In the reverse direction the input
\(\bar u\) is cuspidal. Proposition~3.1 of Wu applies to this
\(j\)-plane, with the roles of the Hermitian and skew-Hermitian
spaces exchanged (equivalently rescale the two forms by reciprocal
trace-zero scalars). Its unnormalized constant term at
\(m(t,h_0)g_0\), where \(t\in\mathrm{GL}_j(\mathbb A_E)\), is a
function of \(h_0,g_0\) times
\(\xi(\det t)|\det t|_{\mathbb A_E}^{d/2}\) for a unitary
splitting character \(\xi\). Take \(t\) to be the scalar \(R I_j\)
at each real place and the identity at finite places and let
\(R\to\infty\). On this positive chamber ray the integral over the
unipotent quotient tends to zero: its representatives can be taken
in a compact set which stays bounded after conjugation by
\(m(t,1)^{-1}\) on the left and \(m(t,1)\) on the right, and the
preceding bound applies. The modulus of the displayed character
grows as \(R^{[F:\mathbb Q]jd}\), forcing its coefficient to
vanish for every \(h_0,g_0\). Thus \(B\) is cuspidal.
\cite[Proposition~3.1]{wuThetaTower}

The group \(U(V')\) has a proper rational parabolic. Indeed the
given real Witt index is \(q>0\); at inert places the perfect
Hermitian lattice has Witt index \(\lfloor a'/2\rfloor\), and at
the remaining finite places it is split. The same norm-form
argument as above supplies a rational isotropic line.
The orthogonal projection in \(L^2([U(V')])\) of the constant
function to the closed cuspidal subspace is right invariant, hence
constant. Its constant term at the line parabolic must vanish, so
this projection is zero. It follows that \(\int B(g)\,dg=0\).
Fubini identifies this with the asserted pairing.
\end{proof}

\begin{lemma}[Removal of the primed zero orbit]\label{tpd:primezero}
Suppose \(p=0<q\). In taking the constant term of the supported
primed integral along \(P\), one can choose a test giving the same
automorphic function for which the \(x=0\) summand vanishes
identically, also after any target translation.
\end{lemma}
\begin{proof}
At one real place average the upper test over \(K_+=U(q)\), the
positive compact block of \(U(V')\). The unchanged source integral
follows by right invariance of Haar measure, compact averaging and
Proposition~\ref{threal:decay}. The full-rank condition at \(v_1\)
is preserved. Mixed evaluation at \(x=0\) is \(U(V')\)-equivariant
to the smaller oscillator for \((V',W_1)\): on the extra variable
the group acts by pullback, using the coherent splitting on decreasing
the partner by a hyperbolic plane. Thus its value on the averaged
test belongs to the \(K_+\)-invariants of that oscillator.

Write \((r_1,t_1)=(d-q-1,q-1)\) for the signature of \(W_1\).
As in the entry calculation of Proposition~\ref{threal:allreal},
the \(K_+\)-vacuum character in the primed oscillator is
\(\det^{b}\), where \(b=(m_\alpha-m_\beta+r_1-t_1)/2>0\).
There are \(r_1\) standard and \(t_1\) dual column variables
for this block. By the Cauchy decomposition the compact types can
only come from
\(\det^b\otimes S_\lambda(\mathbb C^q)
                    \otimes S_\mu(\mathbb C^q)^*\),
with \(\ell(\mu)\leq t_1=q-1\); variables from other rows are
trivial on \(K_+\). An invariant would equate the highest weights
\(\mu=\lambda+(b,\ldots,b)\), impossible at the last entry.
Hence the compact average is zero on the polynomial module, and
by density and continuity on the Schwartz space. Target translations
commute with compact source averaging. The smaller evaluation still
vanishes after applying any source element by equivariance.

On a compact truncation of the integrating source, average the
center of \(N\) first. Additive orthogonality leaves the sum over
the rational null cone, and the zero summand just treated vanishes
pointwise. Taking an exhaustion of the source quotient gives the
constant term as the limit of these nonzero-orbit expressions;
the original kernels are dominated by Proposition~\ref{threal:decay}
uniformly over compact representatives for \([N]\). This formulation
requires no separate integrability assertion for an unrestricted
smaller zero-orbit period.
\end{proof}

\begin{corollary}[Detection for a definite unprimed source]\label{tpd:detect}
If \(p=0,q>0\), two matched polynomial tests give proportional
automorphic functions with scalar \(c\) as soon as their full line
constant terms are proportional with that scalar at every target
translate. For the unprimed form this constant term is the smaller
period of Proposition~\ref{tpd:zero}; for the primed form it consists
only of the nonzero orbit of Lemma~\ref{tpd:primezero}.
\end{corollary}
\begin{proof}
For any \(q\) this follows from Proposition~\ref{threal:detectline}.
We give a direct argument in rational Witt index one.
Suppose \(q=1\).
The rational Witt index of \(W\) is one by its signatures and the
isotropy proved above. Every proper rational parabolic stabilizes a
line and these lines are rationally conjugate by Witt's theorem. The
difference \(D\) of the forms thus has every proper constant term zero.
Here one can see square integrability directly. The seminorm bounds
in Proposition~\ref{threal:decay} give a single moderate growth
exponent for all right archimedean derivatives of \(D\). Indeed
choose the source decay order in that proof once, large enough for
source integration. Applying a right derivative replaces the infinite
test by its derived oscillator action. The orders of the seminorms
used to bound the translated test and the powers of the target
height in those bounds depend on this fixed decay order and the
dimensions. Only the seminorm values of the initial test change.

On a rank-one Siegel set write the target representative as \(tc\)
with positive chamber coordinate \(R\ge1\) and bounded \(c\).
At fixed finite level the compact unipotent quotient in its constant
term reduces to a compact real nilmanifold. Indeed strong
approximation on the unipotent group follows successively for its
additive root spaces, and the finite open stabilizer can be taken
common to the bounded finite components. On this nilmanifold the
function \(n\mapsto D(ntc)\) has mean zero. Let \(\Delta\) be
the sum of negative squares of a left-invariant real basis; Haar
volume is invariant and its kernel consists of constants. Elliptic
Sobolev embedding and the spectral gap on the mean-zero subspace
bound a point value by \(C_j\|\Delta^jD(\cdot\,tc)\|_{L^2}\)
for every sufficiently large \(j\). Each basis derivative becomes
a right derivative on \(D\) with Lie algebra argument
\(\operatorname{Ad}((tc)^{-1})X=O(R^{-1})\), on choosing the
coordinate with smallest positive root \(R\). The uniform growth
bound on the compact set of unipotent representatives now gives
\(D(tc)=O_j(R^{B-2j})\) for a fixed \(B\). Thus \(D\) is rapidly
decreasing and belongs to the cuspidal subspace of \(L^2\).
Propositions~\ref{tpd:zero} and~\ref{tpd:orth} make it orthogonal
to the cuspidal spectrum, hence to itself.
\end{proof}

\begin{lemma}[The Mellin exponent and the good Euler factors]
\label{tpd:mellin}
Choose paired bases \(x,y\) and \(x',y'\) of the source and target
hyperbolic planes, respectively. In the rational mixed polarization
for simultaneous removal of these two planes, take the following
isotropic space to act by modulations:
\[
 S=Ex\otimes_E W\ \oplus\ V'_1\otimes_E Ex'.
\]
Write \(e'\) for evaluation at zero in its dual position coordinates,
and \(\rho\) for the Heisenberg action. With \(K\) the product of the
specified source compact stabilizers, the integral
\[
 T\Phi'(z_1)=\int_K\int_{\mathbb A_E^\times}
   \big[e'\rho(\bar l\, y\otimes x')\omega'(k_0,1)\Phi'\big](z_1)
           \,\nu(l)|l|_E^{k+1}\,d^\times l\,dk_0
\]
is absolutely convergent and defines a smaller Schwartz function.
It preserves the required full-rank support at \(v_1\).
Its target line character is \(\beta(t)|t|_E^{a/2}\).
With local unit groups of multiplicative mass one almost everywhere,
the scalar from the unmodified unramified factors outside any
sufficiently large finite set is the corresponding partial Euler
product for \(L_E(\nu,k+1)\).
\end{lemma}
\begin{proof}
The chosen subspace has dimension \(d+a'-2\) over \(E\);
its symplectic perpendicular modulo itself is the smaller tensor
space \(V'_1\otimes_EW_1\), by the two orthogonal decompositions.
We can choose its dual coordinate space to contain the translation
line \(E(y\otimes x')\), paired nondegenerately with
\(E(x\otimes y')\). Consequently the bracket in the integral is the
restriction of a mixed-model Schwartz function to a linear subspace
with coordinates \((\bar l,z_1)\), allowing a fixed sign or unit
change in the first coordinate. A rational choice of these
coordinates is integral with integral inverse almost everywhere,
where the vacuum is a nonzero multiple of the product lattice
indicator and the compact source fixes it. Choosing each original
basic vector and its image as the smaller basic vector removes the
same multiple on both sides. For a fixed finite test the compact
finite orbit is finite, by smoothness; at infinity the compact
orbit is bounded in every Schwartz seminorm. Finite-place restrictions
are finite sums of products of locally constant compactly supported
functions of \(l\) and \(z_1\) off the basic places. At infinity
all smaller Schwartz seminorms and derivatives in \(z_1\) are
dominated by arbitrarily rapidly decreasing majorants in \(l\).
Since \(k+1>1\), the integrals at zero converge. At a basic place
\(w\) of \(E\) the integral is the geometric series
\(\sum_{j\ge0}\nu_w(\pi_w)^j(Nw)^{-j(k+1)}\) times the smaller
basic vector. The absolute products converge since the Dedekind
zeta series converges at \(k+1\). This proves the analytic and Euler
factor assertions, including differentiation under the integrals.

At \(v_1\), choose the two conjugate idempotent factors of the
split tensor space as position and momentum in the specified primed
matrix model. Order the rows by \((x,V'_1,y)\) and the columns
by \((x',W_1,y')\) on the position side. The complementary momentum
basis pairs \(x\) with \(y\) and \(x'\) with \(y'\). In passing to
the above mixed evaluation, the position vectors belonging to \(S\)
are integrated out by partial Fourier transformation. They are the
first row and the first column of the middle rows. The momentum
vectors belonging to \(S\) require zero in the last row and the last
column of the middle rows. Translation by \(\bar l\,y\otimes x'\)
puts a nonzero entry in the first position of the last row; its other
idempotent component contributes only a Fourier phase. Thus, for
middle position matrix \(Z\), the original test is evaluated solely
on matrices of the form
\[
   \begin{pmatrix} *&*&*\\ *&Z&0\\ b&0&0\end{pmatrix},\qquad
                                      b\in F_{v_1}^{\times}.
\]
Changing the chosen idempotent uses the corresponding conjugate
bases and the same pattern. Clearing the middle entries of the
first column by the last row bounds the full rank by
\(2+\operatorname{rank}Z\). The upper supported test, also after
the compact source action by invertible row operations, vanishes
on all these matrices if \(\operatorname{rank}Z<a'-2\).
Every partial Fourier integral then vanishes at \(Z\), proving
the smaller support assertion.

At infinity an adapted Cartan involution preserves the orthogonal
decomposition into the tensor blocks formed from the two planes
and their complements. In particular the smaller symplectic tensor
is a complex-orthogonal summand for the compatible positive
structure. The Fock vacuum factors over this summand and its
complement; the polynomial module is spanned by finite products of
their creation monomials. Compact source translates of a polynomial
test remain in the finite span of fixed degrees. The evaluation and
translation in the integral take place only in the complement.
The smaller real output is therefore polynomial for polynomial
input, in the coordinates transported by the smaller Cartan
structure. The moments in the complement converge by the same
Gaussian bounds.

Furthermore fixing a nonzero common pair of real matching maps fixes
the infinite component of \(\nu\). Given two allowed systems for
these same maps, their primed source actions on a compact unitary
block at one real place differ by the determinant character
\(\det^{(m_\nu-m_{\widetilde\nu})/2}\), by the doubled splitting
calculation in Proposition~\ref{threal:allreal}. Both actions leave
the same quotient map unchanged. Apply that map to a polynomial
with nonzero image and translate by this compact block. The
determinant character is identically one since the block determinant
ranges over the circle, so \(m_\nu=m_{\widetilde\nu}\) there.
The restriction of each character to \(\mathbb A_F^\times\) fixes
its radial part at infinity. Varying the place proves the claim.
In particular the real Mellin operator itself is fixed across
the allowed characters and actions for fixed real maps, after the
same coordinates and Haar measures are chosen.

The source line Levi \(m(l,1)\) has determinant modulus
\(|l|_E^d\) on \(S\), since it scales the first summand by \(l\).
The corresponding evaluation character is \(\nu(l)|l|_E^{d/2}\).
Indeed first split the source plane, with modulation space
\(Ex\otimes W\) and its prescribed Siegel character \(\nu\),
then split the target plane in \(V'_1\otimes W\). On the source
unipotent radical the line Levi acts on \(a'-2\) coordinates over
\(E\) with determinant modulus \(|l|_E^{a'-2}\), and on the central
\(F\)-coordinate by \(l\bar l\), adding the factor \(|l|_E\).
Its inverse modular character in Iwasawa integration is therefore
\(|l|_E^{1-a'}\); together these give \(d/2+1-a'=k+1\).
The target line has determinant modulus \(|t|_E^{a'-2}\) on \(S\):
its action on the full \(W\) in the first summand has modulus one.
At evaluation its scalar is \(\alpha(t)|t|_E^{(a'-2)/2}\), from
the residual target Siegel action (keep the target splitting
character on reducing the source by a plane). Moving this target
element past the translation replaces \(l\) by \(l/\bar t\).
In the Mellin integral substitute \(l=u\bar t\). Since
\(\nu(\bar t)=\nu(t)^{-1}\), the resulting phase is
\(\alpha(t)/\nu(t)=\beta(t)\) and the exponent is
\((a'-2)/2+k+1=a/2\), as claimed. This computation explains both
the choice of the translation line and the shift by one in this
branch.
\src{src/Development/AbelianVarietyTheory/Definitions/Divisors/CTSmallCone.lean}{57}{The two-plane evaluation, cone and translation};
\src{src/Development/WeilBound/ThetaLifts/SmallCones/ConeLawAB.lean}{115}{The primed two-Levi character}.
\end{proof}

The two fractional-ideal exponents in this branch can also be fixed.
By adaptedness the line \(Ex\) has lattice \(\mathfrak b x\) and
the line \(Ex'\) has lattice \(\mathfrak c x'\). The source lattice
is Hermitian self-dual and the two lines and the orthogonal smaller
lattice split integrally. Pairing with \(\mathfrak b x\) therefore
shows that the lattice of \(Ey\) is \(\bar{\mathfrak b}^{-1}y\).
It follows by taking direct summands in the tensor lattice that the
finite translations \(\bar l\,y\otimes x'\) belonging to it have
\[
                l\in\mathfrak a\widehat{\mathcal O}_E,qquad
                \mathfrak a=\mathfrak b^{-1}\bar{\mathfrak c}.
\]
Both ideals \(\mathfrak b,\mathfrak c\) are prime to the conductor
of \(\nu\). At the chosen half of a ramified split place, this
follows because both paired vectors occur in the specified integral
basis. At the conjugate half, apply duality to the second vector in
that basis; the target's trace codifferent is a unit at the
designated totally split rational primes. At unramified places
the value of \(\nu\) on a uniformizer is independent of its unit
multiple, and \(\nu(\bar l)=\nu(l)^{-1}\). Consequently
\(\nu(\mathfrak a)=\nu(\mathfrak b)^{-1}
\nu(\mathfrak c)^{-1}\). Integrating the line-lattice indicator
outside the conductor contributes this character value, its absolute
norm to the power \(-(k+1)\), and the elementary Euler factors.
The norm power depends only on the fixed lattice and ranks.
Thus the lattice calculation gives \(e_1=e_2=-1\).
\src{src/Development/WeilBound/ThetaLifts/LocalTheory/ConeShifts/IdealId0.lean}{49}{The character on the translated lattice line}.

\begin{proposition}[Unfolding the supported nonzero orbit]
\label{tpd:unfold}
Suppose \(q>0\). In the rational mixed models the
\(\beta|\cdot|_E^{a/2}\) component of the line constant term of
any supported primed test equals
\[
 C_1\,\beta(t)|t|_E^{a/2}
 I'_1\bigl(T(\omega'(1,h)\Phi')\bigr)(h_1).
\]
Here the line Levi element is \(m(t,h_1)\). The difference from
the entire constant term transforms by \(\alpha|\cdot|_E^{a'/2}\).
It vanishes identically if \(p=0\).
For self-dual Fourier changes \(C_1>0\) depends only on the quotient
and Iwasawa measures and the fixed spaces. It is independent of the
characters and the matching maps.
\end{proposition}
\begin{proof}
We first compute at the target identity for a supported test
whose zero evaluation in the target-first model vanishes identically.
Write this condition as \(e_{W,0}\Phi'=0\). When \(p=0\) the real
average of Lemma~\ref{tpd:primezero} supplies such a test and leaves
both the upper integral and \(T\Phi'\) unchanged, the latter by
right invariance of Haar measure on \(K\).
Put \(Y=V'\), \(Y_1=V'_1\). At a fixed
source element \(g\), integration over the target unipotent center
in the target-first model leaves only null vectors in \(Y(E)\).
Their nonzero terms group by their rational lines. Fix one such
line \(Ex\); its nonzero vectors are \(u x\), \(u\in E^\times\).
Having chosen its companion \(y\), polarize the remaining tensor
\(Y\otimes_EW_1\) using \(Ex\otimes_E W_1\) as modulation space
and a fixed position space for \(Y_1\otimes_EW_1\).
For this vector \(ux\), the remaining target radical acts through
the Heisenberg modulation paired with \(ux\otimes w\),
\(w\in W_1(\mathbb A_F)\). Integrating it over the additive quotient
sets the corresponding rational \(W_1\) position to zero, since
\(u\ne0\). This is the character orthogonality in
\cite[(3.15)--(3.20)]{wuThetaTower}. The line contributes
\(\sum_{u\in E^\times}\sum_{z_1} A_g(u,z_1)\), where \(z_1\)
runs over the rational smaller positions and \(A_g\) is the
restriction of the resulting Schwartz test to the indicated positions.
The source element acts on the test before these operations.

There are now two natural polarizations of the removed coordinates.
The modulation space of this last evaluation is
\[
 S_A=Y\otimes Ex'\ \oplus\ Ex\otimes W_1 .
\]
Its common subspace with \(S\) of Lemma~\ref{tpd:mellin} has
codimension one over \(E\) in both. The two remaining dual directions
are \(y\otimes x'\) and \(x\otimes y'\). Thus replacing \(S_A\)
by \(S\) is a partial Fourier transform in a single \(E\)-coordinate
and leaves \(z_1\) unchanged. Explicitly, on restricting all common
evaluation coordinates to zero, its Fourier transform is
\[
 B_g(l,z_1)=[e'\rho(\bar l\,y\otimes x')\omega'(g,1)\Phi'](z_1).
\]
The trace pairing on the two lines is a nondegenerate rational
pairing; a fixed rational change of variable or sign permutes
\(E^\times\) and has adelic modulus one. With self-dual measures
Poisson summation therefore identifies the sums over \(E\) of
\(A_g\) and \(B_g\), separately at each \(z_1\).
Both missing zero terms vanish. For \(A_g(0,z_1)\) this follows
from \(e_{W,0}\Phi'=0\), since this evaluation intertwines the
source action before the second polarization and evaluation.
For \(B_g(0,z_1)\) use the
matrix in Lemma~\ref{tpd:mellin} with \(b=0\): its last row is
identically zero at \(v_1\). This remains so for the test translated
by \(g\), since source translations preserve full-rank support.
Consequently the sums over \(E^\times\) agree. At fixed \(g\)
the original grouping over lines is legitimate by absolute
Schwartz summation and compact integration in the target radical.

To unfold these line sums set \(D_x=\operatorname{Stab}_{U(Y)}(x)\).
It is the semidirect product of its line unipotent radical \(N_Y\)
and \(U(Y_1)\). The source group is transitive on its rational
nonzero null vectors by Witt extension. In the new expression put
\(f(g)=\sum_{z_1} B_g(1,z_1)\). This function is invariant under
\(N_Y(\mathbb A_F)\): an element of this group changes \(y\)
by a vector in \(Ex+Y_1\), so the change in the translation lies
in \(S\). Its Heisenberg commutator is trivial because \(Ex'\)
is isotropic; evaluation itself is invariant under the unipotent.
For the last assertion split the source plane first: the unipotent
acts on the smaller oscillator through Heisenberg operators with
parameters linear in its extra position, and its center by a
quadratic phase, all trivial at zero. The following evaluation on
the smaller target plane preserves this invariance.
The \(U(Y_1)(F)\) action preserves the rational summation.
For the source line Levi \(m_Y(l,g_1)\), moving it past
\(\rho(y\otimes x')\) replaces \(y\) by \(\bar l\,y\).
The simultaneous evaluation therefore gives
\[
 f(m_Y(l,g_1)k_0)=\nu(l)|l|_E^{d/2}\,
 \Theta_1'\bigl(g_1,1;
    e'\rho(\bar l\,y\otimes x')\omega'(k_0,1)\Phi'\bigr).
\]
For rational \(l\) the scalar is one. Hence summing over rational
lines and their nonzero scalar parameters gives exactly the
periodization of \(f\) over \(D_x(F)\backslash U(Y)(F)\).

We check absolute integrability of this periodization at the level
of the signed smaller theta sum. Integrate first on
\(D_x(F)\backslash U(Y)(\mathbb A_F)\). Iwasawa decomposition for
the source line and the chosen maximal compact gives, up to fixed
positive normalization, integration over \([N_Y]\),
\([U(Y_1)]\), \(\mathbb A_E^\times\) and \(K\), with density
\(|l|_E^{1-a'}\). The finite compact is in good position: it is
hyperspecial for the self-dual source lattice at the inert places,
and is an integral general linear compact at the split places.
The quotient \([N_Y]\) is compact with fixed volume. The absolute
integral remaining is bounded by the integral of the absolute
smaller theta kernel against \(|l|_E^{k+1}\). At the distinguished
finite place each smaller test in this integral has full-rank
support by Lemma~\ref{tpd:mellin}. More explicitly the compact
locally constant test in \((l,z_1)\) has only finitely many slices
in the first variable at each exceptional finite place. At \(v_1\)
all slices occurring for invertible \(l\) have the required support.
At infinity all seminorms needed in Proposition~\ref{threal:decay}
for the smaller source are rapidly decreasing in \(l\) and uniformly
bounded near zero and in \(k_0\). The smaller ranks satisfy its
hypotheses, since \(q\) has decreased by one and \(a'\) by two.
The absolute Euler products outside the exceptional primes are
bounded by the corresponding factors of \(\zeta_E(k+1)\), as
in Lemma~\ref{tpd:mellin}. These bounds prove the required
integrability, justify unfolding the periodization, and allow the
smaller period to be interchanged with the Mellin and compact
integrals. This gives \(C_1 I'_1(T\Phi')(1)\).

Write \(A(t)=\alpha(t)|t|_E^{a'/2}\),
\(B(t)=\beta(t)|t|_E^{a/2}\), and let \(J(\Phi')\) be the entire
constant term at the identity. For arbitrary \(p\) and arbitrary
supported \(\Psi\), apply the computation just proved to
\[
   \Phi'_t=\omega'(1,m(t,1))\Psi-A(t)\Psi .
\]
This test satisfies \(e_{W,0}\Phi'_t=0\): in the target-first
model scalar substitution on the \(a'\) coordinates at zero gives
exactly \(A(t)\). Both summands retain the required split support.
Writing \(J_B(\Psi)=C_1 I'_1(T\Psi)(1)\), the result is
\[
 J(\omega'(1,m(t,1))\Psi)-A(t)J(\Psi)
                   =(B(t)-A(t))J_B(\Psi),
\]
using the Mellin covariance in Lemma~\ref{tpd:mellin}. Thus
\(J-J_B\) transforms by \(A\), whereas \(J_B\) transforms by \(B\).
These components are unique since \(a>a'\) and the splitting
characters are unitary. If \(p=0\), averaging any test first shows
directly that \(J=J_B\). Taking \(\Psi=\omega'(1,h)\Phi'\) and
using intertwining for the smaller target factor gives the asserted
formula at \(m(t,h_1)h\). All measure constants occurred in the fixed
Haar decompositions or unipotent quotients.
\end{proof}

\begin{lemma}[A supported representative for the other coefficient]
\label{tpd:so}
Let \(q>0\). For each supported primed test \(\Psi\), one may
replace \(\Psi\) by a supported \(\widetilde\Psi\) having the
same upper period, so that its one-plane zero evaluation
\(e_{W,0}\widetilde\Psi\) is full-rank supported for the unchanged
source \(V'\) and the target \(W_1\). In the coherent one-plane
model the \(\alpha|\cdot|_E^{a'/2}\) component at \(m(t,h_1)\)
is
\[
   \alpha(t)|t|_E^{a'/2}
       I_{V',W_1}(e_{W,0}\widetilde\Psi)(h_1).
\]
For a fixed further right target translate apply the construction
to the already translated upper test.
\end{lemma}
\begin{proof}
Put \(m=a'\) and work at \(v_1\) in the primed row coordinates.
Choose dual target bases with paired first and last columns. In
this model \(e_{W,0}\phi(Z)=\int\phi(u,Z,0)\,du\), with the
corresponding conjugate choice if the position idempotent is
interchanged. In fact one of the removed columns acts by translations
and is integrated, while the paired modulation column must vanish.
An upper test of full row rank therefore maps to a compactly
supported function on the open region \(D=\{\operatorname{rank}
Z\ge m-1\}\). The map is onto that space. At a given \(Z_0\in D\)
choose \(u_0\) completing its columns to rank \(m\), and small
compact open balls in \(u,Z\) on which rank stays full. A unit-mass
indicator in the first ball and a small last-column indicator of
value one at zero lift every test supported on the second ball.
A finite disjoint refinement covers the support of any test in \(D\).

Modulo compactly supported functions on the open full-rank part
of \(D\), these lower tests restrict to all compactly supported
functions on its relatively closed rank \(m-1\) stratum. Extension
follows by compact open local charts and disjoint refinement.
The stratum's source coinvariants are zero by the rank stabilizer
calculation of Proposition~\ref{split:kernel}, applied with
\((d-2)/2=m+k-1>m\). This calculation concerns the stratum itself
and applies with the same unitary phase and the lowered modulus.
Thus, up to a full-rank lower test, \(e_{W,0}\phi\) is a finite sum
of lower source-action differences of tests in \(D\). Lift each
such test through the surjection above, and subtract the
corresponding upper differences from \(\phi\). Equivariance of
evaluation for the unchanged source proves the required support.
Use individual supported local factors and extend by linearity;
at the other places all factors stay fixed. These upper differences
integrate to zero by right translation on the source, with absolute
convergence from Proposition~\ref{threal:decay}.

In the expansion of the upper constant term the zero-vector
contribution of this representative now converges by the
full-rank estimate applied to \((V',W_1)\). Its torus character is
\(\alpha(t)|t|_E^{m/2}\); at \(t=1\) evaluation gives precisely
the smaller displayed period, using the same source measure.
For each fixed source translate the \(A_g(0,z_1)\) of
Proposition~\ref{tpd:unfold} vanishes even though
\(e_{W,0}\widetilde\Psi\) need not be zero. In fact apply the
further source-plane evaluation \(Ex\otimes W_1\) at the split
place to this lower full-rank test. All entries in one row must be
zero and the opposite row is integrated, so every value in the
integrand vanishes. This remains true for any rational null line,
by changing row basis, and for translates by the unchanged source
and smaller target, which preserve matrix rank. The other Poisson
zero term still vanishes on the upper full-rank test. The entire
nonzero orbit therefore has the \(\beta\) law by the same absolutely
convergent unfolding. Uniqueness of the two characters proves the
assertion, including independence of the resulting component from
the choice of the lift of the source differences.
\end{proof}

At a designated split place the upper to lower comparison space is
at most one-dimensional by Proposition~\ref{threal:localcone}.
Its unprimed zero-evaluation map is nonzero, since evaluation is
onto and the lower quotient is nonzero. We next compute the scalar
on one pair of tests with nonzero smaller image. This determines
the comparison on every test at this place.

\begin{proposition}[The ramified primed comparison on all tests]
\label{tpd:ram}
At a designated ramified split place let \(K_0=F_v\) have residue
size \(Q\), let \(\chi=\nu_w\) on its chosen factor of \(E\),
and let \(c>0\) be the conductor exponent of \(\chi\).
Use the prescribed bases and row polarizations, \(\operatorname{vol}
(\mathcal O)=1\) for the additive character of conductor \(\mathcal
O\), and give \(\mathcal O^\times\) and the maximal source compact
masses one. Use the integral-vacuum-normalized evaluations. On
every upper pair matched at this place, the ratio of the lower
primed quotient on \(T_v\phi'\) to the lower unprimed quotient on
\(e_{0,v}\phi\), as a scalar identity of maps, is
\[
   \epsilon_v(-k)=
      Q^{ck}\chi(\pi)^c
       \sum_{r\in(\mathcal O/\pi^c)^\times}
                         \chi(r)^{-1}\psi_v(r/\pi^c).
\]
\end{proposition}
\begin{proof}
Put \(m=a'\) and \(D_\chi(u)=\mathbf1_{\mathcal O^\times}(u)
\chi(u)^{-1}\), extended by zero at \(u=0\).
For an \(n\)-by-\(d\) matrix choose a block of \(n\) pivot columns.
Let its test be \(D_\chi\) of the pivot determinant times the
indicators that the pivot entries belong to \(\mathcal O\) and
all other entries to \(\pi^c\mathcal O\). Take the first \(a\)
columns for \(\phi\), the last \(m\) for \(\phi'\), and the same
rule for \(\phi_1,\phi'_1\) after deleting the first and last of
the \(m\) target columns. The source compact fixes each test:
integral invertible row operations preserve its entry conditions,
and the source factor cancels the inverse determinant character.

The upper tests are matched, as are the two lower tests. To check
this directly, restrict the orbital functions
\eqref{newtheta:big}--\eqref{newtheta:small} to
\(h=\left(\begin{smallmatrix}A&B\\ C&D\end{smallmatrix}\right)
\in\mathrm{GL}_d(\mathcal O)\).
A contributing row transformation must belong to the integral
general linear group: the original row frame is primitive and
the pivot of the transformed test is invertible modulo \(\pi\).
On both sides support requires \(C=0\bmod\pi^c\), in which case
the diagonal blocks are invertible. The determinant character of
the row transformation cancels throughout that integral compact.
The two values are \(\beta_w(\det h)\chi(\det A)\) and
\(\alpha_w(\det h)\chi(\det D)^{-1}\), equal by
\(\det h\equiv\det A\det D\bmod\pi^c\). Iwasawa decomposition
extends equality using their common parabolic law. The common value
at the identity is one. The same calculation in the smaller ranks
gives a nonzero common lower quotient vector.

Unprimed evaluation integrates the \(a\) entries of one nonpivot
column, setting the other removed nonpivot column to zero. Thus
\(e_{0,v}\phi=Q^{-ac}\phi_1\). For the primed calculation write
\(l=(s,t)\in K_0^\times\times K_0^\times\), the first factor
being the chosen place \(w\). Order the pivot rows and columns as
first, middle, last in the prescribed bases. In the conjugate row
polarization the position and modulation parts of
\(\bar l\,y\otimes x'\) are respectively the first-row last-pivot
entry \(s\) and the last-row first-pivot entry \(-t\).
Indeed switching from the ordinary to the conjugate row polarization
interchanges momentum and position with one minus sign; the source
pair has Hermitian off-diagonal entries \(1,1\), the target pair
has skew entries \(1,-1\), and conjugating \(l\) swaps its two
components. These coordinates occupy distinct matrix entries, so
the half-pairing in the Heisenberg action is zero.
\src{src/Development/WeilBound/ThetaLifts/LocalTheory/StandardValuations/AdelicInclusionMapPlaceSection.lean}{228}{The two entries for the specified Fourier signs}.

In this coordinate ordering the primed mixed evaluation integrates
the last row and the last pivot column in the middle rows. The
other removed entries vanish except for the position entry \(s\).
This is the matrix calculation of Lemma~\ref{tpd:mellin}, expressed
in the prescribed conjugate bases. The resulting pivot matrix is
\[
     \begin{pmatrix}0&0&s\\0&D_1&*\\ x&*&*\end{pmatrix},
       \qquad \det=-sx\det D_1 ,
\]
and the Fourier factor is \(\psi_v(-tx)\). The \(a\) integrated
nonpivot entries of the last row have joint volume \(Q^{-ac}\).
The other integrated entries, apart from \(x\), range independently
over \(\mathcal O\) of mass one. Since \(s,x\) and the entries
of \(D_1\) must be integral, the determinant is a unit precisely
when each of \(s,x,\det D_1\) is a unit. The minus sign in the
determinant is removed together with the Fourier minus sign by
\(x\mapsto-x\). The evaluated test at \(l\) is therefore
\[
 Q^{-ac}D_\chi(s)\widehat D_\chi(t)\phi'_1,\qquad
 \widehat D_\chi(t)=
       \int_{\mathcal O^\times}\chi(x)^{-1}\psi_v(tx)\,dx .
\]
The complementary place has character \(\chi^{-1}\), so the Mellin
weight is \(\chi(s)\chi(t)^{-1}|st|^{k+1}\).
Integration in \(s\) gives one. The Fourier transform vanishes
outside \(v(t)=-c\): below this shell additive integration
cancels on cosets modulo \(\pi^c\), and above it multiplicative
translation by a unit trivial on the additive factor but
nontrivial under \(\chi\) forces zero. At \(t=\pi^{-c}u\) its
value is \(\chi(u)Q^{-c}G(\chi,\psi_v)\) by finite residue
summation. Integration on that multiplicative shell therefore
gives \(\chi(\pi)^c Q^{ck}G(\chi,\psi_v)\).
The source compact averaging has no further effect on the chosen
test. The factors \(Q^{-ac}\) on the two sides cancel, and
Proposition~\ref{threal:localcone} extends the resulting scalar
identity to all matched upper tests. The nonvanishing of the
Gauss sum is Lemma~\ref{newtheta:gauss}.
\end{proof}

In the other descent Corollary~\ref{threal:ramA} gives the
corresponding factor \(\epsilon_v(1-k)\) with the same unit sign.

In this unprimed descent the integral line of \(l x\otimes y'\)
has coefficient ideal
\(\mathfrak b\,\mathfrak D_E^{-1}\bar{\mathfrak c}^{-1}\),
where \(\mathfrak D_E^{-1}\) is the absolute codifferent. Indeed
the trace-perfect target lattice has this ideal as its duality
range for the skew-Hermitian form; pairing \(y'\) against
\(\mathfrak c x'\) gives the indicated factor. Since \(E/F\)
is unramified at finite places the absolute codifferent extends
from \(F\), and its character value is one by
\(\nu|_{\mathbb A_F^\times}=1\). The two ideal exponents in the
unprimed branch are therefore \(+1,+1\), and the preceding
Fourier computation gives \(s=+1\) also here.

\begin{proposition}[The comparison at all relatively split good places]
\label{tpd:ordinary}
Let \(v\notin S\) split in \(E/F\), and choose either factor as
the first copy of \(K_0=F_v\). Use the integral-vacuum-normalized
local evaluations, compact mass one, and unit mass one on each
multiplicative group. Write \(z=k+1\) in the primed descent and
\(z=k\) in the unprimed descent. The lower quotient of the local
cone operator is the following scalar times the lower quotient of
the opposite zero evaluation, on every upper matched pair:
\[
   \nu_v(\mathfrak a_v)\,|\mathfrak a_v|_v^z L_{E,v}(\nu,z).
\]
Here \(\mathfrak a_v\) is the corresponding translation-line ideal
computed above, its modulus is that of an idele generating its two
local factors, and \(\nu_v\) is evaluated on such a generator.
The assertion includes \(\nu_w(\pi)=1\).
\end{proposition}
\begin{proof}
The source compact fixes the vacuum at this unramified split place:
in integral row bases its action is linear substitution with
unramified determinant character. In the two-plane evaluation the
translated vacuum maps to the smaller vacuum for
\(l\in\mathfrak a_v\) and to zero otherwise. To see the normalization,
use the two idempotent position and momentum groups, as in
Proposition~\ref{threal:ordinarysplit}. The vacuum is proportional
to the position lattice indicator, and its modulation lattice is
the character dual. Adaptedness decomposes both lattices into the
paired hyperbolic directions and the smaller tensor. Each exchanged
position entry is integrated by partial Fourier transformation at
zero. Translation in the nonintegrated entry leaves its indicator
nonzero precisely on that entry's lattice. Modulation on the
integrated entry gives its original integral on the character-dual
lattice and zero elsewhere. The two entries are distinct, so the
half-pairing vanishes also in residue characteristic two. Together
the two conditions say exactly that the displacement belongs to
the tensor lattice, equivalently \(l\in\mathfrak a_v\). All other
entries give the untranslated evaluation, normalized to the smaller
vacuum. This uses the character dual for the actual additive character,
including its codifferent. Integrating the two multiplicative shells
gives the displayed scalar. The opposite one-plane evaluation sends
vacuum to vacuum.

Let \(u=\nu_w(\pi)\) vary on the unit circle, keeping \(\beta_w\)
fixed and setting \(\alpha_w=\beta_w\nu_w\) with \(\nu_w\) unramified.
On the other split factor take the inverse characters. This family
is local. For \(u\) of infinite order the full source coinvariant
calculations of Propositions~\ref{split:kernel} and~\ref{split:good}
apply over \(K_0\) itself. Their proofs use the local orbital
functions and the infinite-order condition. The quotient-integral
and Jacquet argument of Proposition~\ref{threal:localcone} also uses
only those local facts and the split line decomposition. The smaller
zero evaluation is onto. Thus the two maps from the upper common
quotient to the smaller quotient are proportional for each such
\(u\). The vacuum calculation determines the scalar.

To extend the identity to \(u_0\) on the circle, realize the upper
and lower normalized induced targets by restriction of their functions
to the integral compact. The compact pictures are independent of
\(u\), since all varying inducing characters are unramified. Use
the orbital maps normalized to send both vacua to the spherical
vector of value one. For each fixed test they are continuous in
\(u\) on the circle. For the larger row map this follows by clearing
the determinant denominator as in Proposition~\ref{threal:ordinarysplit}.
The pivot-chart proof there works with the unramified value as a
variable: integration on a ball at zero gives a rational function
with denominator the indicated subproduct of
\(\prod_{j=0}^{n-1}(1-u Q^{j-d/2})\); a pivot of fixed valuation
contributes a Laurent monomial, and the other charts give the
lower-dimensional subproducts. Here \(Q\) is the residue size
and \(n\) the upper row rank. Clearing this denominator and normalizing
the lattice value gives Laurent polynomials in \(u\) up to a nonzero
fixed constant. A basis change with nonunit determinant introduces
a further Laurent monomial. The smaller row map is absolutely
convergent at its parameter and its spherical value has no zero or
pole on the circle. Use the corresponding dimensions and ranks in
the two lower maps. Compact-open invariance reduces their restrictions
to finitely many compact-picture evaluations at fixed level. Thus
continuity includes the normalized residue at a pole.

The zero evaluation is independent of \(u\). The local cone operator
also varies continuously on each test in a fixed finite-dimensional
span of smaller tests. The compact source action is independent of an
unramified change of character, and the restrictions to the two
displacement coordinates are compactly supported locally constant
functions with finitely many slices. On each compact open slice
away from zero the shell integrals are finite sums; near zero the
geometric series converge uniformly since \(z>0\). This applies
independently to the two local factors of \(E\).

For tests \(x,y\) with equal normalized upper images at \(u_0\)
denote the two upper maps by \(P_u,P'_u\). Choose a compact open
subgroup \(J\subset K_d\) fixing their images for every \(u\);
the upper compact actions on the tests are independent of \(u\).
The \(J\)-fixed space of the compact picture is finite-dimensional
and independent of \(u\). By surjectivity of the smaller-row map
\(P'_{u_0}\) from Proposition~\ref{threal:ordinarysplit}, choose
test preimages of a basis and average them over \(J\).
Surjectivity at the critical parameter uses the common irreducible
normalized unitary induction and the nonzero spherical image.
The matrix of \(P'_u\) on the span of these averaged preimages
stays invertible near \(u_0\); denote its inverse by \(H_u\).
The tests
\[
             y(u)=y+H_u(P_u x-P'_u y)
\]
are matched to \(x\) there, and converge to \(y\) in a fixed
finite-dimensional test space. For \(v=v_1\) the argument can
first be made on the full spaces and then restricted to supported
pairs. Apply the equality at infinite-order parameters tending
to \(u_0\) and pass to the limit on the two lower compact pictures.
The computed vacuum scalar is continuous as well. This only uses
unitary parameters.

Finally any allowed irreducible common quotient at this place is
isomorphic to the canonical one, and each map to it for trivial
source is unique up to scalar by the split local Howe theorem.
\cite[Theorem~1 and Corollary~6.3]{minguezHowe}
Equality of the two nonzero vacuum images forces the two scalar
changes to agree, in both the upper and lower models. Thus the
identity holds for every allowed choice of quotients.
\end{proof}

For comparison with the normalized recursion, denote the smaller
periods by \(I_1,I'_1\) and the primed, nonzero-cone test by
\(T\Phi'\). The comparison needed at the identity of the line
torus has three constants. With fixed models and measures they are
specified by the following identities, in which \(b,b'\) denote
the upper components with character \(\beta|\cdot|_E^{a/2}\):
\[
 b=C_0 I_1(e_0\Phi),\qquad b'=C_1 I'_1(T\Phi').
\]
The third requirement is that the following pair be matched on the
smaller spaces, with its second entry supported:
\[
 \big(C_2 F_X(-k)L^{S_0}(\nu,k+1)e_0\Phi,\ T\Phi'\big).
\]
If these identities hold and the smaller scalar is \(c_1\), then
\[
      C_0 c_1 b'=C_2 C_1 F_X(-k)L^{S_0}(\nu,k+1)b .
\]
Proposition~\ref{threal:endB} gives detection of the functions.
The upper scalar under the third requirement is therefore
\[
 c=\frac{c_1}{C F_X(-k)L^{S_0}(\nu,k+1)},\qquad
                  C=\frac{C_2C_1}{C_0}.
\]
Proposition~\ref{tpd:unfold} supplies the second identity for every
\(q>0\). For \(p>0\) on the unprimed side,
Proposition~\ref{threal:coneA} places the entire nonzero orbit in
the other character; the zero orbit gives the first identity by
evaluation, with a convergent smaller unprimed period by the
estimates of Proposition~\ref{threal:decay} with \(d\) decreased
by two. Thus \(C_0=1\) throughout for coherent mixed models. The
remaining test comparison is the third identity on the two smaller
matching quotients with one common normalization.
Propositions~\ref{tpd:ram} and~\ref{tpd:ordinary} give it at the
ramified and all relatively split unramified places, including the
unit sign \(s=+1\) and the factors of the two adapted line ideals.
The vacuum calculation and the partial Fourier transform supply the input for Propositions~\ref{tcr:descend} and \ref{tcr:unique}, which extend the inert identity to all smooth tests and compare the real operators on the polynomial module. Theorem~\ref{tcr:global} assembles their normalized tensor product.
\src{src/Development/WeilBound/ThetaLifts/LocalTheory/StandardValuations/ExistsVol.lean}{162}{Evaluation on a ramified matching pair};
\src{src/Development/WeilBound/ThetaLifts/LocalTheory/LocBOf.lean}{17}{Passage from these tests to local matching}.
The required \(C\) is common to the admissible characters and
the choices of finite matching maps.
\src{src/Development/AbelianVarietyTheory/Definitions/Divisors/CTSmallCone.lean}{239}{The two constant-term evaluations};
\src{src/Development/WeilBound/ThetaLifts/ConeHypotheses/KeyAOf.lean}{37}{Assembly of the three constants};
\src{src/Development/WeilBound/ThetaLifts/ConeHypotheses/PolDense.lean}{46}{The supported smaller matching identity}.

\paragraph{Removing a hyperbolic plane on the given lattice.}
The integral part of the change of model admits the following direct
construction. It will allow the local descent to use a rational plane
while retaining the lattice at every finite place.

\begin{proposition}\label{famtheta:plane}
Let \(E/F\) be the CM extension in the theta tower, unramified at
every finite place. Let \(L\) be a lattice in an \(E\)-space of
dimension \(n\ge3\), with nondegenerate Hermitian or skew-Hermitian
form \(b\), linear in its first argument. Suppose \(L\) is self-dual
for an invertible conjugation-stable fractional ideal \(\mathfrak d\)
and the form has both signs at each infinite place (after multiplication
by an imaginary scalar in the skew case). There are rational isotropic
vectors \(x,y\) with \(b(x,y)=1\) for which
\[
 L=(L\cap Ex)\ \oplus\ (L\cap Ey)\ \perp\ L_0 .
\]
Here the first two lines are paired perfectly with values in
\(\mathfrak d\), and \(L_0\) is self-dual for \(\mathfrak d\).
At any prescribed finite set of split places where
\(\mathfrak d_w=\mathcal O_{E,w}\) on both sides, the two vectors
may simultaneously be taken as local generators of their lines.
\end{proposition}
\begin{proof}
In the skew case multiply \(b\) temporarily by a nonzero element
negated by conjugation. The associated \(F\)-quadratic form
\(v\mapsto b(v,v)\) for the resulting Hermitian space is nondegenerate
of dimension \(2n\). It is indefinite at every real place and isotropic
at every finite place, since quadratic forms over a finite extension
of \(\mathbf Q_p\) of dimension at least five are isotropic (also for
\(p=2\)). Hasse--Minkowski supplies a rational isotropic \(E\)-line
\(U\) for the original form. The intersection \(I=L\cap U\) is
saturated and of rank one. Both \(I\) and \(L/I\) are torsion-free
finite modules over the Dedekind ring \(\mathcal O_E\), hence projective.
In particular \(I\) is a direct summand as a module.

Put \(J=\operatorname{Hom}_{\mathcal O_E}(\bar I,\mathfrak d)\),
where the bar denotes the conjugate module. Self-duality identifies
\(L\) with \(\operatorname{Hom}(\bar L,\mathfrak d)\) by the pairing.
Restriction to \(\bar I\) is onto by the direct summand property.
Choose an integral section \(s:J\to L\) of this restriction,
using projectivity of \(J\). The pairing of \(sJ\) and \(I\)
is perfect in both directions, by biduality of the invertible modules
and conjugation stability of \(\mathfrak d\).
Choose \(\theta\in\mathcal O_E\) of trace one. Indeed the trace ideal
is the unit ideal locally at every finite prime: at a split prime
this is addition, and in an unramified quadratic extension it follows
by reduction from the surjective residue-field trace. This argument
also covers residue characteristic two.
There is a unique linear map \(C:J\to I\) with
\[
 b(Cj,sj')=\theta\, b(sj,sj')\qquad (j,j'\in J)
\]
by the perfect pairing. Hermitian symmetry with its given sign implies
\(b(sj,Cj')=\bar\theta b(sj,sj')\). Since \(b(I,I)=0\), the map
\(s'=s-C\) is another section and \(b(s'J,s'J)=0\).
The Gram pairing on \(I\oplus s'J\) has zero diagonal blocks and
perfect off-diagonal blocks. It is self-dual for \(\mathfrak d\).
For any \(z\in L\), its pairing functional on this submodule is thus
realized by a unique vector inside it. Subtracting that vector gives
the orthogonal projection to \(L_0\). This proves the integral
decomposition and self-duality of the complement. The two summands
of the plane are precisely the full lattice intersections with their
rational spans.

Choose a rational generator \(x\) of \(I\otimes E\); by weak
approximation scale it to generate \(I\) at all the indicated primes.
There is a unique \(y\in(s'J)\otimes E\) with \(b(x,y)=1\). Perfection
and \(\mathfrak d_w=\mathcal O_{E,w}\) make it a generator at the
same places, including the conjugate primes. Complete these two vectors
by bases of \(L_0\) at the split places and put them in the designated
columns. At each real place the plane has signature \((1,1)\);
choosing a Cartan involution on it and on its orthogonal complement
gives the adapted marking. The lower signatures lose one of each sign.
\end{proof}
\src{src/Development/WeilBound/ThetaLifts/TowerModels/PlaneMain.lean}{47}{rational isotropy and integral trace one},
\src{src/Development/WeilBound/ThetaLifts/TowerModels/PlaneSkew.lean}{18}{skew form and codifferent}, and
\src{src/Development/WeilBound/ThetaLifts/TowerModels/TowerCodimensionTwoOf.lean}{32}{marked planes for the two towers}.

\subsection{Local descent on an adapted pair of theta towers}
\label{tcr:main}

Write \(d=a+a'\), \(k=(a-a')/2\), and \(s=p+q\). In the
\(A\)-step take \(q=0<p\), remove a plane from \(V,W\), and put
\(z=k,\ \xi=\alpha|\cdot|_E^{a'/2}\). In the \(B\)-step take
\(q>0\), remove a plane from \(V',W\), and put
\(z=k+1,\ \xi=\beta|\cdot|_E^{a/2}\).
The one-plane evaluation \(e_0\) belongs to the unchanged source
and the two-plane cone operator \(T\) to the changed source. Subscripts
on these maps denote their local factors. The adaptedness assumptions
are those of \S\ref{newtheta:global}: the two integral lattices
split into their two lines and their orthogonal complements, the
finite designated bases have the prescribed columns, and the Cartan
involutions pair the two lines and restrict to the lower involutions.

\begin{lemma}[Square integrability and kernels]\label{tcr:kernels}
Each upper polynomial period, with the stipulated primed support,
belongs to \(L^2([U(W)])\). For any of the matching maps, the period
on each side factors through the restricted tensor product of its
small local theta quotients. A nonzero such period map is injective
on that tensor product.
\end{lemma}
\begin{proof}
The periods are smooth, of uniform moderate growth, finite under the
enveloping center and compact-finite at infinity, as in the proof
of Proposition~\ref{threal:detectline}. Apply the cuspidal-exponent
criterion for square integrability, with the constant term divided
by the square root of the modulus of its parabolic
\cite[I.4.11]{moeglinWaldspurger}. The center of the unitary group
has compact adelic quotient.

Consider a cuspidal exponent on a Levi of a flag of isotropic
subspaces. If the first subspace has dimension \(r>1\), its first
Levi block is \(\mathrm{GL}_r(\mathbb A_E)\). Any cuspidal projection
there is generic and unitary after a central real twist. Here a
cuspidal projection is computed by pairing the flag constant term
on \(M(F)\backslash M(\mathbb A_F)^1\), where all rational
character moduli are one, with a rapidly decreasing cuspidal test
on the Levi \(M\). Uniform moderate growth suffices for convergence;
the central height variables remain free for extracting exponents.
Fix the other cuspidal factors and vary the first cusp vector. A
nonzero pairing gives an equivariant functional with values in
the smooth contragredient of that cuspidal factor, using finite
compact invariance of the local test. An action off the norm-one
subgroup is accounted for by a compensating split-central
translation, hence by the extracted central exponent.
At a completely
split unramified prime with \(\nu_w(\pi)\) of infinite order it has
to be a subquotient of the first block in the Jacquet module of the
upper small quotient. Indeed the full flag constant term can be taken
first along the \(r\)-plane and then in its remaining unitary factor.
Integration there and cuspidal projection commute with the first
block action. The local calculation in Proposition~\ref{threal:detectline}
rules out precisely these subquotients for \(r>1\). This argument
uses the full split source-coinvariant calculation at that one
prime and does not require the period to be in \(L^2\).

If \(r=1\), the two line components are smaller periods by
Propositions~\ref{threal:coneA}, \ref{tpd:unfold} and
Lemma~\ref{tpd:so}, including their zero-orbit formulas. A missing
source plane gives no corresponding nonzero orbit. The same
calculation applies to the remaining flag on each contributing
smaller period. For this calculation rational hyperbolic coordinates
suffice: at the places where the resulting bases are not integral
they change finitely many local tests. At infinity the estimates of
Proposition~\ref{threal:decay} apply to Schwartz tests and their
derivatives. Center-finiteness of the successive coefficient as an
automorphic function follows also by Harish-Chandra projection from
the preceding coefficient. Thus every general linear block of a
contributing cuspidal exponent is a line. If there are \(j\) such
blocks, write \(T_i=|t_i|_E\) for their parameters in flag order.
The real exponents \(b_i\) of the unnormalized term satisfy
\(b_i\le a/2\), since at every stage they are half the rank of one
of the two current sources. The square root of the modulus has
exponents \((d-2i+1)/2\): the line radical for a space of dimension
\(d-2i+2\) has determinant modulus \(T_i^{d-2i+1}\) (its middle
module has dimension \(d-2i\) over \(E\), and its central
\(F\)-line has modulus \(T_i\)). Hence
\[
        b_i-(d-2i+1)/2\ \le\ -(a'-2i+1)/2<0
        \qquad(1\le i\le j\le s).
\]
Every nonzero generating ray of the closed positive chamber
\(\log T_1\ge\cdots\ge\log T_j\ge0\) has strictly negative
normalized exponent. The criterion gives the assertion.

Fix the other local tests. Source invariance factors the resulting
map into \(L^2\) through the big theta module of the trivial
representation, which has finite length and a unique simple quotient
with multiplicity one in its cosocle
\cite{howeTranscending,ganTakedaHowe,minguezHowe}. An admissible
finite-length image in a unitary representation is semisimple.
For a submodule take its Hilbert closure: the projection to a
compact-open fixed space at a finite place or a compact isotypic
space at a real place shows that the closure intersects the
algebraic image in the submodule. These spaces are finite-dimensional,
so orthogonal complements taken on them give an algebraic splitting;
at infinity it is stable under derivatives by skew adjointness on
smooth vectors. Thus each local map kills the kernel of the prescribed
small quotient. At the support place \(v_1\) on the primed side one
uses the supported coinvariants, which map isomorphically to this
same small quotient by Proposition~\ref{split:kernel}. Applying
the universal property of a tensor product to finite sums and then
enlarging the finite set of places gives the factorization.

The kernel of the map on the resulting restricted product is stable
under independent finite-place translations and the real Lie and
compact actions. A submodule containing a nonzero vector contains
the full finite tensor product supporting that vector: the individual
simple factors have scalar endomorphisms, and the algebraic density
theorem gives simplicity of their finite external product. Enlarging
the finite set, using the nonzero spherical vectors at the new
places, gives the whole restricted product. This proves injectivity
whenever the period map is nonzero.
\end{proof}

\begin{lemma}[Nonzero smaller images at infinity]\label{tcr:bump}
Suppose the smaller period on the changed-source side is nonzero
on polynomial tests. Then \(P_{1,\infty}T_\infty\) is nonzero on
upper polynomial tests, where \(P_{1,\infty}\) is its real matching
map. For the other side \(P'_{1,\infty}e_{0,\infty}\) is nonzero
whenever its smaller real quotient is nonzero.
\end{lemma}
\begin{proof}
Choose smaller finite tests and a target argument at which the
smaller period is a nonzero functional \(D\) of the smaller real
test. It is continuous for the Schwartz topology by
Corollary~\ref{threal:limits} and source-invariant. Its restriction
to polynomial tests factors through the smaller real quotient by
Lemma~\ref{tcr:kernels}. For \(e_0\), multiply any smaller Fock
polynomial by the vacuum on the complementary summand, whose
evaluation is nonzero. The positive complex structure preserves
this decomposition by adaptedness. Thus the one-plane evaluation
is onto on polynomials.

For \(T\), use the mixed coordinates with modulation space
\(Y\otimes Ex'\oplus Ex\otimes W_1\), where \(Y\) is the changed
source. Before the second evaluation the first position is
\(v\in Y_\infty\). For a real compact source element \(h\) denote
the evaluation at \(v=u x\), followed by the second evaluation
and \(D\), by \(A_h(u)\), where \(u\in\prod_{w\mid\infty}\mathbb C\).
In these coordinates multiplication of the initial test by a
scalar test \(b(v)\) multiplies \(A_h(u)\) by \(b(h^{-1}ux)\).
Indeed the source acts on the first position by pullback and on
the residual oscillator by its Weil action. With \(h=1,u=1\),
the evaluation onto the smaller Schwartz space is surjective
by extension in the omitted linear coordinates.

In the \(A\)-step the scalar under the integral for \(DT_\infty\)
is \(A_h(u)\) times the everywhere nonzero Mellin density on
\((\mathbb C^\times)^{[F:\mathbb Q]}\). In the \(B\)-step one
instead pairs the one-coordinate Fourier transform of \(A_h\)
at each place with the Mellin density of exponent \(k+1\), by
the partial Fourier identity in Proposition~\ref{tpd:unfold}.
We describe the transformed density on \(u\ne0\). Relative to
additive area on \(\mathbb C\), the original density is a nonzero
constant times
\[
          (l/|l|)^{m_w}|l|^{2k},\qquad |m_w|=m_{\alpha,w}-m_{\beta,w}>d .
\]
Complex conjugation of the displacement or reversal of the Fourier
pairing can reverse the angular weight. The Fourier transform of
this tempered distribution has a smooth nowhere-zero density on
\(\mathbb C^\times\). Indeed it is equivariant under rotations and
positive dilations. On their free transitive orbit such a distribution
is a scalar times the corresponding character times Haar density:
divide by the character and use translation invariance in angle and
logarithmic radius. If its scalar vanished it would be supported
at the origin, hence be a finite sum of derivatives of the delta
distribution. Its inverse Fourier transform would then be a
polynomial of homogeneous degree \(2k\) and angular weight of
absolute value \(>2k\), which is impossible by expanding in
\(l,\bar l\). This argument applies separately at every real place.

Here is why integration over the compact directions cannot cancel
this density. In the first mixed coordinates
\((h,u)\mapsto h^{-1}ux\) maps the compact group times the
nonzero scalars onto the product of punctured null cones. At each
place the two compact unitary groups are transitive on the two
unit spheres, and the lengths of the two components of a null
vector are equal. The fibers are the compact line stabilizers,
acting by \((h,u)\mapsto(jh,\eta u)\) for \(jx=\eta x\).
The factor acting on the smaller oscillator is the smaller compact
source and disappears after \(D\). The scalar on this stabilizer
also cancels: for the unchanged radial measure the source determinant
factor on the line is \(\nu^{-1}(\eta)\) in the \(A\)-step and
\(\nu(\eta)\) in the \(B\)-step, as in the source Levi calculations
of Proposition~\ref{threal:coneA} and Lemma~\ref{tpd:mellin}.
In the first chart the weight on \(u\) is \(\nu(u)\) in the
\(A\)-step. In the other step the Fourier transform carries a
rotation of \(u\) to its dual rotation, so the transformed angular
weight is \(\nu^{-1}(u)\). Explicitly the \(B\)-chart functions
after applying \(D\) satisfy
\[
        B_{jh}(l)=\nu(\eta)B_h(\eta l),\qquad |\eta|=1.
\]
The argument \(\eta l\) follows because the displacement is
\(\bar l\,y\otimes x'\): the Fourier kernel pairs \(u\) with \(l\)
by a real trace of a fixed nonzero multiple of \(ul\).
Substitution in the defining Mellin integral with weight \(\nu(l)\)
cancels the factor \(\nu(\eta)\). This proves invariance of the integrated density
along the compact fibers in either chart.

Choose first a test with \(A_1(1)\ne0\), then multiply it by a
nonnegative bump \(b(v)\) supported sufficiently near \(x\) on
each punctured cone and positive at \(x\). The pullback of that
support has all \(|u_w|\) bounded away from zero and infinity.
The Fourier distribution may therefore be evaluated there using
its smooth density. After compact-fiber integration the remaining
integrand, apart from \(b\), is continuous and nonzero at \(x\);
in a small neighborhood a fixed phase makes its real part positive.
This proves \(DT_\infty\ne0\) on Schwartz tests. Compact averaging,
Fourier transformation, restriction and the absolutely convergent
Mellin moment are continuous in the Schwartz seminorms, as in
Lemma~\ref{tpd:mellin}. Hermite expansions approximate any
Schwartz test by polynomials times the vacuum in these seminorms.
The output for an upper polynomial is again lower polynomial:
decompose it over the smaller positive summand and its orthogonal
complement; the compact orbit spans a finite-dimensional polynomial
space and all evaluation and displacement act in the complement.
Consequently some upper polynomial has \(DT_\infty\ne0\), as claimed.
\end{proof}

\begin{proposition}[Descending the test maps]\label{tcr:descend}
Suppose both smaller period maps are nonzero. After application of
the smaller matching map, each of \(T_v,e_{0,v}\) factors through
the upper matching map, at every inert place and at the product of
the real places (on polynomial vectors there). Both resulting maps
intertwine the target line parabolic with
\(\xi\boxtimes\pi_1\) inflated trivially on its radical.
At infinity this means its Lie algebra and the compact part of its
Levi. The two real maps so obtained are nonzero.
\end{proposition}
\begin{proof}
First produce auxiliary upper tests at every other place for each
of the two operations, such that their smaller images are nonzero.
At infinity use Lemma~\ref{tcr:bump}. At a ramified place use the
standard matched pair of Proposition~\ref{tpd:ram} or
Corollary~\ref{threal:ramA}. At an unramified split place use the
vacuum in Proposition~\ref{tpd:ordinary}. At \(v_1\), if a
supported primed test is required, use the surjection of supported
tests onto the whole upper quotient. All the indicated local maps
at that split place already factor through the upper quotient
there by those propositions (one may compare with the zero test on
the opposite side). In the \(A\)-step the smaller unchanged primed
source is compact at infinity and needs no support for integration.
When the smaller supported space is required for matching, one
can also use the representatives of Lemma~\ref{tpd:so}: the same
matrix proof of its lifting assertion works for \(q=0\).
In the \(B\)-step the changed-source operator itself preserves the
full-rank condition by Lemma~\ref{tpd:mellin}.

At an inert place the following vacuum calculation requires no
factorization. The tensor lattice at the place is self-dual for
the additive-character symplectic pairing. The prime is nondyadic.
In the two-plane symplectic decomposition its summands in the
removed tensor blocks have the following description. Each isotropic
line tensor a residual space pairs perfectly with the opposite
isotropic line tensor that space. The four products of the two pairs
of lines split into two such dual pairs. Adaptedness and self-duality
make these pairs of lattices mutually dual, including any fractional
line ideals. Take the removed modulation space and these opposite
summands as its dual position space. In those positions the basic
vector is the lattice indicator times the smaller basic vector,
up to scalar: invariance under the modulation lattice imposes
support on its position dual and invariance under the full dual
translation lattice forces constancy there. The tensor factor on
the residual space is itself the smaller basic line by the same
Heisenberg characterization. This argument also applies to the
one-plane evaluation. Normalize each evaluation on the vacuum.
The source compact fixes the basic vector, by the unramified
oscillator vacuum calculation in Proposition~\ref{threal:allinert}.
The displacement is a single position coordinate in the respective
chart. Thus shifted evaluation gives precisely the smaller basic
vector times \(1_{\mathfrak a_v}(l)\), where the integral translation
ideal \(\mathfrak a_v\) is the localization of
\[
 \mathfrak b\,\mathfrak D_E^{-1}\bar{\mathfrak c}^{-1}
       \quad(A),\qquad \mathfrak b^{-1}\bar{\mathfrak c}\quad(B).
\]
With multiplicative unit volume one its Mellin integral is
\[
  u_v=\nu_v(b_v)|b_v|_{E_v}^{z}
           (1-\nu_v(\pi_v)|\pi_v|_{E_v}^{z})^{-1},
        \qquad b_v\mathcal O_{E_v}=\mathfrak a_v .
\]
This is nonzero, including when \(\nu_v(\pi_v)=1\). All but
finitely many \(b_v\) are units, and the product of these basic
factors converges absolutely to a nonzero number since \(z>1\).

Now put a vector in the kernel of an upper matching map at the
place under consideration and the chosen auxiliary tests everywhere
else. Lemma~\ref{tcr:kernels} kills its upper period identically.
In particular its \(\xi\) coefficient is zero. This coefficient
is a fixed nonzero Haar multiple of the smaller period on the
tensor of the transformed tests, by the constant-term formulas.
For the primed cone formula for a general upper test, first take
the Levi difference of Proposition~\ref{tpd:unfold}; on the
\(\beta\)-coefficient it acts by the difference of the two
characters at the chosen idele, a nonzero scalar, and on \(T\)
it acts by that same scalar. For a zero evaluation with supported
representative, make the correction at \(v_1\) before reading the
other local factor. By smaller injectivity
(Lemma~\ref{tcr:kernels}) the tensor of lower images vanishes.
All its other factors are nonzero, so the one at the place under
consideration vanishes. This proves the claimed factorizations.

The target equivariance can be read from the same formulas. For
a local element of the line Levi or radical take its right translate
in the upper constant term. At the identity the radical acts
trivially by left invariance, the line acts by \(\xi\), and the
smaller target group translates the smaller period. Subtract the
corresponding two smaller expressions and again cancel the fixed
nonzero auxiliary factors by smaller injectivity. At infinity
one takes compact Levi translates and differentiates the identity
for one-parameter subgroups at zero; these operations preserve the
polynomial domain in the required differentiable sense. Continuity
for passing the derivative through the periods follows from
Proposition~\ref{threal:decay}. For a one-parameter element outside
the compact the identity before differentiation holds for Schwartz
tests by the same absolutely convergent integrals. The nonzero
real images follow from Lemma~\ref{tcr:bump}.
\end{proof}

\begin{proposition}[Inert and real uniqueness]\label{tcr:unique}
Under the assumptions of Proposition~\ref{tcr:descend}, the two
descended maps at an inert place satisfy \(D_{T,v}=u_vD_{0,v}\).
At infinity \(D_{T,\infty}=c_\infty D_{0,\infty}\) for a unique
nonzero scalar. For fixed real models, upper and lower real
matching maps and measures, this number is independent of every
finite character, matching quotient and choice of admissible
Weil actions in the two towers.
\end{proposition}
\begin{proof}
At an inert place the upper spherical vector generates its
irreducible representation. The line parabolic and the lattice
stabilizer satisfy \(G=PK\) by the integral Witt argument in
Proposition~\ref{threal:coneA}. Hence the spherical vector generates
even under \(P\). The comparison of its two images in the preceding
proof gives the identity on the whole representation.

Here is the compact weight calculation at one real place.
Write \(R=d-s,S=s,P=a-p,Q=p,A=m_\alpha,B=m_\beta\), and
\(\Delta=A-B>d\). Put
\[
 c_\pm=\tfrac12(B\pm(P-Q)),\qquad
 h_+=c_++\tfrac12(\Delta-R+S),\quad
 h_-=c_-+\tfrac12(\Delta+R-S).
\]
The bottom \(U(R)\times U(S)\)-type of
Proposition~\ref{threal:allreal} has weights
\[
 ((h_+)^P,(c_+)^{R-P};(h_-)^Q,(c_-)^{S-Q}),
\]
where superscripts denote repetitions. To see the normalization,
start from the two target compact vacuum weights \(\pm(P-Q)/2\);
the source determinant in the joint harmonics adds
\((\Delta-R+S)/2\) in the first \(P\) positions and
\((\Delta+R-S)/2\) in the first \(Q\) positions, and the target
splitting adds \(B/2\) throughout.
Alternatively the bottom-layer formula for the ordered two-block
parabolic is its inducing differential plus
\(\rho(\mathfrak u\cap\mathfrak p)-\rho(\mathfrak u\cap\mathfrak k)\)
in Paul--Trapa's cover normalization: in the first positive block
this correction is \(((S-Q)-(R-P))/2\), in the first negative
block its negative, and in the remaining positive and negative
blocks it is respectively \((P-Q)/2\) and its negative.
The inducing parameter is in the good range by \(\Delta>d\);
the bottom type occurs once by the bottom-layer formula for
cohomological induction (the inducing Levi type is a character).
In ordinary, uncovered normalization one subtracts \(\rho(\mathfrak u)\)
from the inducing differential. The two character exponents then
become \((A-a')/2,(B+a)/2\); their difference is
\((\Delta-d)/2>0\), so the bottom-type theorem applies.
\cite[\S2 and Theorem~4.1]{paulTrapaHowe}
\cite[Theorem~5.3 and Proposition~6.1]{voganZuckerman}

The compact Levi on an isotropic line acts by the same circle
\(\eta\) in one positive and one negative coordinate, and by
\(U(R-1)\times U(S-1)\) on the other coordinates. Indeed a null
vector is a sum of two vectors of equal norm, and an element
carrying it to \(\eta\) times itself acts by \(\eta\) on each.
For a \(U(n)\) weight \(\lambda_1\ge\cdots\ge\lambda_n\) the
\(U(n-1)\)-weights interlace it with multiplicity one, and their
omitted circle weight is \(\sum\lambda_i-\sum\mu_i\). Its maximum
\(\lambda_1\) forces \(\mu_i=\lambda_{i+1}\) for all \(i\);
its minimum \(\lambda_n\) forces \(\mu_i=\lambda_i\).
In the \(A\)-step both high blocks occur and \(h_++h_-=A\) is
the maximum diagonal-circle weight. In the \(B\)-step both low
blocks occur and \(c_++c_-=B\) is the minimum. These are exactly
the respective angular characters of \(\xi\). The extreme block
in the upper bottom type thus consists of a single irreducible
type of the smaller compact, with one entry deleted in each sign.
It is the lower bottom type by the same weights formula:
\(P-Q\) and \(R-S\) do not change. It occurs once in \(\pi_1\).
The compact restrictions of the two parabolic maps on the upper
bottom type are therefore proportional. The Iwasawa Lie algebra
identity \(\mathfrak g=\mathfrak p+\mathfrak k\) and the
Poincar\'e--Birkhoff--Witt theorem show that \(U(\mathfrak p_{\mathbb C})\)
acting on this compact type generates the entire irreducible module.
For the product of real places take the exterior tensor product
of the bottom types and apply the same argument componentwise.
This proves dimension at most one for the required space of maps.
Both are nonzero, so the scalar is unique and nonzero.

For fixed base data the infinite characters are fixed by their
specified exponents \(m_\alpha,m_\beta\). The doubled Siegel law
fixes the resulting real oscillator actions: the ratio of two
Heisenberg implementations with the same doubled Levi law is a
character trivial on that Levi and on the root unipotents, hence
trivial on the doubled group, as in Proposition~\ref{threal:inert}.
The compact integral defining \(T_\infty\) and its Mellin weight
therefore agree for all the allowed systems. The fixed realizations
of the upper and lower real maps give the same proportionality
equation on every polynomially matched real pair in all these
systems. Since some real test has nonzero zero-evaluation image
and its upper image can be matched on the other side, this equation
determines the same number \(c_\infty\) for all systems. If
there are no such systems the uniform comparison is vacuous.
\end{proof}

\begin{theorem}[The scalar for the two undoubled periods]\label{tcr:global}
For every triple of \S\ref{newtheta:global} and every allowed set
of real and finite matching maps, the primed period is nonzero
on some supported test, and there is a unique \(c\ne0\) such that
\(I(\Phi)=cI'(\Phi')\) on all matched pairs. For given adapted upper
and lower geometries and given upper and lower real matching maps,
there is \(C\ne0\) independent of all admissible \(X\), Weil actions
and finite matching maps such that
\[
 \begin{array}{ll}
 c=C\,F_X(1-k)L^{S_0}(\nu,k)c_1,
       &(q=0<p;\ e_1=e_2=+1,\ s_{\rm unit}=+1),\\[2pt]
 c\,L^{S_0}(\nu,k+1)=c_1/(C F_X(-k)),
       &(q>0;\ e_1=e_2=-1,\ s_{\rm unit}=+1).
 \end{array}
\]
Here \(F_X\) is the finite factor of \S\ref{newtheta:global},
and \(c_1\) is the scalar for the specified smaller matching maps.
\end{theorem}
\begin{proof}
First work on an adapted pair of geometries and suppose the
assertions of nonvanishing and proportionality hold in the smaller
rank. In particular both smaller period maps are nonzero: supported
primed vectors are matched by surjectivity of the opposite quotient
and the distinguished lower full-rank quotient. Use polynomial
inputs initially; nonzero polynomial inputs are available by
Corollary~\ref{threal:limits} and density. For \(T\) use the
integral with multiplicative units of volume one and local source
compacts of volume one. Normalize finite mixed evaluations on
their vacua. Rational coherent evaluations differ from this system
by fixed scalars with almost every scalar one
(Proposition~\ref{threal:coneA}). At infinity keep the real
evaluations fixed.

On upper matched tests, the restricted tensor product of lower
images of \(T\) is a scalar times that of \(e_0\). At the real
places use Proposition~\ref{tcr:unique}; at the finite places use
Propositions~\ref{tpd:ordinary}, \ref{tpd:ram},
Corollary~\ref{threal:ramA} and Proposition~\ref{tcr:unique}.
For a finite tensor sum this follows by applying the tensor product
of the descended linear maps to the equality of upper tensors in
the definition of polynomial matching. Outside the finite test set
one must retain the product of the vacuum multipliers: these
converge absolutely and have nonzero product. The translation
ideal contributes \(N(\mathfrak a)^{-z}\nu(\mathfrak a)\).
It is a unit at the ramified factors. As computed after
Lemma~\ref{tpd:mellin} and Proposition~\ref{tpd:ram}, its
character value is \(\nu(\mathfrak b)\nu(\mathfrak c)\) in \(A\)
and its inverse in \(B\). The partial Euler factors at the primes
of \(Q\) are one since \(\nu\) is ramified there. The ramified
factors are the stated Gauss factors, with arguments \(1-k,-k\)
respectively and with unit sign \(+1\) in both cases.
All these factors are nonzero. The remaining factors consist of
\(c_\infty\), the absolute ideal norm power, the fixed changes of
evaluation, and the quotient of the two Haar normalizations in
the orbit formulas. They give a single \(C\ne0\) for the stated
fixed data by Proposition~\ref{tcr:unique}. The Haar measures
are part of the models; no equality of Haar masses in different
dimensions is needed.

Use the smaller identity on these two tensors. In \(A\), the
unprimed \(\xi\)-coefficient is the cone and the primed one is the
zero evaluation, giving the first formula. In \(B\), the primed
coefficient is the cone and the unprimed one is the zero evaluation,
giving the reciprocal in the second formula.
Propositions~\ref{threal:endA} and~\ref{threal:endB} show that
the resulting scalar compares the upper periods at every target
argument. The zero-coefficient side is nonzero for some upper
test: choose its real polynomial and finite factors as in
Proposition~\ref{tcr:descend}, then apply smaller injectivity to
their nonzero tensor. Choose a matching vector on the other side,
using the supported surjection at \(v_1\) when the other side is
primed. The comparison just proved gives nonzero upper periods
on both tests. This also proves uniqueness of \(c\). Passage to
every matched limit is Corollary~\ref{threal:limits}.

To apply this argument by induction on \(p+q\), start at
Theorem~\ref{threal:base}. For the induction step the integral
plane construction of Proposition~\ref{famtheta:plane} supplies
the hyperbolic summands with their ideal lattices in the changing
source and target. On a given rational hyperbolic plane choose
Cartan involutions interchanging the two null lines at infinity
and an arbitrary compatible Cartan on its complement. At the
designated split places take the new adapted integral bases in the
prescribed column order. These choices do not change the rational
spaces, the lattices, the Haar measures or their theta integrals.
The smaller tensor has a rational symplectic basis by symplectic
elimination, giving its usual rational Schr\"odinger model.
Its basic vector is the lattice Heisenberg eigenline. It can also
be obtained by partial evaluation of the upper basic vector in the
adapted coordinates above: the removed mutually dual lattices
give the position indicator in those coordinates, leaving the
smaller eigenline. At a split dyadic place one uses the two isotropic
idempotent halves as in Proposition~\ref{tpd:ordinary}; their
individual quadratic phases vanish. In rational coordinates the
basic lattice gives the ordinary integral indicator almost everywhere.
Use any fixed lower source Haar measures. For the two lower
oscillators use the same pairs of splitting characters
\((\beta,\beta/\alpha)\), \((\alpha,\alpha/\beta)\).
Subtracting two from either rank preserves the quadratic character
restriction required for each unitary splitting. Use the local
splitting and its doubling law
\cite[(1.16) and Proposition~2.2]{harrisKudlaSweet}, and their
adelic product for the two Hecke characters
\cite[\S\S2--3]{wuThetaTower}. On rational elements the character
factors, absolute determinant factors and Weil indices have product
one, giving the rational theta action through Poisson summation.
Equivalently the doubling in a rational diagonal polarization
uses the prescribed determinant character on its Siegel Levi
and evaluation at zero; restricting to the two isometric summands
recovers the chosen single Weil action. All smaller gap inequalities
hold since \(p+q\) decreases by one, \(a'\) decreases by at most
two and \(k\) decreases by at most one. The local matching results
in \S\ref{threal:main} apply at both levels. The induction
hypothesis and the preceding adapted argument now apply.

We finally transport the existence and uniqueness assertion to the
original markings. At infinity choose separately in each source
and commonly in the target real isometries conjugating the original
Cartan involutions to the adapted ones. Their oscillator operators
carry the Gaussian line and its creation derivatives onto the new
ones by Heisenberg covariance. On the period the source translation
disappears by Haar invariance and the common target translation
right translates both functions by the same element. Under this
transport the original and new pairs of real simple quotient maps
differ, after identifying their targets, by one nonzero scalar
for each side, by real Howe uniqueness. The same comparison at
the finitely many remarked split places follows by finite Howe
uniqueness. At all other places one can keep the original maps.
The support condition uses matrix rank in the polarization
specified by the chosen split embedding and is invariant under
invertible row and column changes. Therefore a pair matched with
the original maps is transported to an adapted matched pair after
rescaling one test by the product of these finitely many scalar
ratios. The two upper period functions are proportional also for
the original data, and a supported nonzero matched pair transports
back. This transport is needed only for the existence assertion
for an arbitrary marking. For the uniform recursion with fixed,
already adapted geometries use the direct calculation in the
first part of the proof. The upper period scalar may depend on
the chosen matching maps and on \(X\).
\end{proof}
\src{src/Development/WeilBound/ThetaLifts/LocalTheory/LocAOf.lean}{17}{The unprimed test comparison},
\src{src/Development/WeilBound/ThetaLifts/LocalTheory/LocBOf.lean}{17}{the primed test comparison}, and
\src{src/Development/WeilBound/ThetaLifts/LocalTheory/StepFacEq.lean}{85}{the quantifiers for a single descent constant}.


\subsection{The geometric degree share and its arithmetic cost}\label{newtheta:HS}

Fix one component \(Z\) of a good deep cover and write \(D\) for the
dimension of its flag space, \(T>D\) for the final weight index,
\(\delta\) for the positive geometric degree of its birational image,
\(h\) for its normalized Chow height, \(d_0\) for the rank of the
terminal cusp lattice, and \(g\) for its arithmetic degree before the
associated-bundle transfer. The coefficient field is totally complex.
The fraction \(\binom{T-1}{D}\delta/d_0\) has a geometric meaning.

\begin{lemma}[Independent columns on a transverse section]\label{newtheta:filtration}
Over an infinite splitting field \(L\), let \(y_1,\ldots,y_\delta\)
be nonzero homogeneous-coordinate vectors in \(L^N\) for the points
of a transverse section by \(D\ge1\) hyperplanes \(s_1,\ldots,s_D\).
Let the terminal section space have dimension \(d_0\), and let
\(U_0\supset\cdots\supset U_{T-D}\) be a filtration by subspaces of
it. Suppose that degree-zero sections surject by restriction onto
\(L^\delta\), and that for \(0\le j<T-D\) the symbol on \(U_j\)
has kernel \(U_{j+1}\) and entries indexed by pairs \((\gamma,p)\)
with \(\gamma\in\mathbb N^D,\ |\gamma|=j,\ 1\le p\le\delta\).
For a degree-zero section \(m\), a linear form \(l\), and each such
\(\gamma\), suppose \(s^\gamma l^{T-j}m\in U_j\) and its symbol at
\((\gamma',p)\) is
\(\mathbf1_{\gamma=\gamma'}l(y_p)^{T-j}m(y_p)\). Then
\[
 d_0\ \ge\ B:=\delta\binom{T-1}{D}.
\]
There are \(B\) independent terminal columns of this kind. On the
associated graded their determinant in the point and monomial
coordinates is, up to ordering of the bases,
\[
 \left(\prod_{p=1}^{\delta}l(y_p)\right)^{
       \sum_{j=0}^{T-D-1}\binom{j+D-1}{D-1}(T-j)} .
\]
\end{lemma}
\begin{proof}
Choose \(l\) avoiding the finitely many proper hyperplanes
\(l(y_p)=0\), possible over an infinite field. Lift the standard
coordinate vectors of \(L^\delta\) to \(m_1,\ldots,m_\delta\).
For fixed \(j,\gamma,p\), the symbol of
\(s^\gamma l^{T-j}m_p\) has just one nonzero entry,
\(l(y_p)^{T-j}\). These symbols form a basis of the indicated
coordinate space. A relation between all the chosen columns, inspected
at the least \(j\) with a nonzero coefficient, would therefore give
a nontrivial relation between these independent symbols. The number
of \(\gamma\) at level \(j\) is \(\binom{j+D-1}{D-1}\), by placing
\(D-1\) separators in a row of \(j\) equal marks. Summing from \(0\)
to \(T-D-1\) gives \(\binom{T-1}{D}\) by Pascal's identity.
Multiplication of the displayed diagonal entries gives the determinant.
\end{proof}

This computation takes place over a field. In the integral problem one
needs the dimensions, determinants \emph{and} the norms of the columns.
In the original argument a nonzero polynomial \(Q\) in the \(DN\)
hyperplane coefficients selects transverse sections. On \(Q(s)\ne0\)
the specialized Chow form is a nonzero scalar times
\(\prod_p l(y_p)\) after passing to a splitting field. Degree-zero
restriction, the kernel of the symbol at each level, its action on
monomials, and generation of \(U_{T-D}\) by the images \(s^\gamma V_D\)
for \(|\gamma|=T-D\) all enter. Here \(V_j\) is the space of
weight-\(j\) sections in the chosen path of weights. The geometric
proof uses nefness and bigness of both line bundles and coherent
restriction and vanishing on the modified flag space. A rational
choice of lifts for the symbol bases alone gives no upper bound for
their norms or denominators.
\src{src/Development/AbelianVarietyTheory/Definitions/Divisors/HermSysHypOneLeLambdaZero.lean}{199}{Filtration, residue and multiplication hypotheses};
\src{src/Development/AbelianVarietyTheory/Proofs/DivisorProofs/HilbertSamuelFiltrationAtGoodSections.lean}{51}{Geometric construction of the filtration}.

\paragraph{Finite-degree estimate still needed.}
For the same cover, level, weight path and component, the numerical
statement at the point of use is the following. Put \(P_0\ge1\) for
the family size parameter and \(\kappa\ge0\) for its weight parameter.
There must be \(n,\Lambda,\beta,\Lambda_D,\Lambda_T\ge1\) and \(J\in
\mathbb R\) such that
\begin{align}
 \max(\log n,\log\beta,\log\delta)&\le P_0,&
 \log\Lambda&\le P_0(1+\kappa/T),\nonumber\\
 \max(\log\Lambda_D,\log\Lambda_T)&\le P_0\kappa,&
 J&\le 4DP_0(1+\kappa/T)+4TP_0+2DP_0\kappa/T .
 \label{newtheta:budgets}
\end{align}
Define
\[
 u=1-B/d_0,\qquad
 L_h=\log\Lambda+26\log\!\big(nD\delta(2+|h|+\log\Lambda)\big),
 \qquad G=\tfrac{D+1}{2}(\log n+1)+1 .
\]
The required estimate is \(0\le u\le1\) and
\begin{align}
 \frac{g}{d_0}
 &\ge(1-u)\big[Th-3(D+1)|h|-J-26D^2(L_h+\log T)
          \nonumber\\[-2pt]
 &\hspace{39mm}-(T+D+1)\log\beta-(T+D^2)(G+2)\big]
 \nonumber\\
 &\hspace{5mm}-u\big[T(\log2+L_h+\tfrac12\log\beta)
                    +\log\Lambda_D+\log\Lambda_T\big].
 \label{newtheta:finiteHS}
\end{align}
Here \(d_0>0\). The choice of auxiliary parameters may depend on \(Z\).
For each fixed admissible family and accuracy, the threshold in its
scale \(x\) must work simultaneously for all its late admissible
scales, indices, points, stages, good covers under consideration and
components. In particular the exponent budgets and constants in
\eqref{newtheta:budgets} are common. The subsequent comparisons also
require \((1-z)d_0\le B\) and a transfer cost \(c\ge0\) such that
the degree \(g_M\) at the weight of the actual bundle is at least
\(g-d_0c\).
\src{src/Formalism/LocalToGlobal/IndefiniteForms/IndefiniteTree/RealCubicMeans/LogMeanShares/LogMeanStatements/HSClaim.lean}{16}{Exact finite-degree inequality};
\src{src/Development/AutomorphicArithmetic/ThetaCorrespondence/ThetaPairLifts/ThirdPlaceSeries/ExistsDegreeParametersBound.lean}{34}{Family quantifiers and budgets}.

In the construction, \(n\) counts columns of an integral cusp frame
at the first weight, \(\Lambda\) bounds every column norm at every
embedding, and \(\beta\) bounds the Bergman kernel on every eligible
frame of the component. Both auxiliary weights must define
birational linear systems there. Multiplication by an integral frame
column must preserve integral lattices in all consecutive degrees.
At each embedding, multiplying a complex section by coefficients
\(v\) must have norm at most \(\sqrt\beta\|v\|_{\rm frame}\) times
the section norm. Integral sections spanning the spaces over the
coefficient field must be available with norms bounded at every
embedding by \(\exp(DP_0(1+\kappa/T))\) in degree zero and by
\(\exp(P_0\kappa)\) along the path through \(T\). The latter choice
works for both \(\Lambda_D,\Lambda_T\). The former gives
\(J=DP_0(1+\kappa/T)\). One also needs the normalized height of the
birational image in exactly the frame used for the multiplication
norms. The argument behind \eqref{newtheta:finiteHS} then estimates
the height of a determinant polynomial off the exceptional locus
\(Q=0\) and transports its slope bound back to the component
lattice. We give that argument below. Its geometric hypotheses must
still be realized uniformly on the actual covers; a fixed-variety
Hilbert polynomial specifies no bounds for the multiplication norms.
\src{src/Development/AutomorphicArithmetic/ThetaCorrespondence/ThetaPairLifts/HodgeBundleWeights/ExistsDegreeParametersBound.lean}{33}{Simultaneous frame, birational and short-section hypotheses};
\src{src/Development/AbelianVarietyTheory/Proofs/DivisorProofs/ExplicitArithmeticHilbertSamuelBound.lean}{26}{Determinant and Chow-height steps}.

\paragraph{Construction of the common determinant.}
We can in fact derive \eqref{newtheta:finiteHS} from the geometric
restriction, multiplication and metric conditions just identified.
Here are the exact hypotheses for this derivation. Over a number
field \(K\), let \(V_j\) be finite-dimensional section spaces with
full \(\mathcal O_K\)-lattices \(M_j\) and Hermitian norms at every
embedding. Let \(K^N\) act by commuting linear multiplication operators
of degree one. Assume a spanning frame \(w_1,\ldots,w_n\) of \(K^N\)
has norms at most \(\Lambda\) at each embedding and its multiplications
send \(M_j\) to \(M_{j+1}\) for \(0\le j<T\).
In complex coordinates let the input norm be
\(\|v\|_\sigma=\|A_\sigma^{-1}v\|_2\), \(A_\sigma\in\mathrm{GL}_N(\mathbb C)\).
The operator for \(v\) must have norm at most
\(\sqrt\beta\|v\|_\sigma\) in each of these degrees. For \(j=0,D,T\)
assume \(V_j\) has a \(K\)-basis contained in \(M_j\) with all
embedding norms bounded respectively by \(\Lambda_0,\Lambda_D,\Lambda_T\),
each at least one.

Let \(f(s,l)\) be the irreducible Chow polynomial in \(D+1\)
blocks of \(N\) variables, of degree \(\delta\ge1\) in each block.
For every rational \(s=(s_1,\ldots,s_D)\) outside a fixed hypersurface
\(Q(s)=0\), \(Q\ne0\), require the filtration and symbol rules of
Lemma~\ref{newtheta:filtration} after scalar extension, with
\(f(s,l)\) a nonzero constant times \(\prod_p l(y_p)\). The scalar
extension must preserve independence of rational sections in \(V_T\);
the deepest space is generated by \(s^\gamma V_D\) for
\(|\gamma|=T-D\). These conditions involve every good rational
specialization. The norm and lattice conditions involve the original
field and do not require norm control on the splitting fields.

\begin{proposition}[The filtration on the adjoint model]\label{newtheta:koszul}
Let \(X/K\) be a smooth geometrically integral projective variety of
dimension \(D\) over a number field, \(L\) a globally generated line
bundle with generating coordinates \(b_1,\ldots,b_N\), and \(A,B_2\)
nef and big line bundles. Choose a section
\(\xi\in H^0(X,A\otimes(B_2\otimes L^D)^{-1})\) which is invertible
on a nonempty open \(X^\circ\). Write
\[
 W_j=H^0(X,\omega_X\otimes A\otimes L^j),\qquad
 W_j\subset V_j\subset H^0(X^\circ,\omega_X\otimes A\otimes L^j).
\]
Assume \(V_0=W_0\), the \(V_j\) are stable under multiplication by the
coordinates through degree \(T\), and these inclusions persist after
the field extensions used to split the sections. Let \(D\) linear
combinations \(s_i\) of the \(b_j\) cut out a finite reduced transverse
complete intersection \(Y\subset X^\circ\) of degree \(\delta\).
Then the filtration and residue assertions of
Lemma~\ref{newtheta:filtration}, including deepest-space generation as
used in Lemma~\ref{newtheta:power}, hold over any splitting field of \(Y\).
\end{proposition}
\begin{proof}
For \(M=A\) or \(B_2\), and \(k\ge0\), \(M\otimes L^k\) is nef and
big. Multiplication by a fixed nonzero section of \(L^k\), raised to
the \(n\)th power, embeds \(H^0(X,M^n)\) into
\(H^0(X,(M\otimes L^k)^n)\) for every \(n\), proving bigness.
Kawamata--Viehweg vanishing gives
\(H^i(X,\omega_X\otimes M\otimes L^k)=0\) for \(i>0\). Its smooth
projective characteristic-zero hypotheses hold after embedding \(K\)
into \(\mathbb C\), and coherent cohomology commutes with field extension.
This proves the same vanishing over \(K\) and its finite extensions.
Kawamata, \emph{A generalization of Kodaira--Ramanujam's vanishing theorem}, 1982~\cite{kawamataVanishing}.

Let \(\mathcal I\) be the ideal sheaf generated by the \(s_i\) in
local trivializations of \(L\). The Koszul resolution of
\(\mathcal O_Y\) has terms \((L^{-i})^{\binom Di}\). Twisting it by
\(\omega_X\otimes B_2\otimes L^D\) gives terms with exponents
\(D-i\ge0\), all acyclic for global sections by the preceding
vanishing. The global sections of this complex thus compute the
cohomology on \(Y\), in particular the restriction to \(Y\) is
surjective. Multiplication by \(\xi\) identifies the two bundles
over \(Y\), so the restriction of \(W_0\) is surjective there too.
Similarly, twisting by \(\omega_X\otimes A\otimes L^k\) for \(k\ge D\)
gives exactness at the penultimate term:
\[
 \ker\big(W_k\longrightarrow H^0(Y,(\omega_X A L^k)|_Y)\big)
           =\sum_{i=1}^{D}s_i W_{k-1}.
\]
Here juxtaposition of line bundles denotes tensor product.

Work over a splitting field. At a point of \(Y\), the local equations
of the \(s_i\) form a regular system of parameters by transversality.
Thus the associated graded of the local ring is the polynomial ring
in their initial forms over the residue field. Equivalently, monomials
of total degree \(j\) in the \(s_i\) give the coordinates of the
degree-\(j\) Taylor term. This follows also by passing to the completion,
a formal power series ring in these parameters at a smooth rational
point. Trivialize the line bundles compatibly at each point. The kernel
of this Taylor term on sections vanishing to order \(j\) consists of
those vanishing to order \(j+1\), and the symbol of a product is the
product of its initial forms. This gives the diagonal symbol formula
of Lemma~\ref{newtheta:filtration} with \(y_p\) the vector
of the \(b_i(p)\) in the chosen trivialization.

For \(k-j\ge D\), any section in \(W_k\) vanishing to order \(j\)
on \(Y\) belongs to \(\sum_{|\gamma|=j}s^\gamma W_{k-j}\).
Induct on \(j\), the claim for \(j=0\) being immediate. At the next
step represent the section by
\(\sum_{|\gamma|=j}s^\gamma m_\gamma\) using induction. Vanishing
of its degree-\(j\) Taylor term implies \(m_\gamma|_Y=0\) for
every \(\gamma\), by the independence of the monomials at every
point. The kernel formula expresses \(m_\gamma\) as
\(\sum_i s_i m_{\gamma,i}\). Collecting equal monomials gives the
assertion in degree \(j+1\). The degree used in the kernel formula
is \(k-j\ge D\) throughout.

Set \(k_0=T-D\) and
\(V_*=\sum_{|\gamma|=k_0}s^\gamma V_D\subset V_T\).
Filter \(W_T+V_*\) by order of vanishing along \(Y\), with
\(U_j\) the subspace of sections of order at least \(j\).
Since \(V_*\) vanishes to order \(k_0\), every
\(w+v\in U_{k_0}\) with \(w\in W_T,\ v\in V_*\) has \(w\) vanishing
to that order. The preceding generation formula puts \(w\) in
\(\sum s^\gamma W_D\subset V_*\). It follows that \(U_{k_0}=V_*\).
The degree-zero vectors and their products occurring in the symbol
formula lie in \(W_T\), since \(V_0=W_0\). This verifies the
required symbol and deepest-space relations.
\end{proof}
The two adjoint bundles in the family have exactly the powers \(j\)
and \(j+D\) used here. The section \(\xi\) is required to stay
invertible on the chosen open.
\src{src/Development/AbelianVarietyTheory/Definitions/Divisors/HasBoundedPolesGlobalData.lean}{207}{Adjoint models and canonical multiplier};
\src{src/Development/AbelianVarietyTheory/Proofs/DivisorProofs/VanishesToOneOfStalkMap.lean}{133}{Surjectivity for the second bundle};
\src{src/Development/AbelianVarietyTheory/Proofs/DivisorProofs/RestrictedSectionChartCoords.lean}{180}{First-order kernel for the first bundle}.

\begin{proposition}[Boundary vanishing and the adjoint inclusion]\label{newtheta:modified}
Let \(\bar X\) be a smooth integral projective variety over a number
field, \(B\) a reduced effective Cartier divisor, \(U=\bar X-B\),
and \(D=\dim\bar X\). Let \(M_1,M_2,C\) be line bundles on \(\bar X\).
Assume the canonical comparison on \(U\) extends to an isomorphism
\(C\simeq\omega_{\bar X}(B)\). For \(0\le j\le T\), set
\(N_j=M_2 M_1^{D+j}\), and assume that the given space of cusp forms
in weight \(j\), as a space of sections on \(U\), is exactly the
restriction of \(H^0(\bar X,CN_j(-B))\).
Let \(W_i'\subset H^0(\bar X,M_i)\) be finite-dimensional systems
with no base points on \(U\). Let \(\pi:X\to\bar X\) be a proper
birational morphism from a smooth variety, an isomorphism over \(U\),
on which both base ideals become invertible ideals
\(\mathcal O_X(-E_i)\). Put \(L=\pi^*M_1(-E_1)\),
\(B_2=\pi^*M_2(-E_2)\), \(A=\pi^*(M_2M_1^D)\).
Then \(L,B_2\) are globally generated. There is a section
\(\xi\) of \(A(B_2L^D)^{-1}\) invertible on \(U\).
In the adjoint coordinates on \(U\), the cusp spaces satisfy
\[
 V_j=H^0(X,\omega_X A L^j(jE_1))\big|_U,\qquad
 W_j=H^0(X,\omega_X A L^j)\big|_U\ \subset V_j,\qquad W_0=V_0 .
\]
\end{proposition}
\begin{proof}
The local coefficients of \(W_i'\) generate the base ideal.
After pullback, division by a local generator of that invertible
ideal produces functions generating the unit ideal. These functions
are the local coefficients of sections generating the moving part.
The \(E_i\) are supported over \(B\). The canonical section of
\(\mathcal O_X(E_2+DE_1)=A(B_2L^D)^{-1}\) therefore gives \(\xi\).

Under \(C\simeq\omega_{\bar X}(B)\) the cusp space is
\(H^0(\bar X,\omega_{\bar X}N_j)\) on \(U\). For any line bundle \(N\)
on \(\bar X\), rational pullback identifies
\[
 H^0(\bar X,\omega_{\bar X}N)=H^0(X,\omega_X\pi^*N).
\]
Indeed, pullback of a regular top differential is regular by the
map on differentials. Conversely a top form regular on \(X\) defines
a rational top form on \(\bar X\), since the function fields agree.
At each codimension-one point \(y\) the inverse rational map to \(X\)
extends to \(\operatorname{Spec}\mathcal O_{\bar X,y}\) by properness.
Pulling the form back along this morphism puts it in
\((\omega_{\bar X}N)_y\), as differentials commute with localization.
A rational section of a line bundle on a normal variety regular at
all codimension-one points is regular everywhere, by the intersection
description of a normal noetherian domain inside its fraction field.
This proves the equality. Now \(\pi^*N_j=A L^j(jE_1)\).
Multiplication by the canonical section of \(\mathcal O_X(jE_1)\)
is an inclusion and is invertible over \(U\). It is the identity for
\(j=0\). Using this section to identify the two open bundles gives
the displayed comparisons.
\end{proof}

On the toroidal flag component the first hypothesis above requires
the determinant of the Kodaira--Spencer comparison to have local
coefficient \(u/t\), where \(t\) cuts out the \emph{reduced}
boundary and \(u\) is a unit. In smooth toroidal coordinates where
the boundary is \(x_1\cdots x_b=0\), a basis of logarithmic
differentials is \(dx_i/x_i\) for \(i\le b\) and \(dx_i\) for
\(i>b\). Wedging proves the pole formula if the automorphic
comparison matrix in this basis has unit determinant. A comparison
known only on \(U\) does not give that last unit condition. In one
variable, multiplication by \(x\) is invertible on \(x\ne0\) and
changes the boundary order by one. The second hypothesis identifies
\emph{every} form vanishing on the boundary with an automorphic cusp
form of the specified weight and conversely. These two hypotheses
give the inclusion for every \(j\) on the same modification without
any estimate involving \(j\). Generation of the two systems on \(U\),
bigness of their moving parts where used, and nefness and bigness of
\(A\) in Proposition~\ref{newtheta:koszul} require geometric input.
\src{src/Development/AutomorphicArithmetic/ThetaCorrespondence/ThetaPairLifts/HodgeBundleWeights/ExistsWeightHodgeBundleGenerated.lean}{30}{Both base systems and the common modification};
\src{src/Development/AutomorphicArithmetic/ThetaCorrespondence/ThetaPairLifts/HodgeBundleWeights/ExistsSimpleLogPoleAlongBoundary.lean}{38}{Logarithmic Kodaira--Spencer determinant at the boundary};
\src{src/Development/AbelianVarietyTheory/Definitions/Divisors/ExistsResolProjSpace.lean}{203}{Fixed divisors and the moving adjoint coordinates}.

There is a primary comparison theorem that supplies the unit condition
on the standard PEL model. On the characteristic-zero fiber of a
neat unitary PEL moduli space with a compatible admissible smooth
projective cone decomposition, write \(S\) for its smooth toroidal
compactification and \(B_S\) for the reduced boundary. For the
universal semiabelian family the polarized, endomorphism-balanced
Kodaira--Spencer sheaf is isomorphic by its extended Kodaira--Spencer
morphism to \(\Omega^1_S(\log B_S)\). The reduced normal crossings
boundary and the isomorphism are parts (3) and (4) of the same
theorem. Lan,
\emph{Toroidal compactifications of PEL-type Kuga families},
Definition 1.40 and Theorem 1.41~\cite{lanToroidal}.

Here is the determinant on the flag bundle for a CM field action
with conjugation as Rosati involution. Work over a splitting field
for this action and use the polarization isogeny to identify the two
Hodge bundles over the characteristic-zero base. Write \(H_\tau\)
for the dual of the \(\tau\)-part of the Lie bundle, of rank \(m_\tau\).
In Lan's balanced tensor quotient, the relation for an element
\(a\) of the CM field kills the tensor of the \(\tau\) and
\(\nu\) parts unless \(\nu=\bar\tau\), since some \(a\) has
\(\bar\tau(a)\ne\nu(a)\) otherwise. The polarization symmetry
identifies each surviving ordered pair with its transposed pair.
As \(\tau\ne\bar\tau\), the quotient has precisely one summand
\(H_\tau\otimes H_{\bar\tau}\) per unordered pair. This also follows
by taking the projection onto that summand on one ordering and its
transpose on the other, giving an explicit inverse to the inclusion.

Let \(p:\bar X\to S\) be the bundle of complete increasing flags
in the Lie eigenbundles. Let \(Q_{\tau,i}\) denote the
\emph{dual} of the \(i\)-th graded Lie line, numbered
\(0\le i<m_\tau\). The smooth projection gives an exact sequence
for logarithmic cotangents with terms \(p^*\Omega^1_S(\log B_S)\)
and \(\Omega^1_{\bar X/S}\), as is seen on a trivializing flag
chart by adding the flag coordinates to local coordinates on \(S\).
In the determinant the first term contributes \(m_{\bar\tau}\) to
the exponent of every \(Q_{\tau,i}\), by the tensor determinant
formula above. For the relative tangent the graded factors are
\(\operatorname{Hom}(Q_{\tau,i}^{\vee},Q_{\tau,j}^{\vee})\),
\(i<j\): an infinitesimal deformation of a flag is an
endomorphism modulo those preserving its filtration.
The determinant of the relative \emph{cotangent} thus contributes
\(i-(m_\tau-1-i)=2i+1-m_\tau\). Consequently
\[
 \omega_{\bar X}(p^{-1}B_S)
 \simeq\bigotimes_{\tau,i}
 Q_{\tau,i}^{\,m_{\bar\tau}+2i+1-m_\tau}.
\]
This is the character called \(\mathrm{detTanEx}\) in the source;
the central exponent is zero. Thus when the flag and line bundles
in Proposition~\ref{newtheta:modified} are the indicated PEL bundles
and their canonical comparison on the interior is the determinant
of this map up to a nonzero constant, the required boundary order
follows. Checking this comparison uses the identification of the
original framed moduli problem, flags and character rule with this
model. Multiplication by a varying invertible function on the open
would require a separate calculation of its boundary orders.
No metric constant enters this calculation.
\src{src/Development/AbelianVarietyTheory/Definitions/Divisors/ExistsUniqueIsTanFrameOfIsEmpty.lean}{276}{Canonical weight and determinant exponents}.

On the same standard model, let \(\mathcal L\) be the line bundle of
a flag character and set \(\mathcal W=p_*\mathcal L\). Over a
trivialization of the Hodge bundles this direct image is the constant
bundle with fiber the sections of the character line on the product
of flag varieties: formation of \(H^0\) commutes with the flat base
change from the coefficient field. Its gluing is exactly the change
of Hodge frame in that representation. By the projection formula
\[
 H^0(\bar X,\mathcal L(-p^{-1}B_S))
       = H^0(S,\mathcal W(-B_S)).
\]
Thus the cusp comparison reduces to identifying the coefficient
space with the rational form of the usual holomorphic cusp space in
this representation. Over \(\mathbb C\), the latter has the standard
description by sections of the subcanonical bundle
\(\mathcal W(-B_S)\). Over \(K\), equality of two \(K\)-subspaces
in a common space follows from equality after \(K\hookrightarrow\mathbb C\)
by faithful flatness. This last assertion requires both embeddings
of the spaces in the same ambient space of sections on the open.
It gives no comparison of integral lattices or of Petersson norms.
Lan and Stroh,
\emph{Relative cohomology of cuspidal forms on PEL-type Shimura varieties},
introduction preceding Theorem 1.1~\cite{lanStroh}.

\begin{lemma}[One open set of pencils]\label{newtheta:pencil}
In Proposition~\ref{newtheta:koszul}, suppose further that the \(b_i\)
define a birational morphism \(X\to Z\subset\mathbb P^{N-1}\), and
let \(f\) be the Chow form of the geometrically integral image \(Z\)
of degree \(\delta\). There is a nonzero polynomial \(Q(s)\) over
\(K\) such that at every \(s\) with \(Q(s)\ne0\) the required transverse
section lies in \(X^\circ\) and its Chow specialization is a scalar
times \(\prod_p l(y_p)\) over a splitting field.
\end{lemma}
\begin{proof}
There is a dense smooth open of \(Z\) on which the morphism is an
isomorphism. Shrink it to avoid the image of \(X\setminus X^\circ\);
this image is closed by properness and has dimension at most \(D-1\).
In the affine parameter space of \(D\) hyperplane coefficient vectors
the condition that all hyperplanes pass through a fixed point imposes
\(D\) independent equations. Thus the incidence over the excluded
locus has dimension at most \(DN-1\). Over the retained smooth locus,
the hyperplanes through a fixed point induce arbitrary covectors on
its tangent space, since it is embedded in projective space. Linear
dependence of their \(D\) covectors imposes one further nontrivial
determinant equation. The incidence of nontransversality likewise has
dimension at most \(DN-1\). The closures of their parameter images
are therefore proper closed sets. On the remaining open intersect also
the open of proper linear sections of degree \(\delta\) from the
definition of projective degree. It contains a principal open \(Q\ne0\)
defined over \(K\), since all the conditions descend to \(K\).
Each such section has \(\delta\) distinct geometric points and the
same inverse image on \(X\). For fixed \(s\), the defining incidence
property of the Chow form says exactly that \(f(s,l)=0\) when the
last hyperplane passes through one of these points. The resulting
linear forms in \(l\) are pairwise nonproportional. Each divides the
nonzero polynomial \(f(s,l)\), whose degree in \(l\) is \(\delta\).
Their product therefore equals it up to a nonzero scalar. The points
and scalar lie in a finite extension.
\end{proof}

\begin{lemma}[One determinant divisible by a specified Chow power]\label{newtheta:power}
Set \(R_0=\binom{T-1}{D}\), \(B=R_0\delta\), and
\[
 b_*=\sum_{j=0}^{T-D-1}\binom{j+D-1}{D-1}(T-j)
     =R_0\frac{T+D(D+1)}{D+1}.
\]
There are \(0\le y\le d_0-B\) and a single nonzero determinant
polynomial \(\Delta=f^{b_*}H\) over \(K\). Its columns comprise
\(B\) degree-zero short vectors multiplied \(T\) times,
\(y\) degree-\(D\) short vectors multiplied \(T-D\) times,
and \(d_0-B-y\) degree-\(T\) short vectors. The cofactor \(H\)
is multihomogeneous of total degree
\[
 e=BT+y(T-D)-b_*(D+1)\delta\le T(d_0-B).
\]
\end{lemma}
\begin{proof}
The first equality for \(b_*\) counts the exponent in
Lemma~\ref{newtheta:filtration}. Write \(T-j=D+(T-D-j)\).
Summing the latter factor by counting one copy of each term for
every \(j\le i<T-D\) gives \(\sum_{i=0}^{T-D-1}\binom{i+D}{D}
=\binom{T}{D+1}\). Thus \(b_*=D R_0+\binom{T}{D+1}\),
which gives the stated expression.

Choose a rational \(s^0\) outside \(Q=0\) maximizing the rank \(y\)
of the deepest space. Such a point exists since \(K\) is infinite
and the ranks form a nonempty finite set. At \(s^0\), choose
\(\delta\) vectors \(m_r\) from the short degree-zero basis whose
residues form a basis, and \(y\) monomials \((s^0)^\eta m\),
\(|\eta|=T-D\), with \(m\) in the short degree-\(D\) basis,
that span the deepest space. This is possible since the proposed
generators span after base change. Use all the columns
\(s^\gamma l^{T-j}m_r\), \(|\gamma|=j<T-D\), together with
the selected \(s^\eta m\). At \(s=s^0\) choose a rational \(l\)
with all \(l(y_p)\ne0\); finitely many nonzero linear forms over
an extension cannot vanish everywhere on \(K^N\). The jet symbols
and the deepest columns give \(B+y\) independent vectors by
Lemma~\ref{newtheta:filtration}. Complete them using the short
degree-\(T\) basis. With the choices of short vectors and exponents
now fixed, take the determinant in an arbitrary \(K\)-basis of
\(V_T\). Its specialization just constructed is nonzero.

At any other good \(s\) the \(y\) selected deepest columns are
dependent or span the deepest space, by maximality of \(y\). In
the former case the specialized determinant is identically zero.
In the latter take coordinates adapted to the filtration and complete
them by a basis modulo \(U_0\). The matrix is block triangular in
the successive symbol bases, with the deepest block and the block
modulo \(U_0\) independent of \(l\). For each \(j,\gamma\) the
diagonal symbol block is
\(\operatorname{diag}(l(y_p)^{T-j})_{p}\,(m_r(y_p))_{p,r}\).
The symbol rule and surjectivity of restriction show that each graded
piece has exactly the indicated size. Taking determinants gives
\(\Delta(s,l)=c_s f(s,l)^{b_*}\), allowing \(c_s=0\) and working
over the splitting field.

Choose \(l_0\in K^N\) such that \(f(s,l_0)\) is a nonzero
polynomial in \(s\). Clearing the two specialized identities gives
\(\Delta(s,l) f(s,l_0)^{b_*}=f(s,l)^{b_*}\Delta(s,l_0)\)
on \(Q(s)\ne0\), hence identically by polynomial density over the
infinite field. Since \(f\) is irreducible and has positive
\(l\)-degree, it is coprime to \(f(s,l_0)\). Unique factorization
gives \(f^{b_*}\mid\Delta\). Each column has fixed multidegree, so
the determinant has their sum as multidegree, and the quotient by
a nonzero multihomogeneous factor is multihomogeneous. Total
degrees give the stated \(e\): in particular \(e=y(T-D)-BD(D+1)\)
and \(y\le d_0-B\).
\end{proof}

\paragraph{Height of the determinant.}
For a multiform \(P\ne0\) put \(r=[K:\mathbb Q]\) and define
\[
 \mathcal H(P)=\frac{1}{r}\left[
 \sum_{\mathfrak p}\log\|P(w\cdot)\|_{\mathfrak p,\mathrm{Gauss}}+
 \sum_{\sigma:K\hookrightarrow\mathbb C}
 \int\log|P^\sigma(A_\sigma z)|\,dm(z)\right].
\]
The substitution \(w\) sends an \(n\)-tuple to \(\sum_i w_i t_i\)
in each input block. The finite norm is the maximum coefficient
absolute value with \(|x|_{\mathfrak p}=N\mathfrak p^{-v_{\mathfrak p}(x)}\).
The integral uses independent uniform unit vectors in \(\mathbb C^N\)
in all the blocks. At a finite prime the Gauss norm is multiplicative:
divide each polynomial by a coefficient of largest absolute value and
reduce modulo the prime, where products of nonzero polynomials are
nonzero. At infinity logarithms add, so \(\mathcal H\) is additive
on products. The normalized component height used below satisfies
\[
 h=\frac{\mathcal H(f)}{(D+1)\delta}+\tfrac12 H_{N-1},
 \qquad H_j=\sum_{i=1}^{j}\frac1i,\quad H_0=0.
\]
This is the Chow-height normalization for the \emph{same} integral
frame and Hermitian matrices. Identifying the given component height
with this expression requires its birational presentation.
\src{src/Development/AbelianVarietyTheory/Definitions/Divisors/PointHeight.lean}{30}{Frame-substituted Chow height}.

\begin{lemma}[Multiform and determinant height bounds]\label{newtheta:heights}
Write \(C=\log(n\Lambda)+(\log n+1)/2\). For every nonzero
multiform \(P\) of total degree \(v\),
\(\mathcal H(P)\ge-Cv\). For the polynomial of
Lemma~\ref{newtheta:power},
\[
 g\ge\mathcal H(\Delta)
 -B\log(\beta^{T/2}\Lambda_0)
 -y\log(\beta^{(T-D)/2}\Lambda_D)
 -(d_0-B-y)\log\Lambda_T .
\]
\end{lemma}
\begin{proof}
For a homogeneous polynomial \(p\) of degree \(v\) and a unit vector
\(x\), its sphere mean satisfies
\(\int\log|p(z)|\,dm(z)\ge\log|p(x)|-vH_{N-1}/2\) whenever
\(p(x)\ne0\). Rotate coordinates so \(x\) is the first coordinate vector.
Condition on \(z_1\) and average the common phase of the orthogonal
coordinates. Jensen's inequality for the logarithm of the modulus of
a holomorphic polynomial on a disk gives a lower bound
\(\log|z_1^v p(x)|\). For \(N>1\), \(|z_1|^2\) has density
\((N-1)(1-t)^{N-2}\) on \(0<t<1\): a normalized vector of independent
standard complex Gaussians is uniform on the sphere and its squared
coordinates before normalization are independent exponential variables.
Thus \(-\mathbb E\log |z_1|^2=\int_0^1[1-(1-t)^{N-1}]\,dt/t
=H_{N-1}\). For \(N=1\) the cost is zero. Applying the same phase
averages blockwise gives the assertion for multiforms with \(v\) the
sum of the block degrees. The logarithm is bounded above on the
compact product of spheres, and the lower bound starting at any
nonzero value makes its negative part integrable.

The frame spans, so \(P(w\cdot)\ne0\). Choose any nonzero coefficient
\(c\) of it. At finite places \(|c|_{\mathfrak p}\le
\|P(w\cdot)\|_{\mathfrak p,\mathrm{Gauss}}\). At infinity, Cauchy's
formula bounds \(|\sigma c|\) by the supremum of \(|P^\sigma|\)
on tuples of frame combinations with coefficients of modulus at most
one. Each combination has frame norm at most \(n\Lambda\). By the
sphere estimate and homogeneity its logarithm is at most the sphere
mean plus \(v(\log(n\Lambda)+H_{N-1}/2)\). The product formula for
\(c\) and \(H_{N-1}\le1+\log n\) (as \(N\le n\)) give the first
bound.

Fix a rational basis of \(V_T\), and let \(I\) be its determinant
fractional ideal for \(M_T\), generated by determinants of integral
columns. At every finite prime the coefficients of each constructed
column after frame substitution lie in \(M_T\), by multinomial expansion
and stability under the \(w_i\). All coefficients of \(\Delta(w\cdot)\)
therefore lie in \(I\). Its finite log-norm sum is at most
\(-\log N(I)\). At infinity let \(G_\sigma\) be the positive definite
Gram matrix in the same basis. Hadamard's inequality bounds the volume
\(|\Delta^\sigma|\sqrt{\det G_\sigma}\) by the product of the column
norms. Evaluating with each input on its unit sphere, iterating the
multiplication estimate gives the three norm bounds in the statement.
Integrate the logarithm and sum over embeddings. In this basis the
normalized arithmetic degree is
\(g=-r^{-1}(\log N(I)+\tfrac12\sum_\sigma\log\det G_\sigma)\):
the finite contribution is the local determinant ideal and the infinite
contribution is the logarithmic volume of this basis, with a minus
sign. The formula is basis-independent by the product formula. It
yields the second inequality.
\end{proof}

\begin{corollary}[The effective inequality from the normed filtration]\label{newtheta:effective}
Under the common-determinant hypotheses and height identification above,
\eqref{newtheta:finiteHS} holds with \(J=\log\Lambda_0\).
\end{corollary}
\begin{proof}
Put \(h'=\mathcal H(f)/((D+1)\delta)\). By
Lemma~\ref{newtheta:heights}, \(h'\ge-C\) and
\(\mathcal H(H)\ge-Ce\ge-CT(d_0-B)\).
Also \(b_*\mathcal H(f)=B(T+D(D+1))h'\).
Add these estimates in \(\mathcal H(\Delta)\) and insert the
three column costs in the same lemma. Since \(y\le d_0-B\), all
the costs are nonnegative, and \(u=1-B/d_0\), this gives
\begin{align*}
 \frac{g}{d_0}&\ge(1-u)
 [Th'-D(D+1)C-J-\tfrac T2\log\beta]\\
 &\hspace{8mm}-u[TC+\tfrac T2\log\beta
                        +\log\Lambda_D+\log\Lambda_T].
\end{align*}
Here \(0\le u\le1\) by Lemma~\ref{newtheta:filtration}.
The height normalization gives \(h'\ge h-(1+\log n)/2\).
Since \(n,D,\Lambda,\delta\ge1\),
\(L_h\ge\log\Lambda+26\log(2n)\ge C\).
Thus \(D(D+1)C\le26D^2(L_h+\log T)\), and
\(T(1+\log n)/2\le(T+D^2)(G+2)\).
Enlarging the nonnegative costs, including the term
\(3(D+1)|h|\) in the first bracket, now gives
\eqref{newtheta:finiteHS} term by term. The dependence on \(T,D,n\)
is explicit, so the proof applies simultaneously to all components
carrying the stated inputs.
\end{proof}

In particular, this derivation removes a separate appeal to an effective
Hilbert--Samuel estimate \emph{after} the normed filtration and Chow
height have been identified. No bound on \(s^0\) or on its splitting
field enters Lemma~\ref{newtheta:power}. The controlled norms belong
to the short coefficient vectors of a polynomial in variable \(s,l\).
The unproved geometric step here is producing, on each actual modified
flag component at the common family threshold, the restriction
surjectivity at degree zero, the successive symbols and deepest-space
generation on the complement of \(Q=0\), the stated short bases and
multiplication operator norms, and the same-frame Chow-height
identification. Propositions~\ref{newtheta:koszul} and
Lemma~\ref{newtheta:pencil} give the restriction, symbol and generation
statements once the smooth projective model, adjoint inclusions,
nef and big bundles, nonvanishing multiplier, and birational
pencil are realized there. The numerical estimates on \(P_0,\kappa\)
are needed at that threshold. For example, surjectivity of restriction after
tensoring with the coefficient field does not bound the norms of a
basis of the original integral lattice. The rational filtration and
the metric hypotheses address different necessary parts of this step.
\src{src/Development/AbelianVarietyTheory/Proofs/DivisorProofs/SlopeLowerBoundTopLatticeOffHypersurface.lean}{123}{Choice of a maximal-rank section and common determinant};
\src{src/Development/AbelianVarietyTheory/Proofs/DivisorProofs/MultiformHeightLowerBound.lean}{115}{Multiform cost};
\src{src/Development/AbelianVarietyTheory/Proofs/DivisorProofs/LatticeDegreeGeHeightPolynomialVectorsDet.lean}{91}{Determinant ideal and degree}.

The required uniformity of the first frame can be expressed one level
earlier. At an admissible stage write \(Z_0=f_0r_0L_0\) for the
product of the host parameter, rank parameter and scale, and
\(m_1=\sum_{\tau,k}|u_1(\tau,k)|\) for its first weight size.
The short-generator comparison and the uniform Bergman estimate
require fixed \(C_s,C_b\ge0\), independent of the component and the
late stage, for which
\[
 \log n\le C_s Z_0^3,\quad
 \log\Lambda\le C_s(Z_0^3m_1+Z_0^5),\quad
 \sup_{\text{eligible frames}}\operatorname{Berg}(u_1)
                         \le e^{C_bZ_0^3}.
\]
Here the short generators must span the integral frame lattice with
the recorded embedding norms; the logarithmic estimate for
\(\Lambda\) allows enlarging it to at least one.
At the same stage the size bounds give \(Z_0^{12}\le P_0\) and
\(m_1\le Z_0^2(1+\kappa/T)\).
Consequently, as soon as \(Z_0\ge2C_s+C_b+1\),
\(\log n,\log\beta\le P_0\) with \(\beta=e^{C_bZ_0^3}\),
and \(\log\Lambda\le2C_sZ_0^5(1+\kappa/T)\le
P_0(1+\kappa/T)\). This proves the claimed uniform absorption of
these three costs using a single threshold. The existence of the short
generators with that dependence and the componentwise supremum bound
are the upstream metric estimates. A pointwise bound for any fixed
level with an unspecified level-dependent constant cannot provide the
single threshold.
\src{src/Development/AutomorphicArithmetic/ThetaCorrespondence/ThetaPairLifts/ThirdPlaceSeries/HilbertSamuelBergmanFrameBound.lean}{137}{Short frame and uniform cubic and quintic costs}.

\begin{lemma}[From field-spanning sections to integral generators]\label{newtheta:short}
Let \(K\) be a number field of degree \(d\) and absolute
discriminant \(\Delta\). Let \(V\) be an \(R\)-dimensional \(K\)-space
and \(M\subset V\) a finitely generated \(\mathcal O_K\)-module.
At each embedding \(\sigma\), give \(V\) a seminorm \(N_\sigma\)
with \(N_\sigma(ax)=|\sigma(a)|N_\sigma(x)\).
Assume a subset of \(M\) spans \(V\) over \(K\) and all its elements
satisfy \(N_\sigma(x)\le H\) for every \(\sigma\), where \(H\ge0\).
Then \(M\) has at most \(2dR\) generators over \(\mathcal O_K\), all
with every seminorm at most \(\max(1,d^4R\Delta^{1/d}H)\).
\end{lemma}
\begin{proof}
Use the Euclidean structure
\(\|z\|^2=\sum_\sigma|z_\sigma|^2\) on the real subspace of
\(\mathbb C^d\) where \(z_{\bar\sigma}=\overline{z_\sigma}\).
The dual to the embedded lattice \(\mathcal O_K\) consists of the
coordinatewise conjugates of the images of the inverse different
\(\mathfrak D^{-1}\). Indeed its pairing with \(a\in\mathcal O_K\)
then equals \(\operatorname{Tr}_{K/\mathbb Q}(ay)\); the trace
pairing is nondegenerate, and the integral-valued linear functionals
on the full lattice are rational on its rational span.
For \(0\ne y\in\mathfrak D^{-1}\), the ideal \(y\mathfrak D\)
is integral, so \(|N_{K/\mathbb Q}(y)|\Delta\ge1\).
Arithmetic--geometric mean on the \(d\) squared absolute values gives
\(\|y\|\ge\sqrt d\,\Delta^{-1/d}\).
The Euclidean transference theorem
\(\lambda_d(\Gamma)\lambda_1(\Gamma^*)\le d\), applied to this
full real lattice with this dual, now gives \(d\) independent integers
\(\omega_i\) with \(|\sigma(\omega_i)|\le\sqrt d\,\Delta^{1/d}\)
at every embedding. Here \(\lambda_i\) denotes the least radius
of a closed Euclidean ball containing \(i\) independent lattice vectors.
Banaszczyk,
\emph{New bounds in some transference theorems in the geometry of numbers}, Theorem 2.1~\cite{banaszczyk}.

If \(R=0\) the empty set suffices. Otherwise choose a \(K\)-basis
\(b_j\) from the spanning sections. Then \(e_{ij}=\omega_i b_j\)
are in \(M\), form a rational basis of \(V\), and all their
seminorms are at most \(\sqrt d\,\Delta^{1/d}H\).
As an abelian group \(M\) is finitely generated and torsion free
of rank \(dR\). For each element \(v\) of an integral basis of \(M\),
write \(v=\sum q_{ij}e_{ij}\), \(q_{ij}\in\mathbb Q\), and take
\(v'=v-\sum\lfloor q_{ij}\rfloor e_{ij}\).
The triangle inequality gives
\(N_\sigma(v')\le dR\sqrt d\,\Delta^{1/d}H\).
The \(e_{ij}\) and all these \(v'\) generate \(M\) over \(\mathbb Z\),
hence over \(\mathcal O_K\). Their count and the asserted weaker
fourth-power bound follow.
\end{proof}

Let \(R>0\) be the component rank. Suppose
\[
 \log d,\ \log(\Delta^{1/d})\le391 Z_0^2,\qquad \log R\le C_b Z_0^3.
\]
Let \(H=\exp(C_2(Z_0^3m_1+Z_0^5))\) bound the integral
field-spanning sections, with \(C_2\ge0\).
Lemma~\ref{newtheta:short} gives the required count and height of
integral generators with \(C_s=2000+C_b+C_2\). Indeed
\(\log(2dR)\le\log 2+391 Z_0^2+C_bZ_0^3\), and
\(\log(d^4R\Delta^{1/d}H)\le1955 Z_0^2+C_bZ_0^3+\log H\),
where \(Z_0\ge1\). This removes the need to bound an integral change
of basis for the cusp lattice.
\src{src/Development/AbelianVarietyTheory/Proofs/DivisorProofs/ExistsShortGeneratorsOfShortSpanningFamily.lean}{13}{Totally bounded generators of the whole module};
\src{src/Development/AutomorphicArithmetic/ThetaCorrespondence/ThetaPairLifts/HodgeBundleWeights/ShortGeneratorNormComparison.lean}{25}{Uniform field and rank costs}.

The remaining field-spanning assertion concerns the actual deep
cover, coefficient field, integral lattice and Petersson norm on
each component. In the source the weight \(\lambda\) is decreasing
at every embedding, its gap between the last of the \(p\) coordinates
and the first of the \(q\) coordinates is at least \(8f_0r_0+10\),
\(p,q\ge1\), and \(\log\sum_{\tau,k}|\lambda_{\tau,k}|\le10L_0\).
The auxiliary even theta rank used for this conversion is
\(n_\theta=8(f_0r_0+1)\). The span at that even rank is obtained
from three assertions about a common family of theta series:
field-spanning at the given level, an archimedean bound at
\emph{all} embeddings and components, and a positive integer
clearing the finite denominators of the specified rescaled series.
The costs in the logarithms of the archimedean bound and of the
integer must each be at most a fixed constant times
\(Z_0^2\,\mathrm{sizeK}\), where \(\mathrm{sizeK}\) is the
theta weight size
\(\sum_{\tau}\sum_{0\le i<m_\tau}a_{\tau,i}
 +f_0r_0 n_\theta/2\) for the nonnegative polynomial weights
\(a_{\tau,i}\) of the series. The conversion to \(\lambda\) additionally
uses two twisting sections integral at all finite places, with
controlled archimedean norms, whose product induces the prescribed
change of weights. Qualitative generation over the coefficient field
does not bound either cost. For instance, the lattice
\(\mathbb Z\subset\mathbb Q\) is spanned over \(\mathbb Q\) by
\(1/m\), and its least integer denominator is \(m\), arbitrarily
large even with the usual norm at most one. For the actual theta
family Theorem~\ref{intjet:semi} supplies a common finite multiplier,
and Proposition~\ref{intjet:arch} gives a complex estimate in terms
of the shortest period and the norms of the generator, directions
and polynomial. Proposition~\ref{bfam:harm} provides short data for the polynomial span and
Corollary~\ref{bfam:metric} gives the normalized metric bound under its
stated analytic hypotheses. The final identification of their image
with cusp sections on the given charts follows from Propositions~\ref{bfam:unitread}, \ref{dg:identity} and~\ref{munf:span}. In the final weights one must in addition bound the parallel determinant twist of \eqref{ctwist:degree} uniformly on the actual integral components.
\src{src/Development/AutomorphicArithmetic/ThetaCorrespondence/ThetaPairLifts/HodgeBundleWeights/ExistsGoodCoverAtMinimalLevel.lean}{25}{Three theta-family inputs for a single even weight};
\src{src/Development/AutomorphicArithmetic/ThetaCorrespondence/ThetaPairLifts/HodgeBundleWeights/ShortSpanCompositionClosed.lean}{50}{Change of weight and controlled twisting sections}.

\paragraph{Cancellation of the lattice index in a rescaled series.}
Here the auxiliary integer \(N\) can be arbitrarily large for a
fixed level. Write \(O=\mathcal O_E\), \([E:\mathbb Q]=2f\),
\(r=p+q\), \(M\) for the positive integer level, and
\(\Lambda\subset E^n\) for the self-dual lattice for
\(\langle a,b\rangle=\sum\bar a_i b_i\).
The construction takes \(N\) odd, \(M\mid N\),
\(NO^n\subset\Lambda_0=\Lambda\cap O^n\). On a geometric
characteristic-zero PEL fiber let \(A\) be the abelian variety with
symmetric polarizing line \(L\) and theta section \(s_0\).
Its period lattice \(\Pi\) is projective of rank \(r\) over \(O\).
Put \(d_P=\deg P\) for a nonzero polynomial in the integral harmonic
space of the series. Write \(\mathcal Q\) for the \(n\)-fold level
quotient and \(C_M=|\mathcal Q|=M^{frn}\). In fact the datum's
projective lattice of rank \(r\) is reduced modulo the half-level
ideal \(\mathfrak m\), with \(\mathfrak m\bar{\mathfrak m}=(M)\).
Hence \(\mathrm N\mathfrak m=M^f\), and local freeness, prime by
prime, gives the asserted cardinality. The Fourier weight
\(\omega_{x_0}(u)=\zeta_M^{\langle x_0,u\rangle_M}\) is a unit at
every finite place. Here the exponent is the integral level pairing.

The torsion set \(Z_\Lambda\subset A[N^2]^n\) consists of the
tuples \(z\) satisfying \((N\mu)z=0\) for all rows
\(\mu\in\operatorname{Hom}_O(\Lambda,O)\subset E^{1\times n}\).
The entries of \(N\mu\) are integral. For \(u\in\mathcal Q\), fix the
specified \(N\)-division point \(t_u\in A[N^2]^n\) of its level
point. Form the pullback of \(s_0^{\boxtimes n}\) by multiplication
by \(N\) on \(A^n\), then translate by the normalized theta lift of
\(z\), then by that of \(t_u\). Denote by \(J_P(u,z)\) the jet
evaluation of this section in the given Hodge frame. Explicitly,
for the individual translated sections \(F_b\) in the \(n\) factors
it is
\[
 J_P(u,z)=\sum_{\mathbf k} P_{\mathbf k}
       \prod_{b=1}^{n}\nabla_{\mathbf k_b}F_b .
\]
Here \(\mathbf k_b\) is the column of the monomial's multi-index at
\(b\); \(\nabla_{\mathbf k_b}\) is the prescribed linear theta
jet in that frame. The expression is multilinear in the factors,
hence depends on their tensor. For the raw theta section \(y\)
and the chosen scalar \(s\), the actual finite sum is
\[
 (s y)(x_0)=\frac{s}{C_M\,m_\Lambda N^{d_P}}
     \sum_{u\in\mathcal Q}\omega_{x_0}(u)
                      \sum_{z\in Z_\Lambda}J_P(u,z),\qquad
 m_\Lambda=[\Lambda_0:NO^n]^r .
\]
The notation suppresses the geometric point and Hodge frame. The
auxiliary determinant scalar \(\beta_H\) in the short evaluation
of \(\beta_H^{-1}y\) cancels when one returns to \(y\).
\src{src/Formalism/LocalToGlobal/ProjectiveLattices/ThShort.lean}{244}{Finite torsion evaluation with all three denominators};
\src{src/Development/AbelianVarietyTheory/Definitions/PlaceAdaptation/AdaptStd.lean}{107}{Harmonic evaluation on translated sections}.

\begin{lemma}[Exact cancellation on every level class]\label{newtheta:index}
Put \(K_0=Z_\Lambda\cap A[N]^n\). Then \(|K_0|=m_\Lambda\) and
\(J_P(u,z+v)=J_P(u,z)\) for every \(v\in K_0\). In particular
\(m_\Lambda^{-1}\sum_{z\in Z_\Lambda}J_P(u,z)\) is precisely
the sum over \(Z_\Lambda/K_0\), for each \(u\) separately.
\end{lemma}
\begin{proof}
Work first over \(\mathbb C\). Write \(z=x/N\bmod\Pi^n\).
The defining conditions say that \(\mu x\in\Pi\) for all these
rows, that is, \(x\in\Lambda\otimes_O\Pi\) inside
\((E\otimes_O\Pi)^n\). This last equality can be checked after
localizing at every prime of \(O\), where both lattices are free
and evaluation against the dual basis recovers all coordinates.
Self-duality and \(NO^n\subset\Lambda\) give
\(\Lambda\subset N^{-1}O^n\), so \(N^2z=0\) already holds.
The additional condition \(Nz=0\) is \(x\in\Pi^n\).
Since \(\Pi\) is locally free of rank \(r\), this identifies \(K_0\)
with \((\Lambda_0/NO^n)\otimes_O\Pi\), whose order is \(m_\Lambda\).

Let \(E_L\) be the integral alternating Riemann form of the
polarization, extended rationally. For \(x/N,x'/N\) in this
description the commutator of the two translations on
\((L^{\boxtimes n})^{N^2}\) is
\(\exp(\pm2\pi i\sum_i E_L(x_i,x'_i))\), by the Weil pairing
formula for the Riemann form. For \(x_i=a_i h,\ x'_i=b_i h'\),
with \(a,b\in\Lambda,\ h,h'\in\Pi\), the exponent before
\(\pm2\pi i\) is \(E_L(h,\langle a,b\rangle h')\) by the
Rosati involution. This is integral by self-duality; bilinearity
gives the same result for sums of tensors. Hence the lifts commute
on \(Z_\Lambda\).

For clarity, on odd \(N^2\)-torsion the normalized lift has
\(N^2\)-th power one and inversion sends it to its inverse.
Two such lifts at the same point differ by a scalar \(c\) with
\(c^{N^2}=c^2=1\), hence coincide. On a commuting subgroup the
product satisfies both normalizations, so the normalized lift there
is multiplicative. On \(A[N]^n\), use the canonical identification
of the pullback by \([N]\) with \((L^{\boxtimes n})^{N^2}\).
Translation by \(v\in A[N]^n\) on the pullback lifts by the
identity on the fiber at \([N]x=[N](x+v)\). This lift is
normalized, as its powers and inversion are the corresponding
translations on a pullback. It fixes the pulled-back section.
Consequently translation by \(z+v\) has exactly the same effect
as translation by \(z\) on that section. The subsequent translation
by \(t_u\) and the linear jet evaluation preserve equality.
The identities descend to arbitrary characteristic-zero geometric
fibers by defining the abelian variety, the finitely many points,
the line, the sections and the required finite jets over a finitely
generated subfield of characteristic zero and embedding it into
\(\mathbb C\). Mumford,
\emph{On the equations defining abelian varieties I}, \S1~\cite{mumfordEquations}.
\end{proof}

The scalar \(s\) in the same formula also admits a simultaneous
choice at the finite and infinite places. Write
\(I=[\Lambda:\Lambda_0]\), let \(F\) be the totally imaginary
coefficient field of degree \(h\) containing the coefficient
embeddings \(\tau:E\to F\), and set
\(\mathfrak a=\det(\Lambda_0)\det(\Lambda)^{-1}\), using the
fractional determinant ideals in \(E\). Then
\(\mathfrak a\subset O\) and \(\mathrm N\mathfrak a=I\):
over each discrete valuation ring the inclusion has elementary
divisors and both statements follow by multiplying those divisors.
Put
\[
 \mathfrak b=\prod_\tau
       \big(\tau(\mathfrak a)\mathcal O_F\big)^{m_\tau}.
\]
The finite scalar condition is \(s\in\mathfrak b^{-1}\), since
the quotient of the two type norms of determinants is exactly
\(\mathfrak b\). Each extended ideal in the product has norm
\(I^{h/[E:\mathbb Q]}\), by the residue-degree and ramification
identity above each prime. Since \(\sum m_\tau=fr\), we have
\(\mathrm N\mathfrak b=I^{hr/2}\).
The image of \(\mathfrak b^{-1}\) in \(\mathbb C^{h/2}\) has real
covolume \(2^{-h/2}\sqrt{|\Delta_F|}/\mathrm N\mathfrak b\),
by the discriminant determinant for an integral basis and the
fractional ideal index. Take the product of closed disks of common
radius \(R_s=|\Delta_F|^{1/(2h)} I^{-r/2}\). Its volume
\(\pi^{h/2}R_s^h\) exceeds \(2^h\) times that covolume because
\(\pi>2\). Minkowski's convex body theorem supplies a single nonzero
\(s\in\mathfrak b^{-1}\) with \(|\sigma(s)|\le R_s\) at every
embedding, including the conjugates. At every level prime
\(\ell\mid M\), local standardness of \(\Lambda\) gives
\((\Lambda_0)_\ell=\Lambda_\ell\); hence \(\mathfrak b\) is a unit
at all primes over \(\ell\), and this same \(s\) is integral there.
This proves the determinant rescaling bound on the actual family,
including its index dependence. It applies before the jet estimates.
\src{src/Development/AbelianVarietyTheory/Proofs/DivisorProofs/ExistsSmallScalarReflexIdealSelfDualLattice.lean}{29}{Single rescaling scalar for all embeddings and the determinant ideals}.

In fact \(|Z_\Lambda/K_0|=I^r\) by the same locally free tensor
calculation. Cauchy--Schwarz on the \(C_M I^r\) remaining terms
and \(|\sigma(\omega_{x_0}(u))|=1\) thus give at that frame
\[
 |\sigma(sy)(x_0)|^2
 \le \frac{|\Delta_F|^{1/h}}{C_M N^{2d_P}}
       \sum_{u,\,\bar z}|J_{P^\sigma}(u,\bar z)|^2 .
\]
The right side uses the conjugate geometric data at \(\sigma\).
In particular no separate factor \(I^r\) is lost in estimating the
sum by the squared jet values.

There is an exact test for the finite denominator.
Write \(w=\sum_{\tau,i}a_{\tau,i}\),
\(\delta_E=|\operatorname{disc}E|\), \(S=w+frn/2\), and
\[
 D_M=C_M(M\delta_E)^w.
\]
This is the positive integer \(\mathrm{intZ}\) of the construction,
whose proposed common multiplier is \(D_M^2\).
For any valuation \(v\) of an integral framed test, extended to the
finite torsion field, put
\[
 T_{x_0}=s\sum_{u\in\mathcal Q}\omega_{x_0}(u)
                   \sum_{\bar z\in Z_\Lambda/K_0}J_P(u,\bar z).
\]
At that test the desired integrality is \emph{equivalent} to
\[
 v(T_{x_0})\ge d_P v(N)+v(C_M)-2v(D_M).
\]
Thus at a prime over \(\ell\mid N\) with
\(\ell\nmid M\delta_E\), the required divisibility is by
\(N^{d_P}\) locally, even though \(\ell\) need not divide the level.
At a level prime the allowable loss is exactly
\(2v(D_M)-v(C_M)\) relative to \(d_Pv(N)\).
Integrality of the roots of unity and the geometric coset
cancellation impose no such divisibility on the harmonic jets.
Checking all these tests is stronger than checking an expansion at
one cusp unless a corresponding integral expansion principle for
the chosen lattice is supplied.

\begin{proposition}[The auxiliary integer at good abelian reduction]
\label{newtheta:goodjet}
Let \(\ell>\max(2,d_P)\), \(\ell\nmid M\delta_E\), and
\(\Lambda\otimes\mathbb Z_\ell=O^n\otimes\mathbb Z_\ell\).
At a mixed-characteristic discrete valuation ring \(R\) of residue
characteristic \(\ell\), take a framed test extending to an abelian
scheme equipped integrally with the rigidified symmetric line,
theta section, integral Hodge directions, polarized biextension
and its rigidified identification with \(L\) on the graph of the
specified endomorphism \(w_0\). Then for every allowed \(N\),
every \(u,z\), and the actual translated sections,
\[
                 v(J_P(u,z))\ \ge\ d_P v(N).
\]
The valuation may be extended to define the torsion points.
This proves the preceding denominator test at these places.
\end{proposition}
\begin{proof}
Put \(N_0=\operatorname{lcm}(M,\exp(O^n/\Lambda_0))\), taking
the exponent of the trivial group to be one. Then \(N_0\mid N\),
\(\ell\nmid N_0\), and \(k=N/N_0\) sends \(t_u,z\) to \(kt_u,kz\).
The first maps under \(N_0\) to the same level point. The second
belongs to \(Z_{\Lambda,N_0}\) by the description \(z=x/N\)
in Lemma~\ref{newtheta:index}. Both images are killed by \(N_0^2\).
For either translation by \(t\), lift translation by \(kt\) on
\(L^{N_0^2}\) and pull its line isomorphism back along \([k]\),
using \([k]\circ T_t=T_{kt}\circ[k]\).
The resulting lift at \(t\) is symmetric. Its \(N^2\)-th power is
the identity, since \(N^2t=0\) and \(N_0^2\mid N^2\).
By uniqueness of the normalized odd lift this is precisely the
one used in the test. The rigidified identifications for successive
multiplications compose, so both translations, in their prescribed
order, carry the pullback of \([N_0]^*s_0\) to the pullback of
the corresponding translated section.

Here is the jet calculation, including the frame correction.
In formal coordinates write \(e(x,y)\) for a normalized frame of
the polarized biextension. Its logarithmic correction \(h(x,y)\)
vanishes on the axes. The frame \(e\exp(-h)\) multiplies exactly
in each argument, since the two multiplication factors of \(e\)
are the exponentials of the respective additive defects of \(h\).
Pulling this corrected frame to \((x,w_0x)\) yields the frame in
which the jet of a section of \(L^m\) with raw coefficient \(f\)
is obtained by differentiating \(f\exp(mh(x,w_0x))\).
In changing the raw biextension frame to \(e b\), one can take
\(h+\log b\) for the correction; the raw coefficient changes to
\(f b(x,w_0x)^{-m}\), leaving the corrected germ unchanged.
Any further choice differs by a biadditive Rosati-compatible term,
whose contribution to the harmonic jet vanishes as shown below.
Biadditivity identifies the corrected frame at \((kx,kw_0x)\)
with the \(k^2\)-th tensor power of that at \((x,w_0x)\).
The induced isomorphism of lines is the rigidified one: it is
constructed by iterating the two global biextension laws and is
the identity in the rigidification at zero. In formal logarithm
coordinates \(x\), multiplication by \(k\) is \(x\mapsto kx\).
Each homogeneous derivative of degree \(j\) therefore gains \(k^j\).
The prescribed row sums of every monomial of \(P\) give the same total
degree \(d_P=\sum_{\tau,i}a_{\tau,i}\).
Suppressing \(u\) and displaying both translation points gives the
exact equality on the generic fiber
\[
       J_{P,N}(t_u,z)=k^{d_P} J_{P,N_0}(kt_u,kz).
\]

It remains to check the smaller-level jet over \(R\). All translations
now have order prime to \(\ell\). After finite extension their points
extend by properness. Each translated line is isomorphic to the
original line over \(R\): their quotient represents zero in the
relative Picard scheme by the polarization and \(N_0^2\)-torsion,
and a line from the local base is free. Any integral lift has
\(N_0^2\)-th power a unit. The scalar that changes it to the
normalized lift likewise has valuation zero, by taking that power.
Thus each of the smaller-level sections has integral raw germ in
an integral frame.

Work through total degree \(d_P\). The formal logarithm and its
inverse are integral to this order: integrating the closed invariant
differentials term by term divides each homogeneous component only
by its positive total degree (contract the form with the Euler
vector field), and inversion of a series with identity linear term
requires no division. In these coordinates the group law is
addition. Restrict both arguments \(X,Y\) to the linear combinations
of the integral eigen-directions of the frame, with variables
indexed by \((\tau,i)\). Write \(h\) also for this restriction.
The three series
\[
 \begin{split}
 h(X,Y)+h(X',Y)-h(X+X',Y),\qquad&
 h(X,Y)+h(X,Y')-h(X,Y+Y'),\\
 &h(aX,Y)-h(X,\bar aY)\qquad (a\in O)
 \end{split}
\]
are integral through that order. Indeed they are the logarithms
of the integral multiplication and Rosati adjunction factors of the
biextension frame, each of constant coefficient one; the truncated
logarithm has unit denominators. For a coefficient of \(X^\alpha
Y^\gamma\) with \(|\alpha|\ge2\), split off one variable \(X'_j\)
with \(\alpha_j>0\) in the first defect. Its coefficient is
\(-\alpha_j\) times the one sought, with \(\alpha_j\) a unit.
The second defect handles \(|\gamma|\ge2\). The axes vanish.
In bidegree \((1,1)\), the last defect multiplies the coefficient
in types \((\tau,\nu)\) by \(\tau(a)-\nu(\bar a)\).
For \(\tau\ne\bar\nu\) some \(a\in O\) makes this a unit:
since \(\ell\nmid\delta_E\), the distinct embeddings have distinct
reductions in the common splitting residue field. Thus, to this
order, \(h\) is integral up to a bilinear form supported on
the conjugate pairs.

Write \(X_b\) for the variables in factor \(b\).
After evaluation at \((X_b,w_0X_b)\) and summation over \(b\),
the possible nonintegral part is a linear combination \(Q\) of
\(\sum_b X_{\tau i b}X_{\bar\tau j b}\).
Under the pairing \(P(\partial)F(0)\), multiplication by such a
quadratic is adjoint to \(\sum_b\partial_{\tau i b}
\partial_{\bar\tau j b}\) on \(P\), by the monomial factorial formula.
These operators annihilate \(P\) by its defining harmonic equations.
Consequently the factor \(\exp(N_0^2 Q)\) contributes only its
constant term. The remaining exponentials, raw germs and coefficients
of \(P\) are integral through degree \(d_P\). For the coefficients
of \(P\) this uses local standardness of \(\Lambda\) in the
integral linear generators of its polynomial ring.
We obtain \(v(J_{P,N_0})\ge0\), and \(v(k)=v(N)\).
The same local standardness makes \(\mathfrak b\) a unit at this
place, so \(s\) is integral as well.
\end{proof}
For the standard lattice this works at every such good test with
\(N_0=M\), even when \(\ell\) divides \(N\) to arbitrary order.
We next remove the restrictions on the residue characteristic and
the local lattice at abelian reduction. The archimedean norm
of the smaller-level sections is a separate estimate.
\src{src/Formalism/LocalToGlobal/ProjectiveLattices/RowDeg.lean}{43}{Harmonic trace equations with identity column matrix};
\src{src/ReciprocityTheory/ReductiveGroups/ArithmeticGeometry/BirationalGeometry/ToroidalCompactification/CocycleRightZeroThd.lean}{174}{The two biextension defects and the adjunction factor};
\src{src/Development/AbelianVarietyTheory/Proofs/JetHarmonics/HarmonicValueNormalisedPullbackChainRule.lean}{516}{Normalized isogeny on harmonic evaluations}.

\paragraph{Finite denominators at abelian reduction.}
\label{intjet:sec}
Write \(d=\deg P=\sum a_{\tau i}\), \(\Delta=|\operatorname{disc}E|\),
and \(v\) for the additive valuation of a mixed-characteristic
discrete valuation ring \(R\). Valuations and all the inclusions below
may be tested after finite extension. We retain the integral symmetric
line, its rigidification, the section, the Hodge directions and the
polarized biextension on an abelian scheme \(A/R\) as in
Proposition~\ref{newtheta:goodjet}. The graph of \(w_0\in O_E\)
in that biextension gives the line. In a semiabelian degeneration the
generic abelian scheme need not extend in this form.

\begin{lemma}[Integral Taylor derivatives of harmonic jets]\label{intjet:der}
In the preceding formal chart take an integral raw coefficient,
any nonnegative tensor power of the line, and an integral homogeneous
harmonic polynomial of degree \(d\) at the standard lattice.
Its evaluated jet belongs to \(\Delta^{-\lfloor d/2\rfloor}R\).
The assertion holds in every residue characteristic.
\end{lemma}
\begin{proof}
Work over the fraction field in logarithm coordinates. For a series
\(F\) define its Taylor values to be the mixed derivatives
\(\partial^\alpha F(0)\), with no division by \(\alpha!\).
If \(F\) comes from an integral function on the formal group,
on integral tangent vectors these values lie in \(R\). Indeed
each vector extends by translation to an invariant integral
derivation of the formal group ring. In logarithm coordinates it
is the constant vector field, so applying a sequence of these
derivations and evaluating at zero gives exactly the desired value.
This argument works on a product and on any list of integral
directions, whether or not they form an integral basis.

Series with integral Taylor values form a ring by Leibniz's formula.
An inverse with unit constant term has integral Taylor values by the
same formula and induction. For \(g(0)=1\) and integral Taylor
values, those of \(\log g\) are integral by
\(\partial_i\log g=g^{-1}\partial_i g\). For \(h(0)=0\) with
integral Taylor values, use
\(\partial_i\exp h=(\exp h)\partial_i h\) and induction.
These are identities over the fraction field and give integrality
of derivatives without an estimate for any inverse factorial.

Let \(h(X,Y)\) be the logarithmic correction of the biextension in
the chosen frame. Its two additive defects and its \(O_E\)-adjunction
defect are logarithms of integral units with constant term one, as
in Proposition~\ref{newtheta:goodjet}. For an index \(\alpha\) of
degree at least two write \(\alpha=\alpha'+e_j\), \(|\alpha'|>0\).
Apply \(\partial_X^{\alpha'}\partial_{X'_j}\partial_Y^\beta\) to
\(h(X,Y)+h(X',Y)-h(X+X',Y)\) at zero. The result is
\(-\partial_X^\alpha\partial_Y^\beta h(0)\). This proves its
integrality. Interchanging the two arguments treats \(|\beta|\ge2\).
The two axes vanish. In bidegree \((1,1)\) write \(c_{\tau i,\nu j}\)
for the derivative along eigenvectors of types \(\tau,\nu\).
The adjunction defect gives
\((\tau(a)-\nu(\bar a))c_{\tau i,\nu j}\in R\) for every \(a\in O_E\).
If \(\tau\ne\bar\nu\), then \(\Delta c_{\tau i,\nu j}\in R\).
Indeed choose an integral basis \(b_i\) and its trace-dual basis
\(b_i^\vee\). The embedding matrices are inverse transposes, so
\(\sum_i\tau(b_i^\vee)\sigma(b_i)\) is one for \(\sigma=\tau\)
and zero otherwise. All \(\Delta b_i^\vee\) are integral by the
adjugate of the trace matrix. Apply this equality to
\(\sigma=\tau,\bar\nu\) and subtract.

Thus the gauge on the eigen-directions is the sum of a series with
integral Taylor values, a bilinear form with coefficients in
\(\Delta^{-1}R\) on the nonconjugate pairs, and an unrestricted
bilinear form on the conjugate pairs. On the graph of \(w_0\) these
are quadratic terms, summed over the product factors. The last
quadratics are killed by harmonicity as in
Proposition~\ref{newtheta:goodjet}. In the exponential of the
nonconjugate quadratic each derivative of order \(j\) at zero is
a sum over partitions into pairs of the \(j\) labelled derivative
positions. This follows by differentiating \(\partial_i\exp
Q=(\exp Q)\partial_i Q\); a singleton has zero value and derivatives
of \(Q\) of order greater than two vanish. Each summand costs at
most \(\lfloor j/2\rfloor v(\Delta)\). Leibniz's formula for the
remaining integral factors gives the claimed bound.
\end{proof}

\begin{proposition}[All abelian tests at a locally standard lattice]
\label{intjet:standard}
Suppose \(\Lambda_\ell=(O_E^n)_\ell\) and the integral abelian data
above are given. For every allowed \(N\) and each pair of translation
points in the original order one has
\[
 v(J_{P,N}(u,z))\ \ge\
 d\bigl(v(N)-v(M)\bigr)-\lfloor d/2\rfloor v(\Delta).
\]
In particular \(D_M^2\) clears the rescaled sum at every such test,
including the level primes, \(2\), and primes ramified in \(E\).
\end{proposition}
\begin{proof}
Let \(N_0=\operatorname{lcm}(M,\exp(O_E^n/\Lambda_0))\) as before.
Now \(v(N_0)=v(M)\). The identity
\(J_{P,N}= (N/N_0)^d J_{P,N_0}\) on the generic fibre proved in
Proposition~\ref{newtheta:goodjet} works whether or not \(\ell\mid N_0\).
All torsion points of the smaller test extend after finite extension
by properness. For \(t\in A[N_0^2]\) the line \(T_t^*L^{N_0^2}\)
is integrally isomorphic to \(L^{N_0^2}\). The class of their
ratio in the dual abelian scheme is zero since \(N_0^2t=0\),
and an invertible module on the local base is free. Any integral
translation lift has its \(N_0^2\)-th power in \(R^\times\).
The generic scalar adjusting it to the normalized lift therefore
has valuation zero. This applies to both translations in succession,
also for torsion of residue-characteristic-power order.
The smaller-level section has integral raw germ. Apply
Lemma~\ref{intjet:der}. Local standardness makes \(P\) integral in
the standard columns and \(s\) integral at this place. Sum after the
index cancellation of Lemma~\ref{newtheta:index} and use
\(D_M=C_M(M\Delta)^d\).
\end{proof}

\begin{theorem}[The trace at an arbitrary self-dual lattice at abelian reduction]
\label{intjet:trace}
Suppose the abelian data above extend integrally. Take any
\(\Lambda,N,P,s\) in the rescaled series, with \(\Lambda\) standard
at the primes of \(M\), and the prescribed level points \(t_u\).
At a prime \(\ell\nmid M\),
\[
 v\left(s\sum_{\bar z\in Z_\Lambda/K_0}J_{P,N}(u,\bar z)\right)
 \ \ge\ d\,v(N)-\lfloor d/2\rfloor v(\Delta)
 \qquad(u\in\mathcal Q).
\]
Together with Proposition~\ref{intjet:standard}, this proves the
finite denominator test at every residue characteristic on these
abelian models.
\end{theorem}
\begin{proof}
Let \(O=O_E\), \(B=A^n\), \(B_0=\Lambda_0\otimes_O A\) and
\(B_1=\Lambda\otimes_O A\). A full lattice over the Dedekind
ring \(O\) is finitely generated projective. Its tensor with \(A\)
is obtained by splitting an idempotent on a product of copies of
\(A\). This gives abelian schemes over \(R\). Their Lie algebras
are obtained by taking the same tensor with \(\operatorname{Lie}A\).
Consider the maps induced by
\[
 O^n\xrightarrow{\,N\,}\Lambda_0
 \xrightarrow{\,i\,}\Lambda,\qquad
 \Lambda_0\xrightarrow{\,j\,}O^n
\]
where \(i,j\) are inclusions. Write \(a_0,i,j\) also for the
resulting isogenies, and \(a=i a_0\). Then \(j a_0=[N]\).
These maps are isogenies on every fibre since suitable integer
multiplications factor through them in both directions. They
are finite flat by properness and flatness of a quasi-finite map
between smooth schemes of the same relative dimension.

Give each tensor scheme its rigidified line from the graph of
\(w_0\) in the polarized biextension for the restricted Hermitian
form. Denote the lines on \(B,B_1\) by \(\mathcal L,\mathcal L_1\).
Then \(j^*\mathcal L=i^*\mathcal L_1\) and
\(a^*\mathcal L_1=[N]^*\mathcal L=\mathcal L^{N^2}\)
by rigidified identifications. To see the normalization, on free
modules a matrix \(U\) pulls the biextension for a Hermitian
matrix \(G\) back to that for \(U^\dagger G U\). Expand the
biextension in its two arguments and move each first-argument
entry across by Rosati adjunction. These identifications are
the identity on the axes. Restriction to direct summands proves
the assertion for projective modules as well. This construction
uses the line on the graph, so no choice among symmetric lines
inducing the same polarization is needed. See Amir-Khosravi,
\S\S1.2--1.4~\cite{serreTensor}.

On the geometric generic fibre, the period lattice calculation of
Lemma~\ref{newtheta:index} gives \(\ker a=Z_\Lambda\) and
\(\ker a_0=K_0\). The natural action of \(\ker a\) on the pullback
line is exactly its normalized lift. It has \(N^2\)-th power one,
and inversion on the symmetric pullback sends the lift to its
inverse. These conditions determine the lift since \(N\) is odd.
In particular they agree with the products of the individual
normalized lifts on the \(n\) factors by the same uniqueness.
Put \(t_0=j^*(s_0^{\boxtimes n})\), viewed as a section of
\(i^*\mathcal L_1\). Before the level translation the sum of
the \(|Z_\Lambda/K_0|\) translated sections on \(B\) is
\[
                 a^*\operatorname{Tr}_i(t_0).
\]
Indeed a finite étale isogeny on the generic fibre has trace
equal, after pullback, to the sum over its kernel. The input
on \(B\) was \(a_0^*t_0=[N]^*(s_0^{\boxtimes n})\).

The determinant of \(di\) generates the ideal
\(\mathfrak b R=\prod_\tau
 (\tau(\det\Lambda_0\,\det\Lambda^{-1})R)^{m_\tau}\).
Choose bases of the two lattices over the semilocal Dedekind
ring \(O\otimes\mathbf Z_{(\ell)}\), where they are free.
The determinant on each generic eigenspace is the indicated
embedding of the determinant of the inclusion, repeated \(m_\tau\)
times. The chosen bases give integral changes of basis on the
Lie lattices, so this calculates their determinant ideal even
when the eigenspaces do not split integrally.
Translation of invariant differential frames carries this same
different ideal to every point. For a finite flat map between smooth
schemes of equal relative dimension the different is generated by
the determinant of its differential. Locally in the finite
algebra the trace functional is this generator times a generator
of the dualizing module, by the complete intersection formula
for the different~\cite{tateDifferent}. Trivializing \(\mathcal L_1\)
shows that \(\operatorname{Tr}_i(t_0)\) is divisible by
\(\mathfrak b R\) as a section. Thus \(S=s\operatorname{Tr}_i(t_0)\)
is integral, including when \(v(s)<0\).

It remains to apply the actual final translation. Its point
\(t=t_u\) is killed by a power of the half-level ideal, hence by
a power of \(M\); this extra condition is part of the evaluation
rule. In particular its order \(h\) is prime to \(\ell\).
\src{src/Formalism/LocalToGlobal/ProjectiveLattices/ThShort.lean}{252}{The support of the prescribed division point}.
Set \(b=a(t)\). Rigidify the line
\(\mathcal D=T_b^*\mathcal L_1\otimes\mathcal L_1^{-1}\) at zero
by choosing an integral generator of its fibre. It is an element
of \(\operatorname{Pic}^0(B_1)[h]\), since \(b\mapsto
T_b^*\mathcal L_1\otimes\mathcal L_1^{-1}\) is a homomorphism
to the dual abelian scheme. Thus \(\mathcal D^h\) has an integral
rigidified trivialization. On the formal neighborhood of zero
there is a unique \(h\)-th root frame with prescribed value at
zero and this trivialization as its \(h\)-th power. Coefficient
recursion only divides by the unit \(h\).

The normalized lift at \(t\) identifies the translation of
\(a^*S\) with the pullback of \(T_b^*S\), followed by an
isomorphism \(a^*\mathcal D\to\mathcal O_B\). This isomorphism
is integral by the translation-lift argument of
Proposition~\ref{intjet:standard}. Its \(h\)-th power and the
pullback of the chosen trivialization differ by a global unit
of the abelian scheme, hence by an element of \(R^\times\).
It takes the pullback of the root frame to a constant unit:
over the fraction field a formal unit whose \(h\)-th power is
constant is constant. In particular, in the formal neighborhood
the translated section is, up to a constant unit, \(a^*\) of
an integral germ of \(\mathcal L_1\), obtained by dividing the
germ \(T_b^*S\) by the root frame. This explains why the final
translation does not change the trace divisibility.

To compute the jet of this pullback, choose a split surjection
\(U:O^t\to\Lambda\) with section \(V:\Lambda\to O^t\),
extended rationally, and put \(G=U^\dagger U\). Its entries
are integral by self-duality. View \(U\) as an \(n\)-by-\(t\)
matrix in \(E^n\). Then \(UV=I_n\), \(V^\dagger G V=I_n\),
and \(NV\) has integral entries. The map \(a\) factors as \(NV\)
followed by \(U\). Pull the integral germ back by \(U\); its raw
coefficient in an integral frame of the line for \(G\) is integral.
In the corrected frame the chain rule for \(NV\) gives \(N^d\)
times the differential operator \(P_V\) obtained from \(P\)
by sending a column-direction \(v\) of type \(\tau\) to
\(\tau(V)v\). Indeed in logarithm coordinates the map is
linear on each eigenspace. The corrected biextension frame
multiplies in both arguments and intertwines the adjoint
\(O\)-actions by its defining identities. Expanding the matrix
entries in these identities identifies its pullback with the
corrected frame for \((NV)^\dagger G(NV)=N^2 I_n\).
Thus the comparison differentiates the full section with the
same correction as in the original jet.

The coefficients of \(P_V\) are integral. The defining integral
polynomial ring is generated by the linear directions from
\(\Lambda\), and \(V\Lambda\subset O^t\).
\src{src/Formalism/LocalToGlobal/ProjectiveLattices/ThShort.lean}{235}{Integral polynomial directions of a nonstandard lattice}.
For the Hermitian matrix \(G\) the quadratic ambiguity in the
corrected log frame comes from summing
\(h(X_j,w_0G_{jk}X_k)\) over its columns. For a conjugate pair
of types its contraction on \(P_V\) uses the coefficients
\(\bar\tau(G_{jk})\). By the ordinary chain rule and
\(V^\dagger G V=I_n\), it is the image under \(V\) of the standard
contraction on \(P\); hence it vanishes. All other terms obey
Lemma~\ref{intjet:der} because the entries of \(G\) and \(w_0\)
are integral. This costs at most \(\lfloor d/2\rfloor v(\Delta)\).
The factor \(N^d\) proves the stated inequality.
\end{proof}

\paragraph{The norm of the generator at semiabelian reduction.}
Let \(\mathcal A/R\) be the semiabelian model in an integral test,
with abelian generic fibre \(A\), rigidified cubical line
\(\mathcal L\), integral section \(s_0\), and integral tangent
vectors. Its principal polarization is \(\lambda\). Write
\(\mathcal P\) for the rigidified Poincaré line of \(A\times A^t\).
In this test \(L=(1,\lambda w_0)^*\mathcal P\), with
\(w_0+\bar w_0=1\) and \(\bar w_0\) the Rosati adjoint of \(w_0\).
Indeed these are the graph-line and adjunction conditions of the
datum. The graph identification can be made rigidified by its value
at zero. A rigidified isomorphism of cubical lines on an abelian
variety respects the cube structure by uniqueness.
\src{src/ReciprocityTheory/ReductiveGroups/ArithmeticGeometry/BirationalGeometry/ToroidalCompactification/TwoWeightEqOneAddRoot.lean}{82}{Graph-line condition};
\src{src/Formalism/LocalToGlobal/ProjectiveLattices/VAtVNeBot.lean}{53}{The integral model, eigendirections and theta generator}.

\begin{lemma}[The maximum of the canonical theta norm]\label{intjet:maximum}
On the analytification over the completed algebraic closure, the
canonical norm of \(s_0\) is at most one everywhere. Its maximum
is attained at the point of the skeleton belonging to the identity
component of the special fibre.
\end{lemma}
\begin{proof}
Use the valuation \(v=-\log|\cdot|\) with normalized uniformizer.
One may make a valued extension for this calculation. For the split
Raynaud uniformization let \(Y\) be the period lattice,
\(\Phi:Y\longrightarrow X\) its identification with the character
lattice under the principal polarization, and \(b\) the positive
definite monodromy form on \(Y_{\mathbb R}\).
For any rigidified ample line \(H\) giving \(\lambda\) take its
descent data on the Raynaud extension. Denote the constant of descent
at \(u\in Y\) by \(c_H(u)\) and the period pairing by \(t\), using
the convention
\[
 c_H(u+v)=c_H(u)c_H(v)t(u,\Phi v),\qquad
 -\log\|t(u,\Phi v)\|=b(u,v).
\]
Products in this identity use the biextension identifications of
the fibres. The Poincaré line pulled back by \(1\times\lambda\)
is \(m^*H\otimes p_1^*H^{-1}\otimes p_2^*H^{-1}\), with its
rigidification. Pulling back by the graph of \(w_0\) and taking
the quotient of the three descent constants therefore gives
\[
 c_L(u)=\frac{c_H(u+w_0u)}{c_H(u)c_H(w_0u)}
       =t(u,\Phi w_0u),\qquad
 -\log\|c_L(u)\|=b(u,w_0u)=\tfrac12 b(u,u).
\]
The last equality uses \(b(w_0u,v)=b(u,\bar w_0v)\), by functoriality
of the dual period pairing, and symmetry of \(b\).

Here are the precise standard uniformization results used to pass
from this computation to a metric. The balanced translate in
de Jong--Shokrieh, Proposition~8.3, has descent valuation
\(b(u,u)/2\). Translation by a point of tropical coordinate \(z\)
adds \(b(u,z)\) to the descent valuation (Lemma~8.1).
Thus the balancing point for the present triple has \(z=0\).
It acts trivially on the skeleton. The symmetry of both lines makes
its image a two-torsion point, since \(T_y^*L\simeq T_{-y}^*L\)
implies \(\lambda(2y)=0\). Their canonical metrics correspond under
translation up to a constant from rigidification (Lemma~5.1).
For the balanced translate Proposition~9.1 and the tropical
Riemann theta formula give on \(Y_{\mathbb R}/Y\)
\[
        -\log\|s_0(\nu)\|=\mathrm{constant}
                      +\tfrac12\min_{u\in Y} b(u+\nu,u+\nu).
\]
Its minimum occurs at \(\nu=0\bmod Y\). Corollary~10.4 gives
the same maximum of the norm over the whole analytification.
Proposition~10.3 at the identity component identifies the negative
log norm there with the order of the original section along that
component in its cubical model. This order is nonnegative by
integrality. This last application uses the original rigidification,
so no unspecified metric constant enters the inequality.
See de Jong and Shokrieh, \S\S6--11~\cite{deJongShokrieh}, and the
Raynaud theta descent in Foster, Rabinoff, Shokrieh and Soto,
\S\S3--4~\cite{fosterTheta}.
\end{proof}

\begin{lemma}[Integral traces on identity discs]\label{intjet:disc}
Let \(i:C_0\to C_1\) be an isogeny of abelian generic fibres induced
by an inclusion of full projective \(O_E\)-lattices tensored with
\(\mathcal A\). Write \(\mathfrak d R\) for the determinant ideal
of its differential on the integral Lie algebras. Let \(L_1\)
be a symmetric cubical line on \(C_1\) with its canonical metric
and integral cubical extension to the identity model. Suppose
\(u\) is a section of \(i^*L_1\) bounded by one everywhere
in the pullback metric. Write \(T=\operatorname{Tr}_i u\).
Take \(b\in C_1[h]\), where \(h\) is odd and prime to the residue
characteristic, choose \(c\) with \(hc=b\) and \(h^2c=0\),
and suppose the normalized translation at \(c\) exists on
\([h]^*L_1\). On the formal disc at zero the raw coefficient of
\(\widetilde T_c[h]^*T\) in an integral frame of \([h]^*L_1\)
is divisible by \(\mathfrak d R\).
\end{lemma}
\begin{proof}
Extend the ground field finitely to contain the points of the
relevant finite fibres. The homomorphism of completed identity
discs of the two models is finite flat. For details, each smooth
disc has ring \(R[[X_1,\ldots,X_g]]\). On special fibres the
kernel of \(i\) is finite since an integer multiplication factors
through \(i\) and acts as an isogeny of semiabelian varieties.
Consequently the images of target parameters and a uniformizer
generate an ideal primary to the maximal ideal of the source.
Lifting a basis of the finite quotient by this ideal and iterating
in the ideal gives finite generation over the complete target ring.
These \(g+1\) parameters form a regular sequence in the regular
local source. The local criterion for flatness, successively along
the parameters, gives a finite free module. This finite, finitely
presented flat map is a local complete intersection at the closed
point: its closed fibre is cut out by the regular sequence of \(g\)
specialized parameters in the smooth special-fibre disc. The complete
intersection locus is open, hence contains the whole local spectrum.
The finite module of relative differentials is computed by the
formal Jacobian. Indeed derivations into finite modules are
continuous for the maximal-ideal topology by the product rule and
these modules are complete. In integral translation-invariant
differential bases the Jacobian determinant has constant ideal
\(\mathfrak d R\). Thus the trace of any power series in the source
lies in \(\mathfrak d R[[Y]]\) by the different formula in
Theorem~\ref{intjet:trace}.

The rigid generic fibre of either completed local ring is its open
unit polydisc. The disc map just constructed is surjective there:
the finite free map has positive rank, and points lying above a
point reducing to zero reduce to zero in this local map. More
explicitly the primary-ideal identity gives, for some \(e\),
\(X_j^e\in(\pi,i_1(X),\ldots,i_g(X))\); thus every point of the
finite formal fibre above such a point has \(|X_j|<1\). Finite
base change, or evaluation first on closed subdiscs, gives the
analytic covering. It is étale in characteristic zero, and each
fibre is a coset of the subgroup \(K^\circ\subset\ker i\) lying in
the identity disc. The full preimage over that disc in the abelian
generic fibre is the disjoint union of its translates indexed by
\(\ker i/K^\circ\). Indeed choose a disc preimage of a given
point. All other preimages differ by elements of \(\ker i\), and
two translated discs meet precisely when their translates differ by
an element of \(K^\circ\).

An integral frame of a symmetric cubical line on the identity disc
has canonical norm one. To see the metric comparison, on the generic
fibre of the formal completion along the entire special fibre take
the model metric. This domain is covered by finitely many compact
affinoids and is stable under doubling. The rigidified cubical
isomorphism for \([2]\) implies that the logarithm \(F\) of the
ratio of canonical and model metrics satisfies \(F(2x)=4F(x)\).
It is bounded on this domain, so iteration gives \(F=0\).
An integral formal frame differing from a model frame by a power
series with unit constant term has the same norm on the identity
disc. The canonical metric is functorial for rigidified isogeny
pullbacks by the same doubling argument on the compact generic
fibre.

A finite-order lift \(U\) of translation at any torsion point
preserves the canonical metric on its symmetric line. Choose
\(m>1\) congruent to one modulo the order of the point. Translation
commutes with \([m]\); comparison of the two line isomorphisms
on the proper connected generic fibre differs by a scalar.
The continuous logarithmic distortion therefore satisfies
\(F([m]x)=m^2F(x)+\mathrm{constant}\). Iteration on the compact
space makes \(F\) constant, and the finite order of \(U\) makes it
zero. This also applies to the normalized lift at \(c\).

In the identity frame on the target, view the translated pullback
as the trace above \(y\mapsto b+[h]y\), using that lift to identify
the line. Multiplication by \(h\) is an isomorphism on either
identity disc, by the integral formal inverse function theorem.
For each class in \(i^{-1}(b)/K^\circ\), choose a representative
\(x\). The preimages on its sheet are parametrized by
\(x+[h]z\), with \(i(z)=y\). On this disc the coefficient of
\(u(x+[h]z)\) in the pulled-back translated frame is analytic of
absolute value at most one, by the metric calculation. Every such
function on the open unit polydisc has an integral power series:
apply the Gauss norm bound on closed polydiscs of radii tending to
one, after a valued extension if needed. The coefficients already
belong to the field of definition. The sum over this sheet is the
trace of that power series under the finite free formal map \(i\).
On geometric generic fibres this is the usual sum over the sheet,
which identifies the two analytic traces. Each is divisible by
\(\mathfrak d R\), so their finite sum has the same divisibility.
\end{proof}

\begin{theorem}[Finite denominators at semiabelian tests]\label{intjet:semi}
At every integral test with the stated semiabelian, cubical and
Hodge data, \(D_M^2\) clears the rescaled finite theta sum. The
valuation inequalities of Proposition~\ref{intjet:standard} at
locally standard lattices and of Theorem~\ref{intjet:trace} away
from \(M\) hold also at semiabelian reduction.
\end{theorem}
\begin{proof}
At a locally standard lattice use \(N_0\) and the generic-fibre
identity of Proposition~\ref{intjet:standard}. The norm of the
generator is at most one everywhere by Lemma~\ref{intjet:maximum}.
Its normalized pullback by \([N_0]\) and its two finite-order
translates have the same bound, by functoriality and the isometry
proved in Lemma~\ref{intjet:disc}. In the integral formal frame
at zero their raw germs therefore have integral coefficients by
the polydisc argument. Lemma~\ref{intjet:der} applies on a smooth
semiabelian formal group verbatim: it uses the formal group law,
the integral tangent directions and the cubical adjunction.
This proves the individual inequality, also when \(\ell\mid M\).

For \(\ell\nmid M\) use the notations \(a_0,i,j,a\) in
Theorem~\ref{intjet:trace}. The tensor models are direct summands
of products of the smooth semiabelian scheme. An idempotent on a
smooth group gives a product decomposition by its image and the
image of its complement: these are the closed kernels of the
opposite projections, and addition and the pair of projections
are mutual inverses. Each factor is smooth by the lifting
criterion and has semiabelian fibres as a direct factor of such a
group. The graph-line construction on these models uses only the
integral cubical biextension and its \(O_E\)-adjunction. On the
generic fibre the identities of pullback lines and normalized
translations give, after summing the first translation,
\(a^*\operatorname{Tr}_i(j^*s_0^{\boxtimes n})\), as before.
The section inside this trace has norm at most one everywhere,
since the rigidified line identifications are isometries.

The final point \(t\) has odd order \(h\) prime to \(\ell\).
Take \(t'\) with \(ht'=t\), and set \(c=a(t')\), \(b=a(t)\).
The normalized lift at \(c\) on \([h]^*\mathcal L_1\) exists:
\(h^2c=0\) and \([h]^*\mathcal L_1=\mathcal L_1^{h^2}\).
Set \(T=\operatorname{Tr}_i(j^*s_0^{\boxtimes n})\).
The sections
\[
       [h]^*\widetilde T_t\,a^*T,\qquad
           a^*\widetilde T_c\,[h]^*T
\]
agree up to a root of unity. Both arise from the same translation
\(t'\) on \([h]^*a^*\mathcal L_1\), by pulling back the respective
lifts and using \([h]T_{t'}=T_t[h]\) and \(aT_{t'}=T_ca\).
These two lifts have finite order: a power divisible by the point
order and the order of the original lift is the identity. Their
ratio is a scalar on the proper generic fibre, hence is a root
of unity. Lemma~\ref{intjet:disc} gives divisibility by
\(\mathfrak b R\) for the section on the right before applying
\(a^*\), since the determinant on integral Lie algebras has
exactly the same ideal computed in Theorem~\ref{intjet:trace}.
Multiplication by \(s\in\mathfrak b^{-1}\) gives an integral raw
germ. Apply the integral polynomial and harmonic chain rule for
\(a=U(NV)\) in that theorem. The extra pullback \([h]\) on the
left multiplies a homogeneous corrected jet by \(h^d\), a unit.
The factor \(N^d\) and the loss of at most
\(\lfloor d/2\rfloor v(\Delta)\) follow as claimed. The index
cancellation and the level factor \(D_M=C_M(M\Delta)^d\) finish
the original finite-place test.
\end{proof}


The size of the multiplier itself is unconditional under the stage
bounds \(\log M\le30L_0,\ \log\delta_E\le7fL_0\),
\(f,r,L_0\ge1\). Indeed \(frn+w\le2S,\ w\le S\), so
\[
 \log D_M^2
 =2(frn+w)\log M+2w\log\delta_E
 \le134 f L_0 S\le134(frL_0)^2 S .
\]
This bound is independent of \(N,\Lambda,P,x_0\) for the fixed
weight and level. Theorem~\ref{intjet:semi} supplies the valuation
test for the jets with the scalar \(s\) allowed by the determinant
ideals. In the source the proof uses valuations of theta
derivatives in integral semiabelian charts and a local isogeny on
the primary torsion subgroup.
\src{src/Development/AbelianVarietyTheory/Proofs/DivisorProofs/LogShortThetaSeriesDenominatorBoundLe.lean}{102}{Size of the common clearing integer};
\src{src/Development/AbelianVarietyTheory/Proofs/DivisorProofs/SmulIntZSqNonnegOfSqDvd.lean}{71}{Normalized jet sum at a square divisor of the level};
\src{src/Development/AbelianVarietyTheory/Proofs/RingsOfIntegers/SumHvPrimaryKernel.lean}{221}{Local isogeny and integral formal theta jets}.

\paragraph{An archimedean estimate on each complex theta fibre.}
Fix a complex embedding \(\sigma\) of the coefficient field \(F\),
of degree \(h\) over \(\mathbf Q\). Use the cubical metrics
rigidified to have norm one at zero on the complex graph lines and
Haar probability on each abelian variety. Put \(g=fr=\dim A\),
\(m=gn\), \(d=\deg P\), and \(\Delta=|\operatorname{disc}E|\).
Let \(\lambda>0\) be the shortest nonzero period of \(A\) in the
positive Hermitian metric \(H\) of its principal polarization,
normalized by the Riemann form \(\Im H\). Take an orthogonal Hodge
eigenframe whose individual direction vectors have length at most
\(U\), with \(U>0\). Use the row copies of that frame in \(A^n\).
Let \(\|P^\sigma\|_1\) denote the sum of absolute values of its
coefficients in these directions. The integer \(N\) and the self-dual
Hermitian column lattice \(\Lambda\) can vary at this fixed frame.

\begin{proposition}[A complex fibre bound independent of the auxiliary level]
\label{intjet:arch}
With the rescaling \(s\) chosen above by Minkowski's theorem, put
\[
 C=2fn\Delta,\qquad
 R=\frac{\min(1,\lambda/C)}{4\sqrt m}.
\]
Then the cubically normalized value of the rescaled theta series satisfies
\[
 \frac{|\sigma(sy)(x_0)|}{\|s_0\|_2^n}
 \le |\Delta_F|^{1/(2h)}e^{\pi/2}
       \|P^\sigma\|_1\,d!\,(\sqrt2\,U)^d R^{-(d+m)}.
\]
This estimate is uniform in \(N,\Lambda\), the choice of division
points and the level character.
\end{proposition}
\begin{proof}
Use the isogeny \(a:B=A^n\to B_1=\Lambda\otimes_O A\) from
Theorem~\ref{intjet:trace}; in rational column coordinates its
differential is \(N\) times the identity. Both graph lines on
\(B\) and \(B_1\) induce principal polarizations. For the latter
this follows by tensoring the principal polarization of \(A\) with
the isomorphism from \(\Lambda\) to its conjugate \(O\)-dual
given by the unimodular Hermitian pairing. The equality
\(a^*\mathcal L_1=[N]^*\mathcal L\) is isometric: the two
curvature forms agree and the ratio of metrics on a compact complex
torus is then constant, fixed at zero. A finite-order normalized
translation lift is also an isometry, since the ratio is constant
and a positive real number of finite multiplicative order is one.
All the \(L^2\) pullbacks under isogenies preserve norms because
the pushforward of Haar probability is Haar probability.

Write \(v=[N]^*s_0^{\boxtimes n}\), and \(U_z\) for the normalized
translation. The vectors \(U_zv\), \(z\in Z_\Lambda/K_0\), are
pairwise orthogonal. Indeed \(v\) is fixed by the deck translations
\(b\in B[N]\), and \(U_zv\) is then an eigenvector for each
\(U_b\), with character given by the theta commutator with \(z\).
The annihilator of \(B[N]\) inside \(B[N^2]\) for this commutator
is \(B[N]\) itself. In a symplectic basis for the principal
polarization, the pairing is the standard perfect alternating
pairing modulo \(N^2\), and \(B[N]\) consists of all multiples of
\(N\). Distinct classes modulo \(K_0=Z_\Lambda\cap B[N]\) therefore
give distinct characters of commuting unitaries. The trace identity
of Theorem~\ref{intjet:trace} gives
\[
 a^*T=\sum_{z\in Z_\Lambda/K_0}U_zv,\qquad
 \|T\|_2= I^{r/2}\|s_0\|_2^n,\qquad I=[\Lambda:\Lambda_0].
\]
In particular \(\|\sigma(s)T\|_2\le
|\Delta_F|^{1/(2h)}\|s_0\|_2^n\). This equality of squared norms
uses the original order of the translations. The \(t_u\) act
after this calculation.

We bound the shortest period on \(B_1\). On \(\Lambda\) take the
real trace norm \(\|v\|_{\rm tr}^2=
\operatorname{Tr}_{E/\mathbf Q}\langle v,v\rangle\).
For \(v\ne0\) its squared norm is at least one, since the trace
is a positive rational integer. The trace-dual lattice equals
\(\mathfrak D_E^{-1}\Lambda\), by \(O\)-duality and the definition
of the codifferent. Multiplication by \(\Delta\) sends it into
\(\Lambda\), so its minimum is at least \(1/\Delta\).
Euclidean transference in dimension \(2fn\) gives \(2fn\) independent
vectors of \(\Lambda\) of trace norm at most \(C\)
\cite[Theorem 2.1]{banaszczyk}. Contractions with those vectors
separate every nonzero \(x\in\Lambda\otimes_O\Pi\), where
\(\Pi\) is the period lattice of \(A\): their rational span is the
whole column space over \(\mathbf Q\), and the pairing is perfect
over \(E\). Each contraction belongs to \(\Pi\). Its operator norm
on the complex tangent space with product Hermitian norm is at
most \(C\), by Cauchy--Schwarz on each eigenblock; on a given block
the sum of the squared column coefficients is at most the squared
trace norm. The positive form on \(B_1\) in rational column
coordinates is this same product form by the isometry of pullbacks
and \(da=N\,\mathrm{id}\). Thus every nonzero period on \(B_1\)
has norm at least \(\lambda/C\).

Use orthonormal eigen-coordinates on the universal cover of \(B_1\)
centered at any prescribed point. The polydisc of coordinate radius
\(R\) injects, by the period bound. In it a holomorphic line frame
of squared norm \(\exp(-\pi\|z\|^2)\) exists and is unique up to
unit phase. Indeed its prescribed logarithmic norm has the required
curvature, and a real pluriharmonic discrepancy on the simply
connected chart is the real part of a holomorphic function.
The Euclidean covolume of the periods in these coordinates is one:
the real volume form is \((\Im H_1)^m/m!\) up to orientation,
and its period integral for a principal polarization is one.
The same holds for any translate of the line and any choice of
center. If \(G=\sum c_\alpha z^\alpha\) is the coefficient of a
section of \(L^2\)-norm \(D\) in such a frame, then
\[
 |c_\alpha|^2\,
 \frac{\pi^mR^{2|\alpha|+2m}}{\prod_i(\alpha_i+1)}
 \le \int_{\{|z_i|<R\}}|G|^2\,dz\le e^\pi D^2 .
\]
The first inequality is angular orthogonality followed by radial
integration. The second uses \(mR^2\le1\) and injectivity. Since
\(\prod(\alpha_i+1)\le2^{|\alpha|}\), derivatives of degree \(d\)
in the row frame are bounded by
\(e^{\pi/2}D\,d!(\sqrt2 U)^d R^{-d-m}\).
Here each row direction is a multiple of an orthonormal coordinate
direction. Summing with \(P\) introduces \(\|P^\sigma\|_1\).

To compare this local derivative with the harmonic jet, on the
vector-space cover choose the holomorphic frame of the polarized
biextension with squared norm
\(\exp[-\pi(H(z,w)+H(w,z))]\), normalized at the origin.
Its curvature is that of the polarized biextension, by taking the
pullback under addition and dividing by the two projections.
Addition in either variable and Rosati adjunction preserve this
frame: their metric factors agree by sesquilinearity, and the ratio
is holomorphic of unit modulus and is one on the axes. On the graph
\(w=w_0z\) its squared norm is \(\exp(-\pi H(z,z))\) since
\(w_0+w_0^\dagger=1\). Make this frame comparison first on the
single factor \(A\times A\). Any permitted additive frame there
differs by the exponential of a Rosati-compatible complex bilinear
form, hence only couples conjugate types. On \(B=A^n\) the jet
uses the exterior product of the same single-factor frame in all
columns, followed by its \(N^2\)-th tensor power. Its graph correction
is therefore a sum of scalar multiples of
\(\sum_b z_{\tau i b}z_{\bar\tau j b}\). Each such sum disappears
under the harmonic equations, as in Lemma~\ref{intjet:der}.
For the final translation
by \(t_u\), apply the centered frame at \(a(t_u)\) to the translated
germ. Its pullback has squared frame norm
\(\exp(-\pi N^2\|z\|^2)\); the isometry \(U_{t_u}\) identifies it
with the frame at zero up to unit phase, by uniqueness. The derivative
then gains precisely \(N^d\), which cancels the normalizing
denominator in the series. Finally average over \(u\). Its Fourier
weights have modulus one and the divisor is their number \(C_M\).
This proves the inequality.
\end{proof}

Write \(S=d+m/2\) and \(Z_0=frL_0\). Suppose the frame satisfies
\[
 \max\{\log^+(\lambda^{-1}),\log^+ U,\log^+\|s_0\|_2\}\ll Z_0^2,
 \qquad \log^+\|P^\sigma\|_1\ll Z_0^2S .
\]
Then the log of the absolute bound is \(O(Z_0^2 S)\) under the
stage bounds above.
Here \(\log^+x=\max(0,\log x)\) for \(x>0\); a zero polynomial is
immediate. In fact \(m\le2S\), \(d\le S\),
\(\log^+(d!)\le d\log(1+d)\), \(\log(2S)\ll L_0\),
\(\log\Delta\le7fL_0\) and \(h^{-1}\log|\Delta_F|\ll Z_0^2\).
These bound every factor in the displayed inequality, including
\(\log(1/R)\ll Z_0^2\). To pass from this complex fibre estimate
to the simultaneous section norms requires the normalized algebraic
and analytic readings on every component, short polynomials spanning
the actual space, and the indicated frame control at its half
maxima. The normalized norm of the chosen generator matters when
that generator is identified with a frame of the determinant line.
\src{src/Development/AbelianVarietyTheory/Proofs/CuspFormEvaluations/NormalisedThetaValueBound.lean}{146}{Normalization on the tensor torus};
\src{src/Development/AbelianVarietyTheory/Proofs/DivisorProofs/NormRawThetaSeriesIndexBoundLeSqrt.lean}{18}{Square-root index in the archimedean estimate}.


\paragraph{The component Bergman bound and its volume normalization.}
Fix one good complex embedding of the coefficient field. Write
\(f=[E:\mathbb Q]/2\), \(\Delta=|\operatorname{disc}E|\),
\(R_E=\Delta^{1/(2f)}\), \(r\ge2\) for the Hermitian rank,
\(d\) for the positive host integer and \(M\) for the deep level.
Let \(w\ge1\) bound the sum, over the CM type, of the spreads
\(\max_i\lambda_{\tau i}-\min_i\lambda_{\tau i}\) of the flag weights.
The complex dimension of the product of full flag domains is
\(b=fr(r-1)/2\). The constants in the component estimate are
\begin{align*}
 \rho&=[200r^3(w+f)]^{-1},&
 \varepsilon&=\min\{1,\pi/[2M\sqrt d\,(fw+2)\log2]\},\\
 A&=(1+2(2fr)^2\Delta/\varepsilon)^{2fr^2},&
 V_*&=(15M)^{2fr^2}[3(2f)^2r^4R_E^2]^{6fr^2},\\
 C_{\rm loc}&=\frac{8\,b!}{(\pi\rho^2)^b}\,A,&
 \beta_*&=C_{\rm loc}V_* .
\end{align*}
Here \(V_*\) is the upper bound called \(\mathrm{massLev2}\) in
the Tamagawa comparison. The normalizing mass itself is the sum
of the volumes of the chosen arithmetic fundamental sets.
\src{src/Development/AbelianVarietyTheory/Definitions/Divisors/UnifOfSpaceEq.lean}{103}{Flag radius, return bound and local factor};
\src{src/Development/WeilBound/RepresentationTheoryGL/TubeMeasures/DisjointTubes/MassLev2.lean}{13}{Level-dependent total mass bound};
\src{src/Development/WeilBound/RepresentationTheoryGL/TubeMeasures/DisjointTubes/ConstantSupremum/SupremumConstantLogBound.lean}{14}{Product defining the normalized supremum constant}.

\begin{lemma}[Evaluation on a flag ball]\label{bergvol:local}
At an orthonormal frame \(x\), let \(B_x\) be the lower unipotent
flag chart of Euclidean radius \(\rho\). Write \(d\mu\) for the flag
measure specified below. For every holomorphic flag section \(P\),
\[
        |P(x)|_{\rm Pet}^2\le
        \frac{2b!}{(\pi\rho^2)^b}\int_{B_x}|P|_{\rm Pet}^2\,d\mu .
\]
\end{lemma}
\begin{proof}
In the frame \(x\) each factor has a diagonal Hermitian matrix
\(J\), with the positive signs first and the negative signs last.
Write the chart matrix as \(U=I+Z\) with \(Z\) strictly lower
triangular; the entries over all factors are the Euclidean coordinates.
In particular \(\|Z\|_{\rm op}\le\|Z\|_{\rm HS}\le\rho\).
Let \(D_{\tau i}\) be the absolute determinant of the leading
\(i\)-by-\(i\) Gram block of \(U^*JU\). The endpoint determinants
\(D_{\tau0},D_{\tau r}\) equal one. On the chart the squared norm
and invariant measure in question are
\[
 |P|_{\rm Pet}^2=|P(xU)|^2
 \prod_{\tau}\prod_{i=0}^{r-1} (D_{\tau,i+1}/D_{\tau i})^{-\lambda_{\tau i}},
 \qquad
 d\mu=\left(\prod_{\tau}\prod_{i=0}^{r-1}D_{\tau i}^{-2}\right)dV .
\]
The difference of each Gram block from the corresponding sign matrix
has operator norm at most \(2\rho+\rho^2\le3\rho\). Its singular
values therefore lie in \([1-3\rho,1+3\rho]\), and the block has
the required inertia by continuity. Thus the whole ball belongs to
the flag domain and \(|\log D_{\tau i}|\le6r\rho\).
Subtract the midpoint of the extreme weights from each row of
weights at \(\tau\). This leaves the logarithm of the norm factor
unchanged by telescoping the endpoint determinants. Its absolute value
is at most \(6r^2\rho w\le1/20\); the absolute logarithm of the
density is at most \(12fr^2\rho\le1/10\). At the origin both
factors equal one. The holomorphic mean-value inequality on the
Euclidean ball gives the assertion since its volume is
\((\pi\rho^2)^b/b!\) and \(\exp(1/20+1/10)<2\).
\end{proof}

\begin{lemma}[Overlaps at a frame with no short lattice vector]\label{bergvol:returns}
Suppose the arithmetic group preserves a Hermitian \(O_E\)-lattice
\(\Lambda'\subset E^r\) self-dual for the \(O_E\)-valued form.
Let \(\|v\|_x^2\) be the sum of the positive majorants of that form
over the fixed CM type, obtained by changing the negative signs to
positive in \(x\). If its minimum on \(\Lambda'\setminus\{0\}\)
is at least \(\varepsilon\), at most \(A\) arithmetic elements
return \(B_x\) to itself, meaning \(\gamma B_x\cap B_x\ne\varnothing\).
\end{lemma}
\begin{proof}
Set \(m=2fr\), the real lattice dimension. The Euclidean dual of the
embedded lattice is \(2J\) times its dual for
\(\operatorname{Tr}_{E/\mathbb Q}h\), since that trace is twice
the real inner product after applying the sign involution \(J\).
The trace dual lies in \(\Delta^{-1}\Lambda'\): testing against
\(O_E\Lambda'\) puts each Hermitian pairing in the codifferent,
which is carried into \(O_E\) by \(\Delta\), as seen by inverting
the integral trace matrix. Hence the Euclidean dual minimum in length
is at least \(2\sqrt\varepsilon/\Delta\).
The Euclidean transference inequality
\(\lambda_m(L)\lambda_1(L^*)\le m\) applies to this full rank lattice.
It gives \(m\) independent vectors of length at most
\(m\Delta/(2\sqrt\varepsilon)\). They span \(E^r\) over \(E\):
otherwise their real span would have dimension less than \(m\).
Select an \(E\)-basis from them.
Banaszczyk,
\emph{New bounds in some transference theorems in the geometry of numbers},
\S2~\cite{banaszczyk}.

In the chart the positive subspace is the graph of a matrix \(T\)
with \(\|T\|\le\rho/(1-\rho)\), by inverting its upper block.
Its majorant is between one half and twice the central squared norm.
Indeed singular value decomposition reduces its change to independent
positive-negative planes with graph slope \(t\le\|T\|\); the
eigenvalues of the new majorant there are
\((1+t)/(1-t)\) and its inverse, by using the orthogonal lines
spanned by \((1,t)\) and \((t,1)\). Unpaired directions retain
eigenvalue one. A returning \(\gamma\) carries one chart majorant
isometrically onto another; hence \(\|\gamma v\|_x\le2\|v\|_x\).
All images of the selected basis vectors have length at most
\(m^2\Delta/\sqrt\varepsilon\), allowing the extra factor \(m\ge1\).
Lattice points in a ball of radius \(T_0\) number at most
\((1+2T_0/\sqrt\varepsilon)^m\), by packing disjoint open balls of
radius \(\sqrt\varepsilon/2\). The \(r\) basis images determine
\(\gamma\). Taking the \(r\)-th power of this count yields \(A\).
\end{proof}

\begin{proposition}[From cusp and mass data to a normalized kernel]\label{bergvol:global}
Assume the following on every analytic flag component at the chosen
embedding.
\begin{enumerate}
\item The automorphic sections under consideration are bounded in
Petersson norm. For each nonzero such cusp section there is an
orthonormal frame \(x\) where its norm is at least half the component
supremum and the minimum in Lemma~\ref{bergvol:returns} is at least
\(\varepsilon\), for the same \(\Gamma\)-stable self-dual lattice.
\item There are invariant flag measures in the chart normalization
of Lemma~\ref{bergvol:local} and measurable arithmetic fundamental
sets \(F_j\). With \(V=\sum_j\mu_j(F_j)\), assume \(0<V\le V_*\).
The pushforward of the adelic probability measure on arithmetic
and frame-invariant functions is \(\sum_j d\mu_j|_{F_j}/V\).
\end{enumerate}
Then for each cusp section under consideration the global squared supremum
is bounded by \(\beta_*\) times its squared Petersson norm for that probability. On any nonempty
algebraic component \(Z\) which corresponds to a union of these
analytic components of positive mass, the same bound holds with
the conditional probability on \(Z\). In particular the scalar
evaluation Bergman kernel on \(Z\) is at most \(\beta_*\).
\end{proposition}
\begin{proof}
At the chosen half-maximum frame, project the flag ball onto the
arithmetic quotient. Each fiber in the ball has at most \(A\)
points, since mapping one of them to the others requires elements
which return the ball. Equivalently, cut a measurable fundamental
domain into its arithmetic translates and integrate its invariant
integrand there; the number of copies of a point in the ball is
bounded by the return count, including stabilizers if needed.
Lemma~\ref{bergvol:local} and the half-maximum then give
\(\sup_{x\in F_j}|P(x)|_{\rm Pet}^2\le C_{\rm loc}
\int_{F_j}|P|_{\rm Pet}^2\,d\mu_j\).
Adding the nonnegative integrals and inserting \(V\) proves the
global result. For \(Z\) use the identical calculation summing only
over its analytic components and multiplying the conditional
integral by \(V_Z=\sum_{j\subset Z}\mu_j(F_j)\le V_*\).
Taking the supremum of the evaluation ratio over sections gives
the Bergman bound. Its integral against the conditional probability
bounds the rank as well: in any finite orthonormal system the sum
of the squared pointwise norms equals the squared norm of scalar
evaluation on its span, and each summand integrates to one.
\end{proof}

For clarity, the half-maximum hypothesis admits a quantitative test
which isolates the cusp input. Put \(H=2M\sqrt d\). If at a
half-maximum frame a lattice vector satisfies \(\|v\|_x^2<
\varepsilon\le1\), every conjugate of the algebraic integer
\(h(v,v)\) has absolute value less than one, since it is bounded
by the corresponding majorant. Thus \(h(v,v)=0\) by its norm.
For each such isotropic vector use the flow parameter scaled by
\(k_0=2|\tau(\xi_0)|\), independent of \(\tau\).
In this parameter the input needed is a transvection
period \(0<H'\le H\), acting with trivial scalar multiplier, and its
holomorphic one-variable cusp restriction \(G(\zeta)\). In the
normalization with scalar metric factor one at the center it
suffices that \(G(0)=P(x)\), that \(G\) be \(H'\)-periodic and
holomorphic on a neighborhood of \(\Im\zeta\ge-t\) for
\(t=1/(2\varepsilon)\), and that it tend uniformly to zero as
\(\Im\zeta\to+\infty\). The flag metric on the flow gives
\[
       |G(\zeta)|\le 2\cdot 2^{w/2}|G(0)|
                     \qquad(\Im\zeta=-t).
\]
To verify the metric assertion, write \(u\) for the coordinates of
\(\tau v\) at the orthonormal central frame, with positive columns
first. Since \(\tau\kappa_0=-i k_0\) on the chosen CM type, the
flow in the scaled parameter is \(U=I-i\zeta uu^*J\).
Isotropy gives \(u^*Ju=0\) and
\(U^*JU=J+2\Im\zeta\,Juu^*J\).
For \(S_i=\sum_{a=1}^i J_{aa}|u_a|^2\), the rank-one determinant
formula on each leading block now gives
\(D_{\tau i}=|1+2\Im\zeta\, S_i|\).
The \(S_i\) increase from zero to \(\|u\|^2/2\), then decrease to
zero, since the total positive and negative sums coincide.
Here \(\|u\|^2\) is the local majorant, less than
\(\varepsilon\) by hypothesis. All signed determinants retain their
central signs for \(\Im\zeta>-1/\varepsilon\). Thus the flow stays
inside the flag domain in this entire half-plane and the restriction
of a holomorphic section is holomorphic there.
At \(\Im\zeta=-1/(2\varepsilon)\), the \(D_{\tau i}\) decrease from
one to a number greater than \(1/2\), then increase to one.
The total variation of \(\log D_{\tau i}\) is at most \(2\log2\).
Centering the weights at their midrange in the squared metric formula
of Lemma~\ref{bergvol:local} bounds its absolute logarithm by
\(\sum_\tau\operatorname{spread}(\lambda_\tau)\log2\le w\log2\).
The boundary raw coefficient is therefore bounded by \(2^{w/2}\)
times its Petersson norm, which itself is at most \(2|G(0)|\).

For \(y=\Im\zeta\ge0\), the same variation argument gives a raw
coefficient of size at most \(2|G(0)|(1+\varepsilon y)^{w/2}\).
Thus for \(q=\exp(2\pi i\zeta/H')\) the growth is at most a polynomial
in \(\log(1/|q|)\). By Cauchy's estimate on circles shrinking to
zero, every negative Laurent coefficient vanishes. It suffices to
show that the constant Laurent coefficient vanishes, equivalently
that the average of this restriction along one real period is zero
at any fixed height.

For the analytic cusp space of this construction the constant is
indeed zero. Its definition imposes bounded Petersson norm and,
for every nonzero isotropic \(\xi\in\Lambda\) and every frame,
zero average of the coefficient on the following \(f\)-dimensional
real torus. Put \(F=E^+\); its period lattice in the raw flow
parameters is \(L=m e_\chi\,\iota(O_F)\), where \(\iota\) is the
real embedding vector, \(m\) is the integer transvection period of
the flag datum and \(e_\chi\) kills its finite order multiplier.
This is a full lattice by the nonzero discriminant of \(F\).
The real period in direction \(\iota(a)\) follows from the
rational transvection for \(a\in O_F\) and the multiplier is
trivial after its \(e_\chi\)-th power. Both boundedness and
zero average persist under complex linear combinations, so they
hold on the full span used for the analytic Bergman kernel.
\src{src/ReciprocityTheory/ReductiveGroups/ArithmeticGeometry/FundamentalDomains/ModerateCoefficientPositivity.lean}{58}{Definition of the analytic cusp space and its zero Fourier coefficient}.

Choose an integer \(c_1>0\) with \(\xi=c_1v\in\Lambda\), possible
since \(\Lambda\) is a full lattice in the rational vector space
underlying \(E^r\). Apply the flow in independent raw
parameters \(z_\tau\) to \(\xi\), and call the resulting holomorphic
coefficient \(A(z)\). For \(y_\tau=\Im z_\tau\ge0\), the same
principal minors are \(1+2 k_0 y_\tau S_i\), now with \(S_i\)
computed for \(\tau\xi\). Their variation and bounded Petersson norm
give
\[
 |A(t+iy)|\le C_P\prod_\tau
 (1+k_0 y_\tau\|\tau\xi\|_x^2)^{\operatorname{spread}(\lambda_\tau)/2}
 \qquad(t\in\mathbb R^f).
\]
Here the norm is the single-place majorant. In particular growth is
at most polynomial when one imaginary coordinate tends to infinity
and the others stay zero. For \(\eta\in L^*\) let \(a_\eta(y)\)
be the Fourier coefficient of \(t\mapsto A(t+iy)\) using
\(\exp(-2\pi i\eta\cdot t)\) and probability Haar measure on the
real torus. Holomorphy and integration by parts on this torus give
\(\partial_{y_j}a_\eta=-2\pi\eta_j a_\eta\), hence
\(a_\eta(y)=e^{-2\pi\eta\cdot y}a_\eta(0)\).
The polynomial bound implies \(a_\eta(0)=0\) if any \(\eta_j<0\).
The defining cusp average gives \(a_0(0)=0\).

The smooth periodic function at \(y=0\) has absolutely summable
Fourier coefficients: repeated integration by parts with
\(1-\sum_j\partial_{t_j}^2\) bounds them by
\(C_N(1+|\eta|^2)^{-N}\) for every \(N\).
At each \(y\ge0\) the resulting series represents \(A(t+iy)\)
by Fourier uniqueness. On the diagonal ray its nonzero modes
therefore tend uniformly to zero as the height tends to infinity,
by dominated convergence since \(\sum_j\eta_j>0\) on their support.
Finally \(G(\zeta)=A((\zeta/(k_0c_1^2))_\tau)\), since the
rank-one operator of \(c_1v\) is \(c_1^2\) times that of \(v\).
This proves the required uniform vanishing and the claimed
one-variable constant-term identity.

This exponential identifies the half-plane above
the boundary modulo
periods with a punctured disk up to radius \(e^{2\pi t/H'}\).
The induced function extends with value zero at \(q=0\). Apply
the maximum principle to its quotient by \(q\) and evaluate at
\(q=1\). This would give
\[
  1\le 2^{1+w/2}\exp(-2\pi t/H')
       \le 2^{1+w/2-(fw+2)}<1 .
\]
Thus the cusp test implies the required minimum. Its application
on the actual deep covers needs the same period bound on every
component for the lattice used in the return count.

The period has a direct lattice calculation once the integral model
is identified. In the host normalization put
\(c=\mathrm N(\mathfrak c_K)>0\); the purely imaginary \(\xi_0\)
satisfies \(|\tau(\xi_0)|=1/(c\sqrt d)\), and the transvection
coefficient \(\kappa_0=\pm2\xi_0\) satisfies
\((cd\kappa_0)^2=-4d\). Thus \(cd\kappa_0\) is integral. Suppose
the principal level is imposed on the lattice
\(\Lambda=\mathfrak a\Lambda'\), as in the actual comparison, with
\(\Lambda'\) Hermitian self-dual. For \(v\in\Lambda'\) isotropic
and \(y\in\Lambda\), \(h(y,v)v\in\Lambda\), using the convention
that \(h\) is linear in its first argument. The transvection
\[
        y\longmapsto y+Mcd\kappa_0\,h(y,v)v
\]
therefore preserves \(\Lambda\) and is the identity modulo \(M\).
It is unitary by \(\bar\kappa_0=-\kappa_0\) and \(h(v,v)=0\);
its inverse is obtained by changing the sign. Its raw period is
\(Mcd\), and scaling by \(k_0=2/(c\sqrt d)\) gives precisely
\(H=2M\sqrt d\). The same rational transformation realizes the
period simultaneously at every \(\tau\), since the imaginary
coefficient has the same negative imaginary part \(-k_0\) there.
For the even model weights the required arithmetic
character on these transvections must be trivial. The identification
of these lattices with the actual congruence groups on every deep
component, and the identification of the sections under test with
the stated analytic cusp space, remain necessary for this application.
\src{src/Development/AbelianVarietyTheory/Proofs/CubicalLineBundles/IsIntegralLevelMulKappa.lean}{35}{Integral transvection in the principal congruence subgroup}.

In the total mass comparison, the more precise intermediate upper
bound has the form
\[
   V\le 16^{fr^2} M^{2fr^2}\operatorname{vol}(U(r))^f
                            \operatorname{mass}_I .
\]
Here the compact group volume uses the real normalization in the
adelic tube comparison and \(\operatorname{mass}_I\) is the
definite standard lattice mass for the host. A bound
\(\operatorname{vol}(U(r))\le13^{r^2}\) and
\(\operatorname{mass}_I\le[3(2f)^2r^4R_E^2]^{6fr^2}\)
give \(V\le V_*\) since \(16\cdot13\le15^2\).
The principal-level index factor \(M^{2fr^2}\) itself follows by
reducing integral lattice automorphisms modulo \(M\). Their image
lies in the endomorphism ring of \(\Lambda/M\Lambda\), of order
\(M^{2fr^2}\) by local freeness of rank \(r\) over \(O_E\).
It remains to compare the adelic component tubes, with their disjointness
and volumes, and to identify the pushforward probability in the second
hypothesis. We calculate the compact factor and the definite mass below.
\src{src/ReciprocityTheory/ReductiveGroups/ArithmeticGeometry/FundamentalDomains/MaxBoundPf.lean}{12}{Isotropic flow and periodic cusp restriction};
\src{src/Development/WeilBound/RepresentationTheoryGL/TubeMeasures/DisjointTubes/MeasureSumBounds/SumMuFLe.lean}{110}{Principal level tube comparison for all components};
\src{src/Development/AutomorphicArithmetic/ThetaCorrespondence/FockModels/TamagawaPeterssonLambdaSqNonneg.lean}{371}{Restriction to normalized algebraic components}.

\paragraph{The compact real factor and the definite lattice mass.}
Put \(n=[E:\mathbb Q]=2f\), and write \(O=O_E\). The two numerical
factors in the preceding comparison have the following bounds.
\begin{proposition}\label{bergvol:compactmass}
For the real density specified below and the genus of the standard
positive definite Hermitian \(O\)-lattice of rank \(r\ge1\),
\[
 c_r=\operatorname{vol}(U(r))
     =\prod_{j=1}^r\frac{2\pi^j}{(j-1)!}\le13^{r^2},
 \qquad
 \operatorname{mass}_I
     \le [3n^2r^2R_E^2]^{2nr^2}
     \le[3n^2r^4R_E^2]^{3nr^2}.
\]
\end{proposition}
\begin{proof}
In the skew-Hermitian Lie algebra use the density \(dx\) given by
the imaginary diagonal entries and the real and imaginary parts of
the entries strictly below the diagonal. To compare it with the real
measure of the tube calculation, let \(\lambda\) be the additive
measure for which \((X,Y)\mapsto X+iY\) carries
\(\lambda\otimes\lambda\) to entrywise Lebesgue measure on
\(M_r(\mathbb C)\). Each diagonal block has absolute real determinant
one. Each off-diagonal pair gives the two complex entries
\(z+iw,-\bar z-i\bar w\), a block of absolute real determinant four.
Hence \(\lambda=2^{\binom r2}dx\). The specified real Haar density
at the identity, in exponential coordinates, is
\(2^{-\binom r2}\lambda=dx\).
\src{src/Formalism/LocalRingCardinality/DataCD.lean}{42}{Complexification measure on the real Lie algebra};
\src{src/Formalism/LocalToGlobal/BesselHermitianTheta/MinfDensity.lean}{13}{Real density relative to complexification}.

The map \(U(j)\to S^{2j-1}\), \(g\mapsto ge_1\), has fibre
\(U(j-1)\). At the identity its differential selects the first
column, whose independent coordinates give Euclidean tangent
measure on the sphere at \(e_1\). The remaining coordinates give
exactly the chosen density on \(U(j-1)\). By invariance the same
quotient measure is Euclidean surface measure on the whole sphere.
Starting with \(c_0=1\) gives the displayed product. For \(j\le r\)
each factor is at most \(2\cdot4^r\le8^r\).
In particular compact probability Haar has density
\(2^{-\binom r2}/c_r\) relative to \(\lambda\). For a real form
\(U(p,q)\), \(p+q=r\), the compact-dual convention prescribes this
same scalar relative to its own complexification measure. Thus
multiplication by \(c_r\) converts that convention to the specified
real density at each place.
\src{src/Formalism/LocalRingCardinality/DataCD.lean}{57}{Matching of the real density across signatures};
\src{src/ReciprocityTheory/ReductiveGroups/DiscreteSeries/FormalDegree/VolUFactProdMulLe.lean}{14}{Compact real volume factor}.

For the second assertion, every lattice \(L\) in the standard genus
is Hermitian self-dual. Indeed this holds locally for \(O^r\), and
localization commutes with the dual of a projective lattice. One may
use the conjugation-stable semilocal rings at pairs
\(\{\mathfrak p,\bar{\mathfrak p}\}\) here, including at split primes.
We bound the number of isometry classes of all totally positive
definite self-dual Hermitian \(O\)-lattices of rank \(r\).
Give the underlying real space the scalar product
\(\tfrac12\operatorname{Tr}_{E/\mathbb Q}h(x,y)\). Its squared norm
is \(\sum_\tau h_\tau(x_\tau,x_\tau)\), with one summand per conjugate
pair. For the absolute different \(\mathfrak D\) the Euclidean
dual lattice is \(L^*=2\mathfrak D^{-1}L\). Indeed testing the
pairing on all \(O\)-multiples of each lattice vector gives precisely
the codifferent condition, and \(L\) is Hermitian self-dual.
For \(0\ne y\in\mathfrak D^{-1}L\) we have
\(h(y,y)\in\mathfrak D^{-1}\bar{\mathfrak D}^{-1}\), whence
\(\lvert N_{E/\mathbb Q}h(y,y)\rvert\ge\Delta^{-2}\).
The arithmetic-geometric mean inequality gives
\(\|y\|^2\ge f\Delta^{-2/n}=fR_E^{-2}\).
Thus \(\lambda_1(L^*)\ge 2\sqrt f/R_E\).
By Banaszczyk's Euclidean transference theorem in dimension \(nr\),
\(\lambda_{nr}(L)\lambda_1(L^*)\le nr\)
\cite[Theorem~2.1]{banaszczyk}. Among real-linearly independent
vectors realizing this last minimum we may choose an \(E\)-basis
\(x_1,\ldots,x_r\) with \(\|x_i\|\le T=nrR_E\).

Write \(G=(h(x_i,x_j))\) and \(L_0=\sum O x_i\). The entries of
\(G\) are integral and have modulus at most \(T^2\) at every
embedding. There are at most \((1+2T^2)^n\) choices per entry.
Indeed the points of \(O\) in its Minkowski space have mutual sup distance
at least one by the norm inequality, so open polydiscs of radius
\(1/2\) around these points are disjoint. For a fixed \(G\) the marked
rational Hermitian space and \(L_0\) are fixed up to isometry.
The inclusions \(L_0\subseteq L\subseteq L_0^\#\), where \(\#\)
denotes the Hermitian dual, locate \(L/L_0\) in a finite module
\(A=L_0^\#/L_0\) of order \(\lvert N_{E/\mathbb Q}\det G\rvert
\le T^{2nr}\). The last inequality follows by Gram's determinant
inequality at every embedding.

Each \(O\)-submodule of \(A\) has an ordered generating tuple of
length \(r\). Locally at each prime in the support its inverse image
in \((L_0^\#)_{\mathfrak p}\) is free of rank \(r\) over a discrete
valuation ring, so its image has \(r\) generators. Combine these
tuples coordinate by coordinate under the primary decomposition of
the finite module. Consequently there are at most \(|A|^r\)
intermediate lattices for the fixed Gram matrix. All isometry
classes together number at most
\[
 (1+2T^2)^{nr^2}T^{2nr^2}\le (3T^2)^{2nr^2}.
\]
The definite standard mass is the sum of
\(1/|\operatorname{Aut}_O(L,h)|\) over the classes in its genus.
Each group is finite, since it is a discrete lattice stabilizer
inside a compact real group. The class count therefore proves
the required bound, including for nonfree \(L\).
\src{src/Formalism/SharpRestriction/GlobalSchurBounds/GeneralLinearActions/GapAtLeast.lean}{20}{Definition of the definite standard mass};
\src{src/Formalism/HermitianLattices/StatementMass.lean}{11}{Reciprocal stabilizer weights}.
\end{proof}

These estimates give \(V\le V_*\) once the displayed adelic tube
inequality holds. Applying that inequality to a different signature
still requires the comparison of the adelic covolumes with their
finite lattice stabilizers and the stated real densities.


\begin{corollary}[Size of the component bound]\label{bergvol:size}
Suppose \(L\ge1,\ \log(fr)\le L,\ \log M\le30L,\ \log R_E\le3L,\)
\(d^f\le\Delta\), and \(\log w\le25L\). Then
\[
  1\le\beta_*,\qquad \log\beta_*\le 400 f^2 r^2 L
          \le4000(frL)^3 .
\]
\end{corollary}
\begin{proof}
We have \(\log d\le6L\), \(\log\Delta\le6fL\),
\(\rho^{-1}\le e^{35L}\), and \(\varepsilon^{-1}\le
3fwM\sqrt d\le e^{61L}\). The radius estimate uses
\(\log(w+f)\le\log2+25L\); the bound on \(\varepsilon^{-1}\)
uses \(fw+2\le3fw\) and \(2\log2\le\pi\).
Since \(1\le b\le fr^2/2\) and \(\log b\le2L\), the logarithm
of \(8b!/(\pi\rho^2)^b\) is at most \(3L+36fr^2L\).
The logarithm of the base in \(A\) is at most
\(\log16+2\log(fr)+\log\Delta+\log(\varepsilon^{-1})
\le(66+6f)L\).
In \(V_*\), the first factor contributes at most \(66fr^2L\);
the logarithm of the base of its second factor is bounded by
\(\log3+2\log(2f)+4\log r+2\log R_E\le16L\).
Adding gives \(\log\beta_*\le400 f^2r^2L\). Each factor in
\(\beta_*\) is at least one, in particular \(\pi\rho^2<1\).
The final inequality uses \(f,r,L\ge1\).
\end{proof}
For the cover weight take \(w=2[E:\mathbb Q]\sum_{\tau,i}
|\lambda_{\tau i}|+1\). If the weight sum is zero then \(w=1\).
Otherwise it is an integer at least one; if its logarithm is at
most \(20L\) and \(\log[E:\mathbb Q]\le3L\), then
\(\log w\le\log3+23L\le25L\). The stage's field size
\(\log[E:\mathbb Q]+\log R_E\le3L\) implies both field estimates.
\src{src/Development/AutomorphicArithmetic/ThetaCorrespondence/FockModels/TamagawaPeterssonLambdaSqNonneg.lean}{458}{Cubic cost with the host discriminant bound};
\src{src/Development/AutomorphicArithmetic/ThetaCorrespondence/ThetaPairLifts/HodgeBundleWeights/IntegralCuspSlopePowerBound.lean}{137}{Use of the same bound for the component kernel}.

\paragraph{Short harmonic polynomials on every column lattice.}
Let \(O=O_E\), \(e=[E:\mathbb Q]=2f\), \(R_E=|\Delta_E|^{1/e}\),
and let \(F\) denote the coefficient field of the chosen deep
cover (in particular it contains the labelled images of \(E\)).
Write \(\Phi\) for the chosen CM type. In the following calculation
\(\Lambda\subset E^n\) is the \emph{column} lattice, self-dual for
\(v^*u\). It may vary through the entire allowed genus. The period
lattices of the abelian varieties have rank \(r=p+q\) over \(O\).

\begin{proposition}[A common short family for all embeddings]
\label{bfam:harm}
Assume \(n\ge p+q\). At each \(\tau\in\Phi\) prescribe two
partitions of lengths at most \(p,q\), respectively. Let \(d_P\)
be the sum of their sizes. Consider row highest weight polynomials
in the two matrices of \(p,q\) rows and \(n\) columns at each
\(\tau\), annihilated by all mixed contractions
\(\sum_b\partial_{X_{ib}}\partial_{Y_{jb}}\) at each type.
This space over \(F\) is spanned by a finite set of polynomials
in the integral algebra generated by the linear forms in the rows
with coefficients \(\tau(u)\) in the first block and
\(\tau(\bar v)\) in the second, for \(u,v\in\Lambda\).
Put \(D=(2p+1)d_P,\ h=en,\ T=\max(1,hD\sqrt e\,nR_E)\). Every
polynomial \(P\) in the set satisfies, at every embedding of \(F\),
\[
       \|P^\sigma\|_1\le (r\sqrt n\,T^{2p+1})^{d_P},
 \qquad \#\{P\}\le (D+1)^{hr}.
\]
The norm is the sum of absolute coefficients in the original
column coordinates. For the stage parameters \(L\ge1\),
\(\log(fr)\le L,\ \log R_E\le3L,\ \log S\le20L\), where
\(S=d_P+frn/2\), this gives
\(\log^+\|P^\sigma\|_1\ll rL d_P\), with an absolute constant.
\end{proposition}
\begin{proof}
Use the real Euclidean structure with norm squared
\(\operatorname{Tr}_{E/\mathbb Q}(u^*u)\). Its dual lattice is
\(\mathfrak D_E^{-1}\Lambda\) by self-duality and the trace
description of the codifferent. For \(0\ne v\) in that dual,
\(v^*v\in\mathfrak D_E^{-2}\) since conjugation preserves the
different. Its absolute field norm is at least
\(|\Delta_E|^{-2}\) by inclusion of principal fractional ideals.
All its conjugates are positive. Their arithmetic-geometric mean
inequality gives \(\|v\|_{\rm tr}\ge\sqrt e/R_E\). Euclidean
transference in dimension \(h\) now gives \(h\) independent
\(b_i\in\Lambda\) of trace norm at most \(\sqrt e\,nR_E\).
They form a real basis of the product of the \(f\) complex
column spaces. \cite[Theorem~2.1]{banaszczyk}

Independently choose \(p+q\) vectors \(u_i,w_j\) on the grid
\(\sum_{\ell=1}^h a_\ell b_\ell,\ 0\le a_\ell\le D\),
with integral \(a_\ell\). Write \(U\) for the matrix of the \(u_i\)
as columns, \(G=U^*U\), and set
\[
       v_j=\det(G)w_j-U\,\operatorname{adj}(G)U^*w_j .
\]
The coefficients used to combine the lattice vectors are integral
by their Hermitian pairings. Also \(U^*v_j=0\) by the adjugate
identity. At any embedding the input vectors have Euclidean norm
at most \(T\). For independent columns the displayed operator is
\(\det G\) times an orthogonal projection. By continuity its norm
is at most \(T^{2p}\) even at singular tuples, by Hadamard's
inequality. Hence \(\|\tau v_j\|\le T^{2p+1}\) at every embedding.

At \(\tau\), contract each of the first rows with the ordered list
\(\tau u_i\), and each of the second rows with \(\tau\bar v_j\),
using the ordinary bilinear dot product. Take the determinants of
the first \(i\) rows and columns in each resulting square matrix,
to powers the successive differences of the respective partitions
(use the reversed column order in the standard opposite block).
Multiply over \(i\) and \(\tau\). The determinants give the
specified row character. Every mixed derivative contracts a pair
of mutually orthogonal coefficient vectors and is zero. A linear
form has coefficient norm at most \(\sqrt n\,T^{2p+1}\).
A determinant of order \(i\) costs at most an additional \(i!\le
r^i\), giving the asserted bound. Its coefficients lie in the
stated integral algebra, also for a nonfree projective \(\Lambda\).

To see spanning, a complex linear functional annihilating the
grid polynomials gives a polynomial in the \(hr\) real grid
coordinates of degree at most \(D\): conjugation is real-linear
and each projected vector has degree \(2p+1\). Successive
one-variable interpolation makes it zero for all real choices.
The \(b_i\) allow independent choices at the \(f\) archimedean
places. In particular at each place use the first \(p\) and last
\(q\) columns of any unitary matrix, in the orders of the
prescribed minors. Here \(G=1\) and the projection fixes the
last columns. The obtained products form the orbit of the
product of the standard opposite flag minors under \(U(n)^f\).
Kashiwara and Vergne identify this product as a highest weight
vector and its column representation as the entire harmonic
row highest weight space for the given partitions, with multiplicity
one \cite[unitary case, \S6, Proposition~6.1 and Theorem~6.3]{kashiwara-vergne}.
Taking products gives the assertion at the \(f\) places. Descent
of matrix rank gives spanning over \(F\), since the equations
and grid vectors are defined there. Finally \(n\le2S,\ d_P\le S\)
give \(\log T\ll L\) if \(d_P>0\). For \(d_P=0\) take just \(1\).
\end{proof}
The two sets of rows here are labelled by
\(\tau,\bar\tau\), with the dual rows reversed when passing to
the ordered flag weight. Thus the integral algebra is exactly
the one generated by lattice directions in the harmonic jet of
Theorem~\ref{intjet:semi}. Relabelling the letters by an embedding
permutes the column coordinates and preserves the coefficient
estimate. For each fixed auxiliary integer, lattice and level
tuple the theta sum is linear in \(P\): the jet is the evaluation
of the constant-coefficient operator \(P(\partial)\) and the
translations and finite sums are independent of \(P\). Consequently
the short family has exactly the same image span as all such
harmonic polynomials for these data. This costs no determinant
or denominator from a change of basis.
\src{src/Development/AbelianVarietyTheory/Proofs/CuspFormEvaluations/ExistsShortHarmonicFamilyUnimodular.lean}{96}{The integral algebra and its simultaneous coefficient bounds};
\src{src/Development/AbelianVarietyTheory/Proofs/DivisorProofs/PolynomialCoefficientEmbeddingNormLeOfL1Bound.lean}{399}{Use at the actual deep coefficient field}.

\paragraph{From a half maximum to the polarized period metric.}
The transfer from the lattice in Proposition~\ref{bergvol:global}
to the periods in Proposition~\ref{intjet:arch} must account for
the host fractional ideal. Write the actual period lattice as
\(\Pi=\mathfrak a L'\), where \(L'\subset E^r\) is self-dual
for the row Hermitian form \(h_0\). Write \(\xi\) for its
imaginary polarization coefficient and \(x=|\tau\xi|\),
independent of \(\tau\). The row polarization is
\(\operatorname{Tr}_{E/\mathbb Q}(-\xi h_0)\), with
\((\xi)\mathfrak a\bar{\mathfrak a}\mathfrak D_E=O\).

\begin{lemma}[Comparison of the two minima]\label{bfam:minimum}
Let \({\rm maj}_\tau\) denote the positive majorant of \(h_0\)
at an orthonormal flag. Assume that
\[
        \sum_{\tau\in\Phi}{\rm maj}_\tau(u)\ge\varepsilon
                 \qquad(0\ne u\in L').
\]
Then the shortest nonzero polarized period
has squared length at least \(2\varepsilon/R_E^2\).
\end{lemma}
\begin{proof}
Taking ideal norms gives \(N(\mathfrak a)^{-1/f}=xR_E\).
In the product of \(f\) complex planes, \(\mathfrak a^{-1}\)
has covolume \(2^{-f}\sqrt{|\Delta_E|}/N(\mathfrak a)\).
The product of discs of squared radius \(xR_E^2\) has volume
strictly greater than \(2^{2f}\) times this covolume, as
\(\pi>2\). Minkowski gives \(0\ne b\in\mathfrak a^{-1}\)
with \(|\tau b|^2\le xR_E^2\) simultaneously. For \(v\in\Pi\)
nonzero, apply the hypothesis to \(bv\in L'\).
The polarized squared norm of \(v\) is
\(2x\sum_\tau{\rm maj}_\tau(v)\): in the orthogonal positive
and negative coordinates one conjugates the positive block in the
period convention of Lemma~\ref{dg:gaussian}, and the two terms
in the trace give \(2x\) on both
blocks. This also follows from the imaginary part of the Hermitian
form, which is the displayed trace polarization. The estimate
follows on summing
\({\rm maj}_\tau(bv)=|\tau b|^2{\rm maj}_\tau(v)\).
\end{proof}
\src{src/Development/AbelianVarietyTheory/Proofs/DivisorProofs/FirstMinAtHalfMax.lean}{82}{The inverse ideal and the distinction between the two lattices}.

\begin{corollary}[Simultaneous estimate in the normalized cover metric]
\label{bfam:metric}
Use the framed uniformization of the deep cover at each good embedding
and the self-dual row lattice model above. Assume the analytic cusp
reading and the half-maximum condition of the first hypothesis of
Proposition~\ref{bergvol:global} for the theta sections in question.
Use the parameters of Corollary~\ref{bergvol:size} and
Proposition~\ref{intjet:arch}, with \(S=d_P+frn/2\) and
\(\log S\le20L\). The sections read from the polynomials of
Proposition~\ref{bfam:harm}, multiplied by the reflex-ideal
scalar \(s\) of Theorem~\ref{intjet:trace}, have logarithmic
pointwise supremum at most \(C(frL)^2S\) in the normalized cover
metric, for an absolute \(C\). The same bound holds for the full
cover Petersson norm and every positive-mass component Petersson
norm. Multiplication by the common integer of
Theorem~\ref{intjet:semi} adds only its logarithm.
\end{corollary}
\begin{proof}
At an orthonormal flag all absolute leading Gram determinants are
one. The framed uniformization takes the Hodge directions to the
coordinate vectors and takes the theta generator to one of squared
mean one in the rigidified Gaussian metric on the period torus.
Explicitly its normalization is
\[
  \frac{1}{\operatorname{vol}\mathcal F}\int_{\mathcal F}
  |\theta(z)|^2 e^{-\pi H(z,z)}\,dz=1 .
\]
This is Haar probability on the abelian fibre and therefore
\(\|s_0\|_2=1\) in Proposition~\ref{intjet:arch}. It is part of
the framed uniformization, and changing the scalar of the
generator alone would change this equation. The polarized form
in these coordinates is \(H(z,z)=2x\sum_j|z_j|^2\), so that
\(U=\sqrt{2x}\). The cover metric multiplies the absolute raw
coefficient by \((2\pi x)^{-d_P/2}\) at this flag. In particular
the factor \((\sqrt2 U)^{d_P}\) there becomes
\((2/\pi)^{d_P/2}\), independently of the size of \(x\).
By Lemma~\ref{bfam:minimum} the minimum in the radius in
Proposition~\ref{intjet:arch} satisfies
\(\log^+\lambda^{-1}\le\log R_E+\tfrac12\log(1/\varepsilon)\ll L\).
That proposition, the short coefficient bound above, and
\(\log|\Delta_F|/[F:\mathbb Q]\ll(frL)^2\) give the asserted
estimate at the half maximum. Doubling bounds every frame.

The full cover norm integrates the square of this pointwise
quantity against the probability on the adelic unitary quotient
\(U(W_j)(\mathbb Q)\backslash U(W_j)(\mathbb A)\), writing
the group by restriction of scalars. The component norm uses the
same integrand through the frame map and the component measure
(conditional probability on positive mass). Any such measure of
mass at most one gives the same upper bound by integration.
Embeddings of the other CM orientation are read at their conjugate
good embeddings; the coefficient estimate was simultaneous for both.
This direction of the estimate uses no comparison of two total
adelic volumes.
\end{proof}
\src{src/Formalism/LocalToGlobal/ProjectiveLattices/PetersonNormSubspace.lean}{111}{The squared mean in the framed uniformization};
\src{src/Development/AbelianVarietyTheory/Proofs/DivisorProofs/EvenDomainWeightBound.lean}{68}{The full and component cover norms}.

\paragraph{Descent of the raw rule in characteristic zero.}
One can check that the finite-jet formula gives algebraic values
before identifying it with any classical modular form. Work on
the framed abelian locus over \(F\), with its actual level
structure, eigenframe, theta generator and polarized biextension.

\begin{lemma}\label{bfam:rule}
Assume the column lattice is standard locally at every prime
above the deep level \(M\). The odd-order finite theta sum of Lemma~\ref{newtheta:index}
with the prescribed half-primary division points and harmonic \(P\)
defines a regular rule on these framed points, functorial in
the characteristic-zero base. The complex reading of natural
weight rules on reduced finite-type \(F\)-tests is injective after
extension to \(\mathbb C\), whenever the framed uniformization
surjects onto complex framed points up to the weight action.
\end{lemma}
\begin{proof}
Work locally on the base in logarithm coordinates of the smooth
commutative formal group; they are defined over a \(\mathbb Q\)-algebra
by integrating its invariant differentials, with identity linear term.
On the polarized biextension choose a frame restricting to the
rigidified frames on both axes. Such a frame exists formally
locally: an arbitrary frame at the origin extends over the
completion and its axis restrictions can be divided out.
Use the biextension identification of its pullback by
\(([2],[2])\) with its fourth power. The logarithm of the
discrepancy starts in total degree three. Indeed each first-variable
addition discrepancy vanishes on \(X_1=0,\ X_2=0,\ Y=0\);
the second-variable discrepancy has the analogous property.
Compute dilation by doubling the two variables successively.
Changing the log frame by a degree-\(j\) homogeneous correction
changes the degree-\(j\) dilation discrepancy by \(2^j-4\) times
that correction. Thus all discrepancies of degree at least three
can be removed. Each addition discrepancy of the resulting frame
is now homogeneous of total degree two, by applying simultaneous
dilation to the addition diagram. Its three axis vanishings force
it to be zero.

The Rosati adjunction for an integer \(a\in O\) in this biadditive
frame can differ by an exponential of a bilinear form. In
types \((\tau,\nu)\), let its log coefficient be \(c_a\).
Functoriality for addition and composition of endomorphisms gives
\[
 c_{a+b}=c_a+c_b,\qquad
 c_{ab}=\tau(b)c_a+\nu(\bar a)c_b .
\]
Indeed for composition apply adjunction first to \(a\), with
first input \(bX\), then to \(b\), with second input \(\bar aY\).
If \(\tau=\bar\nu\), this is a derivation on the finite separable
field \(E/\mathbb Q\) with the given coefficient action, hence
zero. If \(\tau\ne\bar\nu\), interchanging \(a,b\) shows
\(c_a\) is a constant multiple of \(\tau(a)-\nu(\bar a)\).
Changing the bilinear log coefficient of the frame removes it.
The remaining ambiguity is Rosati-compatible bilinear. On the
product of \(n\) identical factors it gives the sums of mixed
quadratics killed by harmonic contraction in
Lemma~\ref{intjet:der}. Thus the jet is regular in the base and
independent of these choices. Its evaluation uses only a finite
truncation of these formal constructions and commutes with base
change. The weight action follows by changing the eigenframe in
the polynomial and raising the theta generator's scalar to the
\(n\)-th power.

After an étale base change one may enumerate \(Z_\Lambda\) and
divide each level point by \(N\) in its prescribed half-primary
part. In characteristic zero the torsion of an abelian scheme
is finite étale, and multiplication by \(N\) on the appropriate
primary torsion over geometric points is surjective onto the
prescribed level part. The symmetric normalized lifts on the
odd torsion are algebraic: in the theta group the fibre over
a point is a multiplicative torsor; the two normalizations of
Lemma~\ref{newtheta:index} select a unique lift étale locally.
Indeed first make its \(N^2\)-th power one by taking a root,
then adjust by an \(N^2\)-th root of unity to make inversion
agree with its inverse. The discrepancy is such a root and
squaring is an automorphism on these roots. Uniqueness gives
descent of the lift.

If two choices of division point differ by \(u\), then \(u\)
is \(N\)-torsion supported on the prescribed half above \(M\).
Here the column lattice is standard, so \(u\in K_0\) in the
notation of Lemma~\ref{newtheta:index}, by its period description
at those primes. The Weil pairing between two points in the
same half is trivial: its adjoint involution interchanges the
two relatively prime half ideals, hence their idempotents pair
only across the halves. The translations by \(t\) and \(u\) thus
commute and their normalized product is the lift at \(t+u\).
The vector to which the last translation is applied is fixed by
\(K_0\), by Lemma~\ref{newtheta:index}. This proves independence
of division and of enumeration. The denominator is invertible
in the base, and the level character takes values in \(F\).
The resulting regular scalars agree on the étale double overlap
and descend. On a disconnected base the assertion is checked
componentwise; the same identities hold scheme-theoretically by
the finite étale torsion and the uniqueness of its lifts.

Finally let \(\sum_i c_i f_i\) have vanishing complex reading
at a fixed embedding \(F\to\mathbb C\), with the \(f_i\)
independent over \(F\). On a framed point of any reduced
finite-type test algebra \(R\), it follows that
\(\sum_i c_i\otimes f_i(R)=0\) in \(\mathbb C\otimes_F R\).
Indeed its value is zero at each complex point by naturality,
and the uniformization up to the common weight factor.
This algebra is reduced since
\(F\) is perfect, and its complex points detect zero by the
Nullstellensatz. Expand all \(c_i\) in a basis of their finite
\(F\)-span in \(\mathbb C\). Each coefficient yields an \(F\)-linear
relation among the \(f_i\) on every test, hence is zero. This
proves injectivity.
\end{proof}
\src{src/Formalism/LocalToGlobal/ProjectiveLattices/ThShort.lean}{244}{The faithfully flat evaluation, half-primary division and weight};
\src{src/Formalism/LocalToGlobal/ProjectiveLattices/AlgebraMapCAt.lean}{9}{Reduced finite-type test rings};
\src{src/Development/AbelianVarietyTheory/Proofs/UnitaryExistence/ReadingLinearIndependentAndInjective.lean}{22}{Reading after extension of the coefficient field}.

\paragraph{The operator on row lattices.}
There is a separate numerical reason for the lower bound on the
theta rank. It applies before any denominator or choice of
representatives on the column genus.

\begin{lemma}[Invertibility at the stated theta rank]\label{bfam:gram}
Let \(\mathcal R\) be any set of full lattices of rational rank
\(b=2fr\), self-dual for the fixed rational symplectic polarization.
Set \(m(L,L')=[L:L\cap L']\). Suppose a complex matrix \(C\)
indexed by \(\mathcal R\) has constant diagonal \(A>0\) and
\(|C_{LL'}|\le A\,m(L,L')^{-n/2}\).
If \(n\ge4b+4\), then on bounded lattice-indexed coefficients
\[
     \|A^{-1}C-1\|_\infty\le\pi^2/6-1,\qquad
     \|C^{-1}\|_\infty\le A^{-1}(2-\pi^2/6)^{-1}.
\]
The result holds for coefficients in any complex Banach space
and on any closed subspace preserved by \(C\).
\end{lemma}
\begin{proof}
All self-dual lattices have the same symplectic covolume, hence
\([L':L\cap L']=m(L,L')\). In a rank \(b\) free abelian group
there are \((\ell^b-1)/(\ell-1)\le\ell^b\) subgroups of prime
index \(\ell\), by their kernels in the reduction modulo \(\ell\).
Any index-\(m\) subgroup occurs at the end of a chain with
successive prime indices in a prescribed ordering of the prime
factors of \(m\), by the primary decomposition of the quotient.
There are at most \(m^b\) such chains. The same count holds for
overlattices by ordinary lattice duality. Thus at most \(m^{2b}\)
choices of \(L'\) have this index, by first choosing \(L\cap L'\)
and then its overlattice. Index one gives only \(L\). The sum of
absolute off-diagonal entries in each normalized row is bounded
by \(\sum_{m\ge2}m^{2b-n/2}\le\sum_{m\ge2}m^{-2}\).
The norm-convergent Neumann series proves both conclusions,
including on the indicated subspaces.
\end{proof}

Here is the finite vacuum calculation giving the entry estimate.
For unshifted row lattices \(L,L'\), put
\(I=L\cap L'\). For any locally free rank-\(n\) column lattice
\(\Lambda\) self-dual for its Hermitian form, write
\(A_L=\Lambda\otimes_O L\) for the corresponding tensor lattice
with its trace symplectic form. Then
\[
 A_L\cap A_{L'}=\Lambda\otimes_O I,\qquad
 [A_L:A_L\cap A_{L'}]=m(L,L')^n .
\]
The first identity is flatness and the second follows on choosing
local bases at the finite primes. The two tensor lattices are
self-dual for the symplectic form since the column pairing
is perfect over \(O\). In the Heisenberg representation take
unit vacuum vectors \(v,v'\) with the prescribed characters on
these two lattices. Distinct translates of \(v'\) by representatives
of \(A_L/(A_L\cap A_{L'})\) are orthogonal: on \(A_{L'}\) their
eigencharacters differ by the symplectic commutator, whose
annihilator is precisely \(A_{L'}\). Their inner products with
\(v\) all have the same absolute value, because \(v\) is an
eigenvector on \(A_L\). Bessel's inequality gives
\(|\langle v,v'\rangle|\le m(L,L')^{-n/2}\).
This argument works with any translates of the two vacua, since
translation just changes their eigencharacters.

For the actual finite marks write \(\mathfrak m\) for the nonzero
proper odd ideal prime to \(\bar{\mathfrak m}\), and \(L_*\) for
the reference row lattice. The row set consists of self-dual \(L\)
equal to \(L_*\) at primes dividing \(\mathfrak m\bar{\mathfrak m}\).
The column lattice is standard there. Set
\(B_L=\Lambda\otimes_O(\mathfrak m\bar{\mathfrak m}^{-1}L)\).
The twist ideal times its conjugate is \(O\); hence these \(B_L\)
are again self-dual and their relative index remains \(m(L,L')^n\).
Choose a basis of \(L_*/\mathfrak m L_*\). The marked vector
\(\phi_L\) is the sum of translates of the unit vacuum on \(B_L\).
They are indexed by \(k\in J\subset{\rm GL}_r(O/\mathfrak m)\)
and weighted by the unitary level character. In fact the quotient
used to choose the translating tuples is the reduction modulo
\(\mathfrak m\) of \(\Lambda\otimes_O(2\bar{\mathfrak m}L)\).
Since \(2\bar{\mathfrak m}\) is locally a unit on \(\mathfrak m\),
this gives precisely the \(n\) marked rows modulo \(\mathfrak m L_*\).
The first \(r\) rows give the basis transformed by \(k\), and
the remaining rows vanish modulo \(\mathfrak m\). There is one
class for each such \(k\).

Use the common vacuum semicharacter on intersections. For
translates by \(y,y'\), nonzero inner product implies triviality
of the commutator of \(y-y'\) on \(B_L\cap B_{L'}\), hence
\(y-y'\in (B_L\cap B_{L'})^\#=B_L+B_{L'}\).
Here the duality identity follows by dualizing the sum of the
two self-dual lattices. At every \(\mathfrak p\mid\mathfrak m\)
this sum has rows in \(\mathfrak m L_{*,\mathfrak p}\).
The marked residues coincide, so only pairs with the same \(k\)
contribute. Their diagonal squared norm is \(|J|\); the mixed
inner product is at most \(|J|m(L,L')^{-n/2}\).
This proves precisely the hypotheses of Lemma~\ref{bfam:gram}
for the positive weighted class sum, with \(A=|J|\) times the
total weight. No lower bound on that mass is needed for invertibility.

The same calculation detects the arithmetic group of the marks.
For \(\gamma\in U(h_0)(E^+)\) replace \(B_{L'}\) by \(\gamma B_{L'}\);
the vacuum semicharacter is preserved because it is defined by the
polarization, the fixed \(O\)-scalar \(w_0\) with \(w_0+\bar w_0=1\),
and the unitary action. Nonzero overlap now gives
\(y-b=\gamma(y'-b')\) for some \(b\in B_L,\ b'\in B_{L'}\),
using the same duality identity. At \(\mathfrak p\mid\mathfrak m\)
both adjusted tuples have first \(r\) rows forming an integral
basis of \(L_{*,\mathfrak p}\), by reduction modulo the maximal
ideal. Thus \(\gamma\) preserves the localized reference lattice
and its residue maps the basis marked by \(k'\) to that marked by
\(k\). Its residue is \(k(k')^{-1}\in J\). This proves vanishing
on the graph orbits outside the subgroup specified by these
congruences.

One also obtains the finite vanishing along a nonzero isotropic
row line \(E\xi\). For completeness take a prime
\(\mathfrak p\mid\mathfrak m\), put \(a={\rm ord}_{\mathfrak p}
(\mathfrak m)>0\), and work at this prime and its distinct conjugate.
The trace symplectic pairing identifies \(L_{*,\bar{\mathfrak p}}\)
with the additive character dual lattice of \(L_{*,\mathfrak p}\)
in \(E_{\bar{\mathfrak p}}^r\). This follows by testing the
global self-duality locally on the two conjugate idempotents.
Scale \(\xi_{\bar{\mathfrak p}}\) to a primitive vector \(u\)
of \(L_{*,\bar{\mathfrak p}}\). The element
\(\varpi_{\bar{\mathfrak p}}^{-a}u\) belongs to the shifted row
lattice there but lies outside the additive dual of \(L_{*,\mathfrak p}\).
For any marked tuple \(y\) its first \(r\) rows form a basis at
\(\mathfrak p\). Some integral multiple (in the local integer
ring) of \(\varpi_{\bar{\mathfrak p}}^{-a}u\) therefore pairs
nontrivially under the additive character with one of those rows.
Put this vector in the corresponding column position of a tuple
\(\delta\in B_L\otimes\widehat{\mathbb Z}\) on \(E^n\otimes E\xi\),
and zero at the other components. Equivalently one can choose a
rational \(\delta\) giving the same character: the intersection of
a full lattice with a rational subspace spans a lattice in that
subspace, and its completion is the intersection of the completed
lattice with the completed subspace, by a basis adapted to the
saturated intersection. Approximation modulo a finite-index
sublattice then preserves the character.
On this subspace both the alternating form and the vacuum quadratic
phase vanish, since \(h_0(\xi,\xi)=0\). The Heisenberg translations
on it commute and their action on the marked translate has the
character of this pairing. For each of the finitely many labels
choose such a \(\delta\) with eigenvalue \(\eta\ne1\). The product
of the commuting operators \(T_\delta-\eta\) kills the sum of
marked translates and has nonzero augmentation
\(\prod(1-\eta)\). Hence its image in the coinvariants of the
isotropic row subgroup vanishes. This supplies the vanishing for
non-graph terms whose evaluation factors through those coinvariants.
\src{src/Development/AbelianVarietyTheory/Proofs/SiegelThetaTransformation/ThetaTestFunctionNormAndOverlapBound.lean}{133}{The localized congruence for a surviving pair of labels};
\src{src/Development/AbelianVarietyTheory/Proofs/SiegelThetaTransformation/HermitianThetaLevelDataFormsSystem.lean}{60}{The off-diagonal sum at rank \(8fr+4\)};
\src{src/Development/AbelianVarietyTheory/Proofs/SiegelThetaTransformation/ThetaTestFunctionNoCoinvariantsIsotropic.lean}{216}{Isotropic row coinvariants of the same marked tests}.

\paragraph{Which definite Siegel--Weil formula is needed.}
There is a way to use the ordinary adelic theorem without comparing
the volumes of different genera of self-dual Hermitian lattices.
Use just the genus of \(\Lambda_*=O^n\) in the positive definite
space \(V=E^n\), and require equality with the standard lattice
at a prescribed integer \(B_0\) divisible by the deep level and
the excluded rational primes. Every lattice used in this smaller
family is among the permitted columns. Its positive class mass
still gives the strict inequality in Lemma~\ref{bfam:gram}.

\begin{lemma}[Adelic realization of a lattice family]\label{bfam:genus}
Write \(H=U(V)\) over \(E^+\) and suppose \(n>4r\).
For each such column lattice let
\(\varphi_\Lambda\) be a function on \(V(E^+)^r\) (thus on
\(r\) columns in \(E^n\)). Assume the following conditions for
some positive integer \(R\).
\begin{enumerate}
\item The function is supported in \(R^{-1}\Lambda^r\) and
invariant under \(R\Lambda^r\).
\item If two lattices coincide at all primes over \(RB_0\),
their functions agree on the intersection of their indicated
supports \(R^{-1}\Lambda^r\cap R^{-1}\Lambda'^r\).
\item For a rational isometry \(\gamma\) integral and congruent
to the identity modulo \(B_0\) at its prime divisors,
\(\varphi_{\gamma\Lambda}(\gamma x)=\varphi_\Lambda(x)\).
\end{enumerate}
Then the Schwartz functions obtained by completion form the
translates of one finite adelic Schwartz function in this genus.
For a fixed \(H(E^+\otimes\mathbb R)\)-invariant archimedean
Schwartz function and the splitting in which \(H\) acts by
linear substitution in the doubled Schrödinger model, the
genus average of the theta functions with stabilizer weights,
divided by the sum of those weights, is the anisotropic
Siegel--Weil Eisenstein value for the standard lattice test.
The same is true of the genus average of the corresponding
Eisenstein values themselves.
\end{lemma}
\begin{proof}
Let \(\mathcal U\subset H(\mathbb A_f)\) be unrestricted away
from \(B_0\), and the principal congruence subgroup of the
standard lattice at its divisors. For \(g\in\mathcal U\)
consider \(\Lambda'=g\widehat\Lambda\cap E^n\). Write
\(\widehat\varphi\) for the function on the finite adeles
defined by the same finite quotient
\(R^{-1}\Lambda^r/R\Lambda^r\), extended by zero.
We claim
\(\widehat\varphi_{\Lambda'}(g x)=\widehat\varphi_\Lambda(x)\).
It suffices to check this on \(R^{-1}\widehat\Lambda^r\).
Choose a rational unitary \(\gamma\) sufficiently close to \(g\)
at all primes over \(RB_0\). Weak approximation here follows
from the Cayley transform between the skew-adjoint Lie algebra
and a dense open in \(H\); the latter is locally dense at every
finite place, since its complement is a proper algebraic closed
set in a smooth connected group. Linear weak approximation in
the Lie algebra proves the assertion. Thus \(\gamma\Lambda\)
and \(\Lambda'\) agree at those primes, \(\gamma\) meets the
congruences of the third hypothesis, and \(\gamma x,g x\)
have the same residues modulo the local period lattice there.
The lattice \(R^{-1}(\gamma\Lambda\cap\Lambda')^r\) maps onto
these residues by the Chinese remainder theorem. Choose a
rational \(z\) in it representing them. Outside \(R\) the
support lattice and period lattice coincide. Consequently the
second hypothesis compares \(\varphi_{\Lambda'}(z)\) and
\(\varphi_{\gamma\Lambda}(z)\), and the third identifies the latter
with \(\widehat\varphi_\Lambda(x)\). This proves the claim.

The same weak approximation at \(B_0\) gives
\(H(\mathbb A_f)=H(E^+)\mathcal U\).
Put \(K=\operatorname{Stab}_{\mathcal U}(\widehat\Lambda_*)\).
The compact right cosets on the finite double quotient
\(H(E^+)\backslash H(\mathbb A_f)/K\) correspond to the
lattice orbits with the indicated congruences. Representatives
can all be taken in \(\mathcal U\). Right invariance by \(K\)
follows from the claim. On each coset the compact archimedean
group integrates the invariant archimedean factor, and the
remaining mass is proportional to the reciprocal order of
the stabilizer. Indeed its chart is a quotient of
\(H(E^+\otimes\mathbb R)\times K\) by this finite group;
the numerator Haar volume is common to all cosets. Dividing by
the total mass gives the theta integral for Haar probability.
The doubled split group has Witt index \(r\); the source
dimension is \(n>4r\). In the notation of Gan--Qiu--Takeda the
anisotropic formula is at \(s=(n-r)/2>0\), with unitary
constant one for the probability theta integral
\cite[\S3.1 and Theorem~7.1]{ganQiuTakeda}. Finally each
finite transport has the same Siegel section: the \(H\)-action
commutes with the doubled group and linear substitution fixes
evaluation at zero. This proves the last assertion as well.
\end{proof}
For an archimedean polynomial insertion one takes the invariant
sum over a unitary orthonormal basis of the column representation
in Proposition~\ref{bfam:harm}, paired with its conjugate on
the second copy. This is the identity endomorphism of that
representation, so the resulting doubled Schwartz function is
invariant under the compact column group. Lemma~\ref{bfam:genus}
applies just as it does to the Gaussian. We verify the mixed
finite equivariance directly.
\src{src/Development/AbelianVarietyTheory/Definitions/SpanningFamilies/FirstCoordinateAdmissible.lean}{139}{Equivariance, periodicity and finite locality of the column family};
\src{src/Development/AbelianVarietyTheory/Proofs/SiegelThetaTransformation/AdmissibleDoubledHermitianThetaFamily.lean}{413}{Mixed vacuum comparison for the marked doubled test}.

\begin{lemma}[The finite phase for the projected marked vacua]\label{bfam:phase}
Use the finite adelic Schrödinger representation of the trace
symplectic space underlying \(V\otimes_E E^r\). For each row
lattice \(L_i\), \(i=1,2\), start with any unit vacuum \(v_i\) for
the standard column and the shifted row lattice
\(D_i=\mathfrak m\bar{\mathfrak m}^{-1}L_i\). For the column
\(\Lambda\) take the unit normalization \(u_i^\Lambda\) of the
orthogonal projection of \(v_i\) to the vacuum of
\(A_i^\Lambda=\widehat\Lambda\otimes_{\widehat O}\widehat D_i\).
The vacuum character at each lattice is the prescribed
semicharacter defined by \(w_0\). Form the marked sums
\(\phi_i^\Lambda\) above from these \(u_i^\Lambda\). For
\(h\in H(\mathbb A_f)\) preserving the marked residues (in
particular for \(h\in\mathcal U\) of the preceding proof),
\[
 W_{\phi_1^{h\Lambda},\phi_2^{h\Lambda}}(hx)
       =W_{\phi_1^\Lambda,\phi_2^\Lambda}(x),\qquad
 W_{u,v}(x)=\langle \rho(x)u,v\rangle .
\]
Here \(h\Lambda\) means the lattice with completion
\(h\widehat\Lambda\), and \(\rho\) denotes the Weyl operators.
In particular the finite doubled tests used in the ordinary
adelic theta identity are the translates of a single test for
each pair of row lattices, with their original complex coefficients.
\end{lemma}
\begin{proof}
Choose complementary rational Lagrangians in the underlying
symplectic space. Use dual Haar measures for the finite additive
character \(\psi_f\) and the pairing of the two Lagrangians. With
the convention \(\rho(a,b)u(t)=\psi_f(b(t+a/2))u(t+a)\) and the
inner product linear in its first argument, the map to the doubled
diagonal model is
\[
 W_{u,v}(a,b)=\int u(t+a/2)\overline{v(t-a/2)}\,\psi_f(bt)\,dt .
\]
Indeed one first changes variables from the two position coordinates
to their difference and midpoint, then Fourier transforms the
midpoint. The Jacobian of the first change is of absolute value
one, also at dyadic places. A change of basis for the dual
coordinate changes the corresponding dual Haar. Fourier inversion
and this change of variables give
\[
 \int W_{u,v}\overline{W_{s,t}}=
             \langle u,s\rangle\overline{\langle v,t\rangle}.
\]
The integrals here are integrals of Schwartz functions. In
particular \(W_{u,v}(0)=\langle u,v\rangle\).
In the paper's rational doubled coordinates the symbol is
\(a\,W_{u,v}(-x)\), where \(a\ne0\) is the common change of Haar
normalization. In fact translation of the two inputs by
\(y_1,y_2\) translates the position argument by \(-(y_1-y_2)\)
in those coordinates. Set \(y_1=-x,\ y_2=0\) and evaluate at
the origin. The remaining half-pairing has zero rational trace:
it is a purely imaginary scalar times the real Hermitian norm
of \(x\). Thus the value of the doubled symbol at \(x\) equals
its origin value on \(\rho(-x)u,v\). The origin pairing is the
displayed inner product by partial Fourier transformation. The
reflection commutes with every linear column isometry.

Write \(\epsilon(a)\rho(a)\) for the action fixing a vacuum on a
self-dual lattice \(A\). The formula for \(\epsilon\) depends on
the vector and \(w_0\), so its restrictions to intersections
agree. Explicitly it is the half quadratic character of the trace
pairing of \(a\) with \(w_0 a\), with the sign determined by
\(\psi_f\). On \(A\) it cancels the Weyl multiplication cocycle:
\(w_0+\bar w_0=1\) and the pairing of \(A\) with \(w_0A\) is
integral. The integral of \(\epsilon(a)\rho(a)\) against Haar
probability on \(A\) is the orthogonal projection \(P_A\) to its
fixed space. For a vacuum \(v\) on \(A_*\), \(W_{v,v}\) vanishes
outside \(A_*\): conjugation by \(b\in A_*\) multiplies the
coefficient by the commutator character and \(A_*^\#=A_*\).
Thus
\[
 \|P_Av\|^2=\langle P_Av,v\rangle
      =[A:A\cap A_*]^{-1}.
\]
The norm is positive and the unit projection has positive real
inner product with \(v\). This integral is also the normalized
finite average of the distinct translates used on the rational
lattice: the orbit on a Schwartz vector is finite and each fibre
of its orbit map has equal measure.

For vacua \(u_i\) on \(A_i\), the same commutator, tested on
\(A_1\cap A_2\), gives support in \(A_1+A_2\). At
\(x=a_1+a_2\), \(a_i\in A_i\), move \(\rho(a_1)\) to the
first vacuum and \(\rho(a_2)\) to the second. The result is
\(W_{u_1,u_2}(0)\) times a unit scalar determined by
\(\epsilon(a_i)\) and their Weyl cocycle. In particular its
absolute value is constant throughout the support. The preceding
Fourier identity on the diagonal gives
\[
 |W_{u_1,u_2}(0)|={\rm vol}(A_1+A_2)^{-1/2}.
\]
This volume is independent of the self-dual column: each
self-dual \(A_i\) has volume one and the relative index here
is \([D_2:D_1\cap D_2]^n\).

To compare the \emph{arguments} for columns \(\Lambda,\Lambda_*\),
observe that the intersection of the supports is
\[
 (\widehat\Lambda\cap\widehat\Lambda_*)\otimes_{\widehat O}
                        (\widehat D_1+\widehat D_2).
\]
This can be checked in a local basis of the row sum. Every
element of this intersection therefore has a decomposition
\(a_1+a_2\) with the two \(a_i\) in the corresponding vacuum
lattices for both columns. Its scalar relative to the value at
zero is identical for the two symbols. Their Fourier inner
product is consequently the positive volume of this intersection
times the product of the two values at zero, with the second
conjugated. On the other hand the Fourier identity identifies
this inner product with
\(\langle u_1^\Lambda,v_1\rangle
 \overline{\langle u_2^\Lambda,v_2\rangle}>0\).
The two values at zero have equal modulus, so they are equal.
This applies separately to every column and does not impose a
choice of phase on either initial \(v_i\).

An isometry \(h\) of the column carries a decomposition
\(a_1+a_2\) to \(ha_1+ha_2\). It preserves the half quadratic
character and the Weyl cocycle. Equality at zero now gives the
assertion for the unmarked vacua. For a marked summand translate
by \(y_i\in\widehat\Lambda\otimes 2\bar{\mathfrak m}\widehat L_i\).
Its doubled symbol is the unmarked symbol at a translate of \(x\)
multiplied by the Weyl phase of \(x,y_1,y_2\); the latter is
also invariant under \(h\). The lifts \(hy_i\) give the prescribed
residues for the new column. Any two lifts of the same residue
differ by an element of \(2A_i^{h\Lambda}\), as may be checked at
the primes of \(\mathfrak m,\bar{\mathfrak m}\) and their
complement. Its vacuum character is trivial, and its half
pairing with the chosen lift is integral since both lifts lie in
\(2(h\widehat\Lambda\otimes\widehat L_i)\).
Thus each marked translate is independent of this choice.
Rational lifts with these residues exist by the Chinese remainder
theorem. Summing with the given weights on \(J\times J\) proves
the formula.
\end{proof}
The proof uses the finite additive character trivial on
\(\widehat{\mathbb Z}\) obtained from the rational trace pairing.
It gives the finite diagram also at \(2\), with the specified
semicharacter. In particular the normalized class average of
these genuine doubled symbols is exactly the finite input of
Lemma~\ref{bfam:genus}. The choices of phases for the initial
row vacua are allowed to depend on \(L_i\).
\src{src/Development/AbelianVarietyTheory/Definitions/SpanningFamilies/FirstCoordinateAdmissible.lean}{81}{The sign in the doubled position coordinate};
\src{src/Development/AbelianVarietyTheory/Proofs/SiegelThetaTransformation/DoubledHermitianTestFunctionZeroAndNorm.lean}{275}{The common origin normalization};
\src{src/Development/AbelianVarietyTheory/Definitions/SpanningFamilies/IsSystem.lean}{148}{Vacuum semicharacter and the normalized projection};
\src{src/Development/AbelianVarietyTheory/Proofs/SiegelThetaTransformation/DoubledVacuumEqHermitianCharacter.lean}{554}{Constancy of the mixed scalar across columns}.

In the doubled model a rational maximal isotropic subspace of
\(W\oplus(-W)\) either is the graph of an isometry, or contains an
isotropic row line in \(W\oplus0\). Indeed if its first projection
has a kernel, its second projection is orthogonal to that kernel,
and hence has a nonzero kernel as well. If both projections are
invertible, isotropy is precisely the isometry equation. On a
graph the finite evaluation at zero is the inner product with the
second vacuum transported by the row isometry, since the base
graph has functional \(W_{u,v}(0)\). For a nongraph the evaluation
is invariant under Weyl translations in the indicated row line
tensored with \(V\). In the polarization of this maximal isotropic
subspace they are multiplication by additive characters of the
position variable and each such character is one at zero. Thus
the augmentation argument above kills this orbit term, including
after archimedean polynomial insertion. This accounts for all
rational isotropic subspaces.

\begin{corollary}[The Fourier mark on the identity cover]\label{bfam:mark}
At the deep level used for the raw spanning calculation the
auxiliary map of that calculation is the identity. Its auxiliary
ideal is \(O\), its integer is one and its auxiliary character is
trivial. Thus \(\mathfrak m\) is the same half of the level on
both sides. The subgroup \(J\) of residues is the identity:
its entries differ from the identity by reductions of elements
of \(\mathfrak m\) modulo \(\mathfrak m\). The algebraic
Fourier mark on the half-primary torsion selects exactly the
single translated shifted vacuum with the first \(r\) basis
rows prescribed and the remaining rows zero.
\end{corollary}
\begin{proof}
Put \(A=\widehat\Lambda\otimes\widehat L\) and
\(D=\widehat\Lambda\otimes\mathfrak m\bar{\mathfrak m}^{-1}
\widehat L\), and use a subscript \(*\) for the standard column.
The alternating commutator identifies \(A/(A\cap D)\) with
the character dual of \(D/(A\cap D)\): each induces a character
on the other group, and self-duality of both lattices gives
trivial kernel on both sides. The second group is supported on
\(\bar{\mathfrak m}\), where it is
\((\bar{\mathfrak m}^{-1} L/L)^n\); at those primes the column is
standard. It is the group \(Q\) of the \(n\) half-level torsion
points. To see the identification, use the level isomorphism
on \(M\)-torsion and the two Chinese remainder idempotents of
\(O/(M)\), since \((M)=\mathfrak m\bar{\mathfrak m}\).
The dual group similarly has residues in \((L/\mathfrak m L)^n\).
In these identifications the Fourier coefficient of the chosen
algebraic mark is precisely the pairing character for those
residues; this is the defining property of the dual marking.
One can lift its tuple by \(y\in\widehat\Lambda\otimes
2\bar{\mathfrak m}\widehat L\), since \(2\bar{\mathfrak m}\)
is a unit at \(\mathfrak m\). Let \(v\) be a unit unshifted
vacuum for the standard column and put \(z=P_Av\). The restriction
of \(\epsilon(d)\rho(d)\) to the intersection fixes \(z\).
Consequently averaging these operators on \(z\) amounts to averaging
over \(Q\). Conjugation by \(\rho(y)\) multiplies the summands
by their pairing character. Moreover \(\rho(y)z=z\), because
\(y\in2A\). With the character taken from this conjugation, the
normalized Fourier average is therefore \(\rho(y)P_D P_Av\).
Taking the opposite convention for the additive character conjugates
both the definition of the dual mark and the commutator.
The odd-order canonical Heisenberg lifts agree on \(Q\) with this
action on \(z\). Indeed the lifts on its isotropic subgroup are
additive, inversion acts by negation, and a change of such a lift
is a character of order dividing two on the odd group.

The normalization here can be compared on \emph{all} rows at
once. Pick the initial unit shifted vacuum to be the normalized
\(P_{D_*}v\). At the primes over \(M\), \(A=A_*,D=D_*\);
elsewhere \(D=A,D_*=A_*\). The local projectors at distinct
places commute, so on \(v\) we have
\(P_DP_Av=P_DP_{D_*}v\).
The product statement follows by factoring a vacuum for \(A_*\)
into local vacua. For completeness the vacuum line of a self-dual
lattice is one dimensional: the calculation of \(W_{u,v}\) above
gives \(|\langle u,v\rangle|=1\) for any two unit vectors fixed
with its character. The restricted tensor of local vacua supplies
one such vector, hence any choice differs only by a unit scalar.
Projector norms in Lemma~\ref{bfam:phase} now show that the
coefficient of the unit shifted marked vacuum is the positive
number
\[
 [D_*:D_*\cap A_*]^{-1/2}\,
 [D:D\cap D_*]^{-1/2}.
\]
The first index is \(|O/\mathfrak m|^{nr}\). The second is
\([\widehat\Lambda:\widehat\Lambda\cap\widehat\Lambda_*]^r\),
by using local bases of the shifted row lattice. Both are
independent of the row and its component. Changing the original
shifted unit vacuum for a fixed row multiplies all its projected
unit vacua by the same phase. Finally, the residue conditions on
the last \(n-r\) rows on both sides of the comparison are zero
modulo the identical ideal \(\mathfrak m\).
\end{proof}
\src{src/Formalism/LocalToGlobal/ProjectiveLattices/ThShort.lean}{300}{Identity auxiliary map on the deep model};
\src{src/Development/AbelianVarietyTheory/Definitions/SpanningFamilies/IsBaseClassLevelAdapted.lean}{78}{The two half-level ideals and the dual residue marking};
\src{src/Development/AbelianVarietyTheory/Proofs/SiegelThetaTransformation/LeSupOfCoprime.lean}{100}{Composition of the unshifted and shifted projections}.

\paragraph{Integral change by a flat character.}
Fix an algebraic component \(Z\) of the cover over \(F\).
For a rational trivialization \(t\) of a flat automorphic line on
\(Z\), suppose its pointwise metric is constant \(r_\sigma>0\)
on each complex fibre, with conjugate embeddings giving the same
constant. Suppose \(\rho_w\) is a lower bound for its vertical
orders at every test over the finite prime \(w\) of \(F\), in the
normalization which gives the order \({\rm ord}_w(a)\) to
\(a\in F^\times\). Assume \(\rho\) has finite support and is
integral off a finite set \(B\) of rational primes. Put \(h=[F:\mathbb Q]\)
and
\[
 \delta(t)=h^{-1}\left(\sum_w\rho_w\log Nw-
                         \sum_{\sigma:F\hookrightarrow\mathbb C}
                               \log r_\sigma\right).
\]
There is \(a\in F^\times\) such that \(at\) is integral at all
the tests and
\begin{equation}\label{bfam:flat}
 \max_\sigma\log\|at\|_\sigma
 \le-\delta(t)+\frac{\log|\Delta_F|}{2h}+\sum_{\ell\in B}\log\ell .
\end{equation}
Indeed take \(\mathfrak b=\prod_w w^{\lceil-\rho_w\rceil}\).
Every element of this ideal gives integrality. The rounding cost in
\(\log N\mathfrak b\) over \(\ell\in B\) is at most
\(\sum_{w\mid\ell}\log Nw\le h\log\ell\).
In the Minkowski embedding use the intervals and complex discs
with radii \(R/r_\sigma\) at real and complex places. Their volume
is \(2^{r_1}\pi^{r_2} R^h/\prod_\sigma r_\sigma\).
The covolume of \(\mathfrak b\) is
\(2^{-r_2}\sqrt{|\Delta_F|}\,N\mathfrak b\).
Taking \(R^h=\sqrt{|\Delta_F|}\,N\mathfrak b\prod_\sigma r_\sigma\)
achieves Minkowski's bound, including the equality case by compactness.
It gives \eqref{bfam:flat}.

If \(t\) is a nowhere vanishing section on the characteristic-zero
component, multiplication is a field-linear isomorphism to the
twisted cusp space there whenever the flat trivialization and
its inverse preserve cusp boundary conditions. Extension by zero
from the component and restriction back are inverse on its supported
sections. Thus any field-spanning family of global sections restricts
to a spanning family on \(Z\) under this component decomposition.
Multiplication by \(at\) preserves the vertical inequalities and
multiplies the component Petersson norm by \(|\sigma a|r_\sigma\):
this follows under the integral defining that norm, since the
absolute metric of the multiplier is constant on the component.
The inverse of \(at\) is needed only over \(F\). In particular no
division by a component mass or by a generator of a fractional
ideal is needed in the integral lattice.
\src{src/Development/AbelianVarietyTheory/Proofs/DivisorProofs/ExistsShortIntegralFlatLineTwistingSection.lean}{203}{Minkowski rescaling for fractional vertical orders};
\src{src/Development/AbelianVarietyTheory/Proofs/DivisorProofs/ExistsSpanningFamilyMulTwistingSection.lean}{97}{Component projection and multiplication on the field-spanning family}.

\paragraph{Passage from the marked family to the cover.}
Fix a good embedding of \(F\) and a component
\(k\). Let \(L_k\) be its row lattice after transport of its level
marking into the common rational Hermitian space. Write
\(\Theta_{\Lambda,P}(L_k,x)\) for the corresponding analytic
theta jet on a frame \(x\), using the marked vector of
Corollary~\ref{bfam:mark}, and write \(f_{\Lambda,N,P}\) for the
algebraic rule. Identify the polynomial variables by the
eigenframe and the framed periods. The required reading has the form
\begin{equation}\label{bfam:reading}
  f_{\Lambda,N,P}(k,x)=C_k\,u_{\Lambda,N}\,
                   \Theta_{\Lambda,P}(L_k,x).
\end{equation}
Here \(C_k\ne0\) is independent of \(\Lambda,N,P,x\); the
nonzero \(u_{\Lambda,N}\) is independent of \(k,P,x\).
A different normalization for the initial vacuum for \(L_k\)
can be absorbed in \(C_k\).

\begin{proposition}[The theta reading on a polarized period chart]
\label{bfam:unitread}
Use the principal PEL period chart of Corollary~\ref{bfam:metric}
on one connected component. Use its rigidified identification
\(\mathcal L=(1,w_0)^*\mathcal B\) with the pullback of its
polarized biextension, its given half-level marking and its dual,
and its labelled eigenframe as the tangent coordinates. Use the
identity auxiliary map of Corollary~\ref{bfam:mark}. For the raw
section of Lemma~\ref{newtheta:index}, \eqref{bfam:reading} holds
on unit flags up to a common factor of constant positive absolute
value in place of \(C_k\). With the shifted standard vacua chosen
as in Corollary~\ref{bfam:mark}, one may take
\[
        u_{\Lambda,N}=|Q|^{-1/2}
           [\Lambda:\Lambda\cap O^n]^{r/2}.
\]
For the convention that divides the raw rule by \(\beta_H\),
multiply this number by \(\sigma(\beta_H)^{-1}\).
On the good framed chart, \eqref{bfam:reading} holds on the
connected component with a single \(C_k\ne0\).
\end{proposition}
\begin{proof}
Write \(\Pi\) for the row period lattice and \(g=fr\).
In the Appell--Humbert frame over the Lie algebra the automorphy
factor of \(\mathcal L\) at a period \(v\) has unitary
semicharacter \(\chi(v)\) and Hermitian form \(H\).
The pullback by \((1,w_0)\) of the addition biextension has
Hermitian form \(H(z,w_0z')+H(w_0z,z')=H(z,z')\), by Rosati
adjunction. Its semicharacter is the quotient of the characters
at \((1+w_0)v,v,w_0v\), namely
\(\exp(\pi i E_L(v,w_0v))\) by the semicharacter law. Uniqueness
of the Appell--Humbert data of a rigidified line gives this same
semicharacter for \(\mathcal L\). Its values on periods are signs.

Choose a symplectic \(\mathbb Z\)-basis of \(\Pi\) and denote the
two period matrices in the unit tangent frame by \(P_0,Q_0\).
Set \(h=P_0^{-1}\), \(\Omega=h Q_0\), \(Y=\Im\Omega\). The
Riemann bilinear relations give \(Y>0\), \(\Omega^t=\Omega\)
and \(H(z,z)=(hz)^tY^{-1}\overline{hz}\). The signs in this
basis specify a half characteristic \((\alpha,\eta)\). Put
\[
 T(z)=e^{\pi(hz)^tY^{-1}(hz)/2}
          \sum_{\ell\in\mathbb Z^g+\alpha}
           e^{\pi i\ell^t\Omega\ell+2\pi i\ell^t(hz+\eta)}.
\]
Its multipliers follow by shifting the summation index on imaginary
periods and multiplying each term on real periods. The signs
produce the specified characteristic. The same period laws
force every holomorphic section to be a scalar times \(T\):
remove the first exponential, take the Fourier series in the real
periods, then use the other periods to recover every coefficient
from one. This also proves \(T\ne0\).

For a unit finite vacuum for \(A_*\), the \(n\)-fold Schrödinger
theta in these coordinates is
\[
       d_L(\det h)^{n/2}\prod_{j=1}^n T(z_j)
\]
with a fixed \(d_L\ne0\). In the basis adapted to \(\Pi\) the
integral translations and their dual characters force the finite
test to be supported on the translate
\(\alpha+\widehat{\mathbb Z}^{g}\) in each row; its values there
are fixed multiples of the character prescribed by \(\eta\).
Restriction of this function to the rational points gives exactly
the displayed characteristic series. The determinant is the
weight-\(n/2\) factor for the change \(z\mapsto hz\).
To compare to a fixed initial rational symplectic basis, it suffices
to change basis by a linear substitution, a symmetric shear, or
Fourier exchange. These generate the rational symplectic group by
symplectic elimination. Linear substitution reindexes the rational
sum. For a shear the archimedean and finite quadratic characters
at every rational point are inverse by the product formula. For
exchange, apply Poisson summation with dual Haar measures to the
finite test times the archimedean Gaussian with linear argument.
Completing the square gives in each row the factor
\(\det(-i\Omega)^{-1/2}
e^{-\pi i\zeta^t\Omega^{-1}\zeta}\) for the Fourier transform
with shifted argument \(\zeta\). The quadratic exponential is
absorbed by the change of the holomorphic Gaussian prefactor in
\(T\), and the half determinant by the change of \(\det h\).
Thus the new finite vacuum is multiplied by a fixed scalar
independent of the period point; choices of square root give
only a constant. This is the period-coordinate version of the
Schrödinger theta transformation.

The squared mean of \(T\) on the polarized torus is
\(\det(2Y)^{-1/2}\). In fact write \(hz=t+\Omega s\) in
the unit period cube. The integral over \(t\) kills all unequal
Fourier indices. The remaining integrand, including the metric
\(e^{-\pi H(z,z)}\), gives the translates of
\(e^{-2\pi s^t Y s}\); integrating the sum over the \(s\)-cube
gives its integral over \(\mathbb R^g\).
The squared mean of the \(n\)-fold Schrödinger product is
therefore
\[
    |d_L|^2|\det h|^n\det(2Y)^{-n/2}
      = |d_L|^2\bigl(\det H/2^g\bigr)^{n/2}.
\]
At unit flags \(\det H=(2x_0')^g\), where \(x_0'>0\) denotes
the common absolute polarization scalar called \(x\) in
Corollary~\ref{bfam:metric}; equivalently one may choose
orthonormal flags first and extend by the flag character.
The determinant is constant. The geometric generator has squared
mean one. Hence its \(n\)-fold product is the Schrödinger product
times a scalar of constant positive modulus, independent of all
column data and polynomial insertions.

It remains to check what the geometric translations and
differentiations do to this one product. In these coordinates
put
\[
  (U_t^{(m)}s)(z)=e^{-\pi mH(z+t/2,t)}s(z+t).
\]
It sends \([N]^*s=s(Nz)\) under \(U_t^{(N^2)}\) to
\((U_{Nt}^{(1)}s)(Nz)\). This pullback equality in the
Appell--Humbert frames follows since the factors of automorphy
agree and the isomorphism is rigidified at zero. When
\(N^2t\in2\Pi\), inversion sends this translate to its inverse
and its \(N^2\)-th power acts trivially: the residual period
sign at \(N^2t\) is one. Thus it is exactly the normalized
geometric lift of the rule. Such representatives can always
be chosen since \(N\) is odd. On an individual summand in
period coordinates \(u=a+\Omega b\) the finite Heisenberg
operator on a rational coefficient function is
\[
   v(s)\longmapsto
             e^{2\pi i a^t(s-b/2)}v(s-b).
\]
Reindex \(s\) by \(s+b\) in the characteristic sum. The linear
argument changes from \(\zeta\) to \(\zeta+u\), and the factor
is \(e^{\pi i(b^t\Omega b+2b^t\zeta+a^tb)}\).
Restoring the holomorphic Gaussian this is just \(U_{h^{-1}u}^{(1)}\).
Thus this comparison holds with phases before taking any norm.

For the first sum of the rule write the division points as
\(t=s/N\) with \(s\in\Lambda\otimes\Pi\) modulo \(N\Pi^n\),
as in Lemma~\ref{newtheta:index}. Representatives can be chosen
even in \(\Lambda\otimes\Pi\): its quotient by \(N\Pi^n\) has
odd order since the column is standard at two and \(N\) is odd.
Then \(Ns\in2\Pi^n\) and their semicharacters are trivial.
The finite sum on the initial vacuum is exactly
\(|Z_\Lambda|P_Av\), since \(N\Pi^n\) fixes it with the
prescribed semicharacter. The actions of the second, half-level
sum were computed in Corollary~\ref{bfam:mark}. Equivalently
its translating vectors modulo \(\Lambda\otimes\Pi\) exhaust
\(D/(D\cap A)\): at the primes of \(M\) this is the given
half-level isomorphism, and at its complement the primary torsion
class vanishes. Division representatives may be taken to be zero
modulo \(2N\Pi^n\) at the complement of \(\bar{\mathfrak m}\)
and even in the shifted lattice at \(\bar{\mathfrak m}\), by the
Chinese remainder theorem. Their symmetric lifts consequently
act by the normalized shifts used in that corollary.

In the universal-cover Appell--Humbert frames the constant
biextension frame on two Lie variables is biadditive and
Rosati-compatible. Indeed the factor of automorphy of the
addition biextension is obtained by dividing the factor at
\(z+z'\) by those at \(z,z'\). The resulting cross terms are
additive in the opposite variable, and
\(H(az,z')=H(z,\bar a z')\). Equivalently these statements follow
by taking the rigidified cube trivializations; in this frame the
third differences of the linear and quadratic exponents vanish.
The pullback frame along \((1,w_0)\) is the
Appell--Humbert frame of \(\mathcal L\) just used, by equality
of their Hermitian forms, semicharacters and rigidifications.
Consequently the algebraic jet may be evaluated by ordinary
differentiation on this analytic frame. Any change of compatible
frame is the common bilinear ambiguity killed by harmonicity
in Lemma~\ref{bfam:rule}. Each labelled Hodge derivative is
the corresponding coordinate derivative by the definition of the
framed period chart. Substitution \(z\mapsto Nz\) contributes
\(N^{d_P}\), cancelled by the same denominator in the rule.
All sums are finite except the normally convergent holomorphic
theta series, so these differentiations commute with summation.

The raw column factor is thus \(|Z_\Lambda|/|K_0|\) times the
two square-root indices in Corollary~\ref{bfam:mark}. Setting
\(I=[\Lambda:\Lambda\cap O^n]\), the three factors are
\(I^r,\ |Q|^{-1/2},\ I^{-r/2}\), respectively. This proves
the asserted formula, independent even of the auxiliary \(N\).

For the extension off unit flags it matters that the
Appell--Humbert Gaussian itself need not depend holomorphically
on the period. Work instead in the classical holomorphic
trivialization on the varying torus. The matrices \(P_0,Q_0\)
are holomorphic in the flag coordinate: their entries are either
entries of the inverse flag matrix applied to a fixed period,
or contractions of a flag column against a fixed vector by the
Hermitian form linear in the flag column. Thus \(h,\Omega\) are
holomorphic. The characteristic series without its first Gaussian
defines a holomorphic family of sections on the torus in these
coordinates. Its integral-period automorphy factors, read by
shifting the sum, are holomorphic in \(\Omega,\zeta\); the series
and its derivatives converge normally on compacta, by Gaussian
decay for \(Y\) uniformly positive. Rigidify its line at zero
in this holomorphic trivialization. On each fibre changing to the
Appell--Humbert frame multiplies that function by exactly the
Gaussian in \(T\): it has value one at zero and takes the
classical factors of automorphy to the Appell--Humbert factors.
On the tautological family of period tori this gives the same
relative rigidified polarized line as the pullback of its
biextension by \((1,w_0)\). Indeed their ratio is trivial on
every fibre by the computation of the two Appell--Humbert data;
its
pushforward is a line by cohomology and base change in degree
zero and evaluation is an isomorphism by fibrewise evaluation.
The rigidification trivializes that base line.

Apply the geometric torsion projections and harmonic jet to this
holomorphic characteristic section. These operations are
holomorphic in the period: the torsion points are locally
holomorphic sections of finite étale groups; the formal
biextension frame of Lemma~\ref{bfam:rule} can be constructed
over the local holomorphic base by the same finite-order
operations. The preceding fibre calculation identifies the
result with the analytic marked jet, up to the displayed
nonzero index factors, \(d_L\) and \((\det h)^{n/2}\), all
holomorphic in this chart. Lemma~\ref{dg:gaussian} also proves
this directly by removing the conjugate-type quadratic terms with
the harmonic contractions. This proves holomorphy of the
analytic jet without differentiating a nonholomorphic
\(Y^{-1}\) in the period. Its flag character follows in the
same coordinates. A triangular flag change gives the corresponding
triangular changes in the two labelled tangent blocks (first
columns in one block, last inverse rows in the other). It changes
both period matrices by the same tangent substitution, hence
fixes \(\Omega\); the determinant factor gives the scalar
weight \(n/2\), and the labelled derivatives give the character
of the row-highest polynomial. These are precisely the two parts
of the geometric weight. The geometric reading of a regular
weight rule is holomorphic with that character by the good framed
uniformization (its section property).
Finally all
cross products of the two rows of readings vanish on unit flags
and hence on every flag by the weight and Gram--Schmidt. Where
one analytic reading is nonzero its quotient with the geometric
reading is holomorphic; the quotient has the same absolute value
as on unit flags, hence is constant. Holomorphic continuation
gives the claimed equality everywhere on the component. When
all analytic entries vanish the same reasoning on unit flags
gives zero throughout for the geometric entries as well.
\end{proof}
\src{src/Formalism/LocalToGlobal/ProjectiveLattices/PetersonNormSubspace.lean}{75}{The section property and period identification in a good uniformization};
\src{src/Development/AbelianVarietyTheory/Proofs/SiegelThetaTransformation/PeriodCoordinateThetaOfSignedHeisenbergVacuum.lean}{451}{The half characteristic and determinant in an adapted basis};
\src{src/Development/AbelianVarietyTheory/Proofs/SiegelThetaTransformation/ExistsSymplOfSelfDual.lean}{226}{The probability normalization on the period torus};
\src{src/Development/AbelianVarietyTheory/Proofs/SiegelThetaTransformation/ShortThetaValueComplexPointEqClassicalTheta.lean}{194}{The product of generators and its covariant derivatives};
\src{src/Development/AbelianVarietyTheory/Proofs/SiegelThetaTransformation/ShortSeriesReadingsEqClassicalThetaSeries.lean}{88}{The separated nonzero constants in the actual chart comparison}.

\begin{theorem}[Integral field generation on the marked period charts]\label{bfam:oncharts}
Let \(p,q\ge1,\ m=fr+1,\ n=8m\). Suppose the actual level-\(M\)
PEL cover has at each good embedding a good framed uniformization
by the finitely many marked periods of its common rational
Hermitian space, as in Lemma~\ref{munf:marks} and
Proposition~\ref{bfam:unitread}. Use its identity auxiliary map
and the positive dominant harmonic weight of size \(S=d_P+frn/2\).
Assume the cusp space is read as bounded holomorphic forms of that
weight invariant under the marked principal stabilizers, the raw
rules have the prescribed cusp condition at all good embeddings,
and their readings satisfy the half-maximum hypothesis of
Corollary~\ref{bfam:metric}. Corollary~\ref{crev:good} verifies
these analytic conditions for the selected good covers and their
rule spaces. Use the semiabelian integral tests
of Theorem~\ref{intjet:semi} and the field, level and weight size
bounds of Corollary~\ref{bfam:metric}. Then each algebraic
component has a finite set of integral cusp sections spanning
its cusp space over \(F\). Their logarithmic Petersson norms at
every embedding are at most
\[
                 C(frL)^2 S+2\log D_M .
\]
The same estimate holds for simultaneous generators of the
untwisted full cover with its probability measure.

For the final weight \(\lambda\), put \(W=\sum_{\tau,i}
|\lambda_{\tau,i}|\), \(j_\tau=\lambda_{\tau,p}-1-4m\)
(algebraic weight rows numbered from one, with \(p\) the Lie
multiplicity at \(\tau\)). The first Hermitian block in its period
chart has \(p_0=q\) rows. Use the actual weight decomposition
with \(S\le(1+r)W+rf(1+4m)\), and the key section on the tests of
Theorem~\ref{ctwist:semi}. If \(j\ne0\), assume the horizontal
parallel determinant sections and their product have the common
vertical orders on each local component and the corrected degree
estimate of \eqref{ctwist:degree}, with the stated ramification-one
tests outside \(B\). Their regular trivializations and inverses
with constant pointwise norms preserve the rule cusp conditions
at all good embeddings by Corollary~\ref{crev:twist}. Under the stage bounds
for \(B\), the coefficient field and \(D_M\), the resulting
integral sections span on each component and satisfy
\[
           \log\|f\|_\sigma
            \le C'\bigl((frL)^3 W+(frL)^5\bigr).
\]
If every \(j_\tau=0\), the conclusion holds without the parallel
determinant or degree hypotheses. For \(j\ne0\) the same conclusion
also follows from the one-sided product criterion after
\eqref{ctwist:degree}, provided its lower degree bound gives the
same cost in \eqref{bfam:flat}. On the indefinite principal
family of Proposition~\ref{host:art} with its selected marked
PEL moduli and tests, Theorem~\ref{cmeig:splitmodel} supplies the
common orders and degree for every \(j\). Outside the excluded
primes the original level is hyperspecial, so integrality of
the common order already follows from
Proposition~\ref{cmeig:hyperspecial}.
\end{theorem}
\begin{proof}
First consider the row multiplier. A rational isometry maps the
unit vacuum of a row to a unit scalar times that of the transported
row, by preservation of the specified semicharacter and the
one-dimensional vacuum line. Conjugation also carries the
projectors to those for the transported row, so this scalar is
independent of the column. On the marked stabilizer it is one
at the rank used here. Indeed in a symplectic basis of the
\emph{shifted} row lattice the integral symplectic matrix
preserves its half characteristic, since the isometry commutes
with \(w_0\). For one characteristic theta the residual factor,
after the square root of the Siegel determinant and the quadratic
exponential, is an eighth root of unity. This follows on the
integral symplectic generators. A symmetric integral shear
changes the characteristic by the shear and its half diagonal;
the remaining exponential has argument \(\pi i\) times a
quarter-integer. An integral linear basis change reindexes the
sum, with at most a square root of \(-1\) from its determinant.
Fourier exchange gives \(e^{-\pi i g/4}\) times
\(e^{\pm 2\pi i\alpha^t\eta}\) and the displayed Gaussian
transform in Lemma~\ref{dg:gaussian}'s coordinates.
Integral changes of representatives of the characteristic give
signs, and choices in composing square roots give signs.
Thus the eighth-power assertion persists under composition.
The determinant factors cancel those of the period matrices in
the normalized theta, as in Proposition~\ref{dg:identity}. The
\(n\) identical rows contribute the \(n\)-th power of the
residual factor, equal to one since \(8\mid n\).

The marked translate introduces no extra central character here.
Indeed the residue label is the identity by
Corollary~\ref{bfam:mark}. Both \(y,\gamma y\) lie in
\(2\widehat\Lambda\otimes\bar{\mathfrak m}\widehat L\) for a
stabilizing \(\gamma\), and their residues agree at
\(\mathfrak m\). Their difference lies in \(2D\), by the lift
calculation in Lemma~\ref{bfam:phase}. The vacuum character on
this difference is one. The Weyl phase of the two translations is
also one since their half pairing is integral in the unshifted
self-dual lattice. Hence \(\rho(\gamma y)\) and \(\rho(y)\)
give the very same marked vector. This establishes compatibility
with the invariant geometric readings in the hypotheses. For
arbitrary even \(n\ge8fr+4\) the first conclusion also follows
whenever the geometric and vacuum stabilizer characters match;
divisibility by eight proves this particular matching at the
rank of the final weight composition.

Lemma~\ref{munf:marks} now supplies distinct marked row orbits.
Apply Proposition~\ref{munf:span} with the actual graph identity
of Proposition~\ref{dg:identity}, after dividing each component
by its constant \(C_k\). Equations~\eqref{bfam:reading} and the
nonvanishing of \(u_{\Lambda,N}\) give simultaneous generation.
The sum covers the entire genus of the standard column with its
prescribed level; a single representative need not suffice.
On each chosen column Proposition~\ref{bfam:harm} replaces the
polynomials by its finite spanning set. Lemma~\ref{bfam:rule}
and faithful scalar extension descend generation to \(F\).
Multiplication by the respective reflex scalars and the common
\(D_M^2\) preserves the field span. Theorem~\ref{intjet:semi}
gives integrality, and Corollary~\ref{bfam:metric} gives the
simultaneous norm bounds, including for the component probabilities.

On each component multiply further by its integral key section
of weight \(-4m\) and its integrally rescaled parallel determinant.
Theorem~\ref{ctwist:semi} and equations
\eqref{ctwist:degree}, \eqref{bfam:flat} give the combined
logarithmic cost at most
\[
  2frm\log(2\pi)+12rL\bigl(W+f(1+4m)\bigr)
  +\sum_{\ell\in B}\log\ell+
                    \frac{\log|\Delta_F|}{2[F:\mathbb Q]} .
\]
The last two quantities are \(O((frL)^2)\) under the stage
bounds. Also
\[
 2\log D_M=2(frn+d_P)\log M+2d_P\log\delta_E
                     \ll fL S
\]
by the level and discriminant bounds in
Theorem~\ref{intjet:semi}'s application.
Inserting the stated bound on \(S\) gives the result.
For \(j=0\) the parallel section is identically one and has
neither an order nor a metric cost. If only \(\sum_\tau j_\tau=0\),
formula~\eqref{ctwist:degree} cancels the common degree once
the common finite orders and common corrected degree have been
obtained. Component restriction and extension by zero preserve
the field decompositions and the integral evaluation tests.
\end{proof}
For every selected cover satisfying \(Y.\mathrm{Good}(e,o)\),
Corollary~\ref{crev:good} supplies the marked period, cusp,
half-maximum and measure clauses with the stated stage parameters.
For nonzero parallel weights
Theorem~\ref{cmeig:splitmodel} supplies the local orders and the
arithmetic degree on those marked components. Neither the finite
genus identity nor the positive real reproducing scalar depends
on that degree.
\src{src/Development/AbelianVarietyTheory/Proofs/SiegelThetaTransformation/KernelIdentity.lean}{181}{Doubled identity with the level tests and harmonic insertions};
\src{src/Development/AbelianVarietyTheory/Proofs/SiegelThetaTransformation/HermitianThetaSeriesSpanBoundedForms.lean}{168}{Inversion and reproduction on marked component families};
\src{src/Development/AutomorphicArithmetic/ThetaCorrespondence/ThetaPairLifts/HodgeBundleWeights/ShortSpanCompositionClosed.lean}{69}{The two actual component twists}.

\paragraph{The cusp test on the principal marked charts.}
Retain the period chart and identity Fourier mark of
Propositions~\ref{bfam:unitread} and \ref{bfam:oncharts}.
The following argument uses the mark at one prime over \(M\).
In particular it applies before summing over the definite column genus.

\begin{lemma}\label{crev:mode}
For every rational isotropic row line and every orthonormal frame on
the chart, the real diagonal transvection restriction of each marked
theta jet has mean zero along any period with trivial multiplier.
Its full real cusp-torus average is zero as well.
\end{lemma}
\begin{proof}
Write \(v\ne0\) for the isotropic row and \(I=(Ev)^n\) in the
rational trace symplectic space of the oscillator. Choose a rational
Lagrangian in one row containing \(Ev\), and take its \(n\)-fold sum
\(D\). Complete it to a symplectic polarization \(D\oplus D^*\),
using \(D^*\) for the position variable. Translations by \(a\in
D(\mathbb A_f)\) are multiplication by the finite additive character
of \(z(a)\) at a position \(z\). In particular they act as the identity
at all \(z\) annihilating \(I\).

Let \(\varphi\) be the marked finite vector in these coordinates.
The finite translation calculation preceding Lemma~\ref{bfam:genus}
gives an element of \(I(\mathbb A_f)\) acting on the single marked
vacuum with eigenvalue different from one. Indeed, at a prime of
\(\mathfrak m\) its first \(r\) residues are a basis and a primitive
vector on the conjugate line, divided by the local half-level,
detects one of them. The unshifted character of the isotropic line
is one. It follows pointwise that
\[
                 \varphi(z)=0\qquad (z|_{I(\mathbb A_f)}=0).
\]
This also holds for a sum of the indicated marks, term by term. The
calculation takes place after changing polarization, which intertwines
the Weyl translations by the Fourier and substitution laws in the
proof of Proposition~\ref{bfam:unitread}.

In this polarization the real diagonal flow is a shear. Its
infinitesimal action \(N\) has image \(I_{\mathbb R}\) and kernel
\(I_{\mathbb R}^{\perp}\), and it fixes \(D_{\mathbb R}\) pointwise.
The shear contributes the character with quadratic exponent,
up to the fixed Fourier sign, \(Q(z)=\tfrac12\langle \widetilde
z,N\widetilde z\rangle\), where \(\widetilde z\) is a lift of the
position and the pairing is symplectic. At each place the trace
polarization applied to the transvection
\(\kappa_0 h(y,v)v\) gives a nonzero real multiple of
\(\sum_{j=1}^n|h_\tau(y_j,v)|^2\). The multiples have the same sign:
they are obtained from the product of \(\tau\xi_0\) and
\(\tau\kappa_0=\pm2\tau\xi_0\) with one fixed choice of sign, and
the scaling \(k_0\) of the real parameter is positive. Consequently
\(Q(z)=0\) exactly on the annihilator of \(I_{\mathbb R}\).

More explicitly, on the real flow the first period matrix is fixed
because the shear fixes the chosen first Lagrangian. The second
matrix changes by the first times the real symmetric shear matrix.
Thus \(\Im\Omega\), the holomorphic argument \(hz_j\), the Gaussian
prefactor and \(\det h\) in Proposition~\ref{bfam:unitread} stay
fixed. The change in each summand is exactly its quadratic phase,
also after applying the polynomial in the \(z_j\)-derivatives.
The real restriction therefore has the expansion
\[
            \sum_{\ell\in D^*(\mathbb Q)}
                 \varphi(\ell)\,A_x(\ell)\,e^{2\pi i s Q(\ell)}
\]
with the Fourier sign absorbed in \(Q\). At the fixed frame the
archimedean coefficient is a polynomial times a decaying Gaussian,
and the compact finite support restricts rational positions to a
fixed denominator lattice. The series converges absolutely and
uniformly for real \(s\). If a rational position annihilates \(I\)
at infinity it annihilates \(I\) at all finite places. Hence every
zero-frequency summand vanishes. Averaging over \([0,T]\) and
taking \(T\to\infty\) gives zero by dominated convergence, even
when individual frequencies are not characters of the chosen
period circle. For a periodic sum this is its period mean.
Real translates of the starting frame give the same conclusion.
Integrating along the diagonal circles in the full cusp torus now
gives the torus average by translation invariance of Haar probability.
\end{proof}

\begin{proposition}\label{crev:half}
On each of these principal period charts the raw theta readings at
\(n=8(fr+1)\) are bounded in Petersson norm. For every nonzero
reading each orthonormal frame attaining at least half the component
supremum has row minimum squared at least \(\varepsilon\) of
Proposition~\ref{bergvol:global}. The same minimum conclusion applies
to every nonzero bounded cusp form with trivial transvection multiplier on
the chart. In particular the cusp and half-maximum tests on the raw
rules in Theorem~\ref{bfam:oncharts} follow from its stated marked
period uniformization.
\end{proposition}
\begin{proof}
Write \(a(x)=\min_{0\ne v\in L'}\sum_\tau{\rm maj}_{x,\tau}(v)>0\).
Lemma~\ref{bfam:minimum} and Proposition~\ref{intjet:arch}, at
orthonormalized frames, first give a bound of the form
\[
             \|P(x)\|_{\rm Pet}\le C(1+a(x)^{-1})^B .
\]
Indeed \(\lambda^2\ge 2a(x)/R_E^2\), the theta generator has
mean square one, and the directional norm and its weight correction
are fixed on orthonormal frames. Here \(C,B\) may depend on the
fixed section and stage. Multiplication by its nonzero fixed
reflex scalar does not affect this assertion.

Put \(H=2M\sqrt d\). For \(a=a(x)<\min(1,H^{-1})\) choose a
shortest \(v\). Its integral self-pairing vanishes by the norm
argument following Proposition~\ref{bergvol:global}. The
transvection there acts trivially modulo \(M\) on \(\mathfrak a L'\),
so lies in the actual marked stabilizer by Lemma~\ref{munf:marks}.
Its multiplier on the raw sections is trivial by the eighth-power
and identity-mark calculation in Theorem~\ref{bfam:oncharts},
which uses only the transformation formula. Take \(H\) as period
in the scaled coordinate, and write \(G(\zeta)\) for the raw
coefficient along this flow. It is holomorphic on
\(\Im\zeta>-1/a\) by the leading-minor calculation above.

We give the majorant comparison needed to apply the polynomial bound
before knowing a global supremum. At a single place set
\(l=\|\tau v\|_x^2\le a\). Rotate the positive and negative blocks
so the two components of \(\tau v\) lie on their first coordinate
axes, each of norm \(\sqrt{l/2}\). These rotations preserve the
central majorant. In this plane the new positive line has positive
and negative coordinates \(A=1-i\zeta l/2,\ B_1=-i\zeta l/2\),
with \(|A|^2-|B_1|^2=1+l\Im\zeta\). All orthogonal directions
retain their majorant. The two comparison eigenvalues in the plane
are \(K\) and \(K^{-1}\), where
\[
 K=\frac{(|A|+|B_1|)^2}{1+l\Im\zeta}
    \le\frac{(1+l|\zeta|)^2}{1+l\Im\zeta}.
\]
Indeed a unitary change of phase makes the normalized line coordinates
nonnegative and its orthogonal negative mate interchanges them.
At \(\zeta=s-i/(2a)\), \(0\le s\le H\), this gives \(K\le 13\):
the denominator is at least \(1/2\) and
\(l|\zeta|\le Ha+1/2\le3/2\).
Thus the new row minimum is at least \(a/13\).
The leading-minor variation estimate preceding
Lemma~\ref{bfam:harm} bounds the conversion from Petersson norm
to raw coefficient on this line by \(2^{w/2}\).
It follows that
\[
      \sup_{0\le s\le H}|G(s-i/(2a))|
                    \le 2^{w/2}C(1+13/a)^B .
\]
For \(y=\Im\zeta\ge0\), \(0\le s\le H\), the same comparison
gives \(K\le(1+a(H+y))^2\). The leading-minor formula bounds
the raw metric conversion by \((1+ay)^{w/2}\).
Together these give polynomial growth of \(G(s+iy)\) in \(y\)
for the fixed starting frame. In \(q=\exp(2\pi i\zeta/H)\),
Cauchy's coefficient estimates therefore remove all negative
Laurent powers. Lemma~\ref{crev:mode} removes the constant power.
The maximum principle for \(G/q\) on the disk of radius
\(\exp(\pi/(aH))\) gives
\[
      \|P(x)\|_{\rm Pet}=|G(0)|
       \le 2^{w/2}C(1+13/a)^B e^{-\pi/(aH)} .
\]
This is bounded as \(a\downarrow0\). On the complementary range
the original polynomial estimate is bounded. Hence every nonzero
section has a finite positive supremum and one can choose an
orthonormal frame above half of it. With boundedness established,
the half-maximum argument following Proposition~\ref{bergvol:global}
applies with its original \(\varepsilon\); the qualitative constants
\(B,C\) disappear. For a general bounded cusp form the real torus
and positive-frequency argument there supplies its diagonal zero
mode. Thus the same quantitative conclusion holds for that space
at every weight whose transvection character is trivial.
\end{proof}
The proof uses the short vector of the component under consideration.
It gives the absolute-constant bound of Corollary~\ref{bfam:metric}
for all selected parameters, all permitted column lattices and
auxiliary integers at every good embedding; it requires neither
a common fundamental domain as the rank changes nor a uniform
constant in the preliminary polynomial growth estimate.
\src{src/Development/AbelianVarietyTheory/Proofs/CuspFormEvaluations/ExistsExpansionOfEqSumThFr.lean}{26}{The source's different zero-mode test using positivity of the last row weight};
\src{src/Development/AbelianVarietyTheory/Proofs/DivisorProofs/FirstMinAtHalfMax.lean}{111}{Transfer to the minimum at the selected level}.

\begin{corollary}\label{crev:twist}
On the good marked charts the algebraic rule space defined by
testing all good complex readings requires no additional cusp
comparison for the raw spanning argument. A regular unit section on
a component with constant positive Petersson norm, in mutually dual
multiplicative metrics for it and its inverse, preserves this cusp
condition in both directions. This applies to the key and the
horizontal determinant whenever they exist on that component with
their stated constant metrics.
\end{corollary}
\begin{proof}
In the rule definition the cusp subspace is the intersection of
the regular weight rules with the inverse images, under every good
reading, of the complex spans of the analytic cusp forms. Each such
span is the analytic cusp space itself: holomorphy, transformation
law and vanishing of the torus integrals are linear, and each
finite sum of bounded forms is bounded. Conversely inclusion in the
span holds for each member. Lemma~\ref{bfam:rule},
Lemma~\ref{crev:mode} and Proposition~\ref{crev:half} place the
raw rules in this intersection. For generation one only needs to
express the reading of every member in the span of the raw readings,
which is the conclusion of Proposition~\ref{munf:span} on those
charts. Injectivity after scalar extension in Lemma~\ref{bfam:rule}
then descends the span to the coefficient field.

For the multiplication assertion use a common finite-index
transvection period lattice killing both characters. Passing to
it leaves each normalized constant term unchanged when applied to
a function periodic for the original lattice. Along a rational
isotropic row the positive-ray metric estimate following
Proposition~\ref{bergvol:global} and boundedness imply that
every nonzero Fourier coefficient of either factor has all
coordinates nonnegative, by growing one imaginary coordinate
at a time. At height zero the coefficients are absolutely summable
by smoothness on the real torus. Thus the product constant is the
product of the two constants, by convolution: two vectors with
nonnegative coordinates can sum to zero only if both are zero.
This proves cusp vanishing for the product, and boundedness follows
by multiplicativity of the metric. Apply the same argument to the
inverse section, which is holomorphic and has the reciprocal
constant norm. In particular its Fourier support is tested afresh
using holomorphy in the same upper tube, even though its \emph{weight}
has the opposite sign. The polynomial growth bound uses only the
weight spread. This treats the Fourier condition in the rule space;
a separate use of a subcanonical sheaf on a compactification must
identify its restriction with that space.
\end{proof}
\src{src/Formalism/LocalToGlobal/ProjectiveLattices/CuspFree.lean}{14}{All good embeddings in the definition of the cusp rules};
\src{src/Formalism/LocalToGlobal/ProjectiveLattices/JFDFree.lean}{42}{The span of the bounded holomorphic cusp readings}.

\paragraph{The flag measure of the selected unitary genus.}
Let \(F_0=E^+\). Use the host fields of Proposition~\ref{host:art},
the genus of the self-dual row, and its odd principal level \(M\ge3\).
Write \(K\) for the finite adelic stabilizer and \(K(M)\) for its
principal subgroup. Let \(c_r\) and \({\rm mass}_I\) be as in
Proposition~\ref{bergvol:compactmass}.

\begin{proposition}\label{crev:mass}
For the full adelic packet of marked rows, the sum of the component
flag volumes in the normalization of Lemma~\ref{bergvol:local} is
\[
        V=[K:K(M)]\,\frac{c_r^f\,{\rm mass}_I}{(2\pi)^{fr}}
           \le V_* .
\]
Adelic probability on flag-invariant functions pushes forward to
the sum of the fundamental-set measures divided by \(V\).
Conditioning on a nonempty union of components gives the same
formula with its own total flag volume.
\end{proposition}
\begin{proof}
In the marked period description of Lemma~\ref{munf:marks}, the
finite double cosets of \(G=U(h_0)\) represent the marked rows of
this genus. Indeed finite translates of the completed reference
lattice and its basis give a row by intersection with the rational
space. Conversely choose local isometries to a specified row in
the genus. At primes of \(M\) the extension is split and the mark
is a basis on one half with its pairing-dual on the other. An
invertible matrix over the local residue quotient lifts to an
integral invertible matrix by lifting entries. Thus the local
isometries may match the marks and the stabilizer is \(K(M)\).
In terms of the two algebraic Lie ranks \(p,q\), the real flag
space at a place is \(U(p_0,r-p_0)/U(1)^r\), with \(p_0=q\).

Use the rational form \(J=\operatorname{diag}(1_{p_0},-1_{r-p_0})\) given
by Lemma~\ref{cmeig:row}, and compare with \(G_c=U(1_r)\), the
group of the positive definite standard form. At every finite
place the two self-dual lattices are isometric, including residue
characteristic two, by the local part of that lemma. Keep track of
the Haar densities in this assertion. On the Lie algebra for a
Hermitian matrix \(J\), left translation identifies an invariant
top differential by identifying its tangent at one with the fixed
\(F_0\)-space of skew-Hermitian matrices via \(X\mapsto JX\).
Choose the same rational top differential on that vector space
for both \(J\) and \(1_r\). Under an isometry \(P^*JP=1_r\),
the comparison in these coordinates is \(A\mapsto P^*AP\).
Its determinant on the skew-Hermitian space is
\(N_{E/F_0}(\det P)^r\): extend scalars to the split algebra,
where the two factors multiply an unrestricted \(r\)-by-\(r\)
matrix on its left and right. Its absolute value is one at each
finite place because \(N(\det P)=1/\det J=\pm1\).
Thus the corresponding maximal stabilizers have the same local
volume. No choice of their local bases contributes a discriminant.

The global scalar and convergence factors in the Tamagawa measure
are the same: the base field and dimension agree and the geometric
character lattice in both cases is generated by determinant with
the quadratic Galois action of \(E/F_0\). At each real place the
common gauge in the \(A\)-coordinates is the same scalar times
Lebesgue measure on the imaginary diagonal and the complex
strictly lower entries. Multiplication by \(J\) only changes signs
in these coordinates. Divide out that same scalar for both groups.
Use volume one for the maximal stabilizers at finite places and
the resulting real coordinate densities. In these choices the
total adelic covolumes of the two groups are equal: they are
obtained from Tamagawa measure by a common scalar and both
Tamagawa numbers are two for a nonzero Hermitian space over a
quadratic extension. See \href{https://doi.org/10.1007/BF02391774}
{Weil, \emph{Sur la formule de Siegel dans la th\'eorie des groupes
classiques}}, and Ichino, \emph{A regularized Siegel--Weil formula
for unitary groups}, \S4. For the definite group the real volume is
\(c_r^f\). The finite double cosets at maximal level are its
standard lattice genus: cutting the quotient over such a coset
divides this common real volume by its finite automorphism group.
Summing gives \(c_r^f\,{\rm mass}_I\).

For the indefinite group, the quotient of the real density by
angle measure \(d\theta_1\cdots d\theta_r\) is the stated flag
measure. At the central flag the tangent projection takes the
strictly lower entries of the Lie algebra with Jacobian one.
At a lower unipotent matrix \(u=k b\), with \(k\) unitary and
\(b\) upper triangular, the projected tangent in this frame is
the strictly lower part of \(b(u^{-1}du)b^{-1}\).
The unipotent differential has determinant one in its entries.
Conjugation on matrices modulo upper triangular matrices has
complex determinant \(\prod_{i>j} b_{ii}/b_{jj}\). Its squared
absolute value is \(\prod_{i=1}^{r-1}D_i^{-2}\), using
\(|b_{ii}|^2=D_i/D_{i-1}\) and \(D_0=D_r=1\).
This recovers the chart density at every point.

Over each finite double coset \(gK(M)\), the quotient of
\(G_\infty\times gK(M)\) is by
\(\Gamma_g=G(F_0)\cap gK(M)g^{-1}\). Cutting the real factor
by a fundamental set for \(\Gamma_g\) gives the common finite
volume \(1/[K:K(M)]\) times real quotient volume. There is no
stabilizer on the flag: it would be discrete in a compact torus,
hence finite, and the principal subgroup on the underlying
integral period lattice is torsion-free. For the latter, at an
odd prime dividing \(M\) the powers of a nonidentity
\(1+\ell^a B\), \(a\ge1,\ B\not\equiv0\bmod\ell\), have depth
\(a+v_\ell(t)\) at exponent \(t>0\), by the binomial theorem.
Integrating along the angle tori therefore gives \((2\pi)^{fr}\)
times the flag integral. The same disintegration applies to any
nonnegative measurable function of the flag. Applying it first
to one proves the identity and then dividing by the total mass
proves the probability formula. The congruence index is at most
\(M^{2fr^2}\) by reduction on the row, so
Proposition~\ref{bergvol:compactmass} gives the required bound.
\end{proof}
In particular both analytic tests in
Proposition~\ref{bergvol:global} are available on the selected
principal packet for the bounded analytic cusp space of trivial
transvection character. Using this for an algebraic component
uses its identified framed period charts. For the rule space the
reading is as in Corollary~\ref{crev:twist}; agreement with a
separately defined toroidal sheaf is a geometric comparison.

\begin{corollary}[Analytic estimates on a selected good cover]\label{crev:good}
Take a stage in the selected range on the host of
Proposition~\ref{host:art}, a core with the prescribed host,
polarized datum and level, and the deep cover \(Y\) with
\(Y.\mathrm{Good}(e,o)\). Use its identity auxiliary rule
at \(n=8(fr+1)\), and the parameters and semiabelian tests of
Theorem~\ref{bfam:oncharts}. At every good embedding the raw
rules are cusp rules on this \(Y\). The short integral rules in
that theorem span the untwisted space on the cover and on each
component, with logarithmic Petersson norms at most
\(C(frL)^2S+2\log D_M\) for the stated stage size bounds.
For the component cusp spaces at the frame weights of
Corollary~\ref{bergvol:size} the normalized evaluation kernel
is at most \(\beta_*\).
\end{corollary}
\begin{proof}
Here \(Y.\mathrm{Good}(e,o)\) contains the following data. Its
\code{unf} gives the complex framed period uniformization at each
good embedding, including the section transformation law and the
Haar mean-square-one normalization of the theta generator on unit
frames. In these coordinates the first \(p_0=q\) Hermitian
directions are positive. Their conjugate, reversed tangent block
has Lie rank \(q\), and the other tangent block has Lie rank \(p\)
at the labelled \(\tau\). Thus the pivot \(\lambda_{\tau,p}\)
in Theorem~\ref{bfam:oncharts} is numbered in the algebraic
weights, while the positive block used for the majorant ends at
\(p_0\) in the Hermitian flag.

The \code{modelW} and \code{xiAbs} clauses give the signed host
Gram matrix and the absolute polarization coefficient at every
place. The common lattice embedding and alternating pairing in
\code{adPin} identify the original period lattice, its finite
adelic translates, the half-level mark and the eigenframe
coordinates on these same components. In particular stripping the
fixed host ideal from a translate of the reference period
\(\mathfrak a L'_*\) gives the translated Hermitian self-dual row:
finite unitary matrices commute with \(\mathfrak a\), and
localization detects self-duality. This is \(L'\) in the return
count and in Lemma~\ref{bfam:minimum}. The \code{lev} clause gives
auxiliary integer one on the deep variety; the spanning construction
there uses its identity auxiliary map of Corollary~\ref{bfam:mark}.
It does not replace the original map from the deep variety to the
base by an identity.

The \code{dictU} and \code{levelK} clauses identify the rational
and right-principal-level equivalences with the chart action and
exhaust the real flag orbits. Indeed the adelic frame map is
equivariant modulo arithmetic translates under the rational and
level actions and intertwines the real action; it hits each flag
modulo its upper triangular frame change. Equality of its frames
modulo the arithmetic action forces rational and finite-level
equivalence. The unit frames over each flag form a single orbit
under its diagonal angle group. Together with the marked-period
classification in Lemma~\ref{munf:marks}, this is the packet and
the probability disintegration of Proposition~\ref{crev:mass}.
The \code{unit} clause gives leading Gram determinants of absolute
value one at the adelic frames, as used in their metric reading.

Finally \code{chiTriv} gives trivial arithmetic multiplier at the
dominant natural weights on each good chart. It gives the same
for every integral flag weight: for fixed arithmetic element its
multiplier is multiplicative in addition of weights and never
zero. Take natural weights with \(n=0\) that are one on an
initial segment at one labelled embedding and zero elsewhere.
Consecutive differences give all individual weight coordinates,
with reversed order and a minus sign on the dual block. Their
integral span is the whole flag weight lattice. In particular
there is no character loss in the transvection period for the
frame weights. Apply Propositions~\ref{crev:half} and
\ref{crev:mass}, then Proposition~\ref{bergvol:global}.
The actual all-weight space \code{cuspW} is the intersection of
the regular rules with the good-reading tests of
Corollary~\ref{crev:twist}; at a natural weight it is the ordinary
rule space \code{cusp}, and \code{cuspOn} adds only the component
support condition. Thus Theorem~\ref{bfam:oncharts} applies with
its integral tests and no additional analytic cusp assumption.
\end{proof}
\src{src/Development/AbelianVarietyTheory/Definitions/Divisors/NormNegCmConstantEqAtAllEmbeddings.lean}{187}{The clauses of the selected deep cover};
\src{src/Formalism/LocalToGlobal/ProjectiveLattices/FamilyUniformisationBasePointOrthogonalBis.lean}{97}{The marked finite translates in the period identification};
\src{src/Development/AbelianVarietyTheory/Definitions/Foundations/IsProbabilityMeasureMuPart.lean}{425}{The actual all-weight rule space};
\src{src/Development/AbelianVarietyTheory/Counterexamples/LevelGroupExamples/AdelEqZeroOfOffDom.lean}{77}{Generation of the arithmetic characters by natural weights}.

Constructing such covers at all scales is a separate existence
statement. For example, the sized-cover construction starts with
the pinned core, a stage in range, the host splitting bound at least
five and \(15\) coprime to the base level. It constructs the half-level,
uses PEL moduli representability, forms the coefficient compositum
with a ray class field, and builds the marked uniformization; the
field-size bound uses the ray-class-field estimate. If the comparison
to a prescribed base frame requires \(Y.\mathrm{Match}\), that
condition concerns the same base embedding and base frames and is
not part of \(\mathrm{Good}\). The corollary asserts the estimates
on the selected \(Y\); it does not construct the core, the compatible
family or its toroidal compactifications.
\src{src/Development/AbelianVarietyTheory/Proofs/DivisorProofs/ExistsGoodCoverOverSizedField.lean}{111}{Inputs and construction of the sized good level};
\src{src/Development/AbelianVarietyTheory/Definitions/Divisors/NormNegCmConstantEqAtAllEmbeddings.lean}{207}{The common marked frames included in goodness}.


\subsection{Reproduction on the archimedean frame slab}
\label{arrep:section}

Write $\Psi$ for the CM type, $p=p_0$, $q=r-p$, and $G=G_\tau$
for the Hermitian matrix at $\tau$. The pairing is linear in its
first variable, so the Gram matrix of a column frame $x$ is
$x^tG\bar x$. Put $J=\operatorname{diag}(1_p,-1_q)$. The domain
$\mathcal D$ consists of invertible frames at all $\tau$ whose first $p$
columns span a positive space. This is equivalent to the prescribed
signs of all the leading Gram determinants: after the first $p$
columns, the additional orthogonal directions belong to a negative
definite space. Choose an orthoframe $x_0$ with Gram matrix $J$ at
every place. Let $B$ be the product of the invertible upper triangular
groups, and let $dy$ be ordinary Lebesgue measure in the original
complex matrix coordinates. Set
\[
 S_\delta=\{y:\ |(y_\tau^tG_\tau\bar y_\tau-J)_{ij}|<\delta
                         \text{ for every }\tau,i,j\}.
\]

For each place choose nonnegative integers $c_1,\ldots,c_q$,
$c'_1,\ldots,c'_p$ and write $a_i=\sum_{j\ge i}c_j$,
$a'_i=\sum_{j\ge i}c'_j$. In physical column order the character
$\chi_\lambda(b)=\prod_{\tau,i}b_{\tau,ii}^{\lambda_{\tau,i}}$ has weight
\[
 \lambda_\tau=(-n/2-a'_p,\ldots,-n/2-a'_1,\;
                          n/2+a_1,\ldots,n/2+a_q).
\]
The polynomial letters at each place have $n$ columns and two row
blocks. The $q$ rows are the negative rows in their physical order,
and the $p$ rows are the positive rows in reverse physical order.
A row-highest harmonic polynomial has row degrees $a_i,a'_i$, is
invariant under lower unitriangular row substitutions in each letter
block, and is killed by
$\sum_{k=1}^n\partial_{u_{ik}}\partial_{v_{jk}}$ for every pair
of opposite rows at the same place. These conventions apply to
differential symbols as well as ordinary polynomial variables.

\begin{proposition}\label{arrep:reproduce}
Let $n$ be even, $2r\le n$, and let $P_1,\ldots,P_m$ be any finite
family of such polynomials, with at least one nonzero member. Denote by
$K(x,y)$ the result of applying
$\sum_j P_j(\partial_z)\overline{P_j}(\partial_w)$ at $z=w=0$
to the doubled Gaussian below, where the bar conjugates coefficients.
If $0<\delta\le1/(4(r+1))$, then there is a positive real number
$\rho_\delta$ such that every holomorphic $F$ on $\mathcal D$ with
$F(xb)=\chi_\lambda(b)F(x)$ and bounded absolute value on the
orthoframes satisfies the absolutely convergent identity
\[
         \int_{S_\delta} K(x,y)F(y)\,dy=\rho_\delta F(x)
                               \qquad(x\in\mathcal D).
\]
On orthoframes the pointwise automorphic norm in the bounded-family
construction equals $|F|$. The scalar depends on the specified measure,
archimedean data, family and slab. It applies to every deep level
using those same data.
\end{proposition}

\begin{proof}
First compute what survives the insertion into the Gaussian. Suppress
$\tau$, let $P_+,P_-$ be the diagonal projections to the two physical
blocks, and set $S=-J$, $\kappa=\tau(\kappa_0)$,
$h=\kappa/4$, $t=\pi i\kappa=-\pi\operatorname{Im}\kappa>0$. Put
\[
\begin{aligned}
 L&=P_+\bar y^tG^t+P_-x^{-1},&
 H&=P_+x^tG+P_-\bar y^{-1},\\
 W&=(L^{-1}SL)^tG,& Y&=(hW-(hW)^*)/(2i).
\end{aligned}
\]
The kernel of the bottom rows of $L$ is the positive space spanned
by the first $p$ columns of $x$. The kernel of its top rows is
the negative orthogonal complement of the first $p$ columns of $y$.
They intersect trivially. Interchanging $x,y$ and conjugating gives
invertibility of $H$. The involution $L^{-1}SL$ acts by $-1$ on the
first space and $+1$ on the second. For $v=v_++v_-$ in this
decomposition, the real part of $v^tW\bar v$ is
$-\langle v_+,v_+\rangle+\langle v_-,v_-\rangle$ because the
cross terms have zero real part. Thus $W+W^*$ is negative definite.
As $h/(2i)<0$, the Hermitian matrix $Y=(h/(2i))(W+W^*)$ is positive.

Write $M=L^{-t}GH^{-1}$ and $N=L^{-t}GY^{-1}GH^{-1}$.
For each column $k$ put $\alpha=P_+z_k+P_-w_k$,
$\beta=P_+w_k+P_-z_k$. The actual undifferentiated Gaussian has
prefactor $\prod_\tau(\det L\,\det H)^{-n/2}$ and exponent
$-4\pi\sum_{\tau,k}h_\tau^2\beta^t N\alpha$.
The off-diagonal blocks of $M^{-1}=HG^{-1}L^t$ vanish because
the corresponding products before taking blocks are identity
matrices. Its diagonal blocks are the leading $p$ block of
$x^tG\bar y$ and the trailing $q$ block of $\bar y^{-1}G^{-1}x^{-t}$.
In particular $SM=MS$. Moreover
\[
                     \frac{h}{2i}(SN+NS)=M .
\]
Indeed $W G^{-1}W=G$ and $W^*G^{-1}W^*=G$. For $T=W+W^*$
these give $WG^{-1}T=TG^{-1}W^*$ and therefore
$GT^{-1}W=W^*T^{-1}G$. It follows that
$WT^{-1}G+GT^{-1}W=G$. Multiply this last identity on the left
by $L^{-t}$ and on the right by $H^{-1}$,
using $L^{-t}W=SL^{-t}G$ and $WH^{-1}=GH^{-1}S$.
The latter equality follows from $SM=MS$.

The off-diagonal terms in $\sum_k\beta^t N\alpha$ are linear
combinations of $\sum_kz_{ik}z_{jk}$ and $\sum_kw_{ik}w_{jk}$
for opposite rows $i,j$. Pairing with a polynomial differential
symbol at zero makes multiplication by each such quadratic adjoint
to its mixed Laplacian. Every term containing one of these quadratics
vanishes, also after both sets of derivatives are taken. Only finitely
many degrees of the exponential enter the pairing. Thus one may
differentiate the following explicit expression instead:
\[
 \prod_\tau (\det L\,\det H)^{-n/2}
 \exp\!\left(\sum_{\tau,k}t_\tau
      [\,w_{+,k}^t M_{++}z_{+,k}-z_{-,k}^t M_{--}w_{-,k}\,]\right).
\]
It is holomorphic in $x$ and antiholomorphic in $y$.
The two diagonal blocks of $M(y,x)$ are the adjoints of those
of $M(x,y)$, and the prefactors interchange by complex conjugation.
Consequently $K(y,x)=\overline{K(x,y)}$.

Let $\mathcal U=\prod_\tau\{u:u^tG_\tau\bar u=G_\tau\}$. Under
simultaneous left translation by $u$, $L,H$ acquire right factors
$u^{-1},\bar u^{-1}$, so $M$ and the prefactor are invariant.
Under $x\mapsto xb$, write the two diagonal blocks of $b$ as $A,D$.
Then $L$ acquires left factor $1\oplus D^{-1}$ and $H$ left factor
$A^t\oplus1$. The prefactor gains $(\det D/\det A)^{n/2}$ at this
place. In the exponent the substitutions in $z$ are $A^{-t}$ on
the positive rows and $D$ on the negative rows. In the differential
symbol the resulting row substitutions are $A^{-1}$ and $D^t$.
Both are lower triangular in letter order. The prescribed row
degrees give exactly the remaining factor of $\chi_\lambda(b)$.
This proves equivariance in $x$, and symmetry gives the conjugate
character in $y$.

At $x=y=x_0$ one has $L=x_0^{-1}$, $H=\bar x_0^{-1}$ and $M=J$.
The exponential becomes $\exp(\sum_{\tau,i,k}t_\tau z_{\tau ik}
w_{\tau ik})$. With $P_{j,\nu}$ the monomial coefficients, ordinary
differentiation gives
\[
 C_0=K(x_0,x_0)=
   \left(\prod_\tau|\det x_{0,\tau}|^n\right)
   \sum_{j,\nu}|P_{j,\nu}|^2\,\nu!\,
                       \prod_\tau t_\tau^{|\nu_\tau|}>0 .
\]
Here $|\nu_\tau|$ is the sum of all exponents at that place and
$\nu!$ is the product of their factorials.

We next identify the only holomorphic function that can survive
the compact torus average. In coordinates $g=x_0^{-1}y$, let
$\Delta_{\tau,i}$ be the leading $i$-row principal minor of $g_\tau$.
Define
\[
 H_\lambda(y)=\prod_\tau
    \Delta_{\tau,r}^{\,\lambda_{\tau,r}}
    \prod_{i=1}^{r-1}\Delta_{\tau,i}^{
                              \lambda_{\tau,i}-\lambda_{\tau,i+1}} .
\]
For $r=0$ take the empty product. The weight decreases within each
block. When $0<p<r$ the only internal negative exponent occurs at
$i=p$, where the minor cannot vanish: the projection of a positive
$p$-plane to the standard positive coordinate space has trivial
kernel. Hence $H_\lambda$ is holomorphic throughout $\mathcal D$ and
has value $1$ at $x_0$. Under left diagonal unitary translation in
the $g$ coordinates it has character $\chi_\lambda$, and under right
triangular multiplication it has the same character.

These two transformations determine a holomorphic function up to its
value at $x_0$. Indeed on the big cell near $g=1$ write $g=\ell b$
with $\ell$ lower unitriangular. The two character laws imply
invariance of the restriction to $\ell$ under diagonal unitary
conjugation. On every strictly lower entry $\ell_{ij}$ the simultaneous
circle $\operatorname{diag}(e^{is},e^{2is},\ldots,e^{ris})$ acts with
positive weight $i-j$. The holomorphic power series of the restriction
therefore has only its constant term. Right equivariance determines
the function on a neighborhood, and the identity theorem extends the
conclusion. Indeed indefinite Gram--Schmidt identifies the domain with
the product of the orthoframes and the upper triangular matrices with
positive real diagonal. Both are connected: the unitary group modulo
$U(p)\times U(q)$ is the matrix ball described below, with connected
compact fiber. It follows that
\[
               K(y,x_0)=C_0H_\lambda(y).
\]
For this application the left torus law follows from simultaneous
invariance and right equivariance of $K$ in its second variable.

We check absolute integrability on the slab. The preceding
Gram--Schmidt coordinates write $y=x_0 g b$ with $g$ standard
$J$-unitary and $b$ upper triangular with positive real diagonal.
For our $\delta$, each leading Gram matrix has operator distance
at most $r\delta<1/4$ from the corresponding signed identity. Thus
even the closure of the permitted Gram matrices is compact with all
pivots nonzero and of the required signs. The signed triangular
factorization $b^tJ\bar b$ has a continuous inverse there by the
Gram--Schmidt recursion. The permitted $b$ therefore have compact
closure inside the triangular group. In these smooth coordinates
Lebesgue measure is $dg\,j(b)\,db$, where $dg$ is Haar measure and
$j$ is smooth positive: the original measure is left invariant
under $\mathcal U$ since $|\det u_\tau|=1$. Both $|\chi_\lambda(b)|$
and $j(b)$ are bounded on this compact closure. In particular
$F$ is bounded throughout $S_\delta$.

At a place with $pq>0$ write the standard unitary matrix
$g=\left(\begin{smallmatrix}A&B\\ C&D\end{smallmatrix}\right)$.
All singular values of $A$ and $D$ are at least one by the unitary
relations. The complementary-minor identity and $g^{-1}=Jg^*J$
give $|\det A|=|\det D|$ and bound each leading minor of size
$i>p$ by $|\det D|$ using the complementary trailing minor of
$g^{-1}$. For $i\le p$ the exterior powers of $A$ bound it by
$|\det A|$. As $|\det g|=1$, the product of the minor estimates,
keeping the exact denominator at $i=p$, gives the one-place bound
\[
       |H_{\lambda,\tau}(x_{0,\tau}g)|
          \le |\det A|^{\lambda_{\tau,1}-\lambda_{\tau,r}}
          \le |\det A|^{-n}\le1 .
\]
If either block is empty the group is compact and the function bounded.
For the mixed signature the group modulo the compact block diagonal
subgroup is the ball $Z=C A^{-1}$, $Z^*Z<I_p$, the graph coordinates
of its positive planes. The group is transitive by orthogonal
completion, and $|\det A|=\det(I_p-Z^*Z)^{-1/2}$.
The invariant density on the ball is proportional to
$\det(I_p-Z^*Z)^{-r}\,dZ$. To see this, for an acting matrix with
blocks $a,b,c,d$, set $U=a+bZ$ and $Z'=(c+dZ)U^{-1}$.
Then $I_p-Z'^*Z'=U^{-*}(I_p-Z^*Z)U^{-1}$ and the differential is
$H\mapsto(d-Z'b)HU^{-1}$. Block elimination gives
$|\det U|\,|\det(d-Z'b)|=|\det g_{\rm acting}|=1$, so the real
Jacobian is $|\det U|^{-2r}$. Hence integration of $|\det A|^{-n}$
on the group reduces, up to a finite constant, to integration of
$\det(I_p-Z^*Z)^{n/2-r}$ on a bounded ball. Its exponent is
nonnegative. This proves that $H_\lambda$ is bounded and absolutely
integrable on $S_\delta$, including at $n=2r$. In particular both
$K(x_0,y)F(y)$ and $|K(x_0,y)|^2$ are integrable there.

Finally let $t$ range over the product diagonal unitary torus with
probability Haar measure. Its action on frames is
$t\cdot y=x_0 t x_0^{-1}y$. The holomorphic projection
$F_\lambda(y)=\int\chi_\lambda(t)^{-1}F(t\cdot y)\,dt$ has the two
character laws above and has value $F(x_0)$ at $x_0$. Thus
$F_\lambda=F(x_0)H_\lambda$. The slab and its measure are invariant
under this torus. By absolute convergence and the conjugate torus
law of $K(x_0,y)$, replacing $F$ with $F_\lambda$ in the integral
at $x_0$ leaves its value unchanged. The resulting constant is
\[
                \rho_\delta=\frac{1}{C_0}
                   \int_{S_\delta}|K(x_0,y)|^2\,dy>0 .
\]
Positivity follows from continuity, $K(x_0,x_0)=C_0>0$, and the
open neighborhood of $x_0$ in the slab (with mass one in dimension
zero). Left translation carries the identity and absolute convergence
to every orthoframe by applying it to the translated bounded
holomorphic function. Right equivariance then carries it to every
frame of $\mathcal D$.
\end{proof}

The required invariant finite family can be chosen in these same
variables. At each place form the product of the initial $i$-row
minors to powers $c_i$ in the negative letters and to powers $c'_i$
in the positive letters, using respectively the first $q$ and the
last $p$ columns, nested within each set. Since $n\ge p+q$, these
sets are disjoint. The product over places is a nonzero row-highest
polynomial and all mixed Laplacians kill it. The entire space of
the stated harmonics lies in total degree
$d=\sum_{\tau,i}i(c_{\tau,i}+c'_{\tau,i})$, extending the arrays by
zero. Its dimension is at most $\binom{nr|\Psi|+d}{d}$, the number
of monomials of degree at most $d$. The compact column groups $U(n)^\Psi$ act on opposite letters
by a unitary matrix and its contragredient, preserving the row
conditions and the contractions. An invariant Hermitian product
exists by averaging. An orthonormal basis $P_1,\ldots,P_m$ then has
invariant sum $\sum_j P_j(z)\overline{P_j(w)}$ and positive $C_0$.
Thus the insertion used for spanning bounded forms satisfies the
proposition with its actual weight and with finite $m$.

\src{src/Development/AbelianVarietyTheory/Definitions/SpanningFamilies/IsSystem.lean}{65}{The doubled Gaussian before the two polynomial insertions};
\src{src/Development/AbelianVarietyTheory/Proofs/SiegelThetaTransformation/SetIntegralCompOfStable.lean}{268}{The positive reproducing scalar on the full frame slab}; and
\src{src/Development/AbelianVarietyTheory/Proofs/SiegelThetaTransformation/HermitianThetaSeriesSpanBoundedForms.lean}{168}{Spanning bounded families in the range $4r<n$}.

\subsection{The graph terms of the doubled marked family}
\label{dg:section}

We compute the graph functional in the rational symplectic basis in
which the undoubled sections are defined. Write $g=r|\Psi|$,
$\Omega_x$ for the Siegel period matrix, $Y_x=\operatorname{Im}\Omega_x$,
and $h_x$ for the inverse of the first period matrix. Thus a real
period vector with symplectic coordinates $(u,v)$ has normalized
complex period $u+\Omega_x v$. Use Lebesgue measure on the $ng$ real
position coordinates and mass one for their integral finite lattice.
For one row the Gaussian, apart from the determinant, is
\[
 q_{x,z}(s)=\exp\!\left(\pi i\,s^t\Omega_xs+
          2\pi i\,s^t h_xz+\tfrac{\pi}{2}
                   z^t h_x^tY_x^{-1}h_xz\right).
\]
Put $Q_{x,z}=(\det h_x)^{n/2}\prod_{k=1}^nq_{x,z_k}(s_k)$
and $Q_{x,P}=P(\partial_z)Q_{x,z}|_{z=0}$. These give exactly
the rational theta sums defining $\theta_P(\phi;x)$ with finite
test $\phi$. We use inner products linear in their first variable.

\begin{lemma}\label{dg:gaussian}
For the family and weight of Proposition~\ref{arrep:reproduce},
\[
 \sum_j\langle Q_{x,P_j},Q_{y,P_j}\rangle_\infty
 =c_\infty K(x,y),\qquad
 c_\infty=2^{-ng/2}\prod_{\tau\in\Psi}
     \bigl(|\det G_\tau|\,[-\operatorname{Im}\tau(\kappa_0)]^r
                                                        \bigr)^{n/2}.
\]
This factor is independent of the weight, the polynomials, and both
frames.
\end{lemma}
\begin{proof}
First we show that $Q_{x,P}$ is holomorphic in $x$ with values in
Schwartz space. The physical period map at $\tau$ has components
$\langle x_i,s\rangle$ for $i\le p$ and $(x^{-1}s)_i$ for $i>p$.
Thus $h_x,\Omega_x$ are holomorphic. For the remaining quadratic put
$C=h_x^tY_x^{-1}h_x$. Choose a purely imaginary $\theta\in E$
which generates $E/\mathbf Q$. Such elements exist: writing the
purely imaginary subspace as $\eta E^+$ for any nonzero imaginary
$\eta$, distinct embeddings have distinct restrictions to this
subspace (evaluate first on $\eta$, then on its multiples).
Finitely many proper rational linear subspaces therefore suffice to
exclude all coincidences. Multiplication by $\theta$ in physical
period coordinates is diagonal $D$, with entries $\bar\tau(\theta)$
on the positive rows and $\tau(\theta)$ on the negative rows.
Write its rational blocks in the symplectic basis as
$\left(\begin{smallmatrix}a&b\\ c&d\end{smallmatrix}\right)$.
Applying it to first periods gives $h_xD=(a+\Omega_xc)h_x$.
Multiplication by $\theta$ is skew-adjoint for the alternating trace
pairing and complex-linear on the normalized periods. Hence it is
skew-adjoint for their Hermitian form with matrix $Y_x^{-1}$: the
imaginary part is the alternating pairing, and replacing one argument
by $i$ times that argument recovers the real part. With $A=a+\Omega_xc$,
this gives $A^tY_x^{-1}+Y_x^{-1}\bar A=0$. Consequently
\[
                  D^t C+C D=2i h_x^t c h_x .
\]
Whenever $D_i+D_j\ne0$, $C_{ij}$ is holomorphic by division.
The remaining pairs are exactly the opposite signs at the same
place. Their terms in the product exponential multiply the other
terms by exponentials of mixed row contractions. They disappear
under $P(\partial_z)$ by the mixed Laplacian, as in
Proposition~\ref{arrep:reproduce}. On each compact subset of
$\mathcal D$, $Y_x$ is bounded below by a positive matrix, so the
resulting polynomial Gaussians and their derivatives have locally
uniform Schwartz bounds. This proves holomorphy and hence
sesquiholomorphy of the displayed pairing.

On replacing $x$ by $xb$, with triangular diagonal blocks $A_0,D_0$,
the physical period map acquires left factor $A_0^t\oplus D_0^{-1}$.
Thus $\Omega_x$ is unchanged, and $h_x$ gains right factor
$A_0^{-t}\oplus D_0$. Ordinary differentiation and the row-highest
law give $Q_{xb,P}=\chi_\lambda(b)Q_{x,P}$. It suffices to compare
the two kernels on the diagonal of orthoframes and then on the full
diagonal by triangular equivariance.

For an orthoframe $x$, put $d_\tau=-\operatorname{Im}\tau(\kappa_0)$.
In its basis the physical periods of a vector $s$ are the conjugates
of its positive coefficients and its negative coefficients themselves.
The imaginary part of the Hermitian pairing with diagonal entries
$d_\tau$ in these coordinates is therefore the original alternating
form $\operatorname{Tr}_{E/\mathbf Q}(-\kappa_0\langle s,s'\rangle/2)$.
In normalized period coordinates the same form is
$\operatorname{Im}((u+\Omega_x v)^tY_x^{-1}
\overline{(u'+\Omega_x v')})=v^tu'-u^tv'$. Equality of imaginary
parts gives
$h_x^tY_x^{-1}\bar h_x=\operatorname{diag}(d_\tau 1_r)_\tau$.
Completing the real Gaussian square now gives
\[
 \langle Q_{x,z},Q_{x,\bar w}\rangle_\infty
 = |\det h_x|^n\det(2Y_x)^{-n/2}
          \exp\!\left(\pi\sum_{\tau,i,k} d_\tau z_{\tau ik}w_{\tau ik}\right).
\]
This is the orthoframe Gaussian of Proposition~\ref{arrep:reproduce}
times $c_\infty$, since $|\det x_\tau|^2=|\det G_\tau|^{-1}$.
Differentiate and use triangular equivariance. Sesquiholomorphic
kernels equal on the full diagonal are equal everywhere: the power
series of their difference at each diagonal point has coefficients
detected by independent Wirtinger derivatives on the diagonal,
and $\mathcal D$ is connected.
\end{proof}

\begin{proposition}\label{dg:identity}
Use the marked family in Lemmas~\ref{bfam:genus} and~\ref{bfam:phase},
even $n>4r$, and the column-invariant polynomial family of
Proposition~\ref{arrep:reproduce}. Let $\sum_\Lambda w_\Lambda$
denote the finite genus sum with reciprocal stabilizer weights.
Write $\phi_{L,\Lambda}$ for the marked vector on row $L$,
$\Delta$ for its prescribed rational row group, and
$w(\gamma)\phi_{L,\Lambda}=\mu(\gamma,L)\phi_{\gamma L,\Lambda}$.
For all permissible rows $L_1,L_2$ and all $x,y\in\mathcal D$,
\[
\begin{split}
 &\sum_\Lambda w_\Lambda\sum_j
    \theta_{P_j}(\phi_{L_1,\Lambda};x)
           \overline{\theta_{P_j}(\phi_{L_2,\Lambda};y)}\\
 &\quad=c_\infty\sum_{\gamma\in\Delta}\overline{\mu(\gamma,L_2)}
               K(x,\gamma y)
       \sum_\Lambda w_\Lambda
       \langle\phi_{L_1,\Lambda},\phi_{\gamma L_2,\Lambda}\rangle_f .
\end{split}
\]
The graph sum is absolutely convergent. This is the kernel identity
for this genus, including the phase of each marked row.
\end{proposition}
\begin{proof}
In the doubled oscillator use the rational diagonal polarization.
For a pure tensor $u\otimes\bar v$ the change to this polarization
takes difference and midpoint of the two position coordinates,
followed by Fourier transform in the midpoint. Evaluation at the
origin is therefore $\langle u,v\rangle$ in the self-dual position
measure, both at infinity and at the finite places. In the source
orientation the coordinate function is $W_{u,v}(-s)$ as in
Lemma~\ref{bfam:phase}. The product of the two rational theta sums
is the rational theta sum after this change: the rational pairs of
positions are in bijection with their rational differences and
midpoints, and Poisson summation in the latter variable replaces
it by its rational dual. This proof works for arbitrary Schwartz
inputs. In particular there is no primitivity restriction on the
rational differences. A common rescaling of the finite intertwiner
and the inverse rescaling of the archimedean intertwiner leave the
global evaluations unchanged.

Here it matters that the inserted archimedean test
$\sum_j Q_{x,P_j}\otimes\overline{Q_{y,P_j}}$ is invariant under the
compact column group in the doubled model. In physical period
coordinates a column unitary $u_\tau$ acts by $u_\tau$ in negative
rows and by $\bar u_\tau$ in positive rows. It preserves the
polarization, so its oscillator action fixes the Gaussian vacuum
line at each frame. One can see this by the first-order Gaussian
equations $(\partial_s-2\pi i\Omega_xs)q_{x,0}=0$: the operators
are the positive-polarization annihilators and their common solution
line is spanned by this Gaussian. Its unit multiplier is the same
at any two frames. Indeed their vacuum pairing is nonzero by the
Gaussian integral with matrix $-i(\Omega_x-\bar\Omega_y)$ of
positive real part. To check the action on the inserted vectors,
the Weyl translate of the vacuum with shift $v$ and modulation $u$
is $q_{x,z}\exp(-\pi\zeta^tY_x^{-1}\bar\zeta/2)$, where
$\zeta=u+\Omega_xv=h_xz$. This follows by replacing $s$ with $s+v$
and multiplying by $\exp(2\pi i u^t(s+v/2))$. The column action
preserves the Hermitian norm in this formula. Thus it transforms
the coherent parameters by the stated linear substitutions and the
common unit multiplier. The multipliers cancel between the two
copies, and column invariance of the sum of differential symbols
proves the assertion (the transposes in differentiation are again
paired unitary substitutions).

Apply Lemma~\ref{bfam:genus} to this invariant Schwartz test.
At the Siegel--Weil point the section for a Schwartz function $\Phi$
is $g\mapsto(\omega(g)\Phi)(0)$. Here the doubled group has split
rank $r$, the positive definite column space has dimension $n$ and
the parameter is $s_0=(n-r)/2$. The theta integral uses Haar
probability, giving coefficient one. Multiplying back by the
total mass gives the same weights on both sides.
The splitting character on the Siegel Levi can be taken trivial:
its restriction to $\mathbb A_{E^+}^\times$ is required to be the
$n$-th power of the quadratic extension character, and $n$ is even.
In this splitting the diagonal column group acts by substitution.
The same reference gives the doubled tensor intertwiner and its
value at zero; its general second-factor twist by the splitting
character is trivial here.
\cite[\S\S2.8, 3.1, 11.3 and Theorem~7.1]{ganQiuTakeda}

We specify the graph convention. Label the graph by the rational
row isometry $\gamma$ acting on the second member of the tensor to
give the diagonal functional. Thus for the seed graph one evaluates
at zero, and for this graph one first translates that member by
$\gamma$. In row coordinates a representative is the conjugate of
$\operatorname{diag}(1,\gamma^{-t})$ by the rational change of
polarization with blocks
$\left(\begin{smallmatrix}1&\kappa_0G/4\\-1&\kappa_0G/4\end{smallmatrix}\right)$.
Equivalently the action on the second undoubled theta test is the
rational oscillator action carrying its frame $y$ to $\gamma y$.
To check the scalar, let
$\left(\begin{smallmatrix}a&b\\c&d\end{smallmatrix}\right)$ be
the induced Siegel transformation in the symplectic period basis,
$U=c\Omega_y+d$, $\Omega'=(a\Omega_y+b)U^{-1}$.
The equality of the period maps on $\gamma s$ at $\gamma y$ and on
$s$ at $y$ gives $h_{\gamma y}=U^{-t}h_y$,
$\Omega_{\gamma y}=\Omega'$. The Gaussian action sends
$Q_{y,z}$ to $Q_{\gamma y,z}$ up to a constant phase independent
of $y,z$. Indeed under a real linear coordinate change this follows
by substitution, and under adding a real symmetric matrix to the
period it follows by multiplying by the corresponding quadratic
phase. Under Fourier inversion use the Gaussian integral
\[
 \int_{\mathbb R^g}e^{\pi i s^t\Omega s+2\pi i s^t t}\,ds
       =\det(-i\Omega)^{-1/2}e^{-\pi i t^t\Omega^{-1}t}.
\]
For the Fourier generator $\Omega'=-\Omega^{-1}$, the factor
$e^{-\pi i\zeta^t\Omega^{-1}\zeta}$ combines with the original
quadratic correction to give the one at
$\zeta'=\Omega^{-t}\zeta$, since
$\zeta'^tY'^{-1}\zeta'=\zeta^tY^{-1}\zeta-
2i\zeta^t\Omega^{-1}\zeta$. The determinant in $Q$ absorbs the
period-dependent determinant after taking $n$ rows. These generators
generate the rational symplectic group by symplectic elimination.
Choose the reciprocal scalar on the finite operator. Its action
$w(\gamma)$ is then exactly the one characterized by
$\theta(w(\gamma)\phi;\gamma y,z)=\theta(\phi;y,z)$: for the
Fourier generator this is rational Poisson summation, for rational
linear changes it is reindexing with the product formula, and for
the shear it is cancellation of the rational quadratic phases.
This characterization determines the finite operator, since at
fixed $y$ the distinct rational position frequencies of the theta
series can be separated by Fourier coefficients in $h_yz$ after
removing the common nonzero Gaussian prefactor.

These are also the normalizations of the graph section. Indeed the
diagonal-polarization change intertwines the rational row actions
with their conjugates in the doubled group. For the linear changes,
quadratic phases and Fourier generators this is just covariance
of Weyl translations under these operations. Any scalar ambiguity
between the two intertwiners is independent of the test by
irreducibility of the Schrödinger representation; its global value
is one since both rational actions preserve the nonzero rational
theta distribution by Poisson summation. The choice of another
rational Siegel-parabolic representative likewise preserves the
global evaluation, by the product formula on its determinant and
triviality of its unipotent character at zero. The term on the
graph consequently is precisely
\[
 \langle\phi_{L_1,\Lambda},w(\gamma)\phi_{L_2,\Lambda}\rangle_f\,
          \sum_j\langle Q_{x,P_j},Q_{\gamma y,P_j}\rangle_\infty .
\]
Graphs outside $\Delta$ vanish by the marked support condition.
For $\gamma\in\Delta$ the inner product in the second variable
contributes $\overline{\mu(\gamma,L_2)}$.

For a nongraph there is a nonzero rational isotropic row line
$E\xi$ in the intersection of its maximal isotropic with the first
copy. Indeed either projection being singular forces the other to
be singular by orthogonality and dimension. For every $v\in E^n$
the finite Weyl translation on the first input by $v\otimes\xi$ acts,
after passage to this polarization, by a pure modulation. It takes
value one at the origin: in difference and dual coordinates the
action is $\phi(t-d)\psi_f(\langle b,t-d/2\rangle)$ with $d=0$.
There is no residual scalar, since linear symplectic covariance
preserves the symmetric Weyl operators including their half
pairings. The resulting finite evaluation $\ell$ thus satisfies
$\ell(\rho(v\otimes\xi)\phi_1,\phi_2)=\ell(\phi_1,\phi_2)$.
The marked isotropic relation has finite coefficients $a_v$ with
$\sum a_v\ne0$ and $\sum a_v\rho(v\otimes\xi)\phi_1=0$.
It gives $(\sum a_v)\ell(\phi_1,\phi_2)=0$ before any
archimedean evaluation. This accounts for every nongraph.

Finally the polynomial differentiations are legitimate in this
expansion. On compact sets of complex $z,w$ the archimedean
Gaussians, with every fixed parameter derivative, are bounded in
Schwartz space, by completion of the square and $Y_x,Y_y>0$.
The same holds after the fixed partial Fourier transform.
For a Siegel representative with Iwasawa decomposition $pk$ and
Levi block $a_p$, evaluation is bounded by
$|\det a_p|_{\mathbb A_E}^{n/2}
\sup_k|(\omega(k)\Phi)(0)|$; the supremum is uniform on these
bounded test families by continuity and compactness. The finite
tests lie in finitely many compact orbits under the maximal compact
subgroup (each smooth test has open stabilizer there).
The scalar Siegel Eisenstein majorant converges absolutely for
$\operatorname{Re}s>r/2$, which contains $s_0$ since $n>4r$.
Thus either Cauchy's formula on a slightly larger parameter
polydisc or direct Schwartz differentiation permits the fixed
finite insertions under the sum. The genus formula above was only
needed for the invariant sum of inserted vectors. Gaussian theta
series on the left also converge normally since the finite test
restricts rational positions to a fixed fractional lattice.
Combining the surviving graph evaluations with Lemma~\ref{dg:gaussian}
proves the identity.
\end{proof}

\src{src/Development/AbelianVarietyTheory/Proofs/SiegelThetaTransformation/EisensteinSeriesTermUnitaryGraph.lean}{33}{The graph representative and both translated tests};
\src{src/Development/AbelianVarietyTheory/Proofs/SiegelThetaTransformation/KernelIdentity.lean}{181}{The class-weighted differentiated graph identity}.

\subsection{Marked period components and absolute unfolding}

Fix a rational Hermitian space \(W\) over \(E/F\), \(F=E^+\), its
polarizing trace form and CM type. Take the identity auxiliary cover
and its principal odd level \(M\), with
\((M)=\mathfrak m\bar{\mathfrak m}\) and
\(\mathfrak m+\bar{\mathfrak m}=O_E\).
In the period description a component has a self-dual row lattice
\(L\subset W\), the polarized half-level mark and the connected flag
domain of its prescribed signature. A marking on one half of \(L/ML\)
can equally be specified by its dual basis on the opposite half,
using the fixed primitive \(M\)-th root for the Weil pairing.

\begin{lemma}\label{munf:marks}
Fix a reference row \(L_*\) with a basis of \(L_*/\mathfrak mL_*\).
Every such marked period lattice in \(W\) can be carried by a rational
isometry to a row equal to \(L_*\) at the primes of \(M\), with the
specified basis on the \(\mathfrak m\)-half. Write \(\Delta\) for the
rational unitary elements integral there and acting identically on this
marked half of \(L_*\). Two transported rows lie in the same
\(\Delta\)-orbit exactly when their polarized marked period lattices
are isometric. The stabilizer of each row is its principal-level
arithmetic group.
\end{lemma}
\begin{proof}
Every place \(v\) of \(F\) above the level splits in \(E\), since a
prime fixed by conjugation could not divide the product of the
coprime conjugate halves. The polarization pairs the two summands of
\(L/ML\) perfectly: the CRT idempotents are interchanged by the
adjoint involution, so each summand is isotropic and is the dual of the
other. This proves the equivalence of the two ways of specifying
the marking, including at a prime power. Locally \(U(W)(F_v)\) is
\(\mathrm{GL}_r(F_v)\) on the first factor, acting by the inverse
adjoint on the second. Lift the indicated residue bases in both
first-factor lattices to bases over the DVR. The matrix carrying one
basis to the other carries the lattices onto one another. Its inverse
adjoint carries the second lattices and the dual marks, by perfection.

The resulting conditions on the tuple of local isometries form a
nonempty open set. Rational isometries meet it. Indeed the Cayley
map \(X\mapsto(1-X)(1+X)^{-1}\) identifies an open of the
\(F\)-vector space of skew-adjoint endomorphisms with the unitary
open where \(1+g\) is invertible. Every selected local open meets
this chart: multiply its matrix by a norm-one scalar arbitrarily
close to one, avoiding the finitely many prohibited eigenvalues.
At these split places such scalars are \((t,t^{-1})\).
Weak approximation on the vector space, in a fixed \(F\)-basis,
now gives a single rational isometry. The transported lattices agree
with \(L_*\) at the whole level. Globally, this is equivalent to the
existence of a common denominator prime to the level in the two
lattice inclusions: localize the finite torsion modules
\((L+L_*)/L\) and \((L+L_*)/L_*\), then choose an annihilator
congruent to one at those primes by CRT.

An element of \(\Delta\) carrying one transported lattice to another
preserves the mark on the first half, hence on the dual half.
Conversely a marked isometry between the lattices extends uniquely
by \(E\)-linearity to \(W\); its local action belongs to \(\Delta\).
When the two rows agree this gives precisely the integral isometries
acting identically modulo \(M\). To see the component statement in
period coordinates, a complex framed polarized torus has a period
lattice on its tangent space. An isomorphism lifts uniquely to its
complex-linear derivative on the universal covers, so it induces
the preceding integral isometry. Conversely transporting the
complex structure and flag by an isometry gives isomorphic marked
complex tori. For a fixed marked lattice the positive-plane domains
are matrix balls with connected compact flag fibres; their product
is connected. In a family of complex abelian varieties the integral
homology, its trace pairing and the finite marking are locally
constant in a topological trivialization. This gives exactly the
isometry classes as component indices in the marked period model.
\end{proof}
\src{src/Development/AbelianVarietyTheory/Proofs/SiegelThetaTransformation/ExistsMoveMarked.lean}{133}{Rational transport of the marked basis},
\src{src/Development/AbelianVarietyTheory/Proofs/SiegelThetaTransformation/ComponentEqOfMatchedLatticesMarkings.lean}{74}{identification of marked indices} and
\src{src/Development/AbelianVarietyTheory/Proofs/SiegelThetaTransformation/ExistsComponentsLatticesMovedToLevelFamily.lean}{162}{distinct orbits after simultaneous transport}.

\begin{proposition}[Unfolding the finite span]\label{munf:span}
On the countable set of admissible rows let \(C\) be the equivariant
matrix of Lemma~\ref{bfam:gram} and \(K\) the kernel of
Proposition~\ref{arrep:reproduce}. Suppose the marked periodized
kernel is the finite sum of paired theta series for the indicated
column genus, with a single nonzero normalization. Then every
bounded family of holomorphic forms of the prescribed weight and
unitary multiplier on these row orbits lies in the span of those
theta series. The conclusion also holds simultaneously on any finite
collection of distinct row orbits.
\end{proposition}
\begin{proof}
Write the input as \(F\) and put \(V=C^{-1}F\). Use the norm given
by the supremum over all rows and orthoframes. The Neumann series
converges in this norm and preserves the weight and multiplier,
since \(C\) commutes with the action. It converges locally uniformly
on the full domain by triangular equivariance and Gram--Schmidt,
so \(V\) is holomorphic there. Its absolute value is bounded on all
rows of the slab by one constant \(B\), by the compactness of the
triangular factors in the proof of Proposition~\ref{arrep:reproduce}.
For a fixed output \(a=(L,x)\) put
\[
          f_a(L',y)=C_{LL'}K(x,y)V_{L'}(y).
\]
On rows times the slab we have
\[
 \sum_{L'}\int_{S_\delta}|f_a(L',y)|\,dy
       \le B\sum_{L'}|C_{LL'}|\int_{S_\delta}|K(x,y)|\,dy<\infty .
\]
Thus we may unfold under the countable diagonal action on this
product with its counting and Lebesgue measures. The action
preserves both measures.

Here is an explicit reason a measurable fundamental set exists.
Choose one \(L'\) in each row orbit. Its stabilizer is discrete in
the real unitary group: on an integral basis of the underlying
rank-\(2fr\) lattice it has integral matrix entries. It acts freely
on the invertible frames and properly, since the isometries carrying
one compact set of frames into another form a subset with compact
closure (write \(g=y'y^{-1}\)). Cover the slab by countably many
open patches disjoint from their nonidentity stabilizer translates.
Keep in each patch the complement of the translates of all earlier
patches. Their union meets every stabilizer orbit once and is Borel.
The union on the representative rows is a diagonal fundamental set
\(\mathcal D_0\). Absolute integrability and Tonelli now give
\[
 \sum_{L'}\int_{S_\delta}f_a(L',y)\,dy
       =\int_{\mathcal D_0}\sum_{\gamma\in\Delta}
                           f_a(\gamma L',\gamma y)\,d(L',y).
\]
The left side is \(\rho_\delta(CV)_L(x)=\rho_\delta F_L(x)\).
On the right the multiplier law for \(V\) and the given periodized
identity express the inner sum as \(\sum_{j=1}^s q_j(a)h_j(L',y)\),
where the \(q_j\) form a basis selected from the finitely many theta
series. The functions \(h_j\) are integrable even if a single
uncombined term in the original genus expression has no known
integral. Indeed choose \(s\) evaluations \(a_i\) for which
\((q_j(a_i))\) is invertible: all evaluation vectors span
\(\mathbf C^s\) by independence of the functions. The sums at these
\(s\) points are integrable by the preceding absolute estimate, and
matrix inversion expresses each \(h_j\) as their linear combination.
The zero-dimensional span gives zero directly. Divide by the
positive \(\rho_\delta\) to obtain the asserted finite expansion.
For a collection of orbits extend the families by the equivariant
law from each representative and by zero on the complement. The
stabilizer law makes the extension well-defined and the unitary
multiplier preserves its bound.
\end{proof}
\src{src/Development/AbelianVarietyTheory/Proofs/SiegelThetaTransformation/ExistsFundamentalDomain.lean}{138}{A fundamental set on lattice--frame pairs} and
\src{src/Development/AbelianVarietyTheory/Proofs/SiegelThetaTransformation/HermitianThetaSeriesSpanBoundedForms.lean}{22}{finite rank after absolute unfolding}.

For the rows of Lemma~\ref{munf:marks} in their common rational
Hermitian space, use the graph sum of Proposition~\ref{dg:identity}.
The passage to algebraic cusp sections then uses \eqref{bfam:reading}
on that space with its marked multiplier.


\paragraph{The fourth power of the theta key.}
Write \(g=f(p+q)\). On the abelian locus let \(L\) be the rigidified
symmetric principal theta line, \(P=\pi_*L\), and
\(\omega=\det(\operatorname{Lie} A)^\vee\), where the dual applies to
the determinant. The line \(P\) commutes with base change. Indeed its
fibres have one section and vanishing higher cohomology. A local
generator \(s\) cuts out a relative effective theta divisor
\(\Theta\). The canonical evaluation identifies
\(\mathcal O(\Theta)=L\otimes\pi^*P^{-1}\), hence
\(e^*\mathcal O(\Theta)=P^{-1}\). In particular this identification
also works when the origin lies on \(\Theta\).

\begin{proposition}\label{ctwist:key}
There is a functorial unit section on the abelian locus, up to a
sign on each parity, of \(P^{\otimes8}\otimes\omega^{\otimes4}\).
Its scalar \(k(s,v)\) in a theta frame \(s\) and a tangent basis
\(v\) transforms by
\[
 k(a s,Mv)=a^{-8}(\det M)^4 k(s,v).
\]
In the metric on the cover its absolute norm is \((2\pi)^{2g}\).
In a primitive theta frame and a full integral tangent basis it is
a unit at good reduction, including residue characteristic two.
\end{proposition}
\begin{proof}
Apply Moret--Bailly, Theorems 0.2 and 0.5 and Corollary 1.5
\cite{moret-bailly-key}, to the effective principally polarized
symmetric divisor. His line is
\(\omega\otimes(e^*\mathcal O(\Theta))^{-2}\).
The scalar is taken relative to \(s^8\otimes\eta_v^4\), where
\(\eta_v(v_1,\ldots,v_g)=1\). This gives both transformation
exponents and the assertion over an abelian scheme over a DVR.
The multiplier-line description of Candelori, Theorem 5.1 and
Remarks 5.2--5.3 \cite{candelori-key}, is over
\(\mathbb Z[1/2,i]\).

Here is the metric calculation. On the even characteristic chart
take periods \((1,\Omega)\), \(Y=\operatorname{Im}\Omega>0\),
the usual Riemann theta, and
\(v_i=(2\pi i)^{-1}\partial/\partial z_i\). The analytic
trivialization of Theorem 0.5 gives \(|k|=1\) in these frames.
Write \(N_s\) for the Haar mean of the squared rigidified theta
norm. With \(z=x+\Omega y\), integration in \(x\) annihilates the
cross terms, and summation of the diagonal terms gives
\(\int_{\mathbb R^g}\exp(-2\pi y^tYy)\,dy=(\det(2Y))^{-1/2}\).
Put \(N_v=\operatorname{covol}(\Lambda)/
(\pi^g|\det(v)|^2)\) in the same complex coordinates. Then
\[
 |k|N_s^4 N_v^2=(2\pi)^{2g}.
\]
The two exponents just computed give invariance under changes of
theta and tangent frames. A symplectic basis puts every even
characteristic in the chosen chart. For an odd characteristic take
the product of the pair with itself. This has even parity, the
two metrics multiply, and Corollary 1.5 identifies the product
trivializations up to sign. Positive square roots give the claimed
identity in odd parity as well. On an orthonormal cover frame the
Hermitian matrix of the polarization is \(2\xi I\) and the theta
has \(N_s=1\). Unimodularity of the alternating form gives
\((2\xi)^g\operatorname{covol}(\Lambda)=1\), so
\(N_v=(2\pi\xi)^{-g}\). The cover's Lie factor for the four
determinants is \((2\pi\xi)^{-2g}\), giving the stated norm.
\end{proof}

There is a further integral change of frame even on the abelian
locus. Let \(R\) be a test DVR whose fraction field splits \(E\),
\(T\) its free Lie lattice, and \(T_\tau=T\cap(T_K)_\tau\) of
rank \(m_\tau\). If \(\mathfrak D_E\) is the different, then
\[
 0\le {\rm length}_R(T/\textstyle\bigoplus_\tau T_\tau)
 \le\sum_\tau m_\tau {\rm ord}_R(\tau\mathfrak D_E).
\]
Indeed for the projector \(e_\tau\), the trace-dual basis formula
\(e_\tau=\sum_i b_i\otimes\tau(b_i^*)\) gives
\(\tau(\mathfrak D_E)e_\tau T\subset T_\tau\):
the \(b_i\) are an integral basis and \(b_i^*\in\mathfrak D_E^{-1}\).
Projection embeds the displayed quotient in
\(\bigoplus_\tau e_\tau T/T_\tau\). Each summand is generated by
at most \(m_\tau\) elements and is killed by
\(\tau(\mathfrak D_E)R\). In the combined eigentangent basis
the valuation of the key increases by four times this length.
All normalized valuations follow by dividing by the ramification
index of the test DVR. Consequently \(k^m\) is integral at these
tests for every \(m\ge0\), with logarithmic complex cost
\(2gm\log(2\pi)\).
\src{src/Development/AbelianVarietyTheory/Proofs/DivisorProofs/ThetaKeySectionNormComplexPointsEqTwoPiPow.lean}{124}{The normalized cover metric};
\src{src/Development/AbelianVarietyTheory/Proofs/DivisorProofs/ExistsKeyUnitPowTwistingSectionOnComponent.lean}{135}{The power of the key and the separate root section}.

\begin{theorem}[Integral key at the semiabelian tests]\label{ctwist:semi}
For the rigidified graph line of Lemma~\ref{intjet:maximum},
\(k(s,v)\) has valuation zero at every semiabelian DVR test in
a primitive theta frame and a full integral tangent basis. The
different bound above therefore applies to the combined eigentangent
basis at all finite places. On each component \(k^m\) supplies the
weight-\(-4m\) multiplier with cost \(2gm\log(2\pi)\).
\end{theorem}
\begin{proof}
It suffices to compute an order after completion and finite
extension. Use the connected model, and split the semistable
reduction after extension. Write \(A/K\) for the generic fibre.
In an additional finite extension make \(A[8]\) rational. Let
\(\mathcal A'\) be the open subgroup of the Néron model generated
by its identity component and the sections of \(A[2]\).
The identity component is the given connected semistable model
(SGA~7, IX, Proposition 3.2, also used in
\cite[IV, \S3.4]{moret-bailly-pinceaux}); this identification
also applies to its semistable base change. The component group
of \(\mathcal A'\) is killed by two, and
\(\mathcal A'[2]\) is finite flat. For the latter assertion the
kernel of multiplication by two is quasi-finite flat on a
semistable group: on each geometric fibre multiplication factors
on the connected semiabelian variety through isogenies of its
torus and abelian quotient, and flatness follows from the
fibrewise flatness criterion. Its generic fibre is thus
schematically dense. Its support is contained in the union of
the finitely many sections extending the generic two-torsion,
which is closed and finite over the base. The kernel is proper
by the valuative criterion, hence is finite. This argument allows
nonreduced special two-torsion.

Moret--Bailly, II, Corollary 1.2.8 with \(n=2\), gives a symmetric
cubical extension \(L'\) on \(\mathcal A'\). Its restriction to
the identity component is the rigidified cubical line of the
test by uniqueness. It is relatively ample by VI, Proposition
2.1, and \(K((L')^2)=\mathcal A'[2]\), since both are the flat
closures of the generic principal two-torsion. The special
fibre of the connected model is geometrically integral. A
primitive theta section has order zero there: if it vanished
along the fibre, division by a uniformizer would give another
integral section. Its nonzero reduction persists under residue
field extension. Lemma~\ref{intjet:maximum} gives this section
canonical norm at most one on the entire analytification.
Consequently its order along every additional component \(C\)
of \(\mathcal A'\) is nonnegative, by the cubical model formula
\({\rm ord}_C(s)=-\log\|s(x_C)\|\) at its Shilov point
(\cite[Proposition 10.3]{deJongShokrieh}). Smoothness gives
normality, so the vertical inequalities and generic regularity
extend the section over \(\mathcal A'\). Restriction back to the
identity component proves that the two theta lattices agree.
This reasoning also compares a test with additional components
by restriction to the identity component. The Lie lattices agree
under open immersions.

Now apply Moret--Bailly, VIII, 1.0--1.3(i)
\cite{moret-bailly-pinceaux}. The base is a normal trait, the
abelian locus is schematically dense, \(K((L')^2)\) is finite,
and the rank \(d\) of the theta module is one. The asserted
canonical trivialization of \(\Delta^{a(g)}\),
\(\Delta=(\pi_*L')^2\otimes\omega\), therefore applies over
\(\mathbb Z[1/d]=\mathbb Z\). Denote it by \(u\). On the integral
stack of abelian schemes with their effective symmetric principal
theta divisors, \(k^{a(g)}/u^4\) is a global unit. Its square is
one by Moret--Bailly, Proposition 1.4(iii) and Corollary 1.5
\cite{moret-bailly-key}. Here one passes between rigidified lines
and effective divisors by the evaluation construction preceding
Proposition~\ref{ctwist:key}. Compatibility of \(u\) with
restriction to the generic point shows
\(a(g){\rm ord}(k)=4{\rm ord}(u)=0\) in the two agreed lattices.
Neither extracting an integral root nor bounding \(a(g)\) is
required. The eigentangent change and Proposition~\ref{ctwist:key}
now give both estimates. Taking the power on the component uses
the inverse only on the characteristic-zero space; extension by
its component idempotent gives the supported multiplier.
\end{proof}

The distinction between nonnegative order and order zero can be
checked on the Tate elliptic curve with parameter \(Q\) and odd
divisor \([0]\). In its connected model
\(e^*\mathcal O([0])=\omega^{-1}\), hence
\(\Delta^4=\omega^{12}\). The modular discriminant is a
nowhere-zero section of this line on the odd abelian stack.
It agrees up to sign with the key by the unit calculation above.
In the invariant Tate frame its scalar is
\(Q\prod_{n\ge1}(1-Q^n)^{24}\), of order \(v(Q)\).
More generally, on a connected semistable model without the
graph condition one always has \(P'\subseteq P\), with \(P'\)
the enlarged model's theta lattice. If \(P'=aP\), the same
canonical-power argument on \(\mathcal A'\) gives
\({\rm ord}(k(s,v))=8{\rm ord}(a)\ge0\). The odd Tate example
has \({\rm ord}(a)=v(Q)/8\) after the extension. Its toric
characteristic is \(1/2\): in the half-characteristic theta
series the toric exponents are shifted by \(1/2\).
In the graph case Lemma~\ref{intjet:maximum} removes this
index and gives the upper different bound as well.

On a normal toroidal chart with these primitive theta and
invariant differential lattices the same argument applies to
the trait at each generic boundary divisor in characteristic
zero. The key and its inverse are units there in full Lie
frames. Thus its multiplication preserves vanishing at those
boundary divisors. The scalar of this fourth-power section and
the existence of the auxiliary one-key section are distinct.
The latter asks for a fourth root up to a scalar in \(F^\times\)
on each component. Over a characteristic-zero component with
trivial geometric monodromy on the symplectic four-torsion this
follows also from Candelori, Definition 3.2 and Theorem 5.1
\cite{candelori-key}. The multiplier torsor is obtained from
the isometry torsor of the symplectic four-group with its theta
characteristic by pushing forward along a character into \(\mu_4\).
In a geometric marking both this torsor and the torsor of
fourth roots of the key are constant: their two fourth-power
trivializations differ by a constant on the universal complex
stack, by the unit calculation in \cite{moret-bailly-key} in
genus \(g\ge2\). Descent of a geometrically constant \(\mu_4\)-torsor
on a geometrically connected component amounts to a torsor over
\(F\), i.e. to a class in \(F^\times/(F^\times)^4\).
The deep level in the statement has odd \(M\ge3\) and a separate
auxiliary marking \(A\). The odd level alone gives no full
four-torsion marking. For the original even weight spaces the
unit \(k^m\) above is intrinsically defined.
\src{src/Development/AbelianVarietyTheory/Proofs/DivisorProofs/ExistsKeySectionFourthRootPerComponent.lean}{93}{The separate descent up to a component scalar}.

\paragraph{Determinants on a characteristic-zero component.}
Let \(\Lambda\) be the projective \(\mathcal O_E\)-lattice of a
component, with level group contained in its unitary isometries
acting as the identity on \(\Lambda/M\Lambda\), \(M\ge3\).
Every such isometry has determinant one over \(E\). In fact its
determinant \(u\) preserves the invertible ideal
\(\bigwedge_{\mathcal O_E}^{p+q}\Lambda\), is \(1\bmod M\)
(localize to check this on free bases), and satisfies
\(|\tau u|=1\) at every embedding. If \(u\ne1\),
\(M^{[E:\mathbb Q]}\le |N_{E/\mathbb Q}(u-1)|
\le2^{[E:\mathbb Q]}\), an impossibility.
For the split half-marking, fixing \(\Lambda/\mathfrak m\Lambda\)
first gives the congruence modulo \(\mathfrak m\). Taking the
adjoint with respect to the perfect alternating integral pairing
gives the conjugate-ideal congruence for the inverse, and hence
for the isometry itself. Indeed the dual of
\(\mathfrak m\Lambda\) is
\(\overline{\mathfrak m}^{-1}\Lambda\). The two ideals are
coprime with product \(M\mathcal O_E\), giving the asserted
principal congruence.

On such a classical complex component the determinant of the
\(\tau\)-part of the standard Betti local system thus has a global
horizontal generator, including for noncompact signatures. On the
period domain the polarization gives the Hodge exact sequence
\[
 0\longrightarrow\omega_\tau\longrightarrow
 H^1_{\rm dR,\tau}\longrightarrow\omega_{\bar\tau}^{\vee}
 \longrightarrow0.
\]
Dualizing its determinant gives
\(\det\operatorname{Lie}_\tau/\det\operatorname{Lie}_{\bar\tau}\),
the line of the elementary section \(s_\tau\).
The norm of a horizontal volume is constant: in a unitary
orthonormal basis the absolute determinant of the Hermitian form
on the full space is one, independently of its positive and
negative Hodge subspaces. Changing to the cover normalization
multiplies this by a constant fixed on the component.
For a smooth geometrically connected algebraic component over
\(F\) carrying the indicated abelian scheme and splitting of \(E\),
the Gauss--Manin determinant is regular singular. Full faithfulness
for regular singular connections gives an algebraic horizontal
generator over \(\mathbb C\) (Deligne, II, Theorem 5.9
\cite{deligne-reg-sing}). It descends up to a scalar to \(F\):
horizontal sections are the kernel of the connection on global
sections, and both global-section spaces and the kernel commute
with extension of scalars on this quasi-compact separated
component. The descended one-dimensional space generates the
line everywhere by faithful flatness. This proves rational
triviality with its archimedean constant for these components.

There are two further consequences for a fixed classical component.
Take a smooth characteristic-zero toroidal compactification with
unipotent local monodromy and with the Hodge bundles given by the
canonical extension of the Gauss--Manin bundle. The horizontal
determinant and its inverse extend as units in the ratio of
the two extended Hodge determinants. Indeed, near a normal
crossings boundary with parameters \(z_i\), the canonical
extension in a multivalued horizontal basis is obtained by
multiplying by
\(\exp(-\sum_i(\log z_i)\log T_i/(2\pi i))\), where \(T_i\)
are the local monodromies. Each \(\log T_i\) is nilpotent on
the \(\tau\)-part, so this matrix has determinant one. Taking
determinants of the extended Hodge sequence identifies its
volume line with the ratio in question. Thus multiplication
in either direction preserves the subcanonical cusp ideal
\(\mathcal O(-D_\infty)\) on these charts. The extended Hodge
sequence and the boundary ideal in the classical PEL model are
given in \href{https://www-users.math.umn.edu/~kwlan/articles/cpt-ram-ord.pdf}
{Lan, Proposition 1.4.2.1 and Definition 1.4.2.9}.

On any fixed finite-type model of this compactification and line
over an open of \(\operatorname{Spec}\mathcal O_F\), the same
section is an integral generator after deleting finitely many more
places. To see this, choose a finite affine trivializing cover.
The coordinates of \(s\) and \(s^{-1}\) on the generic fibre and
the finitely many gluing identities all descend after inverting a
single nonzero base integer. On this open the order is zero at
every test whose lattice is the pullback of that model line.
Away from \(\Delta_E\), the embedding idempotents of the split
étale algebra \(\mathcal O_E\otimes\mathcal O_{F,w}\) are integral,
so formation of its Lie eigensummands commutes with pullback.
In particular this gives the order for the abelian tests extending
to the fixed PEL chart with its universal abelian scheme. The
exceptional primes produced by this argument depend on the
component, horizontal generator and model. Controlling all the
semiabelian tests and the total arithmetic degree in a varying
family requires the comparisons below.

\paragraph{One vertical order on one component.}
Fix a component \(Z\), a host embedding \(\tau\), and a finite place
\(w\) of its coefficient field \(F\). Orders on test traits are
normalized to restrict to \({\rm ord}_w\) on \(F\). The following
divisor calculation explains what an integral model of the parallel
line has to establish. The special fibre may have several irreducible
components with unequal multiplicities.

\begin{lemma}[Vertical order on a proper model]\label{ctwist:vertical}
Let \(\mathcal X\) be integral, normal, proper and flat over a DVR
\(R\) with uniformizer \(\pi\) and nonempty connected special fibre.
Let \(\mathcal L\) be an invertible sheaf, trivialized by \(s\) on the
entire generic fibre. Suppose some positive power of \(\mathcal L\)
is trivial on the reduced special fibre. Write
\({\rm div}(\pi)=\sum_i m_i C_i\) and
\({\rm div}(s)=\sum_i a_i C_i\). Then all \(a_i/m_i\) are the same
rational number \(\alpha\). For every local test \(R\to V\) with
\(\operatorname{Spec}V\to\mathcal X\), where \(V\) is a valuation
ring extending \(R\), the coordinate of \(s\) in an integral generator
has order \(\alpha\,v(\pi)\). If \(R=\mathcal O_{F,w}\) and one
such test is a DVR of ramification index one, then \(\alpha\in\mathbb Z\).
\end{lemma}
\begin{proof}
Set \(\alpha=\min_i a_i/m_i\). Choose \(M>0\) divisible by its
denominator and by the order of the restricted line in the Picard
group. The rational section \(u=\pi^{-M\alpha}s^M\) has nonnegative
order along every prime divisor and order zero at a minimizing
component. It extends over \(\mathcal X\) by normality, checked in
local frames by the height-one intersection criterion.
In a trivialization of \(\mathcal L^M\) on the reduced special
fibre, its restriction is a global function nonzero at that
component's generic point. Global functions on a proper reduced
connected scheme over a field form a field: their algebra is
finite-dimensional and reduced, and connectedness excludes
nontrivial idempotents. Hence \(u\) is a unit on the reduced fibre,
and therefore generates at every point of the special fibre by
Nakayama. It also generates on the generic fibre. Pulling
\(s^M=\pi^{M\alpha}u\) to \(V\) proves the assertion. At an
unramified DVR, \(v(\pi)=1\) and the coordinate has integral order.
\end{proof}
The two scheme facts used here are
\href{https://stacks.math.columbia.edu/tag/0AVB}{Stacks Project,
Tag 0AVB} and
\href{https://stacks.math.columbia.edu/tag/0366}{Section 33.9}.
If the special fibre has just one irreducible component, the same
proof dispenses with its Picard group: after clearing the single
ratio both \(u\) and \(u^{-1}\) extend by normality. In the case
of a reduced irreducible fibre \(M=1\) suffices. For a model over
a ramified base \(R/\mathcal O_{F,w}\), divide its orders by the
ramification index before comparing tests. The rational choice of
frame alone cannot determine \(\alpha\): for the lattice
\(\ell^a\mathbb Z_\ell\) in \(\mathbb Q_\ell\), the frame \(1\)
has integral coordinate \(\ell^{-a}\).

Here is the specific model comparison needed to apply this reasoning.
On the good deep cover write \(U\) for the framed uniformization.
For each fixed \((Z,\tau,w)\), applying the vertical divisor lemma
requires comparing the test frame with local charts on a normal
projective flat model after a finite level and base change. The
comparison must lift
\emph{every} admissible test, match its frame coordinate up to a
unit, and use the same normalized order on the base uniformizer.
Triviality on the reduced fibre, or an equivalent constancy assertion
in these compatible charts, is essential on reducible fibres.
On the dependency path of the exported theorem the generator on
the component is obtained using resolution and GAGA. The source
then uses the associated period frame to compute the degree.
The finite correction between the two frame conventions at a fixed
\((\tau,w)\) is \(q\) times the logarithmic difference between the
different and \(d_0\) in the prescribed local absolute value. Adding
this fixed correction preserves equality of the readings on that
component. The final change of the global degree follows by summing
the finite corrections. Lemma~\ref{ctwist:vertical} proves the
vertical divisor step from its stated model hypotheses. The splitting
model below computes the ordinary integral lattices directly.
\src{src/Development/AutomorphicSpectrum/Proofs/GeometryProofs/FlatFirstFromResolutionAndGagaStatement.lean}{8}{Generator from resolution and GAGA};
\src{src/Development/AutomorphicSpectrum/Proofs/GeometryProofs/KroneckerCharacterFrameBound.lean}{23}{Passage to the component period frame};
\src{src/Development/AutomorphicSpectrum/Proofs/GeometryProofs/KroneckerCharacterPeriodFormula.lean}{122}{Finite correction in the ordinary metric};
\src{src/Development/AutomorphicSpectrum/Proofs/DivisibilityProofs/FlatDegreeStOfResKGaga.lean}{23}{Passage from the constant degree to the corrected section}.

\paragraph{Computing the common degree at a CM point.}
Write \(A_\chi=\operatorname{Re}\log\Omega(E,\chi)\). Assume the
local equality just discussed, and the constant archimedean norm
in the prescribed metric. One point with everywhere good abelian
reduction now suffices to read the uncorrected degree on \(Z\).
In the original construction one chooses an orthogonal \(E\)-frame
in the \emph{specified} complex uniformizing component, constructs
its integral CM object over some \(K/F\), and identifies its framed
generic fibre with a point of \(Z\). This choice works for every
\(\tau\) of the chosen CM orientation on that \(Z\). Write \(T_+\)
and \(T_-\) for the integral Lie lattice intersected with the
\(\tau\)- and \(\bar\tau\)-eigenspaces there, with the induced
polarization metrics. Their ranks are \(p,q\). The frame law
identifies the uncorrected line degree of \(s_\tau\) with
\(\widehat{\deg} T_+-\widehat{\deg} T_-\), using normalized
determinant degrees. Indeed integral local bases compute its finite
coordinate by a ratio of exterior determinants. The same ratio in
orthonormal frames computes its norm at infinity. Changing rational
bases adds principal divisors of degree zero. Passing to \(K\)
preserves the answer because
\(\sum_{v\mid w}e(v/w)f(v/w)=[K:F]\), and each complex embedding
has \([K:F]\) extensions. This calculation uses the integral
eigensublattices even at primes dividing the different.
\src{src/Development/AbelianVarietyTheory/Proofs/DivisorProofs/ExistsCmGoodReductionPointOnComponent.lean}{57}{The CM point on the prescribed component};
\src{src/Development/AbelianVarietyTheory/Proofs/DivisorProofs/FlatLineDegreeEqEigenLatticeDegreeSub.lean}{20}{Finite and infinite determinant readings}.

The cancellation of the component's indices can be seen explicitly.
Let \(S_\pm\subseteq T_\pm\) be the full-rank images spanned by the
differentials of the \(p\) and \(q\) CM factors of the two types,
\(I_\pm\) their determinant index ideals in \(\mathcal O_K\), and
\(g_\pm\) their respective polarization Gram matrices over
\(\mathcal O_E\). Put \(G_\pm=|N_{E/\mathbb Q}\det g_\pm|\).
The CM period comparison for the images gives
\[
 \widehat{\deg}S_+=pA_\chi/2-\frac{\log G_+}{2[E:\mathbb Q]},
 \qquad
 \widehat{\deg}S_-=qA_\chi/2-\frac{\log G_-}{2[E:\mathbb Q]} .
\]
For this mixed-isogeny calculation suppose that the finite comparison
supplies an ideal \(\mathfrak n\) of \(E\) satisfying, with ideals
extended via \(\tau\) when necessary,
\[
 \mathfrak n\bar{\mathfrak n}=(\det g_+\det g_-),\qquad
 I_-(\mathfrak n\mathcal O_K)=I_+(\tau\det g_-).
\]
The definite specialization of the index comparison occurs on the
exported dependency path. The displayed mixed ideal identities are
the hypotheses of the following cancellation calculation.
Proposition~\ref{cmeig:degree} below gives an evaluating point on
the indefinite family for which both Lie lattices are integral
direct sums and the mixed ideal hypotheses can be dispensed with.
\src{src/Development/AbelianVarietyTheory/Proofs/DivisorProofs/DegreeImagesEigenLinesCmUnitaryShimuraObject.lean}{122}{Period and Gram determinant of each image};
\src{src/Development/AbelianVarietyTheory/Proofs/DivisorProofs/CmDefiniteLieIndexIdealEqNormIdeal.lean}{20}{Index comparison for the definite specialization};
\src{src/Development/AbelianVarietyTheory/Proofs/DivisorProofs/CmObjectDefiniteLieEigenLatticeDegreeDifference.lean}{21}{Its degree consequence}.

Taking norms gives \(N\mathfrak n=(G_+G_-)^{1/2}\) and
\(NI_+/NI_-=(N\mathfrak n/G_-)^{[K:\tau E]}\). The finite index
adds \(\log NI_\pm/[K:\mathbb Q]\) to the normalized degree in
passing from \(S_\pm\) to \(T_\pm\). Consequently
\[
 \widehat{\deg}T_+-\widehat{\deg}T_-=(p-q)A_\chi/2 .
\]
All dependence on the two Gram matrices has cancelled. This
explains why the degree is independent of \(Z,\tau\) at fixed
host and signature, although the individual local orders may vary.
It also explains why using the CM point \emph{before} proving local
constancy would leave the orders at the other tests undetermined.

\paragraph{The finite degree in the parallel twist.}
Use \(j_\tau=\lambda_{\tau,p}-1-4m\), \(m=fr+1\). Here the
one-based index \(p\) is the multiplicity at \(\tau\) in the
algebraic weight. The positive row block in the period chart has
size \(p_0=q\). In the
source the elementary flat section \(s_\tau\) has entry \(-1\)
on the entire \(\tau\)-block, so the required product is
\(\prod_\tau s_\tau^{-j_\tau}\). For each \(w\) the comparison
needed with its integral model is the equality of the orders at
all tests over \(w\), with common value
\[
 \rho_{\tau w}+D_{\tau w},\qquad
 D_{\tau w}=q\bigl({\rm ord}_w(\tau\mathfrak D_E)
                            -{\rm ord}_w(d_0)\bigr),
\]
where \(d_0\) is the host prime and \(\rho_\tau\) has finite
support. The source defines the corrected degree using
\(\rho_\tau\); denote it by \(-c_1\) if it is independent of
\(\tau,Z\). The norm of the different and the product formula
give
\[
 [F:\mathbb Q]^{-1}\sum_w D_{\tau w}\log Nw
      =q(\log {\rm rd}(E)-\log d_0),\qquad
 {\rm rd}(E)=|\Delta_E|^{1/[E:\mathbb Q]} .
\]
Indeed the norm of the extended different is
\(|\Delta_E|^{[F:E]}\), while that of \(d_0\mathcal O_F\) is
\(d_0^{[F:\mathbb Q]}\). The uncorrected degree of the product
in \eqref{bfam:flat} is therefore exactly
\begin{equation}\label{ctwist:degree}
 \delta\left(\prod_\tau s_\tau^{-j_\tau}\right)
   =\bigl(c_1-q(\log{\rm rd}(E)-\log d_0)\bigr)
       \sum_\tau j_\tau .
\end{equation}
If \(j_\tau=0\) for every \(\tau\), the multiplier is \(1\) and its
integral cost is zero before any comparison. If only
\(\sum_\tau j_\tau=0\), cancellation in \eqref{ctwist:degree}
uses equality of the corrected degrees for the embeddings with
nonzero exponents. For general \(j_\tau\), suppose as in the composition
that \(|c_1|\le9rL\), \(\log{\rm rd}(E)\le3L\), and
\(|\log{\rm rd}(E)-\log d_0|\le\log{\rm rd}(E)\).
The last inequality follows when \(E\) contains the host quadratic
field of absolute discriminant \(d_0\): the discriminant tower
gives \(2\log{\rm rd}(E)\ge\log d_0\).
At each \(w\) above a prime outside the bad set \(B\), a test of
ramification index one makes the common order integral. Take the
product first, then apply \eqref{bfam:flat} over \(F\) once.
Writing \(W=\sum_{\tau,i}|\lambda_{\tau i}|\), the logarithmic
cost for both twists together is at most
\[
 2gm\log(2\pi)+12rL\,\bigl(W+f(1+4m)\bigr)
 +\sum_{\ell\in B}\log\ell
 +\frac{\log|\Delta_F|}{2[F:\mathbb Q]} .
\]
Here \(q\le r\) and \(\sum_\tau|j_\tau|\le W+f(1+4m)\).
The estimate holds for every component and embedding at the same
threshold as the field, level and degree bounds. On the selected
classical family Theorem~\ref{cmeig:splitmodel} gives the
uncorrected degree \((p-q)A_\chi/2\) of \(s_\tau\). Subtracting
the degree of \(D_\tau\) gives the corrected degree
\[
        -c_1=\tfrac12(p-q)A_\chi-q(\log{\rm rd}(E)-\log d_0).
\]
Thus the product coefficient in \eqref{ctwist:degree} is
\(-\tfrac12(p-q)A_\chi\), including zero in balanced signature.
On the chosen family
the bound on \(c_1\) follows numerically from
\(|A_\chi|\le3L\), \(0\le\log d_0\le L\),
\(0\le\log{\rm rd}(E)\le2L\), and \(p+q=r\).
In fact these bounds give \( |c_1|\le(3/2+2)rL\).
\src{src/Development/AbelianVarietyTheory/Proofs/DivisorProofs/FlatDegreeLogFormAbsLe.lean}{55}{Period and discriminant bounds on the selected stage};
\src{src/Development/AbelianVarietyTheory/Definitions/Foundations/ExistsFlatSecRankZero.lean}{751}{The corrected degree in Lie multiplicities};
\src{src/Development/AbelianVarietyTheory/Proofs/DivisorProofs/ExistsTwSecOfConst.lean}{87}{Conversion to the degree of the actual product};
\src{src/Development/AutomorphicSpectrum/Proofs/GeometryProofs/KroneckerCharacterFrameBound.lean}{20}{The period in the degree of the uncorrected frame};
\src{src/Development/AutomorphicArithmetic/ThetaCorrespondence/ThetaPairLifts/HodgeBundleWeights/ShortSpanCompositionClosed.lean}{54}{The algebraic weight index in the parallel exponent}.

There is a weaker sufficient test when only a lower bound for the
product degree is sought. On a common normal proper flat model at
\(w\) carrying the actual lines and tests of Lemma~\ref{ctwist:vertical},
let \(a_{\tau i}\) be the order of \(s_\tau\) along \(C_i\), and use
the base \(\mathcal O_{F,w}\). Set
\[
 b_w(j)=\min_i\frac{-\sum_\tau j_\tau a_{\tau i}}{m_i}.
\]
Then \(b_w(j)\) is a lower bound at \emph{every} test for this
product without any special-fibre triviality assumption. Indeed,
raise the product to a power \(M\) clearing the denominator of
\(b_w(j)\), and divide by \(\pi^{Mb_w(j)}\). Its divisor is
effective, so normality extends it over the model. Pullback to
each test yields the inequality. A generic trait of a minimizing
vertical divisor realizes equality among all trait maps. On a
reduced special fibre \(m_i=1\) and \(b_w(j)\) is integral.
Thus, if these lower bounds have finite support, are integral
outside \(B\), and give a sufficiently large normalized degree
in \eqref{bfam:flat} with the product's actual metric, the same
integral spanning estimate follows. This criterion applies on each
chosen model separately; a bound uniform in the window still has
to be supplied for the indicated weighted minima.

\paragraph{A product CM point on each indefinite marked row.}
For the fields of Proposition~\ref{host:art} there is another
choice of evaluating point. Write \(k=\mathbb Q(\sqrt{-d_0})\),
\(F_0=E^+\), \(O=O_E\), and \(\Phi\) for the embeddings inducing
the chosen embedding of \(k\). For \(\tau\in\Phi\) write \(p,q\)
for the Lie multiplicities at \(\tau,\bar\tau\), respectively.
Write \(p_0\) for the number of positive Hermitian row directions.
The selected period map conjugates the first \(p_0=q\) coordinates
and is \(\tau\)-linear on the last \(r-p_0=p\) coordinates.
Keep the period convention \(\Pi=\mathfrak a L\) of
Lemma~\ref{bfam:minimum}, so \(L\) is Hermitian self-dual.
\src{src/Formalism/LocalToGlobal/ProjectiveLattices/PetersonNormSubspace.lean}{67}{The conjugated positive block in the period map};
\src{src/Formalism/LocalToGlobal/ProjectiveLattices/PetersonNormSubspace.lean}{75}{The positive row index on a good chart};
\src{src/ReciprocityTheory/ReductiveGroups/ArithmeticGeometry/BirationalGeometry/ToroidalCompactification/FlagBasis.lean}{76}{The Lie multiplicities of the two embeddings}.

\begin{lemma}\label{cmeig:row}
If \(p,q\ge1\), every marked row in this family has, up to rational
isometry, underlying lattice
\[
        J e_1\ \perp\ O e_2\ \perp\cdots\perp\ O e_r,\qquad
        J\bar J=O,\qquad
        h(e_i,e_i)=
        \begin{cases}1&i\le p_0,\\-1&i>p_0.\end{cases}
\]
The transported mark need not split into marks of the summands.
\end{lemma}
\begin{proof}
The two imaginary quadratic subfields of discriminants
\(-d_0,-q_0\) in Proposition~\ref{host:art} have coprime
discriminants (here \(q_0\) is the auxiliary prime). At each
rational prime at least one of them is unramified. Its inertia
restriction is trivial, whereas complex conjugation acts on it
nontrivially. Thus \(E/F_0\) is split or unramified at every finite
place, since relative inertia is the intersection with
\(\operatorname{Gal}(E/F_0)\).

At a split place a basis of the first factor and its pairing-dual
in the second identify a self-dual lattice with the standard one.
At an unramified quadratic place its reduction is a nondegenerate
Hermitian space over the quadratic residue extension. There is a
vector of nonzero norm: otherwise polarization of \(h(x+ay,x+ay)\)
and the perfect residue trace pairing would give \(h=0\).
Lifting gives a unit norm, which can be scaled to either sign,
since the norm maps unramified units onto units (use the residue
norm and lift successively with the residue trace). Its integral
orthogonal complement is again self-dual. This gives the local
orthogonal basis by induction, including in residue characteristic
two. At infinity the signatures agree with the displayed form.
For clarity, the rational isometry follows already from
Hasse--Minkowski: the underlying quadratic form \(h(x,x)\)
represents the prescribed first sign at every place, hence globally.
For the latter implication apply isotropy to \(h(x,x)-a t^2\);
if \(t=0\) gives a nonzero isotropic vector, the original quadratic
space contains a hyperbolic plane and represents \(a\).
Split off the Hermitian line of norm \(a\); by Witt cancellation its
complement is locally isometric to the remaining diagonal form.
Induction now identifies the two rational Hermitian spaces.

Put \(L_0=\bigoplus O e_i\), \(G=U(h)\). Choose
\(g\in G(\mathbb A_{F_0,f})\) carrying \(\widehat L_0\) onto
\(\widehat L\) using the local isometries, choosing the identity
where the two lattices coincide. Write \(t=\det g\) and
\(d(t)=\operatorname{diag}(t,1,\dots,1)\). The determinant-one
groups have strong approximation away from the infinite places:
\(\mathrm{SU}(h)\) is simply connected, absolutely almost simple
(\(r\ge2\)), and noncompact at a real place by \(p,q\ge1\).
This is \href{https://www.mathnet.ru/links/02c4d3528af06a248862fb1aa3ba53f9/im2220_eng.pdf}
{Platonov, Theorem A}. Thus the open set
\(g K_0 d(t)^{-1}\cap\mathrm{SU}(h)(\mathbb A_{F_0,f})\)
contains a rational element \(\gamma\): it contains
\(g d(t)^{-1}\), and \(K_0\) is the open stabilizer of
\(\widehat L_0\). Then \(\gamma^{-1}L\) has completion
\(d(t)\widehat L_0\). Its first ideal is
\(J=E\cap t\widehat O\); \(t\bar t=1\) gives \(J\bar J=O\).
Transport the entire given finite marking by \(\gamma^{-1}\).
The resulting polarized marked row is isomorphic to the one
started with, as in Lemma~\ref{munf:marks}. No approximation to
the identity of the marking, or change of its level, is needed.
\end{proof}

\begin{proposition}[The degree on the product point]\label{cmeig:degree}
On each classical marked period component of
Lemma~\ref{cmeig:row} there is a point over a finite extension of
the coefficient field with everywhere good abelian reduction. For
every \(\tau\in\Phi\) its two integral Lie determinants in the
polarization metric of the cover satisfy
\[
 \widehat{\deg}T_+-\widehat{\deg}T_-=
   \frac{p-q}{2}\left(\gamma_{\rm Eul}+
     \log\frac{2\pi}{{\rm rd}(E)}
       -\frac{L'(1,\chi)}{L(1,\chi)}\right).
\]
All finite places, including those ramified in \(E\), occur with
their ordinary integral eigensublattices.
\end{proposition}
\begin{proof}
In the representative of Lemma~\ref{cmeig:row} take the flag of
the positive and negative coordinate axes. The period lattice
is the orthogonal direct sum of the fractional ideals
\(I_1=\mathfrak a J,\ I_i=\mathfrak a\ (i>1)\).
The last \(p\) factors have CM type \(\Phi\), and the first \(q\)
have its conjugate, by the period map above.
Let \(C\) be the elliptic curve with complex lattice \(O_k\),
with the indicated action (conjugate the action for the opposite
type). It is defined with its endomorphisms over a number field
and acquires good reduction everywhere after finite extension:
its CM \(j\)-value is an algebraic integer, hence gives potential
good reduction at the finitely many places of bad reduction.
The Serre construction \(I_i\otimes_{O_k} C\) exists integrally
since \(I_i\) is projective over \(O_k\): an idempotent on a
finite direct power of \(C\) realizes it as a direct factor.
Its periods at the chosen embedding are exactly the images of
\(I_i\) at the extensions of that \(k\)-embedding. The alternating
trace polarization on each factor is induced from a positive
Hermitian form \(\operatorname{Tr}_{E/k}(\beta u\bar v)\).
Indeed divide its purely imaginary trace coefficient (with the
coordinate sign) by the coefficient of the standard principal
elliptic polarization, a signed generator of
\(\mathfrak D_{k/\mathbb Q}^{-1}\); the ratio \(\beta\in E^+\)
is totally positive by the Riemann positivity condition.
Perfection is exactly
\[
       (\beta)\,I_i\bar I_i\,\mathfrak D_{E/k}=O .
\]
This also follows by taking the trace-dual ideal. Hence the
polarized product has exactly the required period lattice and
polarization. For the polarized tensor construction see
\href{https://arxiv.org/abs/1507.07607}{Amir-Khosravi, \S2}.
The transported finite torsion marking is defined after another
finite extension, irrespective of how it mixes factors. It is
needed on the characteristic-zero point; the underlying product
extends over the full ring of integers. In a classical PEL moduli
chart this constructs the point of the prescribed component,
including its canonical rigidified graph line. An additional
finite étale marking cover is handled by lifting the point to
the given component of the cover, whose map to the underlying
connected component is finite étale and surjective.

Here is the integral and metric calculation for a single factor
with ideal \(I\) and tensor coefficient \(\beta\). Enlarge the
number field \(K\) of definition to split \(E\), and put
\(\mathfrak D=\mathfrak D_{E/k}\). In
\(I\otimes_{O_k}O_K\), the intersection with the coordinate line
at an active embedding \(\nu\) is \(\nu(I\mathfrak D)O_K\).
In fact test membership in the lattice by its dual. The trace
dual of \(I\) over \(O_k\) is \(I^{-1}\mathfrak D^{-1}\);
this remains the dual after flat base change. On the split
generic algebra the trace is the sum of coordinates. The vector
supported at \(\nu\) with entry \(z\) thus pairs integrally
exactly when \(z\nu(I^{-1}\mathfrak D^{-1})\subset O_K\), as
asserted. Tensoring with the invertible \(O_K\)-module
\(\operatorname{Lie} C\) gives the actual Lie eigenline of the
abelian scheme, since taking Lie commutes with the idempotent
defining the tensor construction.

On that coordinate its squared metric at \(\sigma:K\to\mathbb C\)
is the squared elliptic Lie metric multiplied by
\(\sigma\nu(\beta)\): diagonalizing the trace Hermitian form in
the embedding coordinates gives coefficients \(\sigma\nu(\beta)\);
the polarization convention multiplies both Hermitian forms by
the same \(\pi\). For \(e=[E:\mathbb Q]\), normalized ideal and
metric degrees therefore give
\[
\begin{split}
 2\widehat{\deg}T_\nu(I)
 &=2\widehat{\deg}\operatorname{Lie}C
   -\frac{2\log N(I\mathfrak D)+\log|N_{E/\mathbb Q}\beta|}{e}\\
 &=2\widehat{\deg}\operatorname{Lie}C-\frac{\log N\mathfrak D}{e}.
\end{split}
\]
The second equality uses perfection, so even the fractional ideal
on the first coordinate disappears. The norm is the absolute
ideal norm and \(N\mathfrak D=|\Delta_E|/d_0^{e/2}\).

For \(C(\mathbb C)=\mathbb C/\Pi_C\) with Euclidean period area
\(a_C\), principal polarization gives the Hermitian coefficient
\(1/a_C\). The tangent metric is \(\pi/a_C\), so its dual on
\(dz\) is \(a_C/\pi=(i/2\pi)\int dz\wedge d\bar z\).
Thus \(-\widehat{\deg}\operatorname{Lie}C\) is exactly the
stable elliptic Faltings height with the normalization of
\href{https://people.math.wisc.edu/~tonghaiyang/KRY2.pdf}
{Kudla--Rapoport--Yang, (10.14)--(10.15), Proposition 10.10}.
Their quadratic CM formula gives
\[
 2\widehat{\deg}\operatorname{Lie}C
 =\gamma_{\rm Eul}+\log(2\pi/\sqrt{d_0})
                  - L'(1,\chi)/L(1,\chi).
\]
This applies to either elliptic orientation. Substitution gives
the same \(2\widehat{\deg}T_\nu(I)=A_\chi\) on every active
coordinate; orthogonality and the integral direct product give
the two determinants by addition over \(p,q\) coordinates.
\end{proof}

In particular no comparison between a mixed isogeny image and
its saturation is needed to evaluate this point. If the vertical
orders on a selected component agree at all tests and its
archimedean norm is constant, the degree is that of this point
by the normalized base-change calculation preceding
\eqref{ctwist:degree}.

\paragraph{Vertical orders at hyperspecial places.}
There is a direct good-reduction comparison for the ordinary Lie
frames of the classical PEL problem. It also distinguishes a prime
in the level from a prime ramified in the endomorphism algebra.

\begin{proposition}\label{cmeig:hyperspecial}
Suppose the characteristic-zero cover on the chosen component
is a classical finite marking level for its principally polarized
\(O_E\)-abelian scheme, and its finite tests use the integral Lie
eigensublattices of connected semiabelian extensions of that same
scheme. Let \(w\mid\ell\) in the splitting coefficient field \(F\),
with \(\ell\nmid 2\Delta_E\). Each rational horizontal determinant
generator has the same normalized order at all semiabelian DVR
tests on the component above \(w\). If its marked level itself is
hyperspecial at \(\ell\), this common order belongs to \(\mathbb Z\).
\end{proposition}
\begin{proof}
First take a neat prime-to-\(\ell\) marking with full hyperspecial
level at \(\ell\). The algebra is unramified at \(\ell\), its
involution is of unitary type and the principal alternation is
perfect. Thus this is a good reduction PEL problem. With smooth
projective cone data it has a proper smooth toroidal model with
universal semiabelian scheme by
\href{https://www.kwlan.org/articles/cpt-PEL-type-thesis-revision-dense.pdf}
{Lan, Theorems 6.4.1.1 and 7.3.3.4}.
Base change to the completed coefficient DVR. The geometric
components of a proper smooth family form a finite étale scheme;
the selected geometrically connected generic component defined
over the coefficient field extends to an open and closed component
with geometrically connected special fibre. That fibre is smooth,
so is reduced and irreducible. The split embedding idempotents of
\(O_E\otimes_{\mathbb Z}O_{F,w}\) are integral. The indicated Lie
determinant ratio is therefore a line bundle on the model.
Its characteristic-zero horizontal generator extends as a unit at
the toroidal boundary by the unipotent determinant calculation
above Lemma~\ref{ctwist:vertical}. On the whole model its divisor
is an integral multiple of the single reduced special fibre.

Every test extends to the toroidal model by properness. The
pullback of its connected semiabelian group is the given connected
semiabelian extension by uniqueness over a trait (SGA~7 I, IX,
Proposition 3.2). This is an isomorphism respecting the endomorphisms
by generic-fibre uniqueness. Hence its Lie eigensummands, including
those of a toric degeneration, are exactly the test lattices.
Pulling back the divisor in an integral local generator gives the
same normalized integer at every such test.

For general finite marking choose another full neat level of order
prime to \(\ell\) and a geometric connected component of the common
refinement lying over the original component. Over a finite field
extension it and both projection maps are defined. Fix one place
of this field above \(w\). The common component is finite étale
and surjective over each of its two characteristic-zero image
components, as follows by forgetting finite torsion bases. Its two
pulled-back determinant generators differ by a scalar in the
field of definition: both are horizontal for the same connection,
so their quotient has zero differential on the smooth geometrically
connected component. On the auxiliary prime-to-\(\ell\) component
the generator exists by determinant-one monodromy at its full neat
level, as proved above Lemma~\ref{ctwist:vertical}. Any test on the
original component first extends to the fixed base place and
then lifts generically to the common component after a finite
extension of its fraction field, by surjectivity. Connected
semiabelian pullbacks agree by the same uniqueness. The order on
the lifted test is the order on the auxiliary component plus the
valuation of the fixed scalar, both independent of the test after
normalization to \(F_w\). This proves the claim without requiring
an integral level-\(\ell\) marking of a semiabelian group.
\end{proof}

\begin{theorem}[Integral determinant comparison at the ramified primes]
\label{cmeig:splitmodel}
For the fields of Proposition~\ref{host:art}, under the same
classical moduli and test identifications as in
Proposition~\ref{cmeig:hyperspecial}, the common normalized order
exists at every finite place. For each indefinite component the
uncorrected arithmetic degree of the elementary horizontal
determinant is therefore \((p-q)A_\chi/2\).
\end{theorem}
\begin{proof}
Fix \(w,\tau\), and use a neat marking away from its residue
characteristic \(\ell\), with maximal level at \(\ell\) stabilizing
the single self-dual lattice (and its periodic translates).
As in Proposition~\ref{cmeig:hyperspecial} we can pass to this
level and to a finite local splitting field. At every place of
\(E^+\) above \(\ell\) the quadratic algebra is split or
unramified by Lemma~\ref{cmeig:row}. These are exactly the
unitary cases (AL) and (AU) of the Pappas--Rapoport splitting
model. Its order \(O_E\) is maximal, and its alternating lattice
has index one in its dual. These are the hypotheses at the
maximal level of \href{https://stephane-bijakowski.perso.math.cnrs.fr/PRgeometry.pdf}
{Bijakowski--Hernandez, \S2.4 and Theorem 2.30}, which gives a
smooth splitting model in these cases. At two the imaginary
quadratic subfield \(k\) already splits by Lemma~\ref{host:different},
so only the linear case (AL) occurs there.

Use the ordering of the \(E\)-embeddings in
\href{https://www.kwlan.org/articles/cpt-spl.pdf}
{Lan, Choices 2.3.1, Definition 2.3.3}: put \(\tau,\bar\tau\)
first as a consecutive pair in their \(E^+\)-place, and order the
other conjugate pairs consecutively. In the split case the pair
is in distinct integral field factors. In the inert unramified
quadratic case its restrictions to the maximal unramified subfield
are distinct, so the two are again separated by integral
idempotents over the splitting base. On these two different
summands the first quotient of the decreasing Hodge filtration
is labelled respectively by \(\tau,\bar\tau\).
In Lan's interleaved ordering the intervening graded piece on
the other summand is zero. In the increasing notation of
Bijakowski--Hernandez these are the two \emph{last} quotients.
Reversing the ordering within each list, with conjugate lists
paired term by term, is permitted in their \S2.4.

Here are the model comparisons used with the smoothness theorem.
For the single periodic self-dual lattice the open splitting model
parametrizes the polarized Hodge flags of the given ranks. Separating
the unramified idempotents turns Lan's interleaved flags into
the flags for the separate unramified embeddings; conversely they
are interleaved by direct sum. The polarization conditions can
also be compared on the common generic fibre, without identifying
their scheme structures on nonflat thickenings. Indeed both flag
problems are closed in the same relative product of flag varieties
over the naive moduli problem, by the endomorphism and pairing
conditions. They give the same unique flags on the generic fibre.
The Bijakowski--Hernandez model is smooth, hence equals the flat
closure of this generic fibre in that flag space, and is normal.
Lan's normalized open model is the normalization of that closure
on the selected generic components: the forgetful map of the
normalized underlying moduli to the naive moduli is finite and
generically open and closed at full lattice-stabilizer level
(Proposition 2.2.11). Finiteness persists after pullback to the
flag space and after normalization over the excellent base
(Definition 2.4.5 and Proposition 2.3.10). Consequently the normalized
open model on those components is the smooth one, by finite
birationality over a normal scheme.

Lan's Theorem 3.4.1 and Proposition 3.4.14 give a proper flat
toroidal compactification with geometrically regular fibres for
smooth cone data: geometric regularity of each fibre is
equivalent to that of the open splitting model. It is thus smooth.
By Proposition 3.1.2 and Definition 3.1.3 the extended filtration
is a filtration by subbundles on the Hodge bundle of the
semiabelian extension (identified with that of its dual by the
principal polarization, with conjugate labels if using the
dual-action convention). In particular the two locally free
first quotients above exist along the whole boundary.

For any trait write \(H\) for its integral invariant
differentials in the relevant unramified summand. If \(H'\) is
the kernel of its first quotient, then \(H'\) is saturated and
its generic fibre is the sum of all the other eigenspaces in
that summand. Hence
\[
   (H/H')^\vee=\{u\in H^\vee:u(H')=0\}
           =H^\vee\cap (H^\vee_{\rm frac})_\tau .
\]
The identical argument on the conjugate summand gives the
second test lattice. This uses the end quotient of the
filtration: in general an intermediate graded piece has a
different lattice. The determinant ratio of these dual quotients
is therefore an invertible sheaf which gives exactly the
ordinary test metric at \(\ell\). Its generic horizontal section
is a unit across the characteristic-zero boundary, as before.
Over the fixed geometric component the smooth proper model
has a single reduced special fibre. Lemma~\ref{ctwist:vertical}
and the common-refinement argument of
Proposition~\ref{cmeig:hyperspecial} give equality of all orders,
including on semiabelian tests. A different ordering or model
can be used for a different pair \(\tau,\bar\tau\), since orders
are calculated pairwise before multiplying sections.

Finally compare with the CM point of Proposition~\ref{cmeig:degree}
at each place and with its archimedean norm. Its finite extension
preserves normalized arithmetic degrees. This gives the claimed
global degree and hence the uncorrected product degree in
\eqref{ctwist:degree}.
\end{proof}

\paragraph{The selected good cover.}
The level condition on \(Y\) fixes its auxiliary marking to have
order one. Its finite marking is the half-marking at the odd integer
\(M\ge3\).
The characteristic-zero comparison can be made from the moduli
property. The universal object is proper and semiabelian, hence
abelian. Its rigidified line is relatively ample and its
polarization kernel is trivial, hence its polarization \(\lambda\)
is principal. The adjunction property for its biextension gives
the specified Rosati involution on \(O_E\). Conversely from these
PEL data pull the normalized Poincar\'e bundle back along
\((1,\lambda w_0)\). Its polarization is
\(\lambda(w_0+\bar w_0)=\lambda\); the rigidification, symmetry
and cube law are the canonical ones. Its biextension is the
one for \(\lambda\). This reconstruction also recovers the
original line from its graph condition. The rigidified isomorphism
is unique: an automorphism is multiplication by a unit on the
abelian scheme, which comes from the base by properness and
geometric connectedness and equals one at the identity.

Over the splitting field the determinant condition is precisely
the \(p,q\) multiplicity condition. On a reduced
characteristic-zero base the Hodge eigensummands are vector bundles
and their ranks can be tested at geometric points. The mark kills
\(\mathfrak m\) and on the complementary summand is the reduction
of a polarized Tate marking. By CRT these restrictions determine
it on the whole lattice. Its compatibility at all torsion orders
in the definition of an object is a condition on the local
lattice genus, not a choice of additional sections. More explicitly,
in a connected complex PEL family the Betti lattice, its pairing,
action and finite mark are locally constant; a compatible marked
Tate isometry at one point transports to every point. Equivalently
at a geometric point the reduction of a compatible tower is a
point of the prescribed compact-open marking orbit. Conversely
any representative of that orbit gives the required compatible
reductions at all \(n\). To check the roots of unity, fix the
primitive \(M\)-th root in the coefficient field and, at a geometric
point, extend it to a compatible system \((\zeta'_n)\). Relative
to the tower \((\zeta_n)\) of a Tate marking there is a unique
\(u\in\widehat{\mathbb Z}^{\times}\) with \(\zeta'_n=\zeta_n^u\)
for all \(n\ge1\). The reference lattice has an
integral similitude of multiplier \(u\) fixing the half-mark.
At \(\ell\mid M\), use \(1\) on the marked half and \(u_\ell\) on
its pairing-dual half. At \(\ell\nmid M\), use a central scalar of
norm \(u_\ell\) at each place of \(F_0\) over \(\ell\), using the
surjectivity of split and unramified unit norms in Lemma~\ref{cmeig:row}.
Precomposing a Tate marking by this similitude gives the prescribed
tower with the same marked sections. At order \(M\) the Weil pairing
with the marked basis and the fixed root then determines the basis
on the opposite half uniquely. This construction works over the
base at order \(M\) by perfect duality of finite étale groups. The
higher towers occur only in the geometric-point existence test.
This identifies the half-marking with the corresponding fine PEL
level on the selected components. After finite extension the
source moduli property gives a classifying map for the classical
universal object on its reduced smooth base, and the classical
moduli property gives the reverse map on the reduced smooth
source base. The reconstructed universal objects are isomorphic,
so both composites are the identity by uniqueness. The maps
preserve the underlying group and polarization, also after
restricting to a specified marked row.
\src{src/ReciprocityTheory/ReductiveGroups/ArithmeticGeometry/BirationalGeometry/ToroidalCompactification/TwoWeightEqOneAddRoot.lean}{131}{The compatible torsion isometries and primitive roots};
\src{src/ReciprocityTheory/ReductiveGroups/ArithmeticGeometry/BirationalGeometry/ToroidalCompactification/TwoWeightEqOneAddRoot.lean}{144}{The two marks, signature and geometric-point existence condition};
\src{src/ReciprocityTheory/ReductiveGroups/ArithmeticGeometry/BirationalGeometry/ToroidalCompactification/PushInvertible.lean}{29}{The universal object and its reduced-base moduli property};
\src{src/Development/AbelianVarietyTheory/Definitions/Divisors/NormNegCmConstantEqAtAllEmbeddings.lean}{187}{The selected level, period charts and polarization scale}.

On a test DVR the definition supplies a semiabelian group whose
generic fibre is cartesian over this universal group. Its fibres
are geometrically connected. The test vectors form a basis for
all integral eigenvectors at each labelled embedding, so they give
exactly the intersections used in Theorem~\ref{cmeig:splitmodel}.
Indeed the images of algebraic integers in the coefficient field
belong to the valuation ring; the equality expressing the
eigenvector property therefore holds integrally if and only if
it holds over the fraction field, by torsion-freeness of the
smooth Lie lattice. The theta generator in the test contributes
no factor to the parallel character. At infinity the comparison
uses the polarization metric \(\pi H\) on Lie vectors; this is
the metric used in Proposition~\ref{cmeig:degree}. Thus the
ordinary parallel determinant comparison follows on the selected
principal family with its given period readings.
\src{src/Formalism/LocalToGlobal/ProjectiveLattices/VAtVNeBot.lean}{53}{Cartesian semiabelian tests and the universal integral eigenbasis condition};
\src{src/ReciprocityTheory/ReductiveGroups/ArithmeticGeometry/SemiabelianSchemes/AbelianSchemeQuotients/IsAbelianScheme.lean}{41}{Connected geometric fibres of the semiabelian group};
\src{src/Development/AbelianVarietyTheory/Definitions/CMLattices/IsDiffOf.lean}{101}{Polarization norm in the CM comparison}.

\begin{corollary}[Short sections in the final weight on a good cover]
\label{cmeig:final}
Fix one good deep cover \(Y\) as in Corollary~\ref{crev:good}
on a selected host, with \(p,q\ge1\). Take the admissible final
weight \(\lambda\) and the level, denominator, field and bad-prime
size bounds of Theorem~\ref{bfam:oncharts}.
On each marked component \(Z\), its regular cusp rule space of
weight \(\lambda\) is spanned over the coefficient field by
integral sections \(s_i\) with
\[
     \log\|s_i\|_{\sigma}\le
          C\bigl((frL)^3 W+(frL)^5\bigr),\qquad
     W=\sum_{\tau,i}|\lambda_{\tau i}|,
\]
at every embedding in the theorem. The same constant works over
these choices under the common stage bounds.
\end{corollary}
\begin{proof}
Start with the integral generators of Corollary~\ref{crev:good}.
Multiply by the key of Theorem~\ref{ctwist:semi} and by the
product with exponents \(-j_\tau\) in \eqref{ctwist:degree}.
Theorem~\ref{cmeig:splitmodel} computes the latter degree at the
actual integral eigenlattices. Outside the excluded primes the
level is hyperspecial, so the common order is integral by
Proposition~\ref{cmeig:hyperspecial}. At the remaining primes
round the product order once as in \eqref{bfam:flat}.
Proposition~\ref{host:period} supplies the uniform numerical
degree bound. Both unit sections and their inverses preserve the
cusp condition by Corollary~\ref{crev:twist}. The final-weight
estimate of Theorem~\ref{bfam:oncharts} now gives the claim.
This argument starts with \(Y\) and the stated numerical bounds.
Selecting such covers compatibly at all stages and identifying the
restriction maps on the modified flags enter separately in the
Hilbert--Samuel estimate.
\end{proof}



Once \(\beta\) is known, the required archimedean multiplication
estimate has no further dimension factor. For the shared Petersson
measure on the component, let \(B(y)\) be the supremum of
\(|u(y)|^2/\|u\|_2^2\) over nonzero sections of the first frame space,
using the fiber norm of its weight. Suppose \(B(y)\le\beta\) on the
support of the measure, and the fiber metric on a product is the
tensor product of the fiber metrics. For any frame section \(u\)
and higher weight section \(v\), pointwise multiplication and
integration give
\[
 \|uv\|_2^2=\int |u(y)|^2|v(y)|^2\,d\mu(y)
       \le \beta\|u\|_2^2\|v\|_2^2.
\]
Using the original frame norm for \(\|u\|_2\) gives the operator
bound in the determinant argument. This requires the quotient
measure and fiber normalizations for the three weights to agree
under tensor product. A bound for the kernel relative to a
separately rescaled measure does not imply the same \(\beta\).
\src{src/Development/AbelianVarietyTheory/Proofs/DivisorProofs/PeterssonNormMulLeSqrtBergmanMulComponent.lean}{79}{Component product estimate}.

\begin{proposition}[Absorption of the component errors]\label{newtheta:absorption}
Assume \eqref{newtheta:budgets} and \eqref{newtheta:finiteHS},
\(1\le D,\ 3D+5\le T,\ \delta\ge1\), and
\[
 b\ge1,\quad \log b\le P_0,\quad TP_0\le\sqrt b,\quad
 z=7(D^2/T+D\rho)\le\tfrac12,\quad \rho\ge0,\quad
 zP_0\kappa\le b .
\]
If \((1-z)d_0\le B\) and \(g_M\ge g-d_0c\), \(c\ge0\), then
\[
 g_M\ge B\big[Th-3(D+1)|h|
       -26(26D^2+2zT)\log(2+|h|)-4000b-2c\big].
\]
\end{proposition}
\begin{proof}
Write \(P=P_0,\ v=\kappa/T\) and \(L_*=27Pv+92P\). Then
\(TP\le b,\ (TP)^2\le b,\ \log D\le\log T\le P/2\).
The expression \(2+|h|+\log\Lambda\) is at most
\((2+|h|)(1+\log\Lambda)\). Since
\(\log(1+\log\Lambda)\le\log\Lambda\), \eqref{newtheta:budgets}
gives \(L_h\le L_*+26\log(2+|h|)\). Specifically, the terms free
of \(h\) are at most \(27P(v+1)+26(P+P/2+P)=L_*\).
The rank comparisons give \(0\le u\le z\le1/2\),
\(d_0(1-u)=B\), and \(d_0u/B=u/(1-u)\le2z\).
The bracket multiplied by \(u\) is nonnegative. Substitution in
\eqref{newtheta:finiteHS} leaves the claimed height coefficient
and the following upper bound on the error independent of \(h\):
\begin{align*}
 J+26D^2(L_*+\log T)&+(T+D+1)\log\beta+(T+D^2)(G+2)\\
 &+2z\big[T(\log2+L_*+\tfrac12\log\beta)
                   +\log\Lambda_D+\log\Lambda_T\big].
\end{align*}
For explicit bounds use
\[
 D^2P\le b/14,\qquad D^2Pv\le b/7,\qquad zTPv\le b .
\]
Indeed \(D^2\le zT/7\le T/14\); the last inequality displayed
is \(zP\kappa\le b\). Since \(D\le D^2\) and \(DP\le TP\le b\),
the first cost \(J\) is at most \(6b/7+8b<9b\).
The second is at most \(26(27/7+92.5/14)b<300b\).
The third is at most \(2TP\le2b\). For the fourth,
\(G+2\le5DP\) by \(D,P\ge1\), and \(T+D^2\le2T\);
hence it is at most \(10TDP\le10(TP)^2\le10b\).
Finally \(zT\log2\le b/2,\ zT L_*\le27b+46b,\)
\(zT\log\beta/2\le b/4\), and
\(z(\log\Lambda_D+\log\Lambda_T)\le2b\).
Thus the last cost is less than \(200b\). Their sum is bounded by
\(4000b\). Since \(d_0\le2B\), the transfer costs at most \(2Bc\).
The argument nowhere requires a sign for \(h\).
\end{proof}

\begin{proposition}[Summing the component degrees]\label{newtheta:aggregate}
Let \(I\) be a finite nonempty set of components. For \(i\in I\) let
\(d_i>0\) be its rank, \(h_i,\ell_i,g_i\) its height, comparison
slope, and transferred degree, and \(B_i>0\) its degree share.
Let \(0\le \rho<1,\ z\ge0,\ D\ge0,\ T\ge3D+5,\ A>0,\ R\ge0\).
Assume for some \(b_0>0\), \(b_0\le B_i\le b_0/(1-\rho)\) and
\(B_i\le d_i\). Suppose
\[
 g_i\ge B_i[Th_i-3(D+1)|h_i|-A\log(2+|h_i|)-R],
 \qquad 2\sum_i g_i/\sum_i d_i\le\mu .
\]
If \(s,c,X\ge0\) satisfy
\[
 h_i+s\ge0,\qquad
 (1-\rho)(\ell_i+s)_+\le h_i+s+c,\qquad
 |I|^{-1}\sum_i\ell_i\ge-X ,
\]
and for \(E\ge1,\ C_1,C_2\ge0\)
\[
 \max\{TX,Tc,s(6D+10+T(z+2\rho))\}\le C_1E,\qquad
 A\log\max(A,1)+R\le C_2E ,
\]
then \(\mu\ge-2(3C_1+C_2+4)E\).
\end{proposition}
\begin{proof}
Set \(K=3D+4\), \(U=T-K>0\), \(B_\Sigma=\sum_i B_i\).
For \(v\ge0\), \(\log t\le t-1\) at \(t=(2+v)/A\) gives
\(A\log(2+v)\le v+2+A\log\max(A,1)\).
Since \(|h_i|\le h_i+2s\), the bracket in the component bound is
at least \(U(h_i+s)-(T+K)s-2-C_2E\).
Using nonnegativity of the positive parts and their unweighted mean,
\[
 \frac{\sum_i B_i(\ell_i+s)_+}{B_\Sigma}
 \ge(1-\rho)|I|^{-1}\sum_i(\ell_i+s)_+
 \ge(1-\rho)(s-X)_+ .
\]
Consequently the weighted average of \(h_i+s\) is at least
\((1-\rho)^2(s-X)_+-c\ge(1-2\rho)s-X-c\).
The last inequality follows from \((s-X)_+\ge s-X\),
\((1-\rho)^2\ge1-2\rho\), and \((1-\rho)^2\le1\).
It follows that
\[
 \frac{\sum_i g_i}{B_\Sigma}
 \ge-[2K+2\rho U]s-U(X+c)-2-C_2E
 \ge-(3C_1+C_2+2)E .
\]
Indeed \(2K+2\rho U\le6D+10+T(z+2\rho)\) and \(U\le T\).
As \(B_\Sigma\le\sum_i d_i\) and the last bound is nonpositive,
the same lower bound holds with \(\sum_i d_i\) in the denominator.
Multiplication by two proves the assertion, including the stated
slightly larger constant.
\end{proof}
\src{src/Development/AbelianVarietyTheory/Proofs/DivisorProofs/UniformLowerBoundOfSums.lean}{16}{Component shares and covering slope}.

\paragraph{The geometric hypotheses on the family used in Section~\ref{main:assembly}.}
Fix one selected scale, a stage in the prescribed range and its
deep cover \(Y\) with coefficient field \(F\), level \(M\), and
component \(Z\). Its auxiliary marking has order \(A=1\) by the
selected good-level condition. In the indefinite short-section
step one has \(p,q\ge1\), \([E:\mathbb Q]/2>2\), \(p+q=r\).
In the notation of the source, the stage hypothesis is
\(Y.\mathrm{Good}(e,o)\), with
\(\log M\le30L\), host field size at most \(3L\), \(\log(fr)\le L\)
and \(o\le391(frL)^2\). The dominant final weight has its block
gap at least \(8fr+10\) and logarithmic size at most \(10L\).
The following applications on this same \(Y,Z\) distinguish the
points where the paper uses geometric theorems of the original proof.

\begin{enumerate}
\item On this \(Y\), Corollary~\ref{crev:good} identifies the
principal framed charts with their marked periods at each good
embedding. The identity Fourier mark kills the zero transvection
mode, and Proposition~\ref{crev:half} places the positive half
maximum on the thick part for the same self-dual row. The theta
line has Haar mean square one. Proposition~\ref{crev:mass}
computes the total flag volume and the density \(1/V\) of adelic
probability, giving the conditional Petersson measure and the
normalized Bergman bound on each component. These are deductions
on the selected good cover. The construction of such covers
with their field-size bounds is the separate existence step
specified after Corollary~\ref{crev:good}.
\src{src/Development/AutomorphicArithmetic/ThetaCorrespondence/ThetaPairLifts/HodgeBundleWeights/ExistsGoodCoverAtMinimalLevel.lean}{29}{The field-spanning family and its two metric estimates};
\src{src/ReciprocityTheory/ReductiveGroups/ArithmeticGeometry/FundamentalDomains/MaxBoundPf.lean}{12}{Cusp bound};
\src{src/Development/AbelianVarietyTheory/Proofs/DivisorProofs/FirstMinAtHalfMax.lean}{111}{Minimum at a positive half maximum on the same component};
\src{src/Development/WeilBound/RepresentationTheoryGL/TubeMeasures/DisjointTubes/MeasureSumBounds/SumMuFLe.lean}{110}{Mass in the chosen normalization}.

\item Use the rigidified cubical graph line on the semiabelian tests
of Theorem~\ref{ctwist:semi}. For the other line the relevant lattice
is the determinant ratio for the two integral Lie eigenspaces.
Theorem~\ref{cmeig:splitmodel} gives the local comparison for every
test on \(Z\), separately for each \(w,\tau\), and evaluates the
degree on its marked product CM point. For each \(w\)
outside the primes dividing \(2d_0|\Delta_E|\operatorname{lcm}(M,A)\),
the source supplies on this \(Z\) an unramified test, hence a test
taking integral values on nonzero evaluations. At those places
the integral order also follows from the single reduced fibre
in Proposition~\ref{cmeig:hyperspecial}. At the bad primes the
single rounding of the product costs at most \(\sum_{\ell\in B}\log
\ell\). The source then applies the component CM degree formula,
the key and the product to the same spanning set. For the
characteristic-zero rule spaces Corollary~\ref{crev:twist} proves
cusp compatibility in both directions for these constant-norm
units.
\src{src/Development/AbelianVarietyTheory/Definitions/Divisors/MulSmulIdemMem.lean}{264}{Exact excluded primes and integral-valued test};
\src{src/Development/AutomorphicArithmetic/ThetaCorrespondence/ThetaPairLifts/HodgeBundleWeights/UnramifiedPointsAtDeepLevel.lean}{11}{Unramified CM tests on this cover};
\src{src/Development/AutomorphicArithmetic/ThetaCorrespondence/ThetaPairLifts/HodgeBundleWeights/ShortSpanCompositionClosed.lean}{74}{Twists on the chosen component}.

\item For the Hilbert--Samuel step take the modified toroidal flag
component for \(Z\) and the two birational weights of its weight
path. Restriction in degree zero must be surjective for the
transverse flags, and the successive symbols and deepest-space
generators must be those of the same restriction. The short integral
generators and column norms of the first frame, the Bergman bound
\(\beta\), the short sections in the subsequent weights, and the
multiplication bound \(\sqrt\beta\|v\|_{\rm frame}\) use that
component's metric. The paper derives the finite-degree inequality
from these hypotheses in Corollary~\ref{newtheta:effective}.
Their simultaneous realization for the family is invoked at the
following source steps.
\src{src/Development/AutomorphicArithmetic/ThetaCorrespondence/ThetaPairLifts/HodgeBundleWeights/ExistsWeightHodgeBundleGenerated.lean}{22}{Toroidal flag and modified bundles on \(Z\)};
\src{src/Development/AbelianVarietyTheory/Proofs/DivisorProofs/HilbertSamuelFiltrationAtGoodSections.lean}{51}{Restriction and filtration};
\src{src/Development/AutomorphicArithmetic/ThetaCorrespondence/ThetaPairLifts/ThirdPlaceSeries/ShortVectorsSpanCuspLattice.lean}{33}{Short vectors at the family weights};
\src{src/Development/AutomorphicArithmetic/ThetaCorrespondence/ThetaPairLifts/ThirdPlaceSeries/HilbertSamuelBergmanFrameBound.lean}{26}{Frame and kernel bounds at the family stage};
\src{src/Development/AutomorphicArithmetic/ThetaCorrespondence/ThetaPairLifts/HodgeBundleWeights/ExistsDegreeParametersBound.lean}{33}{Combination on a common cover and component}.
\end{enumerate}

\subsection{The coefficient field and the two flag systems on a family stage}
\label{zetageo:section}

\begin{proposition}[A coefficient field for the deep level]\label{zetageo:field}
Write \(s(F)=\log[F:\mathbf Q]+\log\operatorname{rd}(F)\), where
\(\operatorname{rd}(F)=|\Delta_F|^{1/[F:\mathbf Q]}\).
At a stage as above let \(F_0\) be the initial coefficient field,
\(K=\mathbf Q(\sqrt{-d})\) the imaginary quadratic field of the
marking, and \(g_0,g_1\) the two host level factors. Suppose that
\(\Delta_K=-d\), that \(3,5\) split in \(K\), and that
\((15,g_0g_1)=1\). Put \(m=g_0g_1\,15^u\) for \(u\ge0\), \(m\ge3\).
There is a choice of a half of the split \(15^u\)-level, and a
field \(F\supset F_0K(\zeta_m)\) containing the CM \(j\)-invariants
and normalized Weber values of the \(m\)-torsion for the full order
\(\mathcal O_K\). It satisfies
\[
 s(F)\le s(F_0)+\log 2+\tfrac32\log d+3\log m.
\]
On the family with \(u=\lfloor3L/\log3\rfloor\), \(L\ge1\),
\(\log g_i\le9L+1,\ \log d\le L,\ s(F_0)\le217Z_0^2\) and
\(Z_0=frL\ge8L\), this gives \(s(F)\le391Z_0^2\).
\end{proposition}
\begin{proof}
Choose one prime over three and one over five. Equivalently, choose
a root of \(-d\) at each prime and combine the roots to a unit
\(a\bmod 15\). The ideal \(I=(15,a+\sqrt{-d})^u\) is coprime
to its conjugate and \(I\bar I=(15^u)\), by reduction at the four
distinct primes. This is the required half level.
Let \(R\) be the ray class field of \(K\) of modulus \(m\mathcal O_K\).
The ray class exact sequence gives
\([R:K]\le h_K|(\mathcal O_K/m)^\times|\le d m^2\).
Indeed each ideal class has an integral representative of norm
at most \(2\sqrt d/\pi\) by Minkowski. There are at most
\(\sum_{n\le\sqrt d}\sum_{a\mid n}1\le\sqrt d(1+\log\sqrt d)\le d\)
such ideals, by the quadratic Euler factors. The conductor of each
ray class character divides \(m\mathcal O_K\). The
conductor--discriminant formula in \(R/K\) thus gives
\(\operatorname{rd}(R)\le\sqrt d\,m\). Conjugation preserves the
modulus, so \(R/\mathbf Q\) is normal. The Artin action on
\(\zeta_m\) factors through this ray class group by the ideal norm.
The main theorem of complex multiplication in its Weber form gives
the stated CM values in \(R\)~\cite{coxCM}.

Take \(F=F_0R\). Degrees are submultiplicative under compositum, and
\(\operatorname{rd}(F)\le\operatorname{rd}(F_0)\operatorname{rd}(R)\).
For the latter inequality, base change of the different gives
\(\mathfrak D_{R/\mathbf Q}\mathcal O_F\subseteq
\mathfrak D_{F/F_0}\). Apply the tower discriminant formula and
take the \([F:\mathbf Q]\)-th root. The two bounds for \(R\) give the
claimed estimate. Since \(15\le3^3\), one has \(u\log15\le9L\) and
\(\log m\le27L+2\). The increase in \(s\) is at most \(90L\)
for \(L\ge1\). In fact \(90L\le(90/64)Z_0^2\), giving
\(s(F)<219Z_0^2\) inside the stated budget. This estimate uses
the same \(L,f,r\) for every component of the stage.
\end{proof}

Propositions~\ref{addh:constants} and~\ref{addh:level} realize the
component constants and the split half-level map on the principal
marked packet over this \(F\). The compatible charts and the integral
comparisons on modified flags enter the construction of the good cover \(Y\).
Their combination and the bounds on the selected stage are recorded
in the \src{src/Development/AbelianVarietyTheory/Proofs/DivisorProofs/ExistsGoodCoverOverSizedField.lean}{111}{construction of the marked cover}
and the \src{src/Development/AbelianVarietyTheory/Proofs/DivisorProofs/DeepExistenceStatement.lean}{29}{uniform field-size calculation}.

\begin{lemma}[One modification and the cost of changing the lattice]
\label{zetageo:lattice}
On a smooth compactified flag component over its coefficient field
suppose the two systems of Proposition~\ref{newtheta:modified} have
no base points on the marked open \(U\). There is one smooth proper
modification, an isomorphism over \(U\), on which both moving systems
are generated by their given sections. If the rational map of either
system is birational onto its image, its moving system on the
modification has that same property.

Let \(M,M'\) be full \(\mathcal O_F\)-lattices in the same space of
sections of rank \(a>0\), with Hermitian metrics at infinity. Put
\(H=M\cap M'\). At each embedding \(\sigma\), let \(d_\sigma\)
be the determinant norm of the identity from the metric on \(M\)
to that on \(M'\). For normalized Arakelov slopes,
\[
 \widehat\mu(M')-\widehat\mu(M)
 =\frac{\log|M'/H|-\log|M/H|-\sum_\sigma\log d_\sigma}
        {a[F:\mathbf Q]} .
\]
Here both embeddings over a complex place are counted.
\end{lemma}
\begin{proof}
Principalize the product of the two base ideals in characteristic
zero, with centres outside \(U\)~\cite{wlodarczyk}. On the resulting
integral variety each factor is invertible: for \(IJ\) invertible
the fractional ideal \(J(IJ)^{-1}\) is an inverse to \(I\).
Divide the pulled-back sections by a local generator of each
factor. Their coefficients now generate the unit ideal. The two
maps agree with the original rational maps on the dense open.
Proposition~\ref{newtheta:modified} then gives the adjoint
inclusion whenever its extension and cusp identifications hold.
For the slope formula, compute the degree of the determinant
of \(H\) first with each metric. Passing from \(H\) to either
superlattice adds the logarithm of its finite index in the
unnormalized degree. Changing the metric subtracts the logarithm
of the determinant norm at each embedding. Divide by the common
rank and field degree.
\end{proof}

In particular equality of the two rational section spaces determines
neither finite index in this formula. For a comparison with a
separately defined integral flag lattice, these indices and the
determinant norms must be evaluated on the same \(Y,Z\) that
produced the sections. This is one place where a pointwise proof
of \eqref{scope:missing} needs more than an archimedean norm estimate.

\begin{proposition}[A finite test for the component mean height]\label{zetageo:mean}
Fix a good cover, a weight \(u\) with the required gap, and the
finite valuation conditions used at that gap. On each component
\(Z\) let \(D\) be the flag dimension, \(\delta_Z>0\) its prescribed
top degree, \(d_F=[F:\mathbf Q]\), and \(L_{Z,m}\) its mixed integral
cusp lattice in weight \(mu\). Give it the supremum of the marked
norms over all good embeddings and frames.
Suppose the following sections belong to these actual lattices and
their products belong to the lattices of the summed weights, with
submultiplicative norms.
\begin{itemize}
\item There are \(D+1\) nonzero sections \(s_0,\ldots,s_D\) in
degree \(m_0\in8\mathbf Z_{>0}\). The ratios \(s_i/s_0\)
\((1\le i\le D)\) are algebraically independent over \(F\) and
the function field of \(Z\) has degree \(r=m_0^D\delta_Z\) over
the field of these ratios.
\item There are \(r\) nonzero sections \(b_j\) in one degree
\(m_0 k_0\) such that \(b_j/s_0^{k_0}\) form a basis for this field
extension. All sections used here have finite positive norm.
\end{itemize}
Then the mixed component volume slope satisfies
\[
 \operatorname{compVolSlope}_Z(u)\ \ge\
 a_Z:=-\frac{\sum_{i=0}^D\log\|s_i\|}{(D+1)m_0}.
\]
If these finite systems are available on every component and
\(|I|^{-1}\sum_{Z\in I}a_Z\ge-X\), they give the required component
mean with error \(X\). The degrees \(m_0,k_0\) and systems may
depend on the weight and the component.
\end{proposition}
\begin{proof}
In degree \(m=m_0(k_0+k)\) take all products \(b_j s^\alpha\) with
\(\alpha_i\ge0,\ \sum_i\alpha_i=k\). Dividing by \(s_0^{k+k_0}\)
and using the two independence assumptions proves that these products
are \(F\)-linearly independent. Multiply them by a fixed integral
basis of \(\mathcal O_F\) over \(\mathbf Z\). We obtain
\(R_k=d_F r\binom{k+D}{D}\) independent integral sections \(v_j\).
In particular \(D!R_k/(d_Fm^D)\to\delta_Z\).
On averaging over all the monomials each exponent has mean
\(k/(D+1)\), by symmetry. Thus, writing
\(\lambda_j=-m^{-1}\log\|v_j\|\), one has
\(\liminf_k R_k^{-1}\sum_j\lambda_j\ge a_Z\).
The fixed sections \(b_j\) and the coefficient basis cost \(O(1/m)\)
in this average.

For any real \(t\), form the integer combinations with coefficients
between zero and \(\lfloor e^{m(t+\lambda_j)}/R_k\rfloor\) inclusive.
They are distinct and have norm at most \(e^{mt}\). Their log count
is at least both zero and
\(mR_k(t+R_k^{-1}\sum_j\lambda_j)-R_k\log R_k\).
The mixed volume \(V_Z(t)\) takes the limsup of the log count on
degrees \(m=8k'\), multiplied by \(D!/(d_Fm^{D+1})\).
Our degrees are a subsequence and \(\log R_k=o(m)\). Hence
\[
 V_Z(t)\ge\delta_Z\max(0,t+a_Z),\qquad
 \inf_{t\in\mathbf R}\bigl(V_Z(t)/\delta_Z-t\bigr)\ge a_Z .
\]
The infimum is the component volume slope. Taking the \(a_Z\) as
witnesses gives the mean assertion. No interchange of a component
average with a limsup is needed.
\end{proof}

For the family's \src{src/Development/AbelianVarietyTheory/Definitions/Divisors/SemiposAtOfPositive.lean}{101}{mean height condition},
the lattices in this test use the height valuations at the primes
selected by the weight and gap and the recorded valuations elsewhere.
The numerical requirement is the displayed mean with
\(X=C\,\mathrm{errX}(s,u,g)\), for a single \(C\) over the admitted
stages and weights. If the \(s_i\) give a finite map to
\(\mathbf P^D\) from a model of degree \(\delta_Z\) in weight \(u\),
the projection formula gives the required degree \(r\).
A finite map also supplies the generic basis from sufficiently high
twists of its coherent direct image on projective space. Pulling that
basis back into the specified cusp spaces and placing it in the mixed
lattices are part of the stated hypotheses. Only the norms of the
\(D+1\) parameter sections enter the final mean.

There is a quantitative way to choose the parameters. Suppose \(s\)
integral sections in degree \(m_0\), of norm at most \(H\), define
a base-point-free map onto a projective variety of dimension \(D\)
and degree \(\delta'\). Its Chow form in \(D+1\) hyperplanes is
nonzero and has degree \(\delta'\) in each block. Polynomial
interpolation in each coefficient supplies hyperplanes with
coefficients in \(\{0,\ldots,\delta'\}\) where this form is nonzero.
Their common center misses the image. Projection from that center
is finite on the image: the pullback of a hyperplane is ample, so
no curve is contracted. Its degree is \(\delta'\) by the projection
formula. The resulting parameter sections have norm at most
\(s\delta'H\). If the original map has generic degree \(g\) with
\(g\delta'=m_0^D\delta_Z\), and the generic basis belongs to the
mixed lattices as above, this gives the explicit lower bound
\(-m_0^{-1}\log(s\delta'H)\) for the component volume slope.

To pass from one stage to the family, first fix a common threshold
at which the host, core-field and size bounds of
Proposition~\ref{zetageo:field} hold for every admitted point and
stage. On the resulting good covers, suppose the marked flag
identifications, integral restrictions, short frames and their
metric bounds specified above hold with uniform constants. Suppose
also the height and rank comparisons following Lemma~\ref{geo:component}
hold there, in particular the lower bound on the component average
of \(\ell_Z\), which can now be supplied by
Proposition~\ref{zetageo:mean}. Proposition~\ref{newtheta:aggregate} then gives
\(\mu_{\rm cover}\ge-2(3C_1+C_2+4)E_P\) at every such stage.
After the stated descent and Gram comparisons,
Proposition~\ref{geo:covolume} absorbs the exponential error with
one threshold depending only on the fixed constants. Only the
finitely many thresholds for uniform estimates are combined here.
The marked integral comparisons and the mean height of the actual
components remain inputs to this application. The formal proof
combines the height inputs at the
\src{src/Development/AutomorphicArithmetic/ThetaCorrespondence/ThetaPairLifts/ThirdPlaceSeries/HeightLowerBoundGeneral.lean}{21}{component height theorem}
and passes from the covers to the grid at the
\src{src/Development/AutomorphicArithmetic/ThetaCorrespondence/ThetaPairLifts/ThirdPlaceSeries/TopCovolG.lean}{49}{covolume descent}.

\subsection{Marked components over the ray class field and their good-prime lattices}
\label{addh:section}

Write \(k=\mathbf Q(\sqrt{-d})\), \(m=g_0g_1 15^u\) and
\(F=F_0R\) as in Proposition~\ref{zetageo:field}, with \(R\)
the ray class field of \(k\) modulo \(mO_k\). Use the chosen embedding
of \(F\) in \(\mathbf C\). The selected stages have \(d\ge7\), \(m\ge3\)
and, for the deep level, \(u\ge1\). Consider the principal marked PEL
packet of Lemma~\ref{munf:marks}, with trivial auxiliary mark and
indefinite signature \(p,q\ge1\).

\begin{proposition}[A rational CM point on each marked component]
\label{addh:constants}
Every geometric component of this packet contains a product CM point
defined over \(F\) with its prescribed mark. When \(15\mid m\), its
underlying abelian variety has good reduction at every finite place
of \(F\). Every component descends to a geometrically integral
\(F\)-component. The algebraic closure of \(F\) in its function
field, and in the function field of its complete Hodge flag bundle,
is \(F\). The same \(F\) works for all rows, with the size bound of
Proposition~\ref{zetageo:field}.
\end{proposition}
\begin{proof}
Choose an elliptic curve \(C_0/F\) with \(j=j(O_k)\) and fix a complex
identification with \(\mathbf C/O_k\). All its endomorphisms are
defined over \(F\). Indeed their actions on the tangent line give an
injective homomorphism in characteristic zero, and their scalars lie
in \(k\subset F\). As \(O_k^\times=\{\pm1\}\), its normalized Weber
function is a degree-two quotient by sign. It is defined over \(F\):
in short Weierstrass coordinates its formula is a nonzero rational
constant times \(g_2g_3 x/(g_2^3-27g_3^2)\), invariant under a
change of scale. Here \(j\ne0,1728\).

Let \(P\) correspond to \(1/m\bmod O_k\). Its Weber value is in
\(R\) by the second main theorem of complex multiplication
\cite[Theorem 11.39]{coxCM}. Every \(\sigma\in\operatorname{Gal}(\bar F/F)\)
therefore sends \(P\) to \(\epsilon(\sigma)P\) with a unique
\(\epsilon(\sigma)\in\{\pm1\}\). The signs define a continuous
character, since \(P\) has order \(m>2\). As \(P\) generates
\(C_0[m]\) over \(O_k/mO_k\), the action on the whole group is that
same sign. Twist \(C_0\) by this character using the automorphism
\([-1]\). On the resulting \(C/F\) the CM action still descends and
every \(m\)-torsion point is \(F\)-rational.

The integral \(j\)-invariant gives potential good reduction at every
finite place \(v\). When \(15\mid m\), choose \(l\in\{3,5\}\)
different from the residue characteristic at \(v\). Inertia has finite
image on \(T_lC\), and its image lies in \(1+lM_2(\mathbf Z_l)\)
by rationality of \(C[l]\). This kernel is torsion-free: if
\(U=1+l^aB\) with \(a\ge1\) and \(B\not\equiv0\bmod l\),
raising \(U\) to a positive exponent \(b\) prime to \(l\) preserves
its depth \(a\), and taking the \(l\)-th power raises that depth
by exactly one, by the binomial theorem for odd \(l\).
Inertia is therefore trivial. The
\href{https://wstein.org/papers/bib/Serre-Tate-Good_Reduction_of_Abelian_Varieties.pdf}
{N\'eron--Ogg--Shafarevich criterion, Serre--Tate, Theorem 1},
gives good reduction over \(F_v\).

For the given marked row, use the rational isometry of
Lemma~\ref{cmeig:row} and the axis point of
Proposition~\ref{cmeig:degree}. Its factors are
\(I_i\otimes_{O_k} C\), with the conjugated action on the
opposite-type factors. Each \(I_i\) is finite projective over
\(O_k\). Its idempotent on a power of \(C\), the multiplication
action of \(O_E\), and the polarized trace pairing used in that
proposition are all defined over \(F\). For the pairing one may
use the principal elliptic polarization to identify the dual
powers: every matrix entry in their homomorphisms is a CM
endomorphism of \(C\). The rigidified graph line of the polarization
is obtained by pullback of the normalized Poincar\'e bundle.
The product has constant \(m\)-torsion since it is a direct factor
of a product of curves with constant \(m\)-torsion. Thus the
transported mark is defined over \(F\), including marks mixing
the different factors. Good reduction of \(C\) also gives good
reduction of these factors, either by their integral idempotents
or by the same Tate-module criterion. For the tensor construction
see \cite[\S2]{serreTensor}.

In the marked period classification this point lies on the
prescribed geometric component. The classification of the
half-mark by the dual basis uses the root \(\zeta_m\in F\), and
the rigidified graph line and its inverse reconstruction commute
with base change. Consequently the fine moduli comparison preceding
Corollary~\ref{cmeig:final} descends to \(F\): after a finite Galois
extension each conjugate classifying map represents the same
functorial operation, so equals the original map by uniqueness.
Smooth geometric components are disjoint. The unique component
through this \(F\)-point is fixed by \(\operatorname{Gal}(\bar F/F)\)
and descends. For its function field \(K_Z\), geometric integrality
implies that \(K_Z\otimes_F\bar F\) is a domain. A nontrivial
finite extension of \(F\) inside \(K_Z\) would yield nontrivial
idempotents after this scalar extension, since \(F\) has
characteristic zero. The flag bundle has a dense affine cell over
\(K_Z\), so its function field is purely transcendental over \(K_Z\)
and has the same algebraic constants. This is the flag field in
\src{src/Formalism/LocalToGlobal/ProjectiveLattices/IsGenFrGenFr.lean}{33}{the generic flag variables}
and the condition used in
\src{src/Development/AbelianVarietyTheory/Proofs/DivisorProofs/ExistsGoodCoverOverSizedField.lean}{143}{the algebraic constants at the marked level}.
\end{proof}

\begin{proposition}[The split half-level map and its metrics]
\label{addh:level}
Over \(F\), start at a fine base level \(m_0=g_0g_1\ge3\) prime to
\(n=15^u\), \(u\ge1\), in the same classical PEL packet.
Write \(I\bar I=nO_E\) for the coprime halves in
Proposition~\ref{zetageo:field}, extended to \(O_E\), choosing \(I\)
as the annihilator of the new marked half. Forgetting the refinement
from level \(m_0n\) to \(m_0\) is finite \'etale and surjective.
On the full inverse image of the packet its degree is
\(\lvert\mathrm{GL}_r(O_E/I)\rvert\), where \(r=p+q\).
On each complex component above a base component, pullback preserves
the supremum norm and the Petersson \(L^2\)-norm formed with the
component probability measure.
\end{proposition}
\begin{proof}
On the fine base the universal abelian scheme has finite \'etale
\(n\)-torsion. The Chinese remainder idempotents split it as
\(A[I]\oplus A[\bar I]\). Both summands are isotropic and paired
perfectly with each other, since the polarization interchanges the
two idempotents under its adjoint. The first is \'etale locally a
free \(O_E/I\)-module of rank \(r\), as follows on geometric fibers
from the prescribed projective period lattice. A basis of this
summand and \(\zeta_n\) uniquely determine the dual mark. The
prime-to-\(n\) marking is retained: in the equation taking the
\(n\)-multiple of a refined torsion section to its old value,
multiplication by \(n\) is invertible on \(A[m_0]\).
Conversely every such basis extends the old level condition. At
each new local factor one lifts its matrix modulo \(I\) to an
invertible integral matrix and uses the inverse adjoint on the
conjugate factor. This supplies the compatible polarized Tate
markings at geometric points. The representing scheme is thus
the basis torsor under the indicated constant finite group.
The map on marks agrees with
\src{src/Development/AbelianVarietyTheory/BlockTheory/DivisorBase/HalfLatticeTypes.lean}{115}{the prescribed map of levels}.

The projection keeps the polarized abelian variety, tangent frame
and flag. The pointwise norm of a pulled-back section is unchanged.
On each complex component the projection is a finite covering
of its image component, and the invariant flag density of
Proposition~\ref{crev:mass} pulls back locally to the same density.
Integration multiplies by the sheet number, as does total component
volume. Division by that volume gives the claimed equality. The
supremum is unchanged by surjectivity.
\end{proof}

\begin{proposition}[Testing the subcanonical lattice at a good prime]
\label{addh:good-lattice}
Fix such a marked component, a place \(v\mid l\) of the splitting
coefficient field with \(l\nmid 2m\Delta_E\), and a smooth projective
toroidal model at its hyperspecial level there. Let \(\mathcal X\)
be its complete Hodge flag bundle over \(O_{F_v}\), \(B\) the
inverse image of the reduced toroidal boundary, and \(\mathcal L\)
the character line of an integral flag weight with central exponent
zero. Take a rational cusp rule whose evaluation on the generic
flag is a section \(s\) of \(\mathcal L(-B)_{F_v}\). Then \(s\)
belongs to \(H^0(\mathcal X,\mathcal L(-B))\) if and only if it
passes all the recorded integral-frame tests over \(v\) on this
component. It suffices to test abelian points over finite unramified
extensions of \(F_v\).
\end{proposition}
\begin{proof}
The smooth toroidal model with its universal semiabelian scheme
and relative normal crossings boundary is the good-reduction
model used in Proposition~\ref{cmeig:hyperspecial}. Its flag bundle
is smooth and proper and its interior is dense in each special
fiber component. The embedding idempotents of \(O_E\) split
integrally at \(v\). For a recorded test, extend its generic PEL
point to the toroidal model by properness. The connected
semiabelian extension agrees with the one in the test by
\href{https://grothendiecksga.com/read/sga7/en/9-3.html}
{SGA~7 I, Expos\'e IX, Proposition 3.2}. Its integral Lie
eigensummands therefore agree as well. The test bases give the
extension of the generic flag, since their successive spans are
direct summands over the test DVR. Evaluating an integral section
of \(\mathcal L(-B)\) in this frame first uses its inclusion into
\(\mathcal L\) and then an integral trivialization. This gives
an integral value, also at a degenerating test. The central
exponent zero removes the scalar theta generator from the value.

Conversely the regular model is normal, and \(s\), as a section of
\(\mathcal L(-B)\), is regular at every horizontal prime divisor
by the generic-fiber hypothesis. If \(s\) fails to extend, it has
order \(-a<0\) along a vertical prime divisor. The special fiber
is reduced, so for a uniformizer \(\pi\) the section \(\pi^a s\)
is a unit in a trivialization near its generic point. Choose a
closed point in this open set lying over the abelian interior.
Its residue field is finite. Smoothness and Hensel lifting give
a point over the unramified complete DVR with that residue field,
reducing to the chosen point. The eigenflags there admit adapted
integral bases. For its abelian scheme the principal polarization
line has a generator of its module of sections over this DVR:
on every fiber its space of sections has dimension one and higher
cohomology vanishes, so cohomology and base change apply.
These choices supply the framed abelian test. Its generic frame
is defined over a reduced finitely generated \(F\)-algebra, by
taking affine coordinates in the finite-type framed moduli space,
as required for evaluating a rule. The value of \(\pi^a s\) is a
unit, giving a test of negative valuation for \(s\). This proves
both assertions.
\end{proof}

Thus on a common rational cusp space with the indicated generic
realization, the recorded and toroidal subcanonical lattices
contribute zero local index in Lemma~\ref{zetageo:lattice} at
these primes. The test convention is the one in
\src{src/Formalism/LocalToGlobal/ProjectiveLattices/VAtVNeBot.lean}{53}{the semiabelian model and its integral eigenframes}.

In particular let \(\pi:Z'\to Z\) be one of the marked component
maps of Proposition~\ref{addh:level}. For a common rational cusp
space \(V\) on \(Z\) and its pullbacks on \(Z'\), the recorded
lattices over \(v\) satisfy
\[
       \pi^* V\ \cap\ L_{{\rm rec},v}(Z')
                    =\pi^*\bigl(V\cap L_{{\rm rec},v}(Z)\bigr).
\]
Indeed a deep test forgets to a base test of the same value.
Conversely, an abelian test detecting a negative valuation at the
base lifts to the chosen geometric component after a finite extension:
that component maps finitely and surjectively onto \(Z\). This extension
can be chosen unramified. The scheme of possible new marks over the
test DVR is the basis torsor of the direct summand \(A[I]\) of its
finite \'etale \(n\)-torsion, because \(l\nmid n\). Its generic
points split over a finite unramified extension. The lifted abelian
frame therefore detects the same negative valuation. This proves
integral descent on pullbacks even when \(l\) divides the degree
of the level map.

The comparison uses the unmodified flag line.
Exceptional divisors of the moving linear systems, the places
dividing \(2m\Delta_E\), and the final signed component height in
Corollary~\ref{recheight:period-degree} retain their own estimates.

\section{Fourier extraction and a uniform arithmetic inequality}\label{ar:section}

The geometric comparison enters this section through an upper bound for a scalar discrepancy. We retain its sign throughout. The reconstruction below uses only finite character orthogonality and elementary real inequalities once the discrepancy and size data are supplied. In particular, it does not estimate every character sum by its absolute value before averaging.

\subsection{The exact discrepancy}

Let \(H\) be a finite abelian group of order \(n\), let \(A_0\le H\), and put
\[
 \mathcal G=(\mathbb Z/2)^2\times H,\qquad
 \Phi=\{(0,s,h):s\in\mathbb Z/2,\ h\in H\},\qquad f=|\Phi|=2n .
\]
Write \(c(r,s,h)=(-1)^r\) and \(P_t(g)=g+(1,1,t)\). The map \(P_t\) is a bijection from \(\mathcal G\setminus\Phi\) to \(\Phi\). At a rational prime \(p\), designate its inertia subgroup by \(I_p\) and a Frobenius representative by \(F_p\). For \(\chi\in\widehat{\mathcal G}\), put
\[
 v_\chi(p^j)=|I_p|^{-1}\sum_{h\in I_p}\chi(jF_p+h).
\]
Thus \(v_\chi\) at prime powers records the Euler factor of the primitive character, with zero value when its restriction to inertia is nontrivial. Let \(G_\chi(y)\) be the ramp sum with these coefficients, centred by \(y\ell_\chi\) if \(\chi\ne1\), and by \(-y\log y+(1+\gamma)y\) if \(\chi=1\). The active host realization identifies \(\ell_\chi\) with the logarithmic derivative at one in the nontrivial case. \src{src/Formalism/SizeWindowBrackets/FirstPrimeNeSecondPrime.lean}{61}{ramp and inertia average}\quad
\src{src/Development/AutomorphicArithmetic/PrimitiveCosets/PrimRead.lean}{48}{primitive character comparison}.

\paragraph{Arithmetic signs and Fourier coordinates.}
In Proposition~\ref{host:art} the signs are ordered as \(J_q,J_d\).
Write \(\iota:H\longrightarrow H_e\) for the frame identification from
additive to multiplicative coordinates. With \(\operatorname{bit}(+1)=0\)
and \(\operatorname{bit}(-1)=1\), the map used here is
\[
  \beta(s_q,s_d,h)=
       \big(\operatorname{bit}(s_d),\operatorname{bit}(s_q),\iota^{-1}(h)\big).
\]
It is a group isomorphism, since each factor map is an isomorphism and
the two binary factors are interchanged. The three relevant images are
\[
\begin{array}{c|c|c}
 &\text{arithmetic label}&\text{Fourier label}\\ \hline
 d\text{-inertia generator}&(+1,-1,1)&(1,0,0)\\
 \text{complex conjugation}&(-1,-1,1)&(1,1,0)\\
 \text{translation by }t &(+1,+1,\iota(t))&(0,0,t).
\end{array}
\]
The arithmetic CM type is cut out by \(s_d=+1\), hence its image is
exactly \(\Phi\). In particular there is one element outside \(\Phi\)
in each coset of \(I_d\). This identifies the summand at the ramified
prime in Lemma~\ref{ar:balance}.
\src{src/Development/AutomorphicArithmetic/PrimitiveCosets/ToEBij.lean}{37}{CM type, inertia, conjugation and translations}.

Here is the additional identity used for the weights. At fixed \(x\)
the design of Section~\ref{main:assembly} assigns a geometric point
\(i(u)\) to each required translate of the common frame. At that point
the geometric CM-type labels \(\tau\) are in bijection with the
arithmetic CM type. Write \(g_{i(u)}(\tau)\) for the arithmetic label.
For each \(1\le j\le m\) set the Fourier weight to the value on the
unique geometric label with the specified Fourier coordinate, and to
one off the image. Uniqueness follows from injectivity of the label
bijection and of \(\beta\). Thus the choice of a preimage gives
\[
 w_{u,j,\beta(g_{i(u)}(\tau))}
                 = w^{\mathrm{geom}}_{i(u),j}(\tau)
                 \qquad\text{for every }\tau .
\]
The integer \(a_t\) attached to \(t\) is likewise common to all these
points. Substituting these identities term by term transports
both the positive parts and the sum of logarithms in the centre.
The source indexes the Fourier layers from zero, so its geometric
layer at index \(j-1\) is \(j\).
\src{src/Formalism/SharpRestriction/GlobalSchurBounds/TieGlobal.lean}{42}{common ranks, levels and integer coefficients},
\src{src/Formalism/SharpRestriction/GlobalSchurBounds/QuadraticLogDerivCongr.lean}{405}{construction on each matched geometric label} and
\src{src/Formalism/SharpRestriction/GlobalSchurBounds/TieGlobal.lean}{84}{pointwise weight compatibility}.

\begin{lemma}[Inertia orthogonality and centring]\label{ar:balance}
Suppose a designated set of split primes has \(F_p\in\Phi\) and \(I_p\subseteq\Phi\), and that the only other ramified prime is \(d\), with \(I_d=\{0,(1,0,0)\}\). Let \(a>0\), and let \(w_\rho>0\) for \(\rho\in\Phi\), with \(a/w_\rho\le N\). Form \(\mathrm{ASK}\) by summing the positive parts \((a/p^j-w_{P_t(g)})_+\) with coefficient \(2\log p\) on unramified inert odd powers \(g=jF_p\), and coefficient \(\log d\) for \(g\notin\Phi\) with \(g-jF_d\in I_d\). Every prime outside the designated split set and \(d\) is unramified, and such a prime is inert precisely when the first coordinate of \(F_p\) equals \(1\). The sum can be truncated to \(p^j\le N\), since \(a/w_\rho\le N\) makes every summand with \(p^j>N\) vanish. Equivalently one can bound the prime and exponent separately by \(N\), as in the source. Then
\begin{align}
 \mathrm{MAIN}
 &=a\left(\log a-1-\gamma+\ell_c-f^{-1}\sum_{\rho\in\Phi}\log w_\rho\right),\nonumber\\
 \mathrm{ASK}-\mathrm{MAIN}
 &= \frac1f\sum_{g\notin\Phi}\sum_{\chi\in\widehat{\mathcal G}}
     \overline{\chi(g)}\,w_{P_t(g)}G_\chi(a/w_{P_t(g)}).
 \label{ar:balanceeq}
\end{align}
The host in the active application supplies the split set as the prime factors of \(q\,QS\).
\end{lemma}
\begin{proof}
Interchange the inertia average and the sum over the dual group to get
\[
 f^{-1}\sum_\chi\overline{\chi(g)}v_\chi(p^j)
 =2\,|I_p|^{-1}\mathbf1_{g-jF_p\in I_p}.
\]
For unramified \(p\), the first coordinate of \(jF_p\) restricts \(j\) to the odd integers when \(p\) is inert. On the split set the indicated coset stays in \(\Phi\). At \(d\) the factor \(2/|I_d|\) is one. Multiply by \(\log p\) and the positive part, and use the prime-power support of the von Mangoldt function. The truncation at \(N\) removes only zero summands. Factoring \(w\) out of \((a/p^j-w)_+\) gives the appropriate ramp at \(a/w\). The trivial character contributes \(a\log a-a\log w-(1+\gamma)a\). For a nontrivial character the remaining term is \(-a\ell_\chi\). The character sum over \(\mathcal G\setminus\Phi\) equals \(f\) for \(1\), \(-f\) for \(c\), and zero otherwise. Finally \(P_t\) reindexes \(\log w\) over \(\Phi\).
\end{proof}
\noindent
\src{src/Development/AutomorphicArithmetic/PrimitiveCosets/hostFactsS.lean}{127}{host split set}\quad
\src{src/Development/AutomorphicArithmetic/DedekindLFunctions/LandauTheorems/ChBalance.lean}{42}{local coefficients}\quad
\src{src/Development/AutomorphicArithmetic/DedekindLFunctions/LandauTheorems/ChBalance.lean}{200}{exact balance}.

\subsection{Extraction from the signed discrepancy}

Choose \(T\subset H\), \(T_2\subset A_0\), with sizes \(n_T,n_2\), and let \(\delta=n_2/n\). At each translate \(u\in A_0\) and layer \(j\in\{1,\dots,m\}\) use weights \(w_{u,j,\rho}\). For \(t\in T\) let \(|a_t/b-x|\le1\). The thick coordinates are \((0,1,s+u)\) for \(s\in T_2\), with weights in \([b,b+\eta+1]\). Other coordinates have weights at most \(2b'\). Assume \(b,b',X'>0\), \(b'X'\le xb\), \(\eta,\theta\ge0\), and suppose the following bounds hold for every character.
\begin{align*}
 |wG_\chi(a/w)-bG_\chi(a/b)|
 &\le 3a|w-b|/\min(w,b),\\
 |G_\chi(y)|&\le y\theta \quad (y\ge X'/3),\\
 6f(x+1)(\eta+1)&\le b/x,\qquad
 2f(x+1)\theta\le 1/x .
\end{align*}
All weights satisfy the positivity and truncation hypotheses of Lemma~\ref{ar:balance}. Lemma~\ref{host:variation} proves the weight Lipschitz estimate, and Proposition~\ref{host:thin} proves the thin-scale bound at the scales of the chosen field. The last two inequalities constrain its geometric parameters. \src{src/Formalism/SizeWindowBrackets/GrowthBoundGivesRiemannHypothesis.lean}{55}{Thin scales and Lipschitz}\quad
\src{src/Formalism/SizeWindowBrackets/ChainSlopeWindows/LipschitzWeightAndMinimum.lean}{957}{weight Lipschitz proof}\quad
\src{src/Formalism/SizeWindowBrackets/SChain.lean}{50}{cost hypotheses}.

Let \(\Psi\) be the characters annihilating \(\{(0,0,h):h\in A_0\}\), let \(o=[H:A_0]\), and let \(K\subset\Psi\) contain \(1\), with \(\chi(1,0,0)=1\) on \(K\). Given \(c_0=(3o+1)^{-1}\), set
\[
 \omega_\chi=\begin{cases}1&\chi|_H=1,\\c_0&\chi|_H\ne1,\end{cases}\qquad
 G_K=\sum_{\chi\in K}\omega_\chi\Re G_\chi(x),\quad
 \Sigma=\sum_{\chi\in\Psi\setminus K}|G_\chi(x)|.
\]
Suppose the counts of \(T\) in zero and nonzero cosets of \(H/A_0\) differ from \(4c_0 n_T\) and \(3c_0 n_T\), respectively, by at most \(n_Tx^{-1+\varepsilon}\). For fixed parameters, assume that the logarithmic derivatives are bounded on \(K\setminus\{1\}\) by a fixed constant, and by a fixed multiple of \(\log x\) on \(\Psi\setminus K\). In the active chapter \(o\) is fixed by \(\Lambda\). The scale derivative bounds used below are \(|G_1(y)-G_1(x)|\le C_d|y-x|\), and, when \(\chi\ne1\), \(|G_\chi(y)-G_\chi(x)|\le|y-x|(\log\max(y,x)+C_d+|\ell_\chi|)\) for \(x,y\ge1\). \src{src/Development/AutomorphicArithmetic/PrimitiveCosets/TopsSum.lean}{75}{coset shares}\quad
\src{src/Development/AutomorphicArithmetic/PrimitiveCosets/charsA.lean}{80}{character data}\quad
\src{src/Formalism/SizeWindowBrackets/ChainSlopeWindows/LipschitzWeightAndMinimum.lean}{966}{scale derivative estimate}.

\begin{lemma}[Extracted inequality]\label{ar:extract}
Assume the foregoing conditions and \(x\) sufficiently large in terms of the fixed constants and \(\varepsilon>0\). If for every \(u\in A_0\) the real discrepancy in \eqref{ar:balanceeq}, summed over \(t,j\), is at most \(E\), then
\[
 (n_Tmb)\frac{\delta}{2}\bigl(G_K-\Sigma-4x^{1/2+\varepsilon}\bigr)
 \le E+2n_Tmb/x.
 \tag{*}\label{ar:scaled}
\]
\end{lemma}
\begin{proof}
In the thick terms replace \(w\) by \(b\). Each term changes by at most \(3(x+1)(\eta+1)\). In the other terms \(a_t/w\ge (x-1)b/(2b')\ge X'/3\) for \(x\ge3\), so deleting them costs at most \((x+1)b\theta\) per term. The numerator has at most \(f\) classes and \(2f\) characters, with prefactor \(1/f\). The two cost hypotheses therefore give total error at most \(2b/x\) per \(t,j\).

Put \(e_1(h)=(1,0,h)\). The preimage of a thick coordinate is \(e_1(s+u-t)\). For \(s\in A_0\),
\[
 |A_0|^{-1}\sum_{u\in A_0}\chi(-e_1(s+u-t))
 =\mathbf1_{\chi\in\Psi}\chi(e_1(t)).
\]
Thus the thick sum averaged over \(u\) equals
\(\frac{mb\delta}{2}\sum_{t\in T}\Re\sum_{\chi\in\Psi}\chi(e_1(t))G_\chi(a_t/b)\).
In the average over \(T\), the regular coset contribution cancels for a nontrivial quotient character, and the excess share \(c_0\) at the zero coset survives. More precisely,
\(\left|\sum_{t\in T}\chi(0,0,t)-\omega_\chi n_T\right|\le o n_T x^{-1+\varepsilon}\)
on \(K\), with exact equality of main terms when the restriction is trivial. For each nontrivial character in \(K\), moving \(a_t/b\) to \(x\) costs \(O(\log x)\) per term by the scale derivative, while the phase discrepancy costs at most \(O(x^{-1+\varepsilon})|G_\chi(x)|=O(x^\varepsilon\log x)\) per term. Here integration from \(1\) gives \(|G_\chi(x)|\le x(\log x+C_d+|\ell_\chi|)\). The trivial character has bounded move cost. For an excluded character, retain the lower bound \(-|G_\chi(x)|\) and subtract its additional \(O(\log x)\) move cost. Since \(|\Psi|=4o\) is fixed, their combined move and phase errors per \(t\) are at most \(4x^{1/2+\varepsilon}\) for large \(x\). For a threshold uniform in the supplied data, one may use \(\log(x+1)\le5x^{1/4}\) and absorb the fixed coefficient times \(|\Psi|\) into another \(x^{1/4}\), as in the source. Summing gives \eqref{ar:scaled}.
\end{proof}
\noindent
\src{src/Development/AutomorphicArithmetic/DedekindLFunctions/LandauTheorems/ChCore.lean}{56}{annihilator projection}\quad
\src{src/Development/AutomorphicArithmetic/DedekindLFunctions/LandauTheorems/ChCore.lean}{75}{thick and thin error}\quad
\src{src/Development/AutomorphicArithmetic/DedekindLFunctions/LandauTheorems/ChCore.lean}{202}{one-character estimate}\quad
\src{src/Development/AutomorphicArithmetic/DedekindLFunctions/LandauTheorems/ChCore.lean}{343}{uniform comparison}.

\begin{proposition}[Four explicit losses]\label{ar:four}
Let \(x\ge3,\varepsilon>0,C\ge0,b>0,a_E\ge0,n>0\), and assume
\[
 \begin{gathered}a_E\le2xb,\qquad x^{1/2-\varepsilon}\le m\le x^{1/2+\varepsilon/2},\\
 xm\le k\le n^2x^{1/2-\varepsilon},\quad
 n\le2n_Tx^{\varepsilon/8},\quad x^{-\varepsilon}\le n_2/n\le1 .\end{gathered}
\]
Suppose \eqref{ar:scaled} holds and
\[
 E\le x^{C\varepsilon}
 \left(n_Ta_E+n_Ta_E\,m^2/k+2nm^2b\,x^\varepsilon+
                         mb\,n_2 k x^{2\varepsilon}/n^2\right).
\]
If \(x^{3\varepsilon/8}\ge26\), then \(G_K\le x^{1/2+(C+3)\varepsilon}+\Sigma\).
\end{proposition}
\begin{proof}
Write \(B=n_Tmb>0\) and \(U=x^{1/2+13\varepsilon/8}\). Divide the four bracket terms by \(B\). They are bounded, respectively, by
\[
 2x/m\le2U,\quad 2xm/k\le2,\quad
 2(n/n_T)mx^\varepsilon\le4U,\quad
 (n_2/n_T)x^{1/2+\varepsilon}\le2x^{1/2+9\varepsilon/8}\le2U .
\]
The fourth bound uses both \(n_2\le n\) and the upper bound for \(k\). Hence \(E/B\le10x^{C\varepsilon}U\). By \eqref{ar:scaled}, after multiplying by two the right side is at most \(22x^{C\varepsilon}U\). Divide by \(\delta\), using \(\delta^{-1}\le x^\varepsilon\), whenever \(G_K-\Sigma-4x^{1/2+\varepsilon}>0\). Otherwise the desired intermediate bound is immediate. Restoring the move error yields \(G_K-\Sigma\le26x^{1/2+(C+21/8)\varepsilon}\), and the numerical threshold absorbs 26.
\end{proof}
\noindent
\src{src/Formalism/SizeWindowBrackets/ErrorFormsWeakWindowsFalse.lean}{133}{four-term bracket}\quad
\src{src/Formalism/SizeWindowBrackets/ChainSlopeWindows/LipschitzWeightAndMinimum.lean}{507}{explicit numerical bound}.

\begin{proposition}[The translate mean and its four primitive twists]\label{ar:mean}
Retain the host and scale assumptions of Lemma~\ref{ar:extract},
including uniform derivative constants. Let \(D_u\) be its summed
real discrepancy and set
\[
 \bar D=|A_0|^{-1}\sum_{u\in A_0}D_u,\qquad
 B=n_Tmb,\qquad \delta=n_2/n>0 .
\]
Then \eqref{ar:scaled} also holds with \(E=\bar D\), without an
upper bound on any individual translate. More precisely, for the
cyclotomic host of Proposition~\ref{host:art}, write
\(G[\xi]=G_{\xi^\ast}(x)\), where \(\xi^\ast\) is the primitive
inducer, and put \(\omega_0=1,\ \omega_a=(3o+1)^{-1}\) for
\(1\le a<o\). Then
\[
\begin{aligned}
 \frac{2\bar D}{B\delta}
  ={}&\sum_{a=0}^{o-1}\omega_a\Re\big(
          G[\psi^a]+G[\psi^aJ_q]-G[\psi^aJ_d]
                        -G[\psi^aJ_qJ_d]\big)+R_x,\\
 |R_x|\le{}&C(1+x^\varepsilon)\log(2x)+\frac{4}{x\delta}.
\end{aligned}
\]
The constant depends only on the fixed character order and the
derivative bounds. It is independent of the window, the field
degree, and the number of translates.
\end{proposition}
\begin{proof}
In the proof of Lemma~\ref{ar:extract} the thick replacement and
thin deletion change each \(D_u\) by at most \(2B/x\) in absolute
value. The averaged thick expression displayed there, divided by
\(B\delta/2\), is exactly
\[
 \frac{1}{n_T}\sum_{t\in T}\Re\sum_{\chi\in\Psi}
                         \chi(1,0,t)\,G_\chi(a_t/b).
\]
All subsequent lower bounds in that proof apply directly to
\(\bar D\), which proves the first assertion. For the two-sided
formula replace \(a_t/b\) by \(x\). The derivative estimates and
\(|a_t/b-x|\le1\) give a total error \(O_o(\log(2x))\), since there
are \(4o\) characters and their derivatives at one satisfy the
stated bounds. For each of them, average its restriction to \(H\)
over \(T\). The average is exactly one for trivial restriction.
Otherwise the sum over the regular cosets vanishes and the excess
share at zero is \((3o+1)^{-1}\), with error at most
\(o\,x^{-1+\varepsilon}\). In this case the Dirichlet character is
nonprincipal and integration of the scale derivative gives
\(|G_\chi(x)|\ll x\log(2x)\) with the same uniform constants.
The error from the coset shares is therefore
\(O_o(x^\varepsilon\log(2x))\).

Finally the characters on this quotient have primitive inducers
of \(\psi^a J_q^iJ_d^j\), with \(0\le a<o\), \(0\le i,j\le1\).
Restriction to \(H\) is trivial exactly for \(a=0\). The factor
\(\chi(1,0,0)=(-1)^j\) is evaluation on \(d\)-inertia by the
coordinate computation above. These give the asserted signs and
weights. The initial error \(2B/x\) becomes \(4/(x\delta)\).
\end{proof}

\begin{corollary}[Size of the discarded part of the mean]\label{ar:meanhost}
On the hosts in Formal input~\ref{red:design}, the remainder in
Proposition~\ref{ar:mean} satisfies
\[
 \left|\frac{2\bar D}{B\delta}-G_{\mathcal F_\psi}(x)-R_x\right|
 \le \frac{12o}{3o+1}\,x^{1/2+\varepsilon}.
\]
In particular, at every sufficiently large such scale,
\(G_{\mathcal F_\psi}(x)\le x^{1/2+4\varepsilon}\) already implies
the mean version of \eqref{scope:missing}. Conversely the mean
version and the remaining conditions of
Theorem~\ref{main:conditional} give \eqref{eq:window}.
\end{corollary}
\begin{proof}
For each retained power there are three discarded terms in the
signed formula. Each has modulus at most \(x^{1/2+\varepsilon}\)
by the choice of the two primes. Summing their weights gives
\(3(1+(o-1)/(3o+1))=12o/(3o+1)<4\). At the actual scales
\(a_E\ge xb\), \(m\le x^{1/2+\varepsilon/2}\) and
\(x^{-\varepsilon}\le\delta\le1\). If \(E_x\) denotes the right
side of \eqref{scope:missing}, its first nonnegative term gives
\[
   \frac{2E_x}{B\delta}\ge
      \frac{2x^{5\varepsilon}a_E}{mb\delta}
           \ge 2x^{1/2+9\varepsilon/2}.
\]
Under the indicated upper bound on the retained sum, the bound
just proved and the estimate for \(R_x\) put \(2\bar D/(B\delta)\)
below this threshold for large \(x\). In the other direction use
Propositions~\ref{ar:mean} and~\ref{ar:four} as in
Section~\ref{scope:account}. Both implications concern the same
data at \(x\). The design assumption used for the numerical
lower bound \(a_E\ge xb\) is in
\src{src/Formalism/SizeWindowBrackets/ErrorFormsWeakWindowsFalse.lean}{30}{the scale restrictions}.
\end{proof}
For \(\psi=1\) the formula has the particularly simple main term
\(G[1]+G[J_q]-G[J_d]-G[J_qJ_d]\); all four terms are real.
The evaluation therefore retains the zeta ramp itself. Obtaining the
required upper bound on the mean by arithmetic means remains an input.

\begin{corollary}[Testing at squares]\label{ar:squares}
For the critical-line implication it suffices to provide the
finite-character, rank, scale and excluded-character hypotheses
of Theorem~\ref{main:conditional}, together with \eqref{scope:average},
at the integer squares in each required window. Conversely, suppose only
the character and scale conditions of Corollary~\ref{ar:meanhost}
are given at every real scale in these windows. If the retained
sum satisfies \(G_{\mathcal F_\psi}(s)\le s^{1/2+3\varepsilon}\)
at their integer squares, then \eqref{scope:average} holds for
the specified data at every real scale of every sufficiently late
such window.
\end{corollary}
\begin{proof}
The constant terms defining \(G_{\mathcal F_\psi}\) depend only on
the fixed powers of \(\psi\). Each ramp \((y/n-1)_+\) is continuous
and piecewise linear, of slope at most \(1/n\). On \([x/2,2x]\),
for \(x\ge 4\), the absolute slope bound for their sum, including
the centring in \eqref{eq:centred}, is therefore
\[
 C_\psi\left(1+\log(2x)+
                     \sum_{2\le n\le 2x}\frac{\log n}{n}\right)
       \le C'_\psi\log^2(2x).
\]
There are finitely many breaks on this interval, so this is also
a Lipschitz bound across them. For every \(x\in[Y,Y^{3/2}]\)
choose \(s=\lfloor\sqrt{x}\rfloor^2\) if this lies in the window,
and the next square otherwise. In the latter case the first square
at least \(Y\) is at most \((\sqrt Y+1)^2\le Y^{3/2}\) for large
\(Y\). Thus \(|s-x|\le 2\sqrt{x}+1\), and
\[
 |G_{\mathcal F_\psi}(s)-G_{\mathcal F_\psi}(x)|
                    \ll_\psi\sqrt{x}\log^2(2x).
\]
For any fixed \(h>0\), an upper bound \(s^{1/2+h}\) at these
squares now gives \(x^{1/2+h+\varepsilon}\) on the whole window
after a uniform increase in its lower threshold. In the first
assertion Corollary~\ref{ar:meanhost} gives \(h=9\varepsilon\)
at the squares. Theorem~\ref{an:zero} applies because
\(\varepsilon\) can be taken arbitrarily small. For the converse
take \(h=3\varepsilon\) and apply Corollary~\ref{ar:meanhost}
to the specified host at each real \(x\). Only the fixed retained
function is interpolated; the hosts and their weights may change
between squares.
\end{proof}

\subsection{Simultaneous transport and the exponent \(9\varepsilon\)}

For fixed \(\Lambda\) and parameters, the active theorem \(\mathrm{core\_gm}\) first chooses \(x_1\), then quantifies over all \(x\ge x_1\), frames, and hosts with compatible \(\mathrm{EAtH}\) and \(\mathrm{CoreAtH}\) data. Its bound is
\begin{equation}\label{ar:active}
 \Lambda.Gm(x)\le x^{1/2+(e_D+e_L+7)\varepsilon}.
\end{equation}
Here the comparison for each translate is \(\mathrm{ch\_top}\) applied using the coverage map of \(\mathrm{CoreAtH}\). The latter ensures that the arithmetic discrepancy in \eqref{ar:balanceeq} and the geometric data share the same \(a_t\), \(w_{u,j,\rho}\), and the same centring term. The source defines \(C_{\mathrm{err}}=C.eDiv+C.eLoss+3\) and proves \(C_{\mathrm{err}}\le e_D+e_L+3\), so Proposition~\ref{ar:four} uses \(C=e_D+e_L+3\). The precise reindexing of retained characters gives \(G_K=\Lambda.Gm\). For excluded characters the host's separate twist estimate gives \(\Sigma\le4o\,x^{1/2+\varepsilon}\). Taking \(x^\varepsilon\ge1+4o\) absorbs this term and gives \eqref{ar:active}. The all-character bounds of the final theorem have not been used. In particular the bounds for the excluded twists are inputs, whereas the retained sum is the output. \src{src/Development/AutomorphicArithmetic/DedekindLFunctions/LandauTheorems/ChTop.lean}{37}{geometric discrepancy}\quad
\src{src/Development/AutomorphicArithmetic/PrimitiveCosets/GmAEq.lean}{61}{retained terms}\quad
\src{src/Development/AutomorphicArithmetic/PrimitiveCosets/SigALe.lean}{38}{excluded twists}\quad
\src{src/Development/AutomorphicArithmetic/PrimitiveCosets/CoreGm.lean}{12}{pointwise assembled estimate}.

The active family construction is simultaneous in a specific order. A closure condition supplies a positive \(\varepsilon\)-threshold and, for each admissible \(\varepsilon\), a size requirement on a family. Its conjunction with a second condition uses the smaller threshold and the supremum of the size requirements. Only then does \(\mathrm{topDesWin}\) invoke design existence, selecting a single \(F_m\) with \(\mathrm{ChosenE}\), \(\mathrm{TuSize}\), and \(\mathrm{DGamSize}\). It uses \(\mathrm{LinkN}\), \(\mathrm{SmallLevelRankEq}\), and \(\mathrm{HypV}\) to obtain \(\mathrm{DIVinN}\), combines this with the division and rank data, and fits the five hypotheses of \(\mathrm{LinkGLOBAL}\) to that same family. The input list is \(\mathrm{LinkG}\), \(\mathrm{LinkS}\), \(\mathrm{LinkIRR}\), \(\mathrm{LinkDIV}\), and \(\mathrm{LinkMu}\). The output is one choice of constants, followed by \(\mathrm{ChainTied}\) for every sufficiently large \(x\) in \(F_m.\mathrm{Win}\). \src{src/Formalism/SharpRestriction/ParamsHeightRankConstant.lean}{160}{closure under conjunction}\quad
\src{src/Development/AutomorphicArithmetic/CharacterBoxes/DesWin.lean}{11}{single design choice}\quad
\src{src/Formalism/SharpRestriction/GlobalTheoremInputs.lean}{158}{global input and output}.

At each such \(x\), the witness contains a map from the entire translate group into the family's point index. The frame, \(m,f,q,d\), and vertex coding are tied across translates, while the local lattices are read from the corresponding point. The common \(T_2\) and the weights are tied by one arithmetic encoding. This is the reason that the pointwise \(\mathrm{ch\_top}\) estimates in Lemma~\ref{ar:extract} refer to the same Fourier data. The source instantiates \((C_{\rm abs},C_0,e_D,e_L)=(2,3,2,0)\), so \eqref{ar:active} reads \(x^{1/2+9\varepsilon}\) on the chosen family's good scales. The statement allows a different family for another \(\varepsilon\). It requires arbitrarily late complete windows for the family chosen at the current \(\varepsilon\). \src{src/Formalism/SharpRestriction/GlobalSchurBounds/TieGlobal.lean}{97}{translate witnesses and identities}\quad
\src{src/Development/AutomorphicArithmetic/ThetaCorrespondence/ThetaPairLifts/ThirdPlaceSeries/TopShape.lean}{52}{active choice of constants}\quad
\src{src/Development/AutomorphicArithmetic/PrimitiveCosets/WindowInequalityOfArtinInputs.lean}{17}{transport onto good scales}\quad
\src{src/Formalism/SharpRestriction/RamificationFrobenius/TotientGaloisClasses/TopWin.lean}{68}{choice of small exponent}.


\section{Assembly on one arithmetic family}\label{main:assembly}

Fix \(\psi\), and let \(\mathcal F_\psi\) be the cyclic family from Section~\ref{intro:main}. A host at scale \(x\) has a field \(E\), a retained subset of characters, a rank \(t\), and a positive scale \(b\). The source requires on that \emph{same} host both the estimates for the discarded characters with constant \(4\) and a bound \(\sqrt b\) on the cohomological length budget for \((E,t)\). The class of a family requires complete windows and the condition at every index of every sufficiently large scale. \src{src/Development/AutomorphicArithmetic/CharacterBoxes/Supply.lean}{22}{window-family class} and \src{src/Development/AutomorphicArithmetic/CharacterBoxes/Supply.lean}{42}{cyclic host request}.

\begin{lemma}[Simultaneous choice]\label{red:simultaneous}
Let a finite collection of family predicates \(P_j(F,H)\) have the following property. For each \(j\) there is \(\epsilon_j>0\) such that at each \(0<\epsilon\le\epsilon_j\) one can choose a size requirement \(K_j\) ensuring \(P_j(F,H)\) for every admissible pair satisfying that requirement. Suppose requirements have finite joins, with satisfaction of a join implying satisfaction of each summand. If, for all sufficiently small \(\epsilon>0\) and every \(K\), a design theorem produces an admissible pair \((F,H)\) satisfying \(K\) and a property \(Q\), then at every such \(\epsilon\) there is one pair satisfying \(Q\) and all the \(P_j\).
\end{lemma}
\begin{proof}
Reduce the range of \(\epsilon\) by taking the minimum of the finitely many positive thresholds and the design threshold. Fix \(\epsilon\) first. Take the join of the requirements furnished at that \(\epsilon\), apply the design theorem once, and evaluate each universal assertion at its output. The hypotheses on \(K\) are exactly those of \(\code{FamP.Large}\). Its three effective conditions compare lower bounds on a first collection of constants, thresholds depending on preceding choices, and thresholds depending on the accumulated choices. The join is the pointwise maximum at each stage. The identity \(\operatorname{Large}(K\vee K')\Leftrightarrow\operatorname{Large}(K)\wedge\operatorname{Large}(K')\) follows componentwise from the defining inequalities. \src{src/Formalism/SharpRestriction/ParamsHeightRankConstant.lean}{147}{join of requirements}, \src{src/Formalism/SharpRestriction/ParamsHeightRankConstant.lean}{163}{conjunction of uniform predicates} and \src{src/Development/AutomorphicArithmetic/CharacterBoxes/DesWin.lean}{24}{application to the selected family}.
\end{proof}

\begin{formalinput}[The family selected for a primitive character]\label{red:design}
Fix a primitive Dirichlet character \(\psi\), and take the window
power \(3/2\). The following is the specialization of the original
family construction and geometric comparison used in this paper.
There is an \(\varepsilon_0>0\) such that, for
\(0<\varepsilon<\varepsilon_0\) and for every staged lower size
requirement of Lemma~\ref{red:simultaneous}, there are one
family \(F\), a nonnegative height exponent \(C_H\), and a threshold
\(x_0\) with the following properties.

\begin{enumerate}
\item The family meets the specified size requirement. Its set of
good scales contains \([Y,Y^{3/2}]\) for arbitrarily large \(Y\).
At each sufficiently large scale, every point of the family carries
a host \((E,t,b)\) realizing \(\mathcal F_\psi\) with its prescribed
weights and the coset shares of Section~\ref{ar:section}.
The discarded primitive twists there obey
\(|G_\chi(x)|\le x^{1/2+\varepsilon}\) and
\(|L'(1,\chi)/L(1,\chi)|\le4\log x\).
On the same host, for every choice of algebraic closure and local
coefficient data, the cohomological lengths
\(h_{\ell,e}\) of Section~\ref{sup:length} are finite for
\(1\le e<t\), and for every finite set \(S\) of primes,
\[
        \sum_{\ell\in S}\log\ell\sum_{1\le e<t}h_{\ell,e}
                \le\sqrt b .
\]
\item At each good \(x\ge x_0\) one has a common frame, primes \(q,d\),
sets \(T,T_2\), integers \(m,f,k\), positive scales as in
Section~\ref{ar:section}, a base point, and a point \(i(u)\) in
the family for each translate. All the scale, share, derivative
and rank conditions of Lemmas~\ref{ar:balance}, \ref{ar:extract}
and Proposition~\ref{ar:four} hold. There is a common array of
integers \(a_t\). The local lattices at \(i(u)\) have the weights
\(w_{u,j,\rho}\) under the coordinate map \(\beta\).
At that same point the geometric prime sum and its centre are
the sums of \(\mathrm{ASK}_{u,t,j}\) and
\(\mathrm{MAIN}_{u,t,j}\) over \(t,j\).
\item Write \(B_0,N_c\) for the geometric budgets on this frame and
\(\mathcal A_u,\mathcal M_u\) for those last two sums. The
divisibility, slope and degree comparisons give
\[
\begin{aligned}
 \mathcal A_u-\mathcal M_u
 &\le x^{(e_{\rm div}+e_{\rm loss}+3)\varepsilon}
       \big[n_T a_E(1+m^2/k)+m(B_0+N_c)\big],\\
 B_0&\le2nmb\,x^\varepsilon,\qquad
 N_c\le n_2 k b\,x^{2\varepsilon}/n^2 ,
\end{aligned}
\]
for every required \(u\), with \(e_{\rm div}\le2,\ e_{\rm loss}\le0\).
The pointwise comparison uses a threshold common to all these data
after \(F\) and its constants have been chosen.
\end{enumerate}
\end{formalinput}

Here are the component results used in this statement. The design
first supplies the cohomological and character conditions as a single
request on each host. Its additional choices control the sizes of
the thick sets, the auxiliary primes, the thin-character errors,
the rank ratios and the number of row words. For the common scale
\(\Gamma\), each row word \(v\) has
\(1/3\le d_v\Gamma^2/D_*\le3\). The number of row words is at least
\(3\Gamma^2x^{C_H}\). The level primes \(T_u\) and the common
stabilizer denominators of Section~\ref{geo:engine} obey the size
bounds of the same design.
\src{src/Development/AutomorphicArithmetic/CharacterBoxes/ExistsPageN.lean}{12}{cohomological budget on all local data},
\src{src/Formalism/SharpRestriction/ParamsHeightRankConstant.lean}{187}{quantifiers of the design},
\src{src/Formalism/SharpRestriction/GlobalSchurBounds/Chosen.lean}{83}{seven simultaneous choices} and
\src{src/Development/AutomorphicArithmetic/ThetaCorrespondence/ThetaPairLifts/ThirdPlaceSeries/TopShape.lean}{52}{specialization to the cyclic request}.

One can check the interface between cohomology and local division
numerically. At a fixed point write \(L=x^\varepsilon\), so that
\(B_0=fmbL\). Let \(H_j\) be the sum of the local error lengths
at the noncentral places in layer \(j\), with their logarithmic norms
and ramification factors. Let \(d_j\) be the small-level rank and
\(D_0\) the rank of the common row space. The local length estimate
and the rank comparison at this point give
\[
 H_j\le d_j\{C(fmbL^2+(\log x)^2N_c)+V\},\qquad d_j=D_0,
\]
where \(V=\sum_\ell\log\ell\sum_{1\le e<t}h_{\ell,e}\).
Taking the supremum of the finite subsums in part 1 gives
\(V\le\sqrt b\). Here \(b\ge1,\ f,m\ge1,\ N_c\ge0\), and eventually
\((\log x)^2\le L\). Hence \(V\le b\le B_0\) and
\[
                  H_j\le D_0(|C|+1)L(B_0+N_c).
\]
Together with coverage of the local twisted multiplicities this
is the input for local division. Full rank and the comparison of
mean slopes convert the divisor estimate into the slope input.
The prime gains, the Gram matrix inequalities, the irreducibility
loss, local division and this slope input are the five premises of
the global comparison. Theorem~\ref{geo:hn} explains the elimination
of the intervening slope. At a fixed translate the two resulting
geometric bounds give part 3, with \(e_{\rm div},e_{\rm loss}\)
the exponents for division and loss. Thus part 3 describes the
precise strength of the cited geometric result used here.
\src{src/Formalism/SharpRestriction/GlobalTheoremInputs.lean}{40}{local length estimate},
\src{src/Development/AutomorphicArithmetic/CharacterBoxes/DivinNOfHypV.lean}{76}{division at one point},
\src{src/Formalism/SharpRestriction/GlobalTheoremInputs.lean}{158}{five global inputs},
\src{src/Formalism/SizeWindowBrackets/ChainSlopeWindows/LevelOfCongruence.lean}{175}{elimination of the slope},
\src{src/Formalism/SizeWindowBrackets/SChain.lean}{33}{the two geometric budgets} and
\src{src/Development/AutomorphicArithmetic/DedekindLFunctions/LandauTheorems/ChTop.lean}{42}{pointwise discrepancy on the selected data}.

\begin{proof}[Derivation of Theorem~\ref{main:window} using the cited arithmetic inputs]
The closed representation theorem of Section~\ref{tr:section} supplies a row of local--global compatibility data. In the final assembly the \emph{same} proof \(h\) occurs twice in the literature record, together with the local correspondence existence theorem, and gives the window input through \(\code{vwin\_of\_row}\). The constructor of the record is \(\langle h,\mathrm{recTExists},h\rangle\). It is important that the row is instantiated. A derivation starting only with the parameterized theorem taking \(h\) as an argument would still have this entire input to prove. The separate theta matching and covolume arguments are explicitly passed at the next layer. \src{src/Development/AutomorphicArithmetic/Foundations/GaloisRepresentations/IntegralStructures/RowSrc.lean}{6}{closed representation input}, \src{src/Development/AutomorphicArithmetic/CharacterBoxes/MainConclusion.lean}{7}{use and discharge of the row} and \src{src/Development/AutomorphicArithmetic/ThetaCorrespondence/ThetaPairLifts/ThirdPlaceSeries/EveryCharacterOfLiteratureOfVWin.lean}{810}{theta and height specializations}.

Apply Section~\ref{sup:section}'s host-existence input for the request of \(\psi\), with \(\lambda=3/2\). Apply Lemma~\ref{red:simultaneous} to the nine closed properties taken together in the design theorem. They control the chain length data, local division, local errors, irreducibility, small-level ranks, rank comparisons, prime gains, Gram estimates and the global comparison. For the selected pair, the bounds on the auxiliary size and the common stabilizer denominator are provided by the design. Its cohomological budget and small-level rank identity enter the division input. The nonnegative prime-length terms have all finite subsums at most \(\sqrt b\), so their sum is at most \(\sqrt b\le b\) for \(b\ge1\). The local length estimate then proves \(\code{DIVinN}\). Joining it with the error and division bounds produces the divisor estimate on the family. The rank comparison yields its slope bound. The global implication takes these estimates together with the same family's Gram, gain and irreducibility data and yields the tied comparison. These specialized arithmetic implications are inputs proved in the pinned source. Section~\ref{geo:engine} gives their lattice and height calculations, and Section~\ref{ar:section} identifies the signed expression to which they apply. \src{src/Development/AutomorphicArithmetic/CharacterBoxes/DivinNOfHypV.lean}{117}{cohomology into division}, \src{src/Development/AutomorphicArithmetic/CharacterBoxes/DesWin.lean}{29}{five inputs to the global comparison} and \src{src/Development/AutomorphicArithmetic/ThetaCorrespondence/ThetaPairLifts/ThirdPlaceSeries/TopShape.lean}{52}{numeric specialization}.

For each sufficiently small \(\varepsilon>0\) we have thus obtained a single family \(F\) and constants such that (i) \(F\) has arbitrarily late windows of power length \(\lambda\), (ii) every sufficiently large point of its good set has tied geometry at every required translate, and (iii) its host satisfies the character bounds. At each such point, the transport to the pointwise theorem identifies the retained sets by injectivity of their common frame coordinates. Equations~\eqref{ar:balanceeq} and \eqref{ar:active} give an \(x_1\) such that
\[
  x\in\operatorname{Win}(F),\quad x\ge x_1
       \quad\Longrightarrow\quad
 G_{\mathcal F_\psi}(x)\le x^{1/2+(e_D+e_L+7)\varepsilon}.
\]
The bound holds for every point in that set, since the threshold in the pointwise comparison precedes the choice of a frame there. At the actual specialization \(e_D=2,e_L=0\). For any \(Y_0\) take a family window starting after \(\max(1,x_1,Y_0)\). Every \(x\) in \([Y,Y^{3/2}]\) belongs to its good set and exceeds the threshold, so this is \eqref{eq:window}. The constant \(9\) thus measures the loss in the exponent per fixed \(\varepsilon\), and the starting point may depend on \(\varepsilon\). \src{src/Development/AutomorphicArithmetic/PrimitiveCosets/LevelGmLeOnFamilyScales.lean}{13}{transport of the tied data}, \src{src/Development/AutomorphicArithmetic/PrimitiveCosets/CoreGm.lean}{55}{application at every translate}, \src{src/Development/AutomorphicArithmetic/CharacterBoxes/ZeroFreeAllOfTopG.lean}{26}{specialization to each primitive character} and \src{src/Development/AutomorphicArithmetic/ThetaCorrespondence/ThetaPairLifts/ThirdPlaceSeries/TopPlug.lean}{14}{window exponent}.
\end{proof}

\subsection{Using the translate mean in the window construction}\label{red:translate}

Fix the character, \(\varepsilon\), the common host at \(x\), and the
weights and dimensions in part 2 of Formal input~\ref{red:design}.
Write \(\mathcal A_u=\sum_{t,j}\mathrm{ASK}_{u,t,j}\),
\(\mathcal M_u=\sum_{t,j}\mathrm{MAIN}_{u,t,j}\),
\(D_u=\mathcal A_u-\mathcal M_u\). The expectation throughout this
subsection is uniform on the subgroup \(A_0\) in
Proposition~\ref{ar:mean}. The original design \code{topDesWin} selects
one family. Its later application in \code{ch\_top} bounds each tied translate. To apply a finite mean lemma, first restrict the tied data on the host group to \(A_0\).
The source file also gives a
\src{src/Formalism/SizeWindowBrackets/ChainSlopeWindows/LevelOfCongruence.lean}{199}{mean form of the slope comparison}.
This alternative takes two mean loss bounds. Its proof first obtains
a geometric inequality for every translate separately. Its text is
present in the pinned source, while the saved classical dependency
walk uses the pointwise declaration on line 175. We give the scalar
deduction and then replace those separate HN comparisons by
Corollary~\ref{trmean:budget}. When the host group is larger than
\(A_0\), a bound for its whole-group mean alone gives no bound for
the mean over this particular subgroup.

Let \(G_u\) be the recorded divisibility gain and \(I_u\) the
local orbit loss (the quantities \(\mathrm{Ga}\) and \(\mathrm{IO}\)
in the slope comparison). Let \(T_{ui}>0\) be the word level.
Use \(a_E,\mathcal L_a,\mathcal L,b',X''\) for the numerical
parameters \(a,L_a,L,b',X''\) of the common comparison and
\(c_s=c_{\mathrm{Sl1}}\ge0\). Put \(A=\sum_t a_t\) and
\[
\begin{aligned}
 C_N&=c_{\rm count}(\mathcal L_a+m^2\mathcal L_a/k),&
 C_M&=c_{\rm slope} n_Ta_E\mathcal L_a(1+m/k),\\
 E_{\rm div}&=c_{\rm div}\big(n_Ta_E\mathcal L_a+
                       m(B_0+N_c)x^{e_{\rm div}\varepsilon}\big),&
 K_0&=A C_N+E_{\rm div}+mc_EB_0+C_M.
\end{aligned}
\]
The word loss and its rank mean are
\[
 \ell_{ui}=c_s I_u+c_s n_Tm(a_E/k+b')\log T_{ui}+c_sX''+1,
 \qquad \ell_u=\sum_i r_{ui}\ell_{ui}/r_u .
\]
The counting bound \(\log T_{ui}\le c_{\rm count}m\mathcal L\) gives
\(\ell_u\le c_s I_u+K_7\), where
\(K_7=c_s n_Tm(a_E/k+b')c_{\rm count}m\mathcal L+c_sX''+1\).
Here and below the size parameters multiplying this counting bound
are nonnegative.

\begin{proposition}[A sufficient mean comparison at one scale]\label{trmean:assembly}
Retain the identification of the descended columns, their original
shifted-source slopes and the target slopes of Section~\ref{geo:translate}
with the data just chosen. Suppose the division and lower slope
comparisons on these data give
\begin{equation}\label{trmean:lower}
 \mathbb E(G_u-\mathcal M_u)
       \le \mathbb E(s_u-\mu_u)+K_0+\mathbb E\ell_u .
\end{equation}
If \(F_x\) is either right side of \eqref{trmean:gap} or
\eqref{trmean:budget-eq} under its stated hypotheses, then
\begin{equation}\label{trmean:test}
 \mathbb E D_u\le K_0+K_7+F_x+
                  \mathbb E(\mathcal A_u-G_u+c_s I_u).
\end{equation}
In particular, bounding this right side by \(E_x(\mathcal D_x)\) in
Section~\ref{scope:account} proves \eqref{scope:average} at that scale.
\end{proposition}
\begin{proof}
To see precisely what supplies \eqref{trmean:lower} in the original
word estimates, write \(W_u=\sum_j f^{-1}\sum_\rho\log w_{u,j,\rho}\)
and \(\mathcal M_u=S_0-AW_u\) by the definition of the centre.
Here \(S_0\) is independent of \(u\). Write \(v_0>0\) for the
archimedean scale in the count and \(\ell_g\) for its additive height
parameter. Set
\[
 S_D=S_0-mA\log(\pi/v_0),\qquad
 \Lambda_u=W_u+m\log(v_0/\pi)+m\ell_g-C_N .
\]
Expansion gives \(S_D-A\Lambda_u+mA\ell_g=\mathcal M_u+AC_N\).
For a word let \(F_u^{\,0}\) be the recorded shift common to its
translate, \(E_{ui}\) the sum of its local exponent corrections,
and \(\mu^M_{ui}\) its base lattice slope. The input inequalities are
\[
\begin{aligned}
 \mu_u+h_{ui}-F_u^{\,0}
    &\le S_D-A\log T_{ui}+E_{\rm div},&
 \log T_{ui}&\ge\Lambda_u,\\
 \mu(S_{ui})&\ge\mu^M_{ui}+G_u+F_u^{\,0}-E_{ui}-\ell_{ui},&
 \mu^M_{ui}&\ge-C_M,\\
 E_{ui}&\le m(A\ell_g+c_EB_0).&&
\end{aligned}
\]
Insert the lower bound on \(F_u^{\,0}-h_{ui}\) from the first line
into the second. Since \(A\ge0\), one may replace \(A\log T_{ui}\)
by \(A\Lambda_u\) in the resulting lower bound. The displayed identity
then gives \(G_u-\mathcal M_u\le\mu(S_{ui})-h_{ui}-\mu_u+
K_0+\ell_{ui}\). Averaging by word rank and then by translate proves
\eqref{trmean:lower}. An independent proof of just its mean would
also suffice. Finally use
\(D_u=(\mathcal A_u-G_u)+(G_u-\mathcal M_u)\), the word counting
bound, and the single slope comparison. This gives \eqref{trmean:test}.
\end{proof}

\begin{corollary}[The original bracket from one signed loss]\label{trmean:loss}
Use the common numerical parameters \((\delta_*,\eta_*,e_*)\) of
the original error bracket, including its large-\(x\) size and
error-window hypotheses. Assume the mass and compression bounds of
Corollary~\ref{trmean:budget} with this \(\delta_*\) and the mean
frame bounds \(e\le e_*,\ \eta\le\eta_*\). Write \(X=x^{c_{\rm HH}}\),
\(H=Xb\), \(b>0\). Suppose \(\mathcal J=T_{\rm out}=0\) and
\(U\le K_3+\mathbb E\ell_u\), with \(K_3\) the variation bound below.
In addition to \eqref{trmean:lower} assume
\begin{equation}\label{trmean:signed}
 \mathbb E(\mathcal A_u-G_u+\kappa I_u)\le(1+\kappa)L_P,\qquad
 \kappa=(1+C_\phi)c_s,\quad
 C_\phi=\frac{4\log(1+X)}{1-\delta_*},
\end{equation}
where \(L_P=c_{\rm loss}n_Tmb\,x^{e_{\rm loss}\varepsilon}\).
Then the original numerical mean bracket bounds \(\mathbb E D_u\).
In particular the same numerical size hypotheses as in Formal
input~\ref{red:design} give the right side of \eqref{scope:average}.
\end{corollary}
\begin{proof}
For \(a\ge0\), \(a\log(1+Xb/a)\le(a+b)\log(1+X)\), with value
zero at \(a=0\). For \(a\ge b\) bound the logarithm directly.
For \(0<a<b\), the left side increases with \(a\), since
\(\log(1+t)\ge t/(1+t)\) for \(t\ge0\). Apply this at \(a=U\).
Combine \eqref{trmean:budget-eq} with \eqref{trmean:lower} and
substitute \(U\le K_3+\mathbb E\ell_u\). The coefficient of
\(\mathbb E\ell_u\) is \(1+C_\phi\). The word counting bound and
\eqref{trmean:signed} now give
\[
 \mathbb E D_u\le K_0+(1+C_\phi)K_7+C_\phi(K_3+b)
           +\frac{(\eta_*+4\sqrt{e_*})H}{1-\delta_*}
           +(1+\kappa)L_P .
\]
The source budget adds the nonnegative term
\(\mathbb E\log|\{i:i\text{ is a word at }u\}|\) for nonempty word
sets. This is the only translate-dependent summand in its final error
bracket. We verify its numerical absorption below.
The two sufficient loss assumptions
in the source are \(\mathbb E(\mathcal A_u-G_u)\le L_P\) and
\(\mathbb E I_u\le L_P\). Multiplying the second by \(\kappa\ge0\)
and adding gives \eqref{trmean:signed}.

For comparison with the geometric hypotheses, the frame bound
\(\sum_i\|A_{ui}^*A_{ui}-I\|_{\rm HS}^2\le r_u/4\) gives
\(|M_u-r_u|\le r_u/2\) by Cauchy--Schwarz. The original target
minimum implies \(T_{\rm out}=0\), and integral columns with their
height budgets give \(\mathcal J=0\). For the source variation put
\[
\begin{aligned}
 V_0&=2mc_EB_0+c_H(n_Ta_E\mathcal L_a+mB_0)+Ac_{\rm count}\log k,\\
 K_3&=c_{\rm slope}n_Ta_E\mathcal L_a(2+m/k)+V_0+c_s mB_0 .
\end{aligned}
\]
The exponent, height and count comparisons give
\(|(E_{ui}+h_{ui})-(E_{ui'}+h_{ui'})|\le V_0\).
Indeed the exponent differences cost \(2mc_EB_0\), the differences
of \(h_{ui}+A\log T_{ui}\) cost the term with \(c_H\), and the
log-level differences cost \(Ac_{\rm count}\log k\).
Subtract the lower word-slope estimate in the preceding proof at
\(i'\) from the upper estimate
\(\lambda^{\max}_{ui}\le\lambda^{M,\max}_{ui}+G_u+F_u^{\,0}
-E_{ui}+c_smB_0\) at \(i\). Use
\(\lambda^{M,\max}_{ui}\le c_{\rm slope}n_Ta_E\mathcal L_a\)
and \(\mu^M_{ui'}\ge-C_M\).
This bounds \(\lambda^{\max}_{ui}-h_{ui}-
(\mu(S_{ui'})-h_{ui'})\) by \(K_3+\ell_{ui'}\).
Average first over \(i'\), maximize over \(i\), and take the
translate mean. This gives exactly the asserted bound on \(U\).

For the numerical absorption put \(y=x^\varepsilon\),
\(y_D=y^{e_{\rm div}},\ y_L=y^{e_{\rm loss}}\), and
\(P=n_Ta_E(1+m^2/k)+m(B_0+N_c)\). Here \(y,y_D,y_L\ge1\).
Let \(K\) bound the seven nonnegative coefficients
\(c_{\rm div},c_E,c_H,c_{\rm count},c_{\rm slope},c_s,c_{\rm loss}\).
The scalar size assumptions imply
\[
\begin{gathered}
 1\le n_T,\quad 1\le b\le a_E,\quad mb\le a_E,\quad
 m\le k,\quad 1\le k\le\min(a_E,x^{e_K}),\\
 A\le2n_Ta_E,\quad kb'\le a_E,\quad X''\le a_E,\quad
 0\le\mathcal L,\mathcal L_a\le2y,\quad B_0,N_c\ge0 .
\end{gathered}
\]
In particular \(mb\le a_E\) follows from \(m\le x^{(1+\varepsilon)/2}
\le x\) and \(xb\le a_E\). Also
\(\delta_*\le1/2,\ (\eta_*+4\sqrt{e_*})H\le Kyn_Ta_E\),
\(\log|\{i\}|\le(k-3)\log2\le a_E\) in each block, and
\(\log x\le y\). Put \(K'=1+K,\ h'=1+c_{\rm HH},\ g'=1+e_K\),
with \(c_{\rm HH},e_K\ge0\). From \(P\ge n_Ta_E\ge1\),
\(\log k\le e_K y\), and the displayed formulas one obtains
\[
\begin{aligned}
 K_0&\le12K y y_D P,&
 K_7&\le5(K')^2yP,&
 K_3+b&\le13K'g'yP,\\
 C_\phi&\le8h'y,& L_P&\le K y_L P.&
\end{aligned}
\]
For \(K_0\), the four terms cost at most
\(4KyP,(2y+y_D)KP,KP,4KyP\), respectively.
For \(K_7\), its term containing the word count is at most
\(4K^2y\,n_Ta_E m^2/k\), and \(c_sX''+1\le K'P\).
For \(K_3\) the base slope costs \(6KyP\), the variation costs
\((5+2e_K)KyP\), and \(c_s mB_0+b\le K'P\).
The estimate for \(C_\phi\) uses \(\log(1+x^{c_{\rm HH}})\le h'y\).
Thus, in units of \((K')^2h'g'y^2y_Dy_LP\), the five terms in the
displayed bound on \(\mathbb E D_u\) and the extra log word count
cost at most \(12,45,104,2,10,1\), respectively.
Their sum is \(174\). The common threshold \(y\ge200(K')^2h'g'\)
therefore gives \(\mathbb E D_u\le y^{3+e_{\rm div}+e_{\rm loss}}P\).
Substituting \(B_0\le2nmbx^\varepsilon\),
\(N_c\le n_2kb x^{2\varepsilon}/n^2\) and the specialized exponents
proves \eqref{scope:average}.
\end{proof}

\begin{corollary}[Room for lattice defects on the actual scales]\label{trmean:reserve}
Retain the arithmetic identifications, \eqref{trmean:lower}, the mass,
frame, variation and numerical bounds of Corollary~\ref{trmean:loss}.
Take \(e_{\rm div}\le2,\ e_{\rm loss}=0\). Allow arbitrary
\(\mathcal J,T_{\rm out}\) and write
\(\Delta=\mathbb E(\mathcal A_u-G_u+\kappa I_u)-(1+\kappa)L_P\).
At the same threshold \(\mathbb E D_u\le y^3P+
\Delta+4\mathcal J/(1-\delta_*)+T_{\rm out}\). In particular,
\begin{equation}\label{trmean:reserve-eq}
 \Delta+\frac{4\mathcal J}{1-\delta_*}+T_{\rm out}
                 \le (y^5-y^3)P
\end{equation}
suffices for \eqref{scope:average}. If \(\mathcal J=T_{\rm out}=0\)
and \(\Delta\le0\), the same window and character hypotheses yield
\(G_{\mathcal F_\psi}(x)\le x^{1/2+7\varepsilon}\) for large \(x\).
\end{corollary}
\begin{proof}
In the preceding proof replace \(12Kyy_DP\) by
\(12K\max(y,y_D)P\), using its four separate estimates. With
\(y_D\le y^2\) and \(y_L=1\), the same six coefficients now bound
the costs in units of \((K')^2h'g'y^2P\). The threshold absorbs
their sum into \(y^3P\). Keeping the actual signed loss in place
of its cap adds exactly \(\Delta\) to the upper bound.
Equation~\eqref{trmean:budget-eq} adds the two indicated nonnegative
terms. The variation bound still concerns the original source. Under
\eqref{trmean:reserve-eq} the resulting bound is \(y^5P\), and the
same bounds on \(B_0,N_c\) give \eqref{scope:average}. Under the last
assumptions use \(y^3P\) directly. Propositions~\ref{ar:mean} and
\ref{ar:four} with \(C=3\) give \(G_K\le x^{1/2+6\varepsilon}+\Sigma\).
The excluded twists satisfy \(\Sigma\le4o\,x^{1/2+\varepsilon}\).
Taking \(y\ge1+4o\) proves the stated window bound.
\end{proof}

The more general
Proposition~\ref{trmean:assembly} also allows the quantile bound
without a uniform lower bound on \(d_u/n_u\).
In both cases the order of selection remains that of
Section~\ref{scope:account}. For each prescribed late window
threshold one needs one complete window, then at every required
scale a single host with tied data and the stated mean bounds.
Proposition~\ref{ar:mean} then performs the four-twist extraction.
The reserve test \eqref{trmean:reserve-eq} gives the original window
exponent \(1/2+9\varepsilon\).
The group size incurs no extra logarithm. The uniform arithmetic
premises, including the recorded gains on the actual descended
columns, remain to be established on that family.

\section{Detection of zeros by one-sided estimates}\label{an:section}

The arithmetic construction controls the positive excursions of a real function on intervals whose starting points may grow arbitrarily far apart. The proof below derives a zero-free half-plane from precisely this information. The proof of Lemma~\ref{an:tails} with an arbitrary exponential growth rate is due to the present reconstruction. The pinned development proves the specialization with growth rate \(2\) by explicit exponential tilts. \src{src/Development/AutomorphicArithmetic/DedekindLFunctions/LandauTheorems/landauTails.lean}{80}{formal concentration estimate}.

\subsection{The polar part}
For a finite family \(\mathcal F=((a_i,\chi_i))_{i\in J}\) as in \eqref{eq:centred}, put \(g(u)=G_{\mathcal F}(e^u)\) for \(u\ge0\). The functions are locally continuous, since each summand in the ramp turns on at zero height at its breakpoint. For \(y\ge1\),
\[
 |P_\chi(y)|\le y\sum_{n\le y}\frac{\log n}{n}
 \ll y(1+\log^2 y).
\]
The bound follows by partial summation, or by grouping the integers in dyadic intervals. The finitely many compensation coefficients are fixed. In particular \(|g(u)|\le C_0e^{2u}\), so \(\Phi(s)=\int_0^\infty g(u)e^{-su}\,du\) converges for \(\Re s>2\).

\begin{proposition}[Polar divisor of the family transform]\label{an:poles}
The Laplace transform \(\Phi\) of \(g\) extends meromorphically to \(\Re s>0\). Its poles in this half-plane are the zeros of the family or of its conjugate characters, counted as a set. They are simple as poles. If \(m_\chi(\rho)\) denotes the zero multiplicity, then
\begin{equation}\label{eq:residue}
 \operatorname*{res}_{s=\rho}\Phi(s)
 =\frac{\sum_{i\in J}a_i(m_{\chi_i}(\rho)+m_{\overline{\chi_i}}(\rho))/2}
             {\rho(1-\rho)} .
\end{equation}
Every such residue is nonzero, and at a real pole it is strictly positive.
\end{proposition}
\begin{proof}
Set \(b=0\) and \(f(s)=L(s,\chi)\) for a nonprincipal character. For the conductor-one principal character set \(b=1\) and take the removable continuation of \(f(s)=(s-1)\zeta(s)\). Write \(p=f'/f\). Nonvanishing of \(L(1,\chi)\) in the first case and \(f(1)=1\) in the second make \(p(1)\) meaningful. The constant term of \(\zeta\) at one gives \(p(1)=\gamma\) in the second case. For \(\Re s>1\) the Mellin transform of one ramp is
\[
 \int_n^\infty (y/n-1)y^{-s-1}dy
 =n^{-s}\big((s-1)^{-1}-s^{-1}\big).
\]
Absolute convergence of \(\sum_n\Lambda(n)n^{-\Re s}\) justifies termwise integration. The logarithmic derivative of the Euler product therefore gives the transform of \(P_\chi\) as \(- (L'/L)(s)/(s(s-1))\). In the principal case the two compensations contribute \(-1/(s-1)^2+(1+\gamma)/(s-1)\). In the other case the compensation contributes \(p(1)/(s-1)\). Combining the fractions yields in both cases
\begin{equation}\label{eq:mellin}
 \Phi_\chi(s):=\int_1^\infty G_\chi(y)y^{-s-1}dy
 =-\frac{p(s)-s p(1)+b(1-s)}{s(s-1)} .
\end{equation}
The numerator vanishes at one and is holomorphic there. At a zero \(\rho\) of order \(m\), factoring \(f(s)=(s-\rho)^m h(s)\), \(h(\rho)\ne0\), gives \(p(s)=m/(s-\rho)+h'(s)/h(s)\). This proves the claimed principal part for a single term. The standard nonvanishing theorem on \(\Re s\ge1\) confines these zeros to \(0<\Re\rho<1\). These classical inputs are the ones used by the formal character-family lemma. \src{src/Development/AutomorphicArithmetic/DedekindLFunctions/LandauTheorems/CharFamilyLFunctionNeZeroOfOneLeRe.lean}{9}{classical boundary nonvanishing}.

On the real \(y\)-axis, conjugating \(G_\chi(y)\) replaces \(\chi\) by \(\bar\chi\), including in the compensating term. Thus \(y=e^u\) gives \(\Phi=\sum_i a_i(\Phi_{\chi_i}+\Phi_{\bar\chi_i})/2\) on the half-plane of absolute convergence and hence on continuation. At a zero in this union at least one of the multiplicities in \eqref{eq:residue} is positive. Since all the weights are positive, the common denominator prevents cancellation even for a nonreal \(\rho\). For real \(\rho\), the denominator is positive as well. Finiteness of the family and isolatedness of zeros give only finitely many poles on each compact subset of the open half-plane. \src{src/Development/AutomorphicArithmetic/DedekindLFunctions/LandauTheorems/poleData.lean}{137}{single-character continuation} and \src{src/Development/AutomorphicArithmetic/DedekindLFunctions/LandauTheorems/poleData.lean}{346}{weighted zero multiplicities}.
\end{proof}

\subsection{Concentration at a common derivative order}
We use the following local integral hypothesis on a measurable real function \(g\). For fixed \(\delta>0\), \(0<\theta<1\), every \(\eta>0\) admits arbitrarily large \(U\) such that, for \(I_U=[U,(1+\delta)U]\),
\begin{equation}\label{eq:mass}
 \int_{I_U}g_+(u)e^{-\theta u}\,du\le e^{\eta U},\qquad g_+=\max(g,0).
\end{equation}
A nonnegative majorant with this integral bound suffices. In the formal definition it is explicitly included as a witness. If \(G_{\mathcal F}(y)\le C y^\theta\) on arbitrarily late \([Y,Y^{1+\delta}]\), then \(U=\log Y\) gives a bound \(\max(0,C)\delta U\) for the integral, so \eqref{eq:mass} holds for every \(\eta>0\). \src{src/Development/AutomorphicArithmetic/DedekindLFunctions/LandauZeroFreeCriterion.lean}{26}{integrated logarithmic window}.

\begin{lemma}[Exponential concentration]\label{an:tails}
Given \(a\ge0\) and \(\delta>0\), one can fix \(\sigma>\max(a,2)\) so that for every sufficiently large \(U\) there is an integer \(k\ge U\) satisfying
\[
 \int_{[0,\infty)\setminus I_U}\frac{u^k e^{-(\sigma-a)u}}{k!}\,du
 \le C_\delta\,\sigma^{-k}.
\]
The constant may depend on the chosen \(a,\delta,\sigma\).
\end{lemma}
\begin{proof}
Write \(b=1+\delta/2\) and \(I(v)=v-1-\log v\) for \(v>0\). The two numbers \(I(1/b)\) and \(I((1+\delta)/b)\) are positive. Take \(c\) equal to half their minimum. Choose \(\sigma>\max(2,a)\) so large that, for \(\lambda=\sigma-a\), \(b\lambda\ge1\) and \(\log(\sigma/\lambda)<c\). Put \(k=\lceil b\lambda U\rceil\) and \(n=k+1\).

Under the change \(t=\lambda u\), the integral in the lemma is \(\lambda^{-n}\) times the probability that a variable of density \(t^{n-1}e^{-t}/(n-1)!\) lies outside \([\lambda U,\lambda(1+\delta)U]\). Its exponential moment at \(h<1\) is \((1-h)^{-n}\), by the gamma integral. For \(v<1\) use \(h=1-1/v<0\) to bound the lower tail at \(nv\), and for \(v>1\) use the same formula with \(h>0\) to bound the upper tail. In either case Markov's inequality gives \(e^{-hnv}(1-h)^{-n}=e^{-nI(v)}\). Here the respective ratios \(\lambda U/n\) and \(\lambda(1+\delta)U/n\) converge to \(1/b\) and \((1+\delta)/b\). They eventually lie on the appropriate sides of one and have rate at least \(c\). The integral is at most \(2\lambda^{-n}e^{-nc}\le(2/\lambda)\sigma^{-k}\). The ceiling gives \(k\ge U\).
\end{proof}

\begin{lemma}[A signed coefficient inequality]\label{an:signed}
Suppose \(|g(u)|\le C_0 e^{a u}\) and \eqref{eq:mass} holds. Put \(A_k(s)=(-1)^k\Phi^{(k)}(s)/k!\) in the half-plane of convergence. With \(\sigma\) fixed as in Lemma~\ref{an:tails}, for every \(\eta>0\) there are arbitrarily large \(U\) and integers \(k\ge U\) for which, simultaneously for every \(\tau\in\mathbb R\),
\begin{equation}\label{eq:signed}
 |A_k(\sigma+i\tau)|+\Re A_k(\sigma)
 \le C_1 e^{\eta U}(\sigma-\theta)^{-k}.
\end{equation}
\end{lemma}
\begin{proof}
Differentiation under the integral is justified locally in \(\Re s>a\) by an integrable exponential times \(u^k\). Write \(w=u^ke^{-\sigma u}/k!\). The left side of \eqref{eq:signed} is bounded by \(\int_0^\infty(|g|+g)w=2\int_0^\infty g_+w\), uniformly in \(\tau\). On \(I_U\), the exponential series gives \(u^ke^{-(\sigma-\theta)u}/k!\le(\sigma-\theta)^{-k}\). On the complement use the growth bound and Lemma~\ref{an:tails} at its chosen order \(k\). Now \(\sigma^{-k}\le(\sigma-\theta)^{-k}\) and \(U\) can be taken positive. This proves \eqref{eq:signed} with a constant independent of the window, order, height and \(\eta\). \src{src/Development/AutomorphicArithmetic/DedekindLFunctions/LandauTheorems/landauWindow.lean}{40}{uniform derivative inequality}.
\end{proof}

\subsection{Isolation and contradiction}
\begin{lemma}[A pole detected on a vertical line]\label{an:isolate}
Let \(Z\subset\{0<\Re s<\sigma\}\) be locally finite in \(\Re s>0\), let \(0<\theta<\sigma\) and suppose some point of \(Z\) lies to the right of \(\theta\). Put \(E=\sigma-\theta\). At least one of the following holds.
\begin{enumerate}
\item Some \(z=\sigma+i\tau\), \(\rho\in Z\) and \(0<r\le E\) satisfy \(d=|z-\rho|<r\), \(B(z,r)\cap Z=\{\rho\}\), and \(B(\sigma,r)\cap Z=\varnothing\).
\item There is a real \(\beta\in Z\) and \(r\le E\) such that \(d=\sigma-\beta<r\) and \(B(\sigma,r)\cap Z=\{\beta\}\).
\end{enumerate}
\end{lemma}
\begin{proof}
Intersect \(Z\) with the compact \(\overline B(\sigma,E)\). Let \(D\) be the smallest distance from \(\sigma\) in this finite intersection, or \(E\) if it is empty. The open ball of radius \(D\) is free of poles. Suppose a point \(w\in Z\) has \(\sigma-\Re w<D\). At \(z_0=\sigma+i\Im w\), its distance is \(\sigma-\Re w\). Choose \(e>0\) so small that \(|z_0-w|+e<D\) and \(D+e<\sigma\). For heights within \(e\) of \(\Im w\), only the finitely many points in \(\overline B(z_0,D+e)\) can have distance \(<D\) from the moving centre. Two distinct such points can be equidistant at at most one height. Indeed, subtracting squared distances gives an affine equation in the height. When the imaginary parts agree, the horizontal distances are unequal since both points are left of the vertical line. Choose a height avoiding all these finitely many ties. Its nearest point has distance \(d<D\). A radius strictly between \(d\) and the minimum of \(D\) and the remaining distances gives the first conclusion, with \(D\) alone used when there is no second point.

Otherwise every horizontal distance is at least \(D\). The assumed point right of \(\theta\) forces \(D<E\), so the minimum is attained. Equality of Euclidean and horizontal distances forces any minimizer to be \(\beta=\sigma-D\) on the real axis. Finiteness gives a radius \(r\) between \(D\) and the next distance, capped by \(E\). This gives the second conclusion. \src{src/Development/AutomorphicArithmetic/DedekindLFunctions/LandauTheorems/LandauDichotomy.lean}{28}{real-centre dichotomy} and \src{src/Development/AutomorphicArithmetic/DedekindLFunctions/LandauTheorems/LandauNearest.lean}{36}{perturbation of height}.
\end{proof}

\begin{theorem}[One-sided windows exclude zeros]\label{an:zero}
For a finite positive admissible family as in Proposition~\ref{an:poles}, \eqref{eq:mass} implies that all its character \(L\)-functions are zero-free in \(\Re s>\theta\), away from the pole of \(\zeta\).
\end{theorem}
\begin{proof}
Assume the contrary. Proposition~\ref{an:poles} gives a pole right of \(\theta\), and all poles satisfy \(0<\Re s<1\). Apply Lemmas~\ref{an:signed} and \ref{an:isolate} using \(a=2\) and the same \(\sigma\). In either case choose \(d<r'<r\). Subtracting the unique principal part on \(B(z,r)\), Cauchy's estimate on the circle of radius \(r'\) gives, with \(c_\rho=\operatorname{res}_{\rho}\Phi\),
\[
 \left|A_k(z)-\frac{c_\rho}{(z-\rho)^{k+1}}\right|
 \le M(r')^{-k}.
\]
In the first case \(\Phi\) is also holomorphic in \(B(\sigma,r)\), giving \(|A_k(\sigma)|\le M'(r')^{-k}\). Hence the left side of \eqref{eq:signed} is at least \(|c_\rho|d^{-k-1}-(M+M')(r')^{-k}\). In the second case \(z=\sigma\), \(\rho=\beta\), and both the residue and \(\sigma-\beta\) are positive real numbers. Then the real part alone gives the lower bound \(c_\beta d^{-k-1}-M(r')^{-k}\). In particular no restriction on the arguments of the complex residues or on the parity of the available orders is required.

Writing \(q=d/r'\in(0,1)\) and using \(r'\le \sigma-\theta\), multiplication of the two bounds by \(d^k\) gives
\[
 0<C_2\le C_3q^k+C_1e^{\eta U}q^k.
\]
Here \(C_2=|c_\rho|/d\), also in the real case. Take \(\eta=\tfrac12\log(1/q)\), then choose windows with \(U\to\infty\). Since \(k\ge U\), the right side tends to zero. This contradiction proves the assertion. \src{src/Development/AutomorphicArithmetic/DedekindLFunctions/LandauTheorems/landauWindow.lean}{83}{formal pole exclusion} and \src{src/Development/AutomorphicArithmetic/DedekindLFunctions/LandauTheorems/ZeroFreeOfRpowWindow.lean}{11}{character-family conclusion}.
\end{proof}

\section{Deduction of the critical-line statements}\label{main:end}

\begin{proof}[Deduction of Theorem~\ref{main:thm} from the arithmetic window]
Suppose first that \(\psi\) is primitive and \(L(\rho,\psi)=0\) at a regular point with \(\beta=\Re\rho>1/2\). Put
\[
 \theta=\tfrac12+\tfrac12\min(\beta-\tfrac12,\tfrac12).
\]
Then \(1/2<\theta<\min(1,\beta)\). Choose
\(0<\varepsilon<\min(\varepsilon_0,(\theta-1/2)/9)\). Since \(x\ge1\), Theorem~\ref{main:window} gives \(G_{\mathcal F_\psi}(x)\le x^\theta\) on arbitrarily late intervals of the prescribed type. This is the hypothesis of Theorem~\ref{an:zero} with \(\delta=1/2\). Because the cyclic family contains \(\psi\) with positive weight, the pole in Proposition~\ref{an:poles} would be detected, a contradiction. Both \(\varepsilon\) and \(\theta\) were fixed before invoking the existence of \(Y\). The analytic criterion needs neither an estimate for every \(x\) nor an effective procedure for locating the next \(Y\). \src{src/Formalism/SharpRestriction/RamificationFrobenius/TotientGaloisClasses/TopWin.lean}{68}{formal choice of exponents} and \src{src/Development/AutomorphicArithmetic/PrimitiveCosets/Imports.lean}{8}{identification of the analytic family}.

Let \(\chi\) modulo \(N\) be induced from the primitive \(\psi\) of conductor \(f\mid N\). By equality of Euler products followed by continuation,
\[
 L(s,\chi)=L(s,\psi)\prod_{p\mid N,\,p\nmid f}(1-\psi(p)p^{-s}).
\]
For \(\Re s>0\) we have \(|\psi(p)p^{-s}|\le p^{-\Re s}<1\) for every prime in the product. At \(s=1\) a nonprincipal \(\chi\) has a nonprincipal inducer, since the trivial character induces the principal character. Thus the product proves the required half-plane assertion with precisely the exclusion in Theorem~\ref{main:thm}. It also reduces every strip zero to one of a primitive character. \src{src/Development/AutomorphicArithmetic/CharacterBoxes/Imprimitive.lean}{5}{change of conductor} and \src{src/Development/AutomorphicArithmetic/CharacterBoxes/ClaimStrip.lean}{23}{transfer of strip zeros}.

Write \(\epsilon(\psi)\in\{0,1\}\) for the parity. In the open strip, the nonzero factor \(\Gamma_{\mathbb R}(s+\epsilon(\psi))\), together with any positive-conductor exponential in the chosen normalization, identifies zeros of \(L(s,\psi)\) with zeros of its completion. If a zero had \(\Re s<1/2\), the completed functional equation would give a zero of the completion for \(\bar\psi=\psi^{-1}\) at \(1-s\), whose real part exceeds \(1/2\). The gamma factor there is nonzero and the point is regular. This contradicts the half-plane result. This proves both strip statements. \src{src/Formalism/GammaFactorNonvanishing/CritLine.lean}{11}{primitive functional equation} and \src{src/Development/AutomorphicArithmetic/CharacterBoxes/ClaimStrip.lean}{59}{all characters in the strip}.

For the zeta assertion first take \(N=1\). If a remaining zero has \(\Re s<1/2\), handle \(s=0\) by \(\zeta(0)=-1/2\). Otherwise \(t=1-s\) has \(\Re t>1/2\), \(t\ne1\). The zeta functional equation in the normalization of the source reads
\[
 \zeta(s)=2(2\pi)^{-t}\Gamma(t)\cos(\pi t/2)\zeta(t).
\]
The exponential and \(\Gamma(t)\) are nonzero. A zero of the cosine would force \(t\) to be an odd integer, hence \(s\) an even integer. In the present half-plane \(s\) would then be zero or a negative even integer, already excluded. Since \(\zeta(t)\ne0\), this rules out the last possibility. It proves Mathlib's zeta statement, including points outside the open strip. \src{src/Formalism/GammaFactorNonvanishing/RhOfRight.lean}{5}{zeta reflection and exceptional points} and \src{src/Supplements/Kernels/MainTheorems.lean}{9}{specialization at modulus one}.
\end{proof}

\section{Machine verification and statement alignment}
\label{cert:section}
\begingroup
\raggedright
\setlength{\emergencystretch}{4em}

The submission pins Lean 4.33.1 in \texttt{lean-toolchain}. Its five public roots occupy \src{src/Main.lean}{4}{\texttt{Main.lean}, lines 4--12}. At the level of the stored declaration, the first root has type \texttt{RiemannHypothesis} and value \texttt{rhOfAll}. The second has type \texttt{RiemannZetaZerosOnLine} and the same value. The third has type \texttt{DirichletLNonvanishingRightOfHalf} and value \texttt{allChi}. The fourth has type \texttt{DirichletLZerosOnLine}, with value given by the strip projection of \texttt{claimAll}. The fifth has type \texttt{PrimitiveDirichletLZerosOnLine} and value \texttt{critAll}. The explicit zeta statement excludes negative even integers and $s=1$. The half-plane statement excludes the pair of the principal character and $s=1$. The two Dirichlet critical-line statements are restricted to the open critical strip. The last proposition additionally restricts to primitive characters. The hypotheses are given in \src{src/Statements.lean}{3}{\texttt{Statements.lean}, lines 3--15}. These identifications use the type and proof-value records of the fixed export and the matching source declarations.

\paragraph{Full independent checks and local dependency audit.}
Sutherland reports that all 175,187 compiled modules passed kernel replay. Nanoda and ConLeche accepted the full export, ConLeche compared all five statements with the challenge, and Lean4Lean checked all the compiled modules, including five reruns with a larger step limit \cite{mit}. The record dated 21 November at 09:30 UTC still listed the full-export Lean-kernel and end-to-end Comparator runs as in progress. Its Nanoda log reports a peak resident set of 88,880,436~kB. The local dependency audit traversed constants in declaration types and values together with inductive, constructor, and recursor metadata. The respective root-inclusive closures contain $1{,}478{,}219$, $1{,}478{,}221$, $1{,}478{,}211$, $1{,}478{,}224$, and $1{,}478{,}214$ declarations. The public LUMI table independently gives the first and third counts. Each closure contains precisely the axioms \texttt{Classical.choice}, \texttt{propext}, and \texttt{Quot.sound}. The local graph audit computes dependency reachability from the stored terms. The full proof-checking verdict comes from the public independent runs. See \url{https://math.mit.edu/~drew/grh/} and \url{https://math.mit.edu/~drew/grh/nanoda.log}.

\paragraph{Independent local replay of original declarations.}
The extraction program under \texttt{verification/} checks the SHA-256 of the canonical compressed export, then copies its decompressed bytes from the first record through the original theorem record for \texttt{Deduction.landau\_window}. The resulting prefix has 1,539,032,259 bytes and 140,192 declaration records. Its SHA-256 is \texttt{dfabcd8e32844354b199a7f360c1f11ff}\allowbreak\texttt{e8d7c0bcf67ab6e04c7d36ddbedb08b}. Three original proof values in the prefix mark the scope. The first is \src{src/Formalism/GammaFactorNonvanishing/CritLine.lean}{29}{\texttt{Formalism.VW.rh\_of\_every}}, whose hypothesis is \texttt{UpperTheory.EveryCharacter}. The second is \src{src/Development/AutomorphicArithmetic/DedekindLFunctions/LandauTheorems/poleData.lean}{248}{\texttt{Deduction.poleData}}, giving pole data for the character-family series. The third is \src{src/Development/AutomorphicArithmetic/DedekindLFunctions/LandauTheorems/landauWindow.lean}{83}{\texttt{Deduction.landau\_window}}, which gives $\Re\rho\leq\theta$ for every $\rho\in Z$ from \texttt{PoleData}, $0<\theta<1$, $\delta>0$, and \texttt{WindowHypU}. The prefix contains the exact original records and ends before all five public roots.

The same prefix contains original proof-value records for \src{src/Development/AutomorphicArithmetic/DedekindLFunctions/LandauTheorems/ChCore.lean}{343}{\texttt{Deduction.ch\_core}}, \src{src/Development/AutomorphicArithmetic/DedekindLFunctions/LandauTheorems/ChTop.lean}{37}{\texttt{Deduction.ch\_top}}, and \src{src/Development/AutomorphicArithmetic/PrimitiveCosets/CoreGm.lean}{12}{\texttt{UpperTheory.Comparison}$_2$\texttt{.core\_gm}}. The arithmetic witnesses and their raw export indices are recorded in \path{verification/arithmetic-prefix-receipt.json}. The auxiliary-box density estimate occurs beyond the cut, as do \texttt{UpperTheory.\allowbreak Top.topDesWin} and \texttt{UpperTheory.\allowbreak Gal.Lim.exists\_trace\_limit}.

On native Windows, Nanoda 0.4.17 from the Lean 4.36.0 distribution checked the unmodified prefix and exited successfully. It reported ``Checked 145327 declarations with no errors.'' The larger figure is Nanoda's internal declaration count, while 140,192 counts raw declaration records, each of which may encode an inductive block with multiple constants. The run used a strict allowlist consisting of the three standard axioms above. The success log, configuration, receipts, and reproduction commands are in \path{verification/}. The largest sampled peak working set for this checker process was 1,602,121,728 bytes. The scope of this local kernel verdict is the prefix.

The bundled ConLeche checker independently accepted the same prefix in verified mode with two workers. It reported 140,188 accepted declarations in its own counting convention and completed the 137,728 deferred checks. Its reported elapsed time was 816.7 seconds. The largest sampled peak working set for this process was 2,928,775,168 bytes.

\paragraph{Fresh finite-window interface.}
The file \path{verification/WindowBridge.lean} develops a finite-set induction over the logarithmic interval and then transports it through the exponential map. For an exceptional set $E\subseteq\mathbb N$ with $|E|\le R$, parameters $0<\eta\le C$ and $X>1$, it yields $Y$ such that, with $J=4(2C/\eta)^R$,
\[
 X^{1/J}\le Y,\qquad Y^2\le X,\qquad
 \forall x\in[Y,Y^2]\ \forall k\in E,\quad k\le x^C\Longrightarrow k\le x^\eta.
\]
Here is the corresponding real-variable argument. Put $B=C/\eta$, $U=\log X$, and send each positive exceptional integer $k$ to $a=\log(k)/C$. If there is an exceptional $a\in(U/(2B),U]$, remove it and recurse with upper bound $U/(2B)$. The recursion loses a factor $2B$ and uses one unit of the cardinal bound. If that interval is empty, take $y=U/2$. For each $u\in[y,2y]$, any exceptional $a\le u$ then has $a\le U/(2B)$, so $Ba\le u$. Exponentiating the selected $y$ yields the claim. The case $k=0$ follows directly from positivity of $x^\eta$.

The constant is the source constant since $4(2B)^R=2^{R+2}B^R$. The conclusion agrees with \src{src/Development/AutomorphicArithmetic/AutomorphicVarieties/LayeredIsomorphisms/Window1.lean}{6}{the logarithmic selection lemma} and \src{src/Development/AutomorphicArithmetic/AutomorphicVarieties/LayeredIsomorphisms/Window.lean}{7}{its power-window form}, using the definitions of \src{src/Development/AutomorphicArithmetic/AutomorphicVarieties/PrimeSetBoundDataOfPrimitiveChar.lean}{30}{\texttt{InRange}} and \src{src/Development/AutomorphicArithmetic/AutomorphicVarieties/LayeredIsomorphisms/SecondBound.lean}{46}{\texttt{Good}}. The cardinal bound on $E$ is an explicit hypothesis, and delivering such a bound is an upstream arithmetic obligation. The fresh file was checked by Lean 4.36.0 against Mathlib release \texttt{v4.36.0}. Its only imports are Mathlib modules. Lean's \texttt{\#print axioms} reports precisely \texttt{propext}, \texttt{Classical.choice}, and \texttt{Quot.sound} for each of the two bridge theorems.
\paragraph{Fresh critical-line implication.}
The separate file \path{verification/CriticalLineBridge.lean} checks the quantifiers in Section~\ref{main:end} using Mathlib alone. For a fixed real function it assumes the arithmetic windows at every sufficiently small positive \(\varepsilon\), and assumes the analytic exclusion at each \(0<\theta<1\) for that same function. The Lean proof picks \(\theta\) from the proposed zero's real part, then chooses one \(\varepsilon\) and transports the bound at every point of the resulting windows. Its conclusion excludes every detected point to the right of one half. A separate declaration checks \(M/r^k<c/d^{k+1}\) eventually for \(0<d<r\), \(c>0\), as used to compare a nonzero residue to a Cauchy remainder. All four fresh declarations passed Lean 4.36.0 with the three standard axioms. The analytic exclusion and the arithmetic supply in the implication retain exactly their stated status as hypotheses.
\paragraph{Finite averaging.}
The accompanying \path{verification/TranslateMean.lean} checks two parts
of Proposition~\ref{ar:mean} in Mathlib. For any homomorphism
\(\iota:A\to G\) from a finite group and any additive character
\(\chi\) of \(G\) with values in a field, it computes
\(\sum_{u\in A}\chi(t+\iota(u))\). The sum is
\(|A|\chi(t)\) when the restricted character is trivial, and zero otherwise.
The second proof transfers a uniform absolute error through a normalized
finite average of real numbers. These are formal checks of the
projection and averaging steps. The signs of the arithmetic characters,
the estimates for the three discarded twists, and the geometric input
are justified at their respective places in the text.
\par\endgroup

\subsection{Proof values of the two deformation counts}\label{audit:countpaths}

The proof value of the cohomological window theorem directly uses both
the rank-three Case I count and the joint rank-four Case II and III
count. The accompanying archive \path{verification/zeta-count-trace.zip}
gives the original names and expression IDs for the steps described
in Section~\ref{zetaring:section}. In particular it records the
ternary single-character ring construction, the quartic test detecting
nonfiniteness, and the specialization killing the tower ideal in the
Case II component argument.

The 96 selected declarations lie in the classical RH dependency tree.
All 99 direct references have two witnesses: a constant in the original
compiled value, and a path from the value root to that constant in the
original serialized export. Of these edges, 85 are also immediate
value edges of the saved breadth-first tree. That tree retains one
parent for each declaration. The archive contains 2,935 original
records, including the complete name chains and the type and value
indices for all 95 definitions and theorems. The remaining declaration
is the constructor \(\mathrm{And.intro}\). The verifier checks every
path and can re-extract the original modules from their pinned
archives. The zeta host calculations and the mixed component mean
criterion are paper arguments.

The following pairs are expression IDs in the unmodified export.
The last column counts child-expression steps inside one stored
value. The full sequences, including the two ends of each path,
are in \path{proof-edges.tsv} in the archive.
\begin{center}\small
\begin{tabular}{@{}lrrr@{}}
\toprule
Reference in the stored value & Value root & Constant leaf & Steps\\
\midrule
Window to ternary count                 & 991824206 & 991824173 & 12\\
Window to joint quartic count           & 991824206 & 991824164 & 7\\
Quartic count to conjunction constructor & 952152235 & 1242      & 14\\
Nonfinite test to ambient point         & 962691518 & 962667923 & 57\\
Case II propagation to its last step    & 962437115 & 962426112 & 157\\
\bottomrule
\end{tabular}
\end{center}

\section{The remaining arithmetic estimate and the scope of the proof}\label{scope:account}

The proof of Theorem~\ref{main:thm} has two logically different dependencies. On its analytic side, equations~\eqref{eq:centred} and \eqref{eq:window}, ordinary facts about Dirichlet \(L\)-functions at and to the right of one, and their functional equations suffice. The derivations of the polar part, concentration, isolation and final contradiction are given in full. Lemma~\ref{an:tails} has a different proof and choice of derivative order from the formal specialization. Its moment identity fixes the gamma shape at \(k+1\), so the probability estimate applies to the integral including its normalizing factor \((\sigma-a)^{-k-1}\). This factor is essential when the bound is compared to \(\sigma^{-k}\).

For clarity, the following is a sufficient arithmetic statement which a reader could establish by a different method. Fix the primitive character and choose any sufficiently small \(\varepsilon>0\). Require that the retained characters give precisely \(\mathcal F_\psi\), with its weights, and that the excluded contribution obeys \(\Sigma\le4o\,x^{1/2+\varepsilon}\). For every lower bound \(Y_0\) one needs a \(Y\ge\max(1,Y_0)\) such that at \emph{each} \(Y\le x\le Y^{3/2}\) there are data as in Section~\ref{ar:section}. The hypotheses on coset shares, scales, auxiliary twists and ranks in Lemmas~\ref{ar:balance} and \ref{ar:extract} and Proposition~\ref{ar:four} must hold. Fix \(o\) from \(\psi\). The derivative constants and the constants in the retained and discarded logarithmic-derivative bounds of Lemma~\ref{ar:extract} must be uniform in \(Y\), in \(x\), and in the data chosen there. This gives one threshold in that lemma for the fixed \(\psi,\varepsilon\), so that an entire window can be placed beyond it. The field, scales and finite sets at \(x\) may depend on \(x\), but for \emph{every} translate \(u\in A_0\) of those same data one needs
\begin{equation}\label{scope:missing}
\begin{split}
 D_u&=\sum_{t\in T}\sum_{j=1}^{m}
       (\mathrm{ASK}_{u,t,j}-\mathrm{MAIN}_{u,t,j})\\
 &\le x^{5\varepsilon}
 \left(n_Ta_E+n_Ta_E\,m^2/k+2n m^2 b x^\varepsilon+
               mb\,n_2 k x^{2\varepsilon}/n^2\right).
\end{split}
\end{equation}
Here \(\mathrm{ASK}\) is the explicit finite prime-power sum in Lemma~\ref{ar:balance}, and \(\mathrm{MAIN}\) its centring term. The exponent \(5\) is \(2+0+3\) at the geometric specialization. Lemma~\ref{ar:extract} and Proposition~\ref{ar:four} give \(G_K\le x^{1/2+8\varepsilon}+\Sigma\). The excluded-character estimate in Section~\ref{ar:section} absorbs \(\Sigma\) at exponent \(1/2+9\varepsilon\) when \(x\) is sufficiently large, and \(G_K=G_{\mathcal F_\psi}\). This proves the claimed sufficiency before invoking the analytic argument. In the formal construction the two auxiliary primes are fixed on the initial interval, and the degree of the CM field can change from one scale to the next.

Writing \(\mathcal D_x\) for the common data just specified and
\(E_x(\mathcal D_x)\) for the right side, the order at fixed
\(\psi,\varepsilon\) is
\[
\begin{gathered}
 \forall\,Y_0\quad \exists\,Y\ge\max(1,Y_0)\quad
 \forall\,x\in[Y,Y^{3/2}]\quad \exists\,\mathcal D_x\\
 \forall\,u\in A_0(\mathcal D_x),\qquad
           D_u(\mathcal D_x)\le E_x(\mathcal D_x).
\end{gathered}
\]
There is a weaker sufficient condition for the paper implication.
Keep this order through the choice of \(\mathcal D_x\), the same
thresholds, and all its remaining conditions. Replace the last line by
\begin{equation}\label{scope:average}
     \frac{1}{|A_0(\mathcal D_x)|}
       \sum_{u\in A_0(\mathcal D_x)}D_u(\mathcal D_x)
                \le E_x(\mathcal D_x).
\end{equation}
Proposition~\ref{ar:mean} proves the extraction inequality under
\eqref{scope:average}, so the proof of Theorem~\ref{main:conditional}
applies unchanged thereafter. The same proposition evaluates the
mean in terms of four primitive twists for each power of \(\psi\).
Corollary~\ref{ar:meanhost} bounds the combined discarded contribution.
Corollary~\ref{ar:squares} gives a discrete sufficient test: once
the mean bound holds at the integer squares in these windows,
continuity and an explicit slope bound for the retained function
extend the analytic estimate to all real scales. The uniform upper
bound at those squares remains an arithmetic hypothesis. Proposition~\ref{trmean:assembly} reduces it to a mean division estimate, one slope comparison and one signed loss. Corollary~\ref{trmean:reserve} gives the available margin for lattice defects at the original numerical scale.

The following describes what the arithmetic proofs establish on the way to \eqref{scope:missing}.
\begin{description}
\item[Prime selection and the common field.] Theorem~\ref{pr:pair} treats the possible exceptional real zero and selects both primes for every required auxiliary twist and every real point in the interval. Proposition~\ref{host:art} identifies the retained and discarded characters. Lemma~\ref{host:different} constructs the lattice at finite places; Propositions~\ref{host:behrend} and \ref{host:thin} give the translated set and the uniform thin-coordinate estimate. These deductions use ordinary mathematics and the stated prime distribution theorems.
\item[Galois representations.] Sections~\ref{tr:section} and \ref{newgal:section} realize the trace as a semisimple representation and give crystalline closure from simultaneous weighted trace congruences. The Frobenius index runs over every lift at a good prime. For a finite set of seed systems this produces a common dense set on which to interpolate. Proposition~\ref{famgal:rank} extracts the exact rank even for disconnected monodromy. Propositions~\ref{newgal:hodge} and \ref{newgal:phi} test the rectangular lattice maps on the integral period operators at the same precision as the trace congruence, allowing unbounded linear-combination coefficients. Applying those propositions at every singular intermediate weight still needs the geometric seed, simultaneous congruences and the two marked integral comparisons. Regular polarized cuspidal systems satisfying the stated local and irreducibility hypotheses have the level and adjoint Selmer conclusions of Propositions~\ref{newgal:level} and~\ref{newgal:adselmer}.
\item[Global theta matching.] Theorem~\ref{tcr:global} proves nonzero proportionality of the two periods on every matched pair in the stated tower, including positive signature and the support restriction at the distinguished prime. Proposition~\ref{famtheta:plane} supplies a hyperbolic plane on the given lattice. Propositions~\ref{tcr:descend} and \ref{tcr:unique} compare the two descent maps on arbitrary smooth inert tests and real polynomial tests. The resulting recursions include the split Euler and ramified Gauss factors. For fixed adapted geometries, real quotient maps and measures, their constant is independent of the allowed finite characters, finite matching maps and compatible Weil actions. The period scalar for an individual system may depend on that system. This settles the period comparison required here. The reproducing identity for short algebraic sections on the covers involves a separate kernel.
\item[Theta denominators and cover metrics.]
Theorem~\ref{intjet:semi} clears the finite denominators on the integral
semiabelian tests. Proposition~\ref{bfam:harm} gives short harmonic
generators on every self-dual column lattice. Lemma~\ref{bfam:phase}
and Corollary~\ref{bfam:mark} fix the finite test and its phase at the
identity auxiliary level. On a polarized period chart the column
factor in the reading is explicit in Proposition~\ref{bfam:unitread}.
Propositions~\ref{dg:identity}, \ref{arrep:reproduce}
and~\ref{munf:span} give the doubled graph identity, its positive real
reproducing scalar and simultaneous generation on the marked period
components. Lemma~\ref{crev:mode} shows how the finite mark kills
the zero Fourier mode of each raw theta section on a rational
isotropic row line. Proposition~\ref{crev:half} proves boundedness
and a lower bound for the shortest vector at every half-maximum
frame. Proposition~\ref{crev:mass} computes the flag volume and
the probability measure on the selected analytic packet.
Theorem~\ref{ctwist:semi} gives the integral theta key on the
rigidified graph-line tests. On the indefinite marked rows of
Lemma~\ref{cmeig:row}, Proposition~\ref{cmeig:degree} evaluates
the ordinary Lie determinants at a product CM point, including
the different at ramified primes. Theorem~\ref{cmeig:splitmodel}
compares the vertical orders on the classical splitting models
at each fixed finite place and embedding.
Proposition~\ref{host:period} bounds the resulting scalar uniformly
as the field and the weights vary. For a selected good cover the
space of regular cusp rules is defined by its readings at the good
embeddings. Corollary~\ref{crev:twist} treats those readings and the
effect of both invertible weight changes on their cusp condition.
Corollary~\ref{cmeig:final} supplies the integral spanning set in
the final weight with its explicit bound on each such good
indefinite cover at the stated stage sizes.
Obtaining the compatible good covers throughout the design and
identifying a separately constructed toroidal space of sections
on the modified flags have their own geometric hypotheses.
The integral restriction maps, short frames and compatible metrics
for the modified flag spaces in the Hilbert--Samuel estimate are
additional geometric inputs on the same components.
\item[Integral cohomology and deformation rings.]
Theorem~\ref{newcoh:lemmaK} gives the low-layer estimate with an absolute
exponent \(32\). In \eqref{sup:lengthcount} the joint arithmetic tests
inject into the image scaled by \(\ell^{d-1}\); the universal
coefficient quotient costs at most the torsion order in degree
\(d-1\). The construction and the prescribed evaluations of the
arithmetic cycle used for this injection are stated with that formula.
Theorem~\ref{famsel:unit} and Proposition~\ref{famsel:stick}
prove the auxiliary Bernoulli unit and its residual vanishing in both
higher-rank counts. Proposition~\ref{famsel:class} gives zero at
all integral depths under its residual class-group test, including
nontrivial even- and odd-order characters. Theorem~\ref{famsel:rate}
bounds the strict lengths over all primes and characters by
\(\exp((16Ne)^{102})\). Corollary~\ref{famsel:subwindow} gives the
required exponential budget for the strict character contribution on
\(16fe\le x^{\varepsilon/204}\), with \(f\) the primitive conductor.
Propositions~\ref{famsel:hostinertia}--\ref{famsel:splitray} also
treat specified conductors and congruence classes at arbitrarily
large twists. For a cubic target modulo seven, every case III entry
vanishes at \(\ell=29\) and for all sufficiently large coefficient
primes in twelve of the twenty-four reduced classes modulo \(84\).
For a target of two-power order, case III is empty, although case II
already occurs for the zeta target (Corollary~\ref{famsel:hostdensity}). Section
\ref{famsel:complement} sums the complementary contributions on the
same host up to \(\lfloor v^{401}\rfloor\). That computation still
uses the original polynomial bound \((\ell N e)^{C_\chi}\) on the
remaining case II and III entries, the case I sum and the indicated
field comparisons.

The ternary endoscopic and odd orthogonal branches admit classical
attachments. In the good quartic problem, Corollary~\ref{hr:imprimitive}
gives cuspidal attachment and rational symplectic adjoint Selmer
vanishing at every non-strongly-irreducible, absolutely irreducible
characteristic-zero point. Proposition~\ref{hr:case3strong} forces
strong irreducibility for absolutely irreducible tests in the third
host case, in either characteristic. Such a restriction on a test
still leaves the strongly irreducible curves to be controlled.
Section~\ref{hr:tower} traces the ternary and quartic seeds,
dimensions, markings and base changes; Lemma~\ref{hr:towerfin}
controls the selected Hecke quotient. Covering every required test
curve still uses the seed existence and the arithmetic transfers at
the specialized points. The one-sided point exclusions also enter
the full length estimate.
\src{src/Development/AutomorphicArithmetic/AutomorphicVarieties/HeightAndVanishing/JointWinOfCounts.lean}{30}{vanishing and counting in the window}.

\end{description}

To illustrate the finite-fibre issue, let \(O\) be a complete discrete valuation ring with uniformizer \(\pi\) and residue field \(k\). In \(R=O[[T]]/(\pi T)\) every characteristic-zero field point has \(T=0\), its generic fibre is the field \(\operatorname{Frac}O\), and \(R/(T)=O\) is finite. But \(R/\pi R=k[[T]]\) has a nonconstant curve. It follows already by reduction modulo \(\pi\) that \(R\) is infinite as an \(O\)-module. This example explains why checking characteristic-zero attachments and the quotient singled out by a Hecke ideal cannot on their own discharge the curve hypothesis in Lemma~\ref{dp:detect}.

\paragraph{One use of a proof term inside the geometry.}
In the family realization a local proof \(h_{\rm onto}\) is obtained from the theta-doubling theorem. The following line in the original uses tactic search to construct the family record. Reading the elaborated expression resolves which premise tactic search supplied. At the application of the row-data theorem the local proof appears once, at argument 21 counting implicit arguments from one. The resulting family record appears once, at argument 16 of the final application. In ordinary terms, for \(m=2n+2\) and a base row \(T\), surjectivity supplies a terminal vector whose descent to level zero equals \(T\). The source constructs the reachable pair of spaces, packages \(T\) into the initial stage, applies surjectivity, and takes underlying values to recover the desired equality. The independent valuation argument in Section~\ref{geo:engine} has its own inputs. The expression slice in the supplement reproduces this binder calculation on the original indices. \src{src/Development/AutomorphicArithmetic/ThetaCorrespondence/ThetaPairLifts/ThirdPlaceSeries/MainTheoremFromPublishedResultsAltAltAltAlt.lean}{25}{the two local proof bindings}, \src{src/Development/AutomorphicArithmetic/ThetaCorrespondence/ThetaPairLifts/ThirdPlaceSeries/FamAllQOfAllTrue.lean}{51}{surjectivity field}, \src{src/Development/AutomorphicArithmetic/ThetaCorrespondence/ThetaPairLifts/ThirdPlaceSeries/NotMemQBase.lean}{37}{row-space surjectivity} and \src{src/Development/AutomorphicArithmetic/ThetaCorrespondence/ThetaPairLifts/ThirdPlaceSeries/LinkOntoOfOntoAll.lean}{49}{lifting a specified base row}.

\paragraph{Mathematical review.}
The review compared the displayed implications with the source at the points where sign, rank, normalization or quantifier order mattered. In the comparison of theta integrals, uniqueness of the scalar uses a test with nonzero integral \emph{and a matching test on the other side}; the matching step for this particular vector is now explicit in Proposition~\ref{geo:theta}. The integer clearing the terminal period denominator is divisible by the stabilizer orders in \emph{all} layers. It need not equal any one of them. The slope compression requires both a bound \(0\le\delta<1\) and a codomain of dimension \(D_0\ge(1-\delta)D\). The determinant calculation in Corollary~\ref{geo:shift-index} fixes the sign and factor of two. In the Poitou--Tate computation of Section~\ref{sup:length}, the discrete dual of \(\mathbf Z_\ell(E+1)\) has twist \(-E\), which fixes the local \(H^0\) term. \src{src/Formalism/GaloisCohomology/Statement.lean}{69}{cyclotomic action on the discrete module}. Other checks included common-factor positivity for complex residues, the strict radius gap in Lemma~\ref{an:isolate}, the power \(\lfloor v^{401}\rfloor\) preceding the cohomological rank cutoff and the family selection before the pointwise Fourier estimates. The accompanying tables record the proof status and the source locations. In particular the rank-four trace condition is checked against the variant actually used in the export, whose automorphy predicate includes additional comparison witnesses. Each source anchor is paired with an enclosing declaration and a path from the classical RH root or, for the distinct strip and restatement assertions, a path from the appropriate sibling root. These paths certify the presence of the declarations. They do not assign a separate theorem to every tactic line.

The English argument proves the critical-line implication from
\eqref{scope:missing} and also from \eqref{scope:average} with the
specified common data. On the available paper arguments both uniform
arithmetic estimates remain open, even when the target is zeta.
The unresolved component, length and integral geometric inputs
therefore prevent an independent paper proof of RH in this edition.
Published kernel checks concern the original exported declarations.


\phantomsection\addcontentsline{toc}{section}{References}
\begin{thebibliography}{99}
\bibitem{announce} L. Alpöge, \emph{Announcement of a proof of the Riemann hypothesis}, 20 November 2026. \href{https://alpo.ge/remark.pdf}{Author's release note}.
\bibitem{source} L. Alpöge with Claude, \emph{Lean formalization of RH and Dirichlet GRH}, source distribution of 20 November 2026, Lean 4.33.1. \href{https://github.com/leanprover-community/riemann-hypothesis/tree/40e56bab026552e089d7afcfeff08d2defb9c005}{Pinned community mirror}.
\bibitem{mathlib} The Mathlib Community, \emph{Riemann zeta function}, L-series library at the revision imported by the submission. \href{https://github.com/leanprover-community/mathlib4/blob/0df444a360eaa60ab8c11dca51a86af692955474/Mathlib/NumberTheory/LSeries/RiemannZeta.lean\#L162}{Definition of \texttt{RiemannHypothesis}}.
\bibitem{mit} A. V. Sutherland, \emph{Verification files for the Lean proof of the Riemann Hypothesis}, status at 09:30 UTC on 21 November 2026. \href{https://math.mit.edu/~drew/grh/}{Verification records, binaries and exports}.
\bibitem{lumi} \texttt{aapolipponen}, \emph{Independent verification of Alpöge's Lean formalization of RH/GRH, run on LUMI}, November 2026. \href{https://github.com/aapolipponen/rh-lean-check-lumi/tree/82d91dbbd04d71f2cb45c3f2898f421b91c1c9d7}{Pinned report and numerical dependency audit}.
\bibitem{chenevier} G. Chenevier, \emph{The \(p\)-adic analytic space of pseudocharacters of a profinite group and pseudorepresentations over arbitrary rings}, §§1.26 and 2. \href{https://arxiv.org/pdf/0809.0415}{Proposition 1.27 and Theorem 2.12}.
\bibitem{maynard} J. Maynard, \emph{On the Brun--Titchmarsh theorem}, Acta Arith. 157 (2013), 249--296. \href{https://doi.org/10.4064/aa157-3-3}{Theorem 3.2}.
\bibitem{tz} J. Thorner and A. Zaman, \emph{Refinements to the prime number theorem for arithmetic progressions}, Math. Z. 306 (2024), article 54. \href{https://arxiv.org/pdf/2108.10878}{Corollary 1.4 and Remark 1.5}.
\bibitem{antLifting} P. B. Allen, J. Newton and J. A. Thorne, \emph{Automorphy lifting for residually reducible $l$-adic Galois representations, II}, Compos. Math. 156 (2020), 2399--2422. \href{https://doi.org/10.1112/S0010437X20007484}{Text}.
\bibitem{serreTensor} {\raggedright Z. Amir-Khosravi, \emph{Serre's tensor construction and moduli of abelian schemes}, Manuscripta Math. 156 (2018), 409--456. \href{https://arxiv.org/abs/1507.07607}{Text}.\par}
\bibitem{banaszczyk} W. Banaszczyk, \emph{New bounds in some transference theorems in the geometry of numbers}, Math. Ann. 296 (1993), 625--635. \href{https://doi.org/10.1007/BF01445125}{Text}.
\bibitem{blggtLocal} T. Barnet-Lamb, T. Gee, D. Geraghty and R. Taylor, \emph{Local-global compatibility for $l=p$, II}, Ann. Sci. Éc. Norm. Supér. (4) 47 (2014), 165--179. \href{https://www.numdam.org/item/10.24033/asens.2212.pdf}{Text}.
\bibitem{blggtPotential} T. Barnet-Lamb, T. Gee, D. Geraghty and R. Taylor, \emph{Potential automorphy and change of weight}, Ann. of Math. (2) 179 (2014), 501--609. \href{https://annals.math.princeton.edu/wp-content/uploads/annals-v179-n2-p03-p.pdf}{Text}.
\bibitem{bayerEmeryHouriet} E. Bayer-Fluckiger, V. Emery and J. Houriet, \emph{Hermitian lattices and bounds in $K$-theory of algebraic integers}, Doc. Math., Extra Volume: Alexander S. Merkurjev's Sixtieth Birthday (2015), 71--83. \href{https://ems.press/content/book-chapter-files/27436}{Text}.
\bibitem{bellaicheChenevier} J. Bella\"iche and G. Chenevier, \emph{Families of Galois representations and Selmer groups}, Astérisque 324 (2009). \href{https://www.numdam.org/item/AST_2009__324__R1_0/}{Text}.
\bibitem{bernsteinInvariant} J. N. Bernstein, \emph{$P$-invariant distributions on $\mathrm{GL}(N)$ and the classification of unitary representations of $\mathrm{GL}(N)$ (non-Archimedean case)}, in \emph{Lie group representations II}, Lecture Notes in Mathematics 1041, Springer, 1983, 50--102. \href{https://www.math.tau.ac.il/~bernstei/Publication_list/publication_texts/Bernstein-P-invar-SLN.pdf}{Text}.
\bibitem{bzgLemma} I. N. Bernstein and A. V. Zelevinsky, \emph{Induced representations of reductive $p$-adic groups. I}, Ann. Sci. École Norm. Sup. (4) 10 (1977), 441--472. \href{https://www.numdam.org/articles/10.24033/asens.1333/}{Text}.
\bibitem{blasiusRam} D. Blasius, \emph{Hilbert modular forms and the Ramanujan conjecture}, Aspects Math. E37 (2006), 35--56. \href{https://arxiv.org/abs/math/0511007}{Text}.
\bibitem{borelStable} A. Borel, \emph{Stable real cohomology of arithmetic groups}, Ann. Sci. École Norm. Sup. (4) 7 (1974), 235--272. \href{https://www.numdam.org/articles/10.24033/asens.1269/}{Text}.
\bibitem{borelSerre} A. Borel and J.-P. Serre, \emph{Corners and arithmetic groups}, Comment. Math. Helv. 48 (1973), 436--491. \href{https://doi.org/10.1007/BF02566134}{Text}.
\bibitem{calegariEven} F. Calegari, \emph{Even Galois representations and the Fontaine--Mazur conjecture}, Invent. Math. 185 (2011), 1--16. \href{https://doi.org/10.1007/s00222-010-0297-0}{Text}.
\bibitem{caraianiNearby} A. Caraiani, \emph{Local-global compatibility and the action of monodromy on nearby cycles}, Duke Math. J. 161 (2012), 2311--2413. \href{https://www.ma.imperial.ac.uk/~acaraian/papers/lgc1.pdf}{Text}.
\bibitem{chenZouMultiplicity} R. Chen and J. Zou, \emph{Arthur's multiplicity formula for even orthogonal and unitary groups}, J. Eur. Math. Soc. 27 (2025), 4769--4843. \href{https://doi.org/10.4171/JEMS/1564}{Text}.
\bibitem{deJongShokrieh} R. de Jong and F. Shokrieh, \emph{Faltings height and Néron--Tate height of a theta divisor}, Compos. Math. 158 (2022), 1--32. \href{https://doi.org/10.1112/S0010437X21007661}{Text}.
\bibitem{ehlsDoubling} E. Eischen, M. Harris, J.-S. Li and C. Skinner, \emph{$p$-adic $L$-functions for unitary groups}, Forum Math. Pi 8 (2020), e9. \href{https://arxiv.org/pdf/1602.01776}{Text}.
\bibitem{ferreroWashington} B. Ferrero and L. C. Washington, \emph{The Iwasawa invariant $\mu_p$ vanishes for abelian number fields}, Ann. of Math. (2) 109 (1979), 377--395. \href{https://annals.math.princeton.edu/1979/190-2/p05}{Journal record}.
\bibitem{fosterTheta} T. Foster, J. Rabinoff, F. Shokrieh and A. Soto, \emph{Non-Archimedean and tropical theta functions}, Math. Ann. 372 (2018), 891--914. \href{https://arxiv.org/abs/1705.08212}{Text}.
\bibitem{ganQiuTakeda} W. T. Gan, Y. Qiu and S. Takeda, \emph{The regularized Siegel--Weil formula (the second term identity) and the Rallis inner product formula}, Invent. Math. 198 (2014), 739--831. \href{https://arxiv.org/pdf/1207.4709}{Text}.
\bibitem{ganTakedaHowe} W. T. Gan and S. Takeda, \emph{A proof of the Howe duality conjecture}, J. Amer. Math. Soc. 29 (2016), 473--493. \href{https://doi.org/10.1090/jams/839}{Text}.
\bibitem{gelanderVolume} T. Gelander, \emph{Homotopy type and volume of locally symmetric manifolds}, Duke Math. J. 124 (2004), 459--515. \href{https://arxiv.org/pdf/math/0111165}{Text}.
\bibitem{grobnerRaghuram} H. Grobner and A. Raghuram, \emph{On some arithmetic properties of automorphic forms of $\mathrm{GL}_m$ over a division algebra}, Int. J. Number Theory 10 (2014), 963--1013. \href{https://homepage.univie.ac.at/harald.grobner/papers/GLn(D).pdf}{Text}.
\bibitem{harrisSpherical} M. Harris, \emph{Cohomological automorphic forms on unitary groups, II: period relations and values of $L$-functions}, in \emph{Harmonic Analysis, Group Representations, Automorphic Forms and Invariant Theory}, Lecture Notes Ser. Inst. Math. Sci. Natl. Univ. Singap. 12 (2007), 89--150. \href{https://webusers.imj-prg.fr/~michael.harris/periodrelations.pdf}{Text}.
\bibitem{kahnTwo} B. Kahn, \emph{The Quillen--Lichtenbaum conjecture at the prime 2}, preprint, 1997. \href{https://webusers.imj-prg.fr/~bruno.kahn/preprints/QL2.pdf}{Text}.
\bibitem{katoZeta} K. Kato, \emph{$p$-adic Hodge theory and values of zeta functions of modular forms}, Astérisque 295 (2004), 117--290. \href{https://www.numdam.org/item/AST_2004__295__117_0/}{Text}.
\bibitem{kawamataVanishing} Y. Kawamata, \emph{A generalization of Kodaira--Ramanujam's vanishing theorem}, Math. Ann. 261 (1982), 43--46. \href{https://doi.org/10.1007/BF01456407}{Text}.
\bibitem{lanToroidal} K.-W. Lan, \emph{Toroidal compactifications of PEL-type Kuga families}, Algebra Number Theory 6 (2012), 885--966. \href{https://msp.org/ant/2012/6-5/ant-v6-n5-p02-p.pdf}{Text}.
\bibitem{lanStroh} K.-W. Lan and B. Stroh, \emph{Relative cohomology of cuspidal forms on PEL-type Shimura varieties}, Algebra Number Theory 8 (2014), 1787--1799. \href{https://www.kwlan.org/articles/rel-cusp-PEL.pdf}{Text}.
\bibitem{liuTorsion} T. Liu, \emph{Torsion $p$-adic Galois representations and a conjecture of Fontaine}, Ann. Sci. École Norm. Sup. (4) 40 (2007), 633--674. \href{https://www.numdam.org/item/10.1016/j.ansens.2007.05.002.pdf}{Text}.
\bibitem{liuLattices} T. Liu, \emph{A note on lattices in semi-stable representations}, Math. Ann. 346 (2010), 117--138. \href{https://www.math.purdue.edu/~tongliu/pub/latticesnew2.pdf}{Text}.
\bibitem{liuFiltration} T. Liu, \emph{Filtration associated to torsion semi-stable representations}, Ann. Inst. Fourier 65 (2015), 1999--2035. \href{https://www.math.purdue.edu/~tongliu/pub/filtrationnew.pdf}{Text}.
\bibitem{liuThetaInert} Y. Liu, \emph{Theta correspondence for almost unramified representations of unitary groups}, J. Number Theory 230 (2022), 196--224. \href{https://doi.org/10.1016/j.jnt.2021.03.027}{Text}.
\bibitem{mathewHurewicz} A. Mathew, \emph{Torsion exponents in stable homotopy and the Hurewicz homomorphism}, Algebr. Geom. Topol. 16 (2016), 1025--1041. \href{https://arxiv.org/pdf/1501.07561}{Text}.
\bibitem{milneADT} J. S. Milne, \emph{Arithmetic Duality Theorems}, second edition, BookSurge, 2006. \href{https://www.jmilne.org/math/Books/adt.html}{Text}.
\bibitem{minguezHowe} A. Mínguez, \emph{Correspondance de Howe explicite: paires duales de type II}, Ann. Sci. Éc. Norm. Supér. (4) 41 (2008), 717--741. \href{https://doi.org/10.24033/asens.2080}{Text}.
\bibitem{miyauchiYamauchi} M. Miyauchi and T. Yamauchi, \emph{A remark on conductor, depth and principal congruence subgroups}, J. Algebra 592 (2022), 424--434. \href{https://arxiv.org/pdf/2107.08130}{Text}.
\bibitem{mumfordEquations} D. Mumford, \emph{On the equations defining abelian varieties. I}, Invent. Math. 1 (1966), 287--354. \href{https://pazuki.perso.math.cnrs.fr/index_fichiers/Mumford66.pdf}{Text}.
\bibitem{newtonThorne} J. Newton and J. A. Thorne, \emph{Adjoint Selmer groups of automorphic Galois representations of unitary type}, J. Eur. Math. Soc. 25 (2023), 1919--1967. \href{https://ems.press/content/serial-article-files/32858}{Text}.
\bibitem{panFM} L. Pan, \emph{The Fontaine--Mazur conjecture in the residually reducible case}, J. Amer. Math. Soc. 35 (2022), 1031--1169. \href{https://doi.org/10.1090/jams/991}{Text}.
\bibitem{paulTrapaHowe} A. Paul and P. E. Trapa, \emph{One-dimensional representations of $U(p,q)$ and the Howe correspondence}, J. Funct. Anal. 195 (2002), 129--166. \href{https://www.math.utah.edu/~ptrapa/char6_4.pdf}{Text}.
\bibitem{prasadVolume} G. Prasad, \emph{Volumes of $S$-arithmetic quotients of semi-simple groups}, Publ. Math. IHÉS 69 (1989), 91--114. \href{https://www.numdam.org/articles/10.1007/BF02698841/}{Text}.
\bibitem{schlessingerArtin} M. Schlessinger, \emph{Functors of Artin rings}, Trans. Amer. Math. Soc. 130 (1968), 208--222. \href{https://math.uchicago.edu/~amathew/schlessingerdef.pdf}{Text}.
\bibitem{shalikaGeneric} J. A. Shalika, \emph{The multiplicity one theorem for $\mathrm{GL}_n$}, Ann. of Math. (2) 100 (1974), 171--193. \href{https://annals.math.princeton.edu/1974/100-2/p01}{Text}.
\bibitem{springerReduction} T. A. Springer, \emph{Reduction theory over global fields}, Proc. Indian Acad. Sci. Math. Sci. 104 (1994), 207--216. \href{https://doi.org/10.1007/BF02830884}{Text}.
\bibitem{howeTranscending} R. Howe, \emph{Transcending classical invariant theory}, J. Amer. Math. Soc. 2 (1989), 535--552. \href{https://www.math.utah.edu/~ptrapa/math-library/howe/howe-transcending-1989.pdf}{Text}.
\bibitem{tateDifferent} The Stacks Project Authors, \emph{Discriminants and differents}, \S49.12, Lemmas 49.12.4--6. \href{https://stacks.math.columbia.edu/tag/0BWB}{Text}.
\bibitem{tateThesis} J. Tate, \emph{Fourier analysis in number fields and Hecke's zeta-functions}, in J. W. S. Cassels and A. Fröhlich (eds.), \emph{Algebraic Number Theory}, Academic Press, 1967, 305--347.
\bibitem{thorneAdjoint} J. A. Thorne, \emph{On the vanishing of adjoint Bloch--Kato Selmer groups of irreducible automorphic Galois representations}, Pure Appl. Math. Q. 18 (2022), 2159--2202. \href{https://arxiv.org/pdf/2207.04925}{Text}.
\bibitem{vanderKallen} W. van der Kallen, \emph{Homology stability for linear groups}, Invent. Math. 60 (1980), 269--295. \href{https://webspace.science.uu.nl/~kalle101/homstab1980.pdf}{Text}.
\bibitem{voevodskyNorm} V. Voevodsky, \emph{On motivic cohomology with $\mathbb Z/l$-coefficients}, Ann. of Math. (2) 174 (2011), 401--438. \href{https://annals.math.princeton.edu/2011/174-1/p11}{Text}.
\bibitem{weibelKbook} C. A. Weibel, \emph{The $K$-book: An introduction to algebraic $K$-theory}, Graduate Studies in Mathematics 145, American Mathematical Society, 2013, Chapter~VI. \href{https://sites.math.rutgers.edu/~weibel/Kbook/Kbook.VI.pdf}{Chapter VI}.
\bibitem{wuThetaTower} C. Wu, \emph{Irreducibility of theta lifting for unitary groups}, J. Number Theory 133 (2013), 3296--3318. \href{https://arxiv.org/abs/1409.0770}{Text}.
\bibitem{yamanaUnitary} S. Yamana, \emph{The Siegel--Weil formula for unitary groups}, Pacific J. Math. 264 (2013), 235--256. \href{https://doi.org/10.2140/pjm.2013.264.235}{Text}.
\bibitem{senCharacters} G. Dospinescu, V. Pa\v{s}k\=unas and B. Schraen,
\emph{Infinitesimal characters in arithmetic families}, Selecta Math.
31 (2025), article 76, \S5.1.
\href{https://doi.org/10.1007/s00029-025-01045-6}{Hodge--Tate cocharacters}.
\bibitem{stickelberger1890} L. Stickelberger,
\emph{Ueber eine Verallgemeinerung der Kreistheilung}, Math. Ann.
37 (1890), 321--367. \href{https://doi.org/10.1007/BF01721360}{Original paper}.
For the integral annihilator in this notation see C. Greither,
\href{https://doi.org/10.1365/s13291-021-00233-5}{Fitting ideals in number theory and arithmetic}, Theorem 1.
\bibitem{arthurClozelBC} J. Arthur and L. Clozel, \emph{Simple Algebras,
Base Change, and the Advanced Theory of the Trace Formula}, Annals of
Mathematics Studies 120 (1989), Chapter III.
\href{https://www.math.toronto.edu/arthur/pdf/30.pdf}{Authors' text}.
\bibitem{gelbartJacquetAdj} S. Gelbart and H. Jacquet,
\emph{A relation between automorphic representations of \(\mathrm{GL}(2)\)
and \(\mathrm{GL}(3)\)}, Ann. Sci. \'{E}c. Norm. Sup\'er. 11 (1978),
471--542, Theorem 9.3.
\href{https://www.numdam.org/article/ASENS_1978_4_11_4_471_0.pdf}{Published text}.
\bibitem{candelori-key} L. Candelori, \emph{The algebraic functional equation of
Riemann's theta function}, Ann. Inst. Fourier 70 (2020), 809--830.
\href{https://doi.org/10.5802/aif.3324}{Published text}.
\bibitem{moret-bailly-key} L. Moret-Bailly, \emph{Sur l'équation fonctionnelle
de la fonction thêta de Riemann}, Compositio Math. 75 (1990), 203--217.
\href{https://www.numdam.org/item/CM_1990__75_2_203_0/}{Published text}.
\bibitem{moret-bailly-pinceaux} L. Moret-Bailly,
\emph{Pinceaux de variétés abéliennes}, Astérisque 129 (1985).
\href{https://www.numdam.org/item/AST_1985__129__1_0/}{Published text}.
\bibitem{deligne-reg-sing} P. Deligne, \emph{Équations différentielles à
points singuliers réguliers}, Lecture Notes in Mathematics 163 (1970),
II, Théorème 5.9. \href{https://publications.ias.edu/node/355}{Author's text}.
\bibitem{moeglinWaldspurger} C. M{\oe}glin and J.-L. Waldspurger,
\emph{Spectral Decomposition and Eisenstein Series}, Cambridge University
Press (1995), I.4.11.
\href{https://doi.org/10.1017/CBO9780511470905}{Publisher's edition}.
\bibitem{kashiwara-vergne} M. Kashiwara and M. Vergne,
\emph{On the Segal--Shale--Weil representations and harmonic polynomials},
Invent. Math. 44 (1978), 1--47, \S6.
\href{https://doi.org/10.1007/BF01389900}{Published text}.
\bibitem{harrisKudlaSweet} M. Harris, S. S. Kudla and W. J. Sweet,
\emph{Theta dichotomy for unitary groups}, J. Amer. Math. Soc. 9
(1996), 941--1004. \href{https://doi.org/10.1090/S0894-0347-96-00198-1}{Published text}.
\bibitem{voganZuckerman} D. A. Vogan, Jr. and G. J. Zuckerman,
\emph{Unitary representations with non-zero cohomology}, Compositio
Math. 53 (1984), 51--90, Theorem 5.3 and Proposition 6.1.
\href{https://www.numdam.org/item/CM_1984__53_1_51_0/}{Published text}.

\end{thebibliography}
\end{document}
